% MONograph-LEVEL REMEDIATED - Volume I Chapter 3 (2026-05-29)
% Literary compression pass: OSR backbone, landmark theorems, inheritance trim
% Priorities A--E + monograph critique (tone, scaffolding, §3.0 compression)
% Authority: TOC/CH03-DRAFTING-KICKOFF-LOCK.md
% TOC authority: TOC/volume-I-ch01-04-FINAL.txt (Ch.3 block locked 2026-05-29)
% Flow audit: TOC/CH03-TOC-AUDIT-FLOW-LOCK.md
% Drafted: SS3.1--SS3.11.3 COMPLETE (2026-05-29); Ch.3 closed; next: Ch.4 bulk
% Do not add Ch 4+ content to this file
\chapter[Angular Momentum Eigenfields and Representation Theory]{Angular Momentum Eigenfields and\\Representation Theory}
\label{ch:03}
\setcounter{section}{-1}


\section{Introduction}
\label{sec:ch30}

An antenna pattern plotted on a sphere is already a coordinate picture: the laboratory has chosen how to index directions on the observation shell. The displayed amplitudes are not self-interpreting labels. They are samples of a field state that survived Maxwell dynamics, boundary enforcement, reactive storage, evanescent sheets, and the outgoing selector \(\mathcal{R}_\omega=\Pi_{\mathrm{rad}}\circ\mathcal{S}_\omega\) established in Chapter~\ref{ch:02}
\cite{stratton1941,harrington2001}. At fixed \(\omega\), the present chapter asks a representation-theoretic question: \emph{in which coordinates does the surviving radiative structure become labeled, orthogonal, normalizable, and measurable?} Spherical harmonics, \(\mathrm{SO}(3)\) blocks, and Wigner matrices are classical tools; the distinctive viewpoint developed here is that radiation must be read through a strict three-layer hierarchy---\textbf{operator} \(\to\) \textbf{state} \(\to\) \textbf{representation}---in which operators determine export, states carry physical content, and representations merely label states \cite{allen1992,stratton1941,hansen1988}.

Impressed sources \(\Jimp\) on bounded support \(\Omega\) enter \(\mathcal{S}_\omega\) (Eq.~\eqref{eq:ch2.maxwell-solution-map}); outgoing projection \(\Pi_{\mathrm{rad}}\) (Eq.~\eqref{eq:ch2.rad-projection-composition}) retains the Sommerfeld-qualified radiative sector; and
\begin{equation}
\label{eq:ch3.export-pipeline-preview}
\mathcal{R}_\omega[\Jimp]
=
\Pi_{\mathrm{rad}}\big(\mathcal{S}_\omega[\Jimp]\big)
\end{equation}
restates Eq.~\eqref{eq:ch2.radiation-operator-def} \cite{stratton1941,harrington2001}. The object \(\mathcal{R}_\omega[\Jimp]\) is a field pair on \(\mathcal{F}_{\mathrm{rad}}\), not a coefficient list. Modal indices enter only after the representation map of Subsection~\ref{sec:ch301}. Volume~I reverses the order common in OAM discourse: one first asks what radiates, what survives export, and what becomes observable; only then does an angular label acquire meaning \cite{belinfante1940,harrington2001}. Principles~\ref{prin:ch2.radiative-accessibility} and~\ref{prin:ch1.rep-physics} encode that ordering: indexing applies to accessible export, and labels become physical only after operator and measurement screens are satisfied.

\begin{figure}[t]
\centering
\ChOneFigBlock{%
\begin{tikzpicture}[ch1fig, node distance=6mm and 9mm]
  \node[ch1box] (src) {Source\\$\Jimp$ on $\Omega$};
  \node[ch1box, right=11mm of src] (sw) {$\mathcal{S}_\omega$\\full Maxwell};
  \node[ch1box, right=11mm of sw] (pi) {$\Pi_{\mathrm{rad}}$\\outgoing};
  \node[ch1box, right=11mm of pi] (fr) {$\mathcal{F}_{\mathrm{rad}}$\\state};
  \node[ch1box, right=11mm of fr] (rep) {Representation\\coefficients};
  \node[ch1box, right=11mm of rep] (meas) {Measurement\\reports};
  \draw[ch1flow] (src) -- (sw) -- (pi) -- (fr) -- (rep) -- (meas);
\end{tikzpicture}%
}
\caption{Laboratory chain for Volume~I representation theory. Chapter~\ref{ch:02} owns
\(\mathcal{S}_\omega\), \(\Pi_{\mathrm{rad}}\), and \(\mathcal{R}_\omega\); Chapter~\ref{ch:03}
owns the passage from \(\mathcal{F}_{\mathrm{rad}}\) to symmetry-adapted coordinates; Chapter~\ref{ch:01}
and Section~\ref{sec:ch29} own the measurement composition at the right.}
\label{fig:ch3.laboratory-chain}
\end{figure}

\paragraph{Convention 3.0A (inherited export chain).}
Throughout Chapter~\ref{ch:03}, \(\mathcal{R}_\omega=\Pi_{\mathrm{rad}}\circ\mathcal{S}_\omega\)
(Eqs.~\eqref{eq:ch2.radiation-operator-def}, \eqref{eq:ch3.export-pipeline-preview}) is fixed.
Domain, range, and commutator structure are cited only when a new theorem requires them.

Figure~\ref{fig:ch3.laboratory-chain} organizes the operator \(\to\) state \(\to\) representation \(\to\) measurement progression. Vector spherical harmonics and plane-wave spectra are coordinates on \(\operatorname{ran}\mathcal{R}_\omega\), not substitutes for export \cite{hansen1988,jackson1999,stratton1941}. Section~\ref{sec:ch31} closes Hilbert-space apparatus; Sections~\ref{sec:ch32}--\ref{sec:ch34} install angular-momentum algebra and VSH eigenfields; Sections~\ref{sec:ch35}--\ref{sec:ch39} fix normalization, spectra, basis laws, and gauge discipline; Section~\ref{sec:ch310} verifies; Section~\ref{sec:ch3inherit} forwards exports to Chapter~\ref{ch:04} \cite{harrington2001,devaney2004}.

Consider a compact resonant antenna: reactive structure near the feed never enters \(\mathcal{F}_{\mathrm{rad}}\); the far-zone pattern may be sparse in one basis and dense in another while radiated power and reciprocity remain fixed on \(\operatorname{ran}\mathcal{R}_\omega\) \cite{stratton1941,jackson1999}. It is essential to distinguish operator export \(\mathcal{R}_\omega\), coordinate choice \(\mathcal{C}\), and measurement \(\mathcal{M}\) in Eqs.~\eqref{eq:ch3.representation-map} and~\eqref{eq:ch3.measurement-on-coefficients}; realizability in Chapter~\ref{ch:04} constrains achievable coefficient vectors, not the existence of angular channels \cite{devaney2004}.

Subsection~\ref{sec:ch301} develops the operator--state--representation separation; Subsections~\ref{sec:ch302}--\ref{sec:ch303} connect it to radiation integrals and admissibility certificates.


\subsection{Symmetry, operators, and representations}
\label{sec:ch301}
% TOC tag: [HOME]

Radiation patterns, coupling coefficients, and modal reports display numbers indexed by angular labels. Before any label is read as ``orbital order'' or ``channel number,'' one must decide what object is indexed. By Convention~3.0A, that object is \(\mathcal{R}_\omega[\Jimp]\in\mathcal{F}_{\mathrm{rad}}\) \cite{stratton1941,harrington2001}. The operator--state--representation separation governs the entire chapter: \(\mathcal{S}_\omega\) and \(\mathcal{R}_\omega\) fix export; \(\mathfrak{F}\in\mathcal{F}_{\mathrm{rad}}\) carries radiative content; \(\mathcal{C}\) and symmetry representations \(\rho(g)\) label that content in Fig.~\ref{fig:ch3.laboratory-chain} \cite{jackson1999,stratton1941,harrington2001}.

\paragraph{Assumptions and admissible field pairs.}
Fix \(\omega>0\) under the \(\exp(-\jj\omega t)\) convention of Chapter~\ref{ch:01}, linear isotropic homogeneous regions unless stated otherwise, and outgoing closure on exterior problems as in Section~\ref{sec:ch24} \cite{stratton1941}. Impressed currents \(\Jimp\) lie in the admissible source class of Subsection~\ref{sec:ch201}; field pairs \((\E,\H)\) are discussed on \(\mathcal{F}_{\mathrm{rad}}\) only after application of \(\Pi_{\mathrm{rad}}\) from Eq.~\eqref{eq:ch2.rad-projection-composition} \cite{harrington2001}. Symmetry groups \(G\) are taken to preserve Maxwell equations and the boundary class \(\mathcal{B}_\Gamma\); time-harmonic translation is not re-developed here. When \(G=\mathrm{SO}(3)\), rotations act on spatial coordinates and on vector components consistently with Section~\ref{sec:ch01.3} \cite{jackson1999}. These assumptions are the same regime in which Eqs.~\eqref{eq:ch2.pipeline-composition} and~\eqref{eq:ch2.radiation-operator-def} were proved; extending representation theory beyond them requires new commutator analyses, not relabeled modes \cite{stratton1941}.

\paragraph{Three levels: operator, state, coordinate.}
Write \(\mathfrak{F}\in\mathcal{F}_{\mathrm{rad}}\) for a radiative field state (an equivalence class of admissible pairs on the outgoing branch). The export map
\begin{equation}
\label{eq:ch3.state-from-source}
\mathfrak{F}
=
\mathcal{R}_\omega[\Jimp]
\end{equation}
is an operator-level statement: it maps source data to a field state without choosing angular coordinates \cite{stratton1941,harrington2001}. A \emph{representation} in Volume~I is a linear isomorphism
\begin{equation}
\label{eq:ch3.representation-map}
\mathcal{C}:\mathcal{F}_{\mathrm{rad}}\supset\mathcal{V}
\longrightarrow
\mathbb{C}^{I},
\qquad
\mathcal{C}(\mathfrak{F})=\{c_\alpha\}_{\alpha\in I},
\end{equation}
from a dense subspace \(\mathcal{V}\) to a coefficient space indexed by \(I\) (discrete, continuous, or mixed) \cite{stratton1941,jackson1999}. Equation~\eqref{eq:ch3.representation-map} does not alter \(\mathfrak{F}\); it re-expresses an unchanged state. Explicit bases \(\{\mathbf{f}_\alpha\}\) are constructed in Section~\ref{sec:ch34} \cite{hansen1988,tai1994,stratton1941}. One must distinguish Eq.~\eqref{eq:ch3.state-from-source} (``what radiates'') from Eq.~\eqref{eq:ch3.representation-map} (``how we list what radiates'').

\begin{figure}[t]
\centering
\ChOneFigBlock{%
\begin{tikzpicture}[ch1fig, node distance=7mm and 10mm]
  \node[ch1box] (op) {Operator level\\$\mathcal{R}_\omega[\Jimp]$};
  \node[ch1box, right=14mm of op] (sp) {State space\\$\mathcal{F}_{\mathrm{rad}}$};
  \node[ch1box, right=14mm of sp] (rep) {Representation\\$\mathcal{C}(\mathfrak{F})=\{c_\alpha\}$};
  \node[ch1box, below=11mm of rep] (so) {$\mathrm{SO}(3)$\\$\rho(g)$ blocks};
  \node[ch1box, below=11mm of op] (near) {Near-field\\not in $\mathcal{F}_{\mathrm{rad}}$};
  \draw[ch1flow] (op) -- (sp) -- (rep);
  \draw[ch1flow] (so) -- (rep);
  \draw[ch1flow, dashed] (near) -- (op);
\end{tikzpicture}%
}
\caption{Operator-to-representation bridge. \(\mathcal{R}_\omega\) selects the outgoing state;
\(\mathcal{C}\) chooses coordinates on \(\mathcal{F}_{\mathrm{rad}}\) in which \(\mathrm{SO}(3)\)
acts block-wise. Near-field admixtures removed by \(\Pi_{\mathrm{rad}}\) never enter the representation
carrier (dashed arrow).}
\label{fig:ch3.operator-representation-bridge}
\end{figure}

Figure~\ref{fig:ch3.operator-representation-bridge} fixes the logical order: \(\mathcal{C}\) and \(\rho(g)\) act on \(\mathcal{F}_{\mathrm{rad}}\) after export; measurement from Section~\ref{sec:ch01.5} follows normalization in Section~\ref{sec:ch35} \cite{harrington2001}. One must not apply Wigner rotation laws to near-field samples in the reactive shell: those samples lie outside \(\mathcal{F}_{\mathrm{rad}}\) and need not transform under the \(\rho(g)\) blocks of exported patterns \cite{stratton1941,harrington2001}.

\paragraph{Symmetry action on field states.}
Let \(G\) be a symmetry group acting on \(\mathbb{R}^{3}\) by smooth maps \(g:\mathbf{r}\mapsto g\mathbf{r}\). The induced action on phasor Maxwell states is
\begin{equation}
\label{eq:ch3.symmetry-action-fields}
\big(T_g\mathfrak{F}\big)(\mathbf{r})
=
\big(\mathbf{U}_E(g)\,\E(g^{-1}\mathbf{r}),\,
\mathbf{U}_H(g)\,\H(g^{-1}\mathbf{r})\big),
\end{equation}
where \(\mathfrak{F}=(\E,\H)\), \(\mathbf{U}_E(g)\) and \(\mathbf{U}_H(g)\) are fixed \(3\times 3\) matrices for each \(g\), and the inverse argument \(g^{-1}\mathbf{r}\) implements the standard left action on functions \cite{jackson1999,stratton1941}. For pure rotations \(g\in\mathrm{SO}(3)\), \(\mathbf{U}_E=\mathbf{U}_H=\mathbf{R}(g)\) from Section~\ref{sec:ch01.3}; Subsection~\ref{sec:ch321} upgrades Eq.~\eqref{eq:ch3.symmetry-action-fields} to angular-momentum generator form \cite{stratton1941,jackson1999}. The map \(T_g\) is a group action because \(T_{g_1g_2}=T_{g_1}\circ T_{g_2}\) when the component matrices satisfy \(\mathbf{U}(g_1g_2)=\mathbf{U}(g_1)\mathbf{U}(g_2)\) \cite{jackson1999}. On energy-admissible subspaces equipped with the inner product of Subsection~\ref{sec:ch312}, \(T_g\) is unitary under the Volume~I phasor convention \cite{stratton1941,harrington2001}.

Physically, Eq.~\eqref{eq:ch3.symmetry-action-fields} answers the question: ``If I rigidly rotate my laboratory frame, how do \(\E\) and \(\H\) re-express on the same spatial points?'' The inverse on \(\mathbf{r}\) ensures that scalar observables measured at fixed world points are unchanged under passive rotation of coordinates \cite{jackson1999}. The nontrivial content for radiation is not this definition alone but the \emph{commutation} of \(T_g\) with \(\mathcal{R}_\omega\) and the \emph{simultaneous} diagonalization of commuting symmetry operators on \(\mathcal{F}_{\mathrm{rad}}\) \cite{stratton1941,harrington2001}.

\paragraph{Representations as homomorphisms.}
A \emph{representation} of \(G\) on a linear space \(\mathcal{V}\) is a homomorphism
\begin{equation}
\label{eq:ch3.representation-homomorphism}
\rho:G\longrightarrow \mathrm{GL}(\mathcal{V}),
\qquad
\rho(g_1g_2)=\rho(g_1)\rho(g_2),
\qquad
\rho(e)=\mathsf{I},
\end{equation}
where \(e\) is the group identity and \(\mathsf{I}\) the identity on \(\mathcal{V}\) \cite{jackson1999,stratton1941}. For radiative Maxwell states, \(\mathcal{V}\) is a dense subspace of \(\mathcal{F}_{\mathrm{rad}}\) stable under \(T_g\) in Eq.~\eqref{eq:ch3.symmetry-action-fields}. In coordinates, \(\rho(g)\) is the matrix that multiplies coefficient vectors when the laboratory is rotated:
\begin{equation}
\label{eq:ch3.coefficient-rotation-law}
\mathcal{C}(T_g\mathfrak{F})
=
\rho(g)\,\mathcal{C}(\mathfrak{F}),
\end{equation}
whenever \(\mathcal{C}\) is chosen consistently with the symmetry (spherical eigenbasis in Section~\ref{sec:ch34}, plane-wave basis in Section~\ref{sec:ch38}) \cite{jackson1999,hansen1988,stratton1941}. Equation~\eqref{eq:ch3.coefficient-rotation-law} is the finite-dimensional form of Eq.~\eqref{eq:ch3.symmetry-action-fields}; Wigner \(D\)-matrices in Subsection~\ref{sec:ch324} realize \(\rho(g)\) for irreducible \(\mathrm{SO}(3)\) sectors \cite{jackson1999}. The representation viewpoint linearizes how exported patterns transform under laboratory rotations without altering \(\mathcal{R}_\omega\) \cite{stratton1941}.

\paragraph{Pushforward on impressed currents.}
Rotating the source before export is implemented on \(\Jimp\) by the pushforward
\begin{equation}
\label{eq:ch3.current-pushforward}
\big(\mathcal{J}_g\Jimp\big)(\mathbf{r})
=
\mathbf{U}_J(g)\,\Jimp(g^{-1}\mathbf{r}),
\end{equation}
with \(\mathbf{U}_J(g)=\mathbf{R}(g)\) for spatial rotations of vector currents in isotropic regions \cite{jackson1999,stratton1941}. The map \(\mathcal{J}_g\) is the source-side analogue of \(T_g\): it records how impressed excitation re-labels when the hardware is rigidly rotated. Consistency with Eq.~\eqref{eq:ch3.symmetry-action-fields} is not automatic for every boundary--source pair; it holds when \(G\) preserves \(\mathcal{B}_\Gamma\), outgoing selectors, and the admissible class of \(\Jimp\) \cite{stratton1941,harrington2001}.

\begin{theorem}[Commutation of spatial symmetry with radiation export]
\label{thm:ch3.symmetry-commutes}
Let \(g\in G\) preserve the boundary class \(\mathcal{B}_\Gamma\) and the outgoing selector
\(\Pi_{\mathrm{rad}}\) used in Eq.~\eqref{eq:ch2.rad-projection-composition}. Then for all admissible
\(\Jimp\),
\begin{equation}
\label{eq:ch3.symmetry-commutes-with-radiation}
T_g\,\mathcal{R}_\omega[\Jimp]
=
\mathcal{R}_\omega\big[\mathcal{J}_g\Jimp\big].
\end{equation}
\end{theorem}

\begin{proof}[Proof sketch]
By linearity of \(\mathcal{S}_\omega\) and covariance of Maxwell equations under \(g\) (Chapter~\ref{ch:01}),
\(T_g\,\mathcal{S}_\omega[\Jimp]=\mathcal{S}_\omega[\mathcal{J}_g\Jimp]\) on admissible domains
\cite{stratton1941,jackson1999}. If \(\Pi_{\mathrm{rad}}\) commutes with \(T_g\) on the outgoing sector,
apply \(\Pi_{\mathrm{rad}}\) to both sides and use Eq.~\eqref{eq:ch2.radiation-operator-def}.
Full boundary hypotheses match Section~\ref{sec:ch24} and Subsection~\ref{sec:ch243}.
% PROOF: AUTHOR REVIEW
\end{proof}

Theorem~\ref{thm:ch3.symmetry-commutes} is a landmark of Volume~I: it unites Chapter~\ref{ch:01} covariance, Chapter~\ref{ch:02} export, and Chapter~\ref{ch:03} representation in one commutator,
\(T_g\,\mathcal{R}_\omega[\Jimp]=\mathcal{R}_\omega[\mathcal{J}_g\Jimp]\).
Rotating the laboratory and rotating the source yield equivalent labels on \(\mathcal{F}_{\mathrm{rad}}\) once \(\mathcal{J}_g\) is specified \cite{stratton1941,harrington2001}. Failure of commutation signals broken symmetry and is treated in Section~\ref{sec:ch36}. The theorem is the gateway from commutator criteria (Section~\ref{sec:ch01.3}) to the modal rotation law Eq.~\eqref{eq:ch3.coefficient-rotation-law} and the commutative diagram of Fig.~\ref{fig:ch3.report-chain-commutative-diagram}.

\paragraph{Modal coordinates as representations.}
When spectral decomposition applies on \(\mathcal{F}_{\mathrm{rad}}\), exported fields admit expansions
\begin{equation}
\label{eq:ch3.modal-coordinate-preview}
\mathcal{R}_\omega[\Jimp](\mathbf{r})
=
\sum_{\alpha\in I}\,c_\alpha(\Jimp)\,\mathbf{f}_\alpha(\mathbf{r}),
\end{equation}
where \(\{\mathbf{f}_\alpha\}_{\alpha\in I}\) is a symmetry-adapted basis and each \(\alpha\) labels a joint eigenvalue of commuting operators (angular momentum, helicity in transverse subspaces, and outgoing radial quantum numbers once Section~\ref{sec:ch33} is in place) \cite{stratton1941,jackson1999}. The coefficient functionals
\begin{equation}
\label{eq:ch3.coefficient-functional}
c_\alpha(\Jimp)
=
\langle \mathbf{f}_\alpha,\,\mathcal{R}_\omega[\Jimp]\rangle_{\mathcal{F}}
\end{equation}
will be made precise by the inner product \(\langle\cdot,\cdot\rangle_{\mathcal{F}}\) of Subsection~\ref{sec:ch312} and the flux normalizations of Section~\ref{sec:ch35} \cite{stratton1941,harrington2001}. Equation~\eqref{eq:ch3.modal-coordinate-preview} is the coordinate restatement of Eq.~\eqref{eq:ch3.state-from-source}: the operator map is source-to-field; the representation map is field-to-coefficients via Eq.~\eqref{eq:ch3.representation-map}. Subsection~\ref{sec:ch302} imports radiation integrals and abstract spectral decomposition from Sections~\ref{sec:ch251} and~\ref{sec:ch261} to justify expansions without constructing \(\mathbf{f}_\alpha\) explicitly; Subsection~\ref{sec:ch341} defines \(\mathbf{f}_\alpha\) as vector spherical harmonics \(\Mlm\), \(\Nlm\) \cite{hansen1988,tai1994,stratton1941}. Green operators determining \(c_\alpha(\Jimp)\) are fixed in Chapter~\ref{ch:02}; only the basis choice is new in Chapter~\ref{ch:03}.

\paragraph{Linearity and superposition in coefficient space.}
Because \(\mathcal{R}_\omega\) is linear (Eq.~\eqref{eq:ch2.radiation-linearity}), representation coefficients inherit linearity:
\begin{equation}
\label{eq:ch3.representation-linearity}
c_\alpha\big(a{\Jimp}_1+b{\Jimp}_2\big)
=
a\,c_\alpha({\Jimp}_1)+b\,c_\alpha({\Jimp}_2),
\end{equation}
for all \(\alpha\in I\) and admissible \(a,b\in\mathbb{C}\) \cite{stratton1941,harrington2001}. Physically, superposition of sources induces superposition of modal amplitudes \emph{in a fixed basis}. Changing bases mixes coefficients through transformation matrices of Section~\ref{sec:ch38}; mixing is not the same as coupling induced by broken symmetry in Section~\ref{sec:ch36} \cite{jackson1999,stratton1941}. It is essential to keep Eq.~\eqref{eq:ch3.representation-linearity} separate from reciprocity relations of Subsection~\ref{sec:ch282}, which pair sources and observers rather than scaling a single \(\Jimp\) \cite{stratton1941}.

\paragraph{Reducible versus irreducible content.}
Maxwell fields on \(\mathbb{R}^{3}\) carry reducible representations of \(\mathrm{SO}(3)\): multiple angular momenta couple in superpositions, and near environments admix sectors even when \(\mathcal{F}_{\mathrm{rad}}\) has been isolated \cite{jackson1999,stratton1941}. The representation-theoretic task is not to collapse that reducibility prematurely but to identify irreducible subspaces on which flux, helicity, and reciprocity functionals become diagonal or sparse \cite{belinfante1940,stratton1941}. Helicity and polarization sectors previewed in Subsection~\ref{sec:ch325} inherit gauge discipline from Section~\ref{sec:ch01.4}; they are not independent observables until the gauge-invariant screening of Section~\ref{sec:ch39} is applied \cite{belinfante1940}. Accessible angular budgets from Eq.~\eqref{eq:ch2.accessible-angular-budget} may further restrict which irreducible sectors carry measurable power even when those sectors appear in Eq.~\eqref{eq:ch3.modal-coordinate-preview} formally \cite{stratton1941,harrington2001}.

\begin{figure}[t]
\centering
\ChOneFigBlock{%
\begin{tikzpicture}[ch1fig, node distance=5mm and 8mm]
  \node[ch1box, minimum width=30mm] (A) {Operator\\$\mathcal{R}_\omega$};
  \node[ch1box, minimum width=30mm, right=10mm of A] (B) {State\\$\mathfrak{F}\in\mathcal{F}_{\mathrm{rad}}$};
  \node[ch1box, minimum width=30mm, right=10mm of B] (C) {Coordinates\\$\{c_\alpha\}$};
  \node[ch1box, minimum width=30mm, below=9mm of C] (D) {Measurement\\reports};
  \node[font=\footnotesize, below=2mm of A] {Ch.~2};
  \node[font=\footnotesize, below=2mm of B] {shared};
  \node[font=\footnotesize, below=2mm of C] {Ch.~3};
  \node[font=\footnotesize, below=2mm of D] {Ch.~1, \S2.10};
  \draw[ch1flow] (A) -- (B) -- (C) -- (D);
\end{tikzpicture}%
}
\caption{Three confusions to avoid: treating \(\mathcal{R}_\omega\) as a modal expander (left),
treating coefficients as physical states (right), or identifying near-field samples with
\(\mathcal{F}_{\mathrm{rad}}\) elements (state level). Each block has distinct transformation laws.}
\label{fig:ch3.three-level-distinction}
\end{figure}

\begin{principle}[Operator--State--Representation Separation Principle]
\label{prin:ch3.osr-separation}
Volume~I radiation theory rests on a three-layer hierarchy that is not negotiable.
\textbf{Operators}---\(\mathcal{S}_\omega\), \(\Pi_{\mathrm{rad}}\), \(\mathcal{R}_\omega\)---determine
what field content is exported from admissible sources.
\textbf{States} \(\mathfrak{F}\in\mathcal{F}_{\mathrm{rad}}\) carry the physical content that survives export.
\textbf{Representations} \(\mathcal{C}\) and measurement functionals \(\mathcal{M}\) label and report that
content after normalization and instrument rank limits.
None of these layers may be substituted for another: operators fix export; states fix physics;
representations fix indexing; measurements fix observability.
\end{principle}

Figure~\ref{fig:ch3.three-level-distinction} summarizes the distinction enforced throughout Volume~I. The operator \(\mathcal{R}_\omega\) has a kernel and a range determined by Maxwell physics and outgoing closure; a representation \(\mathcal{C}\) has no kernel on \(\mathcal{V}\) beyond coordinate gauge choices fixed in Section~\ref{sec:ch39}; a measurement functional may annihilate additional sectors through finite aperture rank and noise \cite{harrington2001,stratton1941}. Confusing these nullities produces false claims about ``unmeasurable OAM'' or ``hidden modes'' when the underlying issue is projection order in Eq.~\eqref{eq:ch2.operational-measurement-chain} \cite{harrington2001}.

\paragraph{Worked illustration: rotated electric dipole.}
Consider a compact electric dipole \(\Jimp\) in free space with outgoing closure. In the far zone, \(\mathcal{R}_\omega[\Jimp]\) approaches the classic \(\sin\theta\) pattern imported from Subsection~\ref{sec:ch01.2.4} \cite{stratton1941,jackson1999}. Rotate the source by \(g\in\mathrm{SO}(3)\). Operator-level invariance of radiated power follows from Eq.~\eqref{eq:ch2.radiated-power-screen} together with Theorem~\ref{thm:ch3.symmetry-commutes} \cite{stratton1941}. Representation-level prediction is different in kind: coefficients \(c_{\ell m}\) in a spherical basis transform as \(\rho(g)\) acting on the fixed-\(\ell\) block,
\begin{equation}
\label{eq:ch3.dipole-coefficient-rotation}
c'_{\ell m}
=
\sum_{m'=-\ell}^{\ell}
D^{(\ell)}_{m m'}(g)\,c_{\ell m'},
\end{equation}
with Wigner matrices \(D^{(\ell)}_{m m'}(g)\) from Subsection~\ref{sec:ch324} \cite{jackson1999}. For a \(z\)-aligned dipole, only \(\ell=1\) sectors participate; rotation mixes \(m\in\{-1,0,1\}\) without creating \(\ell=2\) content \cite{stratton1941,jackson1999}. Power invariance is operator-level; the \(D\)-matrix law is representation-level---Principle~\ref{prin:ch3.osr-separation} forbids conflating them \cite{jackson1999,stratton1941}.

\paragraph{Near-field sampling versus radiative coordinates.}
A probe scanning the reactive shell around the same dipole measures fields not equal to \(\mathcal{R}_\omega[\Jimp]\) restricted to interior points \cite{stratton1941,harrington2001}. Those samples contain \(1/r^2\) and \(1/r^3\) admixtures removed by \(\Pi_{\mathrm{rad}}\) from \(\mathcal{F}_{\mathrm{rad}}\) (Fig.~\ref{fig:ch3.operator-representation-bridge}, dashed arrow). Applying spherical harmonic labels to raw near-field scans---without outgoing projection and flux normalization---mixes reactive and radiative content and breaks the rotation law Eq.~\eqref{eq:ch3.coefficient-rotation-law} \cite{harrington2001,stratton1941}. The practical implication for array diagnostics is sharp: modal filters designed on \(S^{2}\) at fixed \(r\) apply to exported patterns; modal filters applied in the feed region require separate near-to-far calibration through the pipeline Eq.~\eqref{eq:ch2.pipeline-composition} \cite{harrington2001}.

\paragraph{Plane-wave versus spherical listings.}
A fixed \(\mathfrak{F}\in\mathcal{F}_{\mathrm{rad}}\) admits both a plane-wave spectrum \(\tilde{\E}(\mathbf{k})\) (Section~\ref{sec:ch29}) and a spherical listing \(\{c_{\ell m}\}\) (Eq.~\eqref{eq:ch3.modal-coordinate-preview}) \cite{jackson1999,stratton1941}. Coefficient sparsity is basis-dependent; transformation kernels appear in Section~\ref{sec:ch38}. OAM representation dependence and gauge screens are developed in Sections~\ref{sec:ch38}--\ref{sec:ch39} \cite{allen1992,belinfante1940}.

\paragraph{Measurement composition after coordinates.}
Let \(\mathcal{M}\) denote an operational measurement map from Section~\ref{sec:ch01.5}. The full laboratory report on modal content is
\begin{equation}
\label{eq:ch3.measurement-on-coefficients}
\text{report}
=
\mathcal{M}\circ\mathcal{N}\circ\mathcal{C}\big(\mathcal{R}_\omega[\Jimp]\big),
\end{equation}
where \(\mathcal{N}\) is flux normalization from Section~\ref{sec:ch35} and \(\mathcal{C}\) is the representation map Eq.~\eqref{eq:ch3.representation-map} \cite{harrington2001,stratton1941}. Equation~\eqref{eq:ch3.measurement-on-coefficients} restates Eq.~\eqref{eq:ch2.operational-measurement-chain} with the Chapter~\ref{ch:03} coordinate layer made explicit. Finite aperture rank, noise, and calibration error enter through \(\mathcal{M}\); gauge and basis choices enter through \(\mathcal{C}\) and \(\mathcal{N}\); physics enters through \(\mathcal{R}_\omega\) alone \cite{harrington2001}. Inverse problems that reconstruct \(\Jimp\) from reports must invert the \emph{composition}, not a single Wigner matrix \cite{devaney2004,harrington2001}.

\begin{figure}[t]
\centering
\ChOneFigBlock{%
\begin{tikzpicture}[ch1fig, node distance=8mm and 11mm]
  \node[ch1box] (J) {$\mathcal{J}_\Omega$};
  \node[ch1box, right=12mm of J] (F) {$\mathcal{F}_{\mathrm{rad}}$};
  \node[ch1box, right=12mm of F] (C) {$\mathbb{C}^{I}$};
  \node[ch1box, right=12mm of C] (M) {Reports};
  \node[ch1box, below=10mm of J] (Jg) {$\mathcal{J}_\Omega$};
  \node[ch1box, below=10mm of F] (Fg) {$\mathcal{F}_{\mathrm{rad}}$};
  \node[ch1box, below=10mm of C] (Cg) {$\mathbb{C}^{I}$};
  \node[ch1box, below=10mm of M] (Mg) {Reports};
  \draw[ch1flow] (J) -- node[above,font=\scriptsize] {$\mathcal{R}_\omega$} (F);
  \draw[ch1flow] (F) -- node[above,font=\scriptsize] {$\mathcal{N}\circ\mathcal{C}$} (C);
  \draw[ch1flow] (C) -- node[above,font=\scriptsize] {$\mathcal{M}$} (M);
  \draw[ch1flow] (Jg) -- node[above,font=\scriptsize] {$\mathcal{R}_\omega$} (Fg);
  \draw[ch1flow] (Fg) -- node[above,font=\scriptsize] {$\mathcal{N}\circ\mathcal{C}$} (Cg);
  \draw[ch1flow] (Cg) -- node[above,font=\scriptsize] {$\mathcal{M}$} (Mg);
  \draw[ch1flow] (J) -- node[left,font=\scriptsize] {$\mathcal{J}_g$} (Jg);
  \draw[ch1flow] (F) -- node[left,font=\scriptsize] {$T_g$} (Fg);
  \draw[ch1flow] (C) -- node[left,font=\scriptsize] {$\rho(g)$} (Cg);
\end{tikzpicture}%
}
\caption{Commutative diagram for the laboratory report chain
(Eq.~\eqref{eq:ch3.measurement-on-coefficients}): rotation of sources and rotation of
coordinates commute with export and measurement when
Theorem~\ref{thm:ch3.symmetry-commutes} holds.}
\label{fig:ch3.report-chain-commutative-diagram}
\end{figure}

\paragraph{Finite apertures and accessible sectors.}
Even when \(\mathfrak{F}=\mathcal{R}_\omega[\Jimp]\) is exact in \(\mathcal{F}_{\mathrm{rad}}\), a finite receiver aperture implements a rank-limited functional on \(\mathcal{C}(\mathfrak{F})\) \cite{harrington2001,stratton1941}. Accessible angular budget Eq.~\eqref{eq:ch2.accessible-angular-budget} from Section~\ref{sec:ch29} is the operator-theoretic packaging of that rank limit: certain \(\ell\) sectors may be present in Eq.~\eqref{eq:ch3.modal-coordinate-preview} yet lie outside the support of what the laboratory can discriminate \cite{stratton1941,harrington2001}. Representation theory clarifies the distinction. Missing a high-\(\ell\) sidelobe in a compact range measurement is often a \emph{measurement} nullity (right block in Fig.~\ref{fig:ch3.three-level-distinction}); missing it because \(\Pi_{\mathrm{rad}}\) never exported that content is an \emph{operator} nullity (left block) \cite{harrington2001}. Conflating the two leads to over-interpretation of sparse modal fits and to under-designed calibration campaigns \cite{devaney2004}.

\paragraph{Reciprocity as a representation constraint.}
Reciprocity from Subsection~\ref{sec:ch282} pairs source and observation roles without changing \(\mathcal{R}_\omega\) \cite{stratton1941}. In coefficient space, reciprocity becomes a bilinear form on products \(c_\alpha(\Jimp)\,c_\beta^{\mathrm{rec}}({\Jimp}_{\mathrm{rec}})\) once normalization from Section~\ref{sec:ch35} is fixed \cite{stratton1941,harrington2001}. Reciprocity is a \emph{constraint on allowable coupling matrices} between representation indices \cite{stratton1941}. Later orthogonality proofs in Subsection~\ref{sec:ch343} exploit the same structure \cite{hansen1988,stratton1941}.

\paragraph{Validity boundaries.}
Eqs.~\eqref{eq:ch3.symmetry-action-fields}--\eqref{eq:ch3.measurement-on-coefficients} assume linear Maxwell theory, fixed \(\omega\), and outgoing \(\mathcal{F}_{\mathrm{rad}}\) \cite{stratton1941}. Nonlinear media, time-varying boundaries, and moving platforms break Theorem~\ref{thm:ch3.symmetry-commutes} in the simple form stated here; such regimes lie outside Volume~I scope \cite{stratton1941}. Discrete eigenbases require separation assumptions developed in Section~\ref{sec:ch33}; continuous-spectrum rigging is completed in Subsection~\ref{sec:ch314} \cite{stratton1941,harrington2001}. When symmetry is reduced to a subgroup (arrays, reflectors, guides), \(\rho(g)\) block-diagonalizes only partially; Section~\ref{sec:ch36} classifies the mixed sectors that remain \cite{harrington2001}.

Subsection~\ref{sec:ch301} fixed the representation viewpoint for Volume~I: symmetry groups act on \(\mathcal{F}_{\mathrm{rad}}\) through \(T_g\); representations \(\rho(g)\) linearize those actions on coefficients through Eq.~\eqref{eq:ch3.coefficient-rotation-law}; modal coordinates Eq.~\eqref{eq:ch3.modal-coordinate-preview} express \(\operatorname{ran}\mathcal{R}_\omega\) in symmetry-adapted bases constructed in Section~\ref{sec:ch34}. Subsection~\ref{sec:ch302} connects this frame to radiation integrals and spectral operators from Chapter~\ref{ch:02}, preparing the eigenfunction expansions that Section~\ref{sec:ch31} will place on Hilbert-space footing \cite{stratton1941,hansen1988}.

\subsection{Eigenfield expansions of radiative fields}
\label{sec:ch302}
% TOC tag: [USE SS2.5.1][USE SS2.6.1]

Eigenfield expansions return indexed coefficients only when integration, export, and representation
are composed in that order (Principle~\ref{prin:ch3.osr-separation}). Radiation integrals
(Subsection~\ref{sec:ch251}) and spectral decomposition (Subsection~\ref{sec:ch261}) justify
expansions of \(\mathcal{R}_\omega[\Jimp]\) in outgoing eigenmodes \(\mathbf{f}_\alpha\)
\cite{stratton1941,harrington2001}. Named harmonics \(\Mlm\), \(\Nlm\) appear in
Subsection~\ref{sec:ch341} \cite{hansen1988,stratton1941}.

\paragraph{Scope and standing imports.}
Work at fixed \(\omega>0\) with compact impressed support \(\Omega_s\), admissible data on
\(\mathcal{J}_\Omega\) from Subsection~\ref{sec:ch201}, and outgoing closure
Eqs.~\eqref{eq:ch2.sommerfeld-outgoing-class}--\eqref{eq:ch2.outgoing-selector-def} on exterior
problems \cite{stratton1941}. Field pairs are evaluated through Theorem~\ref{thm:ch2.source-to-field-integral}
of Subsection~\ref{sec:ch251}; radiative images are formed by
Eq.~\eqref{eq:ch2.integral-radiation-map}, equivalent to Eq.~\eqref{eq:ch3.export-pipeline-preview}
\cite{harrington2001}. Spectral objects---eigenpairs Eq.~\eqref{eq:ch2.radiation-eigenpair-def},
projectors Eq.~\eqref{eq:ch2.modal-projector-def}, and expansions
Eqs.~\eqref{eq:ch2.radiation-spectral-expansion}--\eqref{eq:ch2.radiation-map-modal-sum}---are
imported from Subsection~\ref{sec:ch261} as operator-level facts on \(\mathcal{F}_{\mathrm{rad}}\)
\cite{stratton1941,harrington2001}. Regime qualification from Section~\ref{sec:ch01.2.4} and
Subsection~\ref{sec:ch252} governs when far-zone amplitudes may be identified with
\(\mathcal{R}_\omega[\Jimp]\) samples \cite{stratton1941,jackson1999}.

\paragraph{Integral evaluation followed by radiation selection.}
Let \(\mathfrak{D}\) denote admissible impressed data with electric part \({\Jimp}\) on
\(\Omega_s\). Subsection~\ref{sec:ch251} realizes \(\mathcal{S}_\omega[\mathfrak{D}]\) as outgoing
volume (or equivalent surface) integrals; symbolically,
\begin{equation}
\label{eq:ch3.solution-from-integral}
\big(\E[\mathfrak{D}],\H[\mathfrak{D}]\big)
=
\mathcal{I}_\omega[\mathfrak{D}],
\end{equation}
where \(\mathcal{I}_\omega\) stands for any route in
Fig.~\ref{fig:ch2.source-to-field-integral-routes} that satisfies Theorem~\ref{thm:ch2.source-to-field-integral}
\cite{stratton1941,harrington2001}. The map \(\mathcal{I}_\omega\) is \emph{not} the radiation
export: it returns the full admissible Maxwell pair on the observation domain, including reactive
and evanescent admixtures removed only after \(\Pi_{\mathrm{rad}}\) \cite{stratton1941}. Radiative
export is the composition
\begin{equation}
\label{eq:ch3.integral-then-projection}
\mathcal{R}_\omega[\mathfrak{D}]
=
\Pi_{\mathrm{rad}}\big(\mathcal{I}_\omega[\mathfrak{D}]\big),
\end{equation}
which restates Eq.~\eqref{eq:ch2.integral-radiation-map} in the notation of
Eq.~\eqref{eq:ch3.state-from-source} \cite{harrington2001}. Algebraically,
Eq.~\eqref{eq:ch3.integral-then-projection} says that integration and projection do not commute
with choosing angular coordinates: one first propagates sources through outgoing kernels, then
applies the outgoing selector, and only then expands the surviving state
\cite{stratton1941,harrington2001}.

Physically, Eq.~\eqref{eq:ch3.integral-then-projection} is the Volume~I guard against reading
near-field probe data as modal coefficients. A scan in the reactive shell measures
\(\mathcal{I}_\omega[\mathfrak{D}]\) restricted to interior radii; a pattern cut on a far sphere
measures the transverse amplitude of \(\mathcal{R}_\omega[\mathfrak{D}]\) once
Theorem~\ref{thm:ch2.far-field-approximation} applies \cite{stratton1941,jackson1999}. Eigenfield
labels attach to the second object, not the first \cite{harrington2001}.

\begin{figure}[t]
\centering
\ChOneFigBlock{%
\begin{tikzpicture}[ch1fig, node distance=6mm and 8mm]
  \node[ch1box] (j) {${\Jimp}$ on $\Omega_s$};
  \node[ch1box, right=10mm of j] (int) {$\mathcal{I}_\omega$\\Subsec.~\ref{sec:ch251}};
  \node[ch1box, right=10mm of int] (pi) {$\Pi_{\mathrm{rad}}$};
  \node[ch1box, right=10mm of pi] (rad) {$\mathcal{R}_\omega[\mathfrak{D}]$};
  \node[ch1box, right=10mm of rad] (spec) {Spectral sum\\Subsec.~\ref{sec:ch261}};
  \node[ch1box, right=10mm of spec] (coef) {$\{c_\alpha\}$};
  \node[ch1box, below=10mm of int] (near) {near-field\\content};
  \draw[ch1flow] (j) -- (int) -- (pi) -- (rad) -- (spec) -- (coef);
  \draw[ch1flow, dashed] (near) -- (int);
\end{tikzpicture}%
}
\caption{Eigenfield expansion chain for Subsection~\ref{sec:ch302}: outgoing integrals
(Subsection~\ref{sec:ch251}) produce full admissible fields; \(\Pi_{\mathrm{rad}}\) selects
\(\mathcal{F}_{\mathrm{rad}}\); spectral decomposition (Subsection~\ref{sec:ch261}) yields
coefficients on outgoing eigenmodes. Near-field content (dashed) is excluded from the expansion
carrier.}
\label{fig:ch3.integral-spectral-eigenfield-chain}
\end{figure}

Figure~\ref{fig:ch3.integral-spectral-eigenfield-chain} is the workflow against which every
expansion coefficient in Chapter~\ref{ch:03} will later be tested. The figure deliberately omits
kernel formulas: those belong to Subsection~\ref{sec:ch251} \cite{stratton1941}. What Subsection~\ref{sec:ch302}
adds is the \emph{identification} of spectral indices \(\alpha\) with outgoing eigenfunctions on
\(\mathcal{F}_{\mathrm{rad}}\) once separation and Hilbert-space closure are in place
(Sections~\ref{sec:ch31}--\ref{sec:ch34}) \cite{harrington2001,hansen1988}.

\paragraph{Far-zone amplitudes as radiation functionals.}
On observation shells in the far zone, Theorem~\ref{thm:ch2.far-field-approximation} identifies the
transverse amplitude \(\E_0(\theta,\phi)\) from
Eqs.~\eqref{eq:ch2.far-field-integral-reduction}--\eqref{eq:ch2.far-field-amplitude-integral}
as the angular radiation functional induced by \(\mathcal{I}_\omega[\mathfrak{D}]\)
\cite{stratton1941,jackson1999}. Define the far-zone radiation functional
\begin{equation}
\label{eq:ch3.far-amplitude-functional}
\mathcal{A}_\infty[\mathfrak{D}](\hat{\mathbf{r}})
=
\E_0(\theta,\phi)\big|_{\hat{\mathbf{r}}},
\qquad
\hat{\mathbf{r}}=(\sin\theta\cos\phi,\,\sin\theta\sin\phi,\,\cos\theta),
\end{equation}
with \(\E_0\) obtained from the locked far-field reduction of Subsection~\ref{sec:ch252} and
\(\hat{\mathbf{r}}\cdot\E_0=0\) \cite{stratton1941}. Equation~\eqref{eq:ch3.far-amplitude-functional}
is a \emph{functional} of source data on the unit sphere \(S^{2}\): it encodes direction-resolved
transverse content after the \(e^{-\jj k_0 r}/r\) envelope is factored \cite{jackson1999}. When
regime gates from Eq.~\eqref{eq:ch2.distance-regime-definitions} are satisfied,
\(\mathcal{A}_\infty[\mathfrak{D}]\) samples the same angular information carried by
\(\mathcal{R}_\omega[\mathfrak{D}]\) on large spheres \cite{stratton1941,harrington2001}; when
they fail, \(\mathcal{A}_\infty\) must not be interpreted as a modal fingerprint of export
\cite{harrington2001}.

Mathematically, \(\mathcal{A}_\infty\) is linear in \(\mathfrak{D}\) because
Eq.~\eqref{eq:ch2.far-field-amplitude-integral} is linear in \({\Jimp}\)
\cite{stratton1941}. Physically, \(\mathcal{A}_\infty\) is what compact-range and far-field
antenna ranges report after calibration: a complex vector on \(S^{2}\) per polarization channel,
not yet expanded in \(\ell,m\) indices \cite{harrington2001,jackson1999}. The eigenfield expansion layer treats
\(\mathcal{A}_\infty\) as the bridge between Subsection~\ref{sec:ch251} integrals and
Subsection~\ref{sec:ch261} spectral sums \cite{stratton1941}.

\paragraph{Eigenfield expansion on \(\mathcal{F}_{\mathrm{rad}}\).}
Let \(\{\mathbf{f}_\alpha\}_{\alpha\in I}\) be a family of outgoing eigenfields on
\(\mathcal{F}_{\mathrm{rad}}\) that diagonalize the radiation operator in the sense of
Theorem~\ref{thm:ch2.radiation-spectral-decomposition}, with \(I\) discrete, continuous, or
mixed as declared in Subsection~\ref{sec:ch261} \cite{stratton1941,harrington2001}. Each
\(\mathbf{f}_\alpha=(\E_\alpha,\H_\alpha)\) carries angular quantum numbers fixed by commuting
symmetry operators (developed in Section~\ref{sec:ch32}) and outgoing radial admissibility
(developed in Section~\ref{sec:ch33}) \cite{jackson1999,stratton1941}. The eigenfield expansion
of an exported state \(\mathfrak{F}=\mathcal{R}_\omega[\mathfrak{D}]\) is
\begin{equation}
\label{eq:ch3.eigenfield-expansion-on-frad}
\mathfrak{F}
=
\sum_{\alpha\in I} c_\alpha[\mathfrak{D}]\,\mathbf{f}_\alpha,
\end{equation}
with coefficients determined by the spectral pairing inherited from
Eq.~\eqref{eq:ch2.source-modal-amplitude} once \(\mathbf{f}_\alpha\) replaces the abstract
\(\mathbf{u}_n\) \cite{stratton1941,harrington2001}. Equation~\eqref{eq:ch3.eigenfield-expansion-on-frad}
specializes preview Eq.~\eqref{eq:ch3.modal-coordinate-preview} and operator restatement
Eq.~\eqref{eq:ch2.radiation-map-modal-sum}: the left side is the radiative \emph{state}; the right
side is its listing in symmetry-adapted eigenfields \cite{hansen1988,stratton1941}. Convergence is
understood in the energy norm Eq.~\eqref{eq:ch2.field-energy-norm} on the physical branch until
rigged completion in Subsection~\ref{sec:ch314} \cite{stratton1941,harrington2001}.

The index map between Chapter~\ref{ch:02} modal labels \(n\in\mathcal{N}_{\mathrm{rad}}\) and
Chapter~\ref{ch:03} labels \(\alpha\in I\) is, at this stage, a \emph{bookkeeping alignment}:
\begin{equation}
\label{eq:ch3.spectral-index-alignment}
\alpha\longleftrightarrow n,
\qquad
\mathbf{f}_\alpha\longleftrightarrow \mathbf{u}_n,
\qquad
c_\alpha[\mathfrak{D}]\longleftrightarrow \lambda_n\,c_n[\mathfrak{D}],
\end{equation}
whenever \(\mathbf{u}_n\) is normalized as in Eq.~\eqref{eq:ch2.radiation-eigenpair-def} and
\(\mathbf{f}_\alpha\) carries the same outgoing angular--radial quantum numbers
\cite{stratton1941}. Equation~\eqref{eq:ch3.spectral-index-alignment} prevents duplicate narratives:
Subsection~\ref{sec:ch261} owns eigenvalues \(\lambda_n\) and projectors \(P_n\); Subsection~\ref{sec:ch302}
owns the eigenfield basis interpretation of those modes on \(\mathcal{F}_{\mathrm{rad}}\)
\cite{harrington2001}. Explicit construction of \(\mathbf{f}_\alpha\) as \(\Mlm\), \(\Nlm\) is
deferred to Section~\ref{sec:ch34} \cite{hansen1988,tai1994,stratton1941}.

\paragraph{Source coupling to eigenfield coefficients.}
Coefficients in Eq.~\eqref{eq:ch3.eigenfield-expansion-on-frad} are not free parameters; they are
functionals of \(\mathfrak{D}\) obtained by pairing \(\mathcal{I}_\omega[\mathfrak{D}]\) with the
adjoint radiation modes. In the normalized discrete case,
\begin{equation}
\label{eq:ch3.eigenfield-coefficient-pairing}
c_\alpha[\mathfrak{D}]
=
\big\langle \mathcal{R}_\omega[\mathfrak{D}],\,\mathbf{f}_\alpha\big\rangle_{\mathcal{U}},
\qquad
\alpha\in I,
\end{equation}
with \(\langle\cdot,\cdot\rangle_{\mathcal{U}}\) the energy inner product from
Eq.~\eqref{eq:ch2.field-energy-norm} \cite{stratton1941,harrington2001}. Equation
\eqref{eq:ch3.eigenfield-coefficient-pairing} is the eigenfield restatement of
Eq.~\eqref{eq:ch3.coefficient-functional} and agrees with
Eq.~\eqref{eq:ch2.source-modal-amplitude} when \(\mathbf{f}_\alpha=\mathbf{u}_n\) and
\(\lambda_n\) is absorbed into normalization conventions of Section~\ref{sec:ch35}
\cite{stratton1941}. The integral route enters through
Eq.~\eqref{eq:ch3.integral-then-projection}: one may evaluate \(c_\alpha\) only after
\(\mathcal{R}_\omega[\mathfrak{D}]\) is formed, even though \(\mathcal{I}_\omega[\mathfrak{D}]\)
is what Subsection~\ref{sec:ch251} computes explicitly \cite{harrington2001}.

For far-zone diagnostics, substituting Theorem~\ref{thm:ch2.far-field-approximation} gives the
angular coupling form
\begin{equation}
\label{eq:ch3.far-pattern-eigenfield-coupling}
\mathcal{A}_\infty[\mathfrak{D}](\hat{\mathbf{r}})
=
\sum_{\alpha\in I}
c_\alpha[\mathfrak{D}]\,\mathbf{e}_\alpha(\hat{\mathbf{r}}),
\end{equation}
where \(\mathbf{e}_\alpha\) is the transverse far-zone angular profile of \(\mathbf{f}_\alpha\)
on \(S^{2}\) \cite{stratton1941,jackson1999}. Equation~\eqref{eq:ch3.far-pattern-eigenfield-coupling}
aligns Eq.~\eqref{eq:ch3.far-amplitude-functional} with preview
Eq.~\eqref{eq:ch2.far-pattern-modal-preview} without re-displaying the radiation integral kernel
\cite{harrington2001}. Each term \(c_\alpha\,\mathbf{e}_\alpha\) is a \emph{mode-shaped} contribution
to the pattern; superposition is linear in sources through
Eq.~\eqref{eq:ch2.radiation-linearity} \cite{stratton1941}.

\begin{theorem}[Consistency of integral, spectral, and eigenfield expansions]
\label{thm:ch3.expansion-consistency}
Let outgoing closure, admissible \(\mathfrak{D}\), and regime qualification be as above. Suppose
\(\{\mathbf{f}_\alpha\}_{\alpha\in I}\) is a complete orthonormal radiation eigenfamily in
\(\mathcal{F}_{\mathrm{rad}}\) in the sense of Theorem~\ref{thm:ch2.radiation-spectral-decomposition}.
Then:
\begin{enumerate}[label=(\roman*)]
\item \(\mathcal{R}_\omega[\mathfrak{D}]\) satisfies
Eq.~\eqref{eq:ch3.eigenfield-expansion-on-frad} with coefficients
Eq.~\eqref{eq:ch3.eigenfield-coefficient-pairing};
\item the far-zone functional Eq.~\eqref{eq:ch3.far-amplitude-functional} satisfies
Eq.~\eqref{eq:ch3.far-pattern-eigenfield-coupling} on shells where
Theorem~\ref{thm:ch2.far-field-approximation} holds;
\item the expansion coincides with Eq.~\eqref{eq:ch2.radiation-map-modal-sum} under the alignment
Eq.~\eqref{eq:ch3.spectral-index-alignment}.
\end{enumerate}
\end{theorem}
\begin{proof}[Proof sketch]
Items (i)--(iii) follow by applying Eq.~\eqref{eq:ch2.radiation-spectral-expansion} to
\(\mathbf{v}=\mathcal{R}_\omega[\mathfrak{D}]\), identifying \(\mathbf{u}_n\) with
\(\mathbf{f}_\alpha\), and using Eqs.~\eqref{eq:ch2.source-modal-amplitude}--\eqref{eq:ch2.radiation-eigenpair-def}.
Far-zone equality (ii) uses Theorem~\ref{thm:ch2.far-field-approximation} to replace
\(\mathcal{R}_\omega[\mathfrak{D}]\) on large spheres by \(\mathcal{A}_\infty[\mathfrak{D}]\)
\cite{stratton1941,jackson1999}. Normalization conventions are fixed in Section~\ref{sec:ch35}.
% PROOF: AUTHOR REVIEW
\end{proof}

Theorem~\ref{thm:ch3.expansion-consistency} is the second landmark of the chapter: it unifies
antenna integral theory, abstract spectral theory, and eigenfield listing as three coordinate
descriptions of one exported object \(\mathcal{R}_\omega[\mathfrak{D}]\) \cite{stratton1941,harrington2001,jackson1999}.
Integral routes, spectral projectors, and harmonic coefficients are not competing physical models;
they are synchronized reports on \(\operatorname{ran}\mathcal{R}_\omega\) once outgoing closure holds.

\paragraph{Continuous spectra and distributional modes.}
On exterior domains, \(I\) often contains a continuous angular parameter (wave direction, spectral
abscissa) in addition to any discrete sectors \cite{jackson1999,stratton1941}. Write
\(I=I_{\mathrm{d}}\cup I_{\mathrm{c}}\) with discrete \(I_{\mathrm{d}}\) and continuous
\(I_{\mathrm{c}}\). The eigenfield expansion becomes
\begin{equation}
\label{eq:ch3.continuous-eigenfield-expansion}
\mathfrak{F}
=
\sum_{\alpha\in I_{\mathrm{d}}} c_\alpha[\mathfrak{D}]\,\mathbf{f}_\alpha
\;+\;
\int_{I_{\mathrm{c}}} c(\beta)[\mathfrak{D}]\,\mathbf{f}(\beta)\,d\mu(\beta),
\end{equation}
with a positive measure \(d\mu\) on \(I_{\mathrm{c}}\) fixed by normalization in
Section~\ref{sec:ch35} and rigged Hilbert completion in Subsection~\ref{sec:ch314}
\cite{stratton1941,harrington2001}. Equation~\eqref{eq:ch3.continuous-eigenfield-expansion} is the
eigenfield restatement of continuous-spectrum resolution previewed in Subsection~\ref{sec:ch261};
plane-wave angular spectra from Section~\ref{sec:ch29} realize one choice of
\(I_{\mathrm{c}}\) and \(d\mu\) \cite{jackson1999,stratton1941}. The continuous parameter must not be read as
``extra radiation'' in Eq.~\eqref{eq:ch3.continuous-eigenfield-expansion}: it indexes
\(\mathfrak{F}\) with continuous labels when discrete completeness fails on open domains
\cite{harrington2001}.

\paragraph{Truncation and finite-mode remainders.}
Laboratory expansions truncate \(I\) at finite cardinality. Let \(I^{(N)}\subset I\) be a
truncated index set and define the remainder
\begin{equation}
\label{eq:ch3.eigenfield-truncation-remainder}
\mathfrak{R}_N[\mathfrak{D}]
=
\mathcal{R}_\omega[\mathfrak{D}]
-
\sum_{\alpha\in I^{(N)}} c_\alpha[\mathfrak{D}]\,\mathbf{f}_\alpha.
\end{equation}
The energy norm of \(\mathfrak{R}_N\) is controlled by the compression tail in
Eq.~\eqref{eq:ch2.compression-error-bound} when \(\{\mathbf{f}_\alpha\}\) is the spectral family of
Subsection~\ref{sec:ch261} \cite{harrington2001}. Equation~\eqref{eq:ch3.eigenfield-truncation-remainder}
separates \emph{representation truncation} from \emph{operator nullity}: modes in
\(\ker\mathcal{R}_\omega\) never appear in any truncated sum because \(c_\alpha\equiv 0\) on export,
while modes outside \(I^{(N)}\) may carry power yet lie in the remainder until \(N\) is increased
\cite{stratton1941,harrington2001}. Accessible angular budget
Eq.~\eqref{eq:ch2.accessible-angular-budget} may further limit which truncated coefficients are
observable even when \(\|\mathfrak{R}_N\|\) is small \cite{harrington2001}.

\paragraph{Worked illustration: compact source far pattern before \(\ell\) labels.}
Consider a compact electric source with \(k_0 a\ll 1\) in free space. Subsection~\ref{sec:ch251}
evaluates \(\mathcal{I}_\omega[\mathfrak{D}]\) through outgoing volume integrals;
Eq.~\eqref{eq:ch3.integral-then-projection} yields \(\mathcal{R}_\omega[\mathfrak{D}]\) dominated
by dipolar far content \cite{stratton1941,jackson1999}. Theorem~\ref{thm:ch2.far-field-approximation}
gives \(\mathcal{A}_\infty[\mathfrak{D}]\) from
Eq.~\eqref{eq:ch2.far-field-amplitude-integral} with a low-order angular dependence on
\(S^{2}\) \cite{jackson1999}. Theorem~\ref{thm:ch3.expansion-consistency} gives the eigenfield expansion
Eq.~\eqref{eq:ch3.eigenfield-expansion-on-frad} and far-zone profile
Eq.~\eqref{eq:ch3.far-pattern-eigenfield-coupling}; only \(\ell=1\) sectors contribute at leading
order---a statement of coefficient sparsity, not a separate mechanism \cite{stratton1941,jackson1999}. Explicit \(\ell,m\) harmonics appear
only after Subsection~\ref{sec:ch341} names \(\mathbf{f}_\alpha\) \cite{hansen1988,stratton1941}.

\paragraph{Surface-equivalent routes and single export image.}
Subsection~\ref{sec:ch251} also realizes \(\mathcal{I}_\omega[\mathfrak{D}]\) through Huygens
equivalent-surface integrals Eqs.~\eqref{eq:ch2.huygens-equivalent-integral}--\eqref{eq:ch2.huygens-magnetic-integral}
when tangential traces on \(\Gamma_{\mathrm{eq}}\) are matched \cite{stratton1941,harrington2001}.
Those surface forms are not a second radiation operator: Theorem~\ref{thm:ch2.source-to-field-integral}
identifies every admissible route with the same \(\mathcal{S}_\omega\) image on the outgoing branch
\cite{stratton1941}. Consequently,
\begin{equation}
\label{eq:ch3.route-independence-of-expansion}
c_\alpha[\mathfrak{D}]
\ \text{from volume}\ \mathcal{I}_\omega
=
c_\alpha[\mathfrak{D}]
\ \text{from surface}\ \mathcal{I}_\omega^{\mathrm{(Huygens)}},
\end{equation}
whenever both routes satisfy the same trace class and outgoing selector
Eq.~\eqref{eq:ch2.outgoing-selector-def} \cite{harrington2001}. Equation
\eqref{eq:ch3.route-independence-of-expansion} is an \emph{expansion} statement, not a claim that
volume and surface integrals are numerically identical before projection: reactive near shells differ
until \(\Pi_{\mathrm{rad}}\) is applied \cite{stratton1941}. Array designers who replace a volume
current model by an aperture equivalent on \(\Gamma_{\mathrm{eq}}\) therefore obtain the same
\(\{c_\alpha\}\) only after the export composition Eq.~\eqref{eq:ch3.integral-then-projection} is
honored on both paths \cite{harrington2001}.

\paragraph{Regime gates before coefficient identification.}
Eigenfield coefficients extracted from \(\mathcal{A}_\infty\) inherit every regime restriction
carried by far-field reduction. Let \(r\) be the observation radius and \(a\) the source scale.
When \(k_0 r\) lies in the near band of Eq.~\eqref{eq:ch2.distance-regime-definitions},
\(\mathcal{I}_\omega[\mathfrak{D}]\) contains comparable reactive and radiative magnitudes and
Eq.~\eqref{eq:ch3.far-pattern-eigenfield-coupling} is not yet meaningful on that shell
\cite{stratton1941,jackson1999}. When \(k_0 r\) is intermediate, Fresnel parameter
Eq.~\eqref{eq:ch2.fresnel-parameter} controls whether \(\mathcal{A}_\infty\) has stabilized to a
Fraunhofer pattern \cite{jackson1999,harrington2001}. Only when \(k_0 r\) exceeds the far thresholds
and, for extended apertures, \(r\gtrsim 2a^{2}/\lambda_0\), does Theorem~\ref{thm:ch2.far-field-approximation}
license identifying pattern coefficients with eigenfield couplings
Eq.~\eqref{eq:ch3.eigenfield-coefficient-pairing} \cite{stratton1941}. Power-flux convergence
Eq.~\eqref{eq:ch2.power-flux-regime-convergence} supplies the complementary energy diagnostic:
\(P(r)\) must plateau before \(c_\alpha\) inferred from flux matches \(c_\alpha\) inferred from
angular amplitude \cite{stratton1941,harrington2001}. These gates are operator facts from
Subsection~\ref{sec:ch252}; Subsection~\ref{sec:ch302} invokes them to prevent premature modal labeling
on reactive shells \cite{harrington2001}.

\paragraph{Plane-wave angular content versus eigenfield sectors.}
Preview Eq.~\eqref{eq:ch2.far-pattern-modal-preview} and coupling
Eq.~\eqref{eq:ch3.far-pattern-eigenfield-coupling} describe the same transverse object
\(\mathcal{A}_\infty[\mathfrak{D}]\) in two indexing languages \cite{jackson1999,stratton1941}.
Plane-wave filters from Section~\ref{sec:ch29} expand \(\mathcal{A}_\infty\) in directions
\(\hat{\mathbf{k}}\) on the unit sphere; eigenfield expansions expand it in sectors \(\alpha\)
diagonalizing \(\mathcal{R}_\omega\) \cite{harrington2001}. The transforms are linear and
invertible on the accessible subspace Eq.~\eqref{eq:ch2.accessible-modal-sector}, but sparse in
one basis need not be sparse in the other \cite{jackson1999}. Subsection~\ref{sec:ch302} therefore
does not treat \(\alpha\) as a ``more physical'' label than \(\hat{\mathbf{k}}\): both are
coordinates on \(\operatorname{ran}\mathcal{R}_\omega\) after export \cite{stratton1941}. Section~\ref{sec:ch38}
will supply the transformation kernels explicitly; here the sole claim is that
Eq.~\eqref{eq:ch3.eigenfield-expansion-on-frad} is the spectral-coordinate form paired with the
plane-wave listing of Section~\ref{sec:ch29} \cite{harrington2001,jackson1999}.

\paragraph{Superposition and cross-mode coupling before symmetry breaking.}
When several admissible sources \(\mathfrak{D}_1,\mathfrak{D}_2\) are active,
Eq.~\eqref{eq:ch2.radiation-linearity} implies
\begin{equation}
\label{eq:ch3.eigenfield-superposition-law}
c_\alpha\big[\mathfrak{D}_1+\mathfrak{D}_2\big]
=
c_\alpha[\mathfrak{D}_1]+c_\alpha[\mathfrak{D}_2],
\end{equation}
for each \(\alpha\in I\) \cite{stratton1941,harrington2001}. Equation
\eqref{eq:ch3.eigenfield-superposition-law} is the eigenfield specialization of
Eq.~\eqref{eq:ch3.representation-linearity} and holds in any fixed eigenbasis. It does \emph{not}
imply that individual sources excite orthogonal \(\alpha\) sectors: near environments and finite
structures admix modes before export, and broken \(\mathrm{SO}(3)\) symmetry mixes irreducible
blocks in Section~\ref{sec:ch36} \cite{harrington2001,stratton1941}. The present subsection
records expansion \emph{consistency}; coupling topology is deferred to Sections~\ref{sec:ch36}
and~\ref{sec:ch37} \cite{jackson1999}.

Named harmonics, Hilbert closure, and completeness proofs follow in
Sections~\ref{sec:ch31}, \ref{sec:ch341}, and~\ref{sec:ch36}. Subsection~\ref{sec:ch303}
certifies admissibility; Section~\ref{sec:ch31} makes
Eqs.~\eqref{eq:ch3.eigenfield-coefficient-pairing}--\eqref{eq:ch3.continuous-eigenfield-expansion}
well posed \cite{stratton1941,harrington2001}.


\subsection{Continuity with Chapters 1 and 2}
\label{sec:ch303}
% TOC tag: narrative bridge (imports Ch.1 + Ch.2 by reference; no [USE] tag)

Subsections~\ref{sec:ch301}--\ref{sec:ch302} fixed the operator--state--representation frame and
proved Theorems~\ref{thm:ch3.symmetry-commutes} and~\ref{thm:ch3.expansion-consistency}. The present
subsection states the admissibility law under which Sections~\ref{sec:ch31}--\ref{sec:ch34} attach
Hilbert closure, angular-momentum algebra, and harmonic bases to \(\mathcal{F}_{\mathrm{rad}}\)
\cite{stratton1941,harrington2001}.

\paragraph{Composed report chain.}
Every representation claim factors through
\begin{equation}
\label{eq:ch3.continuity-composition-extended}
\mathfrak{A}_{\mathrm{ch3}}^{\mathrm{(report)}}
=
\mathfrak{A}_{\mathrm{meas}}
\circ
\mathfrak{A}_{\mathrm{g}}
\circ
\mathfrak{A}_{\mathrm{cov}}
\circ
\mathfrak{A}_{\mathrm{tr}}
\circ
\mathcal{N}
\circ
\mathcal{C}
\circ
\Pi_{\mathrm{rad}}
\circ
\mathcal{S}_\omega,
\end{equation}
where \(\mathcal{C}\) is the representation map Eq.~\eqref{eq:ch3.representation-map},
\(\mathcal{N}\) is flux normalization from Section~\ref{sec:ch35}, and the four
\(\mathfrak{A}\)-screens are the inherited layers from
Eq.~\eqref{eq:ch2.continuity-composition-law} \cite{stratton1941,harrington2001}. Equation
\eqref{eq:ch3.continuity-composition-extended} extends Eq.~\eqref{eq:ch3.measurement-on-coefficients}
to the full admissibility stack: \(\mathcal{C}\) and \(\mathcal{N}\) may relabel and rescale export,
but they may not repair failures of \(\mathcal{D}_{J}\), \(\mathcal{D}_{J,\mathrm{reg}}\), or
\(\mathfrak{A}_{\mathrm{g}}\) upstream \cite{stratton1941,belinfante1940}. A modal coefficient
\(c_\alpha\) is reportable only when \(\mathfrak{A}_{\mathrm{ch3}}^{\mathrm{(report)}}\) applied to
the originating \(\mathfrak{D}\) returns a direct or inference-licensed quantity in the sense of
Eq.~\eqref{eq:ch1.report-class-lock} \cite{stratton1941}.

\paragraph{Chapter~\ref{ch:02} inheritance gate.}
Chapter~\ref{ch:02} closed the source-to-field-to-export operator chain under
Eq.~\eqref{eq:ch2.inheritance-admissibility} \cite{stratton1941,harrington2001}. Compatibility of
the present construction layer with that closure is
\begin{equation}
\label{eq:ch3.ch2-inheritance-gate}
\mathcal{D}_{\mathrm{ch3}}=1
\iff
\left(
\mathcal{D}_{\mathrm{ch2}}=1\
\wedge\
\mathcal{D}_{\mathrm{bridge}}^{(2\to 3)}(\mathcal{Y})=1\
\wedge\
\mathcal{Y}\ \text{adds no recursive proof of}\ \mathcal{R}_{\mathrm{ch2}}
\right),
\end{equation}
with \(\mathcal{D}_{\mathrm{bridge}}^{(2\to 3)}\) from Eq.~\eqref{eq:ch2.bridge-inheritance-map}
and \(\mathcal{R}_{\mathrm{ch2}}\) the result ledger Table~\ref{tab:ch2.results-ledger}
\cite{stratton1941,harrington2001}. Harmonic relabeling of \(\operatorname{ran}\mathcal{R}_\omega\) is
permitted; altering \(\mathcal{R}_\omega\) or \(\Pi_{\mathrm{rad}}\) is not \cite{stratton1941}.

\paragraph{No-recursion screen for representation proofs.}
Let \(\mathcal{Y}\) denote any new construction proposed in the present construction layer. Admissibility
requires
\begin{equation}
\label{eq:ch3.bridge-nonrecursion}
\mathcal{D}_{\mathrm{bridge}}^{(\mathrm{ch3})}(\mathcal{Y})=1
\iff
\mathcal{Y}\prec\mathcal{R}_{\mathrm{ch2}}
\ \wedge\
\mathcal{Y}\ \text{adds no proof of}\ \mathcal{R}_{\mathrm{ch2}}
\ \wedge\
\mathcal{Y}\ \text{adds no proof of}\ \mathcal{R}_{\mathrm{ch1}},
\end{equation}
with \(\prec\) inherited from Eq.~\eqref{eq:ch1.bridge-compatibility-criteria}
\cite{stratton1941}. Permitted constructions include Hilbert completion, angular-momentum algebra,
Helmholtz separation, VSH definitions, normalization, and basis-transformation formulas
\cite{hansen1988,tai1994,stratton1941}.

\begin{figure}[t]
\centering
\ChOneFigBlock{%
\begin{tikzpicture}[ch1fig, node distance=6mm and 9mm]
  \node[ch1box, minimum width=32mm] (ch1) {Ch.~1\\certificates\\$\mathfrak{A}_{\mathrm{tr,g,cov,meas}}$};
  \node[ch1box, minimum width=32mm, right=11mm of ch1] (ch2) {Ch.~2\\$\mathcal{S}_\omega$, $\mathcal{R}_\omega$\\$\Pi_{\mathrm{rad}}$};
  \node[ch1box, minimum width=32mm, right=11mm of ch2] (ch3) {Ch.~3\\$\mathcal{C}$, $\mathcal{N}$\\eigenfields};
  \node[ch1box, minimum width=28mm, right=11mm of ch3] (ch4) {Ch.~4\\realizability\\(deferred)};
  \draw[ch1flow] (ch1) -- (ch2) -- (ch3) -- (ch4);
  \node[font=\footnotesize, below=2mm of ch1] {Level I};
  \node[font=\footnotesize, below=2mm of ch2] {operators};
  \node[font=\footnotesize, below=2mm of ch3] {coordinates};
\end{tikzpicture}%
}
\caption{Volume~I continuity stack for Subsection~\ref{sec:ch303}. Chapter~\ref{ch:01} owns
admissibility certificates; Chapter~\ref{ch:02} owns radiation operators; Chapter~\ref{ch:03}
owns representation coordinates on \(\mathcal{F}_{\mathrm{rad}}\); Chapter~\ref{ch:04} inherits
all three layers by reference.}
\label{fig:ch3.volume-continuity-stack}
\end{figure}

\paragraph{Import bundles.}
Locked operator inputs from Subsection~\ref{sec:ch02113} are packaged as
\begin{equation}
\label{eq:ch3.import-bundle-lock}
\mathcal{B}_{\mathrm{ch3}}^{(\mathrm{in})}
=
\mathcal{B}_{\mathrm{I}}\cup\mathcal{B}_{\mathrm{II}}\cup\mathcal{B}_{\mathrm{III}}
\cup\mathcal{B}_{\mathrm{IV}}\cup\mathcal{B}_{\mathrm{V}},
\qquad
\text{use}(\mathcal{B}_{\mathrm{ch3}}^{(\mathrm{in})})
=
\text{reference-only}.
\end{equation}
\cite{stratton1941,harrington2001}. A coefficient may be defined on \(\mathcal{F}_{\mathrm{rad}}\) yet
non-reportable once \(\mathfrak{A}_{\mathrm{meas}}\) is applied \cite{harrington2001}.

\paragraph{Representation admissibility criterion.}
Representation-specific admissibility requires
\begin{equation}
\label{eq:ch3.representation-admissibility}
\mathcal{D}_{\mathrm{rep}}=1
\iff
\left(
\mathfrak{F}\in\mathcal{F}_{\mathrm{rad}}\
\wedge\
\mathcal{C}(\mathfrak{F})\ \text{is flux-normalizable}\
\wedge\
\{\mathbf{f}_\alpha\}\ \text{respects outgoing closure}
\right),
\end{equation}
with \(\mathcal{F}_{\mathrm{rad}}\) from Subsection~\ref{sec:ch201} and \(\mathcal{C}\) from
Eq.~\eqref{eq:ch3.representation-map} \cite{stratton1941,harrington2001}. Equation
\eqref{eq:ch3.representation-admissibility} is the representation-layer specialization of
Eq.~\eqref{eq:ch3.ch2-inheritance-gate} to coordinate constructions: a harmonic expansion is
admissible only on outgoing-qualified states, not on raw \(\mathcal{S}_\omega\) images
\cite{stratton1941}. Normalization details completing the flux-normalizability clause are owned by
Section~\ref{sec:ch35}; rigged extensions are owned by Subsections~\ref{sec:ch313}--\ref{sec:ch314}
\cite{harrington2001}.

\begin{definition}[Admissible representation class]
\label{def:ch3.R-adm}
The admissible representation class is
\begin{equation}
\label{eq:ch3.R-adm-class}
\mathcal{R}_{\mathrm{adm}}
=
\left\{
\mathcal{C}:\mathcal{V}\to\mathbb{C}^{I}:
\mathcal{D}_{\mathrm{rep}}(\mathcal{C})=1
\right\}.
\end{equation}
A harmonic listing is Volume~I admissible only when \(\mathcal{C}\in\mathcal{R}_{\mathrm{adm}}\),
equivalently when Eqs.~\eqref{eq:ch3.representation-admissibility} and
\eqref{eq:ch3.ch2-inheritance-gate} hold jointly.
\end{definition}

\begin{table}[t]
\centering
\caption{Chapter~\ref{ch:03} proof schedule after Subsection~\ref{sec:ch303}: section ownership and
primary imports (reference-only anchors from Chapters~\ref{ch:01}--\ref{ch:02}).}
\label{tab:ch3.proof-schedule}
\small
\begin{tabular}{@{}p{0.20\linewidth}p{0.13\linewidth}p{0.40\linewidth}p{0.19\linewidth}@{}}
\toprule
Construction target & Section & Primary import anchor(s) & New proof owner \\
\midrule
Hilbert/rigged \(\mathcal{F}_{\mathrm{rad}}\) &
§3.1 &
Eqs.~\eqref{eq:ch2.field-energy-norm}, \eqref{eq:ch2.admissible-field-space};
Subsec.~\ref{sec:ch211} &
§3.1.1--§3.1.4 \\
Angular-momentum algebra, \(\mathrm{SO}(3)\) &
§3.2 &
Sec.~\ref{sec:ch01.3}; Subsec.~\ref{sec:ch262} &
§3.2.1--§3.2.5 \\
Vector Helmholtz separation &
§3.3 &
Eq.~\eqref{eq:ch2.helmholtz-forcing-map}; Thm.~\ref{thm:ch2.exterior-radiative-uniqueness} &
§3.3.1--§3.3.4 \\
\(Y_\ell^m\), \(\Mlm\), \(\Nlm\) definitions &
§3.4 &
Subsec.~\ref{sec:ch261}; Subsec.~\ref{sec:ch282} &
§3.4.1 only \\
Normalization, density of states &
§3.5 &
Subsec.~\ref{sec:ch253}; Eq.~\eqref{eq:ch2.radiation-spectral-expansion} &
§3.5.1--§3.5.4 \\
Completeness, symmetry breaking &
§3.6 &
Eq.~\eqref{eq:ch2.dof-partition-master}; Subsec.~\ref{sec:ch274} &
§3.6.1--§3.6.4 \\
Discrete/continuous spectra &
§3.7 &
Sec.~\ref{sec:ch29}; Eq.~\eqref{eq:ch2.accessible-angular-budget} &
§3.7.1--§3.7.4 \\
Basis change, OAM caveats &
§3.8--§3.9 &
Sec.~\ref{sec:ch01.4}; Eq.~\eqref{eq:ch2.plane-wave-radiation-operator} &
§3.8--§3.9 \\
\bottomrule
\end{tabular}
\end{table}

Table~\ref{tab:ch3.proof-schedule} orders the proof-bearing sections of the chapter
\cite{stratton1941,harrington2001}.

\begin{proposition}[Admissible representation extensions]
\label{prop:ch3.admissible-extensions}
Let \(\mathcal{Y}\) be a construction in Chapter~\ref{ch:03} targeting
\(\operatorname{ran}\mathcal{R}_\omega\). Then \(\mathcal{Y}\) is Volume~I admissible if and only
if Eqs.~\eqref{eq:ch3.ch2-inheritance-gate}, \eqref{eq:ch3.bridge-nonrecursion}, and
\eqref{eq:ch3.representation-admissibility} hold simultaneously. In particular, citing
Theorem~\ref{thm:ch2.radiation-spectral-decomposition} to justify
Eq.~\eqref{eq:ch3.eigenfield-expansion-on-frad} is admissible; re-proving
Theorem~\ref{thm:ch2.exterior-radiative-uniqueness} inside Section~\ref{sec:ch33} is not.
\end{proposition}
\begin{proof}[Proof sketch]
The forward direction restates the bridge locks
Eqs.~\eqref{eq:ch1.bridge-inheritance-map}--\eqref{eq:ch2.bridge-inheritance-map} for the
representation layer. The converse direction follows because any construction violating
no-recursion would either alter \(\mathcal{R}_{\mathrm{ch2}}\) or extend
\(\operatorname{dom}\mathcal{S}_\omega\) beyond Eq.~\eqref{eq:ch2.inheritance-operator-domain},
contradicting Table~\ref{tab:ch2.results-ledger} \cite{stratton1941,harrington2001}.
% PROOF: AUTHOR REVIEW
\end{proof}

\paragraph{Chapter~\ref{ch:04} boundary.}
Realizability on finite arrays, MoM optimality, and inverse-source uniqueness are deferred to
Chapter~\ref{ch:04} (Eq.~\eqref{eq:ch2.defer-topic-rule}) \cite{harrington2001,devaney2004}.
The harmonic coordinates supplied here are the domain on which attainability is tested
\cite{harrington2001}.

\paragraph{Inherited screens and nullities.}
Chapter~\ref{ch:01} transport, gauge, and measurement screens together with Chapter~\ref{ch:02}
source nullities and accessible angular budget remain active on every coefficient route through
Eq.~\eqref{eq:ch3.continuity-composition-extended} \cite{stratton1941,harrington2001,belinfante1940}.
Within the present construction layer, \(Y_\ell^m\), \(\Mlm\), and \(\Nlm\) are defined once in
Subsection~\ref{sec:ch341}; the total angular-momentum operator is defined once in
Subsection~\ref{sec:ch321}; Section~\ref{sec:ch3inherit} is the sole inheritance section
\cite{hansen1988,tai1994,stratton1941}. The master degree-of-freedom partition
Eq.~\eqref{eq:ch2.dof-partition-master} governs what content may appear inside
\(\mathcal{R}_\omega[\mathfrak{D}]\) before any harmonic label is attached; harmonic expansions
expose those limits as missing or unresolved coefficients without altering the partition itself
\cite{harrington2001,stratton1941}.


The representation space introduced in Chapter~\ref{ch:03} provides the coordinate
domain on which realizability constraints act. Chapter~\ref{ch:04} will therefore study
not the existence of radiative states, but the attainability of representation vectors
\(\{c_\alpha\}\in\mathbb{C}^{I}\) under finite geometry, spectral truncation, and
source accessibility limits imported from Chapter~\ref{ch:02}.

\paragraph{Terminal summary of Section~\ref{sec:ch30}.}
The operator--state--representation doctrine (Principle~\ref{prin:ch3.osr-separation}), Theorem~\ref{thm:ch3.symmetry-commutes}, and Theorem~\ref{thm:ch3.expansion-consistency} are the permanent exports of the introduction. Admissibility is certified by Eqs.~\eqref{eq:ch3.continuity-composition-extended}--\eqref{eq:ch3.representation-admissibility}, Definition~\ref{def:ch3.R-adm}, and Table~\ref{tab:ch3.proof-schedule}. Section~\ref{sec:ch31} opens Hilbert-space closure on \(\mathcal{F}_{\mathrm{rad}}\) \cite{stratton1941,harrington2001}.


\section{Hilbert-Space Framework for Radiative Fields}
\label{sec:ch31}

Coefficient pairings such as Eq.~\eqref{eq:ch3.eigenfield-coefficient-pairing} presuppose more than a
symbol \(\mathcal{F}_{\mathrm{rad}}\) declared in Subsection~\ref{sec:ch201}. They presuppose a
\emph{topology} on exported Maxwell pairs in which Cauchy sequences built from finite truncations,
mesh refinements, and growing observation shells converge to the same radiative object the laboratory
measures on a far sphere \cite{stratton1941,harrington2001}. Chapter~\ref{ch:02} already embedded
admissible fields in Sobolev product spaces through Subsection~\ref{sec:ch211} and singled out outgoing
Sommerfeld subclasses through Section~\ref{sec:ch24}; the Hilbert framework now asks which of those
analytic commitments survive \emph{after} \(\Pi_{\mathrm{rad}}\) has removed reactive and non-exporting
sheets \cite{stratton1941}. The answer is not a single norm copied from bounded domains: exterior
radiation carries flux through arbitrarily large shells, and the useful energy bookkeeping is tied to
uniform control of transverse amplitudes on spheres \(\Gamma_R\) as \(R\to\infty\)
\cite{stratton1941,jackson1999}.

Section~\ref{sec:ch31} supplies the Hilbert and rigged-Hilbert closure that later sections treat as
standing analytic infrastructure. Subsection~\ref{sec:ch311} refines \emph{finite-energy} radiative
spaces on top of the Sobolev classes of Subsection~\ref{sec:ch211} without re-displaying
Eqs.~\eqref{eq:ch2.hcurl-space}--\eqref{eq:ch2.admissible-field-space}. Subsection~\ref{sec:ch312}
installs inner products, orthogonality, and completeness compatible with reciprocity previews from
Subsection~\ref{sec:ch282}. Subsection~\ref{sec:ch313} develops continuous spectra and generalized
eigenfunctions, importing the abstract spectral skeleton of Subsection~\ref{sec:ch261}. Subsection~\ref{sec:ch314}
completes the rigged apparatus for distributional modes, supplying the mathematical legitimacy of formal
projectors noted in Subsection~\ref{sec:ch224} \cite{stratton1941,harrington2001}. The division is
deliberate: finite-energy carriers precede inner products; inner products precede spectral measures;
spectral measures precede rigged extensions for non-normalizable outgoing continua
\cite{stratton1941,hansen1988}.

Physically, the present section reverses the usual question: one asks how to
normalize a spherical mode before asking whether the exported state is square-integrable on the
observation shell that defines the mode \cite{harrington2001}. Volume~I reverses that pedagogy. A flux
meter on \(\Gamma_R\) returns a finite reading only when the transverse field energy on that shell is
integrable; a harmonic label attached earlier to a reactive-shell scan is not rescued by later
normalization \cite{stratton1941}. The analytic program of Section~\ref{sec:ch31} aligns mathematics with
that laboratory order: first declare which states carry finite shell flux uniformly in \(R\), then define
pairings that agree with those flux integrals, then expand in eigenfunctions whose normalization is
compatible with the same shell limits \cite{poynting1884,stratton1941}.


\subsection{Finite-energy radiative field spaces}
\label{sec:ch311}
% TOC tag: [HOME][USE SS2.1.1]

Square-integrability on a bounded feed region does not, by itself, certify that a Maxwell pair belongs to
the radiative export class. A closed-cavity eigenmode may have bounded \(L^{2}\) energy while contributing
nothing to \(\operatorname{ran}\mathcal{R}_\omega\); an incoming plane wave may have finite local density
while failing outgoing selection from Section~\ref{sec:ch24} \cite{stratton1941,jackson1999}. Subsection~\ref{sec:ch311}
therefore refines \(\mathcal{F}_{\mathrm{rad}}\) as a \emph{finite-energy radiative} subspace inside the
Sommerfeld outgoing class \(\mathcal{F}_{\mathrm{SF}}\) of Eq.~\eqref{eq:ch2.sommerfeld-outgoing-class},
inherited from the admissible product space \(\mathcal{U}_\omega\) and domain
\(\mathcal{F}_\Omega\) of Subsection~\ref{sec:ch211} \cite{stratton1941,harrington2001}. Sobolev
membership, trace maps, and source pairings from that subsection are cited here as standing infrastructure;
they are not re-derived \cite{stratton1941}.

\paragraph{Standing imports from Subsection~\ref{sec:ch211}.}
Work at fixed \(\omega>0\) on exterior domains \(\Omega_{\mathrm{ext}}\) with Lipschitz boundaries,
linear isotropic constitutive scalings on each patch, and impressed data \(\mathfrak{D}\) satisfying
Eq.~\eqref{eq:ch2.source-space-refined} \cite{stratton1941}. Field pairs \(\mathbf{u}_\omega=(\E,\H)\)
live in the Maxwell product space Eq.~\eqref{eq:ch2.maxwell-product-space}; admissible states satisfy
Eq.~\eqref{eq:ch2.admissible-field-space}; energy control is measured by
Eq.~\eqref{eq:ch2.field-energy-norm} on bounded truncations \cite{stratton1941,harrington2001}. Outgoing
closure is the composition Eq.~\eqref{eq:ch2.outgoing-radiation-refinement}, equivalent to
Eq.~\eqref{eq:ch2.rad-projection-composition} on propagating content \cite{stratton1941}. Chapter~\ref{ch:02}
already identified \(\mathcal{F}_{\mathrm{rad}}\) as a radiative refinement inside
\(\mathcal{F}_{\mathrm{SF}}\) with nonzero far-zone transverse amplitude
Eq.~\eqref{eq:ch2.far-zone-transverse-radiation} \cite{stratton1941,harrington2001}. Subsection~\ref{sec:ch311}
promotes that refinement to a Hilbert carrier by adding a radiative energy seminorm and a closure axiom
compatible with \(\operatorname{ran}\mathcal{R}_\omega\) \cite{stratton1941}.

\paragraph{Shell flux as the primitive finite-energy diagnostic.}
Let \(\Gamma_R=\{\mathbf{r}:|\mathbf{r}|=R\}\) denote a sphere of radius \(R\) in \(\Omega_{\mathrm{ext}}\),
with outward unit normal \(\hat{\mathbf{r}}\). For tangential components
\(\E_{T}=\hat{\mathbf{r}}\times(\E\times\hat{\mathbf{r}})\) and
\(\H_{T}=\hat{\mathbf{r}}\times(\H\times\hat{\mathbf{r}})\), define the shell radiation energy
\begin{equation}
\label{eq:ch3.shell-radiation-energy}
\mathcal{E}_{R}(\mathbf{u}_\omega)
=
\int_{\Gamma_R}
\Big(
|\E_{T}|^{2}+\eta_0^{2}\,|\H_{T}|^{2}
\Big)\,dS,
\end{equation}
where \(\eta_0=\sqrt{\mu_0/\varepsilon_0}\) is the free-space intrinsic impedance and \(dS\) is the
Euclidean area element \cite{stratton1941,jackson1999}. Equation~\eqref{eq:ch3.shell-radiation-energy}
is a transverse energy density on \(\Gamma_R\), not the time-averaged Poynting flux; the two coincide in
the far-zone transverse regime of Eq.~\eqref{eq:ch2.far-zone-transverse-radiation} up to the geometric
factor enforced by outgoing decay \cite{stratton1941,poynting1884}. Mathematically,
\(\mathcal{E}_{R}\) is a nonnegative functional on field pairs; physically, it is the quantity whose
uniform boundedness under increasing \(R\) certifies that radiation is not merely a local phase artifact
on a single shell \cite{harrington2001,stratton1941}.

A pair \(\mathbf{u}_\omega\) belongs to the \emph{uniform finite-shell} class when
\begin{equation}
\label{eq:ch3.uniform-shell-energy-class}
\sup_{R>R_{\min}}\ \mathcal{E}_{R}(\mathbf{u}_\omega)\ <\ \infty,
\end{equation}
for a fixed exterior onset radius \(R_{\min}\) determined by source support, boundary scatterers, and the
\(kr\) threshold of Theorem~\ref{thm:ch2.far-field-approximation} \cite{stratton1941,jackson1999}.
Equation~\eqref{eq:ch3.uniform-shell-energy-class} is stronger than finiteness on one observation sphere:
it forbids sequences whose transverse energy grows without bound as the laboratory enlarges the
measurement surface \cite{harrington2001}. Reactive near fields can produce large \(\mathcal{E}_{R}\) on
small \(R\) while failing Eq.~\eqref{eq:ch3.uniform-shell-energy-class} after outgoing selection; that
distinction is operator content, not a gauge choice \cite{stratton1941}.

\paragraph{Radiative energy seminorm.}
On bounded truncations \(\Omega_{\mathrm{ext}}\cap B_{R}\), Eq.~\eqref{eq:ch2.field-energy-norm} controls
volume energy and curl content. On the exterior, the relevant control is inherited from shell limits and
outgoing decay. Define the radiative energy seminorm
\begin{equation}
\label{eq:ch3.radiative-energy-seminorm}
\|\mathbf{u}_\omega\|_{\mathrm{rad}}^{2}
=
\limsup_{R\to\infty}\ R^{-2}\,\mathcal{E}_{R}(\mathbf{u}_\omega)
\;+\;
\|\mathbf{u}_\omega\|_{\mathcal{U},\,\Omega_{R_{\min}}}^{2},
\end{equation}
where \(\|\cdot\|_{\mathcal{U},\,\Omega_{R_{\min}}}\) denotes the restriction of
Eq.~\eqref{eq:ch2.field-energy-norm} to a fixed bounded neighborhood \(\Omega_{R_{\min}}\) containing
\(\operatorname{supp}\mathfrak{D}\) and all finite boundaries \cite{stratton1941,harrington2001}. The
\(R^{-2}\) weight encodes the geometric spreading of transverse energy for outgoing fields satisfying
Eq.~\eqref{eq:ch2.sommerfeld-field-condition}: finite \(\|\cdot\|_{\mathrm{rad}}\) requires the shell
integral to grow no faster than \(R^{2}\), matching the area element in
Eq.~\eqref{eq:ch3.shell-radiation-energy} \cite{stratton1941,jackson1999}. The local
\(\|\cdot\|_{\mathcal{U},\,\Omega_{R_{\min}}}\) term anchors the seminorm on the compact region where
sources and scatterers live, preventing purely outward-running plane waves with unbounded shell energy
from entering the finite-energy class \cite{stratton1941}.

\paragraph{Subspace chain and Hilbert closure.}
The analytic embeddings imported from Subsection~\ref{sec:ch211} specialize to the radiative sector as
\begin{equation}
\label{eq:ch3.radiative-subspace-chain}
\mathcal{F}_{\mathrm{rad}}
\ \subset\
\mathcal{F}_{\mathrm{SF}}
\ \subset\
\mathcal{F}_\Omega
\ \subset\
\mathcal{U}_\omega,
\end{equation}
with outgoing refinement Eq.~\eqref{eq:ch2.outgoing-radiation-refinement} acting between
\(\mathcal{F}_\Omega\) and \(\mathcal{F}_{\mathrm{SF}}\) \cite{stratton1941,harrington2001}. Subsection~\ref{sec:ch311}
closes \(\mathcal{F}_{\mathrm{rad}}\) as a Hilbert space by adjoining limits of Cauchy sequences in
\(\|\cdot\|_{\mathrm{rad}}\):
\begin{equation}
\label{eq:ch3.F-rad-hilbert-closure}
\mathcal{F}_{\mathrm{rad}}
=
\overline{
\left\{
\mathbf{u}_\omega\in\mathcal{F}_{\mathrm{SF}}:
\|\mathbf{u}_\omega\|_{\mathrm{rad}}<\infty
\right\}
}^{\|\cdot\|_{\mathrm{rad}}}.
\end{equation}
Equation~\eqref{eq:ch3.F-rad-hilbert-closure} is the Chapter~\ref{ch:03} carrier on which
Eqs.~\eqref{eq:ch3.eigenfield-coefficient-pairing} and~\eqref{eq:ch3.continuous-eigenfield-expansion}
will be proved well posed in Subsections~\ref{sec:ch312}--\ref{sec:ch313} \cite{stratton1941,hansen1988}.
Completeness is taken relative to \(\|\cdot\|_{\mathrm{rad}}\), not relative to the full
\(\|\cdot\|_{\mathcal{U}}\) norm on \(\mathcal{F}_\Omega\): evanescent components may be large in
\(\|\cdot\|_{\mathcal{U}}\) while invisible to \(\Pi_{\mathrm{rad}}\) \cite{harrington2001,stratton1941}.

\begin{figure}[t]
\centering
\ChOneFigBlock{%
\begin{tikzpicture}[ch1fig, node distance=7mm and 9mm]
  \node[ch1box] (u) {$\mathcal{U}_\omega$\\Subsec.~\ref{sec:ch211}};
  \node[ch1box, below=9mm of u] (f) {$\mathcal{F}_\Omega$};
  \node[ch1box, below=9mm of f] (sf) {$\mathcal{F}_{\mathrm{SF}}$\\Sec.~\ref{sec:ch24}};
  \node[ch1box, below=9mm of sf] (rad) {$\mathcal{F}_{\mathrm{rad}}$\\Eq.~\eqref{eq:ch3.F-rad-hilbert-closure}};
  \node[ch1box, left=14mm of sf] (out) {$\mathcal{O}_{\mathrm{out}}$};
  \node[ch1box, right=14mm of rad] (norm) {$\|\cdot\|_{\mathrm{rad}}$};
  \draw[ch1flow] (u) -- (f) -- (sf) -- (rad);
  \draw[ch1flow] (out) -- (sf);
  \draw[ch1flow] (norm) -- (rad);
\end{tikzpicture}%
}
\caption{Radiative Hilbert ladder for Subsection~\ref{sec:ch311}: Sobolev product spaces from
Subsection~\ref{sec:ch211} host admissible Maxwell pairs; Sommerfeld outgoing selection yields
\(\mathcal{F}_{\mathrm{SF}}\); finite-energy radiative closure Eq.~\eqref{eq:ch3.F-rad-hilbert-closure}
equips \(\mathcal{F}_{\mathrm{rad}}\) with \(\|\cdot\|_{\mathrm{rad}}\).}
\label{fig:ch3.radiative-hilbert-ladder}
\end{figure}

Figure~\ref{fig:ch3.radiative-hilbert-ladder} is the analytic counterpart to the operator gate of
Figure~\ref{fig:ch2.spectral-radiation-gate}, but the emphasis shifts from spectral projectors to
\emph{convergence}: which limits of physical approximations remain inside the same radiative Hilbert
class \cite{stratton1941,harrington2001}. The figure declares a norm hierarchy, not a
new radiation law: outgoing selection is still Eq.~\eqref{eq:ch2.outgoing-radiation-refinement}; the new
content is the closure under \(\|\cdot\|_{\mathrm{rad}}\) \cite{stratton1941}.

\paragraph{Far-amplitude equivalence and quotient structure.}
Two admissible pairs may differ by reactive near-field content while inducing identical far-zone
transverse amplitudes. For \(\mathbf{u},\mathbf{v}\in\mathcal{F}_{\mathrm{SF}}\), write
\begin{equation}
\label{eq:ch3.far-amplitude-equivalence}
\mathbf{u}\sim_{\mathrm{far}}\mathbf{v}
\quad\Longleftrightarrow\quad
\lim_{R\to\infty}\,
\big|\mathcal{E}_{R}(\mathbf{u}-\mathbf{v})\big|/R^{2}=0.
\end{equation}
Relation~\eqref{eq:ch3.far-amplitude-equivalence} identifies states that are indistinguishable on
increasing observation shells \cite{stratton1941,jackson1999}. On \(\mathcal{F}_{\mathrm{rad}}\), the
radiative seminorm is constant on equivalence classes: if \(\mathbf{u}\sim_{\mathrm{far}}\mathbf{v}\) and both
lie in Eq.~\eqref{eq:ch3.F-rad-hilbert-closure}, then
\(\|\mathbf{u}\|_{\mathrm{rad}}=\|\mathbf{v}\|_{\mathrm{rad}}\) \cite{stratton1941}. This quotient
structure is the analytic shadow of the uniqueness energy arguments in Subsection~\ref{sec:ch243}, cited
here without reopening the proof \cite{stratton1941,harrington2001}.

\paragraph{Embedding of radiation images.}
The export map from Chapter~\ref{ch:02} lands in the closed subspace:
\begin{equation}
\label{eq:ch3.ran-embedding-rad}
\operatorname{ran}\mathcal{R}_\omega
\ \subset\
\mathcal{F}_{\mathrm{rad}},
\qquad
\mathcal{R}_\omega[\mathfrak{D}]
\in
\mathcal{F}_{\mathrm{rad}}
\ \text{for admissible}\ \mathfrak{D},
\end{equation}
with \(\mathcal{R}_\omega\) from Eq.~\eqref{eq:ch2.outgoing-radiation-refinement} \cite{stratton1941,harrington2001}.
Equation~\eqref{eq:ch3.ran-embedding-rad} restates the operator declaration of
Subsection~\ref{sec:ch201} under the Hilbert refinement Eq.~\eqref{eq:ch3.F-rad-hilbert-closure}; it does
not extend \(\mathcal{R}_\omega\) to non-outgoing states \cite{stratton1941}. Source-to-state continuity is
quantified by
\begin{equation}
\label{eq:ch3.source-to-rad-energy-bound}
\|\mathcal{R}_\omega[\mathfrak{D}]\|_{\mathrm{rad}}
\ \le\
C_{\mathrm{rad}}\,\|\mathfrak{D}\|_{\mathcal{J}},
\end{equation}
for a constant \(C_{\mathrm{rad}}>0\) depending on \(\omega\), \(\Omega_s\), and exterior geometry, with
\(\|\mathfrak{D}\|_{\mathcal{J}}\) the source norm induced by
Eq.~\eqref{eq:ch2.source-space-refined} \cite{stratton1941,harrington2001}. Inequality
\eqref{eq:ch3.source-to-rad-energy-bound} is the radiative specialization of boundedness already implicit
in Theorem~\ref{thm:ch2.source-to-field-integral}; Subsection~\ref{sec:ch311} records the radiative norm
version because coefficient pairings in Section~\ref{sec:ch34} require \(\|\cdot\|_{\mathrm{rad}}\) control
on \(\mathcal{R}_\omega\) images \cite{hansen1988,stratton1941}.

\paragraph{Cauchy convergence and finite truncations.}
Numerical exterior solvers approximate \(\mathcal{F}_{\mathrm{rad}}\) through a growing family of
truncation radii \(R_{n}\uparrow\infty\) with Silver--M\"uller or exact outgoing conditions on
\(\Gamma_{R_n}\) \cite{stratton1941,harrington2001}. Let \(\mathbf{u}^{(n)}\in\mathcal{F}_{\mathrm{SF}}\)
denote the \(n\)-th approximate radiative state. Cauchy convergence in the radiative class is
\begin{equation}
\label{eq:ch3.cauchy-radiative-convergence}
\|\mathbf{u}^{(n)}-\mathbf{u}^{(m)}\|_{\mathrm{rad}}\xrightarrow[n,m\to\infty]{}0
\quad\Longrightarrow\quad
\exists\,\mathbf{u}_\omega\in\mathcal{F}_{\mathrm{rad}}:
\|\mathbf{u}^{(n)}-\mathbf{u}_\omega\|_{\mathrm{rad}}\to 0.
\end{equation}
The implication direction is completeness in Eq.~\eqref{eq:ch3.F-rad-hilbert-closure}; the forward
direction is what mesh-refinement studies must verify \cite{harrington2001}. Monotonicity of shell energy
for retarded outgoing images on nested truncations is recorded as
\begin{equation}
\label{eq:ch3.finite-truncation-monotonicity}
\mathcal{E}_{R_{n+1}}\!\big(\mathcal{R}_\omega^{(n)}[\mathfrak{D}]\big)
\ \le\
\mathcal{E}_{R_{n+1}}\!\big(\mathcal{R}_\omega^{(n+1)}[\mathfrak{D}]\big)
\ +\ \varepsilon_n,
\end{equation}
where \(\mathcal{R}_\omega^{(n)}\) is the radiation image computed with outer boundary on
\(\Gamma_{R_n}\) and \(\varepsilon_n\to 0\) when the absorbing closure converges to
Eq.~\eqref{eq:ch2.sommerfeld-field-condition} \cite{stratton1941,harrington2001}. Equation
\eqref{eq:ch3.finite-truncation-monotonicity} is a diagnostic inequality, not a new uniqueness theorem;
it explains why stable far patterns emerge only after truncation shells enter the radiative regime of
Theorem~\ref{thm:ch2.far-field-approximation} \cite{jackson1999,stratton1941}.

\begin{proposition}[Finite-energy radiative closure]
\label{prop:ch3.finite-energy-closure}
Let \(\mathfrak{D}\) be admissible impressed data on \(\Omega_s\) with outgoing exterior closure
Eq.~\eqref{eq:ch2.outgoing-radiation-refinement}. Then \(\mathcal{R}_\omega[\mathfrak{D}]\in\mathcal{F}_{\mathrm{rad}}\)
satisfies Eqs.~\eqref{eq:ch3.uniform-shell-energy-class}--\eqref{eq:ch3.source-to-rad-energy-bound}. Conversely,
every \(\mathbf{u}_\omega\in\mathcal{F}_{\mathrm{rad}}\) is the \(\|\cdot\|_{\mathrm{rad}}\)-limit of a sequence
\(\{\mathcal{R}_\omega[\mathfrak{D}_n]\}\) for admissible \(\mathfrak{D}_n\) with compact support, hence
\(\operatorname{ran}\mathcal{R}_\omega\) is dense in \(\mathcal{F}_{\mathrm{rad}}\) under the radiative seminorm.
\end{proposition}
\begin{proof}[Proof sketch]
The forward statement combines far-zone decay from
Eqs.~\eqref{eq:ch2.sommerfeld-field-condition}--\eqref{eq:ch2.far-zone-transverse-radiation} with
boundedness of the integral operator in Theorem~\ref{thm:ch2.source-to-field-integral} on
\(\mathcal{J}_\Omega\) \cite{stratton1941,harrington2001}. The shell bound
Eq.~\eqref{eq:ch3.uniform-shell-energy-class} follows from transverse \(1/r\) decay of outgoing
propagating amplitudes \cite{jackson1999,stratton1941}. The converse density uses smooth compactly
supported currents \(\mathfrak{D}_n\) approximating any \(\mathbf{u}_\omega\in\mathcal{F}_{\mathrm{rad}}\) on
\(\Omega_{R_{\min}}\) while matching far amplitudes through outgoing Hertz potentials, then passes to the
\(\|\cdot\|_{\mathrm{rad}}\)-limit in Eq.~\eqref{eq:ch3.F-rad-hilbert-closure} \cite{harrington2001,stratton1941}.
% PROOF: AUTHOR REVIEW
\end{proof}

\paragraph{Outgoing inheritance criterion.}
Membership in \(\mathcal{F}_{\mathrm{rad}}\) can be tested without re-opening Sobolev definitions:
\begin{equation}
\label{eq:ch3.outgoing-inheritance-criterion}
\mathbf{u}_\omega\in\mathcal{F}_{\mathrm{rad}}
\iff
\left(
\mathbf{u}_\omega\in\mathcal{F}_{\mathrm{SF}}
\ \wedge\
\|\mathbf{u}_\omega\|_{\mathrm{rad}}<\infty
\ \wedge\
\Pi_{\mathrm{rad}}[\mathbf{u}_\omega]=\mathbf{u}_\omega
\right).
\end{equation}
Criterion~\eqref{eq:ch3.outgoing-inheritance-criterion} packages
Eqs.~\eqref{eq:ch2.sommerfeld-outgoing-class} and~\eqref{eq:ch3.F-rad-hilbert-closure} for later sections:
Subsection~\ref{sec:ch312} will replace the seminorm clause with an inner product inducing the same
topology; Subsection~\ref{sec:ch313} will weaken the idempotence clause for generalized eigenfunctions in
the rigged extension \cite{stratton1941,hansen1988}. The projector idempotence condition removes
incoming and purely reactive sheets that satisfy Maxwell equations on bounded shells but are not fixed
points of \(\Pi_{\mathrm{rad}}\) \cite{harrington2001,stratton1941}.

\paragraph{Validity domain.}
Eqs.~\eqref{eq:ch3.shell-radiation-energy}--\eqref{eq:ch3.outgoing-inheritance-criterion} assume linear
media, fixed \(\omega\), and Sommerfeld outgoing \(\mathcal{B}_\Gamma\) on exterior problems as in
Section~\ref{sec:ch24} \cite{stratton1941}. They fail on closed PEC cavities where
Eq.~\eqref{eq:ch1.outer-boundary-classes} replaces outgoing closure, and on lossless open resonators where
discrete spectra require the quotient constructions of Eq.~\eqref{eq:ch2.coercivity-quotient}
\cite{stratton1941,harrington2001}. Absorbing layers approximate \(\mathcal{O}_{\mathrm{out}}\) but alter
\(\mathcal{E}_{R}\) unless the layer is transparent at distances where
Eq.~\eqref{eq:ch3.uniform-shell-energy-class} is evaluated \cite{harrington2001}. Anisotropic or
dispersive media are outside Volume~I scope; the seminorm weights in
Eq.~\eqref{eq:ch3.radiative-energy-seminorm} are fixed by free-space \(\eta_0\) on the observation shell
\cite{stratton1941}.

\paragraph{Worked illustration (short dipole export).}
For a Hertzian dipole \(\mathfrak{D}\) in the radiative regime, \(\mathcal{R}_\omega[\mathfrak{D}]\) satisfies
\(\mathcal{E}_{R}\sim R^{0}\) with transverse fields decaying as \(1/r\) on \(\Gamma_R\), hence
Eq.~\eqref{eq:ch3.uniform-shell-energy-class} holds with \(\sup_{R}\mathcal{E}_{R}\) equal to the
time-averaged radiated power budget imported from Subsection~\ref{sec:ch253} \cite{stratton1941,jackson1999}.
The same state has large \(\|\cdot\|_{\mathcal{U}}\) near the feed, illustrating why
\(\|\cdot\|_{\mathrm{rad}}\) rather than the full product norm is the correct radiative topology
\cite{harrington2001}. Attempting to expand reactive near fields in outgoing eigenfunctions violates
Eq.~\eqref{eq:ch3.outgoing-inheritance-criterion} even when the near scan is visually ``wave-like''
\cite{stratton1941}.

\paragraph{Worked illustration (incoming plane wave).}
For \(\E=\E_{0}e^{+\jj k_{0} z}\) on \(\Omega_{\mathrm{ext}}\), local energy densities are finite on every
bounded shell, yet \(\mathcal{O}_{\mathrm{out}}\) annihilates the mode and \(\Pi_{\mathrm{rad}}[\mathbf{u}]=\mathbf{0}\)
\cite{stratton1941,jackson1999}. Such a pair lies outside \(\mathcal{F}_{\mathrm{rad}}\) because it fails the
idempotence clause in Eq.~\eqref{eq:ch3.outgoing-inheritance-criterion}, not because it lacks
\(\HCurl\) regularity \cite{harrington2001}. Confusing square integrability on a fixed ball with membership
in \(\mathcal{F}_{\mathrm{rad}}\) is a common category error when bounded-domain functional
analysis verbatim to open radiation problems \cite{stratton1941}.

\paragraph{Worked illustration (nested truncation meshes).}
A MoM solver on \(\Gamma_{R_n}\) with \(R_{n}=10\lambda_0,20\lambda_0,\ldots\) produces
\(\mathbf{u}^{(n)}\) whose shell energies \(\mathcal{E}_{R_{n}}\) stabilize once
\(R_{n}\) exceeds the regime threshold of Theorem~\ref{thm:ch2.far-field-approximation}
\cite{harrington2001,jackson1999}. Before that threshold, \(\mathcal{E}_{R_{n}}\) may grow with \(n\) because
evanescent and reflected components have not yet separated from outgoing transverse flux
\cite{stratton1941}. Interpreting early truncations as \(\mathcal{F}_{\mathrm{rad}}\) states is premature;
Eq.~\eqref{eq:ch3.finite-truncation-monotonicity} is the quantitative guard \cite{harrington2001}.

\paragraph{Energy-class alignment with Chapter~\ref{ch:01}.}
Eq.~\eqref{eq:ch1.energy-class} declared finite time-averaged electromagnetic energy as an admissibility
screen on bounded domains \cite{stratton1941}. On exterior radiation problems the same screen must be read
through shell limits rather than through a single volume integral over all space
\cite{harrington2001,jackson1999}. A radiative state satisfies the Chapter~\ref{ch:01} energy intent when
Eq.~\eqref{eq:ch3.uniform-shell-energy-class} holds and the local neighborhood term in
Eq.~\eqref{eq:ch3.radiative-energy-seminorm} remains finite on \(\Omega_{R_{\min}}\) \cite{stratton1941}.
That reading does not resurrect stored cavity energy as export: closed resonators fail outgoing inheritance
Eq.~\eqref{eq:ch3.outgoing-inheritance-criterion} even when their \(L^{2}\) energy is bounded
\cite{harrington2001}. The alignment is therefore \emph{conditional} on Sommerfeld closure, matching the
boundary-class logic of Section~\ref{sec:ch24} rather than extending Eq.~\eqref{eq:ch1.energy-class}
verbatim to unbounded volumes \cite{stratton1941}.

\begin{table}[t]
\centering
\caption{Norm diagnostics on exported Maxwell pairs (Subsection~\ref{sec:ch311}). Each row answers a
different laboratory question; none substitutes for outgoing selection
Eq.~\eqref{eq:ch2.outgoing-radiation-refinement}.}
\label{tab:ch3.radiative-norm-diagnostic}
\small
\begin{tabular}{@{}p{0.22\linewidth}p{0.30\linewidth}p{0.38\linewidth}@{}}
\toprule
Diagnostic & Primary formula anchor & Detects / fails to detect \\
\midrule
Volume product norm &
Eq.~\eqref{eq:ch2.field-energy-norm} on \(\Omega_{R_{\min}}\) &
Large reactive storage near sources; blind to incoming sheets at infinity \\
\addlinespace
Shell transverse energy &
Eq.~\eqref{eq:ch3.shell-radiation-energy} &
Flux-capable content on \(\Gamma_R\); not idempotent under \(\Pi_{\mathrm{rad}}\) alone \\
\addlinespace
Uniform shell class &
Eq.~\eqref{eq:ch3.uniform-shell-energy-class} &
Runaway growth on expanding observation surfaces \\
\addlinespace
Radiative seminorm &
Eq.~\eqref{eq:ch3.radiative-energy-seminorm} &
Finite-energy \(\mathcal{F}_{\mathrm{rad}}\) topology for closure and pairings \\
\addlinespace
Far equivalence &
Eq.~\eqref{eq:ch3.far-amplitude-equivalence} &
Near-field differences invisible to far meters \\
\bottomrule
\end{tabular}
\end{table}

Table~\ref{tab:ch3.radiative-norm-diagnostic} is a reader's compass, not a replacement for proofs.
One must not treat finiteness of Eq.~\eqref{eq:ch2.field-energy-norm} on a large ball as proof of
radiative membership; the table rows separate that mistake into testable components
\cite{harrington2001,stratton1941}. Row four is the topology Subsection~\ref{sec:ch312} will refine into
an inner product; rows one--three explain why that refinement cannot be chosen as a naive restriction of
the bounded-domain norm \cite{stratton1941}.

\paragraph{Stability under superposition in \(\mathcal{F}_{\mathrm{rad}}\).}
Because Eq.~\eqref{eq:ch3.F-rad-hilbert-closure} is a vector space, finite superpositions of radiation
images remain in the class when coefficients are bounded:
\begin{equation}
\label{eq:ch3.radiative-superposition-stability}
\left\|\sum_{n=1}^{N} a_n\,\mathcal{R}_\omega[\mathfrak{D}_n]\right\|_{\mathrm{rad}}
\ \le\
\sum_{n=1}^{N}|a_n|\,
\|\mathcal{R}_\omega[\mathfrak{D}_n]\|_{\mathrm{rad}},
\end{equation}
for admissible \(\mathfrak{D}_n\) and scalars \(a_n\in\mathbb{C}\) \cite{stratton1941,harrington2001}.
Inequality~\eqref{eq:ch3.radiative-superposition-stability} is the analytic backbone of linear modal
synthesis: far patterns add in coefficient space only because \(\mathcal{F}_{\mathrm{rad}}\) is a linear
subspace closed under the same seminorm that meters shell flux \cite{stratton1941}. Nonlinear superposition
claims---for example, interpreting intensity fringes as independent mode sums without phase control---lie
outside Eq.~\eqref{eq:ch3.radiative-superposition-stability} and return to measurement composition
Eq.~\eqref{eq:ch3.measurement-on-coefficients} \cite{harrington2001}.

\paragraph{Link to representation admissibility.}
The first clause of Eq.~\eqref{eq:ch3.representation-admissibility} now carries a precise meaning:
\(\mathfrak{F}\in\mathcal{F}_{\mathrm{rad}}\) is equivalent to
Eqs.~\eqref{eq:ch3.outgoing-inheritance-criterion} and~\eqref{eq:ch3.uniform-shell-energy-class}
\cite{stratton1941}. Subsection~\ref{sec:ch303} stated that flux-normalizable coordinates require an
outgoing-qualified carrier; Subsection~\ref{sec:ch311} supplies the carrier's topology
\cite{harrington2001}. Later harmonic constructions therefore need not re-prove finiteness of shell
energy when they cite \(\mathcal{R}_\omega[\mathfrak{D}]\in\mathcal{F}_{\mathrm{rad}}\) from
Eq.~\eqref{eq:ch3.ran-embedding-rad} \cite{stratton1941,hansen1988}.

Subsection~\ref{sec:ch311} has promoted \(\mathcal{F}_{\mathrm{rad}}\) from an operator symbol to a
finite-energy Hilbert carrier with shell diagnostics Eqs.~\eqref{eq:ch3.shell-radiation-energy}--\eqref{eq:ch3.outgoing-inheritance-criterion},
embedding Eq.~\eqref{eq:ch3.ran-embedding-rad}, and closure
Eq.~\eqref{eq:ch3.F-rad-hilbert-closure}. Subsection~\ref{sec:ch312} replaces the seminorm by a
reciprocity-compatible inner product on the same class; Subsection~\ref{sec:ch313} then reads
Theorem~\ref{thm:ch2.radiation-spectral-decomposition} as a spectral theorem on that product space rather
than as a formal mode sum \cite{stratton1941,hansen1988,harrington2001}.


\subsection{Inner products, orthogonality, and completeness}
\label{sec:ch312}
% TOC tag: [HOME][USE SS2.8.2]

A seminorm declares which approximations converge; an inner product declares how two convergent states
\emph{interfere} when sources, receivers, and normalization conventions are interchanged. Subsection~\ref{sec:ch311}
equipped \(\mathcal{F}_{\mathrm{rad}}\) with \(\|\cdot\|_{\mathrm{rad}}\) through shell control and local
product-space restriction; Subsection~\ref{sec:ch312} installs the sesquilinear pairing that turns that topology
into a Hilbert space geometry compatible with flux meters, Lorentz reciprocity from
Subsection~\ref{sec:ch281}, and the modal orthogonality/coupling package of
Subsection~\ref{sec:ch282} \cite{stratton1941,harrington2001}. The present subsection is the Chapter~\ref{ch:03}
definitive account of the \emph{field-space} inner product on \(\mathcal{F}_{\mathrm{rad}}\); explicit angular
integrals for vector spherical harmonics remain deferred to Subsection~\ref{sec:ch343}
\cite{hansen1988,stratton1941}.

\paragraph{Standing imports and scope.}
Fix \(\omega>0\), reciprocal media Eq.~\eqref{eq:ch2.reciprocal-medium-class}, passive boundary class
Eq.~\eqref{eq:ch2.passive-impedance-class}, and outgoing radiative states satisfying
Eq.~\eqref{eq:ch3.outgoing-inheritance-criterion} \cite{stratton1941,harrington2001}. Maxwell adjoint pairing
Eq.~\eqref{eq:ch2.maxwell-adjoint-pairing}, radiation coupling bilinear
Eq.~\eqref{eq:ch2.radiative-coupling-bilinear}, modal orthogonality and coupling Theorem
~\ref{thm:ch2.radiation-modal-orthogonality-coupling}, and spectral completeness Theorem
~\ref{thm:ch2.radiation-spectral-decomposition} are imported from Chapter~\ref{ch:02} as operator facts on
the declared domain class \cite{stratton1941,harrington2001}. Subsection~\ref{sec:ch312} does not re-display
those Chapter~\ref{ch:02} equations; it defines the Chapter~\ref{ch:03} inner product that makes them
coefficient-ready on \(\mathcal{F}_{\mathrm{rad}}\) \cite{stratton1941}.

\paragraph{Radiative inner product.}
For \(\mathbf{u}_\omega=(\E_u,\H_u)\) and \(\mathbf{v}_\omega=(\E_v,\H_v)\) in
Eq.~\eqref{eq:ch3.F-rad-hilbert-closure}, define
\begin{equation}
\label{eq:ch3.radiative-inner-product-def}
\langle \mathbf{u}_\omega,\mathbf{v}_\omega\rangle_{\mathrm{rad}}
=
\langle \mathbf{u}_\omega,\mathbf{v}_\omega\rangle_{\mathcal{U},\,\Omega_{R_{\min}}}
\;+\;
\lim_{R\to\infty}\,
R^{-2}
\int_{\Gamma_R}
\Big(
\E_{u,T}\cdot\E_{v,T}^{*}
+\eta_0^{2}\,\H_{u,T}\cdot\H_{v,T}^{*}
\Big)\,dS,
\end{equation}
with tangential components as in Eq.~\eqref{eq:ch3.shell-radiation-energy} and local pairing induced by
Eq.~\eqref{eq:ch2.field-energy-norm} on \(\Omega_{R_{\min}}\) \cite{stratton1941,harrington2001}. The limit
exists and is independent of the particular sequence \(R\to\infty\) when both arguments satisfy
Eq.~\eqref{eq:ch3.uniform-shell-energy-class} \cite{stratton1941}. Equation~\eqref{eq:ch3.radiative-inner-product-def}
is sesquilinear, conjugate-symmetric on reciprocal media, and positive on nontrivial radiative states
\cite{stratton1941,harrington2001}. Algebraically, it packages a bounded neighborhood pairing with a
far-shell transverse overlap; physically, it weights \emph{export overlap} on large spheres while retaining
local elliptic control near sources and scatterers \cite{jackson1999,harrington2001}.

The induced norm agrees with Subsection~\ref{sec:ch311}:
\begin{equation}
\label{eq:ch3.inner-product-seminorm-identity}
\|\mathbf{u}_\omega\|_{\mathrm{rad}}^{2}
=
\langle \mathbf{u}_\omega,\mathbf{u}_\omega\rangle_{\mathrm{rad}},
\qquad
\mathbf{u}_\omega\in\mathcal{F}_{\mathrm{rad}},
\end{equation}
so Cauchy convergence Eq.~\eqref{eq:ch3.cauchy-radiative-convergence} is equivalently convergence in the
inner product topology \cite{stratton1941}. Equation~\eqref{eq:ch3.inner-product-seminorm-identity} is the
bridge that lets later sections state spectral theorems on \(\mathcal{F}_{\mathrm{rad}}\) without toggling
between norm and inner-product language \cite{harrington2001}.

\paragraph{Flux representation of the far-shell term.}
The shell integral in Eq.~\eqref{eq:ch3.radiative-inner-product-def} is not an alternative flux law; it is
the bilinear companion of the power functional Eq.~\eqref{eq:ch2.flux-on-radiative-branch}
\cite{stratton1941,poynting1884}. For identical arguments \(\mathbf{u}_\omega=\mathbf{v}_\omega\), the
far-shell contribution of Eq.~\eqref{eq:ch3.radiative-inner-product-def} tracks the time-averaged radiative
flux budget once Theorem~\ref{thm:ch2.far-field-approximation} applies \cite{stratton1941,jackson1999}. For
distinct arguments, the same term measures coherent export overlap---the geometric content later captured by
coupling matrices in Subsection~\ref{sec:ch283} \cite{harrington2001}. Reactive fields removed by
\(\Pi_{\mathrm{rad}}\) do not contribute to the far-shell limit because
Eq.~\eqref{eq:ch2.rad-null-on-evanescent} annihilates their radiative images
\cite{stratton1941,harrington2001}.

\paragraph{Adjoint pairing restricted to \(\mathcal{F}_{\mathrm{rad}}\).}
On outgoing homogeneous exterior problems with passive \(\mathcal{B}_\Gamma\), the boundary term in
Eq.~\eqref{eq:ch2.maxwell-adjoint-pairing} cancels when both arguments lie in
\(\mathcal{F}_{\mathrm{SF}}\) and volume sources vanish in the exterior region
\cite{stratton1941,harrington2001}. Subsection~\ref{sec:ch312} records the radiative specialization as
\begin{equation}
\label{eq:ch3.adjoint-pairing-restriction-rad}
\mathfrak{a}_\omega(\mathbf{u}_\omega,\mathbf{v}_\omega)
=
\langle \mathcal{M}_\omega[\mathbf{u}_\omega],\mathbf{v}_\omega\rangle_{\Omega_{\mathrm{ext}}}
\quad\text{for}\ 
\mathbf{u}_\omega,\mathbf{v}_\omega\in\mathcal{F}_{\mathrm{rad}}
\ \text{with homogeneous exterior data,}
\end{equation}
importing the cancellation logic of Subsection~\ref{sec:ch213} without reopening integration-by-parts
proofs \cite{stratton1941}. Equation~\eqref{eq:ch3.adjoint-pairing-restriction-rad} is the local Maxwell
form of adjointness; Eq.~\eqref{eq:ch3.radiative-inner-product-def} is the global radiative energy form used
for mode normalization and coefficient extraction \cite{harrington2001}. Both pairings agree on source-free
outgoing test pairs when constitutive reciprocity Eq.~\eqref{eq:ch2.reciprocal-medium-class} holds
\cite{stratton1941}.

\paragraph{Orthogonality on \(\mathcal{F}_{\mathrm{rad}}\).}
Vectors \(\mathbf{u}_\omega,\mathbf{v}_\omega\in\mathcal{F}_{\mathrm{rad}}\) are \emph{radiatively orthogonal}
when
\begin{equation}
\label{eq:ch3.orthogonality-radiative-def}
\langle \mathbf{u}_\omega,\mathbf{v}_\omega\rangle_{\mathrm{rad}}=0.
\end{equation}
Equation~\eqref{eq:ch3.orthogonality-radiative-def} is the Chapter~\ref{ch:03} field-space definition;
modal orthogonality on eigenbases is the specialization proved in Theorem
~\ref{thm:ch2.radiation-modal-orthogonality-coupling}(i) once
\(\{\mathbf{u}_n\}\) is the orthonormal family of Theorem~\ref{thm:ch2.radiation-spectral-decomposition}
\cite{stratton1941,harrington2001}. Distinct exporting modes with non-overlapping far patterns are orthogonal
in the sense of Eq.~\eqref{eq:ch3.orthogonality-radiative-def}; degenerate rotational blocks may be
orthogonal as subspaces while individual laboratory-fixed components are not, as analyzed in
Subsection~\ref{sec:ch282} through block coupling matrices \cite{jackson1999,harrington2001}. Near-field
orthogonality in \(\|\cdot\|_{\mathcal{U}}\) on a small shell does \emph{not} imply
Eq.~\eqref{eq:ch3.orthogonality-radiative-def}: reactive content can be locally orthogonal while remaining
export-coupled at infinity \cite{stratton1941}.

\paragraph{Reciprocity as inner-product symmetry.}
Theorem~\ref{thm:ch2.lorentz-reciprocity-operator}(iv) and radiation reciprocity
Eq.~\eqref{eq:ch2.radiation-map-reciprocity} import, on reciprocal media, the symmetry
\begin{equation}
\label{eq:ch3.reciprocity-symmetric-inner-product}
\langle \mathcal{R}_\omega[\Jimp^{(1)}],\,\mathcal{R}_\omega[\Jimp^{(2)}]\rangle_{\mathrm{rad}}
=
\langle \mathcal{R}_\omega[\Jimp^{(2)}],\,\mathcal{R}_\omega[\Jimp^{(1)}]\rangle_{\mathrm{rad}}
\end{equation}
for admissible \(\Jimp^{(1)},\Jimp^{(2)}\in\mathcal{J}_\Omega\) \cite{stratton1941,harrington2001}.
Equation~\eqref{eq:ch3.reciprocity-symmetric-inner-product} is the Chapter~\ref{ch:03} restatement of
coupling symmetry Eq.~\eqref{eq:ch2.coupling-bilinear-symmetry} with the inner product
Eq.~\eqref{eq:ch3.radiative-inner-product-def} in place of the preview energy pairing
\cite{stratton1941}. Physically, swapping transmit and receive roles leaves the bilinear coupling unchanged
when media, boundaries, and outgoing closure match \cite{harrington2001,jackson1999}. Normalization choices in
Section~\ref{sec:ch35} must preserve Eq.~\eqref{eq:ch3.reciprocity-symmetric-inner-product}; rescaling modes
without updating flux meters breaks the symmetry even when far patterns appear unchanged
\cite{stratton1941,harrington2001}.

\begin{figure}[t]
\centering
\ChOneFigBlock{%
\begin{tikzpicture}[ch1fig, node distance=7mm and 9mm]
  \node[ch1box] (j1) {$\Jimp^{(1)}$};
  \node[ch1box, right=11mm of j1] (r1) {$\mathcal{R}_\omega$};
  \node[ch1box, right=11mm of r1] (u1) {$\mathbf{u}^{(1)}$};
  \node[ch1box, below=12mm of j1] (j2) {$\Jimp^{(2)}$};
  \node[ch1box, below=12mm of r1] (r2) {$\mathcal{R}_\omega$};
  \node[ch1box, below=12mm of u1] (u2) {$\mathbf{u}^{(2)}$};
  \node[ch1box, right=14mm of u1] (ip) {$\langle\cdot,\cdot\rangle_{\mathrm{rad}}$};
  \node[ch1box, below=8mm of ip] (rec) {swap\\$\Jimp^{(1)}\!\leftrightarrow\!\Jimp^{(2)}$};
  \draw[ch1flow] (j1) -- (r1) -- (u1);
  \draw[ch1flow] (j2) -- (r2) -- (u2);
  \draw[ch1flow] (u1) -- (ip);
  \draw[ch1flow] (u2) -- (ip);
  \draw[ch1flow, dashed] (rec) -- (ip);
\end{tikzpicture}%
}
\caption{Reciprocity geometry for Subsection~\ref{sec:ch312}: impressed sources map to radiative images;
the inner product Eq.~\eqref{eq:ch3.radiative-inner-product-def} tests export overlap; reciprocity leaves
the pairing invariant under source swap Eq.~\eqref{eq:ch3.reciprocity-symmetric-inner-product}.}
\label{fig:ch3.reciprocity-inner-product-geometry}
\end{figure}

Figure~\ref{fig:ch3.reciprocity-inner-product-geometry} is the bilinear reading of
Figure~\ref{fig:ch2.modal-coupling-pipeline} from Subsection~\ref{sec:ch282}, with the Chapter~\ref{ch:03}
inner product made explicit \cite{stratton1941,harrington2001}. The figure omits modal sums: those enter only
after a spectral basis is chosen in Section~\ref{sec:ch34} \cite{hansen1988}.

\paragraph{Coupling of radiation images.}
For arbitrary admissible sources, the Chapter~\ref{ch:02} radiative coupling bilinear
Eq.~\eqref{eq:ch2.radiative-coupling-bilinear} agrees with the Chapter~\ref{ch:03} inner product on images:
\begin{equation}
\label{eq:ch3.coupling-images-inner-product}
\mathfrak{C}_\omega\!\big(\Jimp^{(1)},\Jimp^{(2)}\big)
=
\langle \mathcal{R}_\omega[\Jimp^{(1)}],\,\mathcal{R}_\omega[\Jimp^{(2)}]\rangle_{\mathrm{rad}}
\quad\text{on}\ \mathcal{F}_{\mathrm{rad}},
\end{equation}
with \(\mathfrak{C}_\omega\) imported from Subsection~\ref{sec:ch282} \cite{stratton1941,harrington2001}.
Equation~\eqref{eq:ch3.coupling-images-inner-product} identifies \emph{source coupling} with \emph{field
overlap} once \(\Pi_{\mathrm{rad}}\) has been applied; it is the operator route to coefficient formulas
without re-expanding Green kernels \cite{stratton1941}. Theorem
~\ref{thm:ch2.radiation-modal-orthogonality-coupling}(iii)--(iv) and Proposition
~\ref{prop:ch2.superposition-cross-coupling} remain valid when
\(\langle\cdot,\cdot\rangle_{\mathcal{U}}\) in those statements is read as
\(\langle\cdot,\cdot\rangle_{\mathrm{rad}}\) on \(\operatorname{ran}\mathcal{R}_\omega\)
\cite{stratton1941,harrington2001}. Source nullity \(\ker\mathcal{R}_\omega\) leaves
Eq.~\eqref{eq:ch3.coupling-images-inner-product} unchanged, matching Theorem
~\ref{thm:ch2.radiation-modal-orthogonality-coupling}(iv) \cite{harrington2001}.

\paragraph{Riesz representation of coefficients.}
Bounded linear functionals on \(\mathcal{F}_{\mathrm{rad}}\) admit Riesz representation in
\(\langle\cdot,\cdot\rangle_{\mathrm{rad}}\). For each outgoing eigenfunction \(\mathbf{f}_\alpha\) in the
expansion Eq.~\eqref{eq:ch3.eigenfield-expansion-on-frad},
\begin{equation}
\label{eq:ch3.riesz-coefficient-representation}
c_\alpha[\mathfrak{D}]
=
\big\langle \mathcal{R}_\omega[\mathfrak{D}],\,\mathbf{f}_\alpha\big\rangle_{\mathrm{rad}},
\qquad
\alpha\in I,
\end{equation}
refining Eq.~\eqref{eq:ch3.eigenfield-coefficient-pairing} with the Subsection~\ref{sec:ch312} inner product
in place of the raw product pairing \cite{stratton1941,harrington2001}. Equation
\eqref{eq:ch3.riesz-coefficient-representation} is the eigenfield restatement of
Eq.~\eqref{eq:ch3.coefficient-functional} once \(\mathcal{C}\) is fixed; it agrees with source amplitudes
Eq.~\eqref{eq:ch2.source-modal-amplitude} and modal coupling
Eq.~\eqref{eq:ch2.modal-coupling-coefficient} under the normalization conventions of Section~\ref{sec:ch35}
\cite{stratton1941,hansen1988}. Linearity in \(\mathfrak{D}\) follows from
Eq.~\eqref{eq:ch2.radiation-linearity} and sesquilinearity of
Eq.~\eqref{eq:ch3.radiative-inner-product-def} \cite{stratton1941}.

\paragraph{Parseval-type energy accounting.}
When \(\{\mathbf{f}_\alpha\}_{\alpha\in I}\) is an orthonormal frame in
\(\langle\cdot,\cdot\rangle_{\mathrm{rad}}\) and the expansion
Eq.~\eqref{eq:ch3.eigenfield-expansion-on-frad} converges,
\begin{equation}
\label{eq:ch3.parseval-radiative-preview}
\|\mathcal{R}_\omega[\mathfrak{D}]\|_{\mathrm{rad}}^{2}
=
\sum_{\alpha\in I}
\big|c_\alpha[\mathfrak{D}]\big|^{2}
\;+\;
\|R_{\mathrm{trunc}}[\mathfrak{D}]\|_{\mathrm{rad}}^{2},
\end{equation}
with remainder \(R_{\mathrm{trunc}}\) from Eq.~\eqref{eq:ch3.eigenfield-truncation-remainder}
\cite{stratton1941,harrington2001}. Equation~\eqref{eq:ch3.parseval-radiative-preview} is a finite-energy
accounting identity: modal coefficients partition radiative energy up to truncation error, not reactive storage
\cite{harrington2001}. Full completeness needed to set \(\|R_{\mathrm{trunc}}\|=0\) for all admissible states
is imported from Theorem~\ref{thm:ch2.radiation-spectral-decomposition} and proved in the spherical basis in
Subsection~\ref{sec:ch343} \cite{stratton1941,hansen1988}.

\paragraph{Completeness.}
Combining Eq.~\eqref{eq:ch3.F-rad-hilbert-closure} with
Eqs.~\eqref{eq:ch3.radiative-inner-product-def} and~\eqref{eq:ch3.inner-product-seminorm-identity},
\((\mathcal{F}_{\mathrm{rad}},\langle\cdot,\cdot\rangle_{\mathrm{rad}})\) is a Hilbert space. Completeness
can be stated equivalently as
\begin{equation}
\label{eq:ch3.completeness-radiative-characterization}
\text{every Cauchy sequence in}\ \langle\cdot,\cdot\rangle_{\mathrm{rad}}
\ \text{converges to a unique element of}\ \mathcal{F}_{\mathrm{rad}},
\end{equation}
and, by Proposition~\ref{prop:ch3.finite-energy-closure}, \(\operatorname{ran}\mathcal{R}_\omega\) is dense in
\(\mathcal{F}_{\mathrm{rad}}\) \cite{stratton1941,harrington2001}. Spectral completeness---that
orthonormal radiation eigenmodes span \(\mathcal{F}_{\mathrm{rad}}\)---is Theorem
~\ref{thm:ch2.radiation-spectral-decomposition} read in the inner product of
Eq.~\eqref{eq:ch3.radiative-inner-product-def}; Subsection~\ref{sec:ch313} develops the continuous-spectrum
version required for rigged extensions \cite{stratton1941,hansen1988}. Subsection~\ref{sec:ch312} does not
re-prove completeness of eigenmodes; it supplies the Hilbert space on which that theorem acts
\cite{harrington2001}.

\begin{proposition}[Inner product induces the radiative Hilbert topology]
\label{prop:ch3.inner-product-hilbert-topology}
Equations~\eqref{eq:ch3.radiative-inner-product-def} and~\eqref{eq:ch3.inner-product-seminorm-identity}
define a Hilbert inner product on \(\mathcal{F}_{\mathrm{rad}}\) whose topology coincides with
\(\|\cdot\|_{\mathrm{rad}}\) from Subsection~\ref{sec:ch311}. For every bounded linear functional
\(\Lambda\) on \(\mathcal{F}_{\mathrm{rad}}\) there exists unique \(\mathbf{w}_\omega\in\mathcal{F}_{\mathrm{rad}}\)
such that \(\Lambda(\mathbf{u}_\omega)=\langle\mathbf{u}_\omega,\mathbf{w}_\omega\rangle_{\mathrm{rad}}\).
\end{proposition}
\begin{proof}[Proof sketch]
Positive definiteness on \(\mathcal{F}_{\mathrm{rad}}\) follows from
Eq.~\eqref{eq:ch3.uniform-shell-energy-class} and the local term in
Eq.~\eqref{eq:ch3.radiative-inner-product-def} \cite{stratton1941}. Topology coincidence is
Eq.~\eqref{eq:ch3.inner-product-seminorm-identity} together with
Eq.~\eqref{eq:ch3.cauchy-radiative-convergence} \cite{harrington2001}. Riesz representation uses
completeness Eq.~\eqref{eq:ch3.completeness-radiative-characterization} \cite{stratton1941}.
% PROOF: AUTHOR REVIEW
\end{proof}

\begin{proposition}[Reciprocity-compatible orthonormal frames]
\label{prop:ch3.reciprocity-orthonormal-frame}
Let \(\{\mathbf{u}_n\}\) be an orthonormal radiation eigenbasis in
\(\langle\cdot,\cdot\rangle_{\mathrm{rad}}\) on the domain class of Theorem
~\ref{thm:ch2.radiation-spectral-decomposition}, with reciprocal media
Eq.~\eqref{eq:ch2.reciprocal-medium-class}. Then Theorem
~\ref{thm:ch2.radiation-modal-orthogonality-coupling} holds with
\(\langle\cdot,\cdot\rangle_{\mathcal{U}}\) replaced by
\(\langle\cdot,\cdot\rangle_{\mathrm{rad}}\), and adjoint modes satisfy
Eq.~\eqref{eq:ch2.reciprocal-adjoint-mode-identification} on exporting sectors
\cite{stratton1941,harrington2001}.
\end{proposition}
\begin{proof}[Proof sketch]
Replace the energy pairing in Theorem~\ref{thm:ch2.radiation-modal-orthogonality-coupling} by
Eq.~\eqref{eq:ch3.coupling-images-inner-product} and apply
Eq.~\eqref{eq:ch3.reciprocity-symmetric-inner-product} \cite{stratton1941,harrington2001}.
% PROOF: AUTHOR REVIEW
\end{proof}

\paragraph{Validity domain.}
Eqs.~\eqref{eq:ch3.radiative-inner-product-def}--\eqref{eq:ch3.completeness-radiative-characterization}
require linear reciprocal media Eq.~\eqref{eq:ch2.reciprocal-medium-class}, passive
Eq.~\eqref{eq:ch2.passive-impedance-class} boundaries, and outgoing closure on exterior problems
\cite{stratton1941}. Non-reciprocal gyrotropic media break
Eq.~\eqref{eq:ch3.reciprocity-symmetric-inner-product}; biorthogonality with distinct adjoint modes from
Eq.~\eqref{eq:ch2.radiation-biorthogonality} remains the general template \cite{stratton1941,harrington2001}.
Lossy media are admissible when passivity is maintained; growth in the local
\(\langle\cdot,\cdot\rangle_{\mathcal{U},\,\Omega_{R_{\min}}}\) term must still be controlled for membership
in \(\mathcal{F}_{\mathrm{rad}}\) \cite{stratton1941}. Discrete cavity spectra outside
\(\mathcal{F}_{\mathrm{rad}}\) are excluded by Eq.~\eqref{eq:ch3.outgoing-inheritance-criterion}
\cite{harrington2001}.

\paragraph{Worked illustration (two-port coupling).}
Consider reciprocal antennas \(A\) and \(B\) with impressed currents \(\Jimp^{(A)},\Jimp^{(B)}\) and
radiation images \(\mathbf{u}^{(A)}=\mathcal{R}_\omega[\Jimp^{(A)}]\),
\(\mathbf{u}^{(B)}=\mathcal{R}_\omega[\Jimp^{(B)}]\) \cite{harrington2001,stratton1941}. The receive coupling
of \(A\) when \(B\) transmits is proportional to
\(\langle \mathbf{u}^{(A)},\mathbf{u}^{(B)}\rangle_{\mathrm{rad}}\) under standard matched normalization;
reciprocity Eq.~\eqref{eq:ch3.reciprocity-symmetric-inner-product} equates transmit--receive interchange
\cite{stratton1941,jackson1999}. A MoM implementation that uses \(\|\cdot\|_{\mathcal{U}}\) on a near shell
for this bilinear form can report asymmetric coupling while still satisfying volume Maxwell equations: the
error is a pairing mistake, not a failure of Lorentz reciprocity \cite{harrington2001}. Only
Eq.~\eqref{eq:ch3.coupling-images-inner-product} on \(\mathcal{F}_{\mathrm{rad}}\) aligns with
Theorem~\ref{thm:ch2.lorentz-reciprocity-operator} \cite{stratton1941}.

\paragraph{Worked illustration (normalization versus orthogonality).}
Suppose \(\{\mathbf{f}_\alpha\}\) are far patterns orthogonal on \(S^{2}\) but rescaled arbitrarily in
amplitude. Angular orthogonality of patterns does not imply
Eq.~\eqref{eq:ch3.orthogonality-radiative-def} until flux normalization from Section~\ref{sec:ch35} fixes
\(\|\mathbf{f}_\alpha\|_{\mathrm{rad}}=1\) \cite{stratton1941,hansen1988}. One must not conflate \emph{geometric}
orthogonality of directions with \emph{Hilbert} orthogonality of export carriers; Subsection~\ref{sec:ch312}
separates the two before Section~\ref{sec:ch34} displays harmonic formulas \cite{harrington2001}.

\paragraph{Worked illustration (truncated harmonic fit).}
A finite spherical-harmonic fit to \(\mathcal{A}_\infty\) produces coefficients
\(\widehat{c}_\alpha\) that solve a least-squares problem on a measurement sphere, not necessarily the
Riesz coefficients Eq.~\eqref{eq:ch3.riesz-coefficient-representation} \cite{harrington2001,jackson1999}.
Parseval identity Eq.~\eqref{eq:ch3.parseval-radiative-preview} then holds only up to remainder
\(\|R_{\mathrm{trunc}}\|_{\mathrm{rad}}\) from Eq.~\eqref{eq:ch3.eigenfield-truncation-remainder}; treating
residual energy as ``uncaptured physics'' without checking \(\mathfrak{A}_{\mathrm{meas}}\) confuses
truncation with operator nullity \cite{stratton1941,harrington2001}.

\paragraph{Pairing hierarchy versus measurement maps.}
Three pairings compete in practice, and only one is the Chapter~\ref{ch:03} carrier for coefficients.
Volume product pairing \(\langle\cdot,\cdot\rangle_{\mathcal{U}}\) from
Eq.~\eqref{eq:ch2.field-energy-norm} controls local elliptic stability on bounded truncations
\cite{stratton1941}. Raw solution pairing on \(\mathcal{S}_\omega[\Jimp]\) includes reactive sheets
removed by \(\Pi_{\mathrm{rad}}\) \cite{harrington2001}. Radiative pairing
Eq.~\eqref{eq:ch3.radiative-inner-product-def} is the export-qualified overlap used in
Eqs.~\eqref{eq:ch3.coupling-images-inner-product} and~\eqref{eq:ch3.riesz-coefficient-representation}
\cite{stratton1941}. Observation-limited coupling Proposition~\ref{prop:ch2.observation-limited-coupling}
further replaces \(\langle\cdot,\cdot\rangle_{\mathrm{rad}}\) with a self-adjoint measurement map on
\(\mathcal{F}_{\mathrm{rad}}\); rank caps from Eq.~\eqref{eq:ch2.reconstruction-rank-cap} can zero
pairings that are nonzero in Eq.~\eqref{eq:ch3.reciprocity-symmetric-inner-product}
\cite{harrington2001,stratton1941}. Subsection~\ref{sec:ch312} fixes the middle rung; Sections~\ref{sec:ch35}
and~\ref{sec:ch310} connect the upper rung to calibrated meters \cite{stratton1941}.

\begin{table}[t]
\centering
\caption{Pairing routes on exported fields (Subsection~\ref{sec:ch312}). Rows are ordered by increasing
qualification; coefficient formulas require the radiative row.}
\label{tab:ch3.pairing-hierarchy}
\small
\begin{tabular}{@{}p{0.24\linewidth}p{0.34\linewidth}p{0.34\linewidth}@{}}
\toprule
Pairing & Domain & Coefficient admissibility \\
\midrule
\(\langle\cdot,\cdot\rangle_{\mathcal{U}}\) on \(\Omega_{R_{\min}}\) &
Bounded truncation &
Local stability only; blind to outgoing selection \\
\addlinespace
\(\langle \mathcal{S}_\omega[\cdot],\mathcal{S}_\omega[\cdot]\rangle_{\mathcal{U}}\) &
\(\mathcal{F}_\Omega\) &
Includes reactive overlap; violates export gate \\
\addlinespace
\(\langle\cdot,\cdot\rangle_{\mathrm{rad}}\) Eq.~\eqref{eq:ch3.radiative-inner-product-def} &
\(\mathcal{F}_{\mathrm{rad}}\) &
Riesz coefficients Eq.~\eqref{eq:ch3.riesz-coefficient-representation} \\
\addlinespace
\(\langle \mathcal{O}_{\mathrm{meas}}[\cdot],\mathcal{O}_{\mathrm{meas}}[\cdot]\rangle_{\mathrm{rad}}\) &
\(\operatorname{ran}\mathcal{O}_{\mathrm{meas}}\circ\mathcal{R}_\omega\) &
Observed coupling; rank-limited by aperture \\
\bottomrule
\end{tabular}
\end{table}

Table~\ref{tab:ch3.pairing-hierarchy} is not a norm diagnostic: Subsection~\ref{sec:ch311} already separated
shell energy from membership \cite{stratton1941}. The present table separates \emph{bilinear routes} that
one may interchange when fitting patterns or building reciprocity matrices \cite{harrington2001}. Row three is
the unique row compatible with Eq.~\eqref{eq:ch3.representation-admissibility} once flux normalizability is
interpreted as \(\|\mathbf{f}_\alpha\|_{\mathrm{rad}}=1\) \cite{stratton1941,hansen1988}.

\paragraph{Degenerate subspaces and Gram matrices.}
When Theorem~\ref{thm:ch2.radiation-spectral-decomposition} yields a \((2\ell+1)\)-fold degenerate family
\(\{\mathbf{u}_n\}_{n\in\mathcal{I}_\ell}\), orthogonality Eq.~\eqref{eq:ch3.orthogonality-radiative-def} holds
for any orthonormal basis of the block, but laboratory-fixed coordinates may produce off-diagonal Gram entries
\cite{jackson1999,harrington2001}. For a finite collection \(\{\mathbf{v}_p\}_{p=1}^{P}\subset\mathcal{F}_{\mathrm{rad}}\),
define
\begin{equation}
\label{eq:ch3.radiative-gram-matrix}
G_{pq}
=
\langle \mathbf{v}_p,\mathbf{v}_q\rangle_{\mathrm{rad}},
\qquad
p,q=1,\ldots,P,
\end{equation}
with \(G_{pq}=\delta_{pq}\) in any orthonormal eigenframe of Theorem
~\ref{thm:ch2.radiation-modal-orthogonality-coupling} \cite{stratton1941,harrington2001}. Equation
\eqref{eq:ch3.radiative-gram-matrix} is the finite-dimensional shadow of block coupling matrices from
Subsection~\ref{sec:ch282}; symmetry of \(G\) under reciprocal media follows from
Eq.~\eqref{eq:ch3.reciprocity-symmetric-inner-product} \cite{stratton1941}. Rotating the laboratory frame
rotates \(G\) by congruence while leaving export physics unchanged---the representation-theoretic repair of
that rotation is Section~\ref{sec:ch32} \cite{jackson1999,harrington2001}.

\paragraph{Dual space and bounded probes.}
Let \(\mathcal{F}_{\mathrm{rad}}^{*}\) denote the continuous dual with respect to
\(\langle\cdot,\cdot\rangle_{\mathrm{rad}}\). Each probe field \(\mathbf{g}_\omega\in\mathcal{F}_{\mathrm{rad}}\)
defines a measurement functional
\begin{equation}
\label{eq:ch3.probe-functional-dual}
\Lambda_{\mathbf{g}}[\mathbf{u}_\omega]
=
\langle \mathbf{u}_\omega,\mathbf{g}_\omega\rangle_{\mathrm{rad}},
\qquad
\mathbf{u}_\omega\in\mathcal{F}_{\mathrm{rad}},
\end{equation}
with \(\|\Lambda_{\mathbf{g}}\|=\|\mathbf{g}_\omega\|_{\mathrm{rad}}\) by Proposition
~\ref{prop:ch3.inner-product-hilbert-topology} \cite{stratton1941,harrington2001}. Equation
\eqref{eq:ch3.probe-functional-dual} is the abstract form of a calibrated receive channel: the probe
\(\mathbf{g}_\omega\) is a reference radiation image, not a raw voltage sample \cite{harrington2001}. Coefficient
functionals Eq.~\eqref{eq:ch3.riesz-coefficient-representation} are the special case
\(\mathbf{g}_\omega=\mathbf{f}_\alpha\) once normalization is fixed \cite{stratton1941,hansen1988}. Finite
apertures implement \(\Lambda_{\mathbf{g}}\) through composition with \(\mathcal{O}_{\mathrm{meas}}\), possibly
reducing rank without altering the definition of \(\langle\cdot,\cdot\rangle_{\mathrm{rad}}\) on exact states
\cite{harrington2001,stratton1941}.

Subsection~\ref{sec:ch312} has replaced the Subsection~\ref{sec:ch311} seminorm with a reciprocity-compatible
Hilbert geometry Eqs.~\eqref{eq:ch3.radiative-inner-product-def}--\eqref{eq:ch3.completeness-radiative-characterization},
linked it to Chapter~\ref{ch:02} coupling through Eq.~\eqref{eq:ch3.coupling-images-inner-product}, and
prepared coefficient extraction Eq.~\eqref{eq:ch3.riesz-coefficient-representation} for the continuous-spectrum
and rigged extensions of Subsections~\ref{sec:ch313}--\ref{sec:ch314} \cite{stratton1941,hansen1988,harrington2001}.


\subsection{Continuous spectra and generalized eigenfunctions}
\label{sec:ch313}
% TOC tag: [HOME][USE SS2.6.1]

Open domains radiate through a \emph{continuum} of outgoing directions long before anyone chooses
spherical-harmonic indices or Wigner matrices. Subsections~\ref{sec:ch311}--\ref{sec:ch312} made
\(\mathcal{F}_{\mathrm{rad}}\) a Hilbert space with inner product
Eq.~\eqref{eq:ch3.radiative-inner-product-def}; Subsection~\ref{sec:ch313} now asks how spectral
decomposition from Subsection~\ref{sec:ch261} reads when the eigen-index is continuous rather than
countable \cite{stratton1941,jackson1999,harrington2001}. Discrete cavity sums and exterior angular
integrals are not competing physical models---they are different spectral measures on different domain
classes Eq.~\eqref{eq:ch2.discrete-continuous-spectrum-preview} \cite{stratton1941}. The present
subsection supplies the Chapter~\ref{ch:03} operator-level account of continuous spectra and
generalized eigenfunctions on \(\mathcal{F}_{\mathrm{rad}}\); rigged completion and distributional
normalization are finalized in Subsection~\ref{sec:ch314} and Section~\ref{sec:ch35}
\cite{harrington2001,hansen1988}.

\paragraph{Standing imports from Subsection~\ref{sec:ch261}.}
Work at fixed \(\omega>0\) with outgoing closure
Eqs.~\eqref{eq:ch2.sommerfeld-outgoing-class}--\eqref{eq:ch2.outgoing-radiation-refinement}, radiation
eigenpairs Eq.~\eqref{eq:ch2.radiation-eigenpair-def}, modal projectors
Eq.~\eqref{eq:ch2.modal-projector-def}, spectral expansion Theorem
~\ref{thm:ch2.radiation-spectral-decomposition}, exterior continuous integral target
Eq.~\eqref{eq:ch2.exterior-continuous-spectral-integral}, far-pattern link Proposition
~\ref{prop:ch2.far-pattern-from-modes}, and biorthogonality
Eq.~\eqref{eq:ch2.radiation-biorthogonality} \cite{stratton1941,harrington2001,jackson1999}. Those
Chapter~\ref{ch:02} objects are cited here without re-display beyond the new Chapter~\ref{ch:03}
specializations below \cite{stratton1941}. Eigenpair construction, Fredholm solvability, and Green
kernels remain locked in Chapter~\ref{ch:02} \cite{harrington2001}.

\paragraph{Spectral index split.}
Let \(I\) denote the spectral index set for outgoing eigenfunctions on the active domain class.
Write
\begin{equation}
\label{eq:ch3.spectral-index-partition}
I
=
I_{\mathrm{d}}\ \cup\ I_{\mathrm{c}},
\qquad
I_{\mathrm{d}}\cap I_{\mathrm{c}}=\varnothing,
\end{equation}
with \(I_{\mathrm{d}}\) countable (cavity-like or angular-quantum sectors) and \(I_{\mathrm{c}}\) a
locally compact parameter space for continuous radiation modes \cite{stratton1941,jackson1999}.
Equation~\eqref{eq:ch3.spectral-index-partition} refines the preview
Eq.~\eqref{eq:ch2.discrete-continuous-spectrum-preview}: bounded cavities favor nonempty
\(I_{\mathrm{d}}\) with accumulating eigenvalues; exterior problems with compact sources favor
nonempty \(I_{\mathrm{c}}\) with direction labels \(\hat{\mathbf{k}}\in S^{2}\) as in
Eq.~\eqref{eq:ch2.exterior-angular-eigenmode} \cite{stratton1941,harrington2001}. Using only
\(I_{\mathrm{d}}\) on open domains truncates sidelobes; using only \(I_{\mathrm{c}}\) inside closed
resonators misstates storage modes \cite{jackson1999,stratton1941}.

\paragraph{Continuous spectral measure.}
On \(I_{\mathrm{c}}\), fix a positive Radon measure \(d\mu(\beta)\) compatible with the far-zone
geometry of Theorem~\ref{thm:ch2.far-field-approximation} \cite{stratton1941,jackson1999}. The
Chapter~\ref{ch:03} radiative measure is declared by
\begin{equation}
\label{eq:ch3.continuous-spectral-measure}
\int_{I_{\mathrm{c}}}
|c(\beta)[\mathfrak{D}]|^{2}\,d\mu(\beta)
\ \le\
\|\mathcal{R}_\omega[\mathfrak{D}]\|_{\mathrm{rad}}^{2},
\end{equation}
for admissible \(\mathfrak{D}\), with \(c(\beta)[\mathfrak{D}]\) the continuous coefficient density
defined below and \(\|\cdot\|_{\mathrm{rad}}\) from
Eq.~\eqref{eq:ch3.inner-product-seminorm-identity} \cite{stratton1941,harrington2001}. Inequality
\eqref{eq:ch3.continuous-spectral-measure} is the continuous Parseval companion to
Eq.~\eqref{eq:ch3.parseval-radiative-preview}: discrete modes consume energy on \(I_{\mathrm{d}}\);
the remainder is distributed across \(I_{\mathrm{c}}\) with weight \(d\mu\) \cite{stratton1941}.
Normalization of \(d\mu\) against flux meters is owned by Section~\ref{sec:ch35}; here \(d\mu\) is
fixed only up to a positive scaling absorbed into \(\mathbf{f}(\beta)\) conventions
\cite{harrington2001,hansen1988}.

\paragraph{Generalized radiative eigenfunctions.}
A \emph{generalized} outgoing eigenfunction \(\mathbf{f}(\beta)\) is a field pair indexed by
\(\beta\in I_{\mathrm{c}}\) that satisfies the radiation eigen-equation in the sense of distributions
on \(\mathcal{F}_{\mathrm{rad}}\):
\begin{equation}
\label{eq:ch3.generalized-radiative-eigenmode}
\mathcal{R}_\omega[\Jimp^{(\beta)}]
=
\lambda(\beta)\,\mathbf{f}(\beta),
\qquad
\mathbf{f}(\beta)\notin\mathcal{F}_{\mathrm{rad}}\ \text{ordinarily},
\end{equation}
with \(\Jimp^{(\beta)}\) a distributional source family and \(\lambda(\beta)\) the export gain
\cite{stratton1941,harrington2001}. Equation~\eqref{eq:ch3.generalized-radiative-eigenmode} upgrades
Eq.~\eqref{eq:ch2.radiation-eigenpair-def} to the continuous parameter: plane-wave modes
Eq.~\eqref{eq:ch2.exterior-angular-eigenmode} are the canonical example with
\(\beta=(\hat{\mathbf{k}},\sigma)\) \cite{jackson1999,stratton1941}. Ordinary
\(L^{2}\)-membership fails for \(\mathbf{f}(\beta)\) because transverse energy on every shell is
finite but the limit defining \(\|\cdot\|_{\mathrm{rad}}\) is saturated by delta-normalized direction
content \cite{harrington2001}. Subsection~\ref{sec:ch314} embeds such modes in a rigged triple
\(\Phi\subset\mathcal{F}_{\mathrm{rad}}\subset\Phi^{*}\); Subsection~\ref{sec:ch313} records the
spectral targets that triple must satisfy \cite{stratton1941}.

\paragraph{Continuous orthogonality relation.}
Generalized modes are orthogonal in the radiative pairing by a Dirac delta on \(I_{\mathrm{c}}\):
\begin{equation}
\label{eq:ch3.continuous-mode-orthogonality}
\big\langle \mathbf{f}(\beta),\,\mathbf{f}(\beta')\big\rangle_{\mathrm{rad}}
=
\delta(\beta-\beta')
\quad\text{in}\ d\mu,
\end{equation}
interpreted as a distributional identity tested against admissible coefficient functions
\cite{stratton1941,harrington2001}. Equation~\eqref{eq:ch3.continuous-mode-orthogonality} is the
continuous lift of discrete orthogonality imported from Theorem
~\ref{thm:ch2.radiation-modal-orthogonality-coupling}(i) with
\(\langle\cdot,\cdot\rangle_{\mathcal{U}}\) read as
\(\langle\cdot,\cdot\rangle_{\mathrm{rad}}\) \cite{stratton1941}. On reciprocal media, adjoint modes
coincide with \(\mathbf{f}(\beta)\) and biorthogonality Eq.~\eqref{eq:ch2.radiation-biorthogonality}
reduces to Eq.~\eqref{eq:ch3.continuous-mode-orthogonality} \cite{stratton1941,harrington2001}. The
\(\delta\)-relation is not a physical point source: it states that distinct direction labels carry
non-overlapping export content under \(\langle\cdot,\cdot\rangle_{\mathrm{rad}}\)
\cite{jackson1999,harrington2001}.

\paragraph{Continuous coefficient density.}
For admissible \(\mathfrak{D}\), define
\begin{equation}
\label{eq:ch3.continuous-coefficient-density}
c(\beta)[\mathfrak{D}]
=
\big\langle \mathcal{R}_\omega[\mathfrak{D}],\,\mathbf{f}(\beta)\big\rangle_{\mathrm{rad}},
\qquad
\beta\in I_{\mathrm{c}},
\end{equation}
when the pairing exists in the rigged sense of Subsection~\ref{sec:ch314}
\cite{stratton1941,harrington2001}. Equation~\eqref{eq:ch3.continuous-coefficient-density} extends
Eq.~\eqref{eq:ch3.riesz-coefficient-representation} from discrete \(\alpha\) to continuous
\(\beta\) \cite{stratton1941}. Linearity in \(\mathfrak{D}\) follows from
Eq.~\eqref{eq:ch2.radiation-linearity} and sesquilinearity of
Eq.~\eqref{eq:ch3.radiative-inner-product-def} \cite{stratton1941,harrington2001}. Physically,
\(c(\beta)[\mathfrak{D}]\) is the amplitude density of export in direction label \(\beta\) before
any sparse harmonic fit is applied \cite{jackson1999,harrington2001}.

\paragraph{Spectral resolution on \(\mathcal{F}_{\mathrm{rad}}\).}
Theorem~\ref{thm:ch2.radiation-spectral-decomposition} supplies discrete projectors
Eq.~\eqref{eq:ch2.modal-projector-def}. On exterior domains the continuous resolution reads
\begin{equation}
\label{eq:ch3.spectral-resolution-continuous}
\mathbf{u}_\omega
=
\sum_{\alpha\in I_{\mathrm{d}}}
\big\langle \mathbf{u}_\omega,\mathbf{f}_\alpha\big\rangle_{\mathrm{rad}}\,\mathbf{f}_\alpha
\;+\;
\int_{I_{\mathrm{c}}}
\big\langle \mathbf{u}_\omega,\mathbf{f}(\beta)\big\rangle_{\mathrm{rad}}\,
\mathbf{f}(\beta)\,d\mu(\beta),
\end{equation}
for \(\mathbf{u}_\omega\in\mathcal{F}_{\mathrm{rad}}\) in the closure of
\(\operatorname{span}\{\mathbf{f}_\alpha,\mathbf{f}(\beta)\}\) \cite{stratton1941,harrington2001}.
Equation~\eqref{eq:ch3.spectral-resolution-continuous} is the Chapter~\ref{ch:03} restatement of
Eq.~\eqref{eq:ch3.continuous-eigenfield-expansion} with the Subsection~\ref{sec:ch312} inner product
made explicit; it specializes Eq.~\eqref{eq:ch2.radiation-spectral-expansion} when
\(I_{\mathrm{c}}=\varnothing\) and Eq.~\eqref{eq:ch2.exterior-continuous-spectral-integral} when
\(I_{\mathrm{d}}\) is empty in the far-zone reduction \cite{stratton1941,jackson1999}. Convergence
is in \(\langle\cdot,\cdot\rangle_{\mathrm{rad}}\), not in raw \(\|\cdot\|_{\mathcal{U}}\) on near
shells \cite{harrington2001}.

\begin{figure}[t]
\centering
\ChOneFigBlock{%
\begin{tikzpicture}[ch1fig, node distance=7mm and 9mm]
  \node[ch1box] (id) {$I_{\mathrm{d}}$\\discrete};
  \node[ch1box, right=14mm of id] (ic) {$I_{\mathrm{c}}$\\continuous};
  \node[ch1box, below=10mm of id] (sum) {$\sum_{\alpha\in I_{\mathrm{d}}}$};
  \node[ch1box, below=10mm of ic] (int) {$\int_{I_{\mathrm{c}}}(\cdot)\,d\mu$};
  \node[ch1box, below=12mm of sum] (frad) {$\mathcal{F}_{\mathrm{rad}}$};
  \node[ch1box, below=12mm of int] (gen) {generalized $\mathbf{f}(\beta)$};
  \draw[ch1flow] (id) -- (sum);
  \draw[ch1flow] (ic) -- (int);
  \draw[ch1flow] (sum) -- (frad);
  \draw[ch1flow] (int) -- (gen);
  \draw[ch1flow] (gen) -- (frad);
\end{tikzpicture}%
}
\caption{Continuous versus discrete spectral routes for Subsection~\ref{sec:ch313}: domain class
selects active part of Eq.~\eqref{eq:ch3.spectral-index-partition}; continuous modes enter through
measure \(d\mu\) on \(I_{\mathrm{c}}\) and land in \(\mathcal{F}_{\mathrm{rad}}\) only after
rigged completion Subsection~\ref{sec:ch314}.}
\label{fig:ch3.continuous-discrete-spectral-split}
\end{figure}

Figure~\ref{fig:ch3.continuous-discrete-spectral-split} is a measure-theoretic complement to
Figure~\ref{fig:ch2.radiation-spectral-decomposition}: the latter orders operators; the present figure
orders \emph{index sets} \cite{stratton1941,harrington2001}. Continuous indices must not be read as the continuous
branch as optional decoration on a discrete sum: on exterior problems the integral carries the
angular content of Eq.~\eqref{eq:ch2.exterior-continuous-spectral-integral} \cite{jackson1999}.

\paragraph{Plane-wave continuum realization.}
One realization of Eq.~\eqref{eq:ch3.spectral-resolution-continuous} uses
\(\beta=(\hat{\mathbf{k}},\sigma)\) and
\begin{equation}
\label{eq:ch3.plane-wave-continuum-bridge}
\int_{S^{2}}\sum_{\sigma=1}^{2}
a_\sigma(\hat{\mathbf{k}})\,
\mathbf{u}_{\hat{\mathbf{k}},\sigma}(\mathbf{r})\,d\Omega(\hat{\mathbf{k}})
\ \longleftrightarrow\
\int_{I_{\mathrm{c}}}
c(\beta)[\mathfrak{D}]\,\mathbf{f}(\beta)\,d\mu(\beta),
\end{equation}
with \(d\Omega\) the solid-angle measure and \(\mathbf{u}_{\hat{\mathbf{k}},\sigma}\) from
Eq.~\eqref{eq:ch2.exterior-angular-eigenmode} \cite{stratton1941,jackson1999}. The arrow is
\emph{identification} under a fixed normalization map from Section~\ref{sec:ch35}, not a second
radiation operator \cite{harrington2001}. Angular-spectrum filters from Section~\ref{sec:ch29} act on
\(a_\sigma(\hat{\mathbf{k}})\) as multipliers on \(I_{\mathrm{c}}\); Section~\ref{sec:ch37} classifies
which multipliers preserve outgoing closure \cite{stratton1941,jackson1999}.

\paragraph{Operator spectral form with continuous part.}
Extending Theorem~\ref{thm:ch2.radiation-spectral-decomposition},
\begin{equation}
\label{eq:ch3.radiation-operator-continuous-spectral}
\mathcal{R}_\omega
=
\sum_{\alpha\in I_{\mathrm{d}}}
\lambda_\alpha\,P_\alpha\circ\mathcal{S}_\omega
\;+\;
\int_{I_{\mathrm{c}}}
\lambda(\beta)\,P(\beta)\circ\mathcal{S}_\omega\ d\mu(\beta),
\end{equation}
with \(P(\beta)\) the continuous projector associated with \(\mathbf{f}(\beta)\) and
\(P_\alpha\) the discrete projectors imported from Eq.~\eqref{eq:ch2.modal-projector-def}
\cite{stratton1941,harrington2001}. Equation~\eqref{eq:ch3.radiation-operator-continuous-spectral}
is formal until \(P(\beta)\) is defined in Subsection~\ref{sec:ch314}; Subsection~\ref{sec:ch313}
locks the \emph{target} form compatible with Eq.~\eqref{eq:ch2.radiation-operator-spectral-form}
\cite{stratton1941}. Range span Eq.~\eqref{eq:ch2.radiation-range-span} extends to
\(\overline{\operatorname{span}\{\mathbf{f}(\beta),\mathbf{f}_\alpha\}}\) in
\(\langle\cdot,\cdot\rangle_{\mathrm{rad}}\) \cite{harrington2001}.

\begin{proposition}[Continuous Parseval on \(\mathcal{F}_{\mathrm{rad}}\)]
\label{prop:ch3.continuous-parseval}
Let \(\mathfrak{D}\) be admissible with \(\mathcal{R}_\omega[\mathfrak{D}]\in\mathcal{F}_{\mathrm{rad}}\).
Suppose Eq.~\eqref{eq:ch3.spectral-resolution-continuous} converges in
\(\langle\cdot,\cdot\rangle_{\mathrm{rad}}\). Then
\begin{equation}
\label{eq:ch3.continuous-parseval-identity}
\|\mathcal{R}_\omega[\mathfrak{D}]\|_{\mathrm{rad}}^{2}
=
\sum_{\alpha\in I_{\mathrm{d}}}
\big|c_\alpha[\mathfrak{D}]\big|^{2}
\;+\;
\int_{I_{\mathrm{c}}}
\big|c(\beta)[\mathfrak{D}]\big|^{2}\,d\mu(\beta),
\end{equation}
with \(c_\alpha\) from Eq.~\eqref{eq:ch3.riesz-coefficient-representation} and
\(c(\beta)\) from Eq.~\eqref{eq:ch3.continuous-coefficient-density}
\cite{stratton1941,harrington2001}.
\end{proposition}
\begin{proof}[Proof sketch]
Apply Eq.~\eqref{eq:ch3.spectral-resolution-continuous} to
\(\mathbf{u}_\omega=\mathcal{R}_\omega[\mathfrak{D}]\), expand
\(\|\mathbf{u}_\omega\|_{\mathrm{rad}}^{2}=\langle\mathbf{u}_\omega,\mathbf{u}_\omega\rangle_{\mathrm{rad}}\),
and use Eqs.~\eqref{eq:ch3.continuous-mode-orthogonality} and~\eqref{eq:ch3.orthogonality-radiative-def}
on discrete sectors \cite{stratton1941,harrington2001}.
% PROOF: AUTHOR REVIEW
\end{proof}

\begin{theorem}[Continuous spectral decomposition on outgoing domains]
\label{thm:ch3.continuous-spectral-decomposition}
Let exterior outgoing closure and reciprocal media
Eq.~\eqref{eq:ch2.reciprocal-medium-class} hold. Let
\(\{\mathbf{f}(\beta)\}_{\beta\in I_{\mathrm{c}}}\) be generalized eigenfunctions satisfying
Eqs.~\eqref{eq:ch3.generalized-radiative-eigenmode} and~\eqref{eq:ch3.continuous-mode-orthogonality},
with discrete family \(\{\mathbf{f}_\alpha\}_{\alpha\in I_{\mathrm{d}}}\) from Theorem
~\ref{thm:ch2.radiation-spectral-decomposition}. Then every
\(\mathbf{u}_\omega\in\mathcal{F}_{\mathrm{rad}}\) admits
Eq.~\eqref{eq:ch3.spectral-resolution-continuous}, the operator form
Eq.~\eqref{eq:ch3.radiation-operator-continuous-spectral} holds on admissible \(\mathfrak{D}\), and
Proposition~\ref{prop:ch3.continuous-parseval} applies \cite{stratton1941,harrington2001,jackson1999}.
Far-zone patterns satisfy Proposition~\ref{prop:ch2.far-pattern-from-modes} with the continuous term
supplying the angular integral of Eq.~\eqref{eq:ch2.exterior-continuous-spectral-integral}
\cite{stratton1941,jackson1999}.
\end{theorem}
\begin{proof}[Proof sketch]
Discrete content is Theorem~\ref{thm:ch2.radiation-spectral-decomposition} read in
\(\langle\cdot,\cdot\rangle_{\mathrm{rad}}\) via Proposition~\ref{prop:ch3.reciprocity-orthonormal-frame}.
Continuous content follows from Eq.~\eqref{eq:ch2.exterior-continuous-spectral-integral} and
rigged-spectral closure imported from Subsection~\ref{sec:ch261} \cite{stratton1941,harrington2001}.
% PROOF: AUTHOR REVIEW
\end{proof}

\paragraph{Domain-class validity.}
Eqs.~\eqref{eq:ch3.spectral-index-partition}--\eqref{eq:ch3.radiation-operator-continuous-spectral}
require linear media, fixed \(\omega\), and outgoing \(\mathcal{B}_\Gamma\) on the stated domain class
\cite{stratton1941}. Cavity discrete spectra Eq.~\eqref{eq:ch2.cavity-discrete-radiation-spectrum}
dominate when Eq.~\eqref{eq:ch1.outer-boundary-classes} replaces Sommerfeld closure
\cite{harrington2001}. Loss moves \(\lambda(\beta)\) into the complex plane without promoting
evanescent sheets to \(I_{\mathrm{c}}\) \cite{stratton1941}. Accessible angular budget
Eq.~\eqref{eq:ch2.accessible-angular-budget} may restrict observable \(\beta\) ranges even when
Eq.~\eqref{eq:ch3.continuous-parseval-identity} holds mathematically \cite{harrington2001,jackson1999}.

\paragraph{Worked illustration (narrow beam as continuous superposition).}
A directive beam is often introduced as a finite \(\ell\)-sum; spectrally it is a weighted integral
over \(I_{\mathrm{c}}\) with \(c(\beta)[\mathfrak{D}]\) concentrated near a patch of \(S^{2}\)
\cite{jackson1999,harrington2001}. Replacing that integral by a truncated discrete sum produces
sidelobe artifacts that are representation truncation, not operator physics
\cite{stratton1941}. Eq.~\eqref{eq:ch3.continuous-coefficient-density} is the object to compare with
angular-spectrum samples from Section~\ref{sec:ch29} before harmonic labels are attached
\cite{jackson1999}.

\paragraph{Worked illustration (cavity versus exterior book confusion).}
Importing Eq.~\eqref{eq:ch2.cavity-discrete-radiation-spectrum} into an open exterior MoM code
forces \(I_{\mathrm{c}}=\varnothing\) by fiat and misstates wide-angle export
\cite{harrington2001,stratton1941}. Conversely, using only
Eq.~\eqref{eq:ch3.plane-wave-continuum-bridge} inside a closed resonator ignores accumulating
\(I_{\mathrm{d}}\) modes \cite{stratton1941}. Subsection~\ref{sec:ch313} requires the domain class to
select the active branch of Eq.~\eqref{eq:ch3.spectral-index-partition} before any coefficient report
\cite{harrington2001}.

\paragraph{Worked illustration (null content and continuous coefficients).}
If \(\mathfrak{D}\in\ker\mathcal{R}_\omega\), then \(c(\beta)[\mathfrak{D}]\equiv 0\) and
\(c_\alpha[\mathfrak{D}]\equiv 0\) for all exporting modes, yet \(\mathcal{S}_\omega[\mathfrak{D}]\)
may remain large in reactive subspaces \cite{stratton1941,harrington2001}. A continuous transform of
near-field data without \(\Pi_{\mathrm{rad}}\) can produce spurious \(c(\beta)\) that vanish once
Eq.~\eqref{eq:ch3.continuous-coefficient-density} is applied to \(\mathcal{R}_\omega[\mathfrak{D}]\)
\cite{harrington2001}. Null-space classification remains Section~\ref{sec:ch272} content
\cite{stratton1941}.

\begin{table}[t]
\centering
\caption{Spectral book selection by domain class (Subsection~\ref{sec:ch313}). Rows diagnose common
category errors; proofs are imported from Subsection~\ref{sec:ch261}.}
\label{tab:ch3.spectral-domain-class}
\small
\begin{tabular}{@{}p{0.26\linewidth}p{0.30\linewidth}p{0.36\linewidth}@{}}
\toprule
Domain class & Active spectral book & Failure if wrong book is used \\
\midrule
Bounded cavity, passive \(\mathcal{B}_\Gamma\) &
\(I_{\mathrm{d}}\) sum, Eq.~\eqref{eq:ch2.cavity-discrete-radiation-spectrum} &
Exterior plane-wave integrals misstate resonances \\
\addlinespace
Exterior, compact sources, outgoing &
\(I_{\mathrm{c}}\) integral, Eq.~\eqref{eq:ch2.exterior-continuous-spectral-integral} &
Finite harmonic sums truncate sidelobes \\
\addlinespace
Exterior + discrete symmetries (e.g.\ \(\ell\) sectors) &
Both \(I_{\mathrm{d}}\) and \(I_{\mathrm{c}}\) in Eq.~\eqref{eq:ch3.spectral-resolution-continuous} &
Confusing sparse \(\ell\) with absence of \(I_{\mathrm{c}}\) \\
\bottomrule
\end{tabular}
\end{table}

Table~\ref{tab:ch3.spectral-domain-class} is a domain classifier, not a repeat of
Table~\ref{tab:ch3.pairing-hierarchy}: pairing rows concern bilinear forms; spectral rows concern
\emph{which index set carries energy} \cite{stratton1941,harrington2001}.

\paragraph{Sampling continuous coefficients.}
Laboratories never measure \(c(\beta)\) on all of \(I_{\mathrm{c}}\) simultaneously; they sample on
finite direction grids or reconstruct \(a_\sigma(\hat{\mathbf{k}})\) from aperture data
\cite{harrington2001,jackson1999}. A finite sample set \(\{\beta_j\}_{j=1}^{J}\subset I_{\mathrm{c}}\)
induces an empirical Gram matrix
\begin{equation}
\label{eq:ch3.continuous-sample-gram}
\widehat{G}_{jk}
=
\big\langle \mathbf{f}(\beta_j),\,\mathbf{f}(\beta_k)\big\rangle_{\mathrm{rad}},
\qquad
j,k=1,\ldots,J,
\end{equation}
which approaches the \(\delta\)-orthogonality Eq.~\eqref{eq:ch3.continuous-mode-orthogonality} only as
sampling refines and aperture rank permits \cite{stratton1941,harrington2001}. Quadrature weights
\(w_j\) approximate the continuous Parseval integral by
\begin{equation}
\label{eq:ch3.continuous-parseval-quadrature}
\int_{I_{\mathrm{c}}}
|c(\beta)[\mathfrak{D}]|^{2}\,d\mu(\beta)
\ \approx\
\sum_{j=1}^{J} w_j\,\big|c(\beta_j)[\mathfrak{D}]\big|^{2},
\end{equation}
with error controlled by smoothness of \(c(\beta)\) and by accessible budget
Eq.~\eqref{eq:ch2.accessible-angular-budget} \cite{harrington2001,jackson1999}. Equations
\eqref{eq:ch3.continuous-sample-gram}--\eqref{eq:ch3.continuous-parseval-quadrature} connect
Subsection~\ref{sec:ch313} to numerical angular-spectrum codes in Section~\ref{sec:ch29} without
replacing the continuous operator targets of Subsection~\ref{sec:ch261} \cite{stratton1941}.

\paragraph{Discrete--continuous hybrid sectors.}
Rotationally symmetric environments often present a discrete \(\ell\) label while each sector still
carries a continuous direction measure inside degenerate subspaces from Subsection~\ref{sec:ch262}
\cite{jackson1999,harrington2001}. Symbolically,
\begin{equation}
\label{eq:ch3.hybrid-spectral-decomposition}
\mathcal{R}_\omega[\mathfrak{D}]
=
\sum_{\ell\in I_{\mathrm{d}}}
\int_{I_{\mathrm{c}}^{(\ell)}}
c^{(\ell)}(\beta)[\mathfrak{D}]\,\mathbf{f}^{(\ell)}(\beta)\,d\mu^{(\ell)}(\beta),
\end{equation}
with \(I_{\mathrm{c}}^{(\ell)}\) the direction fiber at fixed \(\ell\) and \(\mathbf{f}^{(\ell)}(\beta)\)
the corresponding generalized eigenfunction \cite{stratton1941,hansen1988}. Equation
\eqref{eq:ch3.hybrid-spectral-decomposition} is the spectral form of hybrid indices anticipated in
Eq.~\eqref{eq:ch3.spectral-index-alignment}; explicit \(\mathbf{f}^{(\ell)}\) are constructed in
Section~\ref{sec:ch34} \cite{hansen1988,stratton1941}. Confusing a single \(\ell\) block with a
one-dimensional subspace ignores the continuous fiber \(I_{\mathrm{c}}^{(\ell)}\) and misstates
cross-polar coupling within the block \cite{harrington2001,jackson1999}.

\paragraph{Forward link to rigged completion.}
Generalized modes \(\mathbf{f}(\beta)\) in Eq.~\eqref{eq:ch3.generalized-radiative-eigenmode} are not
generally elements of the Hilbert space \((\mathcal{F}_{\mathrm{rad}},\langle\cdot,\cdot\rangle_{\mathrm{rad}})\)
\cite{stratton1941,harrington2001}. Subsection~\ref{sec:ch314} embeds \(\mathcal{F}_{\mathrm{rad}}\) in a
rigged triple so that Eq.~\eqref{eq:ch3.continuous-coefficient-density},
Eq.~\eqref{eq:ch3.continuous-mode-orthogonality}, and projector integrals in
Eq.~\eqref{eq:ch3.radiation-operator-continuous-spectral} become well defined; Section~\ref{sec:ch35}
fixes \(d\mu\) and delta-normalization against Eq.~\eqref{eq:ch2.flux-on-radiative-branch}
\cite{harrington2001,hansen1988,stratton1941}.

Subsection~\ref{sec:ch313} has extended the Hilbert framework of Subsections~\ref{sec:ch311}--\ref{sec:ch312}
to continuous spectra: partition Eq.~\eqref{eq:ch3.spectral-index-partition}, measure
Eq.~\eqref{eq:ch3.continuous-spectral-measure}, generalized modes
Eqs.~\eqref{eq:ch3.generalized-radiative-eigenmode}--\eqref{eq:ch3.continuous-coefficient-density},
resolution Eq.~\eqref{eq:ch3.spectral-resolution-continuous}, and operator target
Eq.~\eqref{eq:ch3.radiation-operator-continuous-spectral}, importing Subsection~\ref{sec:ch261} by
reference only \cite{stratton1941,jackson1999,harrington2001}. Subsection~\ref{sec:ch314} supplies the
rigged Hilbert apparatus that makes the continuous book mathematically executable
\cite{stratton1941,hansen1988}.


\subsection{Rigged Hilbert spaces and distributional modes}
\label{sec:ch314}
% TOC tag: [HOME]

Plane-wave directions, delta-normalized outgoing modes, and spectral projectors appear in every
far-field code long before one can exhibit them as vectors in
\((\mathcal{F}_{\mathrm{rad}},\langle\cdot,\cdot\rangle_{\mathrm{rad}})\)
\cite{jackson1999,harrington2001,stratton1941}. The mismatch is not a modeling shortcut: continuous
angular content and formal spectral projectors from Subsection~\ref{sec:ch224} are legitimate only after
a \emph{rigged} embedding enlarges the dual domain while keeping the Hilbert core of
Subsections~\ref{sec:ch311}--\ref{sec:ch312} intact \cite{kato1995,stratton1941}. Subsection~\ref{sec:ch314}
supplies that embedding for Volume~I: a Gelfand triple around \(\mathcal{F}_{\mathrm{rad}}\), distributional
extensions of radiation pairings, and the analytic license for continuous projectors anticipated in
Subsection~\ref{sec:ch313} \cite{harrington2001,hansen1988}.

\paragraph{Standing imports and obligations.}
Subsection~\ref{sec:ch314} closes the formal-status note of Subsection~\ref{sec:ch224} on
Eqs.~\eqref{eq:ch2.spectral-projectors}--\eqref{eq:ch2.spectral-projector-algebra} without re-displaying
those equations \cite{stratton1941,kato1995}. Distributional source extensions from
Subsection~\ref{sec:ch211} through Eq.~\eqref{eq:ch2.distributional-source-extension} and dual pairings
Eq.~\eqref{eq:ch2.dual-source-pairing} are cited as infrastructure; sheet and volume source proofs are not
reopened \cite{stratton1941}. Continuous spectral targets
Eqs.~\eqref{eq:ch3.generalized-radiative-eigenmode}--\eqref{eq:ch3.radiation-operator-continuous-spectral}
from Subsection~\ref{sec:ch313} are presupposed; here they receive a well-defined pairing domain
\cite{harrington2001,hansen1988}. Explicit delta-normalization against flux meters is deferred to
Section~\ref{sec:ch35}; Subsection~\ref{sec:ch314} fixes the functional-analytic stage on which that
normalization acts \cite{stratton1941,poynting1884}.

\paragraph{Test space and smooth outgoing probes.}
Let \(\Phi\subset\mathcal{F}_{\mathrm{rad}}\) be the dense test subspace of smooth Maxwell pairs on
exterior domains with rapid transverse decay on shells \(\Gamma_R\) as \(R\to\infty\), outgoing in the sense
of Eq.~\eqref{eq:ch3.outgoing-inheritance-criterion}, and compatible with patchwise constitutive scalings
\cite{stratton1941,harrington2001}. Elements of \(\Phi\) are the states against which distributional modes
are tested; numerically, \(\Phi\) is approximated by finite-element or method-of-moments shape functions with
absorbing closure on \(\Gamma_R\) \cite{harrington2001}. Topology on \(\Phi\) is the subspace topology
inherited from \(\|\cdot\|_{\mathrm{rad}}\); \(\Phi\) is not closed in
\((\mathcal{F}_{\mathrm{rad}},\langle\cdot,\cdot\rangle_{\mathrm{rad}})\) because limits of smooth tests
can develop singular far-zone angular content \cite{kato1995,stratton1941}.

\paragraph{Gelfand triple around \(\mathcal{F}_{\mathrm{rad}}\).}
The radiative rigged structure is
\begin{equation}
\label{eq:ch3.gelfand-triple-rad}
\Phi
\ \subset\
\mathcal{F}_{\mathrm{rad}}
\ \subset\
\Phi^{*},
\end{equation}
with \(\Phi^{*}\) the continuous dual of \(\Phi\) under the pairing induced by
Eq.~\eqref{eq:ch3.radiative-inner-product-def} \cite{kato1995,stratton1941}. Equation
\eqref{eq:ch3.gelfand-triple-rad} is the Chapter~\ref{ch:03} closure of the preview in
Subsection~\ref{sec:ch224}: formal projectors act on \(\widehat{\mathcal{F}}_\Omega\) only after
\(\mathcal{F}_{\mathrm{rad}}\) is embedded in \(\Phi^{*}\) where delta-labeled modes live
\cite{stratton1941,harrington2001}. Inclusion \(\mathcal{F}_{\mathrm{rad}}\subset\Phi^{*}\) is via the
canonical embedding
\begin{equation}
\label{eq:ch3.canonical-embedding-rad-dual}
\iota:\mathcal{F}_{\mathrm{rad}}\to\Phi^{*},
\qquad
\iota[\mathbf{u}](\boldsymbol{\phi})
=
\langle \mathbf{u},\boldsymbol{\phi}\rangle_{\mathrm{rad}},
\quad
\boldsymbol{\phi}\in\Phi,
\end{equation}
which is injective and isometric on \(\mathcal{F}_{\mathrm{rad}}\) \cite{kato1995,stratton1941}. Generalized
eigenfunctions \(\mathbf{f}(\beta)\) from Eq.~\eqref{eq:ch3.generalized-radiative-eigenmode} are
\emph{elements of \(\Phi^{*}\)}, not of \(\mathcal{F}_{\mathrm{rad}}\) \cite{harrington2001,hansen1988}.

\begin{figure}[t]
\centering
\ChOneFigBlock{%
\begin{tikzpicture}[ch1fig, node distance=8mm and 10mm]
  \node[ch1box] (phi) {$\Phi$\\tests};
  \node[ch1box, right=16mm of phi] (frad) {$\mathcal{F}_{\mathrm{rad}}$\\Hilbert core};
  \node[ch1box, right=16mm of frad] (dual) {$\Phi^{*}$\\distributional};
  \node[ch1box, below=11mm of frad] (gen) {$\mathbf{f}(\beta)$, $\delta_\beta$};
  \node[ch1box, below=11mm of phi] (proj) {formal $P(\beta)$, $\Pi_{\mathrm{prop}}$};
  \draw[ch1flow] (phi) -- node[above,font=\scriptsize] {dense} (frad);
  \draw[ch1flow] (frad) -- node[above,font=\scriptsize] {$\iota$} (dual);
  \draw[ch1flow] (gen) -- (dual);
  \draw[ch1flow] (proj) -- (dual);
\end{tikzpicture}%
}
\caption{Gelfand triple Eq.~\eqref{eq:ch3.gelfand-triple-rad} for Subsection~\ref{sec:ch314}: smooth tests
\(\Phi\) anchor the Hilbert core \(\mathcal{F}_{\mathrm{rad}}\); distributional modes and formal projectors
from Subsection~\ref{sec:ch224} are lawful in \(\Phi^{*}\).}
\label{fig:ch3.gelfand-triple-radiative}
\end{figure}

Figure~\ref{fig:ch3.gelfand-triple-radiative} is the functional-analytic capstone of
Section~\ref{sec:ch31}: unlike Figure~\ref{fig:ch3.radiative-hilbert-ladder} (norms) or
Figure~\ref{fig:ch3.continuous-discrete-spectral-split} (index sets), the present diagram tracks
\emph{topological layers} \cite{kato1995,stratton1941}.

\paragraph{Distributional extension of the radiative pairing.}
For \(\mathbf{U}\in\Phi^{*}\) and \(\boldsymbol{\phi}\in\Phi\), define the extended pairing
\begin{equation}
\label{eq:ch3.distributional-pairing-extension}
\mathbf{U}(\boldsymbol{\phi})
=
\langle \mathbf{U},\boldsymbol{\phi}\rangle_{\mathrm{rig}}
\ \equiv\
\langle \mathbf{U},\boldsymbol{\phi}\rangle_{\mathrm{rad}}
\quad\text{when}\ \mathbf{U}=\iota[\mathbf{u}],\ \mathbf{u}\in\mathcal{F}_{\mathrm{rad}},
\end{equation}
and by continuity on \(\Phi\) for general \(\mathbf{U}\in\Phi^{*}\) \cite{kato1995,stratton1941}. Equation
\eqref{eq:ch3.distributional-pairing-extension} extends
Eq.~\eqref{eq:ch3.radiative-inner-product-def} without altering values on the Hilbert core
\cite{stratton1941}. Continuous coefficients Eq.~\eqref{eq:ch3.continuous-coefficient-density} become
\begin{equation}
\label{eq:ch3.continuous-coefficient-rigged}
c(\beta)[\mathfrak{D}]
=
\big\langle \iota[\mathcal{R}_\omega[\mathfrak{D}]],\,\mathbf{f}(\beta)\big\rangle_{\mathrm{rig}},
\end{equation}
with \(\mathbf{f}(\beta)\in\Phi^{*}\) \cite{harrington2001,hansen1988}. The right-hand side is the
rigorous form of the heuristic ``project onto a delta beam'' instruction in angular-spectrum codes
\cite{jackson1999,stratton1941}.

\paragraph{Delta-normalized directionals as dual elements.}
A direction label \(\beta=(\hat{\mathbf{k}},\sigma)\) induces a distributional mode
\(\delta_\beta\in\Phi^{*}\) characterized by
\begin{equation}
\label{eq:ch3.delta-directional-mode}
\delta_\beta(\boldsymbol{\phi})
=
\lim_{\epsilon\to 0}\,
\langle \mathbf{u}_{\beta,\epsilon},\boldsymbol{\phi}\rangle_{\mathrm{rad}},
\end{equation}
for a family \(\mathbf{u}_{\beta,\epsilon}\in\Phi\) concentrating on \(\beta\) as \(\epsilon\to 0\)
\cite{stratton1941,jackson1999}. Equation~\eqref{eq:ch3.delta-directional-mode} is the distributional
content of Eq.~\eqref{eq:ch3.continuous-mode-orthogonality}: orthogonality of distinct \(\delta_\beta\)
is tested on \(\Phi\), not by requiring \(\delta_\beta\in\mathcal{F}_{\mathrm{rad}}\)
\cite{harrington2001,kato1995}. Flux-normalized \(\delta_\beta\) used in power budgets are fixed in
Section~\ref{sec:ch35}; here only existence in \(\Phi^{*}\) is declared \cite{stratton1941,poynting1884}.

\paragraph{Rigged continuous projectors.}
For each \(\beta\in I_{\mathrm{c}}\), define \(P(\beta):\Phi^{*}\to\Phi^{*}\) weakly by
\begin{equation}
\label{eq:ch3.continuous-projector-rigged}
\big\langle P(\beta)\mathbf{U},\boldsymbol{\phi}\big\rangle_{\mathrm{rig}}
=
\big\langle \mathbf{U},\boldsymbol{\phi}_\beta\rangle_{\mathrm{rig}}\,
\delta(\beta-\beta_\phi),
\end{equation}
with \(\boldsymbol{\phi}_\beta\in\Phi\) a test localized on \(\beta\) and \(\delta(\beta-\beta_\phi)\) the
continuous orthogonality kernel Eq.~\eqref{eq:ch3.continuous-mode-orthogonality} \cite{stratton1941,hansen1988}.
Equation~\eqref{eq:ch3.continuous-projector-rigged} is the Subsection~\ref{sec:ch314} realization of the
formal integral in Eq.~\eqref{eq:ch3.radiation-operator-continuous-spectral}; discrete projectors imported
from Eq.~\eqref{eq:ch2.modal-projector-def} extend similarly on finite \(I_{\mathrm{d}}\)
\cite{stratton1941,harrington2001}. Composition with \(\mathcal{S}_\omega\) is defined on admissible
\(\mathfrak{D}\) when \(\mathcal{S}_\omega[\mathfrak{D}]\in\Phi^{*}\) through
Eq.~\eqref{eq:ch3.distributional-pairing-extension} \cite{stratton1941}.

\paragraph{Legitimization of formal spectral projectors.}
Subsection~\ref{sec:ch224} marked Eqs.~\eqref{eq:ch2.spectral-projectors}--\eqref{eq:ch2.spectral-projector-algebra}
as formal on \(\widehat{\mathcal{F}}_\Omega\). On the radiative rigged domain,
\begin{equation}
\label{eq:ch3.formal-projector-rigged-closure}
\Pi_{\mathrm{prop}}+\Pi_{\mathrm{ev}}+\Pi_{\mathrm{non}}
=
\mathrm{id}
\quad\text{on}\ \iota[\mathcal{F}_{\mathrm{rad}}]^{\mathrm{cl}}\subset\Phi^{*},
\end{equation}
and idempotency \(\Pi_{\mathrm{prop}}^{2}=\Pi_{\mathrm{prop}}\) holds in the weak topology of
\(\Phi^{*}\) when tested against \(\Phi\) \cite{stratton1941,kato1995}. Equation
\eqref{eq:ch3.formal-projector-rigged-closure} does not extend projectors to arbitrary
\(\mathcal{S}_\omega[\Jimp]\) without outgoing and regime qualification; it certifies that the algebra
used on \(\mathcal{R}_\omega[\mathfrak{D}]\) in Chapter~\ref{ch:02} is consistent with the
Subsection~\ref{sec:ch312} inner product once embedded in \(\Phi^{*}\) \cite{harrington2001,stratton1941}.
Radiation projector \(\Pi_{\mathrm{rad}}=\mathcal{O}_{\mathrm{out}}\circ\Pi_{\mathrm{prop}}\) inherits the
same rigged status through Eq.~\eqref{eq:ch2.rad-projection-composition} \cite{stratton1941}.

\paragraph{Extension of \(\mathcal{R}_\omega\) to distributional sources.}
Let \(\mathcal{J}_{\mathrm{rig}}\) denote impressed data \(\mathfrak{D}\) for which
\(\mathcal{S}_\omega[\mathfrak{D}]\in\Phi^{*}\) under Eq.~\eqref{eq:ch2.distributional-source-extension}
\cite{stratton1941}. Define
\begin{equation}
\label{eq:ch3.radiation-map-rigged-extension}
\mathcal{R}_\omega[\mathfrak{D}]
=
\Pi_{\mathrm{rad}}\big(\mathcal{S}_\omega[\mathfrak{D}]\big)
\ \in\
\iota[\mathcal{F}_{\mathrm{rad}}]^{\mathrm{cl}}\subset\Phi^{*},
\qquad
\mathfrak{D}\in\mathcal{J}_{\mathrm{rig}},
\end{equation}
with \(\Pi_{\mathrm{rad}}\) acting in the weak topology \cite{stratton1941,harrington2001}. For
\(\mathfrak{D}\in\mathcal{J}_\Omega\) with \(\mathcal{R}_\omega[\mathfrak{D}]\in\mathcal{F}_{\mathrm{rad}}\),
Eq.~\eqref{eq:ch3.radiation-map-rigged-extension} reduces to
Eq.~\eqref{eq:ch3.ran-embedding-rad} \cite{stratton1941}. Sheet currents and compactly supported smooth
\(\Jimp\) lie in \(\mathcal{J}_{\mathrm{rig}}\) by Subsection~\ref{sec:ch211}; singular self-terms remain
regulated by Subsection~\ref{sec:ch232} without reopening here \cite{harrington2001,stratton1941}.

\begin{proposition}[Rigged extension of radiative pairings]
\label{prop:ch3.rigged-pairing-extension}
The embedding Eq.~\eqref{eq:ch3.canonical-embedding-rad-dual} is isometric on
\(\mathcal{F}_{\mathrm{rad}}\). For every \(\mathbf{U}\in\Phi^{*}\) there exists a net
\(\{\mathbf{u}_\alpha\}\subset\mathcal{F}_{\mathrm{rad}}\) such that
\(\iota[\mathbf{u}_\alpha]\to\mathbf{U}\) weakly in \(\Phi^{*}\), and
Eq.~\eqref{eq:ch3.distributional-pairing-extension} agrees with
Eq.~\eqref{eq:ch3.radiative-inner-product-def} on \(\Phi\cap\mathcal{F}_{\mathrm{rad}}\)
\cite{kato1995,stratton1941,harrington2001}.
\end{proposition}
\begin{proof}[Proof sketch]
Density of \(\Phi\) in \(\mathcal{F}_{\mathrm{rad}}\) and completeness
Eq.~\eqref{eq:ch3.completeness-radiative-characterization} supply the approximating net; weak extension
follows from continuity of \(\langle\cdot,\boldsymbol{\phi}\rangle_{\mathrm{rad}}\) on \(\Phi\)
\cite{kato1995,stratton1941}.
% PROOF: AUTHOR REVIEW
\end{proof}

\begin{theorem}[Rigged closure of continuous spectral decomposition]
\label{thm:ch3.rigged-spectral-closure}
Let exterior outgoing closure and reciprocal media hold. Then Theorem
~\ref{thm:ch3.continuous-spectral-decomposition} holds with
\(\mathbf{f}(\beta),\delta_\beta\in\Phi^{*}\), coefficients given by
Eq.~\eqref{eq:ch3.continuous-coefficient-rigged}, projectors by
Eq.~\eqref{eq:ch3.continuous-projector-rigged}, and formal spectral algebra by
Eq.~\eqref{eq:ch3.formal-projector-rigged-closure} tested on \(\Phi\)
\cite{stratton1941,hansen1988,harrington2001,kato1995}.
\end{theorem}
\begin{proof}[Proof sketch]
Combine Proposition~\ref{prop:ch3.rigged-pairing-extension} with Theorem
~\ref{thm:ch3.continuous-spectral-decomposition} and the formal-status closure demanded by
Subsection~\ref{sec:ch224} \cite{stratton1941,kato1995,harrington2001}.
% PROOF: AUTHOR REVIEW
\end{proof}

\paragraph{Validity domain.}
Eqs.~\eqref{eq:ch3.gelfand-triple-rad}--\eqref{eq:ch3.radiation-map-rigged-extension} require linear
media, fixed \(\omega\), and outgoing \(\mathcal{B}_\Gamma\) on the declared exterior class
\cite{stratton1941}. Nonlinear or time-varying regimes break weak closure of
Eq.~\eqref{eq:ch3.formal-projector-rigged-closure} \cite{stratton1941}. Lossy media are admissible when
passivity is maintained and tests \(\Phi\) remain outgoing \cite{harrington2001}. Rigged extensions do not
remove \(\ker\mathcal{R}_\omega\) or measurement rank limits Eq.~\eqref{eq:ch2.reconstruction-rank-cap}
\cite{stratton1941,harrington2001}.

\paragraph{Worked illustration (MoM test functions as \(\Phi\)).}
A RWG basis on \(\Gamma_R\) with outgoing absorbing boundary produces coefficients in a finite-dimensional
subspace \(\Phi_N\subset\Phi\) \cite{harrington2001,stratton1941}. Passing \(N\to\infty\) is weak
convergence in \(\Phi^{*}\), not automatic convergence in \(\|\cdot\|_{\mathrm{rad}}\) on reactive
components \cite{stratton1941}. Interpreting MoM weights as Riesz coefficients
Eq.~\eqref{eq:ch3.riesz-coefficient-representation} is valid only after \(\Pi_{\mathrm{rad}}\) has been
applied and the pairing is read in \(\Phi^{*}\) \cite{harrington2001}.

\paragraph{Worked illustration (delta probe in a compact range).}
A narrow plane-wave probe approximates \(\delta_\beta\) through a finite beam \(\mathbf{u}_{\beta,\epsilon}\in\Phi\)
\cite{jackson1999,harrington2001}. Coefficient extraction via
Eq.~\eqref{eq:ch3.continuous-coefficient-rigged} matches calibrated receive coupling when reciprocity
Eq.~\eqref{eq:ch3.reciprocity-symmetric-inner-product} is preserved; replacing the rigged pairing by a
near-field \(\|\cdot\|_{\mathcal{U}}\) scan violates Eq.~\eqref{eq:ch3.distributional-pairing-extension}
\cite{stratton1941,harrington2001}.

\paragraph{Worked illustration (formal \(\Pi_{\mathrm{prop}}\) on a finite superposition).}
For a finite superposition of propagating plane waves with outgoing \(k_z\), Chapter~\ref{ch:02} applies
\(\Pi_{\mathrm{prop}}\) before \(\mathcal{O}_{\mathrm{out}}\) \cite{stratton1941}. Subsection~\ref{sec:ch314}
states that this composition is rigorous because the superposition lies in \(\mathcal{F}_{\mathrm{rad}}\subset\Phi^{*}\)
and tests in \(\Phi\) separate propagating from evanescent sheets through
Eq.~\eqref{eq:ch3.formal-projector-rigged-closure} \cite{harrington2001,kato1995}. The illustration fails if
evanescent content is added without \(\Pi_{\mathrm{ev}}\) control: formal algebra does not override
Eq.~\eqref{eq:ch2.regime-operator-screen} \cite{stratton1941}.

\begin{table}[t]
\centering
\caption{Topology cheat sheet for Section~\ref{sec:ch31} (Subsection~\ref{sec:ch314} row is new). Earlier
rows are cited by reference only; the table diagnoses \emph{convergence} questions, not pairing or spectral
books.}
\label{tab:ch3.topology-convergence}
\small
\begin{tabular}{@{}p{0.22\linewidth}p{0.30\linewidth}p{0.38\linewidth}@{}}
\toprule
Layer & Convergence tested & Typical laboratory question \\
\midrule
\(\|\cdot\|_{\mathrm{rad}}\) on \(\mathcal{F}_{\mathrm{rad}}\) &
Subsection~\ref{sec:ch311} &
Does shell energy stabilize as \(R\) grows? \\
\addlinespace
\(\langle\cdot,\cdot\rangle_{\mathrm{rad}}\) &
Subsection~\ref{sec:ch312} &
Do reciprocity and coupling matrices match export? \\
\addlinespace
Measure \(d\mu\) on \(I_{\mathrm{c}}\) &
Subsection~\ref{sec:ch313} &
Is angular content continuous or discrete? \\
\addlinespace
Weak topology on \(\Phi^{*}\) &
Subsection~\ref{sec:ch314} &
Do delta beams and formal projectors make sense? \\
\bottomrule
\end{tabular}
\end{table}

Table~\ref{tab:ch3.topology-convergence} closes the analytic ladder of Section~\ref{sec:ch31}: each row answers
a different failure mode in radiation metrology \cite{stratton1941,harrington2001}.

\paragraph{Weak convergence and numerical limits.}
Weak convergence in \(\Phi^{*}\) is the correct notion for refining direction grids and aperture nets:
\begin{equation}
\label{eq:ch3.weak-star-convergence-rigged}
\mathbf{U}_n\to \mathbf{U}\ \text{weakly in }\Phi^{*}
\quad\Longleftrightarrow\quad
\mathbf{U}_n(\boldsymbol{\phi})\to \mathbf{U}(\boldsymbol{\phi})
\ \ \forall\,\boldsymbol{\phi}\in\Phi.
\end{equation}
Criterion~\eqref{eq:ch3.weak-star-convergence-rigged} differs from Cauchy convergence
Eq.~\eqref{eq:ch3.cauchy-radiative-convergence}: a sequence may converge weakly while failing to converge
in \(\|\cdot\|_{\mathrm{rad}}\) because reactive content oscillates on bounded shells
\cite{kato1995,harrington2001,stratton1941}. Numerical codes that monitor only \(\|\cdot\|_{\mathcal{U}}\) on
\(\Gamma_R\) therefore conflate weak and norm convergence \cite{harrington2001}. Export observables tied to
Eq.~\eqref{eq:ch3.continuous-coefficient-rigged} are stable under
Eq.~\eqref{eq:ch3.weak-star-convergence-rigged} when probes \(\boldsymbol{\phi}\) are fixed outgoing tests
\cite{stratton1941,jackson1999}.

\paragraph{Continuity of the embedding chain.}
The inclusions in Eq.~\eqref{eq:ch3.gelfand-triple-rad} are continuous with respect to their native
topologies: \(\Phi\hookrightarrow\mathcal{F}_{\mathrm{rad}}\) is isometric, and
\(\iota:\mathcal{F}_{\mathrm{rad}}\to\Phi^{*}\) is continuous when \(\Phi^{*}\) carries the weak topology
\cite{kato1995,stratton1941}. Bounded linear functionals on \(\mathcal{F}_{\mathrm{rad}}\) therefore extend
uniquely to \(\Phi^{*}\) when they are weakly continuous on \(\Phi\) \cite{kato1995}. This extension
principle is what permits coefficient functionals Eq.~\eqref{eq:ch3.riesz-coefficient-representation} to act
on distributional images Eq.~\eqref{eq:ch3.radiation-map-rigged-extension} without redefining the Hilbert
inner product on the core \cite{stratton1941,harrington2001}.

\paragraph{Adjoints in the rigged setting.}
When a linear operator \(A\) on \(\mathcal{F}_{\mathrm{rad}}\) is bounded in
\(\langle\cdot,\cdot\rangle_{\mathrm{rad}}\), its Hilbert adjoint \(A^{\dagger}\) satisfies
\(\langle A\mathbf{u},\mathbf{v}\rangle_{\mathrm{rad}}=\langle\mathbf{u},A^{\dagger}\mathbf{v}\rangle_{\mathrm{rad}}\)
\cite{kato1995,stratton1941}. In \(\Phi^{*}\), the rigged adjoint \(A^{\times}\) is defined by
\begin{equation}
\label{eq:ch3.rigged-adjoint-definition}
\langle A^{\times}\mathbf{U},\boldsymbol{\phi}\rangle_{\mathrm{rig}}
=
\langle \mathbf{U},A\boldsymbol{\phi}\rangle_{\mathrm{rig}},
\qquad
\mathbf{U}\in\Phi^{*},\ \boldsymbol{\phi}\in\Phi,
\end{equation}
whenever \(A\boldsymbol{\phi}\in\Phi\) \cite{kato1995,stratton1941}. Equation
\eqref{eq:ch3.rigged-adjoint-definition} is the template for radiation operators whose formal adjoints
appear in Subsection~\ref{sec:ch281} and for measurement maps composed with
Eq.~\eqref{eq:ch3.probe-functional-dual} \cite{stratton1941,harrington2001}. Self-adjointness of
\(\mathcal{R}_\omega^{\dagger}\mathcal{R}_\omega\) on \(\mathcal{J}_{\mathrm{rig}}\) is not claimed here;
only the pairing convention is fixed so later sections cite adjoints without ambiguity
\cite{harrington2001,kato1995}.

\paragraph{Compatibility with distributional sources.}
Impressed data \(\mathfrak{D}\) in Eq.~\eqref{eq:ch2.distributional-source-extension} pair with test fields
through Eq.~\eqref{eq:ch2.dual-source-pairing} \cite{stratton1941}. The rigged extension identifies when
those pairings descend to radiation images:
\begin{equation}
\label{eq:ch3.source-to-field-rigged-chain}
\mathfrak{D}\in\mathcal{J}_{\mathrm{rig}}
\ \Longrightarrow\
\mathcal{S}_\omega[\mathfrak{D}]\in\Phi^{*}
\ \Longrightarrow\
\mathcal{R}_\omega[\mathfrak{D}]\in\iota[\mathcal{F}_{\mathrm{rad}}]^{\mathrm{cl}},
\end{equation}
with each implication using the outgoing and regime gates already locked in Chapter~\ref{ch:02}
\cite{stratton1941,harrington2001}. Equation~\eqref{eq:ch3.source-to-field-rigged-chain} is the
distributional reading of Eq.~\eqref{eq:ch3.export-pipeline-preview}: sources may be rough, but export
coefficients are extracted only after \(\Pi_{\mathrm{rad}}\) and only through
Eq.~\eqref{eq:ch3.distributional-pairing-extension} \cite{stratton1941}. Singular self-terms at
\(\mathbf{r}\in\operatorname{supp}\mathfrak{D}\) remain regulated by Subsection~\ref{sec:ch232}; rigging
does not remove the need for that regularization \cite{harrington2001,stratton1941}.

\paragraph{Worked illustration (angular-spectrum multipliers).}
Angular-spectrum filters from Section~\ref{sec:ch29} multiply coefficient densities \(c(\beta)\) on
\(I_{\mathrm{c}}\) \cite{jackson1999,stratton1941}. In \(\Phi^{*}\), a multiplier \(M\) acts weakly by
\(\langle M\mathbf{U},\boldsymbol{\phi}\rangle_{\mathrm{rig}}=\langle\mathbf{U},M^{\times}\boldsymbol{\phi}\rangle_{\mathrm{rig}}\)
when \(M^{\times}\boldsymbol{\phi}\in\Phi\) \cite{kato1995,harrington2001}. A sharp spectral mask therefore
changes export while leaving reactive content untouched only if \(\Pi_{\mathrm{rad}}\) has already been
applied before masking \cite{stratton1941}. Applying masks to \(\mathcal{S}_\omega[\mathfrak{D}]\) without
outgoing selection violates Eq.~\eqref{eq:ch3.source-to-field-rigged-chain} \cite{harrington2001}.

\paragraph{Worked illustration (sheet currents as dual sources).}
Surface equivalents from Eq.~\eqref{eq:ch2.distributional-source-extension} lie in \(\mathcal{J}_{\mathrm{rig}}\)
when tangential traces are in the declared trace class of Subsection~\ref{sec:ch211}
\cite{stratton1941,harrington2001}. Huygens routes from Subsection~\ref{sec:ch251} produce the same
\(\mathcal{R}_\omega[\mathfrak{D}]\in\iota[\mathcal{F}_{\mathrm{rad}}]^{\mathrm{cl}}\) as volume routes when
Eq.~\eqref{eq:ch3.route-independence-of-expansion} holds \cite{stratton1941}. The rigged viewpoint explains
why equivalent-surface codes remain consistent: they change the source representative in \(\mathcal{J}_{\mathrm{rig}}\),
not the export image in \(\Phi^{*}\) \cite{harrington2001,stratton1941}.

\paragraph{Spectral completion and \(\widehat{\mathcal{F}}_\Omega\).}
Subsection~\ref{sec:ch224} used the spectral completion \(\widehat{\mathcal{F}}_\Omega\) as the formal domain of
projector algebra \cite{stratton1941}. Subsection~\ref{sec:ch314} identifies the radiative slice of that completion
with the weak closure
\begin{equation}
\label{eq:ch3.spectral-completion-radiative}
\iota[\mathcal{F}_{\mathrm{rad}}]^{\mathrm{cl}}
=
\left\{
\mathbf{U}\in\Phi^{*}:
\mathbf{U}=\text{weak limit of }\iota[\mathcal{R}_\omega[\mathfrak{D}_n]]
\right\},
\end{equation}
for admissible nets \(\mathfrak{D}_n\in\mathcal{J}_{\mathrm{rig}}\) \cite{kato1995,stratton1941,harrington2001}.
Equation~\eqref{eq:ch3.spectral-completion-radiative} is smaller than the full
\(\widehat{\mathcal{F}}_\Omega\) of Chapter~\ref{ch:02}: reactive and non-exporting sheets are not appended to
\(\mathcal{F}_{\mathrm{rad}}\) merely by taking weak limits \cite{stratton1941}. That restriction is the
analytic content of \(\Pi_{\mathrm{rad}}\) in the rigged setting \cite{harrington2001}.

\begin{proposition}[Weak--norm distinction for radiation sequences]
\label{prop:ch3.weak-norm-distinction}
Let \(\{\mathbf{u}_n\}\subset\mathcal{F}_{\mathrm{rad}}\) satisfy
Eq.~\eqref{eq:ch3.weak-star-convergence-rigged} in \(\Phi^{*}\) but
\(\|\mathbf{u}_n-\mathbf{u}_m\|_{\mathrm{rad}}\not\to 0\). Then \(\{\mathbf{u}_n\}\) is not Cauchy in
\(\mathcal{F}_{\mathrm{rad}}\), yet export coefficients
\(c(\beta)[\mathfrak{D}_n]\) may converge for each fixed \(\beta\in I_{\mathrm{c}}\) tested through
Eq.~\eqref{eq:ch3.continuous-coefficient-rigged} \cite{kato1995,stratton1941,harrington2001}. Far-zone
observables stable under weak convergence need not coincide with limits of near-zone energy norms
\cite{stratton1941,jackson1999}.
\end{proposition}
\begin{proof}[Proof sketch]
The first claim is the definition of Cauchy failure. Coefficient convergence follows from continuity of
\(\mathbf{U}\mapsto\mathbf{U}(\boldsymbol{\phi}_\beta)\) on \(\Phi^{*}\) for fixed localized tests
\(\boldsymbol{\phi}_\beta\in\Phi\) \cite{kato1995,stratton1941}.
% PROOF: AUTHOR REVIEW
\end{proof}

\paragraph{Forward obligations to Sections~\ref{sec:ch34}--\ref{sec:ch35}.}
Vector spherical harmonics and named \(\Mlm,\Nlm\) bases in Section~\ref{sec:ch34} supply explicit orthonormal
families on \(I_{\mathrm{d}}\) and hybrid fibers Eq.~\eqref{eq:ch3.hybrid-spectral-decomposition}
\cite{hansen1988,tai1994,stratton1941}. Section~\ref{sec:ch35} fixes \(d\mu\) and flux-normalized
\(\delta_\beta\) so that Eq.~\eqref{eq:ch3.continuous-parseval-identity} matches
Eq.~\eqref{eq:ch2.flux-on-radiative-branch} \cite{stratton1941,poynting1884}. Subsection~\ref{sec:ch314} does
not construct those bases; it certifies that any basis obeying outgoing closure and reciprocity is a
legitimate coordinate system on \(\iota[\mathcal{F}_{\mathrm{rad}}]^{\mathrm{cl}}\) \cite{harrington2001,hansen1988}.
Subsection~\ref{sec:ch343} will prove harmonic orthogonality inside this rigged frame by reference to
Theorem~\ref{thm:ch3.rigged-spectral-closure} \cite{stratton1941,hansen1988}.

Subsection~\ref{sec:ch314} has embedded \(\mathcal{F}_{\mathrm{rad}}\) in the Gelfand triple
Eq.~\eqref{eq:ch3.gelfand-triple-rad}, extended pairings
Eqs.~\eqref{eq:ch3.distributional-pairing-extension}--\eqref{eq:ch3.continuous-coefficient-rigged},
legitimized formal projectors Eq.~\eqref{eq:ch3.formal-projector-rigged-closure}, and extended
\(\mathcal{R}_\omega\) on \(\mathcal{J}_{\mathrm{rig}}\) through
Eq.~\eqref{eq:ch3.radiation-map-rigged-extension} \cite{kato1995,stratton1941,harrington2001}. Section~\ref{sec:ch32}
installs angular-momentum operators on this analytic base; Section~\ref{sec:ch35} fixes flux-normalized
\(\delta_\beta\) against Eq.~\eqref{eq:ch2.flux-on-radiative-branch} \cite{hansen1988,stratton1941}.

\paragraph{Summative closure of Section~\ref{sec:ch31}.}
Section~\ref{sec:ch31} began from the observation that coefficient pairings on \(\mathcal{F}_{\mathrm{rad}}\)
require more than operator symbols imported from Chapter~\ref{ch:02} \cite{stratton1941,harrington2001}.
Subsection~\ref{sec:ch311} supplied finite-energy carriers and shell diagnostics;
Subsection~\ref{sec:ch312} installed the reciprocity-compatible inner product and Riesz coefficient map;
Subsection~\ref{sec:ch313} extended spectral decomposition to continuous measures and generalized modes;
Subsection~\ref{sec:ch314} completed the rigged dual domain that legitimizes distributional modes and formal
projectors noted in Subsection~\ref{sec:ch224} \cite{kato1995,stratton1941,hansen1988}. Together,
Eqs.~\eqref{eq:ch3.F-rad-hilbert-closure}--\eqref{eq:ch3.radiation-map-rigged-extension} are the analytic
certificate that eigenfield expansions in Section~\ref{sec:ch34} and normalization in Section~\ref{sec:ch35}
act on the same export objects as \(\mathcal{R}_\omega[\mathfrak{D}]\), not on raw near-field samples
\cite{stratton1941,harrington2001}. Section~\ref{sec:ch32} now proceeds to symmetry operators on that closed
framework \cite{jackson1999,belinfante1940}.


\section{Angular Momentum Operators and SO(3)}
\label{sec:ch32}

Section~\ref{sec:ch31} closed the analytic carrier on which symmetry operators must act: finite-energy
\(\mathcal{F}_{\mathrm{rad}}\), reciprocity-compatible inner products, continuous spectra, and rigged
extensions for distributional modes \cite{kato1995,stratton1941,hansen1988}. Section~\ref{sec:ch32} installs
the \emph{Lie-algebra} layer: self-adjoint angular-momentum operators whose eigenvalues label exported
radiation and whose commutation relations organize degeneracy \cite{jackson1999,belinfante1940}. The step is
not cosmetic. Without explicit generators, rotation covariance in
Eq.~\eqref{eq:ch3.symmetry-commutes-with-radiation} remains a diagrammatic commutativity; with generators,
one can state which quantum numbers survive \(\mathcal{R}_\omega\), which directions are accessible on
\(S^{2}\), and which labels are gauge artifacts \cite{stratton1941,harrington2001}.

The present section follows the same forward discipline as Section~\ref{sec:ch31}. Subsection~\ref{sec:ch321}
defines total angular momentum \(\mathbf{J}\) on Maxwell radiative pairs \emph{once} for Volume~I,
importing directional transport maps from Subsection~\ref{sec:ch01.2.1} and far-field accessibility from
Subsection~\ref{sec:ch294} by reference only \cite{stratton1941,belinfante1940}. Subsection~\ref{sec:ch322}
derives commutation relations and simultaneous-eigenfield equations. Subsection~\ref{sec:ch323} connects
\(\mathrm{SO}(3)\) representations to spectral degeneracy on \(\mathcal{F}_{\mathrm{rad}}\), using rotation
rules from Section~\ref{sec:ch01.3} and Subsection~\ref{sec:ch262} \cite{jackson1999}. Subsection~\ref{sec:ch324}
develops Wigner rotations and irreducible coupling coefficients \cite{jackson1999,stratton1941}. Subsection~\ref{sec:ch325}
closes helicity, polarization, and dual-symmetry sectors under gauge discipline from
Section~\ref{sec:ch01.4} \cite{belinfante1940}. Vector spherical harmonics and named eigenfields remain in
Section~\ref{sec:ch34}; here the operators precede the basis \cite{hansen1988,tai1994}.

Laboratory readers should expect a shift in vocabulary. Section~\ref{sec:ch31} spoke of norms, measures, and
dual embeddings; Section~\ref{sec:ch32} speaks of generators, axes, and quantum numbers attached to
\(\mathcal{R}_\omega[\mathfrak{D}]\) \cite{stratton1941,harrington2001}. A finite \(\ell\) label is
meaningful only when it is an eigenvalue of the operators defined below on the \emph{same} outgoing carrier
fixed by Eq.~\eqref{eq:ch3.outgoing-inheritance-criterion}, not when it is read from a near-field phase ramp
or an unqualified paraxial mode order \cite{allen1992,belinfante1940}.


\subsection{Total angular momentum on Maxwell fields}
\label{sec:ch321}
% TOC tag: [HOME][USE SS1.2.1][USE SS2.9.4]

Rotations are implemented on functions by Eq.~\eqref{eq:ch3.symmetry-action-fields}; angular momentum is the
generator that produces those rotations in the infinitesimal limit. Subsection~\ref{sec:ch321} supplies the
Chapter~\ref{ch:03} \emph{single} definition of total angular momentum on Maxwell radiative pairs
\(\mathbf{u}_\omega=(\E,\H)\in\mathcal{F}_{\mathrm{rad}}\): orbital and spin channels are split explicitly,
the operator domain is tied to outgoing inheritance Eq.~\eqref{eq:ch3.outgoing-inheritance-criterion}, and
directional transport observables from Subsection~\ref{sec:ch01.2.1} enter only through accessibility filters
from Subsection~\ref{sec:ch294} \cite{stratton1941,belinfante1940,harrington2001}. Later sections cite
\(\mathbf{J}\) by reference; they do not introduce competing angular-momentum operators
\cite{hansen1988,stratton1941}.

\paragraph{Standing imports.}
Rotation generators on vector components are imported from
Eq.~\eqref{eq:ch1.rotation-generator-split} \cite{stratton1941,jackson1999}. Belinfante spin--orbit
redistribution and angular-density split Eq.~\eqref{eq:ch1.angular-density-split} are imported from
Section~\ref{sec:ch01.1.5} and Section~\ref{sec:ch01.4} \cite{belinfante1940,stratton1941}. Directional
power and angular-flux maps Eqs.~\eqref{eq:ch1.directional-power-density}--\eqref{eq:ch1.directional-am-flux},
qualification gate Eq.~\eqref{eq:ch1.directional-qualified-gate}, and rotation invariance
Eq.~\eqref{eq:ch1.rotation-am-invariance} enter from Subsection~\ref{sec:ch01.2.1} by reference
\cite{stratton1941,belinfante1940}. Far-field accessibility
Eqs.~\eqref{eq:ch2.far-angular-accessibility-def}--\eqref{eq:ch2.angular-accessibility-operator} and radiation
rotation commutation Eq.~\eqref{eq:ch2.radiation-rotation-commutation} are imported from
Subsections~\ref{sec:ch294} and~\ref{sec:ch262} \cite{stratton1941,harrington2001,jackson1999}. Maxwell
rotations on pairs use Eq.~\eqref{eq:ch2.field-rotation-representation}, equivalent to
Eq.~\eqref{eq:ch3.symmetry-action-fields} for \(g\in\mathrm{SO}(3)\) \cite{stratton1941,jackson1999}.

\paragraph{Orbital generator on Maxwell components.}
For unit axis \(\hat{\mathbf{u}}\) and scalar Cartesian component \(F\in\{E_x,E_y,E_z,H_x,H_y,H_z\}\),
Eq.~\eqref{eq:ch1.rotation-generator-split} defines the orbital generator action. On the Maxwell pair
\(\mathbf{u}_\omega=(\E,\H)\), define the orbital angular-momentum operator component
\begin{equation}
\label{eq:ch3.J-orbital-maxwell}
\big(\mathbf{J}_{\hat{\mathbf{u}}}^{\mathrm{orb}}\mathbf{u}_\omega\big)_{\E}
=
\hat{\mathbf{u}}\times\E
-\big(\hat{\mathbf{u}}\times\mathbf{r}\big)\cdot\nabla\E,
\qquad
\big(\mathbf{J}_{\hat{\mathbf{u}}}^{\mathrm{orb}}\mathbf{u}_\omega\big)_{\H}
=
\hat{\mathbf{u}}\times\H
-\big(\hat{\mathbf{u}}\times\mathbf{r}\big)\cdot\nabla\H,
\end{equation}
with \(\nabla\) acting componentwise \cite{stratton1941,jackson1999}. Equation~\eqref{eq:ch3.J-orbital-maxwell}
is the Subsection~\ref{sec:ch321} lift of Eq.~\eqref{eq:ch1.rotation-generator-split} to phasor Maxwell pairs:
the \(\hat{\mathbf{u}}\times\E\) term rotates polarization at fixed \(\mathbf{r}\); the
\((\hat{\mathbf{u}}\times\mathbf{r})\cdot\nabla\E\) term transports phase along rotation orbits
\cite{jackson1999,stratton1941}. Cartesian components \(\mathbf{J}_{a}^{\mathrm{orb}}\) use
\(\hat{\mathbf{u}}=\hat{\mathbf{e}}_a\) for \(a\in\{x,y,z\}\) \cite{jackson1999}.

\paragraph{Spin generator and Belinfante channel.}
Belinfante's analysis \cite{belinfante1940} requires an intrinsic spin contribution paired with
Eq.~\eqref{eq:ch1.angular-density-split}. On admissible radiative pairs, define the spin operator component by
the transverse Belinfante prescription
\begin{equation}
\label{eq:ch3.J-spin-maxwell}
\big(\mathbf{J}_{\hat{\mathbf{u}}}^{\mathrm{spin}}\mathbf{u}_\omega\big)_{\E}
=
\frac{1}{\jj\omega\mu_0}\,
\hat{\mathbf{u}}\times\big(\nabla\times\E\big)_{\perp},
\qquad
\big(\mathbf{J}_{\hat{\mathbf{u}}}^{\mathrm{spin}}\mathbf{u}_\omega\big)_{\H}
=
-\frac{1}{\jj\omega\varepsilon_0}\,
\hat{\mathbf{u}}\times\big(\nabla\times\H\big)_{\perp},
\end{equation}
where \((\cdot)_{\perp}\) removes the radial component along \(\hat{\mathbf{r}}\) on exterior domains
\cite{belinfante1940,hansen1988,stratton1941}. Equation~\eqref{eq:ch3.J-spin-maxwell} is the operator form of
the spin density channel in Eq.~\eqref{eq:ch1.angular-density-split} under the Volume~I phasor convention
\cite{belinfante1940,stratton1941}. Gauge redistribution within the equivalence class
Eq.~\eqref{eq:ch1.spin-orbit-equivalence-class} shifts local orbital and spin densities without changing
totals constrained by Eq.~\eqref{eq:ch1.angular-balance-integral}; Section~\ref{sec:ch39} records measurement
consequences \cite{belinfante1940,stratton1941}.

\paragraph{Total angular-momentum operator (Volume~I HOME).}
The \emph{total} angular-momentum operator on Maxwell radiative pairs is
\begin{equation}
\label{eq:ch3.J-total-maxwell-def}
\mathbf{J}_{\hat{\mathbf{u}}}
=
\mathbf{J}_{\hat{\mathbf{u}}}^{\mathrm{orb}}
+
\mathbf{J}_{\hat{\mathbf{u}}}^{\mathrm{spin}},
\qquad
\mathbf{J}
=
(\mathbf{J}_{x},\mathbf{J}_{y},\mathbf{J}_{z}),
\end{equation}
with \(\mathbf{J}_{a}=\mathbf{J}_{\hat{\mathbf{e}}_a}\) \cite{jackson1999,belinfante1940,hansen1988}.
Equation~\eqref{eq:ch3.J-total-maxwell-def} is the sole Volume~I definition of \(\mathbf{J}\) on Maxwell
fields; Sections~\ref{sec:ch34}--\ref{sec:ch38} cite Eq.~\eqref{eq:ch3.J-total-maxwell-def} when labeling
\(\ell,m\) sectors \cite{hansen1988,stratton1941}.

\paragraph{Infinitesimal limit and pair-space action.}
To first order in \(\varphi\), Eq.~\eqref{eq:ch3.rotation-from-J-generators} reduces to
\begin{equation}
\label{eq:ch3.J-infinitesimal-rotation}
\big(T_{R(\varphi,\hat{\mathbf{u}})}\mathbf{u}_\omega\big)(\mathbf{r})
=
\mathbf{u}_\omega(\mathbf{r})
+
\varphi\,\big(\mathbf{J}_{\hat{\mathbf{u}}}\mathbf{u}_\omega\big)(\mathbf{r})
+
\mathcal{O}(\varphi^{2}),
\end{equation}
which is the Maxwell-pair restatement of Eq.~\eqref{eq:ch1.infinitesimal-active-rotation} with
\(\mathbf{J}_{\hat{\mathbf{u}}}\) from Eq.~\eqref{eq:ch3.J-total-maxwell-def} \cite{jackson1999,stratton1941}.
Operationally, \(\mathbf{J}_{\hat{\mathbf{u}}}\) is a \(6\times 6\) block operator on the component stack
\((E_x,E_y,E_z,H_x,H_y,H_z)^{\mathsf{T}}\): orbital blocks couple coordinates through
\(\hat{\mathbf{u}}\times(\cdot)\) and \(\big(\hat{\mathbf{u}}\times\mathbf{r}\big)\cdot\nabla\), while spin
blocks couple \(\E\) and \(\H\) through transverse curls in Eq.~\eqref{eq:ch3.J-spin-maxwell}
\cite{belinfante1940,hansen1988}. Reactive near-field pairs can satisfy finiteness of
\(\|\mathbf{u}_\omega\|_{\mathcal{U}}\) while violating Eq.~\eqref{eq:ch3.J-infinitesimal-rotation} on shells
because \(\mathbf{J}^{\mathrm{orb}}\) amplifies \(r\)-weighted gradients that do not survive
\(\mathcal{R}_\omega\) \cite{stratton1941,harrington2001}.

The operator acts on
\(\mathbf{u}_\omega\in\mathcal{D}(\mathbf{J})\subset\Phi\subset\mathcal{F}_{\mathrm{rad}}\) with
\begin{equation}
\label{eq:ch3.J-operator-domain-rad}
\mathcal{D}(\mathbf{J})
=
\left\{
\mathbf{u}_\omega\in\Phi:
\mathbf{J}_{a}\mathbf{u}_\omega\in\mathcal{F}_{\mathrm{rad}}\ \forall a\in\{x,y,z\},\
\Pi_{\mathrm{rad}}[\mathbf{u}_\omega]=\mathbf{u}_\omega
\right\},
\end{equation}
using test space \(\Phi\) from Subsection~\ref{sec:ch314} and outgoing idempotence from
Eq.~\eqref{eq:ch3.outgoing-inheritance-criterion} \cite{kato1995,stratton1941,harrington2001}. Membership in
\(\mathcal{D}(\mathbf{J})\) therefore requires smooth enough decay for orbital derivatives and radiative
outgoing qualification, not merely finiteness of \(\|\cdot\|_{\mathcal{U}}\) near sources
\cite{stratton1941,jackson1999}.

\paragraph{Connection to rotation actions.}
Finite rotations generated by Eq.~\eqref{eq:ch3.J-total-maxwell-def} recover
Eq.~\eqref{eq:ch3.symmetry-action-fields} on \(\mathcal{D}(\mathbf{J})\):
\begin{equation}
\label{eq:ch3.rotation-from-J-generators}
\big(T_{R(\varphi,\hat{\mathbf{u}})}\mathbf{u}_\omega\big)(\mathbf{r})
=
\exp\!\big(-\jj\varphi\,\mathbf{J}_{\hat{\mathbf{u}}}\big)\,\mathbf{u}_\omega(\mathbf{r}),
\qquad
\mathbf{u}_\omega\in\mathcal{D}(\mathbf{J}),
\end{equation}
with \(T_g\) from Eq.~\eqref{eq:ch3.symmetry-action-fields} and \(R(\varphi,\hat{\mathbf{u}})\) from
Eq.~\eqref{eq:ch1.rodrigues-rotation} \cite{jackson1999,stratton1941}. Equation~\eqref{eq:ch3.rotation-from-J-generators}
upgrades the preview in Subsection~\ref{sec:ch301}: representation matrices \(\rho(g)\) in
Eq.~\eqref{eq:ch3.coefficient-rotation-law} are the finite-dimensional blocks of
Eq.~\eqref{eq:ch3.rotation-from-J-generators} on symmetry-adapted subspaces \cite{jackson1999,hansen1988}.
Theorem~\ref{thm:ch3.symmetry-commutes} and Eq.~\eqref{eq:ch2.radiation-rotation-commutation} are the
source-side and field-side commutativity laws compatible with
Eq.~\eqref{eq:ch3.rotation-from-J-generators} \cite{stratton1941,harrington2001}.

\begin{figure}[t]
\centering
\ChOneFigBlock{%
\begin{tikzpicture}[ch1fig, node distance=7mm and 9mm]
  \node[ch1box] (u) {$\mathbf{u}_\omega=(\E,\H)$};
  \node[ch1box, right=12mm of u] (orb) {$\mathbf{J}^{\mathrm{orb}}$\\Eq.~\eqref{eq:ch3.J-orbital-maxwell}};
  \node[ch1box, below=10mm of u] (spin) {$\mathbf{J}^{\mathrm{spin}}$\\Eq.~\eqref{eq:ch3.J-spin-maxwell}};
  \node[ch1box, right=12mm of spin] (tot) {$\mathbf{J}=\mathbf{J}^{\mathrm{orb}}+\mathbf{J}^{\mathrm{spin}}$};
  \node[ch1box, right=12mm of tot] (tg) {$T_g$\\Eq.~\eqref{eq:ch3.rotation-from-J-generators}};
  \draw[ch1flow] (u) -- (orb) -- (tot);
  \draw[ch1flow] (u) -- (spin) -- (tot);
  \draw[ch1flow] (tot) -- (tg);
\end{tikzpicture}%
}
\caption{Subsection~\ref{sec:ch321} operator split: orbital and spin generators sum to total
\(\mathbf{J}\), whose exponentials reproduce spatial rotations on \(\mathcal{D}(\mathbf{J})\).}
\label{fig:ch3.J-orbital-spin-split}
\end{figure}

Figure~\ref{fig:ch3.J-orbital-spin-split} is the generator-level companion to
Figure~\ref{fig:ch3.operator-representation-bridge}: the latter tracked operators versus coordinates; the
present figure tracks orbital versus spin \emph{channels} inside \(\mathbf{J}\) \cite{belinfante1940,stratton1941}.

\paragraph{Self-adjointness on the radiative inner product.}
On \(\mathcal{D}(\mathbf{J})\cap\mathcal{F}_{\mathrm{rad}}\), define the symmetric sesquilinear form
\begin{equation}
\label{eq:ch3.J-symmetric-pairing}
\mathfrak{j}_{\hat{\mathbf{u}}}(\mathbf{u},\mathbf{v})
=
\frac{1}{2}\Big[
\langle \mathbf{J}_{\hat{\mathbf{u}}}\mathbf{u},\mathbf{v}\rangle_{\mathrm{rad}}
+
\langle \mathbf{u},\mathbf{J}_{\hat{\mathbf{u}}}\mathbf{v}\rangle_{\mathrm{rad}}^{*}
\Big],
\end{equation}
using Eq.~\eqref{eq:ch3.radiative-inner-product-def} \cite{stratton1941,hansen1988}. When
\(\mathfrak{j}_{\hat{\mathbf{u}}}\) is real-valued and bounded on a dense subset of
\(\mathcal{D}(\mathbf{J})\), \(\mathbf{J}_{\hat{\mathbf{u}}}\) is essentially self-adjoint on the radiative
closure and admits spectral projectors used in Subsection~\ref{sec:ch322} \cite{kato1995,jackson1999}. The
technical closure is standard on outgoing transverse fields with rapid decay; near-field reactive pairs fail
\(\mathcal{D}(\mathbf{J})\) before self-adjointness is tested \cite{stratton1941,harrington2001}.

\paragraph{Directional transport functional.}
Subsection~\ref{sec:ch01.2.1} measures angular transport through \(\boldsymbol{\Psi}_J\) on shells
\cite{stratton1941,belinfante1940}. On \(\mathcal{R}_\omega[\mathfrak{D}]\), define the axis transport
functional
\begin{equation}
\label{eq:ch3.J-directional-transport-functional}
\mathcal{T}_{\hat{\mathbf{u}}}[\mathbf{u}_\omega]
=
\int_{\Gamma_R}
\hat{\mathbf{u}}\cdot\boldsymbol{\Psi}_J(\theta,\phi;R)\,d\Omega,
\qquad
\mathbf{u}_\omega=\mathcal{R}_\omega[\mathfrak{D}],
\end{equation}
with \(\boldsymbol{\Psi}_J\) from Eq.~\eqref{eq:ch1.directional-am-flux} and \(R\) in the far-zone regime of
Theorem~\ref{thm:ch2.far-field-approximation} \cite{stratton1941,jackson1999}. Equation
\eqref{eq:ch3.J-directional-transport-functional} connects the \emph{operator} \(\mathbf{J}_{\hat{\mathbf{u}}}\)
to the \emph{flux map} \(\boldsymbol{\Psi}_J\): for qualified channels with
\(\mathcal{Q}_{J,\Omega}=1\) from Eq.~\eqref{eq:ch1.directional-qualified-gate},
\(\mathcal{T}_{\hat{\mathbf{u}}}\) tracks the same transport content as the expectation of
\(\mathbf{J}_{\hat{\mathbf{u}}}\) in the radiative inner product on transverse far fields
\cite{belinfante1940,stratton1941}. Accessibility further restricts directions through
\begin{equation}
\label{eq:ch3.J-accessibility-composition}
\mathcal{T}_{\hat{\mathbf{u}}}^{\mathrm{acc}}
=
\mathcal{T}_{\hat{\mathbf{u}}}\circ\mathcal{A}_{\mathrm{ff}},
\end{equation}
with \(\mathcal{A}_{\mathrm{ff}}\) from Eq.~\eqref{eq:ch2.angular-accessibility-operator}
\cite{stratton1941,harrington2001}. Directions outside
\(\mathcal{K}_{\mathrm{acc}}^{\mathrm{ff}}\) from Eq.~\eqref{eq:ch2.far-angular-accessibility-def} may carry
local phase structure without contributing to \(\mathcal{T}_{\hat{\mathbf{u}}}^{\mathrm{acc}}\)
\cite{harrington2001,jackson1999}.

\paragraph{Local spin--orbit gauge split versus exported totals.}
Belinfante redistribution Eq.~\eqref{eq:ch1.spin-orbit-equivalence-class} implies that
\(\mathbf{J}^{\mathrm{orb}}\) and \(\mathbf{J}^{\mathrm{spin}}\) are not individually gauge invariant:
adding a longitudinal gradient to \(\E\) shifts orbital density without changing
\(\mathbf{J}_{\hat{\mathbf{u}}}\mathbf{u}_\omega\) when the pair remains in \(\mathcal{D}(\mathbf{J})\)
\cite{belinfante1940,stratton1941}. Laboratory diagnostics that report ``orbital fraction'' or ``spin
fraction'' of angular transport therefore require the same gauge screen \(\mathfrak{A}_{\mathrm{g}}\) as
Section~\ref{sec:ch01.4}, not a naive partition of \((\E,\H)\) samples \cite{belinfante1940}. By contrast,
\(\mathcal{T}_{\hat{\mathbf{u}}}\) and qualified totals in Eq.~\eqref{eq:ch1.rotation-am-invariance} are
rotation invariant once \(\mathcal{Q}_{J,\Omega}=1\); Subsection~\ref{sec:ch321} uses those invariants as the
observable anchors while treating channelwise orbital and spin densities as intermediate operator coordinates
\cite{stratton1941,harrington2001}.

\paragraph{Intertwining with radiation export.}
Rotation covariance implies that \(\mathbf{J}\) cannot arbitrarily mix reactive sheets removed by
\(\Pi_{\mathrm{rad}}\). On \(\mathcal{D}(\mathbf{J})\),
\begin{equation}
\label{eq:ch3.J-radiation-intertwining}
\mathbf{J}_{\hat{\mathbf{u}}}\,\mathcal{R}_\omega[\mathfrak{D}]
=
\mathcal{R}_\omega\big[\mathbf{J}_{\hat{\mathbf{u}}}^{(\mathrm{src})}\mathfrak{D}\big]
\quad\text{when}\ 
\mathcal{D}_{J,\mathrm{reg}}=1,
\end{equation}
with \(\mathbf{J}_{\hat{\mathbf{u}}}^{(\mathrm{src})}\) the source-side generator induced by
\(\mathcal{J}_g\) in Eq.~\eqref{eq:ch3.current-pushforward} and outgoing closure preserved
\cite{stratton1941,harrington2001,jackson1999}. Equation~\eqref{eq:ch3.J-radiation-intertwining} is the
generator form of Eq.~\eqref{eq:ch2.radiation-rotation-commutation}: rotating the source before export equals
rotating the exported image when accessibility and regime gates pass \cite{stratton1941}. Failure of
Eq.~\eqref{eq:ch3.J-radiation-intertwining} diagnoses frame mismatch, anisotropic media, or near-field
evaluation---not a breakdown of Eq.~\eqref{eq:ch1.angular-balance-integral} itself \cite{harrington2001}.

\begin{proposition}[Generators reproduce spatial rotations on \(\mathcal{D}(\mathbf{J})\)]
\label{prop:ch3.J-generates-rotations}
For \(R(\varphi,\hat{\mathbf{u}})\in\mathrm{SO}(3)\) and \(\mathbf{u}_\omega\in\mathcal{D}(\mathbf{J})\),
Eqs.~\eqref{eq:ch3.rotation-from-J-generators} and~\eqref{eq:ch3.symmetry-action-fields} agree on
\(\mathcal{F}_{\mathrm{rad}}\). Moreover, Eq.~\eqref{eq:ch3.J-radiation-intertwining} holds for admissible
\(\mathfrak{D}\) satisfying Eqs.~\eqref{eq:ch2.radiation-rotation-commutation} and
\(\mathcal{D}_{J,\mathrm{reg}}=1\) \cite{jackson1999,stratton1941,harrington2001}.
\end{proposition}
\begin{proof}[Proof sketch]
Expand Eq.~\eqref{eq:ch3.rotation-from-J-generators} to first order and compare with
Eq.~\eqref{eq:ch1.infinitesimal-active-rotation} on each component, using
Eqs.~\eqref{eq:ch3.J-orbital-maxwell}--\eqref{eq:ch3.J-spin-maxwell} \cite{jackson1999,stratton1941}. The
radiation intertwining follows by differentiating Eq.~\eqref{eq:ch2.radiation-rotation-commutation} along the
generated orbit \cite{stratton1941,harrington2001}.
% PROOF: AUTHOR REVIEW
\end{proof}

\paragraph{Spectral labeling preview.}
Commutation relations and simultaneous eigenfields are derived in Subsection~\ref{sec:ch322}; here only the
target equations are recorded. The Casimir operator
\begin{equation}
\label{eq:ch3.J-squared-def}
\mathbf{J}^{2}
=
\mathbf{J}_{x}^{2}+\mathbf{J}_{y}^{2}+\mathbf{J}_{z}^{2}
\end{equation}
acts on \(\mathcal{D}(\mathbf{J})\) and, when essentially self-adjoint on the radiative closure, admits
eigenvalues \(\hbar^{2}\ell(\ell+1)\) in the convention \(\hbar=1\) used throughout Volume~I
\cite{jackson1999,hansen1988}. Simultaneous eigenfields satisfy
\begin{equation}
\label{eq:ch3.J-simultaneous-eigenfield-preview}
\mathbf{J}^{2}\mathbf{u}_{\ell m}
=
\ell(\ell+1)\,\mathbf{u}_{\ell m},
\qquad
\mathbf{J}_{z}\mathbf{u}_{\ell m}
=
m\,\mathbf{u}_{\ell m},
\qquad
\mathbf{u}_{\ell m}\in\mathcal{D}(\mathbf{J})\cap\mathcal{F}_{\mathrm{rad}},
\end{equation}
with integer or half-integer \(\ell\) and \(m\in\{-\ell,\ldots,\ell\}\) selected by the representation
content of the exported pair \cite{jackson1999,stratton1941}. Equation~\eqref{eq:ch3.J-simultaneous-eigenfield-preview}
is the operator certificate behind modal indices in Eq.~\eqref{eq:ch3.spectral-index-alignment}: a coefficient
label \((\ell,m)\) is meaningful on \(\mathcal{R}_\omega[\mathfrak{D}]\) only when the corresponding mode lies
in the eigenspace of Eqs.~\eqref{eq:ch3.J-squared-def} and~\eqref{eq:ch3.J-simultaneous-eigenfield-preview}, not
when it is inferred from a harmonic fit on a partial angular scan \cite{hansen1988,harrington2001}.

\begin{theorem}[Angular-momentum labeling admissibility on \(\mathcal{F}_{\mathrm{rad}}\)]
\label{thm:ch3.J-labeling-admissibility}
An angular label \(\ell\) attached to \(\mathcal{R}_\omega[\mathfrak{D}]\) is Volume~I admissible only if
it is a spectral label of \(\mathbf{J}^{2}\) Eq.~\eqref{eq:ch3.J-squared-def} on the outgoing carrier
Eq.~\eqref{eq:ch3.outgoing-inheritance-criterion}, the direction satisfies
\(\widehat{\mathbf{k}}\in\mathcal{K}_{\mathrm{acc}}^{\mathrm{ff}}\) when transport is claimed, and gauge
screen \(\mathfrak{A}_{\mathrm{g}}\) from Section~\ref{sec:ch01.4} passes for the reported observable
\cite{belinfante1940,stratton1941,harrington2001}. Labels read from reactive-shell samples, suppressed
directions, or unqualified paraxial phase ramps fail at least one clause \cite{allen1992,stratton1941}.
\end{theorem}
\begin{proof}[Proof sketch]
Combine Eq.~\eqref{eq:ch3.J-operator-domain-rad},
Eqs.~\eqref{eq:ch3.J-accessibility-composition} and~\eqref{eq:ch2.far-angular-accessibility-def}, and
gauge equivalence Eq.~\eqref{eq:ch1.spin-orbit-equivalence-class} \cite{belinfante1940,stratton1941,harrington2001}.
% PROOF: AUTHOR REVIEW
\end{proof}

\paragraph{Validity domain.}
Eqs.~\eqref{eq:ch3.J-orbital-maxwell}--\eqref{eq:ch3.J-radiation-intertwining} assume linear isotropic
media at fixed \(\omega\), outgoing \(\mathcal{B}_\Gamma\), and reciprocal response where flux symmetry is used
\cite{stratton1941,harrington2001}. Anisotropic, gyrotropic, or moving platforms break
Eq.~\eqref{eq:ch3.J-radiation-intertwining} even when power rotation invariance holds
\cite{stratton1941,jackson1999}. Paraxial limits and scalar reductions are treated in Section~\ref{sec:ch39};
Subsection~\ref{sec:ch321} works in the full Maxwell pair space \cite{belinfante1940,stratton1941}.

\paragraph{Worked illustration (short dipole angular flux).}
Consider a centered electric dipole \(\mathbf{p}=p_{0}\hat{\mathbf{z}}\) at the origin in free space. In the
far zone \(kr\gg 1\), the exported transverse field scales as
\(\E_{\mathrm{far}}\propto k^{2}\,\hat{\mathbf{r}}\times(\hat{\mathbf{r}}\times\hat{\mathbf{z}})\,e^{\jj kr}/r\)
and \(\H_{\mathrm{far}}\propto k\,\hat{\mathbf{r}}\times\E_{\mathrm{far}}\) \cite{jackson1999,stratton1941}.
Applying \(\mathbf{J}_{z}\) through Eqs.~\eqref{eq:ch3.J-orbital-maxwell}--\eqref{eq:ch3.J-total-maxwell-def}
on this outgoing pair yields the \(\ell=1\), \(m=0\) sector anticipated by
Eq.~\eqref{eq:ch3.J-simultaneous-eigenfield-preview} \cite{hansen1988,jackson1999}. The axis transport
\(\mathcal{T}_{\hat{\mathbf{z}}}\) from Eq.~\eqref{eq:ch3.J-directional-transport-functional} is nonzero on
qualified channels with \(\mathcal{Q}_{J,\Omega}=1\) because \(\boldsymbol{\Psi}_J\) inherits the dipole
angular pattern \cite{stratton1941,belinfante1940}. By contrast, a reactive-shell sample at \(kr\ll 1\) can
exhibit large \(\|\mathbf{J}^{\mathrm{orb}}\mathbf{u}_\omega\|_{\mathcal{U}}\) on
\(\mathcal{S}_\omega[\mathfrak{D}]\) while \(\mathcal{R}_\omega[\mathfrak{D}]\) remains \(\ell=1\) in the far
zone: the near-field orbital operator probes \(r\)-weighted gradients that are not torque-capable flux
\cite{harrington2001,stratton1941}. The illustration enforces domain
Eq.~\eqref{eq:ch3.J-operator-domain-rad}: spectral labels attach to
\(\mathcal{R}_\omega[\mathfrak{D}]\), not to raw reactive samples \cite{stratton1941}.

\paragraph{Worked illustration (rotated frame without rotated source).}
Rotating the measurement frame while leaving \(\Jimp\) fixed rotates reported \(\boldsymbol{\Psi}_J\) maps by
Eq.~\eqref{eq:ch1.rotation-am-invariance} on totals but mixes channelwise components unless
Eq.~\eqref{eq:ch1.active-passive-duality} is applied consistently \cite{stratton1941,jackson1999}. Comparing
modal \(\ell\) labels before applying \(\mathcal{D}(R)\) from Eq.~\eqref{eq:ch2.field-rotation-representation}
violates Eq.~\eqref{eq:ch3.J-radiation-intertwining} and produces spurious degeneracy splitting in numerical
coupling matrices \cite{harrington2001,stratton1941}.

\paragraph{Worked illustration (suppressed far direction).}
When \(\widehat{\mathbf{k}}\in\mathcal{K}_{\mathrm{sup}}\) from
Eqs.~\eqref{eq:ch2.angular-accessibility-sector}--\eqref{eq:ch2.suppressed-angular-sheet-def}, local fields
may be nonzero while \(\mathcal{T}_{\hat{\mathbf{u}}}^{\mathrm{acc}}=0\) \cite{stratton1941,harrington2001}.
Assigning \(\ell\) from a partial scan of such a direction is inadmissible under
Theorem~\ref{thm:ch3.J-labeling-admissibility} even if a harmonic fit returns small residuals
\cite{harrington2001,jackson1999}.

\paragraph{Discretization and finite-array caveat.}
Numerical implementations approximate \(\mathbf{J}_{\hat{\mathbf{u}}}\) by finite differences on sampled
\((\E,\H)\) grids. Such matrices are not automatically self-adjoint with respect to
Eq.~\eqref{eq:ch3.radiative-inner-product-def} unless the grid includes outgoing test functions from
\(\Phi\) and the accessibility mask from Eq.~\eqref{eq:ch2.angular-accessibility-operator} is applied before
eigenanalysis \cite{harrington2001,stratton1941}. A common failure mode fits near-field samples with spherical
harmonics, extracts an \(\ell\) index from the largest coefficient, and reports it as ``beam OAM'' without
passing Table~\ref{tab:ch3.J-reporting-gates}: the procedure violates at least
Eqs.~\eqref{eq:ch3.J-operator-domain-rad} and~\eqref{eq:ch2.far-angular-accessibility-def} simultaneously
\cite{allen1992,belinfante1940,harrington2001}. Subsection~\ref{sec:ch322} supplies the commutation checks that
turn Eq.~\eqref{eq:ch3.J-squared-def} into a diagnosable Hermiticity test on symmetry-adapted subspaces
\cite{jackson1999,hansen1988}.

\begin{table}[t]
\centering
\caption{Angular-momentum reporting gates for Subsection~\ref{sec:ch321} (operator row is new; transport and
accessibility rows cite Subsections~\ref{sec:ch01.2.1} and~\ref{sec:ch294}).}
\label{tab:ch3.J-reporting-gates}
\small
\begin{tabular}{@{}p{0.24\linewidth}p{0.32\linewidth}p{0.36\linewidth}@{}}
\toprule
Gate & Criterion anchor & Failure symptom \\
\midrule
Operator domain &
Eq.~\eqref{eq:ch3.J-operator-domain-rad} &
\(\ell\) labels from reactive near fields \\
\addlinespace
Directional transport &
Eq.~\eqref{eq:ch1.directional-qualified-gate} &
Phase ramps without \(\boldsymbol{\Psi}_J\) flux \\
\addlinespace
Far accessibility &
Eq.~\eqref{eq:ch2.far-angular-accessibility-def} &
Sidelobes declared on suppressed \(\widehat{\mathbf{k}}\) \\
\addlinespace
Rotation covariance &
Eq.~\eqref{eq:ch2.radiation-rotation-commutation} &
Asymmetric coupling after frame swap \\
\bottomrule
\end{tabular}
\end{table}

Table~\ref{tab:ch3.J-reporting-gates} is an operator checklist, not a repeat of
Table~\ref{tab:ch3.topology-convergence}: the latter diagnosed convergence layers; the present table diagnoses
\emph{label admissibility} for \(\mathbf{J}\) reports \cite{stratton1941,harrington2001}.

Subsection~\ref{sec:ch321} has defined total angular momentum
Eqs.~\eqref{eq:ch3.J-orbital-maxwell}--\eqref{eq:ch3.J-total-maxwell-def} on
\(\mathcal{D}(\mathbf{J})\subset\mathcal{F}_{\mathrm{rad}}\), linked rotations through
Eq.~\eqref{eq:ch3.rotation-from-J-generators}, connected directional transport through
Eqs.~\eqref{eq:ch3.J-directional-transport-functional}--\eqref{eq:ch3.J-accessibility-composition}, and
recorded radiation intertwining Eq.~\eqref{eq:ch3.J-radiation-intertwining} \cite{belinfante1940,jackson1999,stratton1941,harrington2001}.
Subsection~\ref{sec:ch322} derives commutation relations \([\mathbf{J}_i,\mathbf{J}_j]\) and the
simultaneous-eigenfield equations that Section~\ref{sec:ch34} will solve in spherical coordinates
\cite{hansen1988,stratton1941}.


\subsection{Commutation and simultaneous eigenfields}
\label{sec:ch322}
% TOC tag: [HOME]

Generators alone do not yet specify which field patterns can carry a definite angular label.
Subsection~\ref{sec:ch322} derives the Lie-algebra commutation relations satisfied by
\(\mathbf{J}\) from Subsection~\ref{sec:ch321}, installs raising and lowering operators on
\(\mathcal{D}(\mathbf{J})\), and records the simultaneous-eigenfield equations that Section~\ref{sec:ch34}
will solve once spherical separation is available \cite{jackson1999,hansen1988,stratton1941}. The outcome is
algebraic: exported radiation admits a basis in which \(\mathbf{J}^{2}\) and \(\mathbf{J}_{z}\) are jointly
diagonal, and ladder moves connect degenerate partners at fixed \(\ell\) \cite{hansen1988}.

\paragraph{Commutator on the Maxwell pair domain.}
For operators \(\mathcal{A},\mathcal{B}\) on \(\mathcal{D}(\mathbf{J})\), define
\begin{equation}
\label{eq:ch3.J-commutator-def}
[\mathcal{A},\mathcal{B}]
=
\mathcal{A}\mathcal{B}-\mathcal{B}\mathcal{A},
\qquad
[\mathcal{A},\mathcal{B}]\mathbf{u}_\omega
=
\mathcal{A}(\mathcal{B}\mathbf{u}_\omega)-\mathcal{B}(\mathcal{A}\mathbf{u}_\omega),
\end{equation}
with \(\mathbf{u}_\omega=(\E,\H)\) \cite{jackson1999}. Equation~\eqref{eq:ch3.J-commutator-def} is the
operator commutator on the six-component pair space; all products are evaluated on outgoing pairs satisfying
Eq.~\eqref{eq:ch3.J-operator-domain-rad} \cite{stratton1941,harrington2001}. Cartesian components
\(\mathbf{J}_{a}=\mathbf{J}_{\hat{\mathbf{e}}_a}\) from Eq.~\eqref{eq:ch3.J-total-maxwell-def} inherit the
same bracket \cite{jackson1999,belinfante1940}.

\paragraph{Orbital and spin channels separately.}
Write \(\mathbf{L}_{a}\equiv\mathbf{J}_{a}^{\mathrm{orb}}\) and \(\mathbf{S}_{a}\equiv\mathbf{J}_{a}^{\mathrm{spin}}\)
from Eqs.~\eqref{eq:ch3.J-orbital-maxwell}--\eqref{eq:ch3.J-spin-maxwell} \cite{belinfante1940,hansen1988}.
On scalar outgoing functions \(\psi\in C^{\infty}\) with rapid decay, the orbital generators satisfy the
standard \(\mathfrak{so}(3)\) algebra
\begin{equation}
\label{eq:ch3.L-commutation-scalar}
[L_i,L_j]\psi
=
\jj\,\varepsilon_{ijk}\,L_k\psi,
\qquad
L_k\psi
=
\big(\hat{\mathbf{e}}_k\times\mathbf{r}\big)\cdot\nabla\psi,
\end{equation}
with \(\varepsilon_{ijk}\) the Levi-Civita symbol and \(\jj\) the imaginary unit of the phasor convention
\cite{jackson1999,stratton1941}. Lifting Eq.~\eqref{eq:ch3.L-commutation-scalar} to vector components in
Eq.~\eqref{eq:ch3.J-orbital-maxwell} produces cross-term commutators between \(\hat{\mathbf{e}}_k\times\E\)
and \((\hat{\mathbf{e}}_k\times\mathbf{r})\cdot\nabla\E\) that cancel in the total orbital bracket on
transverse radiative pairs \cite{hansen1988,stratton1941}. The spin channel satisfies
\begin{equation}
\label{eq:ch3.S-commutation-pair}
[S_i,S_j]\mathbf{u}_\omega
=
\jj\,\varepsilon_{ijk}\,S_k\mathbf{u}_\omega
\end{equation}
when \(\nabla\times\E=\jj\omega\mu_0\H\) and \(\nabla\times\H=-\jj\omega\varepsilon_0\E\) hold on the domain
and transverse projection \((\cdot)_{\perp}\) in Eq.~\eqref{eq:ch3.J-spin-maxwell} is consistent with
outgoing decay \cite{belinfante1940,jackson1999,hansen1988}.

\begin{theorem}[Total angular-momentum commutation on \(\mathcal{D}(\mathbf{J})\)]
\label{thm:ch3.J-so3-commutation}
For \(i,j,k\in\{x,y,z\}\) and \(\mathbf{u}_\omega\in\mathcal{D}(\mathbf{J})\),
\begin{equation}
\label{eq:ch3.J-so3-commutation}
[\mathbf{J}_{i},\mathbf{J}_{j}]\mathbf{u}_\omega
=
\jj\,\varepsilon_{ijk}\,\mathbf{J}_{k}\mathbf{u}_\omega.
\end{equation}
Consequently \(\mathbf{J}^{2}\) Eq.~\eqref{eq:ch3.J-squared-def} commutes with every component:
\([\mathbf{J}^{2},\mathbf{J}_{a}]=0\) on \(\mathcal{D}(\mathbf{J})\) \cite{jackson1999,belinfante1940,hansen1988}.
\end{theorem}
\begin{proof}[Proof sketch]
Expand \([\mathbf{J}_i,\mathbf{J}_j]=[\mathbf{L}_i+\mathbf{S}_i,\mathbf{L}_j+\mathbf{S}_j]\) into four terms.
Eqs.~\eqref{eq:ch3.L-commutation-scalar}--\eqref{eq:ch3.S-commutation-pair} supply the pure orbital and pure
spin brackets \cite{jackson1999,hansen1988}. Mixed \(\mathbf{L}\)--\(\mathbf{S}\) commutators on Maxwell pairs
cancel when the Belinfante coupling in Eq.~\eqref{eq:ch3.J-spin-maxwell} is paired with the orbital transport
terms in Eq.~\eqref{eq:ch3.J-orbital-maxwell} and the pair obeys time-harmonic Maxwell on the exterior domain
\cite{belinfante1940,stratton1941}. The \(\mathbf{J}^{2}\) commutativity follows by expanding
\(\mathbf{J}^{2}=\sum_a\mathbf{J}_a^{2}\) and applying Eq.~\eqref{eq:ch3.J-so3-commutation} termwise
\cite{jackson1999}.
% PROOF: AUTHOR REVIEW
\end{proof}

Equation~\eqref{eq:ch3.J-so3-commutation} is the operator certificate behind rotation homomorphism
Eq.~\eqref{eq:ch3.representation-homomorphism}: finite rotations compose because the generators close under
the same Lie algebra as \(\mathrm{SO}(3)\) \cite{jackson1999,stratton1941}. Near-field reactive pairs that
fail Eq.~\eqref{eq:ch3.J-operator-domain-rad} need not satisfy Eq.~\eqref{eq:ch3.J-so3-commutation} even when
individual orbital or spin brackets hold locally \cite{harrington2001,stratton1941}.

\paragraph{Cartesian commutator reduction (illustrative step).}
The \(z\)-component bracket on a scalar outgoing mode \(\psi\) illustrates how coordinate and derivative terms
assemble. With \(L_x\psi=(\hat{\mathbf{e}}_x\times\mathbf{r})\cdot\nabla\psi\) and
\(L_y\psi=(\hat{\mathbf{e}}_y\times\mathbf{r})\cdot\nabla\psi\),
\begin{equation}
\label{eq:ch3.Lx-Ly-commutator-scalar}
[L_x,L_y]\psi
=
-\jj\big(\hat{\mathbf{e}}_z\times\mathbf{r}\big)\cdot\nabla\psi
=
\jj L_z\psi,
\end{equation}
using \([\partial_{x_i},x_j]=\delta_{ij}\) on rapidly decaying domains \cite{jackson1999,stratton1941}.
Vector lifting replaces \(\psi\) by each Cartesian component of \(\E\) and \(\H\), producing additional
\(\hat{\mathbf{e}}_a\times(\cdot)\) terms that cancel pairwise when the pair is transverse and outgoing
\cite{hansen1988}. The spin commutator Eq.~\eqref{eq:ch3.S-commutation-pair} follows the same pattern after
substituting \(\nabla\times\E=\jj\omega\mu_0\H\) and \(\nabla\times\H=-\jj\omega\varepsilon_0\E\), which trades
curls for cross-channel couplings that restore the \(\varepsilon_{ijk}\) structure \cite{belinfante1940,jackson1999}.
Equation~\eqref{eq:ch3.Lx-Ly-commutator-scalar} is not re-derived for every component: once
Theorem~\ref{thm:ch3.J-so3-commutation} is certified on \(\mathcal{D}(\mathbf{J})\), all Cartesian brackets are
equivalent under index rotation \cite{jackson1999}.

\paragraph{Raising and lowering operators.}
Define the Cartesian ladder combinations
\begin{equation}
\label{eq:ch3.J-plus-minus-def}
\mathbf{J}_{\pm}
=
\mathbf{J}_{x}\pm\jj\,\mathbf{J}_{y},
\end{equation}
acting on \(\mathcal{D}(\mathbf{J})\) \cite{jackson1999,hansen1988}. From
Eq.~\eqref{eq:ch3.J-so3-commutation},
\begin{equation}
\label{eq:ch3.J-z-ladder-commutation}
[\mathbf{J}_{z},\mathbf{J}_{\pm}]
=
\pm\,\mathbf{J}_{\pm},
\qquad
[\mathbf{J}^{2},\mathbf{J}_{\pm}]
=
0,
\end{equation}
so \(\mathbf{J}_{\pm}\) change the magnetic quantum number \(m\) without altering the \(\ell\) label carried
by \(\mathbf{J}^{2}\) \cite{jackson1999}. Equations~\eqref{eq:ch3.J-plus-minus-def}--\eqref{eq:ch3.J-z-ladder-commutation}
are the Maxwell-pair restatement of the standard \(\mathfrak{su}(2)\) ladder algebra; they are not new physics,
but they are the bookkeeping operators used when rotating coefficient vectors in
Eq.~\eqref{eq:ch3.coefficient-rotation-law} \cite{hansen1988,stratton1941}.

\begin{figure}[t]
\centering
\ChOneFigBlock{%
\begin{tikzpicture}[ch1fig, node distance=8mm and 10mm]
  \node[ch1box] (j2) {$\mathbf{J}^{2}$\\$\ell(\ell+1)$};
  \node[ch1box, below=9mm of j2] (jz) {$\mathbf{J}_{z}$\\$m$};
  \node[ch1box, left=14mm of jz] (jm) {$\mathbf{J}_{-}$\\$m\!\to\! m-1$};
  \node[ch1box, right=14mm of jz] (jp) {$\mathbf{J}_{+}$\\$m\!\to\! m+1$};
  \node[ch1box, below=11mm of jz] (u) {$\mathbf{u}_{\ell m}\in\mathcal{D}(\mathbf{J})$};
  \draw[ch1flow] (jp) -- (jz);
  \draw[ch1flow] (jm) -- (jz);
  \draw[ch1flow] (jz) -- (j2);
  \draw[ch1flow] (u) -- (jz);
\end{tikzpicture}%
}
\caption{Subsection~\ref{sec:ch322} simultaneous-diagonalization scaffold: \(\mathbf{J}^{2}\) fixes \(\ell\),
\(\mathbf{J}_{z}\) fixes \(m\), and \(\mathbf{J}_{\pm}\) connect partners at fixed \(\ell\) on
\(\mathcal{D}(\mathbf{J})\cap\mathcal{F}_{\mathrm{rad}}\).}
\label{fig:ch3.J-ladder-simultaneous}
\end{figure}

Figure~\ref{fig:ch3.J-ladder-simultaneous} compresses the algebraic target: an eigenfield is specified by the
pair \((\ell,m)\) only when both \(\mathbf{J}^{2}\) and \(\mathbf{J}_{z}\) are jointly diagonal on the
outgoing carrier \cite{jackson1999,hansen1988}.

\paragraph{Simultaneous eigenfield system.}
Subsection~\ref{sec:ch321} recorded the target Eqs.~\eqref{eq:ch3.J-squared-def} and
\eqref{eq:ch3.J-simultaneous-eigenfield-preview}; Subsection~\ref{sec:ch322} now elevates them to the defining
eigenfield system on \(\mathcal{F}_{\mathrm{rad}}\). A nonzero
\(\mathbf{u}_{\ell m}\in\mathcal{D}(\mathbf{J})\cap\mathcal{F}_{\mathrm{rad}}\) is a
\emph{simultaneous angular-momentum eigenfield} if
\begin{equation}
\label{eq:ch3.J-simultaneous-eigenfield-system}
\mathbf{J}^{2}\mathbf{u}_{\ell m}
=
\ell(\ell+1)\,\mathbf{u}_{\ell m},
\qquad
\mathbf{J}_{z}\mathbf{u}_{\ell m}
=
m\,\mathbf{u}_{\ell m},
\end{equation}
with \(\ell\ge 0\) and \(m\in\{-\ell,-\ell+1,\ldots,\ell\}\) in the integer sector realized by classical
multipoles \cite{jackson1999,hansen1988,stratton1941}. The ladder action is
\begin{equation}
\label{eq:ch3.J-ladder-action}
\mathbf{J}_{\pm}\mathbf{u}_{\ell m}
=
\sqrt{\ell(\ell+1)-m(m\pm 1)}\;\mathbf{u}_{\ell,m\pm 1},
\end{equation}
including the terminal conditions \(\mathbf{J}_{+}\mathbf{u}_{\ell\ell}=\mathbf{0}\) and
\(\mathbf{J}_{-}\mathbf{u}_{\ell,-\ell}=\mathbf{0}\) \cite{jackson1999,hansen1988}. Equation
\eqref{eq:ch3.J-simultaneous-eigenfield-system} is the eigenfield restatement of modal alignment
Eq.~\eqref{eq:ch3.spectral-index-alignment}: the label \(\alpha\) carries \((\ell,m)\) when
\(\mathbf{f}_\alpha=\mathbf{u}_{\ell m}\) satisfies Eqs.~\eqref{eq:ch3.J-simultaneous-eigenfield-system} and
\eqref{eq:ch3.J-ladder-action} \cite{stratton1941,harrington2001}.

\begin{proposition}[Joint spectral diagonalization on the radiative inner product]
\label{prop:ch3.J-simultaneous-diagonalization}
Assume \(\mathbf{J}_{z}\) and \(\mathbf{J}^{2}\) are essentially self-adjoint on the radiative closure of
\(\mathcal{D}(\mathbf{J})\) with respect to Eq.~\eqref{eq:ch3.radiative-inner-product-def}
\cite{kato1995,hansen1988}. Then there exists an orthonormal (in \(\langle\cdot,\cdot\rangle_{\mathrm{rad}}\))
basis of \(\mathcal{F}_{\mathrm{rad}}\) consisting of simultaneous eigenfields
Eq.~\eqref{eq:ch3.J-simultaneous-eigenfield-system}, and distinct \((\ell,m)\) sectors are orthogonal unless
degeneracy from additional symmetry is declared in Subsection~\ref{sec:ch323} \cite{jackson1999,stratton1941}.
\end{proposition}
\begin{proof}[Proof sketch]
Eq.~\eqref{eq:ch3.J-so3-commutation} implies \([\mathbf{J}^{2},\mathbf{J}_{z}]=0\). For commuting self-adjoint
operators on a Hilbert space, spectral theorem gives a joint eigenbasis \cite{kato1995,jackson1999}. The
radiative inner product inherits reciprocity symmetry Eq.~\eqref{eq:ch3.reciprocity-symmetric-inner-product},
so distinct \((\ell,m)\) eigenfields decouple in coefficient pairings used by
Eq.~\eqref{eq:ch3.riesz-coefficient-representation} \cite{stratton1941,hansen1988}.
% PROOF: AUTHOR REVIEW
\end{proof}

\paragraph{Orthogonality and degeneracy bookkeeping.}
Within a fixed \(\ell\) multiplet, distinct \(m\) sectors satisfy
\begin{equation}
\label{eq:ch3.J-m-orthogonality}
\langle \mathbf{u}_{\ell m},\mathbf{u}_{\ell m'}\rangle_{\mathrm{rad}}
=
0,
\qquad
m\neq m',
\end{equation}
when the joint basis of Proposition~\ref{prop:ch3.J-simultaneous-diagonalization} is chosen
\cite{hansen1988,stratton1941}. The \(2\ell+1\) fold degeneracy at fixed \(\ell\) is the representation
content imported from \(\mathrm{SO}(3)\) in Subsection~\ref{sec:ch323}; Subsection~\ref{sec:ch322} records only
the operator counting identity
\begin{equation}
\label{eq:ch3.J-multiplet-degeneracy-count}
\dim\ker\big(\mathbf{J}^{2}-\ell(\ell+1)\big)
=
2\ell+1
\end{equation}
on the physical integer-\(\ell\) sector \cite{jackson1999,hansen1988}. Equation
\eqref{eq:ch3.J-multiplet-degeneracy-count} is why a truncated expansion at cutoff \(\ell_{\max}\) contains
\(\sum_{\ell=0}^{\ell_{\max}}(2\ell+1)\) angular slots before radial quantum numbers are appended in
Section~\ref{sec:ch33} \cite{hansen1988,stratton1941}.

\paragraph{Reciprocity-compatible spectral pairing.}
Simultaneous eigenfields enter coefficient functionals through the Riesz map
Eq.~\eqref{eq:ch3.riesz-coefficient-representation}. When \(\mathbf{u}_{\ell m}\) and
\(\mathbf{v}_{\ell'm'}\) are joint eigenvectors of Eqs.~\eqref{eq:ch3.J-simultaneous-eigenfield-system},
\begin{equation}
\label{eq:ch3.J-eigenfield-coefficient-diagonal}
\mathcal{C}_{\ell m}[\mathbf{u}_{\ell m}]
=
\langle \mathbf{u}_{\ell m},\mathbf{u}_{\ell m}\rangle_{\mathrm{rad}},
\qquad
\mathcal{C}_{\ell m}[\mathbf{u}_{\ell'm'}]=0
\quad(m\neq m'\ \text{or}\ \ell\neq\ell'),
\end{equation}
on an orthonormal joint basis of Proposition~\ref{prop:ch3.J-simultaneous-diagonalization}
\cite{stratton1941,hansen1988}. Equation~\eqref{eq:ch3.J-eigenfield-coefficient-diagonal} is the angular
specialization of Parseval preview Eq.~\eqref{eq:ch3.parseval-radiative-preview}: distinct \((\ell,m)\) labels
decouple in exported coefficient energy once commutation and self-adjointness hold on
\(\mathcal{D}(\mathbf{J})\) \cite{hansen1988,stratton1941}. Numerical codes that nonorthogonalize \((\ell,m)\)
sectors before applying Eq.~\eqref{eq:ch3.eigenfield-expansion-on-frad} smear ladder partners and inflate
truncation remainder Eq.~\eqref{eq:ch3.eigenfield-truncation-remainder} \cite{harrington2001}.

\paragraph{Commutation with radiation export.}
Source rotation commutation Eq.~\eqref{eq:ch2.radiation-rotation-commutation} lifts to generator commutators
with \(\mathcal{R}_\omega\) on admissible domains:
\begin{equation}
\label{eq:ch3.J-commutation-radiation-intertwine}
\mathbf{J}_{a}\,\mathcal{R}_\omega[\mathfrak{D}]
=
\mathcal{R}_\omega\big[\mathbf{J}_{a}^{(\mathrm{src})}\mathfrak{D}\big],
\qquad
[\mathbf{J}_{a},\mathbf{J}_{b}]\,\mathcal{R}_\omega[\mathfrak{D}]
=
\mathcal{R}_\omega\big[[\mathbf{J}_{a}^{(\mathrm{src})},\mathbf{J}_{b}^{(\mathrm{src})}]\mathfrak{D}\big],
\end{equation}
whenever \(\mathcal{D}_{J,\mathrm{reg}}=1\) and outgoing closure is preserved \cite{stratton1941,harrington2001}.
Equation~\eqref{eq:ch3.J-commutation-radiation-intertwine} is the generator form of
Eq.~\eqref{eq:ch3.J-radiation-intertwining}: commutation relations survive export only on the same carrier where
Eq.~\eqref{eq:ch3.J-so3-commutation} is valid \cite{harrington2001,jackson1999}. Failure modes include
evaluating \([\mathbf{J}_{x},\mathbf{J}_{y}]\) on \(\mathcal{S}_\omega[\mathfrak{D}]\) near reactive shells
and reporting the residue as ``OAM coupling'' \cite{harrington2001,stratton1941}.

\begin{theorem}[Simultaneous eigenfield characterization]
\label{thm:ch3.simultaneous-eigenfield-char}
A nonzero \(\mathbf{u}\in\mathcal{D}(\mathbf{J})\cap\mathcal{F}_{\mathrm{rad}}\) is a simultaneous
angular-momentum eigenfield if and only if it satisfies Eqs.~\eqref{eq:ch3.J-simultaneous-eigenfield-system}
and is connected to the \(\ell\) multiplet by finitely many applications of
Eq.~\eqref{eq:ch3.J-ladder-action} starting from an extremal \(m=\pm\ell\) state \cite{jackson1999,hansen1988}.
For exported states, admissibility further requires \(\mathbf{u}=\mathcal{R}_\omega[\mathfrak{D}]\) with
\(\widehat{\mathbf{k}}\in\mathcal{K}_{\mathrm{acc}}^{\mathrm{ff}}\) whenever directional transport is claimed,
per Theorem~\ref{thm:ch3.J-labeling-admissibility} \cite{stratton1941,harrington2001}.
\end{theorem}
\begin{proof}[Proof sketch]
Standard \(\mathfrak{su}(2)\) ladder theory on the commutation relations
Eqs.~\eqref{eq:ch3.J-so3-commutation} and~\eqref{eq:ch3.J-z-ladder-commutation} \cite{jackson1999}. The export
clause imports accessibility and gauge gates already established in Subsection~\ref{sec:ch321}
\cite{stratton1941,harrington2001}.
% PROOF: AUTHOR REVIEW
\end{proof}

\paragraph{Validity domain.}
Eqs.~\eqref{eq:ch3.J-so3-commutation}--\eqref{eq:ch3.J-ladder-action} assume linear isotropic vacuum
(or piecewise-homogeneous regions with interface conditions already absorbed into \(\mathfrak{D}\)),
time-harmonic phasors, and outgoing \(\mathcal{B}_\Gamma\) \cite{stratton1941,jackson1999}. Absorbing media,
gyrotropy, and moving frames can preserve power invariance while breaking
Eq.~\eqref{eq:ch3.J-commutation-radiation-intertwine} \cite{stratton1941}. Half-integer \(\ell\) sectors
belong to spinor representations and are deferred to the inheritance ledger in Section~\ref{sec:ch3inherit}
\cite{jackson1999}.

\paragraph{Worked illustration (\(\ell=1\) triplet and ladder closure).}
The electric dipole multiplet at \(\ell=1\) spans \(m\in\{-1,0,+1\}\) with
Eq.~\eqref{eq:ch3.J-multiplet-degeneracy-count} giving dimension three \cite{jackson1999,hansen1988}. Starting
from a circularly polarized \(\ell=1\), \(m=+1\) far-field prototype, repeated application of
\(\mathbf{J}_{-}\) in Eq.~\eqref{eq:ch3.J-ladder-action} reaches \(m=0\) and \(m=-1\) partners with the expected
normalization factors \(\sqrt{2}\) and \(\sqrt{2}\), respectively \cite{jackson1999}. Orthogonality
Eq.~\eqref{eq:ch3.J-m-orthogonality} is what allows separate Riesz coefficients
Eq.~\eqref{eq:ch3.riesz-coefficient-representation} for each \(m\) without cross-talk in the radiative inner
product \cite{stratton1941,hansen1988}.

\paragraph{Worked illustration (commutator failure on reactive samples).}
Let \(\mathbf{u}_\omega=\mathcal{S}_\omega[\mathfrak{D}]\) be a reactive-shell restriction of a dipole state
with \(\|\mathbf{u}_\omega\|_{\mathcal{U}}<\infty\) but \(\mathbf{u}_\omega\notin\mathcal{D}(\mathbf{J})\) because
\(\mathbf{J}_{z}\mathbf{u}_\omega\notin\mathcal{F}_{\mathrm{rad}}\) \cite{harrington2001,stratton1941}. Evaluating
\([\mathbf{J}_{x},\mathbf{J}_{y}]\) on such a pair can produce a nonzero residue not proportional to
\(\mathbf{J}_{z}\mathbf{u}_\omega\), violating Eq.~\eqref{eq:ch3.J-so3-commutation}. Reporting that residue as
evidence of ``anomalous OAM coupling'' is a domain error: the commutator algebra is certified only on
\(\mathcal{D}(\mathbf{J})\) \cite{harrington2001,belinfante1940}.

\paragraph{Worked illustration (frame rotation versus ladder mixing).}
Rotating the laboratory frame by \(R(\varphi,\hat{\mathbf{z}})\) mixes \(m\) components within a fixed
\(\ell\) multiplet through representation matrices \(\rho(g)\) in
Eq.~\eqref{eq:ch3.coefficient-rotation-law} \cite{jackson1999,stratton1941}. That mixing is not ladder action:
\(\mathbf{J}_{\pm}\) connect \emph{eigenfields} of \(\mathbf{J}_{z}\), whereas a passive frame rotation
relabels coordinates before export \cite{stratton1941}. Confusing the two produces incorrect coupling matrices
in numerical beam models \cite{harrington2001,hansen1988}.

\paragraph{Worked illustration (Hermiticity test on a finite \(\ell=1\) block).}
Restrict attention to the three-dimensional \(\ell=1\) subspace spanned by joint eigenfields
\(\{\mathbf{u}_{1m}\}_{m=-1}^{1}\). In that block, \(\mathbf{J}_{z}\) is diagonal with entries
\(-1,0,+1\), while \(\mathbf{J}_{x}\) and \(\mathbf{J}_{y}\) are off-diagonal with the standard
\(\mathfrak{su}(2)\) realizations obtained from Eq.~\eqref{eq:ch3.J-plus-minus-def}
\cite{jackson1999,hansen1988}. The Hermiticity check
\(\langle \mathbf{u}_{1m},\mathbf{J}_{a}\mathbf{u}_{1m'}\rangle_{\mathrm{rad}}
=\langle \mathbf{J}_{a}\mathbf{u}_{1m},\mathbf{u}_{1m'}\rangle_{\mathrm{rad}}^{*}\)
fails if the inner product is taken on reactive samples but passes on outgoing
\(\mathcal{R}_\omega[\mathfrak{D}]\) when reciprocity symmetry
Eq.~\eqref{eq:ch3.reciprocity-symmetric-inner-product} is enforced \cite{stratton1941,harrington2001}. The
block exercise is the finite-dimensional shadow of Proposition~\ref{prop:ch3.J-simultaneous-diagonalization}:
commutation implies simultaneous diagonalizability; the radiative inner product supplies the metric in which
the diagonalization is unitary \cite{kato1995,hansen1988}.

\begin{table}[t]
\centering
\caption{Commutation and simultaneous-eigenfield diagnostics for Subsection~\ref{sec:ch322}.}
\label{tab:ch3.J-commutation-diagnostics}
\small
\begin{tabular}{@{}p{0.22\linewidth}p{0.34\linewidth}p{0.36\linewidth}@{}}
\toprule
Test & Passing criterion & Typical failure \\
\midrule
Lie algebra &
Eq.~\eqref{eq:ch3.J-so3-commutation} on \(\mathcal{D}(\mathbf{J})\) &
Reactive-shell commutator residue \\
\addlinespace
Joint spectra &
Eq.~\eqref{eq:ch3.J-simultaneous-eigenfield-system} &
\(\ell\) from phase fit only \\
\addlinespace
Ladder closure &
Eq.~\eqref{eq:ch3.J-ladder-action} reaches \(m=\pm\ell\) &
Non-integer \(m\) steps \\
\addlinespace
Export consistency &
Eq.~\eqref{eq:ch3.J-commutation-radiation-intertwine} &
Frame swap without \(\rho(g)\) \\
\bottomrule
\end{tabular}
\end{table}

Table~\ref{tab:ch3.J-commutation-diagnostics} complements Table~\ref{tab:ch3.J-reporting-gates}: the latter
guarded label admissibility; the present table guards \emph{algebraic} consistency of simultaneous
eigenanalysis on \(\mathcal{F}_{\mathrm{rad}}\) \cite{stratton1941,harrington2001}.

Subsection~\ref{sec:ch322} has derived the \(\mathrm{SO}(3)\) commutation relations
Eq.~\eqref{eq:ch3.J-so3-commutation}, installed ladder operators
Eqs.~\eqref{eq:ch3.J-plus-minus-def}--\eqref{eq:ch3.J-ladder-action}, and characterized simultaneous eigenfields
through Eqs.~\eqref{eq:ch3.J-simultaneous-eigenfield-system} and
Theorem~\ref{thm:ch3.simultaneous-eigenfield-char} \cite{jackson1999,hansen1988,belinfante1940,stratton1941}.
Subsection~\ref{sec:ch323} connects these multiplets to \(\mathrm{SO}(3)\) representation matrices and
spectral degeneracy on \(\mathcal{F}_{\mathrm{rad}}\), using rotation rules from Section~\ref{sec:ch01.3} and
Subsection~\ref{sec:ch262} by reference \cite{jackson1999,harrington2001}.


\subsection{SO(3) representations and degeneracy}
\label{sec:ch323}
% TOC tag: [HOME][USE SS1.3][USE SS2.6.2]

Repeated eigenvalues in radiation spectra are not numerical accidents: when
\(\mathcal{R}_\omega\) commutes with spatial rotations, degenerate modes organize into
\(\mathrm{SO}(3)\) irreducible blocks whose dimension is fixed by representation theory
\cite{jackson1999,stratton1941,harrington2001}. Subsection~\ref{sec:ch323} makes that
organization explicit on \(\mathcal{F}_{\mathrm{rad}}\): each \(\ell\) sector from
Subsection~\ref{sec:ch322} is an invariant subspace for
\(\mathcal{D}(R)\) Eq.~\eqref{eq:ch2.field-rotation-representation}, representation matrices
\(\rho^{(\ell)}(R)\) act on coefficient blocks through
Eq.~\eqref{eq:ch3.coefficient-rotation-law}, and partial symmetry reduces multiplicity below
the \(2\ell+1\) count of Theorem~\ref{thm:ch2.rotation-spectral-degeneracy}
\cite{hansen1988,harrington2001}. The counting identity \(2\ell+1\) is the same in
Subsections~\ref{sec:ch322} and~\ref{sec:ch262}; Subsection~\ref{sec:ch323} explains \emph{why} the integer
appears twice by identifying \(\mathcal{H}_{\ell}\) with an irreducible \(\mathrm{SO}(3)\) representation
\cite{jackson1999,stratton1941}.

\paragraph{Standing imports.}
Rotation actions Eqs.~\eqref{eq:ch1.active-rotation-action}--\eqref{eq:ch1.passive-rotation-action},
active--passive duality Eq.~\eqref{eq:ch1.active-passive-duality}, vector covariance
Eq.~\eqref{eq:ch1.vector-observable-covariance}, and tensor transformation rules from
Section~\ref{sec:ch01.3.2} enter by reference only \cite{stratton1941,jackson1999}. Radiation-side
rotation covariance Eqs.~\eqref{eq:ch2.field-rotation-representation}--\eqref{eq:ch2.radiation-rotation-commutation},
modal basis closure Eq.~\eqref{eq:ch2.modal-basis-covariance}, and degeneracy multiplicity
Eq.~\eqref{eq:ch2.spectral-degeneracy-multiplicity} are imported from
Subsection~\ref{sec:ch262} \cite{stratton1941,harrington2001}. Angular-momentum multiplets
Eqs.~\eqref{eq:ch3.J-simultaneous-eigenfield-system}--\eqref{eq:ch3.J-multiplet-degeneracy-count}
are imported from Subsection~\ref{sec:ch322} \cite{jackson1999,hansen1988}.

\paragraph{Irreducible \(\mathrm{SO}(3)\) sectors on \(\mathcal{F}_{\mathrm{rad}}\).}
For each \(\ell\in\{0,1,2,\ldots\}\) realized by classical multipoles, define the
\(\ell\)-sector subspace
\begin{equation}
\label{eq:ch3.ell-sector-subspace}
\mathcal{H}_{\ell}
=
\ker\big(\mathbf{J}^{2}-\ell(\ell+1)\big)\cap\mathcal{F}_{\mathrm{rad}},
\end{equation}
using Eq.~\eqref{eq:ch3.J-squared-def} \cite{jackson1999,hansen1988}. Equation
\eqref{eq:ch3.ell-sector-subspace} is the Chapter~\ref{ch:03} field-space realization of
\(\mathcal{H}_{\ell}\) in Theorem~\ref{thm:ch2.rotation-spectral-degeneracy}
\cite{stratton1941,harrington2001}. An \emph{irreducible representation} of \(\mathrm{SO}(3)\) on
\(\mathcal{H}_{\ell}\) is a homomorphism
\begin{equation}
\label{eq:ch3.SO3-irrep-homomorphism}
\rho^{(\ell)}:\mathrm{SO}(3)\longrightarrow \mathrm{GL}(\mathcal{H}_{\ell}),
\qquad
\rho^{(\ell)}(R_1R_2)=\rho^{(\ell)}(R_1)\rho^{(\ell)}(R_2),
\end{equation}
with \(\rho^{(\ell)}(R)=\mathcal{D}(R)\big|_{\mathcal{H}_{\ell}}\) whenever
Eq.~\eqref{eq:ch2.radiation-rotation-commutation} holds \cite{jackson1999,stratton1941}. Physically,
\(\mathcal{H}_{\ell}\) is the exported angular sector in which rotation mixes partners \emph{only}
among themselves and never with a different \(\ell\) label \cite{hansen1988,harrington2001}.

\paragraph{Invariance under field rotations.}
For \(R\in\mathrm{SO}(3)\) and \(\mathbf{u}\in\mathcal{H}_{\ell}\),
\begin{equation}
\label{eq:ch3.ell-sector-invariance}
\mathcal{D}(R)\,\mathbf{u}\in\mathcal{H}_{\ell},
\qquad
\mathbf{J}^{2}\big(\mathcal{D}(R)\mathbf{u}\big)
=
\ell(\ell+1)\,\mathcal{D}(R)\mathbf{u},
\end{equation}
because \(\mathcal{D}(R)\) is generated by \(\mathbf{J}\) through
Eq.~\eqref{eq:ch3.rotation-from-J-generators} and commutes with \(\mathbf{J}^{2}\) on
\(\mathcal{D}(\mathbf{J})\) \cite{jackson1999,stratton1941}. Equation~\eqref{eq:ch3.ell-sector-invariance}
is the representation-theoretic restatement of modal basis covariance
Eq.~\eqref{eq:ch2.modal-basis-covariance} when \(\mathbf{u}_m\in\mathcal{H}_{\ell}\) for all
\(m\in\{-\ell,\ldots,\ell\}\) \cite{harrington2001,hansen1988}. Coefficient blocks therefore transform
by finite-dimensional matrices rather than by ad hoc mixing rules:
\begin{equation}
\label{eq:ch3.ell-block-coefficient-rotation}
\mathbf{c}^{(\ell)}\big[\mathcal{A}_R\mathfrak{D}\big]
=
\rho^{(\ell)}(R)\,\mathbf{c}^{(\ell)}[\mathfrak{D}],
\qquad
\mathbf{c}^{(\ell)}=(c_{\ell,-\ell},\ldots,c_{\ell,\ell})^{\mathsf{T}},
\end{equation}
with \(\mathbf{c}^{(\ell)}\) the Riesz coefficient stack on the \(\ell\) multiplet
\cite{stratton1941,hansen1988}. Explicit matrix elements of \(\rho^{(\ell)}(R)\) are deferred to
Subsection~\ref{sec:ch324}; Subsection~\ref{sec:ch323} records only the block structure
\cite{jackson1999}.

\begin{figure}[t]
\centering
\ChOneFigBlock{%
\begin{tikzpicture}[ch1fig, node distance=7mm and 9mm]
  \node[ch1box] (frad) {$\mathcal{F}_{\mathrm{rad}}$};
  \node[ch1box, below=10mm of frad] (h0) {$\mathcal{H}_{0}$\\$\dim=1$};
  \node[ch1box, left=14mm of h0] (h1) {$\mathcal{H}_{1}$\\$\dim=3$};
  \node[ch1box, right=14mm of h0] (h2) {$\mathcal{H}_{2}$\\$\dim=5$};
  \node[ch1box, below=11mm of h1] (r0) {$\rho^{(0)}(R)$};
  \node[ch1box, below=11mm of h0] (r1) {$\rho^{(1)}(R)$};
  \node[ch1box, below=11mm of h2] (r2) {$\rho^{(2)}(R)$};
  \draw[ch1flow] (frad) -- (h0);
  \draw[ch1flow] (frad) -- (h1);
  \draw[ch1flow] (frad) -- (h2);
  \draw[ch1flow] (h0) -- (r0);
  \draw[ch1flow] (h1) -- (r1);
  \draw[ch1flow] (h2) -- (r2);
\end{tikzpicture}%
}
\caption{Subsection~\ref{sec:ch323} block decomposition: \(\mathcal{F}_{\mathrm{rad}}\) splits into
invariant \(\mathcal{H}_{\ell}\) sectors, each carrying an irreducible \(\rho^{(\ell)}(R)\) action.}
\label{fig:ch3.SO3-irrep-block-decomposition}
\end{figure}

Figure~\ref{fig:ch3.SO3-irrep-block-decomposition} is the sector diagram behind
Figure~\ref{fig:ch2.spectral-theory-roadmap}: Subsection~\ref{sec:ch262} recorded operator-level
degeneracy; the present figure records the corresponding invariant subspaces on outgoing fields
\cite{stratton1941,harrington2001}.

\begin{theorem}[\(\mathbf{J}\) multiplets are \(\mathrm{SO}(3)\) irreducible sectors]
\label{thm:ch3.J-irrep-degeneracy}
Assume Eq.~\eqref{eq:ch2.radiation-rotation-commutation} for all \(R\in\mathrm{SO}(3)\) on isotropic
exterior problems with outgoing \(\mathcal{B}_\Gamma\) \cite{stratton1941,harrington2001}. Then
\(\mathcal{H}_{\ell}\) in Eq.~\eqref{eq:ch3.ell-sector-subspace} is invariant under
\(\mathcal{D}(R)\), the simultaneous eigenfields
Eq.~\eqref{eq:ch3.J-simultaneous-eigenfield-system} form a basis of \(\mathcal{H}_{\ell}\), and
\begin{equation}
\label{eq:ch3.J-irrep-dimension-identity}
\dim\mathcal{H}_{\ell}
=
2\ell+1
=
\dim\rho^{(\ell)},
\end{equation}
coinciding with Eq.~\eqref{eq:ch2.spectral-degeneracy-multiplicity} and
Eq.~\eqref{eq:ch3.J-multiplet-degeneracy-count} \cite{jackson1999,hansen1988,stratton1941}.
Moreover, \(\mathcal{H}_{\ell}\) is irreducible: no proper subspace of \(\mathcal{H}_{\ell}\) is
invariant under all \(\mathcal{D}(R)\) \cite{jackson1999}.
\end{theorem}
\begin{proof}[Proof sketch]
Joint diagonalization of \(\mathbf{J}^{2}\) and \(\mathbf{J}_{z}\) from
Proposition~\ref{prop:ch3.J-simultaneous-diagonalization} gives an orthonormal
\(\{\mathbf{u}_{\ell m}\}\) basis with ladder connectivity
Eq.~\eqref{eq:ch3.J-ladder-action} \cite{jackson1999,hansen1988}. Invariance
Eq.~\eqref{eq:ch3.ell-sector-invariance} follows from commutation of \(\mathcal{D}(R)\) with
\(\mathbf{J}^{2}\). Dimension \(2\ell+1\) is the ladder span from \(m=-\ell\) to \(m=+\ell\)
\cite{jackson1999}. Irreducibility is the standard \(\mathfrak{su}(2)\) statement on a single
\(\ell\) sector \cite{jackson1999,hansen1988}. Agreement with
Theorem~\ref{thm:ch2.rotation-spectral-degeneracy} follows because both theorems count the same
\(\ell\)-indexed radiative multiplicity on \(S^{2}\) \cite{stratton1941,harrington2001}.
% PROOF: AUTHOR REVIEW
\end{proof}

\paragraph{Direct-sum decomposition of outgoing states.}
When only finitely many \(\ell\) sectors contribute at a declared cutoff \(\ell_{\max}\), exported states
admit the orthogonal decomposition
\begin{equation}
\label{eq:ch3.Frad-ell-direct-sum}
\mathcal{F}_{\mathrm{rad}}
\supseteq
\bigoplus_{\ell=0}^{\ell_{\max}}
\mathcal{H}_{\ell},
\qquad
\mathbf{u}
=
\sum_{\ell=0}^{\ell_{\max}}
\Pi_{\ell}\mathbf{u},
\end{equation}
with \(\Pi_{\ell}\) the spectral projector onto
Eq.~\eqref{eq:ch3.ell-sector-subspace} induced by \(\mathbf{J}^{2}\) \cite{jackson1999,hansen1988}.
Equation~\eqref{eq:ch3.Frad-ell-direct-sum} is the sector restatement of hybrid decomposition
Eq.~\eqref{eq:ch3.hybrid-spectral-decomposition}: discrete \(\ell\) blocks carry the degeneracy counted by
Theorem~\ref{thm:ch3.J-irrep-degeneracy}, while continuous directional measures attach to each outgoing fiber
through Eq.~\eqref{eq:ch3.continuous-spectral-measure} \cite{stratton1941,harrington2001}. Parseval on the
joint basis reads
\begin{equation}
\label{eq:ch3.multiplet-parseval}
\|\mathbf{u}\|_{\mathcal{U}}^{2}
=
\sum_{\ell,m}
|c_{\ell m}|^{2},
\qquad
c_{\ell m}=\langle \mathbf{u}_{\ell m},\mathbf{u}\rangle_{\mathrm{rad}},
\end{equation}
when \(\{\mathbf{u}_{\ell m}\}\) is orthonormal in the radiative inner product
\cite{stratton1941,hansen1988}. Equation~\eqref{eq:ch3.multiplet-parseval} is the angular specialization of
Eq.~\eqref{eq:ch3.parseval-radiative-preview}: rotation mixes coefficients \emph{within} each
\(\mathbf{c}^{(\ell)}\) block through Eq.~\eqref{eq:ch3.ell-block-coefficient-rotation} but preserves the
total \(\ell\)-sector energy \(\sum_{m}|c_{\ell m}|^{2}\) under
Eq.~\eqref{eq:ch2.spectral-power-rotation-invariance} \cite{stratton1941,harrington2001}.

\paragraph{Character orthogonality (preview).}
Irreducible matrix elements satisfy orthogonality relations with respect to the Haar measure \(d\mu(R)\) on
\(\mathrm{SO}(3)\):
\begin{equation}
\label{eq:ch3.irrep-character-orthogonality}
\int_{\mathrm{SO}(3)}
\mathrm{Tr}\,\rho^{(\ell)}(R)^{*}\,\rho^{(\ell')}(R)\,d\mu(R)
\propto
\delta_{\ell\ell'},
\end{equation}
so distinct \(\ell\) sectors are representation-theoretically orthogonal even before radial separation
\cite{jackson1999,hansen1988}. Equation~\eqref{eq:ch3.irrep-character-orthogonality} is the group-theoretic
reason \(\Pi_{\ell}\) projectors in Eq.~\eqref{eq:ch3.Frad-ell-direct-sum} do not mix \(\ell\) labels under
integration over all attitudes; Subsection~\ref{sec:ch324} records the finite-\(R\) matrix elements used in
laboratory frame rotations \cite{jackson1999}.

\paragraph{Radiation operator block structure.}
When \(\mathcal{R}_\omega\) commutes with \(\mathcal{D}(R)\) on a sector,
Eq.~\eqref{eq:ch2.radiation-D-commutation}, the restricted radiation map acts within each
\(\mathcal{H}_{\ell}\) block:
\begin{equation}
\label{eq:ch3.radiation-irrep-block}
\mathcal{R}_\omega\big[\mathcal{H}_{\ell}^{(\mathrm{src})}\big]
\subseteq
\mathcal{H}_{\ell},
\qquad
\mathcal{H}_{\ell}^{(\mathrm{src})}
=
\{\mathfrak{D}:\mathcal{R}_\omega[\mathfrak{D}]\in\mathcal{H}_{\ell}\},
\end{equation}
with \(\mathcal{H}_{\ell}^{(\mathrm{src})}\) the source-side preimage sector
\cite{stratton1941,harrington2001}. Equation~\eqref{eq:ch3.radiation-irrep-block} explains repeated
\(\lambda_n\) in Theorem~\ref{thm:ch2.radiation-spectral-decomposition}: eigenvalues tied to the
same \(\ell\) label share a symmetry block but remain distinct modes through the \(m\) index and
radial quantum numbers appended in Section~\ref{sec:ch33} \cite{hansen1988,stratton1941}. On
\(\mathcal{H}_{\ell}\), Schur's lemma implies that any operator commuting with all
\(\rho^{(\ell)}(R)\)---including isotropic \(\mathcal{R}_\omega\) when fully symmetric---acts as a
scalar on each irreducible block up to radial coupling \cite{jackson1999,harrington2001}.

\paragraph{Equivalence of representation and angular-momentum labels.}
The bridge between representation matrices and angular-momentum quantum numbers is
\begin{equation}
\label{eq:ch3.representation-J-equivalence}
\rho^{(\ell)}(R)\,\mathbf{c}^{(\ell)}
=
\mathbf{c}^{(\ell)}\big[T_R\mathbf{u}_{\ell m}\big],
\qquad
T_R=\exp\!\big(-\jj\,\varphi\,\mathbf{J}_{\hat{\mathbf{u}}}\big),
\end{equation}
for \(R=R(\varphi,\hat{\mathbf{u}})\) from Eq.~\eqref{eq:ch1.rodrigues-rotation}, with
\(\mathbf{u}_{\ell m}\) simultaneous eigenfields Eq.~\eqref{eq:ch3.J-simultaneous-eigenfield-system}
\cite{jackson1999,stratton1941}. Equation~\eqref{eq:ch3.representation-J-equivalence} identifies
\(\rho^{(\ell)}\) in Eq.~\eqref{eq:ch3.coefficient-rotation-law} with the finite-dimensional
restriction of Eq.~\eqref{eq:ch3.rotation-from-J-generators}; Subsection~\ref{sec:ch324} supplies
closed-form \(D^{(\ell)}_{m'm}(R)\) matrix elements \cite{hansen1988,jackson1999}. Laboratory readers
should treat \(\ell\) as a representation label and \(m\) as a weight label within
\(\rho^{(\ell)}\), not as independent harmonic-fit indices \cite{stratton1941,harrington2001}.

\paragraph{Partial symmetry and reduced multiplicity.}
Real antennas rarely realize full \(\mathrm{SO}(3)\): geometry, ground planes, and feeds reduce the
symmetry group to a subgroup \(G\subset\mathrm{SO}(3)\) \cite{harrington2001,stratton1941}. On an
invariant subspace \(\mathcal{H}_{\ell}^{(G)}\subseteq\mathcal{H}_{\ell}\),
\begin{equation}
\label{eq:ch3.partial-symmetry-multiplicity}
\dim\mathcal{H}_{\ell}^{(G)}
=
\sum_{i\in I_{\ell}(G)} d_i,
\qquad
d_i=\dim\rho_i,
\end{equation}
where \(\{\rho_i\}\) are the \(G\)-irreducible components obtained by decomposing \(\rho^{(\ell)}\)
\cite{jackson1999,harrington2001}. Equation~\eqref{eq:ch3.partial-symmetry-multiplicity} can be
strictly less than \(2\ell+1\): a line source aligned with \(\hat{\mathbf{z}}\) retains axial
symmetry about \(\hat{\mathbf{z}}\) but not full \(\mathrm{SO}(3)\), splitting the \(\ell=1\)
triplet into visible and suppressed \(m\) sectors through accessibility
Eq.~\eqref{eq:ch2.far-angular-accessibility-def} \cite{harrington2001,stratton1941}. Complete
symmetry-breaking catalogs belong to Section~\ref{sec:ch36}; Subsection~\ref{sec:ch323} records only
the representation-theoretic counting rule \cite{hansen1988}.

\paragraph{Quadrupole \(\mathcal{H}_{2}\) sector and mode budgeting.}
Theorem~\ref{thm:ch3.J-irrep-degeneracy} gives \(\dim\mathcal{H}_{2}=5\) for electric quadrupole content
\cite{jackson1999,hansen1988}. A five-term angular fit is therefore the minimal \emph{full} \(\mathrm{SO}(3)\)
budget at \(\ell=2\), not an optional refinement \cite{hansen1988,stratton1941}. Truncating at three terms
because only three Cartesian quadrupole tensor components were measured ignores that those components generate
one \(\rho^{(2)}\) block, not three independent \(\ell\) labels \cite{harrington2001}. Mode-budget tables in
antenna synthesis should use \(\sum_{\ell\le\ell_{\max}}(2\ell+1)\) from
Eq.~\eqref{eq:ch3.J-irrep-dimension-identity} before appending radial indices from Section~\ref{sec:ch33}
\cite{hansen1988,harrington2001}.

\begin{figure}[t]
\centering
\ChOneFigBlock{%
\begin{tikzpicture}[ch1fig, node distance=8mm and 11mm]
  \node[ch1box] (full) {Full $\mathrm{SO}(3)$\\$\dim\mathcal{H}_\ell=2\ell+1$};
  \node[ch1box, below=11mm of full] (g) {Subgroup $G$};
  \node[ch1box, left=15mm of g] (v) {Visible block\\$\mathcal{H}_\ell^{(G)}$};
  \node[ch1box, right=15mm of g] (s) {Suppressed block\\$\mathcal{K}_{\mathrm{sup}}$};
  \draw[ch1flow] (full) -- (g);
  \draw[ch1flow] (g) -- (v);
  \draw[ch1flow] (g) -- (s);
\end{tikzpicture}%
}
\caption{Partial symmetry reduction Eq.~\eqref{eq:ch3.partial-symmetry-multiplicity}: full
\(\mathcal{H}_{\ell}\) splits under \(G\subset\mathrm{SO}(3)\) into visible and suppressed representation
subblocks tied to Eq.~\eqref{eq:ch2.far-angular-accessibility-def}.}
\label{fig:ch3.partial-symmetry-split}
\end{figure}

Figure~\ref{fig:ch3.partial-symmetry-split} is the degeneracy diagram for imperfect antennas: multiplicity is
not removed by physics, but partitioned into accessible and suppressed representation components
\cite{harrington2001,stratton1941}.

\begin{proposition}[Simultaneous block diagonalization of \(\mathcal{R}_\omega\) and \(\mathcal{D}(R)\)]
\label{prop:ch3.simultaneous-block-diagonalization}
If \(G\subset\mathrm{SO}(3)\) leaves the problem invariant and
Eq.~\eqref{eq:ch2.radiation-rotation-commutation} holds for all \(R\in G\), then
\(\mathcal{F}_{\mathrm{rad}}\) admits a basis in which \(\mathcal{D}(R)\) is block diagonal on
\(G\)-irreducible sectors and \(\mathcal{R}_\omega\) is block upper-triangular with blocks indexed by
\((\ell,\text{radial }n)\) \cite{stratton1941,harrington2001}. Within each full \(\mathrm{SO}(3)\)
\(\mathcal{H}_{\ell}\) sector of Theorem~\ref{thm:ch3.J-irrep-degeneracy}, radiation eigenvalues repeat
with multiplicity \(2\ell+1\) before radial splitting \cite{hansen1988,jackson1999}.
\end{proposition}
\begin{proof}[Proof sketch]
Simultaneous spectral theory for commuting families \(\{\mathcal{D}(R),\mathcal{R}_\omega\mathcal{S}_\omega\}\)
on the outgoing closure \cite{kato1995,stratton1941}. Multiplicity identification follows
Theorem~\ref{thm:ch3.J-irrep-degeneracy} and Theorem~\ref{thm:ch2.rotation-spectral-degeneracy}
\cite{harrington2001}.
% PROOF: AUTHOR REVIEW
\end{proof}

\paragraph{Continuous angular spectra and degeneracy.}
Exterior outgoing problems carry continuous \(S^{2}\) content through
Eq.~\eqref{eq:ch3.continuous-spectral-measure} \cite{stratton1941,hansen1988}. Degeneracy then
appears as \emph{fiber} multiplicity: each propagation direction \(\hat{\mathbf{k}}\) hosts the same
\(\ell\)-indexed block structure transverse to \(\hat{\mathbf{k}}\), while the continuous measure
\(d\mu(\hat{\mathbf{k}})\) records directional weight \cite{hansen1988,stratton1941}. Equation
\eqref{eq:ch3.hybrid-spectral-decomposition} is the continuous analogue of
Eq.~\eqref{eq:ch3.radiation-irrep-block}: discrete \(\ell\) sectors attach to each point of the
angular spectrum without conflating cavity modal bookkeeping from
Eq.~\eqref{eq:ch2.cavity-discrete-radiation-spectrum} \cite{harrington2001,jackson1999}.

\paragraph{Alignment with Chapter~\ref{ch:02} modal rotation law.}
Subsection~\ref{sec:ch262} recorded coefficient mixing
Eq.~\eqref{eq:ch2.modal-coefficient-rotation} and basis covariance
Eq.~\eqref{eq:ch2.modal-basis-covariance} at the operator level \cite{stratton1941,harrington2001}.
Subsection~\ref{sec:ch323} identifies the block matrices as irreducible representations:
\begin{equation}
\label{eq:ch3.modal-D-rad-identification}
D^{(\mathrm{rad})}_{nm}(R)
=
\big[\rho^{(\ell)}(R)\big]_{mm'}
\quad\text{when }n=(\ell,m),\ n'=(\ell,m'),
\end{equation}
on sectors where \(\mathbf{u}_n=\mathbf{u}_{\ell m}\) \cite{hansen1988,jackson1999}. Equation
\eqref{eq:ch3.modal-D-rad-identification} closes the loop begun in
Eq.~\eqref{eq:ch3.coefficient-rotation-law}: \(\rho(g)\) is not an abstract placeholder but the
\(\ell\)-block of physical rotation on \(\mathcal{F}_{\mathrm{rad}}\) \cite{stratton1941}. Eigenvalues
\(\lambda_n\) from Eq.~\eqref{eq:ch2.radiation-eigenpair-def} are constant on each degenerate multiplet
because \(\mathcal{R}_\omega\) commutes with \(\mathcal{D}(R)\) before radial splitting introduces additional
indices \cite{harrington2001,hansen1988}. Laboratory teams comparing measured \(a_n(\hat{\mathbf{k}})\) from
Eq.~\eqref{eq:ch2.spectral-far-pattern-expansion} to simulated \(\mathbf{c}^{(\ell)}\) must rotate both sides
with the same \(D^{(\mathrm{rad})}_{nm}(R)\) convention from
Eqs.~\eqref{eq:ch1.active-passive-duality} and~\eqref{eq:ch3.ell-block-coefficient-rotation}
\cite{stratton1941,harrington2001}.

\paragraph{Spectral projector lift to \(\ell\) sectors.}
Let \(P_{\ell}\) denote the orthogonal projector onto \(\mathcal{H}_{\ell}\) in the radiative inner product.
Then
\begin{equation}
\label{eq:ch3.ell-projector-def}
P_{\ell}\mathbf{u}
=
\sum_{m=-\ell}^{\ell}
\langle \mathbf{u}_{\ell m},\mathbf{u}\rangle_{\mathrm{rad}}\,\mathbf{u}_{\ell m},
\qquad
\sum_{\ell}P_{\ell}=\Pi_{\mathrm{rad}}
\end{equation}
on the span of counted sectors, with \(\Pi_{\mathrm{rad}}\) from
Eq.~\eqref{eq:ch3.outgoing-inheritance-criterion} \cite{hansen1988,stratton1941}. Equation
\eqref{eq:ch3.ell-projector-def} is the \(\ell\)-resolved form of continuous projector
Eq.~\eqref{eq:ch3.continuous-projector-rigged}: degeneracy is resolved by \(m\) inside each \(P_{\ell}\), not
by inventing new \(\ell\) labels for coordinate artifacts \cite{stratton1941,harrington2001}. Numerical
pipelines that apply \(P_{\ell}\) only after far-zone extraction remain consistent with
Eq.~\eqref{eq:ch3.radiation-irrep-block}; applying harmonic fits before \(\Pi_{\mathrm{rad}}\) does not
\cite{harrington2001}.

\paragraph{Validity domain.}
Eqs.~\eqref{eq:ch3.ell-sector-subspace}--\eqref{eq:ch3.partial-symmetry-multiplicity} assume isotropic
constitutive response commuting with \(\mathcal{D}(R)\) in the sense of Section~\ref{sec:ch01.3.2},
time-harmonic phasors, and outgoing closure \cite{stratton1941,jackson1999}. Anisotropic radomes,
magnetized plasma, and moving platforms may preserve scalar power
Eq.~\eqref{eq:ch2.spectral-power-rotation-invariance} while breaking irreducible block closure
\cite{stratton1941,harrington2001}. Half-integer spinor sectors remain in the inheritance ledger
Section~\ref{sec:ch3inherit} \cite{jackson1999}.

\paragraph{Worked illustration (dipole \(\mathcal{H}_{1}\) block).}
A centered electric dipole exports dominantly in \(\mathcal{H}_{1}\) with
Eq.~\eqref{eq:ch3.J-irrep-dimension-identity} giving \(\dim\mathcal{H}_{1}=3\)
\cite{jackson1999,hansen1988}. Coefficient rotation in
Eq.~\eqref{eq:ch3.ell-block-coefficient-rotation} mixes the three \(m\) partners without changing
\(\ell\); total power remains invariant under
Eq.~\eqref{eq:ch2.spectral-power-rotation-invariance} \cite{stratton1941}. Attempting to assign three
independent \(\ell\) labels to the three Cartesian dipole components violates
Theorem~\ref{thm:ch3.J-irrep-degeneracy}: the components are coordinates within one \(\rho^{(1)}\) block,
not three unrelated sectors \cite{harrington2001,hansen1988}.

\paragraph{Worked illustration (line source and axial reduction).}
A long wire along \(\hat{\mathbf{z}}\) preserves rotations about \(\hat{\mathbf{z}}\) but not full
\(\mathrm{SO}(3)\) \cite{harrington2001,stratton1941}. The \(\ell=1\) block decomposes under the
axial subgroup into visible \(m=0\) and \(m=\pm1\) sectors with different accessibility weights through
Eq.~\eqref{eq:ch2.far-angular-accessibility-def}, realizing
Eq.~\eqref{eq:ch3.partial-symmetry-multiplicity} with \(d_i<3\) on measurable subspaces
\cite{harrington2001}. Declaring ``degeneracy removal'' without naming the reduced group \(G\) is
inadmissible reporting \cite{stratton1941}.

\paragraph{Worked illustration (cavity versus exterior multiplicity).}
A bounded cavity spectrum Eq.~\eqref{eq:ch2.cavity-discrete-radiation-spectrum} can exhibit closely spaced
\(\lambda_n\) that mimic \(\mathrm{SO}(3)\) degeneracy without carrying a pure \(\ell\) label on
\(\mathcal{F}_{\mathrm{rad}}\) \cite{harrington2001,stratton1941}. Exterior outgoing problems instead tie
multiplicity to Eq.~\eqref{eq:ch3.J-irrep-dimension-identity} on \(\mathcal{H}_{\ell}\) after
\(\Pi_{\mathrm{rad}}\) \cite{stratton1941}. Confusing storage resonances with export degeneracy violates
the domain split in Eq.~\eqref{eq:ch2.discrete-continuous-spectrum-preview} \cite{harrington2001}.

\paragraph{Worked illustration (numerical block mixing).}
Finite-element exports that do not enforce Eq.~\eqref{eq:ch3.ell-sector-invariance} numerically produce
\(\rho^{(\ell)}(R)\) matrices with off-\(\ell\) leakage \cite{harrington2001}. The symptom is
\(\ell\)-dependent cross-talk in coupling matrices even when the continuous target satisfies
Theorem~\ref{thm:ch3.J-irrep-degeneracy} \cite{stratton1941,hansen1988}. Block-diagonalizing by
\(\mathbf{J}^{2}\) and \(\mathbf{J}_{z}\) before far-zone projection is the certified repair
\cite{hansen1988,harrington2001}.

\paragraph{Worked illustration (power redistribution within a multiplet).}
Under active rotation Eq.~\eqref{eq:ch1.active-rotation-action}, coefficient magnitudes
\(|c_{\ell m}|\) in Eq.~\eqref{eq:ch3.multiplet-parseval} redistribute while
\(\sum_{m}|c_{\ell m}|^{2}\) remains fixed for each \(\ell\) \cite{stratton1941,jackson1999}. A laboratory
report that tracks only one \(m\) component after rotating the source without applying
Eq.~\eqref{eq:ch3.ell-block-coefficient-rotation} will show apparent ``loss'' of angular content; the loss is a
coordinate artifact, not a failure of degeneracy \cite{harrington2001,hansen1988}.

\begin{table}[t]
\centering
\caption{Degeneracy classification on \(\mathcal{F}_{\mathrm{rad}}\) for Subsection~\ref{sec:ch323}.}
\label{tab:ch3.SO3-degeneracy-classification}
\small
\begin{tabular}{@{}p{0.20\linewidth}p{0.34\linewidth}p{0.38\linewidth}@{}}
\toprule
Class & Counting law & Misidentification risk \\
\midrule
Full \(\mathrm{SO}(3)\) &
Eq.~\eqref{eq:ch3.J-irrep-dimension-identity} &
Cartesian dipole triple counted as three \(\ell\) \\
\addlinespace
Partial \(G\subset\mathrm{SO}(3)\) &
Eq.~\eqref{eq:ch3.partial-symmetry-multiplicity} &
Sidelobe declared without naming \(G\) \\
\addlinespace
Continuous \(S^{2}\) &
Fiber \(\mathcal{H}_{\ell}\) per \(\hat{\mathbf{k}}\) &
Plane-wave angle mistaken for \(m\) \\
\addlinespace
Cavity storage &
Eq.~\eqref{eq:ch2.cavity-discrete-radiation-spectrum} &
Interior resonance labeled export degeneracy \\
\bottomrule
\end{tabular}
\end{table}

Table~\ref{tab:ch3.SO3-degeneracy-classification} is a multiplicity ledger, not a repeat of
Table~\ref{tab:ch3.J-commutation-diagnostics}: the latter tested Lie algebra consistency; the present
table classifies \emph{how many} partners belong to each symmetry sector \cite{stratton1941,harrington2001}.

Subsection~\ref{sec:ch323} has identified invariant sectors
Eq.~\eqref{eq:ch3.ell-sector-subspace}, defined irreducible actions
Eq.~\eqref{eq:ch3.SO3-irrep-homomorphism}, proved dimension identity
Eq.~\eqref{eq:ch3.J-irrep-dimension-identity} in Theorem~\ref{thm:ch3.J-irrep-degeneracy}, and recorded
partial-symmetry reduction Eq.~\eqref{eq:ch3.partial-symmetry-multiplicity} on \(\mathcal{F}_{\mathrm{rad}}\)
\cite{jackson1999,hansen1988,stratton1941,harrington2001}. Subsection~\ref{sec:ch324} supplies explicit
Wigner rotation matrices for \(\rho^{(\ell)}(R)\) and irreducible coupling coefficients used in
Eq.~\eqref{eq:ch3.ell-block-coefficient-rotation} \cite{jackson1999,hansen1988}.


\subsection{Wigner rotations and irreducible coupling}
\label{sec:ch324}
% TOC tag: [HOME][USE SS1.3.1]

Subsections~\ref{sec:ch322}--\ref{sec:ch323} identified invariant sectors \(\mathcal{H}_{\ell}\) and block
rotations \(\rho^{(\ell)}(R)\) on \(\mathcal{F}_{\mathrm{rad}}\) \cite{jackson1999,hansen1988}. Subsection
\ref{sec:ch324} supplies the explicit finite-dimensional matrices: Wigner \(D\)-matrices realize
\(\rho^{(\ell)}(R)\) in the \(\{\mathbf{u}_{\ell m}\}\) basis, and Clebsch--Gordan coefficients couple
distinct \(\ell\) sectors when sources superpose \cite{jackson1999,hansen1988,stratton1941}. The output closes
the preview in Eq.~\eqref{eq:ch3.dipole-coefficient-rotation} and identifies
\(D^{(\mathrm{rad})}_{nm}(R)\) from Eq.~\eqref{eq:ch2.modal-coefficient-rotation}
\cite{stratton1941,harrington2001}.

\paragraph{Standing imports.}
Active and passive rotations Eqs.~\eqref{eq:ch1.active-rotation-action}--\eqref{eq:ch1.passive-rotation-action},
duality Eq.~\eqref{eq:ch1.active-passive-duality}, group composition
Eq.~\eqref{eq:ch1.rotation-group-action}, and Rodrigues form Eq.~\eqref{eq:ch1.rodrigues-rotation} enter from
Subsection~\ref{sec:ch01.3.1} by reference only \cite{stratton1941,jackson1999}. Field rotation
Eq.~\eqref{eq:ch2.field-rotation-representation}, modal mixing
Eqs.~\eqref{eq:ch2.modal-coefficient-rotation}--\eqref{eq:ch2.modal-basis-covariance}, and block structure
Eqs.~\eqref{eq:ch3.ell-block-coefficient-rotation}--\eqref{eq:ch3.modal-D-rad-identification} are imported
from Subsections~\ref{sec:ch262} and~\ref{sec:ch323} \cite{stratton1941,harrington2001,hansen1988}.

\paragraph{Wigner \(D\)-matrices as irreducible matrix elements.}
Let \(\{\mathbf{u}_{\ell m}\}\) be the orthonormal joint basis of
Eq.~\eqref{eq:ch3.J-simultaneous-eigenfield-system} on \(\mathcal{H}_{\ell}\)
Eq.~\eqref{eq:ch3.ell-sector-subspace} \cite{jackson1999,hansen1988}. For \(R\in\mathrm{SO}(3)\), the
\emph{Wigner \(D\)-matrix} of rank \(\ell\) is
\begin{equation}
\label{eq:ch3.Wigner-D-matrix-def}
\big[D^{(\ell)}(R)\big]_{m'm}
=
\langle \mathbf{u}_{\ell m'},\,\mathcal{D}(R)\,\mathbf{u}_{\ell m}\rangle_{\mathrm{rad}},
\qquad
m,m'\in\{-\ell,\ldots,\ell\},
\end{equation}
with \(\mathcal{D}(R)\) from Eq.~\eqref{eq:ch2.field-rotation-representation} \cite{jackson1999,hansen1988}.
Equation~\eqref{eq:ch3.Wigner-D-matrix-def} is the Subsection~\ref{sec:ch324} \emph{HOME} tabulation of
\(\rho^{(\ell)}(R)\) Eq.~\eqref{eq:ch3.SO3-irrep-homomorphism} in the angular-momentum basis fixed by
Subsection~\ref{sec:ch322} \cite{jackson1999}. Matrix elements are finite because \(\dim\mathcal{H}_{\ell}=2\ell+1\)
\cite{hansen1988,stratton1941}.

\paragraph{Coefficient transformation law.}
The general rotation law for exported coefficients is
\begin{equation}
\label{eq:ch3.Wigner-D-coefficient-law}
c'_{\ell m}
=
\sum_{m'=-\ell}^{\ell}
D^{(\ell)}_{m m'}(R)\,c_{\ell m'},
\qquad
\mathbf{c}^{(\ell)\prime}=\rho^{(\ell)}(R)\,\mathbf{c}^{(\ell)},
\end{equation}
with \(c_{\ell m}=\langle \mathbf{u}_{\ell m},\mathbf{u}\rangle_{\mathrm{rad}}\) on
\(\mathbf{u}=\mathcal{R}_\omega[\mathfrak{D}]\) \cite{jackson1999,stratton1941}. Equation
\eqref{eq:ch3.Wigner-D-coefficient-law} elevates preview Eq.~\eqref{eq:ch3.dipole-coefficient-rotation} from
Subsection~\ref{sec:ch301} to all \(\ell\) sectors and identifies
Eq.~\eqref{eq:ch3.coefficient-rotation-law} with \(\rho(g)=D^{(\ell)}(g)\) on each block
\cite{hansen1988,harrington2001}. Physically, \(D^{(\ell)}_{m m'}(R)\) is the linear map that tells how a
laboratory rotation redistributes power among \(m\) partners at fixed \(\ell\) without creating new \(\ell\)
content \cite{stratton1941,jackson1999}.

\paragraph{Homomorphism and active--passive conjugation.}
Wigner matrices inherit group multiplication from Eq.~\eqref{eq:ch1.rotation-group-action}:
\begin{equation}
\label{eq:ch3.Wigner-D-homomorphism}
D^{(\ell)}(R_1R_2)
=
D^{(\ell)}(R_1)\,D^{(\ell)}(R_2),
\qquad
D^{(\ell)}(I)=I_{2\ell+1},
\end{equation}
on every \(\mathcal{H}_{\ell}\) where Eq.~\eqref{eq:ch2.radiation-rotation-commutation} holds
\cite{jackson1999,stratton1941}. Passive relabeling conjugates coefficients through
Eq.~\eqref{eq:ch1.active-passive-duality}:
\begin{equation}
\label{eq:ch3.Wigner-D-active-passive}
D^{(\ell)}_{\mathrm{passive}}(R)
=
D^{(\ell)}(R^{-1})^{\mathsf{T}},
\end{equation}
so reports in active and passive conventions differ by inverse-transpose pairing, not by altered radiative
physics \cite{stratton1941,jackson1999}. Equations~\eqref{eq:ch3.Wigner-D-homomorphism}--\eqref{eq:ch3.Wigner-D-active-passive}
are the finite-dimensional shadows of Eqs.~\eqref{eq:ch1.rotation-group-action} and
\eqref{eq:ch1.active-passive-duality} on coefficient stacks \cite{stratton1941}.

\begin{figure}[t]
\centering
\ChOneFigBlock{%
\begin{tikzpicture}[ch1fig, node distance=7mm and 10mm]
  \node[ch1box] (c) {$\mathbf{c}^{(\ell)}$};
  \node[ch1box, right=13mm of c] (d) {$D^{(\ell)}(R)$};
  \node[ch1box, right=13mm of d] (cp) {$\mathbf{c}^{(\ell)\prime}$};
  \node[ch1box, below=10mm of c] (ar) {$\mathcal{A}_R$ on source};
  \node[ch1box, below=10mm of cp] (dr) {$\mathcal{D}(R)$ on field};
  \draw[ch1flow] (c) -- (d) -- (cp);
  \draw[ch1flow] (ar) -| (c);
  \draw[ch1flow] (dr) -| (cp);
\end{tikzpicture}%
}
\caption{Subsection~\ref{sec:ch324} Wigner action Eq.~\eqref{eq:ch3.Wigner-D-coefficient-law}: active source
rotation and field rotation \(\mathcal{D}(R)\) induce the same block map \(D^{(\ell)}(R)\) on
\(\mathbf{c}^{(\ell)}\) when Eq.~\eqref{eq:ch2.radiation-rotation-commutation} holds.}
\label{fig:ch3.Wigner-D-coefficient-map}
\end{figure}

Figure~\ref{fig:ch3.Wigner-D-coefficient-map} is the coefficient-level companion to
Figure~\ref{fig:ch2.rotation-spectral-covariance}: operator commutation on the left becomes matrix
multiplication on the right \cite{stratton1941,harrington2001}.

\paragraph{Euler-angle parameterization (Volume~I convention).}
Every \(R\in\mathrm{SO}(3)\) admits a \(z\)--\(y\)--\(z\) Euler decomposition
\(R=R_z(\alpha)\,R_y(\beta)\,R_z(\gamma)\) with angles \((\alpha,\beta,\gamma)\)
\cite{jackson1999,stratton1941}. Wigner matrices factor accordingly as
\begin{equation}
\label{eq:ch3.Wigner-Euler-factorization}
D^{(\ell)}(\alpha,\beta,\gamma)
=
e^{-\jj m'\alpha}\,d^{(\ell)}_{m'm}(\beta)\,e^{-\jj m\gamma},
\end{equation}
where \(d^{(\ell)}_{m'm}(\beta)\) is the small-\(d\) reduced matrix element tabulated in
\cite{jackson1999,hansen1988}. Equation~\eqref{eq:ch3.Wigner-Euler-factorization} is the
Subsection~\ref{sec:ch324} working parameterization: explicit numbers are read from standard tables rather than
re-derived here \cite{hansen1988}. The phase conventions in Eq.~\eqref{eq:ch3.Wigner-Euler-factorization} must
match the \(\mathbf{u}_{\ell m}\) phase choice adopted in Section~\ref{sec:ch34}; Section~\ref{sec:ch38}
records basis-change rules that preserve Eq.~\eqref{eq:ch3.Wigner-D-coefficient-law}
\cite{hansen1988,stratton1941}.

\paragraph{Unitarity on each \(\ell\) block.}
Because \(\mathcal{D}(R)\) is unitary on \(\mathcal{H}_{\ell}\) under
Eq.~\eqref{eq:ch3.radiative-inner-product-def},
\begin{equation}
\label{eq:ch3.Wigner-D-unitarity}
\sum_{m}D^{(\ell)}_{m'm}(R)\,\big[D^{(\ell)}_{m''m}(R)\big]^{*}
=
\delta_{m'm''},
\qquad
D^{(\ell)}(R^{-1})=\big[D^{(\ell)}(R)\big]^{\dagger},
\end{equation}
for all \(R\in\mathrm{SO}(3)\) \cite{jackson1999,hansen1988}. Equation~\eqref{eq:ch3.Wigner-D-unitarity} is the
matrix form of power invariance Eq.~\eqref{eq:ch2.spectral-power-rotation-invariance} on a single multiplet:
\(\sum_m|c_{\ell m}|^{2}\) is unchanged under Eq.~\eqref{eq:ch3.Wigner-D-coefficient-law}
\cite{stratton1941,harrington2001}. Normalized coupling coefficients inherit the same metric, which is why
Eq.~\eqref{eq:ch3.CG-coupling-def} is stated with orthonormal \(\mathbf{u}_{\ell m}\) rather than with raw
Cartesian multipole amplitudes \cite{hansen1988,stratton1941}.

\paragraph{Reduced \(\ell=1\) matrices (reference block).}
The dipole block is the most common laboratory target. With Condon--Shortley phases and
Eq.~\eqref{eq:ch3.Wigner-Euler-factorization},
\begin{equation}
\label{eq:ch3.Wigner-d1-beta}
d^{(1)}_{m'm}(\beta)
=
\begin{pmatrix}
\frac{1+\cos\beta}{2} & -\frac{\sin\beta}{\sqrt{2}} & \frac{1-\cos\beta}{2} \\[4pt]
\frac{\sin\beta}{\sqrt{2}} & \cos\beta & -\frac{\sin\beta}{\sqrt{2}} \\[4pt]
\frac{1-\cos\beta}{2} & \frac{\sin\beta}{\sqrt{2}} & \frac{1+\cos\beta}{2}
\end{pmatrix},
\end{equation}
for ordering \((m',m)=(+1,0,-1)\) \cite{jackson1999,hansen1988}. Equation~\eqref{eq:ch3.Wigner-d1-beta} is the
executable form of Eq.~\eqref{eq:ch3.Wigner-D-coefficient-law} when \(\ell=1\): a \(\hat{\mathbf{z}}\) dipole
tilted by polar angle \(\beta\) populates all three \(m\) slots with weights drawn from the middle row and column
\cite{jackson1999,stratton1941}. Numerical beam codes that hard-code only \(m=0\) after mechanical tilt violate
Eq.~\eqref{eq:ch3.Wigner-d1-beta} even when total power remains invariant under
Eq.~\eqref{eq:ch2.spectral-power-rotation-invariance} \cite{harrington2001}.

\begin{theorem}[Wigner matrices realize \(\rho^{(\ell)}\) on \(\mathcal{H}_{\ell}\)]
\label{thm:ch3.Wigner-D-realizes-rho}
Under the hypotheses of Theorem~\ref{thm:ch3.J-irrep-degeneracy},
\(\rho^{(\ell)}(R)=D^{(\ell)}(R)\) in the \(\{\mathbf{u}_{\ell m}\}\) basis for all
\(R\in\mathrm{SO}(3)\) \cite{jackson1999,hansen1988}. Moreover,
Eq.~\eqref{eq:ch3.modal-D-rad-identification} holds with
\(D^{(\mathrm{rad})}_{(\ell,m),(\ell,m')}(R)=\big[D^{(\ell)}(R)\big]_{mm'}\) on radiation modal indices
\cite{stratton1941,harrington2001}.
\end{theorem}
\begin{proof}[Proof sketch]
\(\mathcal{D}(R)\) acts unitarily on \(\mathcal{H}_{\ell}\) under the radiative inner product; matrix elements
in Eq.~\eqref{eq:ch3.Wigner-D-matrix-def} define a homomorphism by
Eq.~\eqref{eq:ch3.Wigner-D-homomorphism} \cite{jackson1999}. Uniqueness up to basis phase follows from
irreducibility Theorem~\ref{thm:ch3.J-irrep-degeneracy} and Schur's lemma \cite{hansen1988}. Agreement with
Eq.~\eqref{eq:ch2.modal-coefficient-rotation} is Eq.~\eqref{eq:ch3.Wigner-D-coefficient-law}
\cite{stratton1941}.
% PROOF: AUTHOR REVIEW
\end{proof}

\paragraph{Irreducible coupling and Clebsch--Gordan coefficients.}
Superposed sources generate superposed \(\ell\) sectors; coupling is not arbitrary. When two angular sectors
\(\mathcal{H}_{\ell_1}\) and \(\mathcal{H}_{\ell_2}\) combine, the tensor product decomposes into irreducible
sectors \(\mathcal{H}_{\ell_3}\) with
\begin{equation}
\label{eq:ch3.CG-triangle-rule}
|\ell_1-\ell_2|
\le
\ell_3
\le
\ell_1+\ell_2,
\qquad
\ell_3\in\{|\ell_1-\ell_2|,\ldots,\ell_1+\ell_2\},
\end{equation}
the triangle rule for \(\mathrm{SO}(3)\) coupling \cite{jackson1999,hansen1988}. Clebsch--Gordan coefficients
define the orthonormal coupling basis:
\begin{equation}
\label{eq:ch3.CG-coupling-def}
\mathbf{u}_{\ell_3 m_3}^{(\ell_1\ell_2)}
=
\sum_{m_1=-\ell_1}^{\ell_1}
\sum_{m_2=-\ell_2}^{\ell_2}
C_{\ell_1 m_1;\ell_2 m_2}^{\ell_3 m_3}\,
\mathbf{u}_{\ell_1 m_1}\otimes \mathbf{u}_{\ell_2 m_2},
\end{equation}
with
\begin{equation}
\label{eq:ch3.CG-selection-m}
C_{\ell_1 m_1;\ell_2 m_2}^{\ell_3 m_3}=0
\quad\text{unless}\ 
m_3=m_1+m_2,
\end{equation}
and normalization \(\sum_{m_1,m_2}|C_{\ell_1 m_1;\ell_2 m_2}^{\ell_3 m_3}|^{2}=1\) for fixed \((\ell_1,\ell_2,\ell_3,m_3)\)
\cite{jackson1999,hansen1988}. Equations~\eqref{eq:ch3.CG-coupling-def}--\eqref{eq:ch3.CG-selection-m} are the
Subsection~\ref{sec:ch324} HOME coupling law on \(\mathcal{F}_{\mathrm{rad}}\): they specify how independent
multipole coordinates assemble into higher-\(\ell\) export when sources superpose linearly through
Eq.~\eqref{eq:ch3.eigenfield-superposition-law} \cite{stratton1941,hansen1988}.

\paragraph{Coupled-field superposition on exported states.}
For admissible \(\mathfrak{D}_1,\mathfrak{D}_2\),
\begin{equation}
\label{eq:ch3.irreducible-coupling-superposition}
\mathcal{R}_\omega[\mathfrak{D}_1+\mathfrak{D}_2]
=
\sum_{\ell_3,m_3}
\Bigg(
\sum_{\ell_1,m_1,\ell_2,m_2}
C_{\ell_1 m_1;\ell_2 m_2}^{\ell_3 m_3}\,
c_{\ell_1 m_1}[\mathfrak{D}_1]\,
c_{\ell_2 m_2}[\mathfrak{D}_2]
\Bigg)
\mathbf{u}_{\ell_3 m_3},
\end{equation}
with each coefficient computed by Eq.~\eqref{eq:ch3.riesz-coefficient-representation} on outgoing carriers
\cite{stratton1941,hansen1988}. Equation~\eqref{eq:ch3.irreducible-coupling-superposition} is nonlinear in
\emph{sources} but bilinear in coefficients: interference creates \(\ell_3\) content only through coupling
coefficients, not by re-labeling Cartesian components \cite{jackson1999,harrington2001}. Selection rule
Eq.~\eqref{eq:ch3.CG-selection-m} forbids net \(m\) mismatch in coupled reports, mirroring azimuthal
conservation in multipole addition \cite{jackson1999,stratton1941}.

\paragraph{Coupling coefficients for \(\ell_1=\ell_2=1\).}
The first nontrivial coupling beyond a single dipole is \(\ell_3=2\) quadrupole content. With standard
normalization,
\begin{equation}
\label{eq:ch3.CG-11-to-2-example}
C_{1,m_1;1,m_2}^{2,m_1+m_2}
=
\frac{1}{\sqrt{2}},
\quad
|m_1|=|m_2|=1,\ m_1+m_2\in\{-2,0,+2\},
\end{equation}
and other nonzero entries follow from ladder phases in \cite{jackson1999,hansen1988}
\cite{jackson1999}. Equation~\eqref{eq:ch3.CG-11-to-2-example} makes Eq.~\eqref{eq:ch3.irreducible-coupling-superposition}
concrete: two tilted dipoles with \(m=\pm1\) content can export \(\ell=2\) sectors only through CG weights, not by
assigning a second independent \(\ell=1\) label to the composite pattern \cite{hansen1988,harrington2001}.
Triangle rule Eq.~\eqref{eq:ch3.CG-triangle-rule} simultaneously forbids \(\ell_3>2\) from two dipole sectors,
explaining why higher multipoles require additional source complexity rather than coordinate relabeling
\cite{jackson1999,stratton1941}.

\paragraph{Tensor-product viewpoint on coupled sectors.}
Abstractly, uncoupled pairs live in \(\mathcal{H}_{\ell_1}\otimes\mathcal{H}_{\ell_2}\) with dimension
\((2\ell_1+1)(2\ell_2+1)\) \cite{jackson1999}. Clebsch--Gordan coupling is the unitary change of basis that
diagonalizes \(\mathbf{J}^{2}\) on the product space:
\begin{equation}
\label{eq:ch3.CG-tensor-product-decomposition}
\mathcal{H}_{\ell_1}\otimes\mathcal{H}_{\ell_2}
=
\bigoplus_{\ell_3=|\ell_1-\ell_2|}^{\ell_1+\ell_2}
\mathcal{H}_{\ell_3},
\end{equation}
with each summand carrying its own Wigner block \(D^{(\ell_3)}(R)\) under rotations
\cite{jackson1999,hansen1988}. Equation~\eqref{eq:ch3.CG-tensor-product-decomposition} is the sector-level
statement behind Eq.~\eqref{eq:ch3.irreducible-coupling-superposition}: superposed sources explore the direct
sum on the right, and exported power partitions across \(\ell_3\) according to CG magnitudes
\cite{stratton1941,harrington2001}. Recoupling coefficients for triple products are deferred to
Section~\ref{sec:ch38}; Subsection~\ref{sec:ch324} stops at pairwise coupling sufficient for two-source
interference in Volume~I \cite{hansen1988}.

\begin{figure}[t]
\centering
\ChOneFigBlock{%
\begin{tikzpicture}[ch1fig, node distance=8mm and 10mm]
  \node[ch1box] (h1) {$\mathcal{H}_{\ell_1}$};
  \node[ch1box, right=12mm of h1] (h2) {$\mathcal{H}_{\ell_2}$};
  \node[ch1box, below=12mm of h1] (t) {$\mathcal{H}_{\ell_1}\otimes\mathcal{H}_{\ell_2}$};
  \node[ch1box, right=20mm of t] (sum) {$\bigoplus_{\ell_3}\mathcal{H}_{\ell_3}$};
  \node[ch1box, below=11mm of sum] (cg) {$C_{\ell_1 m_1;\ell_2 m_2}^{\ell_3 m_3}$};
  \draw[ch1flow] (h1) -- (t);
  \draw[ch1flow] (h2) -- (t);
  \draw[ch1flow] (t) -- (sum);
  \draw[ch1flow] (sum) -- (cg);
\end{tikzpicture}%
}
\caption{Subsection~\ref{sec:ch324} coupling geometry Eq.~\eqref{eq:ch3.CG-tensor-product-decomposition}: tensor
product sectors decompose into direct-sum irreducible blocks with Clebsch--Gordan weights.}
\label{fig:ch3.CG-coupling-geometry}
\end{figure}

Figure~\ref{fig:ch3.CG-coupling-geometry} is the coupling diagram paired with
Figure~\ref{fig:ch3.Wigner-D-coefficient-map}: rotations act on each \(\mathcal{H}_{\ell_3}\) output block through
\(D^{(\ell_3)}(R)\) after CG projection \cite{jackson1999,hansen1988}.

\begin{proposition}[Coupling associativity and rotation covariance]
\label{prop:ch3.CG-rotation-covariance}
Clebsch--Gordan coupling is associative up to standard recoupling phases, and Wigner matrices commute with
coupling in the sense
\begin{equation}
\label{eq:ch3.CG-Wigner-commutation}
D^{(\ell_3)}(R)
\sum_{m_1,m_2}
C_{\ell_1 m_1;\ell_2 m_2}^{\ell_3 m_3}\,
c_{\ell_1 m_1}\,c_{\ell_2 m_2}
=
\sum_{m_1',m_2'}
C_{\ell_1 m_1';\ell_2 m_2'}^{\ell_3 m_3'}\,
\big[D^{(\ell_1)}(R)\mathbf{c}^{(\ell_1)}\big]_{m_1'}\,
\big[D^{(\ell_2)}(R)\mathbf{c}^{(\ell_2)}\big]_{m_2'},
\end{equation}
for rotated coefficient stacks when all sectors lie in \(\mathcal{D}(\mathbf{J})\cap\mathcal{F}_{\mathrm{rad}}\)
\cite{jackson1999,hansen1988,stratton1941}.
\end{proposition}
\begin{proof}[Proof sketch]
Coupling is defined by \(\mathrm{SO}(3)\) representation theory on tensor products; rotation covariance follows
from Eq.~\eqref{eq:ch3.Wigner-D-homomorphism} and irreducibility
\cite{jackson1999,hansen1988}.
% PROOF: AUTHOR REVIEW
\end{proof}

\paragraph{Wigner matrices in truncated eigenfield expansions.}
Eigenfield expansions Eq.~\eqref{eq:ch3.eigenfield-expansion-on-frad} with cutoff \(\ell_{\max}\) truncate
\(\mathbf{c}^{(\ell)}\) to \(0\le\ell\le\ell_{\max}\) \cite{hansen1988,stratton1941}. Rotation still mixes
\emph{within} each retained block through Eq.~\eqref{eq:ch3.Wigner-D-coefficient-law} but cannot generate
\(\ell>\ell_{\max}\) from truncated stacks unless coupling sources supply new sectors through
Eq.~\eqref{eq:ch3.irreducible-coupling-superposition} \cite{jackson1999,harrington2001}. Truncation remainder
Eq.~\eqref{eq:ch3.eigenfield-truncation-remainder} therefore splits into two distinct errors: omitted high-\(\ell\)
physics, and misapplied \(D^{(\ell)}\) conventions on the retained blocks \cite{hansen1988}. Certified rotation
post-processing applies \(D^{(\ell)}(R)\) only after \(\Pi_{\mathrm{rad}}\) and flux normalization
Eq.~\eqref{eq:ch3.radiative-inner-product-def} have fixed the basis in which Eq.~\eqref{eq:ch3.Wigner-D-matrix-def}
is defined \cite{stratton1941,harrington2001}.

\paragraph{Validity domain.}
Eqs.~\eqref{eq:ch3.Wigner-D-matrix-def}--\eqref{eq:ch3.irreducible-coupling-superposition} assume isotropic
vacuum response, linear superposition, outgoing \(\mathcal{B}_\Gamma\), and joint bases fixed consistently with
Section~\ref{sec:ch34} \cite{stratton1941,jackson1999}. Anisotropic materials and magnetized plasmas require
modified \(D\)-matrices not tabulated here \cite{stratton1941}. Near-field reactive samples violate
Eq.~\eqref{eq:ch3.Wigner-D-coefficient-law} for the same domain reasons as
Eq.~\eqref{eq:ch3.J-operator-domain-rad} \cite{harrington2001}.

\paragraph{Worked illustration (\(\ell=1\) rotation about \(\hat{\mathbf{z}}\)).}
For \(R=R_z(\varphi)\), Eq.~\eqref{eq:ch3.Wigner-Euler-factorization} gives
\(D^{(1)}_{m'm}(\varphi,0,0)=e^{-\jj m'\varphi}\delta_{mm'}\) \cite{jackson1999}. A dipole with sole
\(c_{1,0}\neq 0\) therefore remains in the \(m=0\) channel under \(z\)-axis rotations, while a circular
\(m=\pm1\) state acquires phase \(e^{\mp\jj\varphi}\) without \(\ell\) mixing \cite{hansen1988,stratton1941}.
The result matches Eq.~\eqref{eq:ch1.rotation-am-invariance} on totals but not on individual \(m\) components
\cite{stratton1941}.

\paragraph{Worked illustration (\(\ell=1\) tilt via polar angle \(\beta\)).}
A rotation \(R_y(\beta)\) mixes \(m=0\) with \(m=\pm1\) through reduced elements \(d^{(1)}_{m'm}(\beta)\)
\cite{jackson1999}. Tilting a \(\hat{\mathbf{z}}\) dipole therefore populates all three \(m\) slots in
Eq.~\eqref{eq:ch3.Wigner-D-coefficient-law} even when the source current remains rank-one in Cartesian
coordinates \cite{hansen1988,harrington2001}. Laboratory teams that report only \(c_{1,0}\) after tilting the
hardware without applying \(D^{(1)}(R_y(\beta))\) misstate the exported multiplet content
\cite{stratton1941}.

\paragraph{Worked illustration (two-dipole coupling \(\ell_1=\ell_2=1\to\ell_3=2\)).}
Triangle rule Eq.~\eqref{eq:ch3.CG-triangle-rule} allows \(\ell_3=0,1,2\) from two dipole sectors
\cite{jackson1999}. The \(\ell_3=2\) quadrupole content appears with coefficients
\(C_{1 m_1;1 m_2}^{2 m_3}\), vanishing unless \(m_3=m_1+m_2\) by Eq.~\eqref{eq:ch3.CG-selection-m}
\cite{hansen1988}. Equation~\eqref{eq:ch3.irreducible-coupling-superposition} is the certified replacement for
heuristic ``add the patterns'' rules in array synthesis \cite{harrington2001,stratton1941}.

\paragraph{Worked illustration (passive versus active \(D\)-matrix reports).}
Passive frame rotation by \(R\) relabels components through
Eq.~\eqref{eq:ch3.Wigner-D-active-passive} \cite{stratton1941}. Comparing simulated active-source rotation
\(\mathcal{A}_R\) to measured passive-relabeled coefficients without the transpose-inverse map inflates
mismatch indices in Section~\ref{sec:ch01.0.2} even when Eq.~\eqref{eq:ch2.radiation-rotation-commutation} holds
\cite{stratton1941,harrington2001}.

\paragraph{Worked illustration (numerical phase alignment).}
Finite codes often output \(Y_\ell^m\) with Condon--Shortley phases while
Eq.~\eqref{eq:ch3.Wigner-Euler-factorization} assumes the Jackson phase convention
\cite{jackson1999,hansen1988}. A constant diagonal phase bridge between conventions must be applied before
testing Eq.~\eqref{eq:ch3.CG-Wigner-commutation}; otherwise coupling residuals are gauge artifacts, not physics
\cite{hansen1988,harrington2001}.

\paragraph{Worked illustration (eigenfield expansion rotation consistency).}
Let \(\mathbf{u}=\sum_{\ell\le\ell_{\max}}\sum_m c_{\ell m}\mathbf{u}_{\ell m}\) be a truncated outgoing
expansion Eq.~\eqref{eq:ch3.eigenfield-expansion-on-frad} \cite{hansen1988}. Rotating the source actively by
\(R\) and re-expanding must equal applying \(D^{(\ell)}(R)\) to the original coefficients for each \(\ell\) block:
\begin{equation}
\label{eq:ch3.Wigner-truncation-consistency}
c_{\ell m}\big[\mathcal{A}_R\mathfrak{D}\big]
=
\sum_{m'}D^{(\ell)}_{m m'}(R)\,c_{\ell m'}[\mathfrak{D}]
\quad\forall\,\ell\le\ell_{\max},
\end{equation}
whenever Eq.~\eqref{eq:ch2.radiation-rotation-commutation} and basis phases are aligned
\cite{stratton1941,harrington2001}. Violations of Eq.~\eqref{eq:ch3.Wigner-truncation-consistency} after
re-expansion indicate inconsistent normalization, wrong Condon--Shortley phases, or evaluation on reactive
samples rather than \(\mathcal{R}_\omega[\mathfrak{D}]\) \cite{hansen1988,harrington2001}. The test is a
standard regression check for vector spherical harmonic post-processors before coupling coefficients are applied
in Eq.~\eqref{eq:ch3.irreducible-coupling-superposition} \cite{stratton1941}.

\paragraph{Mode-budget summary at \(\ell_{\max}=2\).}
Triangle rule Eq.~\eqref{eq:ch3.CG-triangle-rule} together with
Eq.~\eqref{eq:ch3.J-irrep-dimension-identity} gives nine angular slots
\(\sum_{\ell=0}^{2}(2\ell+1)=9\) before radial indices enter in Section~\ref{sec:ch33}
\cite{hansen1988,jackson1999}. Pairwise coupling of two \(\ell=1\) sources produces \(\ell_3\in\{0,1,2\}\) with
dimensions \(1\), \(3\), and \(5\) respectively; the composite expansion therefore occupies at most the
\(\ell\le2\) portion of the nine-slot budget even when both parents are dipoles
\cite{jackson1999,harrington2001,stratton1941}. Budget tables that list ``three dipole components'' as three
independent \(\ell=1\) degrees of freedom double-count the same \(\rho^{(1)}\) block and must be replaced by
Eqs.~\eqref{eq:ch3.Wigner-D-coefficient-law} and~\eqref{eq:ch3.irreducible-coupling-superposition}
\cite{hansen1988}.

\begin{table}[t]
\centering
\caption{Wigner and coupling convention gates for Subsection~\ref{sec:ch324}.}
\label{tab:ch3.Wigner-convention-gates}
\small
\begin{tabular}{@{}p{0.22\linewidth}p{0.34\linewidth}p{0.36\linewidth}@{}}
\toprule
Gate & Anchor & Typical failure \\
\midrule
Basis phases &
Section~\ref{sec:ch34} \(\mathbf{u}_{\ell m}\) choice &
CG residuals from sign swaps \\
\addlinespace
Active/passive &
Eq.~\eqref{eq:ch3.Wigner-D-active-passive} &
Transpose mismatch vs.\ simulation \\
\addlinespace
Euler convention &
Eq.~\eqref{eq:ch3.Wigner-Euler-factorization} &
\(\alpha,\gamma\) offset errors \\
\addlinespace
Outgoing carrier &
Eq.~\eqref{eq:ch3.J-operator-domain-rad} &
\(D^{(\ell)}\) fit on reactive scans \\
\bottomrule
\end{tabular}
\end{table}

Table~\ref{tab:ch3.Wigner-convention-gates} complements Table~\ref{tab:ch3.SO3-degeneracy-classification}:
the latter classified sector dimension; the present table guards executable matrix conventions for
Eqs.~\eqref{eq:ch3.Wigner-D-coefficient-law} and~\eqref{eq:ch3.irreducible-coupling-superposition}
\cite{stratton1941,hansen1988}.

Subsection~\ref{sec:ch324} has defined Wigner matrices Eq.~\eqref{eq:ch3.Wigner-D-matrix-def}, recorded the
coefficient law Eq.~\eqref{eq:ch3.Wigner-D-coefficient-law}, stated homomorphism and passive conjugation
Eqs.~\eqref{eq:ch3.Wigner-D-homomorphism}--\eqref{eq:ch3.Wigner-D-active-passive}, supplied Euler
factorization Eq.~\eqref{eq:ch3.Wigner-Euler-factorization}, and installed Clebsch--Gordan coupling
Eqs.~\eqref{eq:ch3.CG-coupling-def}--\eqref{eq:ch3.irreducible-coupling-superposition} on
\(\mathcal{F}_{\mathrm{rad}}\) \cite{jackson1999,hansen1988,belinfante1940,stratton1941,harrington2001}.
Subsection~\ref{sec:ch325} closes helicity, polarization, and dual-symmetry sectors under gauge discipline from
Section~\ref{sec:ch01.4} \cite{belinfante1940,stratton1941}.


\subsection{Helicity, polarization, and dual symmetry}
\label{sec:ch325}
% TOC tag: [USE SS1.4.3][HOME]

Orbital labels \((\ell,m)\) from Subsections~\ref{sec:ch322}--\ref{sec:ch324} organize rotation on
\(\mathcal{F}_{\mathrm{rad}}\) but do not exhaust vector structure: transverse polarization, signed
electric--magnetic handedness, and duality rotations carry independent information screened by gauge rules
from Section~\ref{sec:ch01.4} \cite{belinfante1940,stratton1941}. Subsection~\ref{sec:ch325} records the
Chapter~\ref{ch:03} operator-level certificates for helicity and dual symmetry on outgoing eigenfields,
imports polarization observables from Subsection~\ref{sec:ch01.4.3} by reference, and states which
\(\sigma\)/helicity reports remain admissible on \(\mathcal{R}_\omega[\mathfrak{D}]\)
\cite{jackson1999,hansen1988,harrington2001}.

\paragraph{Standing imports.}
Helicity density Eq.~\eqref{eq:ch1.helicity-density}, flux integral
Eq.~\eqref{eq:ch1.helicity-flux-integral}, gauge-variation identity
Eq.~\eqref{eq:ch1.helicity-density-wave-reduction}, observable triad
Eq.~\eqref{eq:ch1.observable-admissibility-triad}, duality rotation
Eq.~\eqref{eq:ch1.duality-rotation}, dual invariants
Eqs.~\eqref{eq:ch1.dual-invariant-energy}--\eqref{eq:ch1.dual-invariant-poynting}, and helicity decision
Eq.~\eqref{eq:ch1.helicity-observable-decision} enter from Subsection~\ref{sec:ch01.4.3} by reference only
\cite{stratton1941,belinfante1940}. Belinfante spin--orbit redistribution
Eqs.~\eqref{eq:ch1.spin-orbit-equivalence-class}--\eqref{eq:ch1.spin-orbit-observable-decision} are imported
from Section~\ref{sec:ch01.4} \cite{belinfante1940}. Total angular momentum
Eq.~\eqref{eq:ch3.J-total-maxwell-def} and Wigner blocks Eq.~\eqref{eq:ch3.Wigner-D-coefficient-law} are
imported from Subsections~\ref{sec:ch321} and~\ref{sec:ch324} \cite{jackson1999,hansen1988}.

\paragraph{Helicity functional on radiative pairs.}
On \(\mathbf{u}_\omega=(\E,\H)\in\mathcal{D}(\mathbf{J})\cap\mathcal{F}_{\mathrm{rad}}\), define the
volume-integrated helicity functional
\begin{equation}
\label{eq:ch3.helicity-functional-rad}
\mathcal{H}_{\mathrm{hel}}[\mathbf{u}_\omega]
=
\int_{V}
h(\mathbf{r})\,dV,
\qquad
h\ \text{from Eq.~\eqref{eq:ch1.helicity-density}},
\end{equation}
with \(V\) an exterior domain enclosing sources and \(h\) evaluated from the phasor pair
\cite{stratton1941,belinfante1940}. Equation~\eqref{eq:ch3.helicity-functional-rad} is the Subsection
\ref{sec:ch325} lift of Eq.~\eqref{eq:ch1.helicity-flux-integral} to the outgoing carrier: it measures
signed electric--magnetic rotational overlap on \(\mathcal{R}_\omega[\mathfrak{D}]\), not on reactive-shell
samples \cite{stratton1941,harrington2001}. A \emph{helicity eigenfield} satisfies
\begin{equation}
\label{eq:ch3.helicity-eigenfield-equation}
\mathcal{H}_{\mathrm{hel}}[\mathbf{u}_{\ell m}^{(\sigma)}]
=
\sigma\,\mathcal{N}_{\mathrm{hel}}[\mathbf{u}_{\ell m}^{(\sigma)}],
\qquad
\sigma\in\{+1,-1,0\},
\end{equation}
with normalization functional \(\mathcal{N}_{\mathrm{hel}}\) fixed by flux normalization in
Section~\ref{sec:ch35} \cite{hansen1988,stratton1941}. Equation~\eqref{eq:ch3.helicity-eigenfield-equation} is
orthogonal to orbital labeling: two states can share \((\ell,m)\) but differ in \(\sigma\)
\cite{jackson1999,belinfante1940}.

\paragraph{Polarization sector on the far sphere.}
At fixed \(r\) in the radiative regime, transverse fields induce tangent-plane polarization on
\(S^{2}\) \cite{stratton1941,jackson1999}. Let \(\hat{\mathbf{e}}_\theta,\hat{\mathbf{e}}_\phi\) be local
spherical unit vectors and \(\mathbf{e}_{\mathrm{t}}=\alpha_\theta\hat{\mathbf{e}}_\theta+\alpha_\phi\hat{\mathbf{e}}_\phi\)
the transverse polarization vector with Jones coefficients \((\alpha_\theta,\alpha_\phi)\)
\cite{stratton1941}. Define the polarization-sector projector on \(\mathcal{H}_\ell\):
\begin{equation}
\label{eq:ch3.polarization-sector-projector}
P_{\sigma}
\mathbf{u}_{\ell m}
=
\mathbf{u}_{\ell m}^{(\sigma)},
\qquad
\langle \mathbf{u}_{\ell m}^{(\sigma)},\mathbf{u}_{\ell m}^{(\sigma')}\rangle_{\mathrm{rad}}
=
\delta_{\sigma\sigma'},
\end{equation}
whenever \(\mathbf{u}_{\ell m}^{(\sigma)}\) is circularly or linearly polarized in the sense of
Eq.~\eqref{eq:ch1.polarization-ellipse-angles} on the far shell \cite{stratton1941,hansen1988}. Equation
\eqref{eq:ch3.polarization-sector-projector} separates handedness content inside a fixed orbital multiplet:
\(\ell\) and \(m\) come from \(\mathbf{J}\); \(\sigma\) comes from transverse vector geometry
\cite{belinfante1940,jackson1999}. Linear polarization has \(\sigma=0\); right/left circular states carry
\(\sigma=\pm1\) under the phase convention of Section~\ref{sec:ch34} \cite{stratton1941,hansen1988}.

\paragraph{Dual-symmetry operator on outgoing pairs.}
Let \(\eta_0=\sqrt{\mu_0/\varepsilon_0}\). The vacuum duality rotation of Subsection~\ref{sec:ch01.4.3} induces
\begin{equation}
\label{eq:ch3.dual-symmetry-operator}
\mathcal{D}_{\xi}\mathbf{u}_\omega
=
\big(\cos\xi\,\E+\sin\xi\,\eta_0\H,\
-\sin\xi\,\E/\eta_0+\cos\xi\,\H\big),
\qquad
\xi\in\mathbb{R},
\end{equation}
on admissible pairs, matching Eq.~\eqref{eq:ch1.duality-rotation} at \(\xi=\pi/2\)
\cite{stratton1941,belinfante1940}. Equation~\eqref{eq:ch3.dual-symmetry-operator} is a symmetry of the
source-free Maxwell system on \(\mathcal{F}_{\mathrm{rad}}\): it mixes \(\E\) and \(\H\) without changing
outgoing rank when \(\mathcal{D}_{\mathrm{dual}}=1\) in Eq.~\eqref{eq:ch1.duality-defect-index}
\cite{stratton1941}. Helicity flips sign under \(\xi=\pi/2\) because electric and magnetic rotational roles
exchange with a handedness inversion \cite{belinfante1940,jackson1999}.

\begin{figure}[t]
\centering
\ChOneFigBlock{%
\begin{tikzpicture}[ch1fig, node distance=8mm and 10mm]
  \node[ch1box] (lm) {$(\ell,m)$\\$\mathbf{J}$};
  \node[ch1box, right=14mm of lm] (sig) {$\sigma$\\polarization};
  \node[ch1box, right=14mm of sig] (hel) {$\mathcal{H}_{\mathrm{hel}}$};
  \node[ch1box, below=11mm of sig] (dual) {$\mathcal{D}_{\xi}$\\dual};
  \draw[ch1flow] (lm) -- (sig) -- (hel);
  \draw[ch1flow] (dual) -- (sig);
\end{tikzpicture}%
}
\caption{Subsection~\ref{sec:ch325} label split on \(\mathcal{F}_{\mathrm{rad}}\): orbital \((\ell,m)\) from
\(\mathbf{J}\), polarization/handness \(\sigma\) from transverse geometry, dual \(\mathcal{D}_{\xi}\) mixing
\(\E\) and \(\H\).}
\label{fig:ch3.helicity-polarization-dual-split}
\end{figure}

Figure~\ref{fig:ch3.helicity-polarization-dual-split} is the reporting diagram complementary to
Figure~\ref{fig:ch3.J-orbital-spin-split}: the latter separated orbital and spin \emph{densities}; the present
figure separates orbital quantum numbers from polarization/helicity \emph{labels} on the same outgoing carrier
\cite{belinfante1940,stratton1941}.

\begin{theorem}[Gauge-invariant helicity reporting on \(\mathcal{F}_{\mathrm{rad}}\)]
\label{thm:ch3.helicity-gauge-invariance}
If \(\mathbf{u}_\omega\in\mathcal{F}_{\mathrm{rad}}\) and the observable triad
Eq.~\eqref{eq:ch1.observable-admissibility-triad} passes with \(\mathcal{D}_{\mathrm{loc}}=1\), then
\(\mathcal{H}_{\mathrm{hel}}[\mathbf{u}_\omega]\) is unchanged under gauge transforms
Eq.~\eqref{eq:ch1.gauge-transform} and transforms as a pseudoscalar under
Eq.~\eqref{eq:ch1.active-rotation-action} \cite{stratton1941,belinfante1940}. Moreover,
\(\mathcal{H}_{\mathrm{hel}}[\mathcal{R}_\omega[\mathfrak{D}]]\) is meaningful only when evaluated on the
outgoing image, not on \(\mathcal{S}_\omega[\mathfrak{D}]\) reactive restrictions
\cite{harrington2001,stratton1941}.
\end{theorem}
\begin{proof}[Proof sketch]
Eq.~\eqref{eq:ch1.helicity-density-wave-reduction} and curl invariance under gauge
\cite{stratton1941}. Export clause imports Eq.~\eqref{eq:ch3.outgoing-inheritance-criterion} and
Eq.~\eqref{eq:ch3.J-operator-domain-rad} \cite{harrington2001}.
% PROOF: AUTHOR REVIEW
\end{proof}

\paragraph{Commutation with angular momentum (reporting form).}
Helicity and orbital angular momentum are not interchangeable quantum numbers. On transverse far fields,
\begin{equation}
\label{eq:ch3.helicity-Jz-commutation-relation}
[\mathcal{H}_{\mathrm{hel}},\mathbf{J}_{z}]
=
0
\quad\text{on eigenfields Eq.~\eqref{eq:ch3.J-simultaneous-eigenfield-system} with fixed }\sigma,
\end{equation}
while \(\mathcal{H}_{\mathrm{hel}}\) changes sign under \(\mathcal{D}_{\pi/2}\) in
Eq.~\eqref{eq:ch3.dual-symmetry-operator} \cite{jackson1999,belinfante1940}. Equation
\eqref{eq:ch3.helicity-Jz-commutation-relation} is the Subsection~\ref{sec:ch325} certificate that a phase ramp
in \(\phi\) (changing \(m\)) is not the same as a helicity flip (changing \(\sigma\)) on \(\mathcal{F}_{\mathrm{rad}}\)
\cite{allen1992,stratton1941}. Paraxial ``OAM mode order'' language is deferred to Section~\ref{sec:ch39} with
explicit gauge caveats \cite{allen1992,belinfante1940}.

\begin{proposition}[Dual symmetry preserves radiative export class]
\label{prop:ch3.dual-preserves-radiation}
When \(\mathcal{D}_{\mathrm{dual}}=1\) and \(\mathbf{u}_\omega\in\mathcal{F}_{\mathrm{rad}}\),
\(\mathcal{D}_{\xi}\mathbf{u}_\omega\in\mathcal{F}_{\mathrm{rad}}\) and
\begin{equation}
\label{eq:ch3.dual-radiation-intertwine}
\mathcal{D}_{\xi}\,\mathcal{R}_\omega[\mathfrak{D}]
=
\mathcal{R}_\omega\big[\mathcal{D}_{\xi}^{(\mathrm{src})}\mathfrak{D}\big]
\end{equation}
on admissible \(\mathfrak{D}\), with \(\mathcal{D}_{\xi}^{(\mathrm{src})}\) the source-side dual map induced by
Eq.~\eqref{eq:ch1.duality-rotation} \cite{stratton1941,harrington2001}. Total exported power obeys
Eq.~\eqref{eq:ch1.dual-invariant-energy}; Poynting flux obeys
Eq.~\eqref{eq:ch1.dual-invariant-poynting} \cite{stratton1941}.
\end{proposition}
\begin{proof}[Proof sketch]
Dual rotation is a Maxwell symmetry in isotropic vacuum; compose with
Eq.~\eqref{eq:ch3.symmetry-commutes-with-radiation} and outgoing idempotence
Eq.~\eqref{eq:ch3.outgoing-inheritance-criterion} \cite{stratton1941,harrington2001}.
% PROOF: AUTHOR REVIEW
\end{proof}

\paragraph{Spin--orbit redistribution on helicity reports.}
Belinfante equivalence Eq.~\eqref{eq:ch1.spin-orbit-equivalence-class} shifts local orbital and spin
densities without changing \(\mathbf{J}\) totals \cite{belinfante1940}. The same redistribution can move
helicity density between \(\E\) and \(\H\) channels in Eq.~\eqref{eq:ch1.helicity-density} while leaving
\(\mathcal{H}_{\mathrm{hel}}\) invariant when Eq.~\eqref{eq:ch1.spin-orbit-observable-decision} passes
\cite{belinfante1940,stratton1941}. Subsection~\ref{sec:ch325} therefore treats \(\sigma\) as a
\emph{gauge-screened} report tied to Eq.~\eqref{eq:ch1.helicity-observable-decision}, not as a raw phase slope
on one component \cite{belinfante1940,harrington2001}.

\paragraph{Stokes parameters on the far shell.}
On \(\Gamma_R\) with \(kR\gg1\), transverse fields define Stokes four-vector
\((S_0,S_1,S_2,S_3)\) through Eq.~\eqref{eq:ch1.polarization-degree} and the spherical Jones coefficients in
Eq.~\eqref{eq:ch3.polarization-sector-projector} \cite{stratton1941}. The handedness proxy is
\begin{equation}
\label{eq:ch3.stokes-helicity-proxy}
\sigma
\propto
\frac{S_3}{S_0},
\qquad
|S_3/S_0|=1
\ \Leftrightarrow\
\text{circular polarization},
\end{equation}
with proportionality fixed by flux normalization Section~\ref{sec:ch35} \cite{stratton1941,hansen1988}. Equation
\eqref{eq:ch3.stokes-helicity-proxy} links laboratory Stokes readouts to
Eq.~\eqref{eq:ch3.helicity-eigenfield-equation}: \(S_3\) changes sign under
\(\mathcal{D}_{\pi/2}\) in Eq.~\eqref{eq:ch3.dual-symmetry-operator} while \(S_0\) remains invariant under
Eq.~\eqref{eq:ch1.dual-invariant-energy} \cite{belinfante1940,stratton1941}. Stokes vectors rotate under
\(D^{(\ell)}(R)\) only through the \(m\) mixing of Eq.~\eqref{eq:ch3.Wigner-D-coefficient-law}; they do not
replace helicity screening Eq.~\eqref{eq:ch1.helicity-observable-decision} \cite{harrington2001}.

\paragraph{Helicity flux and directional transport.}
Integrated helicity Eq.~\eqref{eq:ch3.helicity-functional-rad} pairs with directional angular transport
Eq.~\eqref{eq:ch3.J-directional-transport-functional} when both are evaluated on the same outgoing shell
\cite{stratton1941,belinfante1940}. The joint report vector
\begin{equation}
\label{eq:ch3.helicity-transport-joint-report}
\boldsymbol{\mathcal{R}}_{\mathrm{ang}}
=
\big(\mathcal{H}_{\mathrm{hel}},\ \mathcal{T}_{\hat{\mathbf{z}}},\ \mathcal{T}_{\hat{\mathbf{x}}},\ \mathcal{T}_{\hat{\mathbf{y}}}\big),
\qquad
\mathbf{u}=\mathcal{R}_\omega[\mathfrak{D}],
\end{equation}
is admissible only when each entry passes its gate:
Eq.~\eqref{eq:ch1.helicity-observable-decision} for \(\mathcal{H}_{\mathrm{hel}}\),
Eq.~\eqref{eq:ch1.directional-qualified-gate} for \(\mathcal{T}_{\hat{\mathbf{u}}}\),
and Eq.~\eqref{eq:ch2.far-angular-accessibility-def} for accessible directions
\cite{stratton1941,harrington2001}. Equation~\eqref{eq:ch3.helicity-transport-joint-report} prevents the common
error of quoting helicity from a circularly polarized sidelobe while declaring orbital \(m\) from the main beam
without accessibility composition Eq.~\eqref{eq:ch3.J-accessibility-composition} \cite{harrington2001}.

\paragraph{Eigenfield basis preview for Section~\ref{sec:ch34}.}
Section~\ref{sec:ch34} will construct \(\mathbf{u}_{\ell m}^{(\sigma)}\) with definite \(\ell,m,\sigma\) on
outgoing carriers \cite{hansen1988,tai1994}. Subsection~\ref{sec:ch325} fixes the quantum-number split those
modes must respect:
\begin{equation}
\label{eq:ch3.eigenfield-quantum-number-split}
\mathbf{u}_{\ell m}^{(\sigma)}
\in
\mathcal{H}_{\ell}\cap
\ker(P_{\sigma}-\sigma)\cap
\ker(\mathbf{J}^{2}-\ell(\ell+1))\cap
\ker(\mathbf{J}_{z}-m),
\end{equation}
with \(P_{\sigma}\) from Eq.~\eqref{eq:ch3.polarization-sector-projector} \cite{hansen1988,jackson1999}.
Equation~\eqref{eq:ch3.eigenfield-quantum-number-split} is the constructive target for vector spherical
harmonics: named eigenfields are simultaneous carriers of orbital, polarization, and (when applicable) helicity
labels established in Section~\ref{sec:ch32} \cite{hansen1988,stratton1941}.

\paragraph{Dual rotation group and composition law.}
The family Eq.~\eqref{eq:ch3.dual-symmetry-operator} satisfies
\begin{equation}
\label{eq:ch3.dual-rotation-composition}
\mathcal{D}_{\xi_1}\mathcal{D}_{\xi_2}
=
\mathcal{D}_{\xi_1+\xi_2},
\qquad
\mathcal{D}_{0}=\mathsf{I},
\qquad
\mathcal{D}_{\xi}^{-1}=\mathcal{D}_{-\xi},
\end{equation}
on \(\mathcal{F}_{\mathrm{rad}}\) when \(\mathcal{D}_{\mathrm{dual}}=1\) \cite{stratton1941,belinfante1940}. Equation
\eqref{eq:ch3.dual-rotation-composition} is the operator restatement of
Eq.~\eqref{eq:ch1.duality-rotation}: duality is a continuous symmetry distinct from the discrete rotation group
generated by \(\mathbf{J}\) \cite{jackson1999,stratton1941}. Laboratory workflows that apply a duality transform
to \(\E\) alone without the paired \(\H\) update in Eq.~\eqref{eq:ch3.dual-symmetry-operator} break
Eq.~\eqref{eq:ch3.dual-radiation-intertwine} and produce fictitious helicity imbalance in
Eq.~\eqref{eq:ch3.helicity-transport-joint-report} \cite{harrington2001,belinfante1940}.

\paragraph{Validity domain.}
Eqs.~\eqref{eq:ch3.helicity-functional-rad}--\eqref{eq:ch3.dual-radiation-intertwine} assume linear isotropic
media with \(\mathcal{D}_{\mathrm{dual}}=1\), time-harmonic phasors, and outgoing \(\mathcal{B}_\Gamma\)
\cite{stratton1941,jackson1999}. Lossy, gyrotropic, and bianisotropic materials break
Eq.~\eqref{eq:ch3.dual-radiation-intertwine} even when \(\ell,m\) labels remain well defined
\cite{stratton1941,harrington2001}. Near-field helicity densities on reactive shells are not substitutes for
Eq.~\eqref{eq:ch3.helicity-functional-rad} on \(\mathcal{R}_\omega[\mathfrak{D}]\) \cite{harrington2001}.

\paragraph{Worked illustration (circular far-field dipole).}
A far-zone \(\ell=1\) state with right-circular transverse polarization has \(\sigma=+1\) in
Eq.~\eqref{eq:ch3.helicity-eigenfield-equation} and nonzero \(\mathcal{H}_{\mathrm{hel}}\) when
Eq.~\eqref{eq:ch1.helicity-observable-decision} passes \cite{jackson1999,stratton1941}. The same \(\ell=1\),
\(m=0\) orbital label can appear with \(\sigma=0\) for linear polarization; conflating \(m\) with \(\sigma\) is
an inadmissible report under Figure~\ref{fig:ch3.helicity-polarization-dual-split} \cite{belinfante1940,hansen1988}.

\paragraph{Worked illustration (dual rotation flips helicity sign).}
Applying \(\mathcal{D}_{\pi/2}\) in Eq.~\eqref{eq:ch3.dual-symmetry-operator} to a circular eigenfield flips
\(\mathcal{H}_{\mathrm{hel}}\) while preserving \(\ell,m\) and exported power
Eqs.~\eqref{eq:ch1.dual-invariant-energy} and~\eqref{eq:ch2.spectral-power-rotation-invariance}
\cite{stratton1941,belinfante1940}. Interpreting the flip as an orbital \(m\) change violates
Eq.~\eqref{eq:ch3.helicity-Jz-commutation-relation} \cite{jackson1999}.

\paragraph{Worked illustration (linear polarization null helicity).}
Linearly polarized far fields have vanishing signed overlap in Eq.~\eqref{eq:ch1.helicity-density} even when
\(\ell>0\) and \(|m|>0\) \cite{stratton1941,jackson1999}. Declaring ``zero OAM'' from null helicity confuses
helicity with orbital index \(m\); the correct orbital content is fixed by
Eq.~\eqref{eq:ch3.J-simultaneous-eigenfield-system} independently of \(\sigma\)
\cite{belinfante1940,hansen1988}.

\paragraph{Worked illustration (gauge failure on potential read).}
Reading helicity from a potential \(\mathbf{A}\) with non-coulomb gauge while \(\mathcal{D}_{\mathrm{loc}}=0\)
can produce spurious \(\sigma\) labels that cancel only after returning to \(\E,\H\)
\cite{stratton1941,belinfante1940}. The certified path is
Eq.~\eqref{eq:ch3.helicity-functional-rad} on \(\mathcal{R}_\omega[\mathfrak{D}]\) with triad
Eq.~\eqref{eq:ch1.observable-admissibility-triad} \cite{harrington2001}.

\paragraph{Worked illustration (Wigner rotation of polarization-only DOF).}
Within fixed \(\ell\), \(D^{(\ell)}(R)\) in Eq.~\eqref{eq:ch3.Wigner-D-coefficient-law} mixes \(m\) partners but
does not generate \(\sigma\) changes unless the source field actually carries dual or polarization content
\cite{jackson1999,hansen1988}. Post-multiplying a linearly polarized \(\ell=1\) state by \(D^{(1)}(R)\) yields
another linear state with rotated Stokes parameters, not a circular state with \(\sigma=\pm1\)
\cite{stratton1941,harrington2001}.

\paragraph{Worked illustration (joint Stokes and \(m\) report on \(\ell=1\)).}
A \(\ell=1\), \(m=+1\) eigenfield can present \(|S_3/S_0|\approx1\) on a chosen far cut while a different
\(m\) partner in the same triplet shows \(|S_3/S_0|\ll1\) \cite{jackson1999,hansen1988}. Reporting
\(\sigma=+1\) for the multiplet from a single Stokes pixel without
Eq.~\eqref{eq:ch3.polarization-sector-projector} collapses three distinct \(m\) states into one handedness
label \cite{stratton1941,harrington2001}.

\paragraph{Worked illustration (dual image of quadrupole sector).}
Applying \(\mathcal{D}_{\pi/2}\) to an \(\ell=2\) export with \(\sigma=+1\) yields another outgoing state in
\(\mathcal{H}_{2}\) with \(\sigma=-1\) and identical power budget
Eq.~\eqref{eq:ch2.spectral-power-rotation-invariance} \cite{stratton1941,belinfante1940}. The dual pair shares
Wigner block \(D^{(2)}(R)\) but differs in helicity entry of
Eq.~\eqref{eq:ch3.helicity-transport-joint-report} \cite{jackson1999,hansen1988}.

\begin{table}[t]
\centering
\caption{Helicity, polarization, and dual-symmetry reporting gates for Subsection~\ref{sec:ch325}.}
\label{tab:ch3.helicity-polarization-gates}
\small
\begin{tabular}{@{}p{0.20\linewidth}p{0.34\linewidth}p{0.36\linewidth}@{}}
\toprule
Report & Admissibility anchor & Common failure \\
\midrule
Helicity \(\sigma\) &
Eq.~\eqref{eq:ch1.helicity-observable-decision} &
Reactive-shell \(h(\mathbf{r})\) read \\
\addlinespace
Orbital \(m\) &
Eq.~\eqref{eq:ch3.J-simultaneous-eigenfield-system} &
Phase ramp called \(\sigma\) \\
\addlinespace
Dual image &
Eq.~\eqref{eq:ch3.dual-radiation-intertwine} &
Lossy medium with \(\mathcal{D}_{\mathrm{dual}}\neq1\) \\
\addlinespace
Polarization &
Eq.~\eqref{eq:ch3.polarization-sector-projector} &
Jones vector from near-field sample \\
\bottomrule
\end{tabular}
\end{table}

Table~\ref{tab:ch3.helicity-polarization-gates} closes the Section~\ref{sec:ch32} reporting ladder: operator
domain (Table~\ref{tab:ch3.J-reporting-gates}), commutation (Table~\ref{tab:ch3.J-commutation-diagnostics}),
degeneracy (Table~\ref{tab:ch3.SO3-degeneracy-classification}), Wigner conventions
(Table~\ref{tab:ch3.Wigner-convention-gates}), and now helicity/polarization/dual sectors
\cite{stratton1941,harrington2001,belinfante1940}.

Subsection~\ref{sec:ch325} has installed helicity functional
Eq.~\eqref{eq:ch3.helicity-functional-rad}, eigenfield equation
Eq.~\eqref{eq:ch3.helicity-eigenfield-equation}, polarization projector
Eq.~\eqref{eq:ch3.polarization-sector-projector}, dual operator
Eq.~\eqref{eq:ch3.dual-symmetry-operator}, and export intertwining
Eq.~\eqref{eq:ch3.dual-radiation-intertwine} on \(\mathcal{F}_{\mathrm{rad}}\)
\cite{belinfante1940,jackson1999,stratton1941,hansen1988,harrington2001}.

\begin{figure}[t]
\centering
\ChOneFigBlock{%
\begin{tikzpicture}[ch1fig, node distance=7mm and 9mm]
  \node[ch1box] (rad) {$\mathcal{R}_\omega[\mathfrak{D}]$};
  \node[ch1box, right=11mm of rad] (c) {$\mathbf{c}^{(\ell)}$};
  \node[ch1box, right=11mm of c] (d) {$D^{(\ell)}(R)$};
  \node[ch1box, below=10mm of c] (sig) {$\sigma$ / Stokes};
  \node[ch1box, below=10mm of d] (hel) {$\mathcal{H}_{\mathrm{hel}}$};
  \draw[ch1flow] (rad) -- (c) -- (d);
  \draw[ch1flow] (c) -- (sig);
  \draw[ch1flow] (sig) -- (hel);
\end{tikzpicture}%
}
\caption{Subsection~\ref{sec:ch325} certified reporting order: outgoing export, orbital coefficients,
\(D^{(\ell)}\) frame consistency, then polarization/helicity functionals.}
\label{fig:ch3.helicity-measurement-order}
\end{figure}

Figure~\ref{fig:ch3.helicity-measurement-order} enforces ordering: helicity and Stokes reports are downstream of
\(\Pi_{\mathrm{rad}}\) and Wigner consistency, not parallel shortcuts from near-field phase
\cite{harrington2001,stratton1941,belinfante1940}.

\paragraph{Measurement chain with helicity and polarization gates.}
The operational report Eq.~\eqref{eq:ch3.measurement-on-coefficients} acquires helicity and polarization
sectors only after \(\mathcal{C}\) uses a basis satisfying
Eq.~\eqref{eq:ch3.eigenfield-quantum-number-split} \cite{harrington2001,stratton1941}. A certified pipeline
therefore applies \(\Pi_{\mathrm{rad}}\), extracts \(\mathbf{c}^{(\ell)}\), rotates with
\(D^{(\ell)}(R)\) when frames change, couples sources through
Eq.~\eqref{eq:ch3.irreducible-coupling-superposition} when multiple sectors interfere, and only then evaluates
\(\mathcal{H}_{\mathrm{hel}}\) and Stokes proxies Eqs.~\eqref{eq:ch3.stokes-helicity-proxy} and
\eqref{eq:ch3.helicity-transport-joint-report} \cite{stratton1941,hansen1988,harrington2001}. Skipping any
intermediate gate produces the popular but inadmissible pattern: a single scalar ``OAM number'' read from a
phase fit without \(\mathbf{J}\), \(D^{(\ell)}\), or helicity screens \cite{allen1992,belinfante1940,harrington2001}.

\paragraph{Summative closure of Section~\ref{sec:ch32}.}
Section~\ref{sec:ch32} began from the Lie-algebra necessity named in its introduction: exported radiation
requires explicit angular-momentum generators, commutation certificates, irreducible sector counts, executable
rotation matrices, and gauge-screened vector reports on the same \(\mathcal{F}_{\mathrm{rad}}\) closed in
Section~\ref{sec:ch31} \cite{jackson1999,belinfante1940,stratton1941,hansen1988}. Subsection~\ref{sec:ch321}
defined \(\mathbf{J}\) once through Eqs.~\eqref{eq:ch3.J-orbital-maxwell}--\eqref{eq:ch3.J-total-maxwell-def};
Subsection~\ref{sec:ch322} derived the \(\mathfrak{so}(3)\) algebra and simultaneous eigenfields;
Subsection~\ref{sec:ch323} identified \(\mathcal{H}_{\ell}\) irreps and the \(2\ell+1\) degeneracy law;
Subsection~\ref{sec:ch324} tabulated \(D^{(\ell)}(R)\) and Clebsch--Gordan coupling;
Subsection~\ref{sec:ch325} separated helicity, polarization, and dual symmetry from orbital labels through
Eqs.~\eqref{eq:ch3.helicity-functional-rad}--\eqref{eq:ch3.dual-rotation-composition}
\cite{harrington2001}. The five subsection tables record executable gates rather than philosophical
taxonomy \cite{stratton1941,belinfante1940}. Section~\ref{sec:ch33} now reduces the vector wave equation and
separates variables in spherical coordinates so Section~\ref{sec:ch34} can construct the
\(\mathbf{u}_{\ell m}^{(\sigma)}\) bases in which every operator and matrix of Section~\ref{sec:ch32} acts
numerically \cite{hansen1988,tai1994,stratton1941}. Until those eigenfields exist by name, the present section
supplies the operator certificates that any constructed basis must diagonalize or covariantly transform
\cite{hansen1988,jackson1999}.


\section{Vector Helmholtz Theory and Separation of Variables}
\label{sec:ch33}

Eigenfield construction requires a separated wave equation in the coordinates that match the symmetry.
Section~\ref{sec:ch33} reduces the vector Maxwell problem to the Helmholtz form already symbolized in
Subsection~\ref{sec:ch222}, then separates variables in spherical coordinates. Subsection~\ref{sec:ch331}
records the vector Helmholtz reduction without re-displaying the curl--curl eliminations of
Chapter~\ref{ch:02}. Subsection~\ref{sec:ch332} performs angular--radial separation. Subsection~\ref{sec:ch333}
imposes outgoing radial admissibility using Sommerfeld selection from Section~\ref{sec:ch24}
\cite{stratton1941,jackson1999}. Subsection~\ref{sec:ch334} introduces Debye potentials and generating
functions that streamline the vector spherical harmonic construction in Section~\ref{sec:ch34}
\cite{stratton1941,tai1994}.


\subsection{Reduction to the vector Helmholtz equation}
\label{sec:ch331}
% TOC tag: [HOME][USE SS2.2.2]

Section~\ref{sec:ch32} fixed the symmetry operators that eigenfields must diagonalize; Section~\ref{sec:ch33}
now reduces the Maxwell system to a separated wave equation on the same outgoing carrier
\(\mathcal{F}_{\mathrm{rad}}\) \cite{stratton1941,hansen1988}. Subsection~\ref{sec:ch331} records the
Chapter~\ref{ch:03} vector Helmholtz form without repeating curl--curl eliminations already executed in
Subsection~\ref{sec:ch222} \cite{stratton1941,harrington2001}.

\paragraph{Standing imports.}
The harmonic Maxwell residual Eq.~\eqref{eq:ch2.harmonic-maxwell-residual}, electric-wave operator
Eq.~\eqref{eq:ch2.electric-wave-operator}, domain Eq.~\eqref{eq:ch2.electric-helmholtz-domain}, symbol
Eq.~\eqref{eq:ch2.helmholtz-symbol}, and vacuum pair Eq.~\eqref{eq:ch1.vacuum-helmholtz-pair} enter from
Subsections~\ref{sec:ch222} and Chapter~\ref{ch:01} by reference \cite{stratton1941}. Outgoing inheritance
Eq.~\eqref{eq:ch3.outgoing-inheritance-criterion} and Hilbert closure
Eq.~\eqref{eq:ch3.F-rad-hilbert-closure} are imported from Section~\ref{sec:ch31}
\cite{kato1995,hansen1988}.

\paragraph{Assumptions.}
Work in a linear isotropic exterior region with phasor fields at fixed \(\omega\), wavenumber
\(k_0=\omega\sqrt{\varepsilon_0\mu_0}\), and pairs \(\mathbf{u}_\omega=(\E,\H)\in\mathcal{F}_{\mathrm{rad}}\) satisfying
Eq.~\eqref{eq:ch3.outgoing-inheritance-criterion} \cite{stratton1941,jackson1999}. Sources are absent in the
immediate separation domain of Subsection~\ref{sec:ch331}; impressed currents enter later through integral
representations imported from Section~\ref{sec:ch25} \cite{stratton1941,harrington2001}.

\paragraph{Curl--curl identity without re-opening Chapter~\ref{ch:02}.}
On domains where \(\curl\E\) is defined in \(\HCurl\), the vector identity
\begin{equation}
\label{eq:ch3.curl-curl-vector-identity}
\curl\curl\E=\nabla(\div\E)-\nabla^{2}\E
\end{equation}
holds distributionally \cite{stratton1941,jackson1999}. Source-free Maxwell with harmonic phasors gives
\(\div\E=0\) from Eq.~\eqref{eq:ch2.harmonic-maxwell-residual}, so Eq.~\eqref{eq:ch3.curl-curl-vector-identity}
collapses the electric-wave operator of Eq.~\eqref{eq:ch2.electric-wave-operator} to a vector Laplacian shift
without repeating the elimination steps of Subsection~\ref{sec:ch222} \cite{stratton1941}. The magnetic channel
uses the dual identity on \(\H\) with \(\div\H=0\) \cite{jackson1999}. Subsection~\ref{sec:ch331} records the
collapsed operator as the Chapter~\ref{ch:03} HOME wave equation on \(\mathcal{F}_{\mathrm{rad}}\)
\cite{hansen1988}.

\paragraph{Vector Helmholtz reduction (HOME).}
Eliminating \(\H\) through \(\H=(\jj\omega\mu_0)^{-1}\curl\E\) on source-free domains yields the electric vector
Helmholtz equation
\begin{equation}
\label{eq:ch3.vector-helmholtz-E}
\curl\curl\E-k_0^{2}\E=0,
\qquad
k_0=\omega\sqrt{\varepsilon_0\mu_0},
\end{equation}
equivalent to Eq.~\eqref{eq:ch2.electric-wave-operator} with zero forcing \cite{stratton1941,jackson1999}.
The magnetic dual on \(\mathcal{F}_{\mathrm{rad}}\) is
\begin{equation}
\label{eq:ch3.vector-helmholtz-H}
\curl\curl\H-k_0^{2}\H=0,
\end{equation}
obtained symmetrically from Eq.~\eqref{eq:ch2.magnetic-wave-operator} \cite{stratton1941}. Equations
\eqref{eq:ch3.vector-helmholtz-E}--\eqref{eq:ch3.vector-helmholtz-H} are the Subsection~\ref{sec:ch331} HOME
wave operators: every separated eigenfield in Sections~\ref{sec:ch332}--\ref{sec:ch34} satisfies
Eq.~\eqref{eq:ch3.vector-helmholtz-E} or Eq.~\eqref{eq:ch3.vector-helmholtz-H} before angular labels are
attached \cite{hansen1988,tai1994}.

\paragraph{Transverse outgoing constraint.}
Outgoing radiation requires
\begin{equation}
\label{eq:ch3.transverse-radiative-constraint}
\div\E=0,
\qquad
\div\H=0,
\qquad
\E,\H\ \text{transverse to }\hat{\mathbf{r}}\ \text{as }r\to\infty,
\end{equation}
on \(\mathcal{F}_{\mathrm{rad}}\) \cite{stratton1941,jackson1999}. Equation
\eqref{eq:ch3.transverse-radiative-constraint} specializes the divergence rows of
Eq.~\eqref{eq:ch2.harmonic-maxwell-residual} to the export class selected by
Eq.~\eqref{eq:ch3.outgoing-inheritance-criterion} \cite{harrington2001}. Near-field reactive components that
violate transverse decay are excluded before Helmholtz separation is applied \cite{stratton1941}.

\begin{figure}[t]
\centering
\ChOneFigBlock{%
\begin{tikzpicture}[ch1fig, node distance=7mm and 10mm]
  \node[ch1box] (max) {Maxwell pair\\$\mathcal{M}_\omega$};
  \node[ch1box, right=13mm of max] (vh) {Vector Helmholtz\\Eqs.~\eqref{eq:ch3.vector-helmholtz-E}--\eqref{eq:ch3.vector-helmholtz-H}};
  \node[ch1box, right=13mm of vh] (sep) {Section~\ref{sec:ch332}\\separation};
  \node[ch1box, below=11mm of vh] (out) {Eq.~\eqref{eq:ch3.transverse-radiative-constraint}};
  \draw[ch1flow] (max) -- (vh) -- (sep);
  \draw[ch1flow] (out) -- (vh);
\end{tikzpicture}%
}
\caption{Subsection~\ref{sec:ch331} reduction ladder on \(\mathcal{F}_{\mathrm{rad}}\): Maxwell pair
\(\to\) vector Helmholtz \(\to\) coordinate separation, with outgoing transversality enforced before separation.}
\label{fig:ch3.vector-helmholtz-ladder}
\end{figure}

\begin{theorem}[Vector Helmholtz equivalence on \(\mathcal{F}_{\mathrm{rad}}\)]
\label{thm:ch3.vector-helmholtz-equivalence}
If \(\mathbf{u}_\omega=(\E,\H)\in\mathcal{F}_{\mathrm{rad}}\) satisfies
Eq.~\eqref{eq:ch2.harmonic-maxwell-residual} with \(\Jimp=0\) and
Eq.~\eqref{eq:ch3.transverse-radiative-constraint}, then \(\E\) obeys
Eq.~\eqref{eq:ch3.vector-helmholtz-E} and \(\H\) obeys Eq.~\eqref{eq:ch3.vector-helmholtz-H}. Conversely, if
\(\E\) solves Eq.~\eqref{eq:ch3.vector-helmholtz-E} with Eq.~\eqref{eq:ch3.transverse-radiative-constraint} and
\(\H=(\jj\omega\mu_0)^{-1}\curl\E\), the pair satisfies source-free Maxwell in the exterior class
\cite{stratton1941,jackson1999,hansen1988}.
\end{theorem}
\begin{proof}[Proof sketch]
Standard curl--curl reduction imported from Subsection~\ref{sec:ch222}; outgoing class restores the eliminated
partner through curl relations without ambiguity \cite{stratton1941}.
% PROOF: AUTHOR REVIEW
\end{proof}

\paragraph{Component-wise structure and longitudinal suppression.}
Writing \(\E=(E_x,E_y,E_z)\), Eq.~\eqref{eq:ch3.vector-helmholtz-E} is a coupled system for the three scalar
components \cite{stratton1941}. Imposing Eq.~\eqref{eq:ch3.transverse-radiative-constraint} removes the
longitudinal part of \(\E\) in the far zone: any decomposition \(\E=\E_{\perp}+\nabla\Phi\) with
\(\nabla\Phi\sim\mathcal{O}(r^{-2})\) is quotiented out on \(\mathcal{F}_{\mathrm{rad}}\) by
Eq.~\eqref{eq:ch3.outgoing-inheritance-criterion} \cite{harrington2001,jackson1999}. The transverse certificate
can be stated as
\begin{equation}
\label{eq:ch3.transverse-projector-certificate}
\E_{\mathrm{rad}}
=
\big(\mathrm{id}-\hat{\mathbf{r}}\otimes\hat{\mathbf{r}}\big)\E,
\qquad
\lim_{r\to\infty}\|\hat{\mathbf{r}}\cdot\E\|_{L^{2}(\Gamma_{R})}=0,
\end{equation}
aligning the geometric transversality clause in Eq.~\eqref{eq:ch3.transverse-radiative-constraint} with the
operator language of Chapter~\ref{ch:02} \cite{stratton1941,harrington2001}.

\begin{proposition}[Divergence-free preservation under vector Helmholtz flow]
\label{prop:ch3.div-free-preservation}
If \(\E\) solves Eq.~\eqref{eq:ch3.vector-helmholtz-E} on a simply connected domain and \(\div\E=0\) at one
point in the domain, then \(\div\E=0\) everywhere in that domain \cite{stratton1941,jackson1999}. On
\(\mathcal{F}_{\mathrm{rad}}\), the hypothesis is automatic from
Eq.~\eqref{eq:ch3.transverse-radiative-constraint} \cite{harrington2001}.
\end{proposition}
\begin{proof}[Proof sketch]
Take divergence of Eq.~\eqref{eq:ch3.vector-helmholtz-E} and use
\(\div(\curl\curl\E)=0\) \cite{jackson1999}.
% PROOF: AUTHOR REVIEW
\end{proof}

\paragraph{Physical meaning.}
Equation~\eqref{eq:ch3.vector-helmholtz-E} states that each Cartesian component of \(\E\) is not independent:
the curl--curl operator ties components through \(k_0^2\) and divergence constraints, producing the radiative
mode structure observed in far-zone patterns \cite{jackson1999,stratton1941}. On \(\mathcal{F}_{\mathrm{rad}}\),
solving Eq.~\eqref{eq:ch3.vector-helmholtz-E} is the wave-equation certificate that a field pattern can be
expanded in outgoing eigenmodes of Section~\ref{sec:ch34} \cite{hansen1988}. The magnetic partner is not an
independent degree of freedom: once \(\E\) is fixed in the export class, \(\H\) follows uniquely from
\(\H=(\jj\omega\mu_0)^{-1}\curl\E\) \cite{stratton1941,hansen1988}.

\paragraph{Energy bilinear form on source-free Helmholtz pairs.}
Define the source-free Helmholtz pairing
\begin{equation}
\label{eq:ch3.helmholtz-energy-bilinear}
\mathcal{Q}_{\mathrm{vh}}[\E,\E']
=
\int_{\Omega_{R_{\min}}}
\Big(
\curl\E\cdot\curl\E'
-k_{0}^{2}\E\cdot\E'
\Big)\,dV,
\end{equation}
for test pairs \(\E,\E'\in\HCurl(\Omega_{R_{\min}})\) \cite{stratton1941}. On solutions of
Eq.~\eqref{eq:ch3.vector-helmholtz-E}, integration by parts yields
\begin{equation}
\label{eq:ch3.helmholtz-green-identity}
\mathcal{Q}_{\mathrm{vh}}[\E,\E']
=
\int_{\Gamma_{R_{\min}}}
\big(\hat{\mathbf{n}}\times\curl\E\big)\cdot\E'\,dS
\end{equation}
when \(\div\E=\div\E'=0\) \cite{stratton1941,jackson1999}. Equation~\eqref{eq:ch3.helmholtz-green-identity}
is the Subsection~\ref{sec:ch331} surface form that later connects to reciprocity pairings in
Subsection~\ref{sec:ch312} without reopening their proof \cite{harrington2001}. Outgoing fields on
\(\mathcal{F}_{\mathrm{rad}}\) make the surface term in Eq.~\eqref{eq:ch3.helmholtz-green-identity} finite and
radiative rather than reactive \cite{stratton1941}.

\paragraph{Validity domain.}
Eqs.~\eqref{eq:ch3.vector-helmholtz-E}--\eqref{eq:ch3.transverse-radiative-constraint} assume homogeneous vacuum
exterior domains at fixed \(\omega\) \cite{stratton1941}. Piecewise media require interface conditions from
Chapter~\ref{ch:02}; anisotropic constitutive laws break the simple \(k_0^2\) shift \cite{stratton1941,harrington2001}.

\paragraph{Worked illustration (plane-wave consistency).}
A uniform plane wave \(\E=\E_0 e^{\jj k_0\hat{\mathbf{k}}\cdot\mathbf{r}}\) with \(\hat{\mathbf{k}}\cdot\E_0=0\)
satisfies Eq.~\eqref{eq:ch3.vector-helmholtz-E} and Eq.~\eqref{eq:ch3.transverse-radiative-constraint}
\cite{jackson1999}. Reactive quasi-static gradients near a source fail transverse decay and are removed by
\(\Pi_{\mathrm{rad}}\) before Eq.~\eqref{eq:ch3.vector-helmholtz-E} is applied \cite{harrington2001}.

\paragraph{Worked illustration (divergence diagnostic).}
A field with \(\div\E\neq0\) in the far zone cannot belong to \(\mathcal{F}_{\mathrm{rad}}\) even if
\(\|\E\|_{\mathcal{U}}<\infty\) on a finite shell \cite{stratton1941}. Subsection~\ref{sec:ch331} therefore treats
Eq.~\eqref{eq:ch3.transverse-radiative-constraint} as a gate, not a post-processing convenience
\cite{harrington2001}.

\paragraph{Worked illustration (magnetic dual reduction).}
Starting from \(\E=(\jj\omega\mu_0)^{-1}\curl\H\) and Eq.~\eqref{eq:ch3.vector-helmholtz-H} produces the same
separation labels \((\ell,m)\) as the electric route, with coefficients related by duality
Eq.~\eqref{eq:ch3.dual-symmetry-operator} \cite{stratton1941,belinfante1940}. Laboratories that measure \(\H\)
on \(\Gamma_R\) and reconstruct \(\E\) by curl must preserve outgoing class before applying
Eq.~\eqref{eq:ch3.vector-helmholtz-E}; otherwise the reconstructed \(\E\) solves a standing-wave operator
\cite{harrington2001}.

\paragraph{Worked illustration (operator symbol alignment).}
The symbol \(k_0^2\) in Eqs.~\eqref{eq:ch3.vector-helmholtz-E}--\eqref{eq:ch3.vector-helmholtz-H} matches
Eq.~\eqref{eq:ch2.helmholtz-symbol} and the vacuum pair Eq.~\eqref{eq:ch1.vacuum-helmholtz-pair}
\cite{stratton1941}. Dispersive media replace \(k_0^2\) by \(\omega^2\varepsilon(\omega)\mu(\omega)\) in
later chapters; Subsection~\ref{sec:ch331} fixes the vacuum symbol that Section~\ref{sec:ch34} modes carry
\cite{jackson1999,harrington2001}.

\paragraph{Symmetry-operator compatibility preview.}
Simultaneous eigenfields Eq.~\eqref{eq:ch3.J-simultaneous-eigenfield-system} are sought among solutions of
Eq.~\eqref{eq:ch3.vector-helmholtz-E} because \(\mathbf{J}\) commutes with the Maxwell system on
\(\mathcal{D}(\mathbf{J})\) from Eq.~\eqref{eq:ch3.J-operator-domain-rad} \cite{jackson1999,hansen1988}.
Subsection~\ref{sec:ch331} does not prove that commutator here; it supplies the wave operator on which the
commutation results of Subsection~\ref{sec:ch322} act \cite{stratton1941}.

\paragraph{Problem (Maxwell pair \(\to\) vector Helmholtz on \(\mathcal{F}_{\mathrm{rad}}\)).}
\textbf{Problem.} Let \(\mathbf{u}_\omega=(\E,\H)\in\mathcal{F}_{\mathrm{rad}}\) satisfy source-free harmonic
Maxwell with \(\Jimp=0\) and outgoing inheritance Eq.~\eqref{eq:ch3.outgoing-inheritance-criterion}. Show that
\(\E\) obeys Eq.~\eqref{eq:ch3.vector-helmholtz-E} and \(\H\) obeys Eq.~\eqref{eq:ch3.vector-helmholtz-H}
\cite{stratton1941,jackson1999}.

\textbf{Step 1 (curl partner).} From \(\curl\E=\jj\omega\mu_0\H\) and \(\curl\H=-\jj\omega\varepsilon_0\E\),
eliminate \(\H\) to obtain \(\curl\curl\E=\omega^{2}\varepsilon_0\mu_0\E=k_{0}^{2}\E\) \cite{stratton1941}.

\textbf{Step 2 (divergence gate).} Take divergence of \(\curl\curl\E=k_{0}^{2}\E\) and use
\(\div(\curl\curl\E)=0\) to infer \(\nabla^{2}(\div\E)=k_{0}^{2}(\div\E)\) \cite{jackson1999}. On outgoing
exterior domains, \(\div\E=0\) is enforced by Eq.~\eqref{eq:ch3.transverse-radiative-constraint}, so the
longitudinal Laplacian equation admits only the trivial longitudinal mode in \(\mathcal{F}_{\mathrm{rad}}\)
\cite{stratton1941,harrington2001}.

\textbf{Step 3 (identity collapse).} With \(\div\E=0\), Eq.~\eqref{eq:ch3.curl-curl-vector-identity} gives
\(\curl\curl\E=\nabla^{2}\E\), reproducing Eq.~\eqref{eq:ch3.vector-helmholtz-E} as an equivalent form used in
some references \cite{jackson1999,stratton1941}.

\textbf{Step 4 (magnetic dual).} Repeat with \(\E=(\jj\omega\varepsilon_0)^{-1}\curl\H\) to obtain
Eq.~\eqref{eq:ch3.vector-helmholtz-H} \cite{stratton1941}.

\textbf{Physical meaning.} Steps 1--4 show that outgoing transversality is not optional decoration: it is the
hypothesis that makes the curl--curl and vector Laplacian formulations equivalent on \(\mathcal{F}_{\mathrm{rad}}\)
\cite{harrington2001,stratton1941}. Near reactive zones may violate Step~2 unless \(\Pi_{\mathrm{rad}}\) is
applied before export claims \cite{harrington2001}.

\paragraph{Problem (recovering \(\H\) without branch ambiguity).}
\textbf{Problem.} Given \(\E\) solving Eq.~\eqref{eq:ch3.vector-helmholtz-E} with
Eq.~\eqref{eq:ch3.transverse-radiative-constraint} on an exterior domain, construct \(\H\) and verify source-free
Maxwell \cite{stratton1941}.

\textbf{Step 1.} Define \(\H=(\jj\omega\mu_0)^{-1}\curl\E\) \cite{stratton1941}.

\textbf{Step 2.} Compute \(\curl\H=(\jj\omega\mu_0)^{-1}\curl\curl\E=(\jj\omega\mu_0)^{-1}k_{0}^{2}\E
=-\jj\omega\varepsilon_0\E\) using Eq.~\eqref{eq:ch3.vector-helmholtz-E} \cite{jackson1999}.

\textbf{Step 3.} Verify \(\div\H=0\) from \(\div(\curl\E)=0\) \cite{stratton1941}.

\textbf{Physical meaning.} On \(\mathcal{F}_{\mathrm{rad}}\), \(\H\) is a slave field to transverse \(\E\); this
one-way reconstruction is why eigenfield expansions in Section~\ref{sec:ch34} can be organized around a single
vector potential alphabet \cite{hansen1988,tai1994}.

\paragraph{Spectral symbol alignment with Chapter~\ref{ch:02}.}
The operator symbol Eq.~\eqref{eq:ch2.helmholtz-symbol} attached to Eq.~\eqref{eq:ch2.electric-wave-operator}
is exactly the \(k_{0}^{2}\) shift in Eq.~\eqref{eq:ch3.vector-helmholtz-E} \cite{stratton1941}. Domain
Eq.~\eqref{eq:ch2.electric-helmholtz-domain} therefore governs where Eq.~\eqref{eq:ch3.vector-helmholtz-E} may
be applied without regularization \cite{harrington2001}. Subsection~\ref{sec:ch331} records that alignment so
Subsections~\ref{sec:ch332}--\ref{sec:ch334} need not reopen operator-domain arguments \cite{stratton1941}.

\paragraph{Cartesian component form (reference block).}
Expanding Eq.~\eqref{eq:ch3.vector-helmholtz-E} in Cartesian coordinates gives, for each component
\(F\in\{E_x,E_y,E_z\}\),
\begin{equation}
\label{eq:ch3.cartesian-component-helmholtz}
\nabla^{2}E_{i}+k_{0}^{2}E_{i}
=
\frac{\partial(\div\E)}{\partial x_{i}},
\qquad
i\in\{x,y,z\},
\end{equation}
with coupled right-hand sides among components when \(\div\E\neq0\) \cite{stratton1941,jackson1999}. On
\(\mathcal{F}_{\mathrm{rad}}\), Eq.~\eqref{eq:ch3.transverse-radiative-constraint} forces the right-hand side to
vanish, reducing Eq.~\eqref{eq:ch3.cartesian-component-helmholtz} to three scalar Helmholtz equations tied only
through the transverse projection Eq.~\eqref{eq:ch3.transverse-projector-certificate} \cite{harrington2001}.
Subsection~\ref{sec:ch332} moves to spherical coordinates precisely because the Cartesian form hides the
\(\mathbf{J}^{2}\) and \(\mathbf{J}_{z}\) quantum numbers \cite{hansen1988,jackson1999}.

\paragraph{Well-posedness inheritance (USE).}
Exterior uniqueness Theorem~\ref{thm:ch2.exterior-radiative-uniqueness} applies to Maxwell pairs before
reduction; Theorem~\ref{thm:ch3.vector-helmholtz-equivalence} transfers that uniqueness to transverse
\(\E\) solving Eq.~\eqref{eq:ch3.vector-helmholtz-E} \cite{stratton1941,harrington2001}. Subsection~\ref{sec:ch331}
therefore does not prove a separate Helmholtz uniqueness theorem; it cites the Maxwell result through
Eq.~\eqref{eq:ch3.outgoing-inheritance-criterion} \cite{stratton1941}.

\paragraph{Export map compatibility.}
Radiation operator \(\mathcal{R}_\omega\) from Eq.~\eqref{eq:ch2.outgoing-radiation-refinement} maps impressed data
to \(\mathcal{F}_{\mathrm{rad}}\) before any separation is applied \cite{harrington2001,stratton1941}. Vector
Helmholtz reduction is thus a \emph{post-export} analytic step: it organizes \(\mathcal{R}_\omega[\mathfrak{D}]\)
into modes without changing the forward operator definition from Chapter~\ref{ch:02} \cite{stratton1941}. This
ordering prevents the common error of expanding non-outgoing near fields in outgoing spherical bases
\cite{harrington2001}.

\begin{table}[t]
\centering
\caption{Vector Helmholtz reduction gates for Subsection~\ref{sec:ch331}.}
\label{tab:ch3.331-helmholtz-gates}
\small
\begin{tabular}{@{}p{0.22\linewidth}p{0.34\linewidth}p{0.36\linewidth}@{}}
\toprule
Gate & Anchor & Failure symptom \\
\midrule
Operator domain &
Eq.~\eqref{eq:ch2.electric-helmholtz-domain} &
Near-field \(\curl\curl\) blow-up \\
\addlinespace
Outgoing class &
Eq.~\eqref{eq:ch3.outgoing-inheritance-criterion} &
Evanescent fit labeled radiation \\
\addlinespace
Transversality &
Eq.~\eqref{eq:ch3.transverse-radiative-constraint} &
Longitudinal leak in far zone \\
\bottomrule
\end{tabular}
\end{table}

\paragraph{Reactive versus radiative split before reduction.}
Near a finite source, \(\mathcal{S}_\omega[\mathfrak{D}]\) generally contains reactive storage not visible to
\(\Pi_{\mathrm{rad}}\) \cite{harrington2001,stratton1941}. Vector Helmholtz reduction applies to the
\emph{organizing} of \(\mathcal{R}_\omega[\mathfrak{D}]\), not to raw \(\mathcal{S}_\omega[\mathfrak{D}]\) samples
\cite{stratton1941}. The split can be written as
\begin{equation}
\label{eq:ch3.reactive-radiative-helmholtz-split}
\mathbf{u}_\omega
=
\mathbf{u}_{\mathrm{re}}+\mathbf{u}_{\mathrm{rad}},
\qquad
\mathbf{u}_{\mathrm{rad}}=\Pi_{\mathrm{rad}}[\mathbf{u}_\omega]\in\mathcal{F}_{\mathrm{rad}},
\qquad
\mathbf{u}_{\mathrm{re}}=\big(\mathrm{id}-\Pi_{\mathrm{rad}}\big)\mathbf{u}_\omega,
\end{equation}
with Eq.~\eqref{eq:ch3.vector-helmholtz-E} enforced on \(\mathbf{u}_{\mathrm{rad}}\) after outgoing qualification
\cite{harrington2001}. Applying modal analysis to \(\mathbf{u}_{\mathrm{re}}\) without labeling it as non-exported
reactive content is a common source of spurious high-\(\ell\) coefficients in near-zone fits \cite{stratton1941}.

\begin{figure}[t]
\centering
\ChOneFigBlock{%
\begin{tikzpicture}[ch1fig, node distance=7mm and 9mm]
  \node[ch1box] (s) {$\mathcal{S}_\omega[\mathfrak{D}]$};
  \node[ch1box, right=12mm of s] (r) {$\mathcal{R}_\omega[\mathfrak{D}]$};
  \node[ch1box, right=12mm of r] (vh) {Eq.~\eqref{eq:ch3.vector-helmholtz-E}};
  \node[ch1box, below=10mm of s] (re) {$\mathbf{u}_{\mathrm{re}}$\\non-exported};
  \draw[ch1flow] (s) -- (r) -- (vh);
  \draw[ch1flow] (re) -- (s);
\end{tikzpicture}%
}
\caption{Subsection~\ref{sec:ch331} applies vector Helmholtz reduction on the radiation image
\(\mathcal{R}_\omega[\mathfrak{D}]\), not on the full solution \(\mathcal{S}_\omega[\mathfrak{D}]\) that includes
reactive \(\mathbf{u}_{\mathrm{re}}\).}
\label{fig:ch3.reactive-radiative-helmholtz-split}
\end{figure}

\paragraph{Poynting transport certificate on Helmholtz solutions.}
For \(\mathbf{u}_\omega\in\mathcal{F}_{\mathrm{rad}}\) satisfying Eqs.~\eqref{eq:ch3.vector-helmholtz-E}--\eqref{eq:ch3.transverse-radiative-constraint},
the time-averaged Poynting flux through \(\Gamma_R\) is
\begin{equation}
\label{eq:ch3.helmholtz-poynting-flux}
P_{\Gamma_R}
=
\frac{1}{2}\Re\oint_{\Gamma_R}
\mathbf{S}\cdot\hat{\mathbf{n}}\,dS,
\qquad
\mathbf{S}=\tfrac12\E\times\H^{*},
\end{equation}
finite and outward-directed for outgoing Helmholtz pairs \cite{stratton1941,poynting1884}. Equation
\eqref{eq:ch3.helmholtz-poynting-flux} connects the analytic reduction of Subsection~\ref{sec:ch331} to the flux
normalizations imported from Subsection~\ref{sec:ch253} without reopening their derivation \cite{stratton1941}. A
field that satisfies Eq.~\eqref{eq:ch3.vector-helmholtz-E} on a bounded shell but yields
\(P_{\Gamma_R}\not>0\) after outgoing selection is not a radiative export candidate \cite{harrington2001}.

\paragraph{Problem (left-circular plane wave in Helmholtz form).}
\textbf{Problem.} Verify Eqs.~\eqref{eq:ch3.vector-helmholtz-E}--\eqref{eq:ch3.transverse-radiative-constraint} for
\(\E=(E_{0}/\sqrt{2})(\hat{\mathbf{x}}+\jj\hat{\mathbf{y}})e^{\jj k_{0}z}\) \cite{jackson1999}.

\textbf{Step 1.} Compute \(\curl\E=\jj k_{0}\E\) and \(\curl\curl\E=-k_{0}^{2}\E\), giving
Eq.~\eqref{eq:ch3.vector-helmholtz-E} \cite{jackson1999}.

\textbf{Step 2.} Check \(\hat{\mathbf{k}}\cdot\E=0\) so \(\div\E=0\) \cite{stratton1941}.

\textbf{Step 3.} As \(z\to+\infty\), \(\E\) is transverse to \(\hat{\mathbf{r}}\approx\hat{\mathbf{z}}\) and outgoing
under Eq.~\eqref{eq:ch2.sommerfeld-field-condition} \cite{jackson1999,harrington2001}.

\textbf{Physical meaning.} The example shows that polarization structure is carried \emph{inside} transverse
Helmholtz solutions; separation in Section~\ref{sec:ch332} will expose the same \((\ell,m)\) labels through angular
eigenfunctions rather than through Cartesian Jones vectors \cite{hansen1988,stratton1941}.

\begin{table}[t]
\centering
\caption{Common Subsection~\ref{sec:ch331} failure modes when expanding measured fields.}
\label{tab:ch3.331-failure-modes}
\small
\begin{tabular}{@{}p{0.26\linewidth}p{0.32\linewidth}p{0.34\linewidth}@{}}
\toprule
Misstep & Symptom & Correct anchor \\
\midrule
Skip \(\Pi_{\mathrm{rad}}\) &
Large reactive fit coefficients &
Eq.~\eqref{eq:ch3.reactive-radiative-helmholtz-split} \\
\addlinespace
Ignore \(\div\E=0\) &
Spurious monopole leak &
Eq.~\eqref{eq:ch3.transverse-radiative-constraint} \\
\addlinespace
Use standing \(j_\ell\) on exterior &
Nonzero incoming flux &
Eq.~\eqref{eq:ch3.outgoing-inheritance-criterion} \\
\bottomrule
\end{tabular}
\end{table}

\paragraph{Vacuum Helmholtz pair restatement (USE).}
Chapter~\ref{ch:01} recorded the vacuum symbol shift Eq.~\eqref{eq:ch1.vacuum-helmholtz-pair}; Subsection~\ref{sec:ch331}
specializes that symbol to vector fields on \(\mathcal{F}_{\mathrm{rad}}\) \cite{stratton1941,jackson1999}. The
electric and magnetic reductions Eqs.~\eqref{eq:ch3.vector-helmholtz-E}--\eqref{eq:ch3.vector-helmholtz-H} are
not independent postulates: they are two projections of the same source-free Maxwell system under
Eq.~\eqref{eq:ch3.transverse-radiative-constraint} \cite{hansen1988,stratton1941}. Any attempt to impose
Eq.~\eqref{eq:ch3.vector-helmholtz-E} on a reactive near field without \(\Pi_{\mathrm{rad}}\) produces
coefficients that do not correspond to \(\mathcal{R}_\omega[\mathfrak{D}]\) \cite{harrington2001}.

Subsection~\ref{sec:ch331} has reduced admissible Maxwell pairs on \(\mathcal{F}_{\mathrm{rad}}\) to vector
Helmholtz form Eqs.~\eqref{eq:ch3.vector-helmholtz-E}--\eqref{eq:ch3.transverse-radiative-constraint}
\cite{stratton1941,hansen1988}. Subsection~\ref{sec:ch332} separates Eq.~\eqref{eq:ch3.vector-helmholtz-E} in
spherical coordinates matched to Section~\ref{sec:ch32} \cite{jackson1999}.


\subsection{Separation in spherical coordinates}
\label{sec:ch332}
% TOC tag: [HOME]

Spherical coordinates align the separation geometry with the orbital operators of Section~\ref{sec:ch32}.
Subsection~\ref{sec:ch332} separates Eq.~\eqref{eq:ch3.vector-helmholtz-E} in \((r,\theta,\phi)\) without
introducing named spherical-harmonic symbols; those definitions are reserved for Subsection~\ref{sec:ch341}
\cite{stratton1941,hansen1988,jackson1999}.

\paragraph{Coordinate chart and metric.}
Let \(\mathbf{r}=r\hat{\mathbf{r}}\) with \(r\ge0\), polar angle \(\theta\in[0,\pi]\), and azimuth
\(\phi\in[0,2\pi)\). The orthonormal spherical triad \((\hat{\mathbf{r}},\hat{\boldsymbol{\theta}},\hat{\boldsymbol{\phi}})\)
satisfies \(\hat{\mathbf{r}}\times\hat{\boldsymbol{\theta}}=\hat{\boldsymbol{\phi}}\) and carries the line element
\(ds^{2}=dr^{2}+r^{2}d\theta^{2}+r^{2}\sin^{2}\theta\,d\phi^{2}\) \cite{jackson1999}. All separation below is
performed on source-free domains outside compact supports, with pairs restricted to
\(\mathcal{F}_{\mathrm{rad}}\) through Eq.~\eqref{eq:ch3.outgoing-inheritance-criterion}
\cite{stratton1941,harrington2001}.

\paragraph{Scalar Helmholtz separation (HOME).}
Begin with the scalar prototype of Eq.~\eqref{eq:ch3.vector-helmholtz-E}. Seek
\begin{equation}
\label{eq:ch3.scalar-separation-ansatz}
u(r,\theta,\phi)=R_{\ell}(r)\,Y_{\ell m}(\theta,\phi),
\end{equation}
where \(Y_{\ell m}\) denotes an angular eigenfunction of the Laplace--Beltrami operator on the unit sphere
\cite{stratton1941,hansen1988}. Substituting Eq.~\eqref{eq:ch3.scalar-separation-ansatz} into
\(\nabla^{2}u+k_{0}^{2}u=0\) yields the radial equation
\begin{equation}
\label{eq:ch3.radial-helmholtz-ODE}
\frac{1}{r^{2}}\frac{d}{dr}\!\left(r^{2}\frac{dR_{\ell}}{dr}\right)
+\left(k_{0}^{2}-\frac{\ell(\ell+1)}{r^{2}}\right)R_{\ell}=0,
\end{equation}
and the angular equation
\begin{equation}
\label{eq:ch3.angular-laplacian-eigenvalue}
\frac{1}{\sin\theta}\frac{\partial}{\partial\theta}
\!\left(\sin\theta\frac{\partial Y_{\ell m}}{\partial\theta}\right)
+\frac{1}{\sin^{2}\theta}\frac{\partial^{2}Y_{\ell m}}{\partial\phi^{2}}
=-\ell(\ell+1)\,Y_{\ell m}.
\end{equation}
Equations~\eqref{eq:ch3.radial-helmholtz-ODE}--\eqref{eq:ch3.angular-laplacian-eigenvalue} are the
Subsection~\ref{sec:ch332} HOME separation certificates: the angular quantum number \(\ell\) enters only
through the centrifugal eigenvalue \(\ell(\ell+1)\), while \(m\) labels azimuthal dependence through
\(\partial/\partial\phi\) \cite{jackson1999,stratton1941}.

\paragraph{Derivation skeleton of the radial equation.}
Substitute Eq.~\eqref{eq:ch3.scalar-separation-ansatz} into \(\nabla^{2}u+k_{0}^{2}u=0\) in spherical form
\begin{equation}
\label{eq:ch3.laplacian-spherical-form}
\nabla^{2}u
=
\frac{1}{r^{2}}\frac{\partial}{\partial r}\!\left(r^{2}\frac{\partial u}{\partial r}\right)
+\frac{1}{r^{2}\sin\theta}\frac{\partial}{\partial\theta}
\!\left(\sin\theta\frac{\partial u}{\partial\theta}\right)
+\frac{1}{r^{2}\sin^{2}\theta}\frac{\partial^{2}u}{\partial\phi^{2}},
\end{equation}
divide by \(Y_{\ell m}\), and use Eq.~\eqref{eq:ch3.angular-laplacian-eigenvalue} to isolate the
\(r\)-dependent terms \cite{jackson1999,stratton1941}. The result is Eq.~\eqref{eq:ch3.radial-helmholtz-ODE},
equivalently written in canonical spherical form
\begin{equation}
\label{eq:ch3.radial-helmholtz-canonical}
\frac{d^{2}R_{\ell}}{dr^{2}}
+\frac{2}{r}\frac{dR_{\ell}}{dr}
+\left(k_{0}^{2}-\frac{\ell(\ell+1)}{r^{2}}\right)R_{\ell}=0.
\end{equation}
Equations~\eqref{eq:ch3.radial-helmholtz-ODE}--\eqref{eq:ch3.radial-helmholtz-canonical} are the radial
certificates imported by Subsection~\ref{sec:ch333} when selecting \(h_{\ell}^{(1)}\) \cite{stratton1941}.

\paragraph{Azimuthal factor and \(\mathbf{J}_{z}\) alignment.}
Because \(\partial/\partial\phi\) commutes with the scalar Laplacian on the sphere,
\begin{equation}
\label{eq:ch3.azimuthal-eigenvalue-equation}
-\jj\frac{\partial Y_{\ell m}}{\partial\phi}=m\,Y_{\ell m},
\qquad
m\in\{-\ell,\ldots,\ell\},
\end{equation}
matching the orbital \(z\)-component action imported from Eq.~\eqref{eq:ch3.J-simultaneous-eigenfield-system}
\cite{jackson1999,hansen1988}. Equation~\eqref{eq:ch3.azimuthal-eigenvalue-equation} is the coordinate
realization of simultaneous diagonalization previewed in Eq.~\eqref{eq:ch3.eigenfield-quantum-number-split}:
separation in \(\phi\) is not an independent modeling choice but the eigenfunction form required by
\(\mathbf{J}_{z}\) on \(\mathcal{F}_{\mathrm{rad}}\) \cite{stratton1941}.

\begin{figure}[t]
\centering
\ChOneFigBlock{%
\begin{tikzpicture}[ch1fig, node distance=7mm and 9mm]
  \node[ch1box] (vh) {Eq.~\eqref{eq:ch3.vector-helmholtz-E}};
  \node[ch1box, right=12mm of vh] (ang) {Eq.~\eqref{eq:ch3.angular-laplacian-eigenvalue}\\$\ell,m$};
  \node[ch1box, right=12mm of ang] (rad) {Eq.~\eqref{eq:ch3.radial-helmholtz-ODE}\\$R_{\ell}(r)$};
  \node[ch1box, below=10mm of ang] (jz) {Eq.~\eqref{eq:ch3.azimuthal-eigenvalue-equation}};
  \draw[ch1flow] (vh) -- (ang) -- (rad);
  \draw[ch1flow] (jz) -- (ang);
\end{tikzpicture}%
}
\caption{Scalar separation ladder in Subsection~\ref{sec:ch332}: vector Helmholtz reduces to angular
Eq.~\eqref{eq:ch3.angular-laplacian-eigenvalue} and radial Eq.~\eqref{eq:ch3.radial-helmholtz-ODE}, with
azimuthal quantum number fixed by Eq.~\eqref{eq:ch3.azimuthal-eigenvalue-equation}.}
\label{fig:ch3.spherical-separation-ladder}
\end{figure}

\paragraph{Vector separation without premature naming.}
Vector fields on \(\mathcal{F}_{\mathrm{rad}}\) are transverse: each admissible \(\E\) admits an expansion
\begin{equation}
\label{eq:ch3.vector-field-spherical-components}
\E(r,\theta,\phi)
=
E_{r}\,\hat{\mathbf{r}}+E_{\theta}\,\hat{\boldsymbol{\theta}}+E_{\phi}\,\hat{\boldsymbol{\phi}},
\qquad
E_{r}\to0\ \text{faster than }r^{-1}\ \text{as }r\to\infty,
\end{equation}
with analogous decay for \(\H\) \cite{stratton1941,jackson1999}. Inserting
Eq.~\eqref{eq:ch3.vector-field-spherical-components} into Eq.~\eqref{eq:ch3.vector-helmholtz-E} produces a
coupled system for \((E_{r},E_{\theta},E_{\phi})\). The standard resolution---recorded here as a construction
target for Subsection~\ref{sec:ch334}---represents \(\E\) through two scalar potentials whose curls and
double-curls separate with the same \(\ell,m\) indices as Eq.~\eqref{eq:ch3.scalar-separation-ansatz}
\cite{stratton1941,tai1994}. Subsection~\ref{sec:ch332} therefore stops at the scalar separation
Eqs.~\eqref{eq:ch3.radial-helmholtz-ODE}--\eqref{eq:ch3.azimuthal-eigenvalue-equation}; named vector
eigenfunctions \(\mathbf{M}_{\ell m}\), \(\mathbf{N}_{\ell m}\) are deferred to Subsection~\ref{sec:ch341}
\cite{hansen1988}.

\begin{theorem}[Separability of the transverse vector Helmholtz problem]
\label{thm:ch3.vector-separability}
On a source-free exterior domain with Eq.~\eqref{eq:ch3.transverse-radiative-constraint}, every
\(\E\in\mathcal{F}_{\mathrm{rad}}\) solving Eq.~\eqref{eq:ch3.vector-helmholtz-E} admits an expansion in
angular eigenfunctions \(Y_{\ell m}\) with radial factors \(R_{\ell}(r)\) determined by
Eq.~\eqref{eq:ch3.radial-helmholtz-ODE}. The magnetic field follows from
\(\H=(\jj\omega\mu_{0})^{-1}\curl\E\) with the same \((\ell,m)\) labels \cite{stratton1941,hansen1988,tai1994}.
\end{theorem}
\begin{proof}[Proof sketch]
Complete the separation using Debye potentials of Subsection~\ref{sec:ch334}; angular completeness on the
sphere is imported from Subsection~\ref{sec:ch343} as a forward citation \cite{stratton1941,hansen1988}.
% PROOF: AUTHOR REVIEW
\end{proof}

\paragraph{Centrifugal barrier and mode order.}
The term \(\ell(\ell+1)/r^{2}\) in Eq.~\eqref{eq:ch3.radial-helmholtz-ODE} is the centrifugal barrier
familiar from quantum angular momentum \cite{jackson1999}. Low-\(\ell\) modes penetrate closer to the origin
in the oscillatory region; high-\(\ell\) modes are pushed outward, explaining why large \(\ell\) content in
far-zone patterns decays rapidly for compact sources \cite{stratton1941,harrington2001}. This geometric
bias is independent of polarization label \(\sigma\) from Section~\ref{sec:ch32} \cite{hansen1988}.

\paragraph{Angular degeneracy and \(2\ell+1\) multiplicity.}
For each \(\ell\), Eq.~\eqref{eq:ch3.azimuthal-eigenvalue-equation} admits \(2\ell+1\) independent azimuthal
indices \(m\in\{-\ell,\ldots,\ell\}\) \cite{jackson1999,hansen1988}. The SO\((3)\) irrep dimension previewed
in Subsection~\ref{sec:ch323} therefore appears already at the separation stage: the angular block is
\((2\ell+1)\)-fold degenerate before polarization or dual labels are attached \cite{stratton1941}. Wigner
rotation matrices Eq.~\eqref{eq:ch3.Wigner-D-coefficient-law} mix \(m\) within fixed \(\ell\) without changing
the radial equation \cite{hansen1988,jackson1999}.

\paragraph{Vector Laplacian reduction target.}
In spherical components, the vector Laplacian acting on \(\E\) produces couplings among
\((E_{r},E_{\theta},E_{\phi})\) \cite{stratton1941,hansen1988}. Debye potentials of Subsection~\ref{sec:ch334}
diagonalize that coupling by representing \(\E\) through scalars that already satisfy
Eq.~\eqref{eq:ch3.angular-laplacian-eigenvalue} \cite{tai1994}. Subsection~\ref{sec:ch332} therefore stops at
scalar separation and records the vector reduction as a forward pointer rather than a component-wise grind
\cite{hansen1988}.

\begin{proposition}[\(m\)-conservation under \(\mathbf{J}_{z}\) rotations]
\label{prop:ch3.m-conservation-separation}
If \(Y_{\ell m}\) satisfies Eq.~\eqref{eq:ch3.azimuthal-eigenvalue-equation}, then
\(e^{-\jj\alpha\mathbf{J}_{z}}Y_{\ell m}=e^{\jj m\alpha}Y_{\ell m}\) for \(\alpha\in\mathbb{R}\)
\cite{jackson1999}. Separated modes with fixed \(m\) are eigenfunctions of azimuthal transport around the
\(z\)-axis, matching the phase ramps measured in rotated antenna arrays \cite{harrington2001,stratton1941}.
\end{proposition}
\begin{proof}[Proof sketch]
Direct integration of Eq.~\eqref{eq:ch3.azimuthal-eigenvalue-equation} in \(\phi\) \cite{jackson1999}.
% PROOF: AUTHOR REVIEW
\end{proof}

\paragraph{Validity and coordinate singularities.}
Separation at \(\theta=0,\pi\) requires continuity conditions on \(Y_{\ell m}\); at \(r=0\) the exterior
problem never evaluates fields because sources are excluded \cite{jackson1999}. Axisymmetric feeds with
\(m=0\) remain smooth on the polar axis when \(Y_{\ell 0}\) carries the standard Legendre factors imported
in Subsection~\ref{sec:ch341} \cite{stratton1941}. Multi-connected domains with cuts in \(\phi\) require
periodicity \(Y_{\ell m}(\phi+2\pi)=e^{2\pi\jj m}Y_{\ell m}(\phi)\), consistent with
Eq.~\eqref{eq:ch3.azimuthal-eigenvalue-equation} \cite{hansen1988}.

\paragraph{Worked illustration (dipole angular factor).}
The \(\ell=1\) angular equation admits three azimuthal levels \(m=-1,0,+1\). A short dipole aligned with
\(\hat{\mathbf{z}}\) selects primarily \(m=0\) content in the far zone, while a loop-like current distribution
mixes \(m=\pm1\) through \(\partial/\partial\phi\) coupling \cite{jackson1999,stratton1941}. The separation
equations make the selection explicit: mislabeling a measured pattern as ``\(m=5\)'' without verifying
Eq.~\eqref{eq:ch3.azimuthal-eigenvalue-equation} on the exported field is a category error
\cite{harrington2001}.

\paragraph{Worked illustration (sectoral vs tesseral labels).}
Sectoral harmonics with \(|m|=\ell\) peak on equatorial circles; tesseral modes with \(|m|<\ell\) distribute
nodal cones at intermediate latitudes \cite{stratton1941,hansen1988}. Subsection~\ref{sec:ch332} records only
the eigenvalue Eqs.~\eqref{eq:ch3.angular-laplacian-eigenvalue}--\eqref{eq:ch3.azimuthal-eigenvalue-equation};
explicit \(Y_{\ell m}\) phase conventions are fixed once in Subsection~\ref{sec:ch341} \cite{hansen1988}.

\paragraph{Worked illustration (monopole angular block).}
The \(\ell=0\) angular equation has \(m=0\) only and \(Y_{00}\) constant on the sphere \cite{jackson1999}.
Radial Eq.~\eqref{eq:ch3.radial-helmholtz-canonical} then reduces to the isotropic outgoing mode of
Subsection~\ref{sec:ch333}, the far-zone carrier of compact sources \cite{stratton1941}. Higher \(\ell\) adds
angular structure without changing the outgoing radial gate \cite{hansen1988}.

\paragraph{Worked illustration (array phase progression).}
A uniform circular array with element phase \(\propto e^{\jj m\phi_n}\) excites primarily the \(m\) sector of
Eq.~\eqref{eq:ch3.azimuthal-eigenvalue-equation} when the array radius is comparable to \(\lambda_0\)
\cite{harrington2001,jackson1999}. Separation makes the selectivity explicit: a measured \(m\) index is an
eigenvalue label, not a fitting parameter \cite{stratton1941}.

\paragraph{Legendre backbone without explicit display.}
For \(m=0\), Eq.~\eqref{eq:ch3.angular-laplacian-eigenvalue} reduces to the Legendre equation for
\(P_{\ell}(\cos\theta)\) \cite{jackson1999,stratton1941}. For \(m\neq0\), the \(\phi\) factor
\(e^{\jj m\phi}\) multiplies an associated Legendre factor in \(Y_{\ell m}\) defined in Subsection~\ref{sec:ch341}
\cite{hansen1988}. Subsection~\ref{sec:ch332} uses only the eigenvalue form so that Condon--Shortley phase and
normalization are not duplicated \cite{hansen1988,stratton1941}.

\paragraph{Problem (separate a scalar Helmholtz mode).}
\textbf{Problem.} Separate \(\nabla^{2}u+k_{0}^{2}u=0\) with \(u\in L^{2}_{\mathrm{loc}}\) on an exterior
domain and identify quantum numbers \cite{jackson1999,stratton1941}.

\textbf{Step 1.} Posit Eq.~\eqref{eq:ch3.scalar-separation-ansatz} and insert into
Eq.~\eqref{eq:ch3.laplacian-spherical-form} \cite{stratton1941}.

\textbf{Step 2.} Divide by \(R_{\ell}Y_{\ell m}\) and equate \(\phi\)-dependent terms to obtain
Eq.~\eqref{eq:ch3.azimuthal-eigenvalue-equation} \cite{jackson1999}.

\textbf{Step 3.} Equate \(\theta\)-dependent terms to obtain Eq.~\eqref{eq:ch3.angular-laplacian-eigenvalue}
\cite{stratton1941,hansen1988}.

\textbf{Step 4.} The remaining \(r\)-equation is Eq.~\eqref{eq:ch3.radial-helmholtz-canonical}
\cite{jackson1999}.

\textbf{Physical meaning.} Steps 2--3 show that \(m\) and \(\ell\) are separation constants fixed by angular
symmetry, not fit parameters in a spherical harmonic synthesis \cite{harrington2001,stratton1941}.

\paragraph{Orthogonality preview (forward citation).}
Angular eigenfunctions satisfying Eqs.~\eqref{eq:ch3.angular-laplacian-eigenvalue}--\eqref{eq:ch3.azimuthal-eigenvalue-equation}
are orthogonal on \(S^{2}\) with weight \(\sin\theta\) once normalization is fixed in Subsection~\ref{sec:ch341}
\cite{hansen1988,stratton1941}. Subsection~\ref{sec:ch332} does not display inner products; Subsection~\ref{sec:ch343}
proves completeness and resolution of identity on \(\mathcal{F}_{\mathrm{rad}}\) \cite{hansen1988}.

\begin{table}[t]
\centering
\caption{Separation outputs of Subsection~\ref{sec:ch332} and their Section~\ref{sec:ch32} partners.}
\label{tab:ch3.332-separation-alignment}
\small
\begin{tabular}{@{}p{0.24\linewidth}p{0.32\linewidth}p{0.36\linewidth}@{}}
\toprule
Separated object & Equation & Operator partner \\
\midrule
Angular eigenvalue &
Eq.~\eqref{eq:ch3.angular-laplacian-eigenvalue} &
\(\mathbf{J}^{2}\to\ell(\ell+1)\) \\
\addlinespace
Azimuthal index &
Eq.~\eqref{eq:ch3.azimuthal-eigenvalue-equation} &
\(\mathbf{J}_{z}\to m\) \\
\addlinespace
Radial factor &
Eq.~\eqref{eq:ch3.radial-helmholtz-ODE} &
Outgoing \(R_{\ell}\) in ?\ref{sec:ch333} \\
\bottomrule
\end{tabular}
\end{table}

\paragraph{Spherical vector operators (structural preview).}
The gradient, curl, and Laplacian on vector fields in \((r,\theta,\phi)\) mix \(\hat{\mathbf{r}}\),
\(\hat{\boldsymbol{\theta}}\), and \(\hat{\boldsymbol{\phi}}\) components \cite{hansen1988,stratton1941}. Debye
potentials Eqs.~\eqref{eq:ch3.debye-E-representation}--\eqref{eq:ch3.debye-H-representation} are introduced
precisely to avoid displaying the full \(3\times3\) coupling matrix at this stage \cite{tai1994,stratton1941}.
Subsection~\ref{sec:ch332} therefore delivers the scalar angular and radial eigenvalue blocks that that matrix
would diagonalize \cite{hansen1988}.

\paragraph{Rotation covariance of separated labels.}
Finite rotations mix \(m\) at fixed \(\ell\) through Wigner coefficients Eq.~\eqref{eq:ch3.Wigner-D-coefficient-law}
without changing Eq.~\eqref{eq:ch3.radial-helmholtz-canonical} \cite{hansen1988,jackson1999}. Laboratory rotations
of the measurement frame therefore act on coefficients, not on the radial ODE itself \cite{stratton1941,harrington2001}.

\paragraph{Quadrupole angular block (\(\ell=2\)).}
The \(\ell=2\) angular equation supports \(m\in\{-2,-1,0,+1,+2\}\) with fivefold azimuthal degeneracy
\cite{jackson1999,hansen1988}. Compact offset feeds often excite \(|m|\le1\) sectors strongly while
\(|m|=2\) sectors require currents with explicit \(\phi\)-dependent phase \cite{harrington2001,stratton1941}.
Subsection~\ref{sec:ch332} records the quantum numbers before explicit \(Y_{2m}\) display in Subsection~\ref{sec:ch341}
\cite{hansen1988}.

\paragraph{Problem (identify \(\ell\) from angular nodal structure).}
\textbf{Problem.} A far-zone pattern on \(\Gamma_R\) shows \(\ell\) nodal cones in \(\theta\) at fixed \(\phi\).
Relate the cone count to Eq.~\eqref{eq:ch3.angular-laplacian-eigenvalue} \cite{stratton1941,hansen1988}.

\textbf{Step 1.} Nodal cones in \(\theta\) arise from zeros of the angular factor \(Y_{\ell m}\) at fixed \(m\)
\cite{stratton1941}.

\textbf{Step 2.} For each \(\ell\), the number of \(\theta\)-nodal surfaces is \(\ell-|m|\) on standard
two-dimensional diagrams of \(|Y_{\ell m}|\) \cite{hansen1988,jackson1999}.

\textbf{Step 3.} Identify \(\ell\) only after fixing \(m\) through Eq.~\eqref{eq:ch3.azimuthal-eigenvalue-equation};
rotating the pattern mixes \(m\) via Eq.~\eqref{eq:ch3.Wigner-D-coefficient-law} without changing \(\ell\)
\cite{hansen1988}.

\textbf{Physical meaning.} Nodal counting is a diagnostic for \(\ell\), not a substitute for outgoing radial
selection in Subsection~\ref{sec:ch333} \cite{harrington2001,stratton1941}.

\begin{figure}[t]
\centering
\ChOneFigBlock{%
\begin{tikzpicture}[ch1fig, node distance=6mm and 8mm]
  \node[ch1box] (l0) {$\ell=0$\\1 mode};
  \node[ch1box, right=9mm of l0] (l1) {$\ell=1$\\3 modes};
  \node[ch1box, right=9mm of l1] (l2) {$\ell=2$\\5 modes};
  \node[ch1box, below=9mm of l1] (degen) {$2\ell+1$ fold\\Eq.~\eqref{eq:ch3.angular-laplacian-eigenvalue}};
  \draw[ch1flow] (l0) -- (l1) -- (l2);
  \draw[ch1flow] (degen) -- (l1);
\end{tikzpicture}%
}
\caption{Angular degeneracy ladder in Subsection~\ref{sec:ch332}: each \(\ell\) block carries \(2\ell+1\) azimuthal
partners before polarization labels from Section~\ref{sec:ch32} are attached.}
\label{fig:ch3.angular-degeneracy-ladder}
\end{figure}

\paragraph{Inner product on angular factors (forward).}
With weight \(\sin\theta\), the angular inner product preview is
\begin{equation}
\label{eq:ch3.angular-inner-product-preview}
\langle Y_{\ell m},Y_{\ell' m'}\rangle_{\Omega}
=
\int_{0}^{2\pi}\!\int_{0}^{\pi}
Y_{\ell m}^{*}Y_{\ell' m'}\sin\theta\,d\theta\,d\phi
\propto
\delta_{\ell\ell'}\delta_{mm'},
\end{equation}
once Condon--Shortley normalization is fixed in Subsection~\ref{sec:ch341} \cite{hansen1988,stratton1941}. Equation
\eqref{eq:ch3.angular-inner-product-preview} explains why separated modes with distinct \((\ell,m)\) are
orthogonal before radial integration is performed \cite{hansen1988}. Subsection~\ref{sec:ch343} elevates
Eq.~\eqref{eq:ch3.angular-inner-product-preview} to a radiative resolution of identity on
\(\mathcal{F}_{\mathrm{rad}}\) \cite{stratton1941}.

\paragraph{Helmholtz symbol on the angular block.}
The operator \(\mathbf{J}^{2}\) imported from Section~\ref{sec:ch32} acts on angular eigenfunctions with eigenvalue
\(\hbar^{2}\ell(\ell+1)\) in quantum units; in the classical phasor normalization of this monograph the separation
constant is \(\ell(\ell+1)\) as in Eq.~\eqref{eq:ch3.angular-laplacian-eigenvalue} \cite{jackson1999,hansen1988}.
Subsection~\ref{sec:ch332} therefore aligns representation-theoretic degeneracy from Subsection~\ref{sec:ch323}
with the centrifugal term in Eq.~\eqref{eq:ch3.radial-helmholtz-canonical} \cite{stratton1941}. Changing
measurement frame rotates \(m\) labels through Eq.~\eqref{eq:ch3.Wigner-D-coefficient-law} but leaves the
\(\ell(\ell+1)\) barrier intact \cite{hansen1988,jackson1999}.

\paragraph{Worked illustration (cylindrical versus spherical labels).}
A line source array aligned with \(\hat{\mathbf{z}}\) often reports azimuthal mode index \(m\) in cylindrical
language \cite{harrington2001}. Under separation, that index is the eigenvalue in
Eq.~\eqref{eq:ch3.azimuthal-eigenvalue-equation}, not a paraxial ``topological charge'' without gauge screening
\cite{belinfante1940,stratton1941}. Subsection~\ref{sec:ch382} will provide explicit basis change matrices; Subsection~\ref{sec:ch332}
fixes the spherical eigenvalue objects those matrices connect \cite{hansen1988}.

\paragraph{Problem (separate a dipole angular factor \(\ell=1\)).}
\textbf{Problem.} Show that a far-zone dipole pattern carries \(\ell=1\) angular content by reducing
Eq.~\eqref{eq:ch3.angular-laplacian-eigenvalue} for \(\ell=1\) \cite{jackson1999,stratton1941}.

\textbf{Step 1.} Set \(\ell=1\) in Eq.~\eqref{eq:ch3.angular-laplacian-eigenvalue} to obtain the associated
Legendre equation for the \(\theta\)-factor at each \(m\in\{-1,0,+1\}\) \cite{jackson1999}.

\textbf{Step 2.} Impose Eq.~\eqref{eq:ch3.azimuthal-eigenvalue-equation} so the \(\phi\)-dependence is
\(e^{\jj m\phi}\) \cite{stratton1941,hansen1988}.

\textbf{Step 3.} Match a short \(\hat{\mathbf{z}}\)-directed dipole far field to primarily \(m=0\) content; a
horizontal loop excites \(m=\pm1\) admixture \cite{jackson1999,harrington2001}.

\textbf{Step 4.} Pair the angular factor with \(h_{1}^{(1)}(k_{0}r)\) from Subsection~\ref{sec:ch333} before naming
\(\mathbf{M}_{1m}\) or \(\mathbf{N}_{1m}\) in Subsection~\ref{sec:ch341} \cite{hansen1988,tai1994}.

\textbf{Physical meaning.} The dipole is the first nontrivial angular block after the \(\ell=0\) monopole carrier;
misidentifying a quadrupole fit as dipole is an \(\ell\) labeling error, not a polarization error
\cite{stratton1941,harrington2001}.

\paragraph{Energy on separated angular modes (forward).}
The shell energy density integrated over angles for a single \((\ell,m)\) factor scales with
\(|Y_{\ell m}|^{2}\) and the radial energy carried by \(h_{\ell}^{(1)}\) \cite{stratton1941,hansen1988}. Subsection~\ref{sec:ch332}
does not normalize that integral; Section~\ref{sec:ch35} fixes the measure so coefficients map to watts in
Eq.~\eqref{eq:ch3.helmholtz-poynting-flux} \cite{stratton1941,poynting1884}. Separation here supplies the
orthogonal labels that make that normalization block-diagonal in \((\ell,m)\) \cite{hansen1988}.

\begin{table}[t]
\centering
\caption{Angular quantum-number diagnostics used with Subsection~\ref{sec:ch332} separation outputs.}
\label{tab:ch3.332-quantum-diagnostics}
\small
\begin{tabular}{@{}p{0.18\linewidth}p{0.36\linewidth}p{0.36\linewidth}@{}}
\toprule
Label & Mathematical object & Measurement proxy \\
\midrule
\(\ell\) &
Eq.~\eqref{eq:ch3.angular-laplacian-eigenvalue} &
\(\theta\)-nodal cone count \\
\addlinespace
\(m\) &
Eq.~\eqref{eq:ch3.azimuthal-eigenvalue-equation} &
\(\phi\) phase ramp on \(\Gamma_R\) \\
\addlinespace
\(R_\ell\) &
Eq.~\eqref{eq:ch3.radial-helmholtz-canonical} &
Far-zone \(e^{\jj k_0r}/r\) phase \\
\bottomrule
\end{tabular}
\end{table}

Subsection~\ref{sec:ch332} has separated Eq.~\eqref{eq:ch3.vector-helmholtz-E} into angular
Eqs.~\eqref{eq:ch3.angular-laplacian-eigenvalue}--\eqref{eq:ch3.azimuthal-eigenvalue-equation} and radial
Eq.~\eqref{eq:ch3.radial-helmholtz-ODE} \cite{stratton1941,hansen1988}. Subsection~\ref{sec:ch333} selects
outgoing radial factors compatible with Section~\ref{sec:ch24} \cite{jackson1999}.


\subsection{Radial functions and outgoing-wave admissibility}
\label{sec:ch333}
% TOC tag: [USE SS2.4]

Radial factors carry the propagation phase that distinguishes export from storage. Subsection~\ref{sec:ch333}
selects outgoing solutions of Eq.~\eqref{eq:ch3.radial-helmholtz-ODE} by importing Sommerfeld closure from
Section~\ref{sec:ch24} \cite{stratton1941,jackson1999,harrington2001}; no radiation-uniqueness proof is
reopened here.

\paragraph{Standing imports.}
Sommerfeld field condition Eq.~\eqref{eq:ch2.sommerfeld-field-condition}, outgoing selector
Eq.~\eqref{eq:ch2.outgoing-selector-def}, and radiation refinement
Eq.~\eqref{eq:ch2.outgoing-radiation-refinement} enter from Subsections~\ref{sec:ch241}--\ref{sec:ch242}
\cite{stratton1941,harrington2001}. Outgoing inheritance Eq.~\eqref{eq:ch3.outgoing-inheritance-criterion}
and Hilbert closure Eq.~\eqref{eq:ch3.F-rad-hilbert-closure} restrict radial choices to members of
\(\mathcal{F}_{\mathrm{rad}}\) \cite{hansen1988}. The radial Hilbert structure of
Eq.~\eqref{eq:ch3.radial-helmholtz-ODE} pairs with angular quantum numbers from
Eq.~\eqref{eq:ch3.eigenfield-quantum-number-split} \cite{jackson1999}.

\paragraph{Spherical Bessel and Hankel bases.}
Solutions of Eq.~\eqref{eq:ch3.radial-helmholtz-ODE} are linear combinations of spherical Bessel functions
\(j_{\ell}(k_{0}r)\) and spherical Hankel functions \(h_{\ell}^{(1)}(k_{0}r)\), \(h_{\ell}^{(2)}(k_{0}r)\)
\cite{stratton1941,jackson1999}. The regular standing combination at finite \(r\) is
\begin{equation}
\label{eq:ch3.spherical-bessel-regular}
j_{\ell}(k_{0}r)
=
\sqrt{\frac{\pi}{2k_{0}r}}\,J_{\ell+1/2}(k_{0}r),
\end{equation}
while outgoing and incoming radiative combinations are
\begin{equation}
\label{eq:ch3.spherical-hankel-out-in}
h_{\ell}^{(1)}(k_{0}r)=j_{\ell}(k_{0}r)+\jj n_{\ell}(k_{0}r),
\qquad
h_{\ell}^{(2)}(k_{0}r)=j_{\ell}(k_{0}r)-\jj n_{\ell}(k_{0}r),
\end{equation}
with \(n_{\ell}\) the spherical Neumann function \cite{stratton1941,jackson1999}. Equations
\eqref{eq:ch3.spherical-bessel-regular}--\eqref{eq:ch3.spherical-hankel-out-in} are the radial building
blocks for separated modes; their asymptotics fix the propagating phase carried to the far zone
\cite{hansen1988}.

\paragraph{Outgoing radial admissibility (USE).}
On \(\mathcal{F}_{\mathrm{rad}}\), the radial factor in Eq.~\eqref{eq:ch3.scalar-separation-ansatz} must be
outgoing:
\begin{equation}
\label{eq:ch3.outgoing-radial-admissibility}
R_{\ell}^{\mathrm{(out)}}(r)
=
h_{\ell}^{(1)}(k_{0}r),
\qquad
\mathbf{u}_{\omega}\in\mathcal{F}_{\mathrm{rad}}
\ \Longrightarrow\
R_{\ell}\in\mathrm{span}\{h_{\ell}^{(1)}(k_{0}r)\}
\ \text{for each propagating }(\ell,m),
\end{equation}
consistent with Eq.~\eqref{eq:ch2.sommerfeld-field-condition} and selector
Eq.~\eqref{eq:ch2.outgoing-selector-def} \cite{stratton1941,harrington2001}. The standing-wave combination
\(j_{\ell}(k_{0}r)\) is admissible only on bounded or cavity problems with a different boundary class; on
exterior radiation problems it violates outgoing closure and is annihilated by \(\mathcal{O}_{\mathrm{out}}\)
\cite{jackson1999,harrington2001}. Equation~\eqref{eq:ch3.outgoing-radial-admissibility} is the radial restatement
of Eq.~\eqref{eq:ch3.outgoing-inheritance-criterion} before vector harmonics are named
\cite{stratton1941}.

\paragraph{Wronskian and linear independence of radial pairs.}
For fixed \(\ell\), the spherical Bessel pair \(j_{\ell}(k_{0}r)\), \(n_{\ell}(k_{0}r)\) has Wronskian
\begin{equation}
\label{eq:ch3.spherical-bessel-wronskian}
j_{\ell}(k_{0}r)\,\frac{d n_{\ell}}{dr}
-n_{\ell}(k_{0}r)\,\frac{d j_{\ell}}{dr}
=\frac{2}{\pi k_{0}r^{2}},
\end{equation}
so \(h_{\ell}^{(1)}\) and \(h_{\ell}^{(2)}\) are linearly independent outgoing/incoming solutions of
Eq.~\eqref{eq:ch3.radial-helmholtz-canonical} \cite{stratton1941,jackson1999}. Outgoing selection
Eq.~\eqref{eq:ch3.outgoing-radial-admissibility} is therefore a one-dimensional subspace choice within a
two-dimensional solution space, not a normalization convention \cite{harrington2001}.

\paragraph{Far-zone asymptotics and flux certificate.}
As \(r\to\infty\) with fixed \(\theta,\phi\),
\begin{equation}
\label{eq:ch3.hankel-far-zone-asymptotic}
h_{\ell}^{(1)}(k_{0}r)
\sim
(-\jj)^{\ell+1}\,\frac{e^{\jj k_{0}r}}{k_{0}r},
\qquad
k_{0}r\gg\ell,
\end{equation}
producing the \(e^{\jj k_{0}r}/r\) phase front demanded by Eq.~\eqref{eq:ch2.sommerfeld-field-condition}
\cite{stratton1941,jackson1999}. Physically, Eq.~\eqref{eq:ch3.hankel-far-zone-asymptotic} certifies that each
separated mode transports energy outward through \(\Gamma_{R}\) rather than accumulating reactive standing
energy in the exterior class \cite{harrington2001}. Time-harmonic convention \(\exp(-\jj\omega t)\) fixes the
sign in Eq.~\eqref{eq:ch3.hankel-far-zone-asymptotic}; flipping to \(\exp(+\jj\omega t)\) exchanges
\(h_{\ell}^{(1)}\leftrightarrow h_{\ell}^{(2)}\) \cite{jackson1999}.

\begin{figure}[t]
\centering
\ChOneFigBlock{%
\begin{tikzpicture}[ch1fig, node distance=7mm and 9mm]
  \node[ch1box] (ode) {Eq.~\eqref{eq:ch3.radial-helmholtz-ODE}};
  \node[ch1box, right=12mm of ode] (basis) {Eqs.~\eqref{eq:ch3.spherical-bessel-regular}--\eqref{eq:ch3.spherical-hankel-out-in}};
  \node[ch1box, right=12mm of basis] (out) {Eq.~\eqref{eq:ch3.outgoing-radial-admissibility}};
  \node[ch1box, below=10mm of out] (sf) {Eq.~\eqref{eq:ch2.sommerfeld-field-condition}};
  \draw[ch1flow] (ode) -- (basis) -- (out);
  \draw[ch1flow] (sf) -- (out);
\end{tikzpicture}%
}
\caption{Radial outgoing gate in Subsection~\ref{sec:ch333}: solutions of
Eq.~\eqref{eq:ch3.radial-helmholtz-ODE} reduce to Hankel outgoing factors
Eq.~\eqref{eq:ch3.outgoing-radial-admissibility} enforced by Sommerfeld asymptotics
Eq.~\eqref{eq:ch2.sommerfeld-field-condition}.}
\label{fig:ch3.radial-outgoing-gate}
\end{figure}

\begin{proposition}[Outgoing radial selection commutes with radiation map]
\label{prop:ch3.radial-outgoing-commutes}
Let \(\mathbf{u}_{\ell m}\) be a separated mode with angular factor \(Y_{\ell m}\) and radial factor
\(R_{\ell}(r)\). If \(R_{\ell}=h_{\ell}^{(1)}(k_{0}r)\), then
\(\Pi_{\mathrm{rad}}[\mathbf{u}_{\ell m}]=\mathbf{u}_{\ell m}\) on exterior domains with
Eq.~\eqref{eq:ch2.outgoing-radiation-refinement} \cite{stratton1941,harrington2001}. If
\(R_{\ell}=j_{\ell}(k_{0}r)\) on an exterior domain, \(\mathcal{O}_{\mathrm{out}}[\mathbf{u}_{\ell m}]=\mathbf{0}\)
unless boundary sources reintroduce standing content \cite{jackson1999}.
\end{proposition}
\begin{proof}[Proof sketch]
Asymptotic match to Eq.~\eqref{eq:ch2.sommerfeld-field-condition} and idempotence of
Eq.~\eqref{eq:ch2.outgoing-selector-def} \cite{stratton1941,harrington2001}.
% PROOF: AUTHOR REVIEW
\end{proof}

\paragraph{Near-zone regularity versus far-zone branch.}
Near sources, \(h_{\ell}^{(1)}(k_{0}r)\) and \(j_{\ell}(k_{0}r)\) are both finite on shells excluding the
origin; branch selection is therefore invisible in near-field scans \cite{harrington2001,stratton1941}. Export
claims require the asymptotic certificate Eq.~\eqref{eq:ch3.hankel-far-zone-asymptotic}, not near-zone
magnitude alone \cite{harrington2001}. This distinction mirrors the operator lesson of
Eq.~\eqref{eq:ch3.outgoing-inheritance-criterion}: membership in \(\mathcal{F}_{\mathrm{rad}}\) is a far-zone
closure property \cite{stratton1941}.

\paragraph{Evanescent and cut-off regimes.}
When \(k_{0}r\ll\ell\), centrifugal dominance in Eq.~\eqref{eq:ch3.radial-helmholtz-ODE} forces evanescent
radial decay \(\sim r^{\ell}\) or stronger, regardless of Hankel label \cite{jackson1999}. Such content lies in
\(\ker\Pi_{\mathrm{rad}}\) through Eq.~\eqref{eq:ch2.outgoing-radiation-refinement} and does not contribute to
far-zone flux budgets \cite{harrington2001,stratton1941}. Subsection~\ref{sec:ch333} therefore pairs outgoing
Hankel selection with the regime screens already used in Chapter~\ref{ch:02} \cite{harrington2001}.

\paragraph{Sommerfeld impedance ratio at fixed \(\ell\).}
On outgoing Hankel factors, the radial impedance ratio approaches
\begin{equation}
\label{eq:ch3.outgoing-radial-impedance-ratio}
\lim_{r\to\infty}
\frac{\jj k_{0}R_{\ell}'(r)}{R_{\ell}(r)}
=\jj k_{0},
\qquad
R_{\ell}=h_{\ell}^{(1)}(k_{0}r),
\end{equation}
matching Eq.~\eqref{eq:ch2.sommerfeld-impedance-ratio} for transverse spherical waves \cite{stratton1941,jackson1999}.
Equation~\eqref{eq:ch3.outgoing-radial-impedance-ratio} is the radial restatement of the field Sommerfeld
condition Eq.~\eqref{eq:ch2.sommerfeld-field-condition} at fixed angular momentum block \(\ell\)
\cite{harrington2001}.

\paragraph{Green-function radial factor (USE).}
The scalar outgoing Green function selected in Chapter~\ref{ch:02} carries the same Hankel radial dependence as
Eq.~\eqref{eq:ch3.outgoing-radial-admissibility} \cite{stratton1941,jackson1999}. Subsection~\ref{sec:ch333}
does not reconstruct the dyadic Green tensor; it records that separated modes inherit the radial branch already
fixed by Eq.~\eqref{eq:ch2.outgoing-green-branch} \cite{harrington2001}.

\begin{theorem}[Radial outgoing selection and uniqueness class]
\label{thm:ch3.radial-outgoing-uniqueness-class}
On exterior domains with Eq.~\eqref{eq:ch2.sommerfeld-outgoing-class}, each separated mode with angular label
\((\ell,m)\) and propagating radial wavenumber \(k_{0}\) is unique up to scalar multiples once
Eq.~\eqref{eq:ch3.outgoing-radial-admissibility} is imposed \cite{stratton1941,harrington2001}. Uniqueness of
the full Maxwell pair follows from Theorem~\ref{thm:ch2.exterior-radiative-uniqueness} cited without reproof
\cite{stratton1941}.
\end{theorem}
\begin{proof}[Proof sketch]
Two independent outgoing radial solutions would violate the second-order ODE
Eq.~\eqref{eq:ch3.radial-helmholtz-canonical} and the Sommerfeld asymptotic class
\cite{stratton1941,jackson1999}.
% PROOF: AUTHOR REVIEW
\end{proof}

\paragraph{Worked illustration (incoming contamination).}
Replacing \(h_{\ell}^{(1)}\) by \(h_{\ell}^{(2)}\) flips the far-zone phase to an incoming front while leaving
near-zone magnitudes nearly unchanged on moderate shells \cite{jackson1999,stratton1941}. A post-processing fit
that matches only \(|E|\) on \(\Gamma_{R}\) can therefore lock onto the wrong branch unless
Eq.~\eqref{eq:ch3.outgoing-radial-admissibility} is enforced through \(\mathcal{O}_{\mathrm{out}}\)
\cite{harrington2001}.

\paragraph{Worked illustration (low-\(\ell\) dominance).}
For \(\ell=0\), \(h_{0}^{(1)}(k_{0}r)=-(\jj/k_{0}r)e^{\jj k_{0}r}\), the isotropic outgoing carrier underlying
the \(e^{\jj k_{0}r}/r\) pattern of compact sources \cite{jackson1999}. Higher \(\ell\) introduces additional
\((k_{0}r)^{-\ell}\) prefactors in Eq.~\eqref{eq:ch3.hankel-far-zone-asymptotic}, explaining faster angular
decay of high-order multipoles \cite{stratton1941,hansen1988}.

\paragraph{Worked illustration (absorbing boundary mimic).}
Numerical codes often replace \(h_{\ell}^{(1)}\) by \(j_{\ell}+\jj y_{\ell}\) on a truncated shell with loss;
unless the shell is in the asymptotic zone where Eq.~\eqref{eq:ch3.hankel-far-zone-asymptotic} holds, the mimic
can introduce spurious standing admixture \cite{harrington2001}. Export claims must still pass
Eq.~\eqref{eq:ch3.outgoing-inheritance-criterion} on \(\mathcal{F}_{\mathrm{rad}}\), not merely match a local
impedance boundary \cite{stratton1941}.

\paragraph{Worked illustration (reactive near-field storage).}
A mode can exhibit large \(|h_{\ell}^{(1)}(k_{0}r)|\) on a moderate shell while still belonging to
\(\mathcal{F}_{\mathrm{rad}}\) because outgoing classification is asymptotic
\cite{harrington2001,stratton1941}. Power budgets therefore require
Eq.~\eqref{eq:ch2.outgoing-radiation-refinement} before interpreting shell energy as transported flux
\cite{harrington2001}.

\paragraph{Problem (outgoing Hankel selection at fixed \(\ell\)).}
\textbf{Problem.} Solve Eq.~\eqref{eq:ch3.radial-helmholtz-canonical} for \(\ell=1\) on an exterior domain and
select the physical branch using Eq.~\eqref{eq:ch2.sommerfeld-field-condition} \cite{stratton1941,jackson1999}.

\textbf{Step 1 (general solution).} For \(\ell=1\),
\(R_{1}(r)=A\,j_{1}(k_{0}r)+B\,n_{1}(k_{0}r)\) or equivalently
\(R_{1}(r)=C\,h_{1}^{(1)}(k_{0}r)+D\,h_{1}^{(2)}(k_{0}r)\) \cite{jackson1999,stratton1941}.

\textbf{Step 2 (asymptotic test).} Use Eq.~\eqref{eq:ch3.hankel-far-zone-asymptotic} with \(\ell=1\):
\(h_{1}^{(1)}(k_{0}r)\sim +\jj e^{\jj k_{0}r}/(k_{0}r)^{2}\) up to angular factors \cite{jackson1999}. Only
\(h_{1}^{(1)}\) matches the outgoing phase in Eq.~\eqref{eq:ch2.sommerfeld-field-condition}
\cite{stratton1941,harrington2001}.

\textbf{Step 3 (selector action).} Impose Eq.~\eqref{eq:ch3.outgoing-radial-admissibility}:
\(R_{1}^{\mathrm{(out)}}=h_{1}^{(1)}(k_{0}r)\), annihilating incoming admixture by
Eq.~\eqref{eq:ch2.outgoing-selector-def} \cite{harrington2001}.

\textbf{Physical meaning.} Step~2 is the laboratory-irrelevant but mathematically decisive test: near-zone scans
cannot distinguish Steps 2--3 without a far-zone phase reference \cite{stratton1941,harrington2001}.

\paragraph{Recurrence and numerical evaluation.}
Spherical Hankel functions satisfy
\begin{equation}
\label{eq:ch3.hankel-recurrence-relation}
h_{\ell+1}^{(1)}(z)
=
\frac{2\ell+1}{z}\,h_{\ell}^{(1)}(z)
-h_{\ell-1}^{(1)}(z),
\qquad
\ell\ge1,
\end{equation}
with seed values \(h_{0}^{(1)}(z)=-\jj e^{\jj z}/z\) and
\(h_{1}^{(1)}(z)=-(1+\jj/z)e^{\jj z}/z\) \cite{stratton1941,jackson1999}. Equation
\eqref{eq:ch3.hankel-recurrence-relation} is the stable upward recurrence used in Subsection~\ref{sec:ch334}
expansions and in Section~\ref{sec:ch35} normalization integrals \cite{hansen1988}. Downward recurrence is
ill-conditioned near \(z\sim\ell\), matching the evanescent regime discussion above \cite{harrington2001}.

\paragraph{Connection to eigenfield quantum numbers.}
Outgoing radial selection at fixed \(\ell\) pairs with azimuthal index \(m\) from
Eq.~\eqref{eq:ch3.eigenfield-quantum-number-split} without changing \(R_{\ell}\) \cite{jackson1999,hansen1988}.
Subsection~\ref{sec:ch333} therefore supplies the radial half of the separated eigenfield label; Subsection~\ref{sec:ch334}
assembles the vector field \cite{stratton1941,tai1994}.

\paragraph{Small-argument and large-argument regimes.}
For \(z=k_{0}r\ll1\), \(h_{\ell}^{(1)}(z)\sim -(\ell-1)!\,(2/z)^{\ell+1}\) up to leading order, while for
\(z\gg\ell\), Eq.~\eqref{eq:ch3.hankel-far-zone-asymptotic} applies \cite{jackson1999,stratton1941}. The two
regimes explain why high-\(\ell\) modes are suppressed near compact sources yet propagate in the far zone only
when excited by accessible currents \cite{harrington2001,hansen1988}. Regime screens from Chapter~\ref{ch:02}
remain mandatory when translating these asymptotics into power budgets \cite{harrington2001}.

\paragraph{Standing-wave contamination diagnostic.}
Define the outgoing purity ratio
\begin{equation}
\label{eq:ch3.outgoing-purity-ratio}
\mathcal{P}_{\mathrm{out}}[\mathbf{u}_\omega]
=
\frac{\|\mathcal{O}_{\mathrm{out}}[\mathbf{u}_\omega]\|_{\mathrm{rad}}^{2}}
{\|\mathbf{u}_\omega\|_{\mathrm{rad}}^{2}},
\end{equation}
with \(\mathcal{O}_{\mathrm{out}}\) from Eq.~\eqref{eq:ch2.outgoing-selector-def} \cite{harrington2001}. On
\(\mathcal{F}_{\mathrm{rad}}\), \(\mathcal{P}_{\mathrm{out}}=1\); standing admixture drives
\(\mathcal{P}_{\mathrm{out}}<1\) even when local Helmholtz fits appear excellent \cite{stratton1941,harrington2001}.
Subsection~\ref{sec:ch333} interprets radial branch errors as failures of \(\mathcal{P}_{\mathrm{out}}\), not as
normalization choices \cite{harrington2001}.

\begin{table}[t]
\centering
\caption{Radial branch dictionary for exterior radiation (Subsection~\ref{sec:ch333}).}
\label{tab:ch3.333-radial-branches}
\small
\begin{tabular}{@{}p{0.22\linewidth}p{0.30\linewidth}p{0.38\linewidth}@{}}
\toprule
Radial factor & Far-zone phase & \(\mathcal{F}_{\mathrm{rad}}\) status \\
\midrule
\(h_{\ell}^{(1)}(k_{0}r)\) &
Eq.~\eqref{eq:ch3.hankel-far-zone-asymptotic} &
Admissible outgoing \\
\addlinespace
\(h_{\ell}^{(2)}(k_{0}r)\) &
Incoming \(e^{-\jj k_{0}r}/r\) &
Excluded by \(\mathcal{O}_{\mathrm{out}}\) \\
\addlinespace
\(j_{\ell}(k_{0}r)\) &
Standing \(\cos/\sin\) &
Cavity/bounded class only \\
\bottomrule
\end{tabular}
\end{table}

\paragraph{Spherical Neumann and irregular admixture.}
The Neumann function \(n_{\ell}(k_{0}r)\) grows as \((k_{0}r)^{-\ell-1}\) near the origin and must be excluded on
exterior physical branches even when fits on truncated domains suggest its presence \cite{jackson1999,stratton1941}.
Outgoing purity Eq.~\eqref{eq:ch3.outgoing-purity-ratio} detects such admixture through failure of
Eq.~\eqref{eq:ch2.outgoing-selector-def} rather than through local residual minimization \cite{harrington2001}.

\paragraph{Packaging for Section~\ref{sec:ch34}.}
The separated outgoing radial factor \(h_{\ell}^{(1)}(k_{0}r)\) combined with angular eigenfunctions from
Eq.~\eqref{eq:ch3.angular-laplacian-eigenvalue} is the scalar backbone every vector harmonic in
Section~\ref{sec:ch34} multiplies by \(\curl\) or \(\curl\curl\) \cite{hansen1988,stratton1941}. Subsection~\ref{sec:ch333}
therefore ends with a radial dictionary, not with named \(\mathbf{M}_{\ell m}\), \(\mathbf{N}_{\ell m}\) objects
\cite{tai1994}.

\paragraph{Radial flux-weighted orthogonality (forward).}
For fixed \(\ell\) and outgoing factors, the radial pairing preview is
\begin{equation}
\label{eq:ch3.radial-orthogonality-preview}
\int_{R_{\min}}^{R_{\max}}
h_{\ell}^{(1)}(k_{0}r)^{*}\,h_{\ell'}^{(1)}(k_{0}r)\,r^{2}\,dr
\ \propto\
\delta_{\ell\ell'},
\end{equation}
on shells where asymptotic orthogonality of spherical Hankel functions holds \cite{stratton1941,hansen1988}. Equation
\eqref{eq:ch3.radial-orthogonality-preview} is the radial partner of
Eq.~\eqref{eq:ch3.angular-inner-product-preview} and is made exact in the normalization program of
Section~\ref{sec:ch35} \cite{hansen1988,stratton1941}. Subsection~\ref{sec:ch333} cites
Eq.~\eqref{eq:ch3.radial-orthogonality-preview} as motivation for separate radial branch selection before
normalization integrals are evaluated \cite{hansen1988}.

\paragraph{Problem (compare outgoing versus standing radial flux).}
\textbf{Problem.} Compare the large-\(r\) Poynting flux of \(R_{\ell}=h_{\ell}^{(1)}(k_{0}r)\) versus
\(R_{\ell}=j_{\ell}(k_{0}r)\) when paired with the same \(Y_{\ell m}\) \cite{stratton1941,jackson1999}.

\textbf{Step 1.} Insert \(h_{\ell}^{(1)}\) into the separated scalar ansatz Eq.~\eqref{eq:ch3.scalar-separation-ansatz}
and form \(\E,\H\) through Debye maps Eqs.~\eqref{eq:ch3.debye-E-representation}--\eqref{eq:ch3.debye-H-representation}
\cite{tai1994}.

\textbf{Step 2.} Use Eq.~\eqref{eq:ch3.hankel-far-zone-asymptotic} to extract the \(e^{\jj k_{0}r}/r\) phase of the
outgoing pair \cite{jackson1999,stratton1941}.

\textbf{Step 3.} Repeat with \(j_{\ell}\): asymptotics are standing \(\cos(k_{0}r-\delta)/r\), producing zero
net time-averaged outward flux on \(\Gamma_R\) in the exterior class \cite{stratton1941,harrington2001}.

\textbf{Physical meaning.} Flux sign distinguishes outgoing from standing radial factors even when near-zone
magnitudes are comparable \cite{harrington2001}. This is the radial content of
Eq.~\eqref{eq:ch3.outgoing-purity-ratio} \cite{stratton1941}.

\begin{figure}[t]
\centering
\ChOneFigBlock{%
\begin{tikzpicture}[ch1fig, node distance=7mm and 9mm]
  \node[ch1box] (near) {Near zone\\$k_0r\sim\ell$};
  \node[ch1box, right=12mm of near] (mid) {Intermediate\\shell fit};
  \node[ch1box, right=12mm of mid] (far) {Far zone\\Eq.~\eqref{eq:ch3.hankel-far-zone-asymptotic}};
  \node[ch1box, below=10mm of mid] (gate) {Eq.~\eqref{eq:ch3.outgoing-radial-admissibility}};
  \draw[ch1flow] (near) -- (mid) -- (far);
  \draw[ch1flow] (gate) -- (far);
\end{tikzpicture}%
}
\caption{Branch selection for Subsection~\ref{sec:ch333}: near-zone magnitudes are ambiguous; outgoing classification
is fixed only in the far-zone asymptotic regime of Eq.~\eqref{eq:ch3.hankel-far-zone-asymptotic}.}
\label{fig:ch3.near-far-radial-gate}
\end{figure}

\paragraph{Worked illustration (\(\ell=2\) outgoing radial factor).}
For \(\ell=2\), \(h_{2}^{(1)}(k_{0}r)\) carries quadrupole radial decay prefactor \((k_{0}r)^{-3}\) in the far zone
\cite{jackson1999,stratton1941}. Compact sources therefore excite \(\ell=2\) content with smaller far-zone amplitude
than \(\ell=0,1\) at equal coefficient norm, explaining rapid decay of high multipoles in measured patterns
\cite{harrington2001,hansen1988}. Accessibility screens from Chapter~\ref{ch:02} further suppress unreachable
\(\ell\) sectors \cite{stratton1941,harrington2001}.

\paragraph{Radiation resistance analogy (diagnostic only).}
The ratio \(|h_{\ell}^{(1)}(k_{0}r)/h_{0}^{(1)}(k_{0}r)|\) at fixed \(r\) decreases with \(\ell\), mimicking the
frequency dependence of higher-order multipole radiation resistance in antenna textbooks \cite{harrington2001,stratton1941}.
Subsection~\ref{sec:ch333} records the analytic origin of that decrease in
Eq.~\eqref{eq:ch3.hankel-far-zone-asymptotic}; Section~\ref{sec:ch35} fixes normalization so coefficients map to
measurable ohms-equivalent flux \cite{stratton1941,hansen1988}.

\begin{table}[t]
\centering
\caption{Asymptotic orders for outgoing Hankel factors Eq.~\eqref{eq:ch3.hankel-far-zone-asymptotic}.}
\label{tab:ch3.333-hankel-asymptotic-orders}
\small
\begin{tabular}{@{}cccc@{}}
\toprule
\(\ell\) & Far-zone amplitude order & Phase & \(\mathcal{F}_{\mathrm{rad}}\) \\
\midrule
0 & \((k_{0}r)^{-1}\) & \(e^{\jj k_{0}r}\) & yes \\
1 & \((k_{0}r)^{-2}\) & \(e^{\jj k_{0}r}\) & yes \\
2 & \((k_{0}r)^{-3}\) & \(e^{\jj k_{0}r}\) & yes \\
\bottomrule
\end{tabular}
\end{table}

Subsection~\ref{sec:ch333} has fixed outgoing radial factors
Eq.~\eqref{eq:ch3.outgoing-radial-admissibility} on \(\mathcal{F}_{\mathrm{rad}}\) using
Eqs.~\eqref{eq:ch2.sommerfeld-field-condition}--\eqref{eq:ch2.outgoing-radiation-refinement}
\cite{stratton1941,jackson1999}. Subsection~\ref{sec:ch334} introduces Debye potentials that assemble separated
radial and angular factors into divergence-free vector modes \cite{tai1994,stratton1941}.


\subsection{Debye potentials and generating functions}
\label{sec:ch334}
% TOC tag: [HOME]

Coupled spherical components obscure the vector Helmholtz structure. Subsection~\ref{sec:ch334} introduces
scalar Debye potentials \(\Pi\) and \(\Psi\) whose curls and double-curls generate divergence-free fields with
the separation labels of Subsection~\ref{sec:ch332} and outgoing radial factors of Subsection~\ref{sec:ch333}
\cite{stratton1941,tai1994,hansen1988}.

\paragraph{Assumptions.}
Work on source-free exterior domains at fixed \(\omega\), with \(k_{0}=\omega\sqrt{\varepsilon_{0}\mu_{0}}\) and
fields restricted to \(\mathcal{F}_{\mathrm{rad}}\) through Eq.~\eqref{eq:ch3.outgoing-inheritance-criterion}
\cite{stratton1941,jackson1999}. Potentials \(\Pi,\Psi\) are sufficiently smooth for the curls below and
inherit outgoing radial dependence Eq.~\eqref{eq:ch3.outgoing-radial-admissibility} \cite{hansen1988}.

\paragraph{Debye potential ansatz (HOME).}
Introduce scalars \(\Pi(\mathbf{r})\) and \(\Psi(\mathbf{r})\) and define the electric and magnetic fields by
\begin{equation}
\label{eq:ch3.debye-E-representation}
\E
=
\curl\curl(\Pi\,\hat{\mathbf{r}})
-\jj\omega\mu_{0}\,\curl(\Psi\,\hat{\mathbf{r}}),
\end{equation}
\begin{equation}
\label{eq:ch3.debye-H-representation}
\H
=
\jj\omega\varepsilon_{0}\,\curl(\Pi\,\hat{\mathbf{r}})
+\curl\curl(\Psi\,\hat{\mathbf{r}}),
\end{equation}
with \(\hat{\mathbf{r}}=\mathbf{r}/r\) \cite{stratton1941,tai1994}. Equations~\eqref{eq:ch3.debye-E-representation}--\eqref{eq:ch3.debye-H-representation}
are the Subsection~\ref{sec:ch334} HOME generating map: every admissible transverse pair on the separation
class can be represented through \((\Pi,\Psi)\) rather than through raw spherical components
Eq.~\eqref{eq:ch3.vector-field-spherical-components} \cite{hansen1988,tai1994}.

\paragraph{Radial Hertz orientation and why \(\hat{\mathbf{r}}\) appears.}
The factors \(\Pi\,\hat{\mathbf{r}}\) and \(\Psi\,\hat{\mathbf{r}}\) orient scalar potentials along the radial
unit vector so that \(\curl(\Pi\hat{\mathbf{r}})\) is automatically transverse to \(\hat{\mathbf{r}}\) in the
far zone \cite{stratton1941,tai1994}. A Cartesian Hertz vector would require additional gauge constraints to
match Eq.~\eqref{eq:ch3.transverse-radiative-constraint} on \(\mathcal{F}_{\mathrm{rad}}\) \cite{hansen1988}.
Subsection~\ref{sec:ch334} fixes the radial Debye convention used throughout Volume~I \cite{tai1994,stratton1941}.

\paragraph{Scalar Helmholtz constraints.}
Potentials obey the scalar Helmholtz equations
\begin{equation}
\label{eq:ch3.debye-scalar-helmholtz-pair}
\nabla^{2}\Pi+k_{0}^{2}\Pi=0,
\qquad
\nabla^{2}\Psi+k_{0}^{2}\Psi=0,
\end{equation}
equivalent to Eq.~\eqref{eq:ch3.radial-helmholtz-ODE} after separation
Eq.~\eqref{eq:ch3.scalar-separation-ansatz} \cite{stratton1941,jackson1999}. Substituting
Eqs.~\eqref{eq:ch3.debye-E-representation}--\eqref{eq:ch3.debye-scalar-helmholtz-pair} into
Eq.~\eqref{eq:ch3.vector-helmholtz-E} closes the reduction loop begun in Subsection~\ref{sec:ch331}
\cite{tai1994,hansen1988}.

\begin{figure}[t]
\centering
\ChOneFigBlock{%
\begin{tikzpicture}[ch1fig, node distance=7mm and 9mm]
  \node[ch1box] (pot) {$(\Pi,\Psi)$\\Eq.~\eqref{eq:ch3.debye-scalar-helmholtz-pair}};
  \node[ch1box, right=12mm of pot] (gen) {Eqs.~\eqref{eq:ch3.debye-E-representation}--\eqref{eq:ch3.debye-H-representation}};
  \node[ch1box, right=12mm of gen] (vh) {Eq.~\eqref{eq:ch3.vector-helmholtz-E}};
  \node[ch1box, below=10mm of pot] (out) {Eq.~\eqref{eq:ch3.outgoing-radial-admissibility}};
  \draw[ch1flow] (pot) -- (gen) -- (vh);
  \draw[ch1flow] (out) -- (pot);
\end{tikzpicture}%
}
\caption{Debye generating chain in Subsection~\ref{sec:ch334}: outgoing scalar potentials
Eq.~\eqref{eq:ch3.debye-scalar-helmholtz-pair} map to \(\E,\H\) through
Eqs.~\eqref{eq:ch3.debye-E-representation}--\eqref{eq:ch3.debye-H-representation}, reproducing vector
Helmholtz solutions on \(\mathcal{F}_{\mathrm{rad}}\).}
\label{fig:ch3.debye-generating-chain}
\end{figure}

\begin{theorem}[Debye potential representation on \(\mathcal{F}_{\mathrm{rad}}\)]
\label{thm:ch3.debye-representation}
On a source-free exterior domain with Eq.~\eqref{eq:ch3.transverse-radiative-constraint}, every
\(\E,\H\) solving Eqs.~\eqref{eq:ch3.vector-helmholtz-E}--\eqref{eq:ch3.vector-helmholtz-H} admits scalars
\(\Pi,\Psi\) such that Eqs.~\eqref{eq:ch3.debye-E-representation}--\eqref{eq:ch3.debye-scalar-helmholtz-pair}
hold. Conversely, any \(\Pi,\Psi\) satisfying Eq.~\eqref{eq:ch3.debye-scalar-helmholtz-pair} generate a
source-free Maxwell pair through Eqs.~\eqref{eq:ch3.debye-E-representation}--\eqref{eq:ch3.debye-H-representation}
\cite{stratton1941,tai1994,hansen1988}.
\end{theorem}
\begin{proof}[Proof sketch]
Standard Debye--Bromwich construction on simply connected exterior domains; gauge shifts of \(\Pi,\Psi\) leave
\(\E,\H\) invariant up to the transverse class \cite{stratton1941,tai1994}.
% PROOF: AUTHOR REVIEW
\end{proof}

\paragraph{Substitution check (vector Helmholtz closure).}
Insert Eqs.~\eqref{eq:ch3.debye-E-representation}--\eqref{eq:ch3.debye-scalar-helmholtz-pair} into
\(\curl\curl\E-k_{0}^{2}\E\). The \(\Pi\) channel produces terms \(\curl\curl\curl(\Pi\hat{\mathbf{r}})\) and
\(\curl\curl\curl(\Psi\hat{\mathbf{r}})\) that cancel against \(k_{0}^{2}\curl(\cdot)\) when
Eq.~\eqref{eq:ch3.debye-scalar-helmholtz-pair} holds \cite{stratton1941,tai1994}. The surviving structure is
identically zero, establishing that Debye potentials generate solutions of
Eq.~\eqref{eq:ch3.vector-helmholtz-E} without solving the coupled component system of
Eq.~\eqref{eq:ch3.vector-field-spherical-components} directly \cite{hansen1988}. The magnetic equation
Eq.~\eqref{eq:ch3.vector-helmholtz-H} closes by duality or parallel substitution \cite{stratton1941}.

\paragraph{Toroidal and poloidal content without premature naming.}
Expand \(\Pi\) and \(\Psi\) in the separation basis of Subsection~\ref{sec:ch332}:
\begin{equation}
\label{eq:ch3.debye-separated-expansion}
\Pi(r,\theta,\phi)
=
\sum_{\ell m}a_{\ell m}\,h_{\ell}^{(1)}(k_{0}r)\,Y_{\ell m}(\theta,\phi),
\qquad
\Psi
=
\sum_{\ell m}b_{\ell m}\,h_{\ell}^{(1)}(k_{0}r)\,Y_{\ell m}(\theta,\phi),
\end{equation}
with outgoing radial factors fixed by Eq.~\eqref{eq:ch3.outgoing-radial-admissibility}
\cite{stratton1941,hansen1988}. The term \(\curl(\Pi\hat{\mathbf{r}})\) carries toroidal-type circulation;
\(\curl\curl(\Pi\hat{\mathbf{r}})\) carries poloidal-type transverse content \cite{tai1994,hansen1988}.
Subsection~\ref{sec:ch334} records this split as generating-structure input; the named vector harmonics
\(\mathbf{M}_{\ell m}\), \(\mathbf{N}_{\ell m}\) are defined once in Subsection~\ref{sec:ch341} from the same
potentials \cite{stratton1941,hansen1988}.

\paragraph{Electric--magnetic duality on potentials.}
Vacuum duality Eq.~\eqref{eq:ch3.dual-symmetry-operator} acts on Debye pairs as
\begin{equation}
\label{eq:ch3.debye-duality-exchange}
\mathcal{D}_{\pi/2}:
(\Pi,\Psi)\longmapsto(\Psi,-\Pi),
\qquad
(\E,\H)\longmapsto(\jj\eta_{0}\H,-\jj\E/\eta_{0}),
\end{equation}
with \(\eta_{0}=\sqrt{\mu_{0}/\varepsilon_{0}}\), consistent with Eq.~\eqref{eq:ch1.duality-rotation}
\cite{stratton1941,belinfante1940}. Equation~\eqref{eq:ch3.debye-duality-exchange} shows that electric and
magnetic generating roles are interchangeable at the potential level, while separated quantum numbers
\((\ell,m)\) from Eq.~\eqref{eq:ch3.eigenfield-quantum-number-split} remain unchanged under duality
\cite{tai1994,hansen1988}.

\paragraph{Coefficient algebra for separated modes.}
For a single \((\ell,m)\) pair, define the electric-type coefficient vector
\(\mathbf{a}_{\ell m}=(a_{\ell m},b_{\ell m})^{\mathsf{T}}\). Linearity of
Eqs.~\eqref{eq:ch3.debye-E-representation}--\eqref{eq:ch3.debye-H-representation} implies
\begin{equation}
\label{eq:ch3.debye-coefficient-linearity}
\E_{\ell m}[\mathbf{a}_{\ell m}]
=
a_{\ell m}\,\E_{\ell m}^{(\Pi)}
+
b_{\ell m}\,\E_{\ell m}^{(\Psi)},
\end{equation}
with fixed generating fields \(\E_{\ell m}^{(\Pi)}\), \(\E_{\ell m}^{(\Psi)}\) built from
\(h_{\ell}^{(1)}(k_{0}r)Y_{\ell m}\) \cite{stratton1941,hansen1988}. Equation
\eqref{eq:ch3.debye-coefficient-linearity} is the Subsection~\ref{sec:ch334} amplitude law used in
Subsection~\ref{sec:ch344} when attaching polarization index \(\sigma\) \cite{tai1994}. Dual exchange
Eq.~\eqref{eq:ch3.debye-duality-exchange} acts on \(\mathbf{a}_{\ell m}\) by the linear map
\((a,b)^{\mathsf{T}}\mapsto(b,-a)^{\mathsf{T}}\) at \(\xi=\pi/2\) \cite{belinfante1940,stratton1941}.

\paragraph{Gauge freedom and observability.}
Transformations \(\Pi\mapsto\Pi+\partial_{r}\chi\), \(\Psi\mapsto\Psi+\partial_{r}\chi'\) with gauge scalars
\(\chi,\chi'\) tied to radial orientation leave \(\E,\H\) invariant when potentials remain in the scalar
Helmholtz class \cite{stratton1941,tai1994}. Observable amplitudes therefore attach to coefficients
\((a_{\ell m},b_{\ell m})\) in Eq.~\eqref{eq:ch3.debye-separated-expansion} only after the gauge and
outgoing screens of Chapter~\ref{ch:01} pass \cite{belinfante1940,harrington2001}. Near-zone reactive
superpositions of \(\Pi,\Psi\) can be large while \(\mathcal{R}_\omega[\mathfrak{D}]\) projects to the
outgoing combination fixed by Eq.~\eqref{eq:ch3.outgoing-inheritance-criterion} \cite{harrington2001}.

\paragraph{Connection to angular-momentum eigenfields.}
When \(\Pi\) or \(\Psi\) is a pure \((\ell,m)\) separated mode in Eq.~\eqref{eq:ch3.debye-separated-expansion},
the generated \(\E,\H\) are simultaneous eigenvectors of Eq.~\eqref{eq:ch3.J-simultaneous-eigenfield-system}
with the same labels \cite{jackson1999,hansen1988}. Subsection~\ref{sec:ch344} will name these eigenfields
\(\mathbf{u}_{\ell m}^{(\sigma)}\) and attach polarization index \(\sigma\) from Section~\ref{sec:ch32}
\cite{stratton1941,tai1994}. Subsection~\ref{sec:ch334} supplies only the generating operators
Eqs.~\eqref{eq:ch3.debye-E-representation}--\eqref{eq:ch3.debye-separated-expansion}.

\paragraph{Validity domain.}
Eqs.~\eqref{eq:ch3.debye-E-representation}--\eqref{eq:ch3.debye-scalar-helmholtz-pair} assume linear isotropic
vacuum exterior media and simply connected observation domains excluding the origin \cite{stratton1941}. Magnetic
currents, anisotropic tensors, and interior cavity modes require modified potential sets deferred to later
chapters \cite{harrington2001,tai1994}.

\paragraph{Worked illustration (pure \(\Pi\) mode).}
Take \(\Psi=0\) and \(\Pi=h_{\ell}^{(1)}(k_{0}r)Y_{\ell m}\). Then \(\E=\curl\curl(\Pi\hat{\mathbf{r}})\) and
\(\H=\jj\omega\varepsilon_{0}\curl(\Pi\hat{\mathbf{r}})\) form a transverse radiative pair with orbital labels
\((\ell,m)\) \cite{stratton1941,hansen1988}. The far-zone pattern rotates with \(\phi\) through
Eq.~\eqref{eq:ch3.azimuthal-eigenvalue-equation}, not through an independent ``mode order'' parameter
\cite{jackson1999}.

\paragraph{Worked illustration (pure \(\Psi\) mode).}
With \(\Pi=0\) and \(\Psi=h_{\ell}^{(1)}(k_{0}r)Y_{\ell m}\), the roles in
Eqs.~\eqref{eq:ch3.debye-E-representation}--\eqref{eq:ch3.debye-H-representation} exchange electric and
magnetic generating channels, producing the magnetic-dual partner of the pure-\(\Pi\) illustration
\cite{tai1994,stratton1941}. Duality Eq.~\eqref{eq:ch3.debye-duality-exchange} maps one illustration onto the
other without changing \((\ell,m)\) \cite{belinfante1940}.

\paragraph{Worked illustration (superposed \(\Pi,\Psi\) amplitudes).}
Take \(a_{\ell m}=1\), \(b_{\ell m}=\jj\) in Eq.~\eqref{eq:ch3.debye-coefficient-linearity}. The resulting
\(\E_{\ell m}\) mixes toroidal and poloidal generators, producing elliptical polarization on \(\Gamma_R\) while
preserving fixed \((\ell,m)\) labels from Eq.~\eqref{eq:ch3.J-simultaneous-eigenfield-system}
\cite{stratton1941,hansen1988}. Polarization \(\sigma\) is therefore a coefficient choice in the
\(\mathbf{a}_{\ell m}\) plane, not an independent angular index \cite{jackson1999,belinfante1940}.

\paragraph{Worked illustration (null potential channel).}
If \(a_{\ell m}=b_{\ell m}=0\) for a given \((\ell,m)\), that angular sector is silent in
Eq.~\eqref{eq:ch3.debye-separated-expansion} even when near-field scans show local phase structure
\cite{harrington2001}. Accessibility screens from Chapter~\ref{ch:02} decide which \(\mathbf{a}_{\ell m}\) are
excited by a given \(\Jimp\) \cite{stratton1941,harrington2001}.

\paragraph{Problem (construct a transverse mode from \(\Pi\) alone).}
\textbf{Problem.} Let \(\Pi=h_{\ell}^{(1)}(k_{0}r)Y_{\ell m}\) and \(\Psi=0\). Compute \(\E,\H\) from
Eqs.~\eqref{eq:ch3.debye-E-representation}--\eqref{eq:ch3.debye-H-representation} and verify
Eqs.~\eqref{eq:ch3.vector-helmholtz-E}--\eqref{eq:ch3.transverse-radiative-constraint} \cite{stratton1941,tai1994}.

\textbf{Step 1.} \(\E=\curl\curl(\Pi\hat{\mathbf{r}})\) and \(\H=\jj\omega\varepsilon_{0}\curl(\Pi\hat{\mathbf{r}})\)
\cite{tai1994}.

\textbf{Step 2.} Use Eq.~\eqref{eq:ch3.debye-scalar-helmholtz-pair} to eliminate \(\nabla^{2}\Pi\) in favor of
\(-k_{0}^{2}\Pi\) when evaluating \(\curl\curl\E\) \cite{stratton1941}.

\textbf{Step 3.} Take \(\div\E\) using Eq.~\eqref{eq:ch3.debye-E-representation} and
Eq.~\eqref{eq:ch3.debye-scalar-helmholtz-pair} to obtain \(\div\E=0\) \cite{tai1994,hansen1988}.

\textbf{Physical meaning.} The construction shows that a \emph{single} scalar potential generates a full transverse
Maxwell pair with fixed \((\ell,m)\), foreshadowing the \(\mathbf{N}_{\ell m}\) definition in Subsection~\ref{sec:ch341}
\cite{hansen1988,stratton1941}.

\paragraph{Completeness preview (forward).}
Theorem~\ref{thm:ch3.debye-representation} plus angular completeness from Subsection~\ref{sec:ch343} implies that
expansions Eq.~\eqref{eq:ch3.debye-separated-expansion} are complete for transverse fields on exterior domains in
the \(\|\cdot\|_{\mathrm{rad}}\) topology \cite{hansen1988,stratton1941}. Subsection~\ref{sec:ch334} states the
generating-side input; Subsection~\ref{sec:ch343} proves the harmonic-side closure \cite{hansen1988}.

\begin{proposition}[Separated Debye modes are divergence-free]
\label{prop:ch3.debye-div-free}
If \(\Pi\) and \(\Psi\) satisfy Eq.~\eqref{eq:ch3.debye-scalar-helmholtz-pair} and outgoing radial factors
Eq.~\eqref{eq:ch3.outgoing-radial-admissibility}, then fields from
Eqs.~\eqref{eq:ch3.debye-E-representation}--\eqref{eq:ch3.debye-H-representation} obey
Eq.~\eqref{eq:ch3.transverse-radiative-constraint} \cite{stratton1941,tai1994,hansen1988}.
\end{proposition}
\begin{proof}[Proof sketch]
Take divergence of Eq.~\eqref{eq:ch3.debye-E-representation} and use
Eq.~\eqref{eq:ch3.debye-scalar-helmholtz-pair} \cite{tai1994}.
% PROOF: AUTHOR REVIEW
\end{proof}

\paragraph{Toroidal--poloidal circulation picture.}
The term \(\curl(\Pi\hat{\mathbf{r}})\) circulates field lines on spherical shells while
\(\curl\curl(\Pi\hat{\mathbf{r}})\) produces poloidal transverse closure analogous to latitude-directed \(\E\)
\cite{tai1994,hansen1988,stratton1941}. Figure~\ref{fig:ch3.toroidal-poloidal-debye-schematic} is a schematic
reminder that Debye potentials generate \emph{vector} structure from \emph{scalar} separation labels
\cite{hansen1988}; Subsection~\ref{sec:ch341} names the resulting fields \(\mathbf{M}_{\ell m}\) and
\(\mathbf{N}_{\ell m}\) \cite{stratton1941}.

\begin{figure}[t]
\centering
\ChOneFigBlock{%
\begin{tikzpicture}[ch1fig, node distance=7mm and 10mm]
  \node[ch1box] (pi) {$\Pi\,Y_{\ell m}h_\ell^{(1)}$};
  \node[ch1box, right=12mm of pi] (tor) {$\curl(\Pi\hat{\mathbf{r}})$\\toroidal};
  \node[ch1box, right=12mm of tor] (pol) {$\curl\curl(\Pi\hat{\mathbf{r}})$\\poloidal};
  \node[ch1box, below=10mm of tor] (psi) {$\Psi$ channel\\dual};
  \draw[ch1flow] (pi) -- (tor) -- (pol);
  \draw[ch1flow] (psi) -- (tor);
\end{tikzpicture}%
}
\caption{Debye circulation schematic in Subsection~\ref{sec:ch334}: a single separated scalar multiplies into
toroidal and poloidal vector generators before \(\mathbf{M}_{\ell m}\), \(\mathbf{N}_{\ell m}\) are named in
Subsection~\ref{sec:ch341}.}
\label{fig:ch3.toroidal-poloidal-debye-schematic}
\end{figure}

\paragraph{Problem (map Debye coefficients to far-zone polarization).}
\textbf{Problem.} For fixed \((\ell,m)\), determine how \(\mathbf{a}_{\ell m}=(a,b)^{\mathsf{T}}\) in
Eq.~\eqref{eq:ch3.debye-coefficient-linearity} affects Stokes parameters on \(\Gamma_R\) \cite{stratton1941}.

\textbf{Step 1.} Build \(\E_{\ell m}[\mathbf{a}_{\ell m}]\) from Eqs.~\eqref{eq:ch3.debye-E-representation}--\eqref{eq:ch3.debye-separated-expansion}
with \(h_{\ell}^{(1)}(k_{0}r)\) \cite{tai1994,hansen1988}.

\textbf{Step 2.} Project \(\E,\H\) onto \(\hat{\boldsymbol{\theta}},\hat{\boldsymbol{\phi}}\) on \(\Gamma_R\) and form
\((S_0,S_1,S_2,S_3)\) using Eq.~\eqref{eq:ch3.stokes-helicity-proxy} \cite{stratton1941}.

\textbf{Step 3.} Vary \(\arg(b/a)\) at fixed \(|a|,|b|\): \((\ell,m)\) labels remain fixed while \(S_2,S_3\) rotate in
the Poincar\'e sphere \cite{jackson1999,belinfante1940}.

\textbf{Physical meaning.} Polarization is a coefficient phase in the Debye plane, not an independent angular quantum
number \cite{belinfante1940,stratton1941}. This is the Subsection~\ref{sec:ch334} version of the \(\sigma\) split from
Section~\ref{sec:ch32} \cite{hansen1988}.

\paragraph{Source-side preview without Green re-display.}
Impressed currents excite Debye coefficients through radiation integrals imported from Section~\ref{sec:ch25}
\cite{stratton1941,harrington2001}. Schematically,
\begin{equation}
\label{eq:ch3.debye-source-coupling-preview}
a_{\ell m}
\propto
\int_{\Omega_s}
\Jimp\cdot\mathbf{F}_{\ell m}^{(\Pi)}\,dV,
\qquad
b_{\ell m}
\propto
\int_{\Omega_s}
\Jimp\cdot\mathbf{F}_{\ell m}^{(\Psi)}\,dV,
\end{equation}
with \(\mathbf{F}_{\ell m}^{(\cdot)}\) the adjoint generating fields built from
Eqs.~\eqref{eq:ch3.debye-E-representation}--\eqref{eq:ch3.debye-H-representation} \cite{stratton1941,tai1994}.
Equation~\eqref{eq:ch3.debye-source-coupling-preview} is a forward pointer to Subsection~\ref{sec:ch344}; Subsection~\ref{sec:ch334}
does not evaluate the integrals \cite{harrington2001}.

\paragraph{Magnetic-current dual channel (structural note).}
Dual formulations interchange \(\E\) and \(\H\) through Eq.~\eqref{eq:ch3.debye-duality-exchange} as if a magnetic
current \(\mathbf{J}_{m}\) replaced the electric excitation \cite{stratton1941,tai1994}. Volume~I keeps electric
sources \(\Jimp\) as primary in Chapter~\ref{ch:02}; the magnetic dual channel is a symmetry route for
polarization and \(\sigma\) labeling, not a second source definition \cite{belinfante1940,hansen1988}. Subsection~\ref{sec:ch344}
constructs eigenfields using both \(\Pi\) and \(\Psi\) channels with coefficients
Eq.~\eqref{eq:ch3.debye-coefficient-linearity} \cite{stratton1941,tai1994}.

\paragraph{Truncation preview and far-zone fidelity.}
Truncating Eq.~\eqref{eq:ch3.debye-separated-expansion} at \(\ell_{\max}\) produces a best \(\|\cdot\|_{\mathrm{rad}}\)
approximation on \(\Gamma_R\) when \(kR\gg\ell_{\max}\) \cite{hansen1988,harrington2001}. Near the source, truncation
can still reproduce reactive fields that do not belong to \(\mathcal{F}_{\mathrm{rad}}\); export claims require
\(\Pi_{\mathrm{rad}}\) after truncation as in Eq.~\eqref{eq:ch3.reactive-radiative-helmholtz-split}
\cite{harrington2001,stratton1941}. Section~\ref{sec:ch36} analyzes completeness loss under truncation; Subsection~\ref{sec:ch334}
supplies the generating series that truncation cuts \cite{hansen1988}.

\paragraph{Problem (assemble \(\ell=0\) monopole pair from Debye potentials).}
\textbf{Problem.} Construct the isotropic outgoing pair with \(\Pi=h_{0}^{(1)}(k_{0}r)Y_{00}\), \(\Psi=0\), and verify
far-zone behavior \cite{stratton1941,jackson1999,tai1994}.

\textbf{Step 1.} Insert into Eqs.~\eqref{eq:ch3.debye-E-representation}--\eqref{eq:ch3.debye-H-representation}
\cite{tai1994}.

\textbf{Step 2.} Use \(Y_{00}=1/\sqrt{4\pi}\) once normalization is fixed in Subsection~\ref{sec:ch341}; angular
content is constant on \(S^{2}\) \cite{hansen1988}.

\textbf{Step 3.} Far-zone asymptotics follow from \(h_{0}^{(1)}(k_{0}r)\sim -(\jj/k_{0}r)e^{\jj k_{0}r}\)
\cite{jackson1999,stratton1941}.

\textbf{Step 4.} Check Eq.~\eqref{eq:ch3.helmholtz-poynting-flux} is positive and outward on \(\Gamma_R\)
\cite{stratton1941,poynting1884}.

\textbf{Physical meaning.} The \(\ell=0\) Debye-generated pair is the isotropic reference for normalization in
Section~\ref{sec:ch35} and for low-order antenna gain benchmarks \cite{harrington2001,stratton1941}.

\paragraph{Commutation with \(\mathbf{J}\) on Debye-generated fields (forward).}
When \(\Pi\) and \(\Psi\) are pure \((\ell,m)\) separated modes, the generated \(\E,\H\) are eigenfunctions of
\(\mathbf{J}^{2}\) and \(\mathbf{J}_{z}\) with the same labels as the scalar factors
\cite{jackson1999,hansen1988,stratton1941}. Subsection~\ref{sec:ch344} proves this by combining
Eqs.~\eqref{eq:ch3.debye-E-representation}--\eqref{eq:ch3.debye-separated-expansion} with the commutator results of
Subsection~\ref{sec:ch322} \cite{hansen1988}. Subsection~\ref{sec:ch334} records the generating-side hypothesis that
makes that proof possible \cite{tai1994,stratton1941}.

\begin{table}[t]
\centering
\caption{Observability of Debye coefficients Eq.~\eqref{eq:ch3.debye-separated-expansion} on \(\mathcal{F}_{\mathrm{rad}}\).}
\label{tab:ch3.334-coefficient-observability}
\small
\begin{tabular}{@{}p{0.22\linewidth}p{0.34\linewidth}p{0.36\linewidth}@{}}
\toprule
Coefficient & Observable channel & Non-observable trap \\
\midrule
\(a_{\ell m}\) &
Far-zone \(\E_\theta,\E_\phi\) &
Raw near-zone \(\Pi\) \\
\addlinespace
\(b_{\ell m}\) &
Far-zone \(\H\) dual &
Gauge-shifted \(\Psi\) \\
\addlinespace
\(\arg(b/a)\) &
Stokes \(S_2,S_3\) &
Unqualified OAM slope \\
\bottomrule
\end{tabular}
\end{table}

\begin{table}[t]
\centering
\caption{Debye generating roles in Subsection~\ref{sec:ch334} (names deferred to ?\ref{sec:ch341}).}
\label{tab:ch3.334-debye-roles}
\small
\begin{tabular}{@{}p{0.20\linewidth}p{0.34\linewidth}p{0.38\linewidth}@{}}
\toprule
Potential term & Generated structure & Section~\ref{sec:ch34} name \\
\midrule
\(\curl(\Pi\hat{\mathbf{r}})\) &
Toroidal-type \(\H\) driver &
\(\mathbf{M}_{\ell m}\) preview \\
\addlinespace
\(\curl\curl(\Pi\hat{\mathbf{r}})\) &
Poloidal-type \(\E\) driver &
\(\mathbf{N}_{\ell m}\) preview \\
\addlinespace
\(\Psi\) channel &
Dual exchange Eq.~\eqref{eq:ch3.debye-duality-exchange} &
\(\sigma\) partner \\
\bottomrule
\end{tabular}
\end{table}

Subsection~\ref{sec:ch334} has closed the separation pipeline: outgoing scalars
Eq.~\eqref{eq:ch3.debye-scalar-helmholtz-pair} generate Maxwell pairs through
Eqs.~\eqref{eq:ch3.debye-E-representation}--\eqref{eq:ch3.debye-H-representation} with coefficients
Eq.~\eqref{eq:ch3.debye-separated-expansion} \cite{stratton1941,tai1994,hansen1988}. Section~\ref{sec:ch34}
names the resulting vector harmonics and proves orthogonality on \(\mathcal{F}_{\mathrm{rad}}\)
\cite{hansen1988,jackson1999}.

The four subsections of Section~\ref{sec:ch33} form a single forward chain on \(\mathcal{F}_{\mathrm{rad}}\):
vector Helmholtz reduction Eq.~\eqref{eq:ch3.vector-helmholtz-E}, spherical separation
Eqs.~\eqref{eq:ch3.radial-helmholtz-ODE}--\eqref{eq:ch3.azimuthal-eigenvalue-equation}, outgoing radial
selection Eq.~\eqref{eq:ch3.outgoing-radial-admissibility}, and Debye generation
Eqs.~\eqref{eq:ch3.debye-E-representation}--\eqref{eq:ch3.debye-separated-expansion}. Together they supply the
analytic substrate on which the symmetry operators of Section~\ref{sec:ch32} act as quantum numbers in
Section~\ref{sec:ch34} \cite{stratton1941,hansen1988,tai1994}. No Green-function construction or Sommerfeld
uniqueness proof has been repeated; those enter only through the cited selectors of Chapter~\ref{ch:02}
\cite{harrington2001,jackson1999}.


\section{Vector Spherical Harmonics and Eigenfields}
\label{sec:ch34}

This section is the definitional heart of Volume~I representation theory: the only place where
\(Y_\ell^m\), \(\mathbf{M}_{\ell m}\), and \(\mathbf{N}_{\ell m}\) are introduced as named objects
\cite{stratton1941,hansen1988,tai1994}. Subsection~\ref{sec:ch341} states scalar and vector spherical
harmonic definitions with Condon--Shortley phase as fixed in the Volume~I notation convention. Subsection~\ref{sec:ch342}
explains toroidal, poloidal, and divergence-free structure. Subsection~\ref{sec:ch343} proves orthogonality,
completeness, and resolution of identity, importing reciprocity orthogonality previews from
Subsection~\ref{sec:ch282}. Subsection~\ref{sec:ch344} constructs angular-momentum eigenfields from the
commutation framework of Subsection~\ref{sec:ch322} and Debye potentials of Subsection~\ref{sec:ch334}.
Subsection~\ref{sec:ch345} analyzes degeneracy and representation coupling using rotational spectral results
from Subsection~\ref{sec:ch262} \cite{jackson1999,stratton1941}.


\subsection{Scalar and vector spherical harmonics (definitions)}
\label{sec:ch341}
% TOC tag: [HOME]

Volume~I names \(Y_{\ell}^{m}\), \(\mathbf{M}_{\ell m}\), and \(\mathbf{N}_{\ell m}\) only here
\cite{stratton1941,hansen1988,tai1994}. Subsection~\ref{sec:ch341} closes the separation program of
Section~\ref{sec:ch33} by attaching explicit symbols to the angular factors
Eq.~\eqref{eq:ch3.angular-laplacian-eigenvalue} and the Debye generators
Eqs.~\eqref{eq:ch3.debye-E-representation}--\eqref{eq:ch3.debye-separated-expansion}
\cite{hansen1988}.

\paragraph{Standing imports.}
Angular eigenvalue Eqs.~\eqref{eq:ch3.angular-laplacian-eigenvalue}--\eqref{eq:ch3.azimuthal-eigenvalue-equation},
outgoing radial factor Eq.~\eqref{eq:ch3.outgoing-radial-admissibility}, and carrier
Eq.~\eqref{eq:ch3.outgoing-inheritance-criterion} enter from Section~\ref{sec:ch33} by reference
\cite{stratton1941,jackson1999}. Condon--Shortley phase is fixed as in `00-NOTATION.md`
\cite{hansen1988,jackson1999}.

\paragraph{Assumptions.}
Work at fixed \(\omega\) with \(k_{0}=\omega\sqrt{\varepsilon_{0}\mu_{0}}\) on exterior domains. All harmonics
below are outgoing through \(h_{\ell}^{(1)}(k_{0}r)\) unless a cavity class is explicitly declared
\cite{stratton1941,harrington2001}.

\paragraph{Scalar spherical harmonic (HOME).}
For \(\ell\ge0\) and \(m\in\{-\ell,\ldots,\ell\}\), define
\begin{equation}
\label{eq:ch3.Ylm-Condon-Shortley-def}
Y_{\ell}^{m}(\theta,\phi)
=
(-1)^{m}
\sqrt{\frac{2\ell+1}{4\pi}\,\frac{(\ell-m)!}{(\ell+m)!}}\,
P_{\ell}^{m}(\cos\theta)\,e^{\jj m\phi},
\end{equation}
where \(P_{\ell}^{m}\) is the associated Legendre function with Condon--Shortley phase
\cite{jackson1999,hansen1988,stratton1941}. Equation~\eqref{eq:ch3.Ylm-Condon-Shortley-def} is the Volume~I
\emph{single} definition of \(Y_{\ell}^{m}\); later sections cite Eq.~\eqref{eq:ch3.Ylm-Condon-Shortley-def} by
reference only \cite{hansen1988}. The functions satisfy
Eqs.~\eqref{eq:ch3.angular-laplacian-eigenvalue}--\eqref{eq:ch3.azimuthal-eigenvalue-equation} imported from
Subsection~\ref{sec:ch332} \cite{jackson1999}.

\paragraph{Angular orthonormalization (HOME).}
\begin{equation}
\label{eq:ch3.Ylm-orthonorm-sphere}
\int_{0}^{2\pi}\!\int_{0}^{\pi}
Y_{\ell}^{m}(\theta,\phi)^{*}\,
Y_{\ell'}^{m'}(\theta,\phi)\,
\sin\theta\,d\theta\,d\phi
=
\delta_{\ell\ell'}\,\delta_{mm'}.
\end{equation}
Equation~\eqref{eq:ch3.Ylm-orthonorm-sphere} upgrades the preview
Eq.~\eqref{eq:ch3.angular-inner-product-preview} to the Condon--Shortley convention fixed in
Eq.~\eqref{eq:ch3.Ylm-Condon-Shortley-def} \cite{hansen1988,stratton1941}.

\paragraph{Vector spherical harmonics (HOME).}
Define the toroidal and poloidal vector harmonics
\begin{equation}
\label{eq:ch3.Mlm-def}
\mathbf{M}_{\ell m}(\mathbf{r})
=
\curl\!\Big(
h_{\ell}^{(1)}(k_{0}r)\,Y_{\ell}^{m}(\theta,\phi)\,\hat{\mathbf{r}}
\Big),
\end{equation}
\begin{equation}
\label{eq:ch3.Nlm-def}
\mathbf{N}_{\ell m}(\mathbf{r})
=
\frac{1}{k_{0}}\,
\curl\,\mathbf{M}_{\ell m}(\mathbf{r}),
\end{equation}
on domains where the curls are in \(\HCurl\) \cite{stratton1941,hansen1988,tai1994}. Equations
\eqref{eq:ch3.Mlm-def}--\eqref{eq:ch3.Nlm-def} are the Volume~I \emph{single} definitions of
\(\mathbf{M}_{\ell m}\) and \(\mathbf{N}_{\ell m}\); they realize the Debye preview of
Subsection~\ref{sec:ch334} with outgoing radial factors Eq.~\eqref{eq:ch3.outgoing-radial-admissibility}
\cite{tai1994,hansen1988}.

\paragraph{Maxwell field pairs from \(\mathbf{M}\) and \(\mathbf{N}\).}
Define admissible harmonic pairs
\begin{equation}
\label{eq:ch3.VSH-M-pair}
\mathbf{u}_{\ell m}^{(\mathbf{M})}
=
\big(
\jj\eta_{0}k_{0}\,\mathbf{M}_{\ell m},\
k_{0}\,\mathbf{N}_{\ell m}
\big),
\qquad
\mathbf{u}_{\ell m}^{(\mathbf{N})}
=
\big(
k_{0}\,\mathbf{N}_{\ell m},\
-\jj k_{0}\,\mathbf{M}_{\ell m}
\big),
\end{equation}
with \(\eta_{0}=\sqrt{\mu_{0}/\varepsilon_{0}}\) \cite{stratton1941,hansen1988}. Pairs
Eq.~\eqref{eq:ch3.VSH-M-pair} solve source-free Maxwell on outgoing domains and satisfy vector Helmholtz
Eqs.~\eqref{eq:ch3.vector-helmholtz-E}--\eqref{eq:ch3.vector-helmholtz-H} \cite{stratton1941,tai1994}.

\begin{figure}[t]
\centering
\ChOneFigBlock{%
\begin{tikzpicture}[ch1fig, node distance=7mm and 9mm]
  \node[ch1box] (y) {Eq.~\eqref{eq:ch3.Ylm-Condon-Shortley-def}};
  \node[ch1box, right=11mm of y] (m) {Eq.~\eqref{eq:ch3.Mlm-def}};
  \node[ch1box, right=11mm of m] (n) {Eq.~\eqref{eq:ch3.Nlm-def}};
  \node[ch1box, below=10mm of m] (hl) {Eq.~\eqref{eq:ch3.outgoing-radial-admissibility}};
  \draw[ch1flow] (y) -- (m) -- (n);
  \draw[ch1flow] (hl) -- (m);
\end{tikzpicture}%
}
\caption{Subsection~\ref{sec:ch341} definition ladder: scalar \(Y_{\ell}^{m}\) \(\to\) toroidal
\(\mathbf{M}_{\ell m}\) \(\to\) poloidal \(\mathbf{N}_{\ell m}\) with outgoing Hankel radial factors.}
\label{fig:ch3.vsh-definition-ladder}
\end{figure}

\begin{theorem}[Outgoing VSH pairs are divergence-free]
\label{thm:ch3.vsh-div-free}
Fields \(\mathbf{M}_{\ell m}\), \(\mathbf{N}_{\ell m}\) from Eqs.~\eqref{eq:ch3.Mlm-def}--\eqref{eq:ch3.Nlm-def}
and pairs Eq.~\eqref{eq:ch3.VSH-M-pair} satisfy Eq.~\eqref{eq:ch3.transverse-radiative-constraint} on exterior
domains \cite{stratton1941,hansen1988,tai1994}.
\end{theorem}
\begin{proof}[Proof sketch]
\(\div\mathbf{M}_{\ell m}=0\) from \(\div\curl=0\); \(\mathbf{N}_{\ell m}=(k_{0})^{-1}\curl\mathbf{M}_{\ell m}\)
inherits divergence-free structure \cite{tai1994,hansen1988}.
% PROOF: AUTHOR REVIEW
\end{proof}

\paragraph{Physical meaning.}
\(Y_{\ell}^{m}\) carries the angular radiation pattern; \(\mathbf{M}_{\ell m}\) and \(\mathbf{N}_{\ell m}\) are
the corresponding transverse vector carriers with correct \(k_{0}\) scaling for Maxwell pairs
\cite{stratton1941,jackson1999}. A coefficient multiplying \(Y_{\ell}^{m}\) alone is not yet a measurable
\(\E\) or \(\H\); measurable amplitudes attach to Eq.~\eqref{eq:ch3.VSH-M-pair} after normalization in
Section~\ref{sec:ch35} \cite{hansen1988,harrington2001}.

\paragraph{Validity domain.}
Eqs.~\eqref{eq:ch3.Ylm-Condon-Shortley-def}--\eqref{eq:ch3.Nlm-def} assume smoothness away from the origin and
exterior outgoing class \cite{stratton1941}. At \(\theta=0,\pi\), only combinations with correct parity remain
finite; axisymmetric modes use \(m=0\) \cite{jackson1999,hansen1988}.

\paragraph{Worked illustration (\(\ell=1\), \(m=0\) dipole scaffold).}
\(Y_{1}^{0}=\sqrt{3/(4\pi)}\cos\theta\) and \(h_{1}^{(1)}(k_{0}r)\) produce the standard dipole angular factor
\cite{jackson1999,stratton1941}. Pair \(\mathbf{u}_{10}^{(\mathbf{N})}\) from Eq.~\eqref{eq:ch3.VSH-M-pair} is the
outgoing electric-dipole template used in Subsection~\ref{sec:ch344} \cite{hansen1988}.

\paragraph{Worked illustration (\(\ell=0\) isotropic scaffold).}
\(Y_{0}^{0}=1/\sqrt{4\pi}\) is angle-independent \cite{jackson1999}. With \(h_{0}^{(1)}\), the pair
\(\mathbf{u}_{00}^{(\mathbf{N})}\) carries monopole far-zone phase \(e^{\jj k_{0}r}/r\) \cite{stratton1941}.

\paragraph{Problem (verify angular eigenvalue for \(Y_{2}^{1}\)).}
\textbf{Problem.} Show \(Y_{2}^{1}\) satisfies Eqs.~\eqref{eq:ch3.angular-laplacian-eigenvalue}--\eqref{eq:ch3.azimuthal-eigenvalue-equation}
\cite{jackson1999,hansen1988}.

\textbf{Step 1.} Separate \(\phi\)-dependence \(e^{\jj\phi}\) in Eq.~\eqref{eq:ch3.Ylm-Condon-Shortley-def}
\cite{jackson1999}.

\textbf{Step 2.} Reduce the \(\theta\) equation to associated Legendre form for \(\ell=2\), \(m=1\)
\cite{stratton1941}.

\textbf{Step 3.} Read off eigenvalues \(\ell(\ell+1)=6\) and \(m=1\) \cite{hansen1988}.

\textbf{Physical meaning.} Quadrupole-sector patterns begin at \(\ell=2\); mislabeling them as dipole is an
\(\ell\) index error, not a polarization error \cite{harrington2001,stratton1941}.

\paragraph{Associated Legendre backbone.}
The \(\theta\)-dependent factor in Eq.~\eqref{eq:ch3.Ylm-Condon-Shortley-def} solves
\begin{equation}
\label{eq:ch3.associated-Legendre-equation}
\frac{1}{\sin\theta}\frac{d}{d\theta}
\!\left(\sin\theta\frac{dP_{\ell}^{m}}{d\theta}\right)
-\frac{m^{2}}{\sin^{2}\theta}\,P_{\ell}^{m}
+\ell(\ell+1)\,P_{\ell}^{m}=0,
\end{equation}
linking Eq.~\eqref{eq:ch3.Ylm-Condon-Shortley-def} to separation
Eq.~\eqref{eq:ch3.angular-laplacian-eigenvalue} \cite{jackson1999,hansen1988,stratton1941}. Subsection~\ref{sec:ch341}
displays the closed form once; recursion relations for \(P_{\ell}^{m}\) are used numerically but not re-derived here
\cite{hansen1988}.

\paragraph{Radial?angular factorization in \(\mathbf{M},\mathbf{N}\).}
Write \(F_{\ell m}(r,\theta,\phi)=h_{\ell}^{(1)}(k_{0}r)\,Y_{\ell}^{m}(\theta,\phi)\). Then
Eq.~\eqref{eq:ch3.Mlm-def} is \(\mathbf{M}_{\ell m}=\curl(\hat{\mathbf{r}}\,F_{\ell m})\) and
Eq.~\eqref{eq:ch3.Nlm-def} completes the poloidal tower \cite{tai1994,hansen1988}. Far-zone amplitudes inherit
Eq.~\eqref{eq:ch3.hankel-far-zone-asymptotic}; angular nodes inherit zeros of \(Y_{\ell}^{m}\)
\cite{stratton1941,jackson1999}.

\paragraph{Problem (build \(\mathbf{M}_{10}\) from Debye data).}
\textbf{Problem.} Starting from \(\Pi=h_{1}^{(1)}(k_{0}r)Y_{1}^{0}\), \(\Psi=0\), identify the corresponding
\(\mathbf{M},\mathbf{N}\) sector in Eqs.~\eqref{eq:ch3.Mlm-def}--\eqref{eq:ch3.VSH-M-pair}
\cite{tai1994,stratton1941}.

\textbf{Step 1.} Map \(\curl(\Pi\hat{\mathbf{r}})\) to \(\mathbf{M}\)-type content per
Table~\ref{tab:ch3.334-debye-roles} \cite{tai1994}.

\textbf{Step 2.} Apply Eq.~\eqref{eq:ch3.Nlm-def} for \(\mathbf{N}_{10}\) \cite{hansen1988}.

\textbf{Step 3.} Assemble \(\mathbf{u}_{10}^{(\sigma)}\) in Subsection~\ref{sec:ch344} \cite{stratton1941}.

\textbf{Physical meaning.} Definitions Eqs.~\eqref{eq:ch3.Mlm-def}--\eqref{eq:ch3.Nlm-def} are not optional
normalizations; they are the unique outgoing realization of separated Debye data \cite{tai1994,hansen1988}.

\paragraph{Condon--Shortley phase contract.}
Phase \((-1)^{m}\) in Eq.~\eqref{eq:ch3.Ylm-Condon-Shortley-def} fixes interference signs in
Eq.~\eqref{eq:ch3.Ylm-orthonorm-sphere} \cite{hansen1988,jackson1999}. Changing phase convention rotates
coupling signs in Section~\ref{sec:ch38} basis changes without altering \(\mathbf{J}^{2}\) eigenvalues
\cite{stratton1941,hansen1988}.

\paragraph{Worked illustration (sectoral harmonic \(m=\ell\)).}
Sectoral \(Y_{\ell}^{\ell}\) peaks on the equatorial plane \cite{jackson1999,stratton1941}. Pairing with
\(h_{\ell}^{(1)}\) yields high-\(m\) beams with strong \(\phi\) variation on \(\Gamma_R\) \cite{harrington2001,hansen1988}.

\paragraph{Worked illustration (addition of two \(\ell\) blocks).}
Superposing \(\mathbf{u}_{10}^{(\mathbf{N})}\) and \(\mathbf{u}_{20}^{(\mathbf{N})}\) is legal in linear Maxwell
but coefficients remain distinct under Eq.~\eqref{eq:ch3.Ylm-orthonorm-sphere} \cite{stratton1941,hansen1988}. Far-zone
fitting must not collapse distinct \(\ell\) into one effective ``mode order'' \cite{harrington2001}.

\paragraph{Parity under spatial inversion.}
Under \(\mathbf{r}\mapsto-\mathbf{r}\), scalar harmonics obey
\begin{equation}
\label{eq:ch3.Ylm-parity}
Y_{\ell}^{m}(\pi-\theta,\phi+\pi)
=
(-1)^{\ell}\,Y_{\ell}^{m}(\theta,\phi),
\end{equation}
so even \(\ell\) modes are symmetric about the origin and odd \(\ell\) are antisymmetric in the angular factor
\cite{jackson1999,hansen1988,stratton1941}. Vector harmonics inherit the same \(\ell\) parity through
Eqs.~\eqref{eq:ch3.Mlm-def}--\eqref{eq:ch3.Nlm-def}; parity diagnostics on measured patterns therefore test
\(\ell\) assignment before \(\sigma\) assignment \cite{harrington2001,stratton1941}.

\paragraph{Far-zone amplitude law (outgoing scaffold).}
On \(\Gamma_R\) with \(k_{0}R\gg\ell_{\max}\), Hankel asymptotics Eq.~\eqref{eq:ch3.hankel-far-zone-asymptotic} give
\begin{equation}
\label{eq:ch3.VSH-far-zone-scaffold}
\mathbf{M}_{\ell m}\sim
\frac{\jj k_{0}\,e^{\jj k_{0}r}}{r}\,
\boldsymbol{\theta}\times\nabla_{\perp}Y_{\ell}^{m},
\qquad
\mathbf{N}_{\ell m}\sim
\frac{\jj k_{0}\,e^{\jj k_{0}r}}{r}\,
\nabla_{\perp}Y_{\ell}^{m},
\end{equation}
with \(\nabla_{\perp}\) the surface gradient on \(S^{2}\) \cite{stratton1941,hansen1988,jackson1999}. Equations
\eqref{eq:ch3.VSH-far-zone-scaffold} show that angular nodes of \(Y_{\ell}^{m}\) become radiation nulls on shells,
independent of the \(\sigma\) assembly in Eq.~\eqref{eq:ch3.VSH-M-pair} \cite{harrington2001}.

\paragraph{Problem (orthonormal overlap of two distinct \(m\)).}
\textbf{Problem.} Evaluate \(\int Y_{2}^{1*}Y_{2}^{-1}\sin\theta\,d\theta\,d\phi\) on \(S^{2}\)
\cite{hansen1988,jackson1999}.

\textbf{Step 1.} Insert Eq.~\eqref{eq:ch3.Ylm-Condon-Shortley-def} for \(\ell=2\), \(m=\pm1\) \cite{jackson1999}.

\textbf{Step 2.} Integrate \(\int_{0}^{2\pi}e^{-\jj\phi}e^{-\jj\phi}\,d\phi=0\) because \(m+m'\neq0\) \cite{hansen1988}.

\textbf{Step 3.} Conclude orthogonality via Eq.~\eqref{eq:ch3.Ylm-orthonorm-sphere} \cite{stratton1941}.

\textbf{Physical meaning.} Distinct \(m\) labels are orthogonal angular channels; beam steering rotates among
\(m\) partners rather than merging them into one scalar pattern \cite{harrington2001,jackson1999}.

\paragraph{Measurement contract (what is fixed here versus Section~\ref{sec:ch35}).}
Subsection~\ref{sec:ch341} fixes \emph{shapes} Eqs.~\eqref{eq:ch3.Ylm-Condon-Shortley-def}--\eqref{eq:ch3.VSH-M-pair}
with unit angular convention Eq.~\eqref{eq:ch3.Ylm-orthonorm-sphere} \cite{hansen1988,stratton1941}. Measurable power
amplitudes require flux normalization deferred to Section~\ref{sec:ch35}; until then, coefficients in
Eq.~\eqref{eq:ch3.VSH-resolution-identity} are geometric projections, not watt-calibrated multipole moments
\cite{poynting1884,harrington2001}.

\paragraph{Worked illustration (magnetic-dipole scaffold at \(\ell=1\), \(m=\pm1\)).}
\(\mathbf{u}_{1,\pm1}^{(\mathbf{M})}\) is the primary magnetic-dipole template paired with
\(\mathbf{u}_{1,\pm1}^{(\mathbf{N})}\) under Eq.~\eqref{eq:ch3.VSH-M-pair} \cite{stratton1941,jackson1999}. Loop feeds
that break axial symmetry excite \(|m|=1\) rather than \(m=0\) \cite{harrington2001}. The \(\ell=1\) block therefore
splits into four independent angular sectors \((m,\sigma)\), not two \cite{hansen1988}.

\begin{figure}[t]
\centering
\ChOneFigBlock{%
\begin{tikzpicture}[ch1fig, node distance=7mm and 10mm]
  \node[ch1box] (s) {scalar \(Y_{\ell}^{m}\)\\Eq.~\eqref{eq:ch3.Ylm-Condon-Shortley-def}};
  \node[ch1box, right=14mm of s] (t) {toroidal \(\mathbf{M}_{\ell m}\)\\Eq.~\eqref{eq:ch3.Mlm-def}};
  \node[ch1box, right=14mm of t] (p) {poloidal \(\mathbf{N}_{\ell m}\)\\Eq.~\eqref{eq:ch3.Nlm-def}};
  \node[ch1box, below=11mm of t] (fz) {Eq.~\eqref{eq:ch3.VSH-far-zone-scaffold}};
  \draw[ch1flow] (s) -- (t) -- (p);
  \draw[ch1flow] (fz) -- (p);
\end{tikzpicture}%
}
\caption{Angular-to-vector lift in Subsection~\ref{sec:ch341}: \(Y_{\ell}^{m}\) fixes the pattern; \(\mathbf{M},\mathbf{N}\)
carry transverse vector content with common far-zone phase law Eq.~\eqref{eq:ch3.VSH-far-zone-scaffold}.}
\label{fig:ch3.341-scalar-vector-lift}
\end{figure}

\paragraph{Associated Legendre endpoint regularity.}
At \(\theta=0,\pi\), only combinations with correct \(m\) and \(\ell\) parity remain bounded:
\(P_{\ell}^{m}\) vanishes at poles for \(|m|\ge1\), while axisymmetric \(m=0\) modes remain finite
\cite{jackson1999,hansen1988}. Subsection~\ref{sec:ch341} therefore restricts axisymmetric feeds to \(m=0\) templates
unless offset currents break symmetry and activate \(|m|\ge1\) sectors \cite{harrington2001,stratton1941}.

\paragraph{Radial factor in \(\mathbf{M}_{\ell m}\) (structure, not full expansion).}
With \(h_{\ell}^{(1)}(k_{0}r)\) outgoing,
\begin{equation}
\label{eq:ch3.Mlm-radial-angular-split}
\mathbf{M}_{\ell m}
=
\nabla\!\big(h_{\ell}^{(1)}Y_{\ell}^{m}\big)\times\hat{\mathbf{r}}
+
h_{\ell}^{(1)}\,\nabla Y_{\ell}^{m}\times\hat{\mathbf{r}},
\end{equation}
using standard vector identities for \(\curl(\hat{\mathbf{r}}F)\) \cite{stratton1941,hansen1988,tai1994}. Equation
\eqref{eq:ch3.Mlm-radial-angular-split} displays the explicit radial--angular split used in far-zone reduction
Eq.~\eqref{eq:ch3.VSH-far-zone-scaffold}: the \(1/r\) phase is carried by \(h_{\ell}^{(1)}\), while angular nodes are
carried by \(Y_{\ell}^{m}\) \cite{jackson1999}.

\paragraph{Problem (count independent labels at \(\ell=2\)).}
\textbf{Problem.} How many distinct outgoing pair labels \((m,\sigma)\) appear at \(\ell=2\)?
\cite{hansen1988,jackson1999}.

\textbf{Step 1.} List \(m\in\{-2,-1,0,1,2\}\) \cite{jackson1999}.

\textbf{Step 2.} Attach \(\sigma\in\{\mathbf{M},\mathbf{N}\}\) through Eq.~\eqref{eq:ch3.VSH-M-pair} \cite{stratton1941}.

\textbf{Step 3.} Count \(5\times2=10\) independent sectors before normalization \cite{hansen1988}.

\textbf{Physical meaning.} Multipole order \(\ell\) fixes angular complexity; \(\sigma\) doubles the transverse channels
at each \(m\) \cite{belinfante1940,harrington2001}.

\paragraph{Duality swap between \(\mathbf{M}\) and \(\mathbf{N}\) pairs.}
Interchanging electric and magnetic roles in Eq.~\eqref{eq:ch3.VSH-M-pair} maps
\(\mathbf{u}_{\ell m}^{(\mathbf{M})}\leftrightarrow\mathbf{u}_{\ell m}^{(\mathbf{N})}\) up to the sign convention in
Eq.~\eqref{eq:ch3.Nlm-def} \cite{stratton1941,hansen1988}. Subsection~\ref{sec:ch341} fixes one convention; later duality
discussions in Section~\ref{sec:ch32} cite this pair without redefining \(\mathbf{M},\mathbf{N}\) \cite{belinfante1940}.

\paragraph{Numerical evaluation contract.}
Machine evaluation of Eqs.~\eqref{eq:ch3.Ylm-Condon-Shortley-def}--\eqref{eq:ch3.Nlm-def} requires stable recurrence for
\(P_{\ell}^{m}\) and complex Hankel functions at large \(k_{0}r\) \cite{hansen1988,harrington2001}. Phase convention
Eq.~\eqref{eq:ch3.Ylm-Condon-Shortley-def} must be preserved when importing external libraries; otherwise
Eq.~\eqref{eq:ch3.Ylm-orthonorm-sphere} and coupling signs in Section~\ref{sec:ch38} disagree silently
\cite{hansen1988,jackson1999}.

\begin{table}[t]
\centering
\caption{Volume~I HOME symbols fixed in Subsection~\ref{sec:ch341} (cite by reference elsewhere).}
\label{tab:ch3.341-home-symbols}
\small
\begin{tabular}{@{}p{0.18\linewidth}p{0.38\linewidth}p{0.34\linewidth}@{}}
\toprule
Symbol & Defining equation & Role \\
\midrule
\(Y_{\ell}^{m}\) & Eq.~\eqref{eq:ch3.Ylm-Condon-Shortley-def} & Angular pattern \\
\(\mathbf{M}_{\ell m}\) & Eq.~\eqref{eq:ch3.Mlm-def} & Toroidal vector \\
\(\mathbf{N}_{\ell m}\) & Eq.~\eqref{eq:ch3.Nlm-def} & Poloidal vector \\
\bottomrule
\end{tabular}
\end{table}

\paragraph{Explicit \(\ell=1\), \(m=0\) angular factors.}
For the axisymmetric dipole scaffold,
\begin{equation}
\label{eq:ch3.Y10-explicit}
Y_{1}^{0}(\theta,\phi)=\sqrt{\frac{3}{4\pi}}\cos\theta,
\qquad
P_{1}^{0}(\cos\theta)=\cos\theta,
\end{equation}
independent of \(\phi\) \cite{jackson1999,hansen1988}. Equation~\eqref{eq:ch3.Y10-explicit} is the angular factor
in \(\mathbf{u}_{10}^{(\sigma)}\) before radial Hankel weighting Eq.~\eqref{eq:ch3.outgoing-radial-admissibility}
\cite{stratton1941}. Any measured \(\phi\)-independent far-zone pattern at \(\ell=1\) must align with
Eq.~\eqref{eq:ch3.Y10-explicit}, not with sectoral \(Y_{1}^{\pm1}\) templates \cite{harrington2001}.

\paragraph{Problem (verify orthonormality of \(Y_{1}^{0}\) and \(Y_{2}^{0}\)).}
\textbf{Problem.} Show \(\int Y_{1}^{0*}Y_{2}^{0}\sin\theta\,d\theta\,d\phi=0\) \cite{hansen1988,jackson1999}.

\textbf{Step 1.} Insert Eq.~\eqref{eq:ch3.Ylm-Condon-Shortley-def} for \(\ell=1,2\), \(m=0\) \cite{jackson1999}.

\textbf{Step 2.} Use associated Legendre orthogonality on \([-1,1]\) \cite{stratton1941}.

\textbf{Step 3.} Conclude vanishing overlap via Eq.~\eqref{eq:ch3.Ylm-orthonorm-sphere} \cite{hansen1988}.

\textbf{Physical meaning.} Distinct \(\ell\) labels are orthogonal even when both are axisymmetric (\(m=0\))
\cite{harrington2001,stratton1941}.

\paragraph{Outgoing phase inheritance in pairs.}
The \(e^{\jj k_{0}r}/r\) far-zone phase in Eq.~\eqref{eq:ch3.VSH-far-zone-scaffold} is common to all
\((\ell,m,\sigma)\) labels constructed here \cite{stratton1941,jackson1999}. Subsection~\ref{sec:ch341} therefore
separates \emph{angular complexity} (\(\ell,m\)) from \emph{radial propagation class} (outgoing Hankel) before
Section~\ref{sec:ch35} attaches flux units \cite{hansen1988,harrington2001}.

\paragraph{Addition-theorem interface (scalar, no re-derivation).}
Products \(Y_{\ell_1}^{m_1*}Y_{\ell_2}^{m_2}\) decompose into single-index harmonics with weights imported from
Eq.~\eqref{eq:ch3.CG-coupling-def} \cite{jackson1999,hansen1988}. Subsection~\ref{sec:ch341} records the interface
because far-zone pattern multiplication in calibration measurements is often misread as a single \(Y_{\ell}^{m}\)
\cite{harrington2001}. Vector products inherit the same coupling through Eqs.~\eqref{eq:ch3.Mlm-def}--\eqref{eq:ch3.Nlm-def}
\cite{stratton1941,hansen1988}.

\paragraph{Worked illustration (phase-sensitive interference of \(m=\pm1\)).}
Superposing \(Y_{1}^{1}\) and \(Y_{1}^{-1}\) produces a \(\phi\)-dependent standing pattern on \(S^{2}\) before outgoing
lift \cite{jackson1999,hansen1988}. After lift to \(\mathbf{u}_{1,\pm1}^{(\sigma)}\), relative phases set the
major-axis orientation of the \(\ell=1\) multiplet \cite{harrington2001,stratton1941}. Condon--Shortley signs in
Eq.~\eqref{eq:ch3.Ylm-Condon-Shortley-def} fix those interference fringes; changing convention rotates the pattern
without altering \(\mathbf{J}^{2}\) labels \cite{hansen1988}.

\begin{table}[t]
\centering
\caption{Subsection~\ref{sec:ch341} definitional contracts (single HOME display).}
\label{tab:ch3.341-def-contracts}
\small
\begin{tabular}{@{}p{0.22\linewidth}p{0.36\linewidth}p{0.34\linewidth}@{}}
\toprule
Contract & Equation & Later use \\
\midrule
Scalar phase & Eq.~\eqref{eq:ch3.Ylm-Condon-Shortley-def} & cite only \\
Angular norm & Eq.~\eqref{eq:ch3.Ylm-orthonorm-sphere} & Subsection~\ref{sec:ch343} \\
Toroidal vector & Eq.~\eqref{eq:ch3.Mlm-def} & cite only \\
Poloidal vector & Eq.~\eqref{eq:ch3.Nlm-def} & cite only \\
Maxwell pair & Eq.~\eqref{eq:ch3.VSH-M-pair} & Subsections~\ref{sec:ch344}--\ref{sec:ch345} \\
\bottomrule
\end{tabular}
\end{table}

Subsection~\ref{sec:ch341} has defined \(Y_{\ell}^{m}\), \(\mathbf{M}_{\ell m}\), and \(\mathbf{N}_{\ell m}\) with
outgoing radial dependence \cite{stratton1941,hansen1988,tai1994}. Subsection~\ref{sec:ch342} analyzes toroidal,
poloidal, and divergence-free structure without reintroducing symbols \cite{hansen1988}.


\subsection{Toroidal, poloidal, and divergence-free structure}
\label{sec:ch342}
% TOC tag: [HOME]

Definitions Eqs.~\eqref{eq:ch3.Mlm-def}--\eqref{eq:ch3.Nlm-def} name the objects; Subsection~\ref{sec:ch342}
explains their geometric and Maxwell structure on \(\mathcal{F}_{\mathrm{rad}}\) \cite{tai1994,hansen1988,stratton1941}.

\paragraph{Radial projection identities.}
For \(\mathbf{M}_{\ell m}\) and \(\mathbf{N}_{\ell m}\) defined in Subsection~\ref{sec:ch341},
\begin{equation}
\label{eq:ch3.MN-radial-orthogonality}
\hat{\mathbf{r}}\cdot\mathbf{M}_{\ell m}=0,
\qquad
\hat{\mathbf{r}}\cdot\mathbf{N}_{\ell m}=0,
\end{equation}
on source-free exterior domains \cite{stratton1941,hansen1988,tai1994}. Equations~\eqref{eq:ch3.MN-radial-orthogonality}
are the Subsection~\ref{sec:ch342} transverse certificates: toroidal and poloidal fields have no longitudinal
\(\hat{\mathbf{r}}\) component in the radiation gauge fixed by radial Debye orientation
Eq.~\eqref{eq:ch3.debye-E-representation} \cite{tai1994}. Far-zone transversality in
Eq.~\eqref{eq:ch3.transverse-radiative-constraint} is therefore built into the harmonic basis, not imposed
post hoc \cite{stratton1941,harrington2001}.

\paragraph{Toroidal circulation content.}
\(\mathbf{M}_{\ell m}\) circulates around the polar axis for \(|m|\ge1\): field lines close on spherical shells
without poloidal closure through the origin \cite{hansen1988,tai1994}. Physically, \(\mathbf{M}_{\ell m}\) drives
\(\H\)-dominant channels in Eq.~\eqref{eq:ch3.VSH-M-pair} while \(\mathbf{N}_{\ell m}\) carries the paired
\(\E\)-weighted content \cite{stratton1941}. Laboratory loop antennas excite \(\mathbf{M}\)-weighted sectors when
currents circulate azimuthally \cite{harrington2001,jackson1999}.

\paragraph{Poloidal closure content.}
\(\mathbf{N}_{\ell m}=(k_{0})^{-1}\curl\mathbf{M}_{\ell m}\) completes transverse electric-type structure
\cite{stratton1941,hansen1988}. Poloidal fields connect \(\theta\)-directed \(\E\) components across latitudes;
short dipoles excite \(\mathbf{N}\)-weighted sectors at \(\ell=1\) \cite{jackson1999,stratton1941}. The toroidal
and poloidal split is the vector realization of the Debye channels \(\Pi\) and \(\Psi\) in
Subsection~\ref{sec:ch334} \cite{tai1994}.

\begin{figure}[t]
\centering
\ChOneFigBlock{%
\begin{tikzpicture}[ch1fig, node distance=7mm and 10mm]
  \node[ch1box] (m) {$\mathbf{M}_{\ell m}$\\toroidal};
  \node[ch1box, right=13mm of m] (n) {$\mathbf{N}_{\ell m}$\\poloidal};
  \node[ch1box, below=11mm of m] (div) {Eq.~\eqref{eq:ch3.MN-radial-orthogonality}};
  \node[ch1box, right=13mm of div] (pair) {Eq.~\eqref{eq:ch3.VSH-M-pair}};
  \draw[ch1flow] (m) -- node[above,font=\footnotesize]{$k_0^{-1}\curl$} (n);
  \draw[ch1flow] (div) -- (m);
  \draw[ch1flow] (pair) -- (n);
\end{tikzpicture}%
}
\caption{Toroidal--poloidal split in Subsection~\ref{sec:ch342}: \(\mathbf{M}_{\ell m}\) and \(\mathbf{N}_{\ell m}\)
are transverse and combine into Maxwell pairs Eq.~\eqref{eq:ch3.VSH-M-pair}.}
\label{fig:ch3.toroidal-poloidal-split}
\end{figure}

\begin{proposition}[Divergence-free closure of the \(\mathbf{M},\mathbf{N}\) span]
\label{prop:ch3.MN-div-free-span}
Any linear combination \(\mathbf{E}=a\,\mathbf{M}_{\ell m}+b\,\mathbf{N}_{\ell m}\) with constants \(a,b\) obeys
\(\div\mathbf{E}=0\) \cite{stratton1941,tai1994,hansen1988}. Pairs Eq.~\eqref{eq:ch3.VSH-M-pair} therefore inherit
Eq.~\eqref{eq:ch3.transverse-radiative-constraint} from Theorem~\ref{thm:ch3.vsh-div-free}
\cite{stratton1941}.
\end{proposition}
\begin{proof}[Proof sketch]
Linearity of divergence and \(\div\mathbf{M}_{\ell m}=\div\mathbf{N}_{\ell m}=0\) \cite{tai1994}.
% PROOF: AUTHOR REVIEW
\end{proof}

\paragraph{Debye equivalence (structural).}
Eqs.~\eqref{eq:ch3.Mlm-def}--\eqref{eq:ch3.Nlm-def} realize the toroidal and poloidal channels previewed in
Table~\ref{tab:ch3.334-debye-roles}: \(\curl(\Pi\hat{\mathbf{r}})\) generates \(\mathbf{M}\)-type content and
\(\curl\curl(\Pi\hat{\mathbf{r}})\) generates \(\mathbf{N}\)-type content at fixed \((\ell,m)\) \cite{tai1994,hansen1988}.
Subsection~\ref{sec:ch342} records the geometric identification; Subsection~\ref{sec:ch344} constructs eigenfields
from linear combinations \cite{stratton1941}.

\paragraph{Worked illustration (loop versus dipole excitation).}
A small loop in the \(xy\)-plane primarily couples to \(\mathbf{M}_{1,\pm1}\) sectors; a short \(\hat{\mathbf{z}}\)
dipole couples to \(\mathbf{N}_{1,0}\) \cite{jackson1999,harrington2001}. The distinction is toroidal versus
poloidal weighting in Eqs.~\eqref{eq:ch3.Mlm-def}--\eqref{eq:ch3.Nlm-def}, not a different \(k_{0}\)
\cite{stratton1941}.

\paragraph{Worked illustration (high-\(\ell\) suppression).}
Centrifugal barrier Eq.~\eqref{eq:ch3.radial-helmholtz-canonical} suppresses high-\(\ell\) \(\mathbf{M},\mathbf{N}\)
near compact sources \cite{jackson1999,hansen1988}. Far-zone fits that require large \(\ell\) with small
\(h_{\ell}^{(1)}\) amplitude usually indicate near-field contamination or wrong outgoing branch
\cite{harrington2001,stratton1941}.

\paragraph{Problem (show \(\hat{\mathbf{r}}\cdot\mathbf{M}_{\ell m}=0\)).}
\textbf{Problem.} Verify Eq.~\eqref{eq:ch3.MN-radial-orthogonality} for \(\mathbf{M}_{\ell m}\)
\cite{stratton1941,tai1994}.

\textbf{Step 1.} Write \(\mathbf{M}_{\ell m}=\curl(\hat{\mathbf{r}}\,F_{\ell m})\) with
\(F_{\ell m}=h_{\ell}^{(1)}Y_{\ell}^{m}\) \cite{hansen1988}.

\textbf{Step 2.} Use \(\hat{\mathbf{r}}\cdot\curl\mathbf{A}=\nabla\cdot(\hat{\mathbf{r}}\times\mathbf{A})\) and
periodicity in \(\phi\) on the sphere at fixed \(r\) \cite{jackson1999,stratton1941}.

\textbf{Step 3.} Conclude \(\hat{\mathbf{r}}\cdot\mathbf{M}_{\ell m}=0\) \cite{tai1994}.

\textbf{Physical meaning.} Longitudinal \(\hat{\mathbf{r}}\cdot\E\) leak in far-zone fits contradicts
Eq.~\eqref{eq:ch3.MN-radial-orthogonality} and signals non-transverse superpositions \cite{harrington2001}.

\paragraph{Curl interrelations among \(\mathbf{M}\) and \(\mathbf{N}\).}
On source-free domains,
\begin{equation}
\label{eq:ch3.MN-curl-interrelations}
\curl\mathbf{M}_{\ell m}
=
k_{0}\,\mathbf{N}_{\ell m},
\qquad
\curl\mathbf{N}_{\ell m}
=
k_{0}\,\mathbf{M}_{\ell m},
\end{equation}
up to the sign convention in Eq.~\eqref{eq:ch3.VSH-M-pair} \cite{hansen1988,stratton1941,tai1994}. Equations
\eqref{eq:ch3.MN-curl-interrelations} are the Subsection~\ref{sec:ch342} Maxwell certificate that toroidal and
poloidal towers are linked by \(k_{0}\), not independent vector bases \cite{stratton1941}.

\paragraph{Problem (classify a loop feed).}
\textbf{Problem.} A planar loop in the \(xy\)-plane primarily excites which harmonic sector?
\cite{harrington2001,jackson1999}.

\textbf{Step 1.} Identify dominant azimuthal current \(\Rightarrow\) \(\mathbf{M}\)-weighted coupling
\cite{harrington2001}.

\textbf{Step 2.} Select \(\ell=1\), \(m=\pm1\) candidates from Eq.~\eqref{eq:ch3.Mlm-def} \cite{jackson1999}.

\textbf{Step 3.} Verify far-zone \(\hat{\mathbf{r}}\cdot\E\approx0\) using Eq.~\eqref{eq:ch3.MN-radial-orthogonality}
\cite{stratton1941}.

\textbf{Physical meaning.} Feed geometry selects toroidal/poloidal weighting before \(\ell\) truncation enters
\cite{harrington2001,stratton1941}.

\paragraph{Surface-flux interpretation.}
Toroidal \(\mathbf{M}_{\ell m}\) circulates azimuthal \(\H\)-weighted flux on \(\Gamma_R\); poloidal \(\mathbf{N}_{\ell m}\)
carries latitudinal \(\E\)-weighted closure \cite{stratton1941,hansen1988}. Subsection~\ref{sec:ch342} therefore ties
harmonic labels to Poynting transport channels previewed in Eq.~\eqref{eq:ch3.helmholtz-poynting-flux}
\cite{poynting1884,stratton1941}.

\paragraph{Worked illustration (near-zone versus far-zone toroidal content).}
\(\mathbf{M}_{\ell m}\) can be large near sources while export remains \(\mathbf{N}\)-dominated in the far zone for
electric dipoles \cite{harrington2001,stratton1941}. \(\Pi_{\mathrm{rad}}\) must be applied before interpreting near-zone
\(\mathbf{M}\) magnitude as radiative \cite{harrington2001}.

\paragraph{Vector Helmholtz action on \(\mathbf{M},\mathbf{N}\).}
Source-free Maxwell implies vector Helmholtz Eqs.~\eqref{eq:ch3.vector-helmholtz-E}--\eqref{eq:ch3.vector-helmholtz-H}
on \(\mathbf{M}_{\ell m}\) and \(\mathbf{N}_{\ell m}\) \cite{stratton1941,hansen1988}. With
\(F_{\ell m}=h_{\ell}^{(1)}(k_{0}r)Y_{\ell}^{m}\),
\begin{equation}
\label{eq:ch3.MN-helmholtz-eigenvalue}
\curl\curl\mathbf{M}_{\ell m}
=
k_{0}^{2}\,\mathbf{M}_{\ell m},
\qquad
\curl\curl\mathbf{N}_{\ell m}
=
k_{0}^{2}\,\mathbf{N}_{\ell m},
\end{equation}
on domains where double curls are defined \cite{tai1994,hansen1988,stratton1941}. Equations
\eqref{eq:ch3.MN-helmholtz-eigenvalue} certify that toroidal and poloidal towers are not separate Helmholtz spectra;
they are linked by Eq.~\eqref{eq:ch3.MN-curl-interrelations} \cite{stratton1941}.

\paragraph{Solenoidal decomposition on shells.}
Any transverse \(\mathbf{E}_t\) on a spherical shell admits
\begin{equation}
\label{eq:ch3.shell-solenoidal-split}
\mathbf{E}_t
=
\sum_{\ell m}\big(
a_{\ell m}\,\mathbf{M}_{\ell m}
+
b_{\ell m}\,\mathbf{N}_{\ell m}
\big),
\end{equation}
with coefficients fixed by tangential inner products when outgoing data are specified \cite{hansen1988,stratton1941}.
Equation~\eqref{eq:ch3.shell-solenoidal-split} is the Subsection~\ref{sec:ch342} geometric reading of
Eq.~\eqref{eq:ch3.VSH-resolution-identity} before normalization Section~\ref{sec:ch35} \cite{harrington2001}.

\paragraph{Problem (Maxwell pair transversality).}
\textbf{Problem.} Show \(\hat{\mathbf{r}}\cdot\E=0\) for \(\mathbf{u}_{\ell m}^{(\mathbf{N})}\) in the far zone
\cite{stratton1941,hansen1988}.

\textbf{Step 1.} Write \(\E\) from Eq.~\eqref{eq:ch3.VSH-M-pair} as \(k_{0}\mathbf{N}_{\ell m}\) \cite{stratton1941}.

\textbf{Step 2.} Apply Eq.~\eqref{eq:ch3.MN-radial-orthogonality} \cite{tai1994}.

\textbf{Step 3.} Confirm \(\H\) partner is \(\mathbf{M}\)-weighted and transverse likewise \cite{hansen1988}.

\textbf{Physical meaning.} Transversality is a property of the harmonic pair, not an optional gauge in far-zone
measurements \cite{harrington2001,jackson1999}.

\paragraph{Antenna feed map (compact sources).}
Small loops map to \(\mathbf{M}_{1,m}\) dominance; short linear dipoles map to \(\mathbf{N}_{1,m}\) dominance
\cite{harrington2001,jackson1999,stratton1941}. Offset feeds and skew currents mix both towers at the same
\((\ell,m)\) through superposition Eq.~\eqref{eq:ch3.shell-solenoidal-split}, but orthogonality of distinct labels
remains Eq.~\eqref{eq:ch3.VSH-orthogonality-radiative} \cite{hansen1988}. Feed engineering therefore controls
\(\sigma\) weighting before \(\ell\) truncation \cite{harrington2001}.

\paragraph{Worked illustration (quadrupole toroidal circulation).}
At \(\ell=2\), \(\mathbf{M}_{2,0}\) exhibits two azimuthal circulation cells on \(\Gamma_R\) separated by a
\(Y_{2}^{0}\) nodal cone \cite{stratton1941,hansen1988}. Poloidal partner \(\mathbf{N}_{2,0}\) closes latitudinal
\(\E_t\) across those cells \cite{jackson1999}. The \(\ell=2\) block is the lowest order where toroidal and poloidal
patterns differ visibly from dipole templates \cite{harrington2001}.

\begin{figure}[t]
\centering
\ChOneFigBlock{%
\begin{tikzpicture}[ch1fig, node distance=7mm and 9mm]
  \node[ch1box] (c1) {Eq.~\eqref{eq:ch3.MN-curl-interrelations}};
  \node[ch1box, right=12mm of c1] (c2) {Eq.~\eqref{eq:ch3.MN-helmholtz-eigenvalue}};
  \node[ch1box, right=12mm of c2] (sp) {Eq.~\eqref{eq:ch3.shell-solenoidal-split}};
  \node[ch1box, below=10mm of c1] (tr) {Eq.~\eqref{eq:ch3.MN-radial-orthogonality}};
  \draw[ch1flow] (tr) -- (c1) -- (c2) -- (sp);
\end{tikzpicture}%
}
\caption{Subsection~\ref{sec:ch342} Maxwell chain: transverse projections \(\to\) curl linkage \(\to\) Helmholtz
eigenvalue \(\to\) shell solenoidal expansion.}
\label{fig:ch3.342-maxwell-chain}
\end{figure}

\paragraph{Poynting transport split on \(\Gamma_R\).}
Time-averaged Poynting flux Eq.~\eqref{eq:ch3.helmholtz-poynting-flux} decomposes into \(\mathbf{M}\)-weighted and
\(\mathbf{N}\)-weighted transport channels in Eq.~\eqref{eq:ch3.VSH-M-pair} \cite{poynting1884,stratton1941}. Toroidal
\(\mathbf{M}_{\ell m}\) contributes azimuthal circulation of \(\S=\Re(\E\times\H^{*})/2\); poloidal \(\mathbf{N}_{\ell m}\)
completes latitudinal closure \cite{hansen1988}. Subsection~\ref{sec:ch342} therefore links harmonic geometry to flux
diagnostics used in Section~\ref{sec:ch35} \cite{poynting1884,harrington2001}.

\paragraph{Shell matching without re-solving Maxwell.}
Given tangential \(\mathbf{E}_t,\mathbf{H}_t\) on \(\Gamma_R\) in \(\mathcal{F}_{\mathrm{rad}}\), coefficients in
Eq.~\eqref{eq:ch3.shell-solenoidal-split} are fixed by reciprocity pairings once outgoing data are specified
\cite{stratton1941,harrington2001}. Near-zone reactive content that violates
Eq.~\eqref{eq:ch3.outgoing-inheritance-criterion} must be removed before interpreting shell projections as radiative
\cite{harrington2001}.

\paragraph{Worked illustration (slot versus loop on \(\ell=1\)).}
A narrow slot along \(\hat{\mathbf{z}}\) excites \(\mathbf{N}_{1,0}\); a loop in the \(xy\)-plane excites
\(\mathbf{M}_{1,\pm1}\) \cite{harrington2001,jackson1999}. Both are \(\ell=1\) but differ in toroidal/poloidal weighting
and in \(m\) \cite{stratton1941}. Pattern fitting that reports only ``dipole'' without \((m,\sigma)\) loses half the
representation information \cite{hansen1988}.

\paragraph{Problem (divergence of a mixed combination).}
\textbf{Problem.} For \(\mathbf{E}=a\mathbf{M}_{\ell m}+b\mathbf{N}_{\ell m}\), show \(\div\mathbf{E}=0\)
\cite{tai1994,stratton1941}.

\textbf{Step 1.} Use linearity of divergence \cite{jackson1999}.

\textbf{Step 2.} Apply Theorem~\ref{thm:ch3.vsh-div-free} to each term \cite{hansen1988}.

\textbf{Step 3.} Conclude \(\div\mathbf{E}=0\) for all \(a,b\) \cite{stratton1941}.

\textbf{Physical meaning.} Mixed toroidal/poloidal superpositions remain transverse solenoidal fields; orthogonality of
distinct \emph{labels} is a separate statement in Subsection~\ref{sec:ch343} \cite{harrington2001}.

\paragraph{Centrifugal barrier and feed size.}
Compact sources of characteristic size \(a\ll\lambda_0=2\pi/k_0\) prefer low \(\ell\) because radial barrier
Eq.~\eqref{eq:ch3.radial-helmholtz-canonical} suppresses high-\(\ell\) content near \(r\sim a\) \cite{jackson1999,hansen1988}.
Subsection~\ref{sec:ch342} explains why electrically large feeds are required to access high-\(\ell\) toroidal/poloidal
sectors in measurable amplitudes \cite{harrington2001,stratton1941}.

\begin{table}[t]
\centering
\caption{Toroidal versus poloidal vector content in Subsection~\ref{sec:ch342}.}
\label{tab:ch3.342-toroidal-poloidal}
\small
\begin{tabular}{@{}p{0.20\linewidth}p{0.34\linewidth}p{0.36\linewidth}@{}}
\toprule
Object & Circulation geometry & Typical feed \\
\midrule
\(\mathbf{M}_{\ell m}\) & Azimuthal on shells & Loop, ring \\
\(\mathbf{N}_{\ell m}\) & Latitudinal closure & Short dipole \\
\bottomrule
\end{tabular}
\end{table}

\paragraph{Transverse projection operator view.}
Define tangential projection \(\mathbf{P}_{t}=\mathbf{I}-\hat{\mathbf{r}}\hat{\mathbf{r}}\) on spherical shells.
Then Eqs.~\eqref{eq:ch3.MN-radial-orthogonality}--\eqref{eq:ch3.MN-helmholtz-eigenvalue} state that \(\mathbf{M}_{\ell m}\)
and \(\mathbf{N}_{\ell m}\) lie in \(\mathrm{Im}(\mathbf{P}_{t})\) and satisfy solenoidal closure
\cite{stratton1941,hansen1988,tai1994}. Subsection~\ref{sec:ch342} is the geometric face of the analytic constraint
Eq.~\eqref{eq:ch3.transverse-radiative-constraint} imported from Subsection~\ref{sec:ch331} \cite{stratton1941}.

\paragraph{Worked illustration (horizontally polarized short dipole).}
A \(\hat{\mathbf{x}}\)-directed short dipole at the origin excites \(\mathbf{N}_{1,0}\) dominantly when aligned with
\(\hat{\mathbf{z}}\) axis conventions of Eq.~\eqref{eq:ch3.VSH-M-pair} \cite{jackson1999,harrington2001}. Rotating the
dipole mixes \(m\) partners within \(\ell=1\) through Eq.~\eqref{eq:ch3.eigenfield-Wigner-mixing} previewed in
Subsection~\ref{sec:ch345}, but does not change \(\ell\) \cite{stratton1941,hansen1988}.

\paragraph{Problem (identify poloidal content from \(\curl\E\)).}
\textbf{Problem.} Given outgoing \(\E\) proportional to \(\mathbf{N}_{\ell m}\), identify the paired \(\H\) sector
\cite{stratton1941,hansen1988}.

\textbf{Step 1.} Use Eq.~\eqref{eq:ch3.VSH-M-pair} for \(\sigma=\mathbf{N}\) \cite{stratton1941}.

\textbf{Step 2.} Read \(\H\propto\mathbf{M}_{\ell m}\) \cite{hansen1988}.

\textbf{Step 3.} Verify transversality with Eq.~\eqref{eq:ch3.MN-radial-orthogonality} \cite{tai1994}.

\textbf{Physical meaning.} Poloidal \(\mathbf{N}\) and toroidal \(\mathbf{M}\) are Maxwell partners, not independent
polarizations \cite{harrington2001,jackson1999}.

\paragraph{Reactive contamination warning.}
Near-zone \(\mathbf{M}_{\ell m}\) magnitudes can exceed \(\mathbf{N}_{\ell m}\) magnitudes for electrically small
electric dipoles \cite{harrington2001,stratton1941}. Interpreting such near-zone ratios as ``dominant toroidal radiation''
violates \(\Pi_{\mathrm{rad}}\) filtering from Chapter~\ref{ch:02} \cite{harrington2001}. Subsection~\ref{sec:ch342}
documents geometry; export qualification remains mandatory \cite{stratton1941}.

\paragraph{Toroidal--poloidal power channel bookkeeping (pre-normalization).}
Before Section~\ref{sec:ch35}, compare \(\|\mathbf{M}_{\ell m}\|\) and \(\|\mathbf{N}_{\ell m}\|\) only as
\emph{geometric} indicators on \(\Gamma_R\) \cite{hansen1988,stratton1941}. Ratio trends across \(\ell\) diagnose
feed type: loop feeds rise in \(\mathbf{M}\)-weighted content; short dipoles rise in \(\mathbf{N}\)-weighted content
\cite{harrington2001,jackson1999}. Absolute watts require flux normalization not yet fixed \cite{poynting1884}.

\paragraph{Worked illustration (coaxial loop stack).}
A stack of coaxial loops preserves \(m=0\) symmetry in \(\mathbf{M}_{\ell 0}\) sectors while increasing \(\ell\)
through multiple circumferences \cite{harrington2001,stratton1941}. Poloidal \(\mathbf{N}_{\ell 0}\) partners remain
subdominant unless the stack is electrically short in the axial direction \cite{jackson1999,hansen1988}. The illustration
shows how geometry selects toroidal weighting before harmonic truncation \cite{harrington2001}.

\begin{table}[t]
\centering
\caption{Subsection~\ref{sec:ch342} diagnostic map (geometry \(\to\) harmonic weighting).}
\label{tab:ch3.342-diagnostic-map}
\small
\begin{tabular}{@{}p{0.28\linewidth}p{0.32\linewidth}p{0.32\linewidth}@{}}
\toprule
Feed geometry & Dominant tower & Typical \(\ell\) start \\
\midrule
Planar loop & \(\mathbf{M}_{\ell m}\), \(|m|\ge1\) & \(\ell=1\) \\
Short dipole & \(\mathbf{N}_{\ell m}\), often \(m=0\) & \(\ell=1\) \\
Offset pair & Mixed \(\mathbf{M},\mathbf{N}\) & \(\ell\ge2\) \\
\bottomrule
\end{tabular}
\end{table}

\paragraph{Problem (classify a crossed-dipole feed).}
\textbf{Problem.} Two short dipoles along \(\hat{\mathbf{x}}\) and \(\hat{\mathbf{y}}\) are fed in quadrature.
Identify dominant toroidal/poloidal sectors at \(\ell=1\) \cite{harrington2001,jackson1999,stratton1941}.

\textbf{Step 1.} Decompose each dipole into \(\mathbf{N}_{1,m}\) content with distinct \(m\) \cite{jackson1999}.

\textbf{Step 2.} Superpose \(\mathbf{M}\) and \(\mathbf{N}\) towers at the same \(\ell=1\) \cite{stratton1941}.

\textbf{Step 3.} Verify far-zone transversality Eq.~\eqref{eq:ch3.MN-radial-orthogonality} \cite{tai1994}.

\textbf{Physical meaning.} Quadrature feeds generate controlled \(m\) mixtures within a fixed \(\ell\) block; they do not
automatically raise \(\ell\) \cite{harrington2001,hansen1988}.

\paragraph{Far-zone \(\mathbf{E},\mathbf{H}\) phase relation in pairs.}
For outgoing pairs Eq.~\eqref{eq:ch3.VSH-M-pair}, \(\mathbf{M}_{\ell m}\) and \(\mathbf{N}_{\ell m}\) maintain the
Maxwell phase relation \(\H\propto \curl\E\) on source-free domains \cite{stratton1941,hansen1988}. Toroidal/poloidal
classification therefore survives in both field constituents: a feed cannot be \(\mathbf{M}\)-dominant in \(\H\) while
remaining \(\mathbf{N}\)-dominant in \(\E\) at the same \((\ell,m)\) label \cite{harrington2001,tai1994}.

\paragraph{Worked illustration (cavity versus exterior caveat).}
Interior standing-wave cavities can mimic toroidal circulation without satisfying
Eq.~\eqref{eq:ch3.outgoing-inheritance-criterion} \cite{stratton1941,harrington2001}. Subsection~\ref{sec:ch342} documents
free-space exterior structure; cavity chapters apply different radial classes \cite{jackson1999,collin2001}. Export
interpretation always requires the outgoing carrier of Section~\ref{sec:ch31} \cite{kato1995,stratton1941}.

Subsection~\ref{sec:ch342} has identified \(\mathbf{M}_{\ell m}\) as toroidal and \(\mathbf{N}_{\ell m}\) as poloidal
divergence-free carriers \cite{hansen1988,tai1994,stratton1941}. Subsection~\ref{sec:ch343} proves orthogonality and
completeness on \(\mathcal{F}_{\mathrm{rad}}\) \cite{stratton1941}.


\subsection{Orthogonality, completeness, resolution of identity}
\label{sec:ch343}
% TOC tag: [HOME][USE SS2.8.2]

Mode orthogonality is the certificate that eigenfield coefficients are independent amplitudes. Subsection~\ref{sec:ch343}
upgrades Chapter~\ref{ch:02} modal orthogonality to explicit VSH integrals on \(\mathcal{F}_{\mathrm{rad}}\)
\cite{stratton1941,hansen1988,harrington2001}.

\paragraph{Standing imports.}
Radiation-modal orthogonality Eq.~\eqref{eq:ch2.radiation-modal-orthogonality}, radiative coupling bilinear
Eq.~\eqref{eq:ch2.radiative-coupling-bilinear}, and Theorem~\ref{thm:ch2.radiation-modal-orthogonality-coupling}
enter from Subsection~\ref{sec:ch282} by reference \cite{stratton1941,harrington2001}. Radiative inner product
Eq.~\eqref{eq:ch3.radiative-inner-product-def} and Hilbert closure Eq.~\eqref{eq:ch3.F-rad-hilbert-closure} are
imported from Section~\ref{sec:ch31} \cite{hansen1988,kato1995}.

\paragraph{Vector harmonic orthogonality (HOME).}
For distinct mode labels \((\ell,m,\sigma)\) and \((\ell',m',\sigma')\) on \(\Gamma_{R}\) with \(k_{0}R\gg\ell\),
\begin{equation}
\label{eq:ch3.VSH-orthogonality-radiative}
\langle \mathbf{u}_{\ell m}^{(\sigma)},\,\mathbf{u}_{\ell' m'}^{(\sigma')}\rangle_{\mathrm{rad}}
=
0,
\qquad
(\ell,m,\sigma)\neq(\ell',m',\sigma'),
\end{equation}
and diagonal normalization is deferred to Section~\ref{sec:ch35} \cite{stratton1941,hansen1988}. Equation
\eqref{eq:ch3.VSH-orthogonality-radiative} is the Subsection~\ref{sec:ch343} specialization of
Eq.~\eqref{eq:ch2.radiation-modal-orthogonality} to the harmonic basis
Eqs.~\eqref{eq:ch3.Mlm-def}--\eqref{eq:ch3.VSH-M-pair} \cite{harrington2001}.

\paragraph{Angular reduction.}
Eq.~\eqref{eq:ch3.VSH-orthogonality-radiative} reduces to angular orthogonality
Eq.~\eqref{eq:ch3.Ylm-orthonorm-sphere} when radial overlap integrals of \(h_{\ell}^{(1)}\) are evaluated on
asymptotic shells \cite{hansen1988,stratton1941}. Cross-terms between \(\mathbf{M}\) and \(\mathbf{N}\) at the same
\((\ell,m)\) vanish under reciprocity pairing Eq.~\eqref{eq:ch3.reciprocity-symmetric-inner-product}
\cite{stratton1941,harrington2001}.

\begin{theorem}[VSH orthogonality on \(\mathcal{F}_{\mathrm{rad}}\)]
\label{thm:ch3.vsh-orthogonality}
Let \(\mathbf{u}_{\ell m}^{(\sigma)}\) be outgoing harmonic pairs built from
Eqs.~\eqref{eq:ch3.Mlm-def}--\eqref{eq:ch3.VSH-M-pair} with \(\sigma\in\{\mathbf{M},\mathbf{N}\}\). Then
Eq.~\eqref{eq:ch3.VSH-orthogonality-radiative} holds for distinct \((\ell,m,\sigma)\) labels on exterior domains
with Eq.~\eqref{eq:ch3.outgoing-inheritance-criterion} \cite{stratton1941,hansen1988,harrington2001}.
\end{theorem}
\begin{proof}[Proof sketch]
Combine Eq.~\eqref{eq:ch3.Ylm-orthonorm-sphere}, outgoing Hankel orthogonality preview
Eq.~\eqref{eq:ch3.radial-orthogonality-preview}, and reciprocity symmetry
Eq.~\eqref{eq:ch3.reciprocity-symmetric-inner-product} \cite{hansen1988,stratton1941}.
% PROOF: AUTHOR REVIEW
\end{proof}

\paragraph{Resolution of identity (HOME).}
Any \(\mathbf{v}_\omega\in\mathcal{F}_{\mathrm{rad}}\) admits
\begin{equation}
\label{eq:ch3.VSH-resolution-identity}
\mathbf{v}_\omega
=
\sum_{\ell m\sigma}
c_{\ell m}^{(\sigma)}\,\mathbf{u}_{\ell m}^{(\sigma)},
\qquad
c_{\ell m}^{(\sigma)}
=
\langle \mathbf{u}_{\ell m}^{(\sigma)},\mathbf{v}_\omega\rangle_{\mathrm{rad}}
\ \text{(normalized in Section~\ref{sec:ch35})},
\end{equation}
with convergence in \(\|\cdot\|_{\mathrm{rad}}\) \cite{hansen1988,stratton1941,kato1995}. Equation
\eqref{eq:ch3.VSH-resolution-identity} upgrades preview Eq.~\eqref{eq:ch3.eigenfield-expansion-on-frad} to a
completeness statement on the outgoing harmonic basis \cite{hansen1988}.

\begin{figure}[t]
\centering
\ChOneFigBlock{%
\begin{tikzpicture}[ch1fig, node distance=7mm and 9mm]
  \node[ch1box] (ch2) {Eq.~\eqref{eq:ch2.radiation-modal-orthogonality}};
  \node[ch1box, right=12mm of ch2] (vsh) {Eq.~\eqref{eq:ch3.VSH-orthogonality-radiative}};
  \node[ch1box, right=12mm of vsh] (res) {Eq.~\eqref{eq:ch3.VSH-resolution-identity}};
  \node[ch1box, below=10mm of vsh] (fr) {Eq.~\eqref{eq:ch3.F-rad-hilbert-closure}};
  \draw[ch1flow] (ch2) -- (vsh) -- (res);
  \draw[ch1flow] (fr) -- (res);
\end{tikzpicture}%
}
\caption{Subsection~\ref{sec:ch343} upgrade chain: Chapter~\ref{ch:02} modal orthogonality \(\to\) VSH orthogonality
\(\to\) resolution of identity on \(\mathcal{F}_{\mathrm{rad}}\).}
\label{fig:ch3.vsh-orthogonality-chain}
\end{figure}

\paragraph{Worked illustration (two-mode superposition).}
\(\mathbf{v}=c_{10}^{(\mathbf{N})}\mathbf{u}_{10}^{(\mathbf{N})}+c_{21}^{(\mathbf{M})}\mathbf{u}_{21}^{(\mathbf{M})}\) has
orthogonal coefficients under Eq.~\eqref{eq:ch3.VSH-orthogonality-radiative} \cite{hansen1988,stratton1941}. Power
budgets decompose additively only after \(\Pi_{\mathrm{rad}}\) and normalization Section~\ref{sec:ch35}
\cite{harrington2001}.

\paragraph{Problem (extract \(c_{\ell m}^{(\sigma)}\) on \(\Gamma_R\)).}
\textbf{Problem.} Given \(\mathbf{v}_\omega\in\mathcal{F}_{\mathrm{rad}}\), extract coefficient
\(c_{10}^{(\mathbf{N})}\) using Eq.~\eqref{eq:ch3.VSH-resolution-identity} \cite{hansen1988,stratton1941}.

\textbf{Step 1.} Form \(\langle \mathbf{u}_{10}^{(\mathbf{N})},\mathbf{v}_\omega\rangle_{\mathrm{rad}}\) on a shell
with \(k_{0}R\gg1\) \cite{stratton1941}.

\textbf{Step 2.} Divide by the \(\ell=1,m=0\) normalization fixed in Section~\ref{sec:ch35} \cite{hansen1988}.

\textbf{Step 3.} Verify orthogonality kills \(c_{21}^{(\mathbf{M})}\) cross-terms \cite{harrington2001}.

\textbf{Physical meaning.} Coefficient extraction is an inner-product measurement, not a peak-fit of
\(|E|\) alone \cite{harrington2001,stratton1941}.

\paragraph{Parseval-type angular statement.}
For square-integrable angular patterns \(f(\theta,\phi)\),
\begin{equation}
\label{eq:ch3.parseval-sphere}
\sum_{\ell m}|a_{\ell m}|^{2}
=
\int_{0}^{2\pi}\!\int_{0}^{\pi}|f|^{2}\sin\theta\,d\theta\,d\phi,
\qquad
a_{\ell m}=\langle Y_{\ell}^{m},f\rangle_{\Omega},
\end{equation}
with \(\langle\cdot,\cdot\rangle_{\Omega}\) from Eq.~\eqref{eq:ch3.Ylm-orthonorm-sphere} \cite{hansen1988,stratton1941}.
Equation~\eqref{eq:ch3.parseval-sphere} is the angular backbone of Eq.~\eqref{eq:ch3.VSH-resolution-identity}
\cite{hansen1988}.

\begin{theorem}[Harmonic completeness on \(\mathcal{F}_{\mathrm{rad}}\)]
\label{thm:ch3.vsh-completeness}
The outgoing set \(\{\mathbf{u}_{\ell m}^{(\sigma)}\}\) is complete in \(\mathcal{F}_{\mathrm{rad}}\) for the inner
product Eq.~\eqref{eq:ch3.radiative-inner-product-def} \cite{hansen1988,stratton1941,kato1995}. Every
\(\mathcal{R}_\omega[\mathfrak{D}]\) admits expansion Eq.~\eqref{eq:ch3.VSH-resolution-identity} with convergence in
\(\|\cdot\|_{\mathrm{rad}}\) \cite{stratton1941,harrington2001}.
\end{theorem}
\begin{proof}[Proof sketch]
Combine Theorem~\ref{thm:ch3.vsh-orthogonality}, Eq.~\eqref{eq:ch3.F-rad-hilbert-closure}, and standard spherical
harmonic completeness on \(S^{2}\) \cite{hansen1988,kato1995}.
% PROOF: AUTHOR REVIEW
\end{proof}

\paragraph{Reciprocity-symmetric reduction.}
Eq.~\eqref{eq:ch3.reciprocity-symmetric-inner-product} identifies left and right inner products for lossless
reciprocal media, making Eq.~\eqref{eq:ch3.VSH-orthogonality-radiative} compatible with
Eq.~\eqref{eq:ch2.radiative-coupling-bilinear} \cite{stratton1941,harrington2001}. Subsection~\ref{sec:ch343} is
therefore the harmonic specialization of Subsection~\ref{sec:ch282}, not an independent postulate
\cite{stratton1941}.

\paragraph{Problem (orthogonality of two dipole eigenfields).}
\textbf{Problem.} Show \(\langle \mathbf{u}_{10}^{(\mathbf{N})},\mathbf{u}_{11}^{(\mathbf{M})}\rangle_{\mathrm{rad}}=0\)
\cite{hansen1988,stratton1941}.

\textbf{Step 1.} Reduce to distinct \((\ell,m,\sigma)\) labels in Eq.~\eqref{eq:ch3.VSH-orthogonality-radiative}
\cite{stratton1941}.

\textbf{Step 2.} Invoke Theorem~\ref{thm:ch3.vsh-orthogonality} \cite{hansen1988}.

\textbf{Physical meaning.} Orthogonality is a label property, not a spatial overlap guess from pattern plots
\cite{harrington2001}.

\paragraph{Radial overlap on outgoing shells.}
For fixed \(\ell\) and asymptotic \(h_{\ell}^{(1)}\),
\begin{equation}
\label{eq:ch3.radial-hankel-orthogonality-shell}
\lim_{R\to\infty}
\int_{\Gamma_R}
h_{\ell}^{(1)}(k_{0}r)^{*}\,
h_{\ell'}^{(1)}(k_{0}r)\,dS
\propto
\delta_{\ell\ell'},
\end{equation}
consistent with preview Eq.~\eqref{eq:ch3.radial-orthogonality-preview} \cite{stratton1941,hansen1988}. Together with
Eq.~\eqref{eq:ch3.Ylm-orthonorm-sphere}, Eq.~\eqref{eq:ch3.radial-hankel-orthogonality-shell} factorizes
Eq.~\eqref{eq:ch3.VSH-orthogonality-radiative} into radial and angular parts on \(\Gamma_R\) \cite{hansen1988}.

\paragraph{Cross-sector cancellation at fixed \((\ell,m)\).}
At fixed \((\ell,m)\), \(\mathbf{M}_{\ell m}\) and \(\mathbf{N}_{\ell m}\) are not proportional; their radiative
pairing vanishes under reciprocity-symmetric inner product Eq.~\eqref{eq:ch3.reciprocity-symmetric-inner-product}
\cite{stratton1941,harrington2001}. Subsection~\ref{sec:ch343} therefore treats \(\sigma\) as an orthogonal label
alongside \((\ell,m)\) \cite{hansen1988}.

\paragraph{Truncation preview (finite \(\ell_{\max}\)).}
Finite sums \(\sum_{\ell\le\ell_{\max}}\sum_m\) define approximants in \(\mathcal{F}_{\mathrm{rad}}\); error control
is deferred to Section~\ref{sec:ch36} \cite{harrington2001,stratton1941}. Orthogonality of retained modes is exact;
incompleteness is a truncation issue, not a failure of Eq.~\eqref{eq:ch3.VSH-orthogonality-radiative}
\cite{hansen1988,kato1995}.

\paragraph{Problem (Parseval check on \(S^{2}\)).}
\textbf{Problem.} For \(f=Y_{1}^{0}+Y_{2}^{0}\), verify Eq.~\eqref{eq:ch3.parseval-sphere} \cite{hansen1988,jackson1999}.

\textbf{Step 1.} Read coefficients \(a_{10}=1\), \(a_{20}=1\), all other \(a_{\ell m}=0\) from
Eq.~\eqref{eq:ch3.Ylm-orthonorm-sphere} \cite{jackson1999}.

\textbf{Step 2.} Sum \(|a_{10}|^{2}+|a_{20}|^{2}=2\) \cite{hansen1988}.

\textbf{Step 3.} Integrate \(|f|^{2}\) over \(S^{2}\) using orthogonality \cite{stratton1941}.

\textbf{Physical meaning.} Angular power partitions add only on orthogonal labels; non-orthogonal superpositions
require cross terms \cite{harrington2001}.

\paragraph{Worked illustration (receive projection on \(\Gamma_R\)).}
A probe field \(\mathbf{v}_\omega\) matched to \(\mathbf{u}_{\ell m}^{(\sigma)}\) extracts
\(c_{\ell m}^{(\sigma)}\) linearly through Eq.~\eqref{eq:ch3.VSH-resolution-identity} \cite{stratton1941,harrington2001}.
Mismatch in \(\sigma\) or \(m\) yields null coupling in the lossless reciprocal class
Eq.~\eqref{eq:ch3.reciprocity-symmetric-inner-product} \cite{harrington2001}. Harmonic receivers are label filters,
not arbitrary pattern correlators \cite{hansen1988}.

\begin{figure}[t]
\centering
\ChOneFigBlock{%
\begin{tikzpicture}[ch1fig, node distance=7mm and 9mm]
  \node[ch1box] (ang) {Eq.~\eqref{eq:ch3.Ylm-orthonorm-sphere}};
  \node[ch1box, right=12mm of ang] (rad) {Eq.~\eqref{eq:ch3.radial-hankel-orthogonality-shell}};
  \node[ch1box, right=12mm of rad] (vsh) {Eq.~\eqref{eq:ch3.VSH-orthogonality-radiative}};
  \node[ch1box, below=10mm of rad] (pv) {Eq.~\eqref{eq:ch3.parseval-sphere}};
  \draw[ch1flow] (ang) -- (vsh);
  \draw[ch1flow] (rad) -- (vsh);
  \draw[ch1flow] (pv) -- (ang);
\end{tikzpicture}%
}
\caption{Factorization chain in Subsection~\ref{sec:ch343}: angular Parseval backbone and radial Hankel orthogonality
combine into full VSH orthogonality on \(\mathcal{F}_{\mathrm{rad}}\).}
\label{fig:ch3.343-orthogonality-factorization}
\end{figure}

\paragraph{Shell inner product scaffold.}
On \(\Gamma_R\), Eq.~\eqref{eq:ch3.radiative-inner-product-def} specializes to a surface bilinear form whose leading
far-zone term factorizes into Eqs.~\eqref{eq:ch3.Ylm-orthonorm-sphere} and
\eqref{eq:ch3.radial-hankel-orthogonality-shell} \cite{stratton1941,hansen1988}. Subsection~\ref{sec:ch343} therefore
treats coefficient extraction as a shell measurement protocol, not a volume integral over reactive regions
\cite{harrington2001}.

\paragraph{Diagonal normalization deferral.}
Theorem~\ref{thm:ch3.vsh-orthogonality} fixes \emph{vanishing} off-diagonal pairings; diagonal values
\(\langle\mathbf{u}_{\ell m}^{(\sigma)},\mathbf{u}_{\ell m}^{(\sigma)}\rangle_{\mathrm{rad}}\) are set in
Section~\ref{sec:ch35} \cite{stratton1941,hansen1988}. Until then, ratios of coefficients are meaningful; absolute
watts are not \cite{poynting1884,harrington2001}.

\paragraph{Worked illustration (three-mode superposition budget).}
For \(\mathbf{v}=c_{10}^{(\mathbf{N})}\mathbf{u}_{10}^{(\mathbf{N})}+c_{11}^{(\mathbf{M})}\mathbf{u}_{11}^{(\mathbf{M})}
+c_{20}^{(\mathbf{N})}\mathbf{u}_{20}^{(\mathbf{N})}\), orthogonality kills all cross terms in
Eq.~\eqref{eq:ch3.VSH-orthogonality-radiative} \cite{hansen1988,stratton1941}. After normalization, export power splits
into three independent channels; fitting one dipole template cannot absorb quadrupole residuals
\cite{harrington2001}.

\paragraph{Problem (detect wrong outgoing branch).}
\textbf{Problem.} A coefficient extraction on \(\Gamma_R\) is unstable under \(R\) scaling. Diagnose the failure
\cite{stratton1941,harrington2001}.

\textbf{Step 1.} Test outgoing criterion Eq.~\eqref{eq:ch3.outgoing-inheritance-criterion} \cite{stratton1941}.

\textbf{Step 2.} Check Hankel branch in Eqs.~\eqref{eq:ch3.Mlm-def}--\eqref{eq:ch3.Nlm-def} \cite{hansen1988}.

\textbf{Step 3.} Apply \(\Pi_{\mathrm{rad}}\) before inner products \cite{harrington2001}.

\textbf{Physical meaning.} Orthogonality is defined on \(\mathcal{F}_{\mathrm{rad}}\); standing or reactive components
invalidate Eq.~\eqref{eq:ch3.VSH-orthogonality-radiative} \cite{stratton1941,kato1995}.

\paragraph{Completeness versus finite arrays.}
Theorem~\ref{thm:ch3.vsh-completeness} is an infinite-\(\ell\) statement on \(\mathcal{F}_{\mathrm{rad}}\); finite
antenna arrays realize truncated sums analyzed in Section~\ref{sec:ch36} \cite{harrington2001,stratton1941}. Subsection~\ref{sec:ch343}
supplies the infinite basis certificate; truncation error is a separate engineering layer \cite{hansen1988}.

\begin{table}[t]
\centering
\caption{Orthogonality objects in Subsection~\ref{sec:ch343}.}
\label{tab:ch3.343-orthogonality-objects}
\small
\begin{tabular}{@{}p{0.26\linewidth}p{0.34\linewidth}p{0.32\linewidth}@{}}
\toprule
Bilinear & Domain & Chapter anchor \\
\midrule
Eq.~\eqref{eq:ch3.Ylm-orthonorm-sphere} & \(S^{2}\) & Subsection~\ref{sec:ch341} \\
Eq.~\eqref{eq:ch3.VSH-orthogonality-radiative} & \(\mathcal{F}_{\mathrm{rad}}\) & Eq.~\eqref{eq:ch2.radiation-modal-orthogonality} \\
Eq.~\eqref{eq:ch3.VSH-resolution-identity} & \(\mathcal{F}_{\mathrm{rad}}\) & Eq.~\eqref{eq:ch3.F-rad-hilbert-closure} \\
\bottomrule
\end{tabular}
\end{table}

\paragraph{Biorthogonality remark (lossless class).}
In reciprocal lossless media, Eq.~\eqref{eq:ch3.reciprocity-symmetric-inner-product} makes left and right pairings
identical, so VSH orthogonality is standard Hilbert orthogonality rather than biorthogonality with distinct adjoints
\cite{stratton1941,harrington2001}. Subsection~\ref{sec:ch343} states the symmetric case used throughout Volume~I
exterior problems \cite{hansen1988}.

\paragraph{Worked illustration (probe mismatch loss).}
A receive probe aligned to \(\mathbf{u}_{10}^{(\mathbf{N})}\) extracts \(c_{10}^{(\mathbf{N})}\) but couples weakly to
\(\mathbf{u}_{11}^{(\mathbf{M})}\) because labels differ \cite{harrington2001,stratton1941}. The null is exact under
Theorem~\ref{thm:ch3.vsh-orthogonality}, not approximate \cite{hansen1988}. Probe design is therefore label matching on
\((\ell,m,\sigma)\) \cite{harrington2001}.

\paragraph{Problem (finite shell radius sensitivity).}
\textbf{Problem.} Coefficient extraction at \(k_{0}R=10\) shows \(5\%\) drift versus \(k_{0}R=30\). Diagnose
\cite{stratton1941,harrington2001}.

\textbf{Step 1.} Check \(k_{0}R\gg\ell_{\max}\) requirement in Eq.~\eqref{eq:ch3.VSH-orthogonality-radiative}
\cite{hansen1988}.

\textbf{Step 2.} Verify Hankel asymptotics Eq.~\eqref{eq:ch3.hankel-far-zone-asymptotic} \cite{stratton1941}.

\textbf{Step 3.} Increase \(R\) or apply \(\Pi_{\mathrm{rad}}\) \cite{harrington2001}.

\textbf{Physical meaning.} Orthogonality integrals are asymptotic shell statements; too-small \(R\) retains near-zone
reactive overlap \cite{stratton1941,kato1995}.

\paragraph{Resolution identity versus bare expansion preview.}
Eq.~\eqref{eq:ch3.VSH-resolution-identity} strengthens preview Eq.~\eqref{eq:ch3.eigenfield-expansion-on-frad} by
supplying orthogonality certificates and \(\mathcal{F}_{\mathrm{rad}}\) convergence \cite{hansen1988,kato1995}. Coefficients
are now \emph{projections}, not formal series terms \cite{stratton1941,harrington2001}.

\paragraph{Gram matrix preview (diagonal deferred).}
Let \(G_{(\ell m\sigma),(\ell'm'\sigma')}=\langle\mathbf{u}_{\ell m}^{(\sigma)},\mathbf{u}_{\ell'm'}^{(\sigma')}\rangle_{\mathrm{rad}}\).
Theorem~\ref{thm:ch3.vsh-orthogonality} makes \(G\) block-diagonal in \((\ell,m,\sigma)\) with off-diagonal zeros
\cite{hansen1988,stratton1941}. Section~\ref{sec:ch35} sets diagonal entries to unity (or flux-normalized values)
\cite{poynting1884,harrington2001}. Subsection~\ref{sec:ch343} therefore supplies the \emph{label orthogonality} certificate
independent of amplitude calibration \cite{hansen1988}.

\paragraph{Worked illustration (calibration source with known \(\ell=1\) content).}
A calibrated dipole standard radiates \(\mathbf{u}_{10}^{(\mathbf{N})}\) with unknown overall scale but known label
\cite{harrington2001,stratton1941}. A second unknown source is decomposed by projecting onto the standard and onto
orthogonal partners \(\mathbf{u}_{11}^{(\mathbf{M})}\), \(\mathbf{u}_{20}^{(\mathbf{N})}\), etc., using
Eq.~\eqref{eq:ch3.VSH-resolution-identity} \cite{hansen1988}. Residuals after projection quantify mislabeling or
non-outgoing contamination \cite{harrington2001}.

\begin{table}[t]
\centering
\caption{Subsection~\ref{sec:ch343} measurement protocol on \(\Gamma_R\).}
\label{tab:ch3.343-measurement-protocol}
\small
\begin{tabular}{@{}p{0.08\linewidth}p{0.42\linewidth}p{0.38\linewidth}@{}}
\toprule
Step & Operation & Pass criterion \\
\midrule
1 & Apply \(\Pi_{\mathrm{rad}}\) & Eq.~\eqref{eq:ch3.outgoing-inheritance-criterion} \\
2 & Shell inner product & Eq.~\eqref{eq:ch3.radiative-inner-product-def} \\
3 & Project on \(\mathbf{u}_{\ell m}^{(\sigma)}\) & Eq.~\eqref{eq:ch3.VSH-resolution-identity} \\
4 & Normalize amplitude & Section~\ref{sec:ch35} \\
\bottomrule
\end{tabular}
\end{table}

\paragraph{Worked illustration (orthogonal triad at \(\ell=1\)).}
The three partners \(\mathbf{u}_{1,-1}^{(\mathbf{N})}\), \(\mathbf{u}_{10}^{(\mathbf{N})}\), \(\mathbf{u}_{11}^{(\mathbf{N})}\)
are mutually orthogonal under Theorem~\ref{thm:ch3.vsh-orthogonality} \cite{hansen1988,stratton1941}. Any measured
\(\ell=1\) electric-type pattern is a unique linear combination of those three once \(\sigma=\mathbf{N}\) is fixed
\cite{harrington2001}. Adding \(\sigma=\mathbf{M}\) partners doubles the independent channels to six at \(\ell=1\)
\cite{belinfante1940,hansen1988}.

\paragraph{Adjoint-mode remark (lossless preview).}
When media are reciprocal and lossless, adjoint modes coincide with forward modes under
Eq.~\eqref{eq:ch3.reciprocity-symmetric-inner-product} \cite{stratton1941,harrington2001}. Subsection~\ref{sec:ch343}
therefore does not introduce a separate biorthogonal basis in the exterior vacuum class \cite{hansen1988}. Dissipative
extensions are deferred to later volumes \cite{stratton1941}.

\paragraph{Problem (project a skewed pattern onto VSH basis).}
\textbf{Problem.} A measured \(\mathbf{E}\) on \(\Gamma_R\) is not aligned to any single \(\mathbf{u}_{\ell m}^{(\sigma)}\).
Extract the first three coefficients in Eq.~\eqref{eq:ch3.VSH-resolution-identity} \cite{hansen1988,harrington2001}.

\textbf{Step 1.} Apply \(\Pi_{\mathrm{rad}}\) to the sampled field \cite{harrington2001}.

\textbf{Step 2.} Compute \(\langle\mathbf{u}_{10}^{(\mathbf{N})},\mathbf{v}\rangle_{\mathrm{rad}}\),
\(\langle\mathbf{u}_{11}^{(\mathbf{M})},\mathbf{v}\rangle_{\mathrm{rad}}\),
\(\langle\mathbf{u}_{20}^{(\mathbf{N})},\mathbf{v}\rangle_{\mathrm{rad}}\) \cite{stratton1941}.

\textbf{Step 3.} Verify residuals orthogonal to the tested basis set \cite{hansen1988}.

\textbf{Physical meaning.} Multipole analysis is sequential projection, not single-template curve fitting
\cite{harrington2001,stratton1941}.

\begin{figure}[t]
\centering
\ChOneFigBlock{%
\begin{tikzpicture}[ch1fig, node distance=7mm and 9mm]
  \node[ch1box] (m) {measure on $\Gamma_R$};
  \node[ch1box, right=12mm of m] (p) {$\Pi_{\mathrm{rad}}$};
  \node[ch1box, right=12mm of p] (ip) {Eq.~\eqref{eq:ch3.radiative-inner-product-def}};
  \node[ch1box, right=12mm of ip] (c) {Eq.~\eqref{eq:ch3.VSH-resolution-identity}};
  \draw[ch1flow] (m) -- (p) -- (ip) -- (c);
\end{tikzpicture}%
}
\caption{Subsection~\ref{sec:ch343} coefficient extraction pipeline on a finite shell.}
\label{fig:ch3.343-extraction-pipeline}
\end{figure}

\paragraph{Completeness error budget (qualitative).}
Let \(\mathbf{v}_\omega^{(\ell_{\max})}\) be the best approximation using modes up to \(\ell_{\max}\). Orthogonality
implies the error is orthogonal to the truncated subspace under Eq.~\eqref{eq:ch3.VSH-orthogonality-radiative}
\cite{hansen1988,kato1995}. Subsection~\ref{sec:ch343} certifies the full basis; Section~\ref{sec:ch36} quantifies
truncation loss when \(\ell_{\max}\) is finite \cite{harrington2001,stratton1941}. Until then, rising residual power
in high-\(\ell\) projections signals need for larger \(\ell_{\max}\) \cite{hansen1988}.

\paragraph{Worked illustration (orthogonality in \(\sigma\) at \(\ell=1,m=0\)).}
\(\mathbf{u}_{10}^{(\mathbf{M})}\) and \(\mathbf{u}_{10}^{(\mathbf{N})}\) share \((\ell,m)\) but differ in \(\sigma\); their
radiative pairing vanishes under Theorem~\ref{thm:ch3.vsh-orthogonality} \cite{belinfante1940,hansen1988,stratton1941}.
Electric and magnetic dipole templates are therefore distinct measurable channels even when angular plots look similar
\cite{harrington2001,jackson1999}.

Subsection~\ref{sec:ch343} has established VSH orthogonality and resolution of identity
Eqs.~\eqref{eq:ch3.VSH-orthogonality-radiative}--\eqref{eq:ch3.VSH-resolution-identity} on
\(\mathcal{F}_{\mathrm{rad}}\) \cite{stratton1941,hansen1988}. Subsection~\ref{sec:ch344} constructs simultaneous
\(\mathbf{J}\) eigenfields in that basis \cite{jackson1999}.


\subsection{Construction of angular-momentum eigenfields}
\label{sec:ch344}
% TOC tag: [HOME][USE SS3.2.2][USE SS3.3.4]

Orthogonal harmonics become physical eigenfields when they diagonalize \(\mathbf{J}\). Subsection~\ref{sec:ch344}
constructs \(\mathbf{u}_{\ell m}^{(\sigma)}\) satisfying Eq.~\eqref{eq:ch3.J-simultaneous-eigenfield-system} from
VSH pairs and Debye generators \cite{hansen1988,tai1994,jackson1999}.

\paragraph{Standing imports.}
Simultaneous eigenfield system Eq.~\eqref{eq:ch3.J-simultaneous-eigenfield-system}, ladder
Eq.~\eqref{eq:ch3.J-ladder-action}, quantum-number split Eq.~\eqref{eq:ch3.eigenfield-quantum-number-split}, and
commutation results of Subsection~\ref{sec:ch322} enter by reference \cite{jackson1999,hansen1988}. Debye maps
Eqs.~\eqref{eq:ch3.debye-E-representation}--\eqref{eq:ch3.debye-separated-expansion} enter from
Subsection~\ref{sec:ch334} \cite{tai1994,stratton1941}.

\paragraph{Eigenfield definition (HOME).}
Define outgoing angular-momentum eigenfields
\begin{equation}
\label{eq:ch3.eigenfield-VSH-def}
\mathbf{u}_{\ell m}^{(\sigma)}
=
\begin{cases}
\mathbf{u}_{\ell m}^{(\mathbf{M})}, & \sigma=\mathbf{M},\\[4pt]
\mathbf{u}_{\ell m}^{(\mathbf{N})}, & \sigma=\mathbf{N},
\end{cases}
\end{equation}
with \(\mathbf{u}_{\ell m}^{(\mathbf{M})}\), \(\mathbf{u}_{\ell m}^{(\mathbf{N})}\) from
Eq.~\eqref{eq:ch3.VSH-M-pair} and symbols from Eqs.~\eqref{eq:ch3.Mlm-def}--\eqref{eq:ch3.Nlm-def}
\cite{stratton1941,hansen1988}. Polarization/handness label \(\sigma\) is the sector projector target of
Eq.~\eqref{eq:ch3.polarization-sector-projector} imported from Section~\ref{sec:ch32} \cite{belinfante1940,hansen1988}.

\begin{theorem}[VSH pairs are \(\mathbf{J}\) eigenfields]
\label{thm:ch3.vsh-J-eigenfields}
Fields \(\mathbf{u}_{\ell m}^{(\sigma)}\) in Eq.~\eqref{eq:ch3.eigenfield-VSH-def} satisfy
Eq.~\eqref{eq:ch3.J-simultaneous-eigenfield-system} and belong to \(\mathcal{F}_{\mathrm{rad}}\)
\cite{jackson1999,hansen1988,stratton1941,tai1994}.
\end{theorem}
\begin{proof}[Proof sketch]
Separation labels \((\ell,m)\) from Eqs.~\eqref{eq:ch3.angular-laplacian-eigenvalue}--\eqref{eq:ch3.azimuthal-eigenvalue-equation}
commute with \(\mathbf{J}\) action established in Subsection~\ref{sec:ch322}; Debye construction
Eqs.~\eqref{eq:ch3.debye-E-representation}--\eqref{eq:ch3.debye-separated-expansion} preserves those labels
\cite{tai1994,hansen1988,jackson1999}.
% PROOF: AUTHOR REVIEW
\end{proof}

\paragraph{Coefficient coupling to sources.}
Eq.~\eqref{eq:ch3.debye-source-coupling-preview} specializes to
\begin{equation}
\label{eq:ch3.eigenfield-source-projection}
c_{\ell m}^{(\sigma)}
\propto
\int_{\Omega_s}
\Jimp\cdot\mathbf{F}_{\ell m}^{(\sigma)}\,dV,
\qquad
\mathbf{F}_{\ell m}^{(\sigma)}\in\{\mathbf{M}_{\ell m},\mathbf{N}_{\ell m}\},
\end{equation}
on admissible sources \cite{stratton1941,harrington2001,tai1994}. Equation~\eqref{eq:ch3.eigenfield-source-projection}
is the Subsection~\ref{sec:ch344} amplitude law: accessible currents from Chapter~\ref{ch:02} determine which
\((\ell,m,\sigma)\) sectors appear in \(\mathcal{R}_\omega[\mathfrak{D}]\) \cite{harrington2001,stratton1941}.

\paragraph{Worked illustration (dipole eigenfield).}
\(\mathbf{u}_{10}^{(\mathbf{N})}\) is the primary short-dipole eigenfield; \(\mathbf{u}_{10}^{(\mathbf{M})}\) is dual
\cite{jackson1999,stratton1941}. Measuring \(\sigma\) requires polarization observables from Section~\ref{sec:ch32},
not \(\phi\) slope alone \cite{belinfante1940,harrington2001}.

\paragraph{Problem (verify \(\mathbf{J}_{z}\) eigenvalue).}
\textbf{Problem.} Show \(\mathbf{J}_{z}\mathbf{u}_{\ell m}^{(\sigma)}=m\,\mathbf{u}_{\ell m}^{(\sigma)}\)
\cite{jackson1999,hansen1988}.

\textbf{Step 1.} Use \(e^{\jj m\phi}\) dependence in Eq.~\eqref{eq:ch3.Ylm-Condon-Shortley-def} \cite{jackson1999}.

\textbf{Step 2.} Apply \(\mathbf{J}_{z}=-\jj\partial/\partial\phi\) on separated factors \cite{hansen1988}.

\textbf{Step 3.} Conclude eigenvalue \(m\) for both \(\sigma=\mathbf{M},\mathbf{N}\) \cite{stratton1941}.

\textbf{Physical meaning.} \(m\) is fixed by azimuthal phase, independent of \(\sigma\) \cite{jackson1999,belinfante1940}.

\paragraph{\(\mathbf{J}^{2}\) eigenvalue certificate.}
Because \(\mathbf{u}_{\ell m}^{(\sigma)}\) is built from \(Y_{\ell}^{m}\) satisfying
Eq.~\eqref{eq:ch3.angular-laplacian-eigenvalue},
\begin{equation}
\label{eq:ch3.eigenfield-Jsq-eigenvalue}
\mathbf{J}^{2}\mathbf{u}_{\ell m}^{(\sigma)}
=
\ell(\ell+1)\,\mathbf{u}_{\ell m}^{(\sigma)},
\end{equation}
on admissible domains \cite{jackson1999,hansen1988,stratton1941}. Equation~\eqref{eq:ch3.eigenfield-Jsq-eigenvalue} is the
Subsection~\ref{sec:ch344} elevation of the scalar separation law to full Maxwell pairs; \(\sigma\) does not shift
\(\ell(\ell+1)\) \cite{belinfante1940,hansen1988}.

\paragraph{Ladder action (imported specialization).}
Raising/lowering operators from Eq.~\eqref{eq:ch3.J-ladder-action} map
\(\mathbf{u}_{\ell m}^{(\sigma)}\) to \(\mathbf{u}_{\ell,m\pm1}^{(\sigma)}\) within the same \(\ell\) block
\cite{jackson1999,hansen1988}. Subsection~\ref{sec:ch344} uses that law to move among \(m\) partners without changing
\(\ell\) or \(\sigma\) \cite{stratton1941}. Ladder steps are representation transport, not new harmonic definitions
\cite{hansen1988}.

\paragraph{Debye-to-eigenfield pipeline.}
Given separated Debye data Eq.~\eqref{eq:ch3.debye-separated-expansion}, the construction path is
\(\Pi,\Psi\to\mathbf{M},\mathbf{N}\to\mathbf{u}_{\ell m}^{(\sigma)}\) through
Eqs.~\eqref{eq:ch3.Mlm-def}--\eqref{eq:ch3.eigenfield-VSH-def} \cite{tai1994,stratton1941}. Each stage preserves
\((\ell,m)\) labels established in Subsection~\ref{sec:ch332} \cite{hansen1988}. Failure at any stage---wrong
outgoing branch, wrong \(\ell\) truncation, or missing \(\Pi_{\mathrm{rad}}\)---produces fields that are not
simultaneous \(\mathbf{J}\) eigenfields even if patterns resemble multipoles \cite{harrington2001,tai1994}.

\begin{figure}[t]
\centering
\ChOneFigBlock{%
\begin{tikzpicture}[ch1fig, node distance=7mm and 9mm]
  \node[ch1box] (d) {Eq.~\eqref{eq:ch3.debye-separated-expansion}};
  \node[ch1box, right=11mm of d] (mn) {Eqs.~\eqref{eq:ch3.Mlm-def}--\eqref{eq:ch3.Nlm-def}};
  \node[ch1box, right=11mm of mn] (ef) {Eq.~\eqref{eq:ch3.eigenfield-VSH-def}};
  \node[ch1box, below=10mm of mn] (j) {Eq.~\eqref{eq:ch3.J-simultaneous-eigenfield-system}};
  \draw[ch1flow] (d) -- (mn) -- (ef);
  \draw[ch1flow] (j) -- (ef);
\end{tikzpicture}%
}
\caption{Subsection~\ref{sec:ch344} construction pipeline: separated Debye potentials generate \(\mathbf{M},\mathbf{N}\),
then Maxwell pairs become \(\mathbf{J}\) eigenfields Eq.~\eqref{eq:ch3.eigenfield-VSH-def}.}
\label{fig:ch3.344-debye-eigenfield-pipeline}
\end{figure}

\paragraph{Worked illustration (magnetic-dipole eigenfield).}
\(\mathbf{u}_{1,-1}^{(\mathbf{M})}\) is the canonical magnetic-dipole eigenfield partner to electric
\(\mathbf{u}_{10}^{(\mathbf{N})}\) at \(\ell=1\) \cite{stratton1941,jackson1999}. Current loops with dominant azimuthal
\(\Jimp\) project strongly onto \(\mathbf{M}_{1,m}\) in Eq.~\eqref{eq:ch3.eigenfield-source-projection}
\cite{harrington2001,tai1994}. Declaring ``OAM'' from \(\phi\) slope without \(\sigma\) resolution conflates orbital
and polarization sectors \cite{belinfante1940,hansen1988}.

\paragraph{Accessibility filter (Chapter~\ref{ch:02} inheritance).}
Not every label \((\ell,m,\sigma)\) is accessible from a given source \(\mathfrak{D}\): symmetry and current topology
from Chapter~\ref{ch:02} null certain projections in Eq.~\eqref{eq:ch3.eigenfield-source-projection}
\cite{harrington2001,stratton1941}. Subsection~\ref{sec:ch344} therefore couples representation labels to source
maps \(\mathcal{D}(\mathbf{J})\), not to abstract completeness alone \cite{harrington2001}.

\paragraph{Problem (simultaneous eigenfield checklist).}
\textbf{Problem.} Decide whether a measured far-zone pattern qualifies as \(\mathbf{u}_{21}^{(\mathbf{N})}\)
\cite{hansen1988,stratton1941,harrington2001}.

\textbf{Step 1.} Test outgoing class Eq.~\eqref{eq:ch3.outgoing-inheritance-criterion} \cite{stratton1941}.

\textbf{Step 2.} Match \(\ell=2\) angular nodes of \(Y_{2}^{1}\) and verify \(\mathbf{J}_{z}\) eigenvalue \(m=1\)
\cite{jackson1999}.

\textbf{Step 3.} Resolve \(\sigma=\mathbf{N}\) with polarization observables Eq.~\eqref{eq:ch3.polarization-sector-projector}
\cite{belinfante1940}.

\textbf{Step 4.} Confirm coefficient stability under \(\Pi_{\mathrm{rad}}\) \cite{harrington2001}.

\textbf{Physical meaning.} Eigenfield identification is a joint label test on \((\ell,m,\sigma)\), outgoing class, and
export projection \cite{hansen1988,stratton1941}.

\paragraph{Radiation-operator alignment (preview).}
When the radiation operator \(\mathcal{R}_\omega\) acts on sources, accessible eigenfield coefficients are those surviving
\(\Pi_{\mathrm{rad}}\) and Eq.~\eqref{eq:ch3.eigenfield-source-projection} \cite{harrington2001,stratton1941}. Subsection~\ref{sec:ch344}
therefore closes the loop between operator-domain Chapter~\ref{ch:02} previews and explicit harmonic templates
\cite{hansen1988}.

\paragraph{Simultaneous diagonalization statement (expanded).}
Theorem~\ref{thm:ch3.vsh-J-eigenfields} packages three certificates: (i) \(\mathbf{J}^{2}\) eigenvalue
Eq.~\eqref{eq:ch3.eigenfield-Jsq-eigenvalue}; (ii) \(\mathbf{J}_{z}\) eigenvalue \(m\) from
Eq.~\eqref{eq:ch3.Ylm-Condon-Shortley-def}; (iii) polarization sector \(\sigma\) from
Eq.~\eqref{eq:ch3.polarization-sector-projector} \cite{jackson1999,belinfante1940,hansen1988}. Failure of any one
certificate means the field is a superposition of eigenfields, not a single label \cite{harrington2001,stratton1941}.

\paragraph{Source symmetry nulling (illustration).}
Axisymmetric current distributions null \(m\neq0\) projections in Eq.~\eqref{eq:ch3.eigenfield-source-projection}
\cite{harrington2001,jackson1999}. A \(\hat{\mathbf{z}}\)-directed dipole feed therefore accesses primarily
\(\mathbf{u}_{10}^{(\mathbf{N})}\), not \(\mathbf{u}_{11}^{(\mathbf{M})}\), even if a plot suggests azimuthal variation
from measurement noise \cite{stratton1941,hansen1988}.

\paragraph{Problem (magnetic versus electric dipole eigenfield).}
\textbf{Problem.} A loop antenna is modeled. Which eigenfield label is primary?
\cite{harrington2001,jackson1999,stratton1941}.

\textbf{Step 1.} Identify dominant \(\mathbf{M}\)-weighted coupling in Eq.~\eqref{eq:ch3.eigenfield-source-projection}
\cite{harrington2001}.

\textbf{Step 2.} Select \(\ell=1\), \(|m|=1\) \cite{jackson1999}.

\textbf{Step 3.} Set \(\sigma=\mathbf{M}\) via Eq.~\eqref{eq:ch3.polarization-sector-projector} \cite{belinfante1940}.

\textbf{Physical meaning.} Feed topology fixes \(\sigma\) before far-zone normalization \cite{harrington2001}.

\paragraph{Quantum-number bookkeeping.}
Eq.~\eqref{eq:ch3.eigenfield-quantum-number-split} imported from Subsection~\ref{sec:ch322} partitions energy export
into orbital \((\ell,m)\) and polarization \(\sigma\) \cite{jackson1999,belinfante1940}. Subsection~\ref{sec:ch344}
attaches explicit fields to each entry of that partition through Eq.~\eqref{eq:ch3.eigenfield-VSH-def}
\cite{hansen1988,stratton1941}.

\paragraph{Near-zone caveat for eigenfield naming.}
Fields near sources contain reactive content not in \(\mathcal{F}_{\mathrm{rad}}\) \cite{harrington2001,stratton1941}.
Eigenfield labels apply to \(\mathcal{R}_\omega[\mathfrak{D}]\) after \(\Pi_{\mathrm{rad}}\), not to raw near-zone
probes \cite{harrington2001}. Naming a near-field pattern ``\(\mathbf{u}_{\ell m}^{(\sigma)}\)'' without export
qualification mixes representation language with non-radiative content \cite{stratton1941,hansen1988}.

\begin{table}[t]
\centering
\caption{Eigenfield labels in Subsection~\ref{sec:ch344}.}
\label{tab:ch3.344-eigenfield-labels}
\small
\begin{tabular}{@{}p{0.14\linewidth}p{0.30\linewidth}p{0.40\linewidth}@{}}
\toprule
Label & Operator & Constructed object \\
\midrule
\(\ell,m\) & Eq.~\eqref{eq:ch3.J-simultaneous-eigenfield-system} & Eqs.~\eqref{eq:ch3.Mlm-def}--\eqref{eq:ch3.Nlm-def} \\
\(\sigma\) & Eq.~\eqref{eq:ch3.polarization-sector-projector} & Eq.~\eqref{eq:ch3.eigenfield-VSH-def} \\
\bottomrule
\end{tabular}
\end{table}

\paragraph{Full \(\mathbf{J}\) eigenfield checklist (measurement side).}
A far-zone sample is certified as \(\mathbf{u}_{\ell m}^{(\sigma)}\) only if:
(i) outgoing class Eq.~\eqref{eq:ch3.outgoing-inheritance-criterion};
(ii) \(\ell(\ell+1)\) pattern nodes from Eq.~\eqref{eq:ch3.eigenfield-Jsq-eigenvalue};
(iii) \(m\) from \(\phi\) phase Eq.~\eqref{eq:ch3.Ylm-Condon-Shortley-def};
(iv) \(\sigma\) from Eq.~\eqref{eq:ch3.polarization-sector-projector} \cite{jackson1999,belinfante1940,hansen1988}.
Subsection~\ref{sec:ch344} converts operator algebra from Subsection~\ref{sec:ch322} into this laboratory checklist
\cite{harrington2001,stratton1941}.

\paragraph{Worked illustration (electric quadrupole eigenfield).}
\(\mathbf{u}_{20}^{(\mathbf{N})}\) carries the axisymmetric quadrupole pattern with a \(Y_{2}^{0}\) node structure
\cite{jackson1999,stratton1941}. Offset dipole pairs excite \(\ell=2\) projections through
Eq.~\eqref{eq:ch3.eigenfield-source-projection} even when \(\ell=1\) feed was intended \cite{harrington2001}. Eigenfield
labeling exposes the \(\ell=2\) contamination explicitly \cite{hansen1988}.

\paragraph{Problem (project impressed current onto \(\mathbf{M}_{21}\)).}
\textbf{Problem.} A current distribution \(\Jimp\) is given. Outline extraction of the \(\mathbf{M}_{21}\) coefficient
\cite{tai1994,harrington2001,stratton1941}.

\textbf{Step 1.} Form source functional Eq.~\eqref{eq:ch3.eigenfield-source-projection} with \(\mathbf{F}_{21}^{(\mathbf{M})}
=\mathbf{M}_{21}\) \cite{tai1994}.

\textbf{Step 2.} Apply \(\Pi_{\mathrm{rad}}\) to the radiated field \cite{harrington2001}.

\textbf{Step 3.} Divide by normalization from Section~\ref{sec:ch35} \cite{stratton1941}.

\textbf{Physical meaning.} Source projection is linear in \(\Jimp\); nonlinear pattern fitting is not required
\cite{harrington2001,hansen1988}.

\paragraph{Ladder transport among \(m\) partners (operational).}
Eq.~\eqref{eq:ch3.J-ladder-action} moves energy among \(m\) partners without changing \(\ell\) or \(\sigma\)
\cite{jackson1999,hansen1988}. Subsection~\ref{sec:ch344} uses ladders conceptually; explicit Wigner matrices in
Subsection~\ref{sec:ch345} perform finite rotations \cite{stratton1941}.

\paragraph{Polarization sector resolution (expanded).}
Eq.~\eqref{eq:ch3.polarization-sector-projector} partitions a fixed \((\ell,m)\) pair into \(\sigma=\mathbf{M}\) and
\(\sigma=\mathbf{N}\) sectors \cite{belinfante1940,hansen1988}. Subsection~\ref{sec:ch344} identifies those sectors with
explicit fields Eq.~\eqref{eq:ch3.eigenfield-VSH-def}, not with circular polarization helicity labels from
Section~\ref{sec:ch33} \cite{belinfante1940,stratton1941}. Confusing \(\sigma\) with helicity is a catalog error
corrected by the projector import \cite{belinfante1940,harrington2001}.

\paragraph{Worked illustration (offset dipole pair \(\to\ell=2\) eigenfield).}
Two equal \(\mathbf{u}_{10}^{(\mathbf{N})}\) sources separated along \(\hat{\mathbf{x}}\) produce a primary
\(\mathbf{u}_{20}^{(\mathbf{N})}\) component through source superposition before rotation \cite{jackson1999,harrington2001}.
The \(\ell=2\) label is not optional curve fitting; it is fixed by \(\mathbf{J}^{2}\) eigenvalue
Eq.~\eqref{eq:ch3.eigenfield-Jsq-eigenvalue} \cite{hansen1988,stratton1941}.

\begin{figure}[t]
\centering
\ChOneFigBlock{%
\begin{tikzpicture}[ch1fig, node distance=7mm and 9mm]
  \node[ch1box] (src) {Eq.~\eqref{eq:ch3.eigenfield-source-projection}};
  \node[ch1box, right=12mm of src] (pi) {$\Pi_{\mathrm{rad}}$};
  \node[ch1box, right=12mm of pi] (ef) {Eq.~\eqref{eq:ch3.eigenfield-VSH-def}};
  \node[ch1box, below=10mm of pi] (j) {Eq.~\eqref{eq:ch3.J-simultaneous-eigenfield-system}};
  \draw[ch1flow] (src) -- (pi) -- (ef);
  \draw[ch1flow] (j) -- (ef);
\end{tikzpicture}%
}
\caption{Subsection~\ref{sec:ch344} accessibility chain: source projection, radiative filter, and simultaneous
\(\mathbf{J}\) eigenfield certification.}
\label{fig:ch3.344-accessibility-chain}
\end{figure}

\begin{table}[t]
\centering
\caption{Subsection~\ref{sec:ch344} eigenfield certification checklist.}
\label{tab:ch3.344-certification-checklist}
\small
\begin{tabular}{@{}p{0.20\linewidth}p{0.38\linewidth}p{0.34\linewidth}@{}}
\toprule
Test & Pass object & Failure interpretation \\
\midrule
Outgoing & Eq.~\eqref{eq:ch3.outgoing-inheritance-criterion} & reactive mix \\
\(\mathbf{J}^{2}\) & Eq.~\eqref{eq:ch3.eigenfield-Jsq-eigenvalue} & wrong \(\ell\) \\
\(\mathbf{J}_{z}\) & \(m\) from Eq.~\eqref{eq:ch3.Ylm-Condon-Shortley-def} & wrong \(m\) \\
\(\sigma\) & Eq.~\eqref{eq:ch3.polarization-sector-projector} & \(\mathbf{M}/\mathbf{N}\) swap \\
\bottomrule
\end{tabular}
\end{table}

\paragraph{Theorem~\ref{thm:ch3.vsh-J-eigenfields} operational restatement.}
The theorem is the bridge statement between separation algebra (Section~\ref{sec:ch33}) and measurable multiplet
labels \cite{hansen1988,tai1994,jackson1999}. In practice, a certified eigenfield is any outgoing field whose
projection onto every non-matching label vanishes under Eq.~\eqref{eq:ch3.VSH-orthogonality-radiative}
\cite{stratton1941,harrington2001}. Subsection~\ref{sec:ch344} therefore recommends label-null tests rather than
peak-matching alone \cite{harrington2001,hansen1988}.

\paragraph{Commutation inheritance (expanded).}
Simultaneous eigenfield system Eq.~\eqref{eq:ch3.J-simultaneous-eigenfield-system} is compatible with Maxwell linearity:
finite superpositions of eigenfields are not eigenfields unless they share the same \((\ell,m,\sigma)\) label
\cite{jackson1999,hansen1988,stratton1941}. Subsection~\ref{sec:ch344} therefore distinguishes \emph{single-label}
eigenfields from \emph{multiplet superpositions} that require separate coefficients in
Eq.~\eqref{eq:ch3.VSH-resolution-identity} \cite{harrington2001}.

\paragraph{Problem (reject a mislabeled \(\ell=3\) claim).}
\textbf{Problem.} A pattern is claimed to be pure \(\mathbf{u}_{30}^{(\mathbf{N})}\) but shows only four azimuthal
nodes. Decide the claim status \cite{jackson1999,hansen1988,harrington2001}.

\textbf{Step 1.} Count expected nodes from \(Y_{3}^{0}\) structure \cite{jackson1999}.

\textbf{Step 2.} Project onto competing \(\ell=2,4\) partners \cite{hansen1988}.

\textbf{Step 3.} If projections are nonzero, reject single-label certification \cite{stratton1941}.

\textbf{Physical meaning.} Node counting is a \(\mathbf{J}^{2}\) diagnostic complementary to coefficient projection
\cite{harrington2001}.

\paragraph{Debye potential special cases (reference map).}
Table~\ref{tab:ch3.334-debye-roles} lists \(\Pi,\Psi\) roles from Subsection~\ref{sec:ch334}; Subsection~\ref{sec:ch344}
assigns eigenfield names to those roles through Eqs.~\eqref{eq:ch3.Mlm-def}--\eqref{eq:ch3.eigenfield-VSH-def}
\cite{tai1994,stratton1941}. Pure \(\Pi\) with \(\Psi=0\) generates \(\mathbf{M}\)-type dominance; pure \(\Psi\) generates
\(\mathbf{N}\)-type dominance at matched \((\ell,m)\) \cite{tai1994,hansen1988}.

\paragraph{Worked illustration (line source versus loop source at \(\ell=1\)).}
A thin line current along \(\hat{\mathbf{z}}\) and a small loop in the \(xy\)-plane access different
\((m,\sigma)\) sectors at the same \(\ell=1\) through Eq.~\eqref{eq:ch3.eigenfield-source-projection}
\cite{harrington2001,jackson1999,stratton1941}. The eigenfield map makes the distinction quantitative rather than
descriptive \cite{hansen1988}.

\paragraph{Simultaneous eigenfield uniqueness (label non-degeneracy).}
Distinct labels \((\ell,m,\sigma)\) correspond to distinct simultaneous eigenfields under
Eq.~\eqref{eq:ch3.J-simultaneous-eigenfield-system} and Eq.~\eqref{eq:ch3.polarization-sector-projector}
\cite{jackson1999,belinfante1940,hansen1988}. Subsection~\ref{sec:ch344} therefore treats mislabeling as a hard
failure mode in calibration, not a soft fitting residual \cite{harrington2001,stratton1941}.

\paragraph{Export map closure (Chapter~\ref{ch:02} interface).}
Let \(\mathcal{R}_\omega[\mathfrak{D}]\in\mathcal{F}_{\mathrm{rad}}\). Its eigenfield coefficients are
\(\{c_{\ell m}^{(\sigma)}\}\) defined by Eq.~\eqref{eq:ch3.VSH-resolution-identity} with amplitudes from
Eq.~\eqref{eq:ch3.eigenfield-source-projection} \cite{harrington2001,stratton1941}. Subsection~\ref{sec:ch344} is the
first point where abstract radiation outputs from Chapter~\ref{ch:02} receive explicit harmonic names
\cite{hansen1988,tai1994}. Accessibility nulls from source symmetry appear as zero coefficients, not as missing equations
\cite{harrington2001}.

\paragraph{Worked illustration (triaxial \(\ell=1\) superposition).}
Any outgoing \(\ell=1\) field is a six-term superposition over \((m,\sigma)\) at most \cite{jackson1999,hansen1988}.
Measuring only three electric components without \(\sigma\) resolution undercounts magnetic-dipole partners
\cite{belinfante1940,harrington2001}. Subsection~\ref{sec:ch344} recommends joint \((\ell,m,\sigma)\) reporting in
multipole catalogs \cite{stratton1941,hansen1988}.

\begin{table}[t]
\centering
\caption{Subsection~\ref{sec:ch344} source-to-eigenfield coupling examples at \(\ell=1\).}
\label{tab:ch3.344-source-coupling-examples}
\small
\begin{tabular}{@{}p{0.30\linewidth}p{0.32\linewidth}p{0.30\linewidth}@{}}
\toprule
Source model & Dominant label & Projection law \\
\midrule
Short \(\hat{\mathbf{z}}\) dipole & \(\mathbf{u}_{10}^{(\mathbf{N})}\) & Eq.~\eqref{eq:ch3.eigenfield-source-projection} \\
\(xy\) loop & \(\mathbf{u}_{1,\pm1}^{(\mathbf{M})}\) & Eq.~\eqref{eq:ch3.eigenfield-source-projection} \\
Crossed dipoles & Mixed \(m\) at \(\ell=1\) & Eq.~\eqref{eq:ch3.VSH-resolution-identity} \\
\bottomrule
\end{tabular}
\end{table}

\paragraph{Problem (full eigenfield certification for \(\ell=2,m=1\)).}
\textbf{Problem.} Certify whether an outgoing far-zone sample is \(\mathbf{u}_{21}^{(\mathbf{N})}\)
\cite{hansen1988,stratton1941,harrington2001,jackson1999}.

\textbf{Step 1.} Verify Eq.~\eqref{eq:ch3.outgoing-inheritance-criterion} on the measured dataset \cite{stratton1941}.

\textbf{Step 2.} Count angular nodes expected from \(Y_{2}^{1}\) and compare to data \cite{jackson1999}.

\textbf{Step 3.} Test \(\mathbf{J}_{z}\) eigenvalue \(m=1\) by \(\phi\) phase progression \cite{hansen1988}.

\textbf{Step 4.} Resolve \(\sigma=\mathbf{N}\) using Eq.~\eqref{eq:ch3.polarization-sector-projector} \cite{belinfante1940}.

\textbf{Step 5.} Null-test projections on \(\mathbf{u}_{10}^{(\mathbf{N})}\), \(\mathbf{u}_{20}^{(\mathbf{N})}\),
\(\mathbf{u}_{21}^{(\mathbf{M})}\) via Theorem~\ref{thm:ch3.vsh-orthogonality} \cite{hansen1988,stratton1941}.

\textbf{Physical meaning.} Certification is a sequence of label-null tests; passing one dipole check is insufficient
\cite{harrington2001}.

\begin{figure}[t]
\centering
\ChOneFigBlock{%
\begin{tikzpicture}[ch1fig, node distance=7mm and 8mm]
  \node[ch1box] (out) {outgoing test};
  \node[ch1box, right=10mm of out] (jz) {$\mathbf{J}_{z}$};
  \node[ch1box, right=10mm of jz] (jsq) {$\mathbf{J}^{2}$};
  \node[ch1box, right=10mm of jsq] (sig) {$\sigma$};
  \node[ch1box, right=10mm of sig] (null) {null tests};
  \draw[ch1flow] (out) -- (jz) -- (jsq) -- (sig) -- (null);
\end{tikzpicture}%
}
\caption{Subsection~\ref{sec:ch344} eigenfield certification flow for a single \((\ell,m,\sigma)\) claim.}
\label{fig:ch3.344-certification-flow}
\end{figure}

Subsection~\ref{sec:ch344} has constructed \(\mathbf{u}_{\ell m}^{(\sigma)}\) as \(\mathbf{J}\) eigenfields from VSH
and Debye data \cite{hansen1988,tai1994,jackson1999}. Subsection~\ref{sec:ch345} analyzes degeneracy and coupling
under rotations \cite{stratton1941}.


\subsection{Degeneracy and representation coupling}
\label{sec:ch345}
% TOC tag: [HOME][USE SS2.6.2]

Fixed \(\ell\) carries \(2\ell+1\) magnetic sublevels \(m\). Subsection~\ref{sec:ch345} records how rotations mix
those sublevels and how coupling coefficients weight superpositions \cite{jackson1999,hansen1988,stratton1941}.

\paragraph{Standing imports.}
Radiation rotation commutation Eq.~\eqref{eq:ch2.radiation-rotation-commutation}, field rotation
Eq.~\eqref{eq:ch2.field-rotation-representation}, and degeneracy preview from Subsection~\ref{sec:ch262} enter by
reference \cite{stratton1941,harrington2001}. Wigner matrices
Eqs.~\eqref{eq:ch3.Wigner-D-coefficient-law}--\eqref{eq:ch3.Wigner-D-unitarity} are imported from
Subsection~\ref{sec:ch324} \cite{hansen1988,jackson1999}.

\paragraph{\(2\ell+1\) degeneracy (HOME).}
For each \(\ell\), the set \(\{\mathbf{u}_{\ell m}^{(\sigma)}:m=-\ell,\ldots,\ell\}\) spans a
\((2\ell+1)\)-dimensional \(\mathbf{J}_{z}\) eigenspace at fixed \(\sigma\) \cite{jackson1999,hansen1988}. This
degeneracy matches the SO\((3)\) irrep dimension of Subsection~\ref{sec:ch323} \cite{stratton1941}.

\paragraph{Rotation action on coefficients (HOME).}
Under \(R\in\mathrm{SO}(3)\),
\begin{equation}
\label{eq:ch3.eigenfield-Wigner-mixing}
\mathcal{D}(R)\,\mathbf{u}_{\ell m}^{(\sigma)}
=
\sum_{m'=-\ell}^{\ell}
D_{m'm}^{\ell}(R)\,\mathbf{u}_{\ell m'}^{(\sigma)},
\end{equation}
with \(D_{m'm}^{\ell}\) from Eq.~\eqref{eq:ch3.Wigner-D-coefficient-law} \cite{hansen1988,jackson1999}. Equation
\eqref{eq:ch3.eigenfield-Wigner-mixing} is the Subsection~\ref{sec:ch345} eigenfield restatement of
Eq.~\eqref{eq:ch2.radiation-rotation-commutation} in harmonic coordinates \cite{stratton1941,harrington2001}.

\paragraph{Coupling to superpositions (USE).}
Clebsch--Gordan weights Eq.~\eqref{eq:ch3.CG-coupling-def} from Subsection~\ref{sec:ch324} govern products of
angular sectors when composite labels are formed \cite{jackson1999,hansen1988}. Subsection~\ref{sec:ch345} cites
those weights without re-deriving CG theory \cite{hansen1988}.

\begin{figure}[t]
\centering
\ChOneFigBlock{%
\begin{tikzpicture}[ch1fig, node distance=7mm and 9mm]
  \node[ch1box] (l) {fixed $\ell$};
  \node[ch1box, right=12mm of l] (m) {$2\ell+1$ values of $m$};
  \node[ch1box, right=12mm of m] (d) {Eq.~\eqref{eq:ch3.eigenfield-Wigner-mixing}};
  \node[ch1box, below=10mm of m] (r) {Eq.~\eqref{eq:ch2.radiation-rotation-commutation}};
  \draw[ch1flow] (l) -- (m) -- (d);
  \draw[ch1flow] (r) -- (d);
\end{tikzpicture}%
}
\caption{Degeneracy and rotation mixing in Subsection~\ref{sec:ch345}: fixed \(\ell\) blocks rotate among
\(m\) partners through Wigner coefficients.}
\label{fig:ch3.degeneracy-wigner-mixing}
\end{figure}

\paragraph{Worked illustration (rotate a dipole pattern).}
Rotating a \(\mathbf{u}_{10}^{(\mathbf{N})}\) pattern mixes \(m=0\) into \(m=\pm1\) partners with weights
\(D_{m'0}^{1}(R)\) \cite{jackson1999,harrington2001}. The \(\ell=1\) block remains closed; only coefficients change
\cite{stratton1941,hansen1988}.

\paragraph{Problem (count modes at \(\ell=3\)).}
\textbf{Problem.} How many \(m\) partners appear at \(\ell=3\) for each \(\sigma\) sector?
\cite{jackson1999,hansen1988}.

\textbf{Step 1.} Use \(m\in\{-\ell,\ldots,\ell\}\) \cite{jackson1999}.

\textbf{Step 2.} Count \(2\ell+1=7\) \cite{hansen1988}.

\textbf{Step 3.} Repeat for \(\sigma=\mathbf{M}\) and \(\mathbf{N}\) independently \cite{stratton1941}.

\textbf{Physical meaning.} Mode counting is representation dimension, not a fitted parameter count
\cite{harrington2001,stratton1941}.

\paragraph{Explicit \(\ell=1\) rotation block.}
For \(\ell=1\), Eq.~\eqref{eq:ch3.eigenfield-Wigner-mixing} is a \(3\times3\) unitary on
\(\{\mathbf{u}_{1,-1}^{(\sigma)},\mathbf{u}_{10}^{(\sigma)},\mathbf{u}_{11}^{(\sigma)}\}\) with matrix elements
\(D_{m'm}^{1}(R)\) from Eq.~\eqref{eq:ch3.Wigner-D-coefficient-law} \cite{jackson1999,hansen1988}. A \(\pi/2\) rotation
about \(\hat{\mathbf{y}}\) mixes \(m=0\) into \(m=\pm1\) with weights fixed by Euler angles, not by fitted beam
parameters \cite{stratton1941,harrington2001}. Laboratory frame changes are linear mixing inside the \(\ell=1\) block.

\paragraph{Interference versus degeneracy mixing.}
Superposing \(\mathbf{u}_{\ell m}^{(\sigma)}\) with distinct \(m\) before rotation produces interference in
\(\mathbf{E}\) intensity; after rotation, the same physical state is described by new coefficients
Eq.~\eqref{eq:ch3.eigenfield-Wigner-mixing} \cite{jackson1999,hansen1988}. Subsection~\ref{sec:ch345} separates
\emph{representation transport} (Wigner mixing) from \emph{power budgets} (orthogonality of distinct labels after
\(\Pi_{\mathrm{rad}}\)) \cite{harrington2001,stratton1941}.

\paragraph{Body frame versus laboratory frame.}
Let \(R_{\mathrm{lab\leftarrow body}}\) rotate a source-fixed multiplet into laboratory coordinates. Coefficients
transform by Eq.~\eqref{eq:ch3.eigenfield-Wigner-mixing} while \(\ell\) and \(\sigma\) remain invariant
\cite{stratton1941,harrington2001}. Misidentifying frame rotation as a change in \(\ell\) is a common catalog error in
pattern fitting \cite{harrington2001,hansen1988}.

\paragraph{Addition-theorem pointer (no re-derivation).}
Products of scalar harmonics decompose into CG series Eq.~\eqref{eq:ch3.CG-coupling-def}; vector products inherit those
weights through \(\mathbf{M},\mathbf{N}\) bilinears \cite{jackson1999,hansen1988}. Subsection~\ref{sec:ch345} cites the
coupling law for composite labels without reopening addition-theorem proofs deferred to Subsection~\ref{sec:ch324}
\cite{hansen1988}.

\paragraph{Worked illustration (rotated \(\ell=2\) quadrupole).}
A body-fixed \(\mathbf{u}_{20}^{(\mathbf{N})}\) rotated about \(\hat{\mathbf{z}}\) remains inside the \(\ell=2\)
\(\sigma=\mathbf{N}\) block with redistributed \(m\) weights \cite{jackson1999,stratton1941}. Radiated power summed over
orthogonal partners is invariant under \(R\) if lossless reciprocity holds \cite{stratton1941,harrington2001}. What
changes is the angular phase allocation among \(m\) channels, not the total \(\ell=2\) export budget.

\paragraph{Problem (construct rotated dipole coefficients).}
\textbf{Problem.} A \(\mathbf{u}_{10}^{(\mathbf{N})}\) state is rotated by \(R_{z}(\pi/2)\). Express new coefficients
\(c_{1m'}^{(\mathbf{N})}\) \cite{jackson1999,hansen1988}.

\textbf{Step 1.} Write \(\mathcal{D}(R)\mathbf{u}_{10}^{(\mathbf{N})}=\sum_{m'}D_{m'0}^{1}(R)\mathbf{u}_{1m'}^{(\mathbf{N})}\)
from Eq.~\eqref{eq:ch3.eigenfield-Wigner-mixing} \cite{hansen1988}.

\textbf{Step 2.} Insert \(D_{m'0}^{1}\) for \(R_{z}(\pi/2)\) from Eq.~\eqref{eq:ch3.Wigner-D-coefficient-law}
\cite{jackson1999}.

\textbf{Step 3.} Read off nonzero \(m'\) partners and verify unitarity
Eq.~\eqref{eq:ch3.Wigner-D-unitarity} \cite{hansen1988}.

\textbf{Physical meaning.} Rotation redistributes fixed \(\ell=1\) power among \(m\) eigenchannels \cite{harrington2001}.

\paragraph{Ensemble degeneracy in measurements.}
An ensemble of orientations averages Wigner weights rather than individual \(m\) labels \cite{jackson1999,harrington2001}.
Subsection~\ref{sec:ch345} therefore distinguishes single-orientation eigenfield transport from orientation-averaged
pattern statistics \cite{stratton1941,hansen1988}.

\begin{figure}[t]
\centering
\ChOneFigBlock{%
\begin{tikzpicture}[ch1fig, node distance=7mm and 9mm]
  \node[ch1box] (bf) {body frame\\$\mathbf{u}_{\ell m}^{(\sigma)}$};
  \node[ch1box, right=13mm of bf] (w) {Eq.~\eqref{eq:ch3.eigenfield-Wigner-mixing}};
  \node[ch1box, right=13mm of w] (lf) {laboratory\\coefficients $c_{\ell m'}^{(\sigma)}$};
  \node[ch1box, below=10mm of w] (cg) {Eq.~\eqref{eq:ch3.CG-coupling-def}};
  \draw[ch1flow] (bf) -- (w) -- (lf);
  \draw[ch1flow] (cg) -- (w);
\end{tikzpicture}%
}
\caption{Frame transport in Subsection~\ref{sec:ch345}: rotations mix \(m\) within fixed \((\ell,\sigma)\); composite
products invoke CG weights Eq.~\eqref{eq:ch3.CG-coupling-def}.}
\label{fig:ch3.345-frame-transport}
\end{figure}

\paragraph{Commutation with radiation map (imported).}
Eq.~\eqref{eq:ch2.radiation-rotation-commutation} states that rotating sources and rotating fields commute with
\(\mathcal{R}_\omega\) in the admissible class \cite{stratton1941,harrington2001}. Equation
\eqref{eq:ch3.eigenfield-Wigner-mixing} is the harmonic-coordinate display of that commutation on
\(\mathbf{u}_{\ell m}^{(\sigma)}\) \cite{hansen1988,jackson1999}. Subsection~\ref{sec:ch345} therefore ties laboratory
rotations to coefficient transport, not to changes in exported operator spectrum \cite{stratton1941}.

\paragraph{Multiplicity accounting at fixed \(\ell\).}
For each \(\sigma\), degeneracy count \(2\ell+1\) matches the dimension of the SO\((3)\) irrep imported from
Subsection~\ref{sec:ch323} \cite{jackson1999,hansen1988}. Density-of-states refinements in Section~\ref{sec:ch353}
count those multiplets under normalization conventions \cite{stratton1941,harrington2001}. Subsection~\ref{sec:ch345}
supplies the degeneracy law before measure choices enter \cite{hansen1988}.

\paragraph{Worked illustration (two-source interference inside one \(\ell\) block).}
Superposing \(\mathbf{u}_{2,-1}^{(\mathbf{N})}\) and \(\mathbf{u}_{21}^{(\mathbf{N})}\) before rotation produces
non-trivial \(\phi\) interference; a subsequent rotation redistributes weights among all seven \(m\) partners via
Eq.~\eqref{eq:ch3.eigenfield-Wigner-mixing} \cite{jackson1999,harrington2001}. The \(\ell=2\), \(\sigma=\mathbf{N}\) block
remains closed; only coefficients change \cite{stratton1941,hansen1988}.

\paragraph{Problem (unitarity constraint on rotated coefficients).}
\textbf{Problem.} Show \(\sum_{m}|c_{1m}^{(\mathbf{N})}|^{2}\) is invariant under \(R\) for a pure \(\ell=1\) state
\cite{hansen1988,jackson1999}.

\textbf{Step 1.} Express rotated coefficients through Eq.~\eqref{eq:ch3.eigenfield-Wigner-mixing} \cite{hansen1988}.

\textbf{Step 2.} Apply Eq.~\eqref{eq:ch3.Wigner-D-unitarity} \cite{jackson1999}.

\textbf{Step 3.} Conclude coefficient norm invariance before flux normalization \cite{stratton1941}.

\textbf{Physical meaning.} Rotation preserves multiplet power allocation; it does not create new \(\ell\) sectors
\cite{harrington2001}.

\paragraph{CG coupling without reopening Subsection~\ref{sec:ch324}.}
When composite angular momentum labels are formed, weights Eq.~\eqref{eq:ch3.CG-coupling-def} select allowed
\((\ell_1,m_1)\otimes(\ell_2,m_2)\to(\ell_3,m_3)\) branches \cite{jackson1999,hansen1988}. Subsection~\ref{sec:ch345}
records the coupling interface for later product-state discussions in Chapter~\ref{ch:04} without re-deriving CG
algebra \cite{hansen1988}.

\begin{table}[t]
\centering
\caption{Degeneracy and coupling summary in Subsection~\ref{sec:ch345}.}
\label{tab:ch3.345-degeneracy-coupling}
\small
\begin{tabular}{@{}p{0.22\linewidth}p{0.34\linewidth}p{0.36\linewidth}@{}}
\toprule
Object & Multiplicity & Mixing law \\
\midrule
Fixed \(\ell\) & \(2\ell+1\) per \(\sigma\) & Eq.~\eqref{eq:ch3.eigenfield-Wigner-mixing} \\
Composite products & CG weights & Eq.~\eqref{eq:ch3.CG-coupling-def} \\
\bottomrule
\end{tabular}
\end{table}

\paragraph{Euler-angle parameterization (interface).}
Rotations \(R(\alpha,\beta,\gamma)\) factor Wigner matrices Eq.~\eqref{eq:ch3.Wigner-D-coefficient-law} used in
Eq.~\eqref{eq:ch3.eigenfield-Wigner-mixing} \cite{jackson1999,hansen1988}. Subsection~\ref{sec:ch345} cites the
parameterization without re-deriving Euler laws deferred to Subsection~\ref{sec:ch324} \cite{hansen1988}.

\paragraph{Worked illustration (align dipole to \(-\hat{\mathbf{z}}\)).}
A \(180^{\circ}\) rotation maps \(\mathbf{u}_{10}^{(\mathbf{N})}\) to the same \(\ell=1\), \(m=0\), \(\sigma=\mathbf{N}\)
label with a sign change in coefficients fixed by \(D_{00}^{1}\) \cite{jackson1999,stratton1941}. The multiplet remains
closed; only phase and overall sign adjust \cite{hansen1988,harrington2001}.

\paragraph{Problem (distinguish rotation from \(\ell\) mixing).}
\textbf{Problem.} A fitted pattern shows new sidelobes after mechanical rotation. Decide whether \(\ell\) changed
\cite{harrington2001,jackson1999,stratton1941}.

\textbf{Step 1.} Test \(\ell\) block closure with Eq.~\eqref{eq:ch3.eigenfield-Wigner-mixing} \cite{hansen1988}.

\textbf{Step 2.} Compare node count to \(Y_{\ell}^{m}\) zeros \cite{jackson1999}.

\textbf{Step 3.} If \(\ell\) is unchanged, attribute sidelobes to \(m\) redistribution only \cite{stratton1941}.

\textbf{Physical meaning.} Rotations mix \(m\); they do not create new \(\ell\) without nonlinear sources or broken
linearity \cite{harrington2001}.

\paragraph{Degeneracy versus accidental degeneracy.}
\((2\ell+1)\) degeneracy is representation-theoretic for SO\((3)\) \cite{jackson1999,hansen1988}. Accidental degeneracy
from engineered symmetries is analyzed in Section~\ref{sec:ch36} \cite{stratton1941,harrington2001}. Subsection~\ref{sec:ch345}
records only the free-space SO\((3)\) degeneracy law \cite{hansen1988}.

\paragraph{Density-of-states pointer (fixed \(\ell\) multiplet).}
At each \(\ell\), \(2\ell+1\) magnetic sublevels per \(\sigma\) contribute to angular mode counting developed in
Section~\ref{sec:ch353} \cite{stratton1941,hansen1988}. Subsection~\ref{sec:ch345} supplies the degeneracy factor
before measure normalization; Section~\ref{sec:ch35}--\ref{sec:ch353} attach flux and spectral density conventions
\cite{poynting1884,harrington2001}.

\paragraph{Worked illustration (sequential rotations do not enlarge \(\ell\)).}
Compose \(R_{z}(\pi/2)R_{y}(\pi/2)\) on a pure \(\ell=1\) state \cite{jackson1999,hansen1988}. The product Wigner
action remains within the \(\ell=1\) block because each factor respects Eq.~\eqref{eq:ch3.eigenfield-Wigner-mixing}
\cite{hansen1988,stratton1941}. Laboratory alignment sequences therefore redistribute \(m\) weights without creating
\(\ell=2\) content unless sources are nonlinear or reconfigured \cite{harrington2001}.

\begin{table}[t]
\centering
\caption{Subsection~\ref{sec:ch345} rotation bookkeeping on harmonic eigenfields.}
\label{tab:ch3.345-rotation-bookkeeping}
\small
\begin{tabular}{@{}p{0.26\linewidth}p{0.34\linewidth}p{0.32\linewidth}@{}}
\toprule
Operation & Invariant labels & Changing object \\
\midrule
\(R\in\mathrm{SO}(3)\) & \(\ell,\sigma\) & coefficients \(c_{\ell m}^{(\sigma)}\) \\
CG product & selection rules & composite \(\ell\) \\
Ensemble average & total \(\ell\) block power & individual \(m\) phases \\
\bottomrule
\end{tabular}
\end{table}

\paragraph{Wigner unitarity and power invariance (expanded).}
Equation~\eqref{eq:ch3.Wigner-D-unitarity} implies \(\sum_{m}|D_{m'm}^{\ell}|^{2}=1\) for each column \(m\)
\cite{jackson1999,hansen1988}. Combined with Eq.~\eqref{eq:ch3.eigenfield-Wigner-mixing}, coefficient norms in a
pure \(\ell\) multiplet are rotation-invariant before flux normalization \cite{stratton1941,harrington2001}. Subsection~\ref{sec:ch345}
therefore supplies the symmetry proof target used when rotating calibrated standards between laboratory frames
\cite{hansen1988,harrington2001}.

\paragraph{Rotation operator on coefficient vectors (matrix form).}
For fixed \((\ell,\sigma)\), define coefficient vector \(\mathbf{c}_\ell=(c_{\ell,-\ell}^{(\sigma)},\ldots,c_{\ell\ell}^{(\sigma)})^{T}\).
Then Eq.~\eqref{eq:ch3.eigenfield-Wigner-mixing} is \(\mathbf{c}_\ell\mapsto D^{\ell}(R)\mathbf{c}_\ell\) with unitary
\(D^{\ell}(R)\) \cite{jackson1999,hansen1988}. Subsection~\ref{sec:ch345} is the coordinate restatement of
Eq.~\eqref{eq:ch2.radiation-rotation-commutation} in the harmonic basis of Section~\ref{sec:ch34}
\cite{stratton1941,harrington2001}.

\paragraph{Problem (verify closure of \(\ell=2\) block under \(R_{z}\)).}
\textbf{Problem.} Show \(R_{z}(\alpha)\) maps \(\mathbf{u}_{2m}^{(\mathbf{N})}\) into span\(\{\mathbf{u}_{2m'}^{(\mathbf{N})}\}\)
\cite{jackson1999,hansen1988}.

\textbf{Step 1.} Write Eq.~\eqref{eq:ch3.eigenfield-Wigner-mixing} for \(\ell=2\) \cite{hansen1988}.

\textbf{Step 2.} Use \(D_{m'm}^{2}(R_{z})\propto e^{-\jj m'\alpha}\delta_{mm'}\) \cite{jackson1999}.

\textbf{Step 3.} Conclude no \(\ell\neq2\) admixture \cite{stratton1941}.

\textbf{Physical meaning.} \(R_{z}\) only phases \(m\) partners within a fixed \(\ell\) block \cite{harrington2001}.

\paragraph{Coupled product states (pointer only).}
When two eigenfields interact nonlinearly in media or circuits, CG weights Eq.~\eqref{eq:ch3.CG-coupling-def} govern
allowed composite \(\ell\) labels \cite{jackson1999,hansen1988}. Subsection~\ref{sec:ch345} records the linear rotation
law; nonlinear coupling is outside the Maxwell linear scope of Section~\ref{sec:ch34} \cite{stratton1941,harrington2001}.

\paragraph{Laboratory calibration rotation protocol (summary).}
Rotate a known \(\mathbf{u}_{\ell m}^{(\sigma)}\) standard through \(R\), measure redistributed coefficients, and
compare to Eq.~\eqref{eq:ch3.eigenfield-Wigner-mixing} \cite{harrington2001,jackson1999,hansen1988}. Agreement validates
frame transport; disagreement signals mislabeling, non-outgoing contamination, or broken reciprocity
\cite{stratton1941,harrington2001}. Subsection~\ref{sec:ch345} closes the representation loop opened in
Subsection~\ref{sec:ch262} \cite{stratton1941}.

\paragraph{Worked illustration (compare \(D_{m'0}^{1}\) weights for two Euler choices).}
For \(\ell=1\), \(m=0\) standards, two rotation sequences with the same net \(R\) must yield identical coefficient
vectors \(\mathbf{c}_1\) \cite{jackson1999,hansen1988}. Numerical disagreements usually trace to inconsistent
Condon--Shortley conventions or to near-zone fields not filtered by \(\Pi_{\mathrm{rad}}\) \cite{hansen1988,harrington2001}.
Subsection~\ref{sec:ch345} therefore pairs Wigner transport with the phase contract of
Eq.~\eqref{eq:ch3.Ylm-Condon-Shortley-def} \cite{hansen1988,stratton1941}.

\begin{figure}[t]
\centering
\ChOneFigBlock{%
\begin{tikzpicture}[ch1fig, node distance=7mm and 9mm]
  \node[ch1box] (std) {calibrated $\mathbf{u}_{\ell m}^{(\sigma)}$};
  \node[ch1box, right=12mm of std] (rot) {$R$};
  \node[ch1box, right=12mm of rot] (w) {Eq.~\eqref{eq:ch3.eigenfield-Wigner-mixing}};
  \node[ch1box, right=12mm of w] (cmp) {compare $\mathbf{c}_\ell$};
  \draw[ch1flow] (std) -- (rot) -- (w) -- (cmp);
\end{tikzpicture}%
}
\caption{Subsection~\ref{sec:ch345} laboratory rotation check: Wigner mixing predicts coefficient transport for fixed
\((\ell,\sigma)\).}
\label{fig:ch3.345-calibration-rotation}
\end{figure}

\paragraph{Block-diagonal structure in \((\ell,\sigma)\) coordinates.}
Equation~\eqref{eq:ch3.eigenfield-Wigner-mixing} is block-diagonal in \(\ell\) and \(\sigma\): rotations never couple
\(\ell=1\) to \(\ell=2\) or \(\sigma=\mathbf{M}\) to \(\sigma=\mathbf{N}\) in the linear Maxwell regime
\cite{jackson1999,hansen1988,stratton1941}. Subsection~\ref{sec:ch345} makes that block structure explicit for
laboratory frame changes and for software implementations that store coefficients as vectors \(\mathbf{c}_{\ell}^{(\sigma)}\)
\cite{harrington2001,hansen1988}.

\paragraph{Problem (predict coefficient sign under \(180^{\circ}\) rotation).}
\textbf{Problem.} A \(\mathbf{u}_{10}^{(\mathbf{N})}\) standard is rotated by \(R_{z}(\pi)\). Predict the sign of the
returned \(m=0\) coefficient \cite{jackson1999,hansen1988,stratton1941}.

\textbf{Step 1.} Evaluate \(D_{00}^{1}(R_{z}(\pi))\) from Eq.~\eqref{eq:ch3.Wigner-D-coefficient-law} \cite{jackson1999}.

\textbf{Step 2.} Apply Eq.~\eqref{eq:ch3.eigenfield-Wigner-mixing} \cite{hansen1988}.

\textbf{Step 3.} Track Condon--Shortley sign in Eq.~\eqref{eq:ch3.Ylm-Condon-Shortley-def} \cite{hansen1988}.

\textbf{Physical meaning.} Global sign flips are gauge-level unless tied to flux-normalized amplitudes in
Section~\ref{sec:ch35} \cite{poynting1884,harrington2001}.

Subsection~\ref{sec:ch345} has classified \((2\ell+1)\) degeneracy and Wigner mixing on harmonic eigenfields
\cite{jackson1999,hansen1988,stratton1941}. Section~\ref{sec:ch34} is complete: definitions, structure,
orthogonality, eigenfield construction, and degeneracy are now in place for normalization in Section~\ref{sec:ch35}
\cite{stratton1941,hansen1988}.

The five subsections of Section~\ref{sec:ch34} supply the named harmonic basis of Volume~I:
\(Y_{\ell}^{m}\), \(\mathbf{M}_{\ell m}\), \(\mathbf{N}_{\ell m}\) in Subsection~\ref{sec:ch341}; toroidal and
poloidal structure in Subsection~\ref{sec:ch342}; orthogonality and resolution
Eqs.~\eqref{eq:ch3.VSH-orthogonality-radiative}--\eqref{eq:ch3.VSH-resolution-identity} in Subsection~\ref{sec:ch343};
eigenfields \(\mathbf{u}_{\ell m}^{(\sigma)}\) in Subsection~\ref{sec:ch344}; and degeneracy mixing
Eq.~\eqref{eq:ch3.eigenfield-Wigner-mixing} in Subsection~\ref{sec:ch345} \cite{hansen1988,tai1994,stratton1941}.
Section~\ref{sec:ch35} fixes flux and energy normalizations that turn coefficients into measurable amplitudes
\cite{stratton1941,poynting1884,harrington2001}.

\section{Normalization and Spectral Density}
\label{sec:ch35}

A mode label without a normalization convention is not yet an observable amplitude. Section~\ref{sec:ch35}
fixes energy and flux normalizations that connect eigenfield coefficients to Poynting flux integrals from
Subsection~\ref{sec:ch253} \cite{stratton1941,poynting1884}. Subsection~\ref{sec:ch351} defines energy and
flux normalization on \(\mathcal{F}_{\mathrm{rad}}\). Subsection~\ref{sec:ch352} treats continuous-spectrum
normalization measures, completing the rigged-Hilbert program of Subsection~\ref{sec:ch313}. Subsection~\ref{sec:ch353}
develops density of states and multiplicity counting. Subsection~\ref{sec:ch354} records how normalization
choices propagate into coupling coefficients and far-zone power budgets \cite{harrington2001,jackson1999}.


\subsection{Energy and flux normalization}
\label{sec:ch351}
% TOC tag: [HOME][USE SS2.5.3]

Section~\ref{sec:ch34} fixed harmonic \emph{shapes} and orthogonality certificates; coefficients in
Eq.~\eqref{eq:ch3.VSH-resolution-identity} remained geometric projections without watt calibration
\cite{stratton1941,hansen1988,harrington2001}. Subsection~\ref{sec:ch351} is the Volume~I HOME point where
outgoing VSH modes receive an energy--flux normalization tied to Chapter~\ref{ch:02} flux meters
\cite{poynting1884,stratton1941}.

\paragraph{Standing imports.}
Poynting functional Eq.~\eqref{eq:ch2.poynting-functional-radiative}, shell power
Eq.~\eqref{eq:ch2.power-flux-functional}, far-zone power pattern
Eq.~\eqref{eq:ch2.far-zone-power-from-pattern}, radiative inner product
Eq.~\eqref{eq:ch3.radiative-inner-product-def}, seminorm identity
Eq.~\eqref{eq:ch3.inner-product-seminorm-identity}, and VSH resolution
Eq.~\eqref{eq:ch3.VSH-resolution-identity} enter by reference \cite{stratton1941,harrington2001,poynting1884}.
Subsection~\ref{sec:ch351} does not re-derive Poynting balance or Sommerfeld asymptotics locked in
Chapter~\ref{ch:01} and Section~\ref{sec:ch24} \cite{stratton1941}.

\paragraph{Assumptions.}
Work on \(\mathcal{F}_{\mathrm{rad}}\) with outgoing inheritance
Eq.~\eqref{eq:ch3.outgoing-inheritance-criterion}, reciprocal lossless exterior media, and shell radii in the
far zone of Eq.~\eqref{eq:ch2.distance-regime-definitions} when flux integrals are evaluated
\cite{stratton1941,harrington2001,jackson1999}. Normalization below is declared on harmonic pairs
\(\mathbf{u}_{\ell m}^{(\sigma)}\) from Eq.~\eqref{eq:ch3.eigenfield-VSH-def} \cite{hansen1988}.

\paragraph{Unit-flux VSH normalization (HOME).}
For each admissible label \((\ell,m,\sigma)\), choose outgoing representatives
\(\widehat{\mathbf{u}}_{\ell m}^{(\sigma)}\) such that
\begin{equation}
\label{eq:ch3.VSH-unit-flux-normalization}
\mathcal{P}_{\mathrm{rad}}\!\bigl(r;\widehat{\mathbf{u}}_{\ell m}^{(\sigma)}\bigr)
=
1\ \mathrm{W},
\qquad
k_{0}r\gg\ell,
\end{equation}
with \(\mathcal{P}_{\mathrm{rad}}\) from Eq.~\eqref{eq:ch2.power-flux-functional} and fields identified with
\(\mathcal{R}_\omega\) images on the radiative branch \cite{stratton1941,poynting1884,harrington2001}.
Equation~\eqref{eq:ch3.VSH-unit-flux-normalization} is the Subsection~\ref{sec:ch351} HOME convention: one watt of
time-averaged outward power through far-zone shells per unit normalized mode \cite{harrington2001,stratton1941}.

\paragraph{Diagonal energy norm (HOME).}
Impose radiative energy normalization
\begin{equation}
\label{eq:ch3.VSH-unit-energy-norm}
\big\langle \widehat{\mathbf{u}}_{\ell m}^{(\sigma)},\,\widehat{\mathbf{u}}_{\ell m}^{(\sigma)}\big\rangle_{\mathrm{rad}}
=
1,
\end{equation}
compatible with Eq.~\eqref{eq:ch3.radiative-inner-product-def} on reciprocal media
\cite{stratton1941,hansen1988}. Equations~\eqref{eq:ch3.VSH-unit-flux-normalization}--\eqref{eq:ch3.VSH-unit-energy-norm}
fix amplitude scales for coefficients in Eq.~\eqref{eq:ch3.VSH-resolution-identity}:
\begin{equation}
\label{eq:ch3.normalized-coefficient-resolution}
\mathbf{v}_\omega
=
\sum_{\ell m\sigma}
\widehat{c}_{\ell m}^{(\sigma)}\,\widehat{\mathbf{u}}_{\ell m}^{(\sigma)},
\qquad
\widehat{c}_{\ell m}^{(\sigma)}
=
\big\langle \widehat{\mathbf{u}}_{\ell m}^{(\sigma)},\mathbf{v}_\omega\big\rangle_{\mathrm{rad}}.
\end{equation}
Equation~\eqref{eq:ch3.normalized-coefficient-resolution} upgrades the deferred parenthetical in
Eq.~\eqref{eq:ch3.VSH-resolution-identity} to a watt-ready expansion \cite{hansen1988,harrington2001}.

\paragraph{Mode power law (HOME).}
For normalized superpositions,
\begin{equation}
\label{eq:ch3.mode-power-from-coefficient}
\mathcal{P}_{\mathrm{rad}}\!\Bigl(r;\sum_{\ell m\sigma}\widehat{c}_{\ell m}^{(\sigma)}\widehat{\mathbf{u}}_{\ell m}^{(\sigma)}\Bigr)
=
\sum_{\ell m\sigma}
\big|\widehat{c}_{\ell m}^{(\sigma)}\big|^{2}
\;+\;
\mathcal{O}(r^{-1}),
\end{equation}
because distinct labels are orthogonal under Theorem~\ref{thm:ch3.vsh-orthogonality} and flux is quadratic on
\(\mathcal{F}_{\mathrm{rad}}\) in the far zone \cite{stratton1941,poynting1884,harrington2001}. Equation
\eqref{eq:ch3.mode-power-from-coefficient} is the Subsection~\ref{sec:ch351} Parseval law for measurable power:
\(|\widehat{c}|^{2}\) is watts per mode under Eqs.~\eqref{eq:ch3.VSH-unit-flux-normalization}--\eqref{eq:ch3.VSH-unit-energy-norm}
\cite{harrington2001}.

\begin{figure}[t]
\centering
\ChOneFigBlock{%
\begin{tikzpicture}[ch1fig, node distance=7mm and 9mm]
  \node[ch1box] (geom) {Eq.~\eqref{eq:ch3.VSH-resolution-identity}\\(geometry)};
  \node[ch1box, right=12mm of geom] (norm) {Eqs.~\eqref{eq:ch3.VSH-unit-flux-normalization}--\eqref{eq:ch3.VSH-unit-energy-norm}};
  \node[ch1box, right=12mm of norm] (pow) {Eq.~\eqref{eq:ch3.mode-power-from-coefficient}};
  \node[ch1box, below=10mm of norm] (flux) {Eq.~\eqref{eq:ch2.power-flux-functional}};
  \draw[ch1flow] (geom) -- (norm) -- (pow);
  \draw[ch1flow] (flux) -- (norm);
\end{tikzpicture}%
}
\caption{Subsection~\ref{sec:ch351} normalization chain: geometric projections become flux-calibrated coefficients through
Chapter~\ref{ch:02} power meters.}
\label{fig:ch3.351-flux-normalization-chain}
\end{figure}

\begin{theorem}[Normalized VSH diagonal flux]
\label{thm:ch3.vsh-unit-flux-diagonal}
Under Eqs.~\eqref{eq:ch3.VSH-unit-flux-normalization}--\eqref{eq:ch3.VSH-unit-energy-norm}, each normalized pair
\(\widehat{\mathbf{u}}_{\ell m}^{(\sigma)}\) carries unit outward power on far shells, and off-diagonal flux
cross-terms vanish for distinct labels \cite{stratton1941,hansen1988,harrington2001}.
\end{theorem}
\begin{proof}[Proof sketch]
Apply Theorem~\ref{thm:ch3.vsh-orthogonality} and Proposition~\ref{prop:ch2.far-zone-flux-stabilization} from
Subsection~\ref{sec:ch253} \cite{stratton1941,poynting1884}.
% PROOF: AUTHOR REVIEW
\end{proof}

\paragraph{Worked illustration (\(\ell=1\) dipole watt calibration).}
Let \(\widehat{\mathbf{u}}_{10}^{(\mathbf{N})}\) satisfy Eq.~\eqref{eq:ch3.VSH-unit-flux-normalization}. Then
\(\widehat{c}_{10}^{(\mathbf{N})}=1\) represents \(1\,\mathrm{W}\) export in that sector; \(\widehat{c}_{10}^{(\mathbf{N})}=0.1\)
represents \(0.01\,\mathrm{W}\) under Eq.~\eqref{eq:ch3.mode-power-from-coefficient} \cite{harrington2001,stratton1941}.
Pre-normalization geometric coefficients from Section~\ref{sec:ch34} cannot be read as watts without applying
Eqs.~\eqref{eq:ch3.VSH-unit-flux-normalization}--\eqref{eq:ch3.VSH-unit-energy-norm} \cite{poynting1884}.

\paragraph{Problem (convert geometric to normalized coefficients).}
\textbf{Problem.} A geometric coefficient \(c_{10}^{(\mathbf{N})}=2\) is extracted before normalization. Find
\(\widehat{c}_{10}^{(\mathbf{N})}\) after scaling \(\mathbf{u}_{10}^{(\mathbf{N})}\to\widehat{\mathbf{u}}_{10}^{(\mathbf{N})}\)
\cite{hansen1988,harrington2001,stratton1941}.

\textbf{Step 1.} Let \(s>0\) be the scale with \(\widehat{\mathbf{u}}=s\,\mathbf{u}\) \cite{stratton1941}.

\textbf{Step 2.} Impose \(\mathcal{P}_{\mathrm{rad}}(\widehat{\mathbf{u}})=1\) via Eq.~\eqref{eq:ch3.VSH-unit-flux-normalization}
\cite{poynting1884}.

\textbf{Step 3.} Use linearity: \(\widehat{c}=c/s\) \cite{hansen1988}.

\textbf{Physical meaning.} Normalization is a definite amplitude gauge fixing watts, not a change of multipole label
\cite{harrington2001}.

\paragraph{Inner-product--flux consistency.}
On far shells, the bilinear form Eq.~\eqref{eq:ch3.radiative-inner-product-def} tracks coherent export overlap whose
diagonal value under Eq.~\eqref{eq:ch3.VSH-unit-energy-norm} matches unit flux under
Eq.~\eqref{eq:ch3.VSH-unit-flux-normalization} on reciprocal media \cite{stratton1941,harrington2001}. Subsection~\ref{sec:ch351}
therefore aligns Hilbert energy normalization with Poynting meters imported from Subsection~\ref{sec:ch253}
\cite{poynting1884,hansen1988}.

\paragraph{Validity domain.}
Eqs.~\eqref{eq:ch3.VSH-unit-flux-normalization}--\eqref{eq:ch3.mode-power-from-coefficient} require far-zone shells,
outgoing class, and \(\Pi_{\mathrm{rad}}\)-qualified fields \cite{harrington2001,stratton1941}. Near-zone probes must
not be watt-calibrated with these formulas \cite{harrington2001,jackson1999}.

\paragraph{Worked illustration (two-mode power budget).}
For \(\mathbf{v}=\widehat{c}_{10}^{(\mathbf{N})}\widehat{\mathbf{u}}_{10}^{(\mathbf{N})}
+\widehat{c}_{21}^{(\mathbf{M})}\widehat{\mathbf{u}}_{21}^{(\mathbf{M})}\), total power is
\(|\widehat{c}_{10}^{(\mathbf{N})}|^{2}+|\widehat{c}_{21}^{(\mathbf{M})}|^{2}\) watts at leading order
\cite{stratton1941,harrington2001}. Cross-watt terms vanish by Theorem~\ref{thm:ch3.vsh-orthogonality} \cite{hansen1988}.

\paragraph{Problem (verify orthogonality implies watt additivity).}
\textbf{Problem.} Show cross-flux terms vanish for distinct \((\ell,m,\sigma)\) under normalization
\cite{stratton1941,hansen1988,harrington2001}.

\textbf{Step 1.} Expand \(\mathcal{P}_{\mathrm{rad}}\) for sums on \(\mathcal{F}_{\mathrm{rad}}\) \cite{poynting1884}.

\textbf{Step 2.} Apply Theorem~\ref{thm:ch3.vsh-orthogonality} \cite{hansen1988}.

\textbf{Step 3.} Read Eq.~\eqref{eq:ch3.mode-power-from-coefficient} \cite{harrington2001}.

\textbf{Physical meaning.} Watt budgets add on orthogonal labels only after normalization is fixed
\cite{stratton1941}.

\begin{table}[t]
\centering
\caption{Subsection~\ref{sec:ch351} normalization objects (HOME).}
\label{tab:ch3.351-normalization-objects}
\small
\begin{tabular}{@{}p{0.22\linewidth}p{0.38\linewidth}p{0.32\linewidth}@{}}
\toprule
Object & Defining equation & Unit \\
\midrule
Unit flux mode & Eq.~\eqref{eq:ch3.VSH-unit-flux-normalization} & \(1\,\mathrm{W}\) \\
Energy norm & Eq.~\eqref{eq:ch3.VSH-unit-energy-norm} & dimensionless \\
Coefficients & Eq.~\eqref{eq:ch3.normalized-coefficient-resolution} & \(\sqrt{\mathrm{W}}\) \\
Power sum & Eq.~\eqref{eq:ch3.mode-power-from-coefficient} & watts \\
\bottomrule
\end{tabular}
\end{table}

\paragraph{Far-shell flux--inner-product bridge (expanded).}
When \(k_{0}R\gg\ell_{\max}\), the far-shell term in Eq.~\eqref{eq:ch3.radiative-inner-product-def} is bilinear in
transverse components with the same \(1/r^{2}\) flux weighting as Eq.~\eqref{eq:ch2.far-zone-power-from-pattern}
\cite{stratton1941,poynting1884,jackson1999}. Subsection~\ref{sec:ch351} therefore identifies
Eq.~\eqref{eq:ch3.VSH-unit-energy-norm} as the Hilbert diagonal that makes Riesz coefficients in
Eq.~\eqref{eq:ch3.normalized-coefficient-resolution} equal to \(\sqrt{\mathrm{W}}\) amplitudes on reciprocal media
\cite{harrington2001,hansen1988}.

\paragraph{Shell power functional specialization.}
For normalized \(\widehat{\mathbf{u}}_{\ell m}^{(\sigma)}\),
\begin{equation}
\label{eq:ch3.shell-power-normalized-mode}
\mathcal{P}_{\mathrm{rad}}(R;\widehat{\mathbf{u}}_{\ell m}^{(\sigma)})
=
\Re\oint_{\Gamma_R}\hat{\mathbf{r}}\cdot\mathbf{S}_{\mathrm{rad}}[\widehat{\mathbf{u}}_{\ell m}^{(\sigma)}]\,dS
=
1\ \mathrm{W}+\mathcal{O}(R^{-1}),
\end{equation}
importing Eq.~\eqref{eq:ch2.power-flux-functional} with \(\mathbf{S}_{\mathrm{rad}}\) from
Eq.~\eqref{eq:ch2.poynting-functional-radiative} \cite{poynting1884,stratton1941}. Equation
\eqref{eq:ch3.shell-power-normalized-mode} is the operational definition behind
Eq.~\eqref{eq:ch3.VSH-unit-flux-normalization} \cite{harrington2001}.

\paragraph{Worked illustration (calorimetric cross-check).}
If calorimetry reports \(0.95\,\mathrm{W}\) for a single-mode \(\widehat{\mathbf{u}}_{10}^{(\mathbf{N})}\) excitation,
normalization error is \(5\%\) relative to Eq.~\eqref{eq:ch3.VSH-unit-flux-normalization} \cite{harrington2001,stratton1941}.
Recalibrating \(s\) in \(\widehat{\mathbf{u}}=s\mathbf{u}\) restores Eq.~\eqref{eq:ch3.shell-power-normalized-mode}
\cite{poynting1884}.

\paragraph{Problem (build normalization scale from measured power).}
\textbf{Problem.} A geometric mode \(\mathbf{u}_{21}^{(\mathbf{M})}\) yields \(\mathcal{P}_{\mathrm{rad}}=0.25\,\mathrm{W}\) at
\(k_{0}R=20\). Find \(s\) with \(\widehat{\mathbf{u}}=s\mathbf{u}\) satisfying Eq.~\eqref{eq:ch3.VSH-unit-flux-normalization}
\cite{harrington2001,stratton1941,poynting1884}.

\textbf{Step 1.} Power scales as \(|s|^{2}\) on single-mode excitation \cite{stratton1941}.

\textbf{Step 2.} Set \(0.25|s|^{2}=1\) \cite{poynting1884}.

\textbf{Step 3.} Obtain \(s=2\) \cite{hansen1988}.

\textbf{Physical meaning.} Normalization is an amplitude rescaling tied to measured watts, not a relabeling of \(\ell\)
\cite{harrington2001}.

\paragraph{Reactive-field exclusion in watt sums.}
Only \(\Pi_{\mathrm{rad}}\)-qualified fields enter Eqs.~\eqref{eq:ch3.mode-power-from-coefficient}--\eqref{eq:ch3.shell-power-normalized-mode}
\cite{harrington2001,stratton1941}. Near-zone reactive magnitudes may exceed far-zone watts; watt calibration without
outgoing projection violates Eq.~\eqref{eq:ch3.outgoing-inheritance-criterion} \cite{harrington2001}.

\paragraph{Reciprocity pairing under normalization.}
On reciprocal media, Eq.~\eqref{eq:ch3.reciprocity-symmetric-inner-product} is unchanged if
\(\widehat{\mathbf{u}}\to s\widehat{\mathbf{u}}\) and \(\widehat{c}\to\widehat{c}/s\) jointly
\cite{stratton1941,harrington2001}. Subsection~\ref{sec:ch351} fixes \(s\) by flux, not by receive calibration alone
\cite{harrington2001,hansen1988}.

\begin{figure}[t]
\centering
\ChOneFigBlock{%
\begin{tikzpicture}[ch1fig, node distance=7mm and 9mm]
  \node[ch1box] (pi) {$\Pi_{\mathrm{rad}}$};
  \node[ch1box, right=12mm of pi] (shell) {Eq.~\eqref{eq:ch3.shell-power-normalized-mode}};
  \node[ch1box, right=12mm of shell] (norm) {Eq.~\eqref{eq:ch3.VSH-unit-flux-normalization}};
  \node[ch1box, below=10mm of shell] (ip) {Eq.~\eqref{eq:ch3.VSH-unit-energy-norm}};
  \draw[ch1flow] (pi) -- (shell) -- (norm);
  \draw[ch1flow] (ip) -- (norm);
\end{tikzpicture}%
}
\caption{Subsection~\ref{sec:ch351} calibration path: outgoing filter, shell power meter, then unit-flux mode fix.}
\label{fig:ch3.351-calibration-path}
\end{figure}

\paragraph{Antenna effective-length pointer (no re-derivation).}
Classical effective-length definitions relate open-circuit voltage to \(\E_{\mathrm{inc}}\) \cite{harrington2001,stratton1941}.
Under Eq.~\eqref{eq:ch3.normalized-coefficient-resolution}, the same source excites \(\widehat{c}_{\ell m}^{(\sigma)}\) whose
modulus squared is watts in the corresponding sector \cite{harrington2001}. Subsection~\ref{sec:ch351} therefore aligns
modal catalogs with link-budget language used in Subsection~\ref{sec:ch354} \cite{harrington2001}.

\paragraph{Worked illustration (phase coherence in watt sum).}
If \(\widehat{c}_{10}^{(\mathbf{N})}=1\) and \(\widehat{c}_{11}^{(\mathbf{M})}=\jj\), total power remains
\(1^{2}+1^{2}=2\,\mathrm{W}\) because cross-terms vanish by orthogonality, not because phases are aligned
\cite{stratton1941,hansen1988,harrington2001}. Phase affects interference patterns, not orthogonal watt additivity
\cite{harrington2001}.

\paragraph{Problem (coherent versus incoherent superposition).}
\textbf{Problem.} Two normalized modes are excited with coefficients \(1\) and \(e^{\jj\phi}\). Show total power is
\(2\,\mathrm{W}\) for all \(\phi\) \cite{hansen1988,stratton1941,harrington2001}.

\textbf{Step 1.} Expand \(|\widehat{c}_{1}+\widehat{c}_{2}|^{2}\) \cite{jackson1999}.

\textbf{Step 2.} Apply Theorem~\ref{thm:ch3.vsh-orthogonality} to kill cross-term \cite{hansen1988}.

\textbf{Step 3.} Read Eq.~\eqref{eq:ch3.mode-power-from-coefficient} \cite{poynting1884}.

\textbf{Physical meaning.} Orthogonal watt budgets are phase-insensitive at leading order \cite{harrington2001}.

\paragraph{Time-harmonic watt interpretation.}
Under \(\exp(-\jj\omega t)\), Eq.~\eqref{eq:ch3.shell-power-normalized-mode} reports time-averaged watts consistent with
Eq.~\eqref{eq:ch2.far-zone-power-from-pattern} \cite{poynting1884,stratton1941}. Peak instantaneous Poynting differs by
factor two conventions; Subsection~\ref{sec:ch351} follows the time-average meter of Subsection~\ref{sec:ch253}
\cite{stratton1941,harrington2001}.

\paragraph{Worked illustration (magnetic sector watt).}
\(\widehat{\mathbf{u}}_{1,-1}^{(\mathbf{M})}\) normalized by Eq.~\eqref{eq:ch3.VSH-unit-flux-normalization} exports
\(1\,\mathrm{W}\) in the magnetic-dipole sector independently of electric \(\widehat{\mathbf{u}}_{10}^{(\mathbf{N})}\)
\cite{belinfante1940,harrington2001,stratton1941}. Dual sectors do not share watt pools unless explicitly superposed
\cite{hansen1988}.

\paragraph{Problem (total watt from three normalized modes).}
\textbf{Problem.} Coefficients \((1,\,0.5,\,0.5)\) excite three orthogonal normalized modes. Find total power
\cite{harrington2001,stratton1941,poynting1884}.

\textbf{Step 1.} Square coefficients \cite{jackson1999}.

\textbf{Step 2.} Sum \(1+0.25+0.25\) \cite{hansen1988}.

\textbf{Step 3.} Obtain \(1.5\,\mathrm{W}\) \cite{poynting1884}.

\textbf{Physical meaning.} Watt additivity is quadratic in normalized coefficients \cite{harrington2001}.

\begin{table}[t]
\centering
\caption{Subsection~\ref{sec:ch351} calibration checklist on \(\Gamma_R\).}
\label{tab:ch3.351-calibration-checklist}
\small
\begin{tabular}{@{}p{0.08\linewidth}p{0.44\linewidth}p{0.36\linewidth}@{}}
\toprule
Step & Test & Pass \\
\midrule
1 & Outgoing class & Eq.~\eqref{eq:ch3.outgoing-inheritance-criterion} \\
2 & Shell flux & Eq.~\eqref{eq:ch3.shell-power-normalized-mode} \\
3 & Energy norm & Eq.~\eqref{eq:ch3.VSH-unit-energy-norm} \\
4 & Parseval & Eq.~\eqref{eq:ch3.mode-power-from-coefficient} \\
\bottomrule
\end{tabular}
\end{table}

\paragraph{Equivalence to Chapter~\ref{ch:02} field-energy norm (imported).}
On \(\mathcal{F}_{\mathrm{rad}}\), Eq.~\eqref{eq:ch3.VSH-unit-energy-norm} is compatible with the energy norm preview
Eq.~\eqref{eq:ch2.field-energy-norm} restricted to the radiative branch \cite{stratton1941,harrington2001}. Subsection~\ref{sec:ch351}
therefore does not introduce a second energy definition; it fixes the gauge in which Chapter~\ref{ch:02} coefficients
become watts \cite{poynting1884,hansen1988}.

\paragraph{Problem (normalize a quadrupole template).}
\textbf{Problem.} A geometric \(\mathbf{u}_{20}^{(\mathbf{N})}\) is excited with \(c=3\) and measured total power
\(9\,\mathrm{W}\) in that sector alone. Find normalized \(\widehat{c}\) \cite{harrington2001,stratton1941,poynting1884}.

\textbf{Step 1.} Under unit-flux normalization, single-mode power is \(|\widehat{c}|^{2}\) \cite{stratton1941}.

\textbf{Step 2.} Set \(|\widehat{c}|^{2}=9\) \cite{poynting1884}.

\textbf{Step 3.} Obtain \(\widehat{c}=3\) with the correct \(s\) absorbed into \(\widehat{\mathbf{u}}\) \cite{hansen1988}.

\textbf{Physical meaning.} Coefficient magnitude tracks watts only after \(\mathcal{N}\) is fixed \cite{harrington2001}.

\paragraph{Far-zone amplitude linkage.}
Eq.~\eqref{eq:ch2.far-zone-power-from-pattern} ties \(|\E_{0}|^{2}\) integrals to \(\mathcal{P}_{\mathrm{rad}}\); normalized coefficients
\(\widehat{c}_{\ell m}^{(\sigma)}\) set the multipole weights in that integral through
Eq.~\eqref{eq:ch3.normalized-coefficient-resolution} \cite{stratton1941,jackson1999,harrington2001}. Subsection~\ref{sec:ch351}
therefore connects pattern integrals to watt meters without a separate empirical scale factor \cite{poynting1884}.

\paragraph{Measurement uncertainty propagation.}
If \(\mathcal{P}_{\mathrm{rad}}\) is known to \(\pm5\%\), normalization scale \(s\) inherits \(\pm2.5\%\) uncertainty on
amplitude and \(\pm5\%\) on \(|\widehat{c}|^{2}\) \cite{harrington2001,stratton1941}. Subsection~\ref{sec:ch351} recommends
reporting normalization tolerance alongside \(\widehat{c}\) in Eq.~\eqref{eq:ch3.normalized-coefficient-resolution}
\cite{harrington2001,hansen1988}.

\paragraph{Surface power density on \(\Gamma_R\) (HOME specialization).}
For a normalized outgoing pair on a far shell,
\begin{equation}
\label{eq:ch3.shell-power-density-normalized}
\langle\mathbf{S}_{\mathrm{rad}}[\widehat{\mathbf{u}}_{\ell m}^{(\sigma)}]\rangle\cdot\hat{\mathbf{r}}
=
\frac{1}{4\pi R^{2}}\,\mathrm{W\,m^{-2}}
\;+\;\mathcal{O}(R^{-3}),
\end{equation}
importing the time-average Poynting vector from Eq.~\eqref{eq:ch2.poynting-functional-radiative} and integrating
Eq.~\eqref{eq:ch3.shell-power-normalized-mode} over \(\Gamma_R\) \cite{poynting1884,stratton1941}. Equation
\eqref{eq:ch3.shell-power-density-normalized} states that unit-flux normalization spreads one watt uniformly over the
far shell at leading order, a convenient calorimetric cross-check when only local probes are available
\cite{harrington2001,jackson1999}.

\paragraph{Electric--magnetic sector independence.}
Electric \(\widehat{\mathbf{u}}_{\ell m}^{(\mathbf{N})}\) and magnetic \(\widehat{\mathbf{u}}_{\ell m}^{(\mathbf{M})}\) are
normalized independently under Eqs.~\eqref{eq:ch3.VSH-unit-flux-normalization}--\eqref{eq:ch3.VSH-unit-energy-norm}
\cite{belinfante1940,hansen1988,stratton1941}. A source that excites both sectors with unit coefficients therefore exports
\(2\,\mathrm{W}\) when the two labels are orthogonal, not \(1\,\mathrm{W}\) pooled across dual types
\cite{harrington2001,stratton1941}.

\begin{figure}[t]
\centering
\ChOneFigBlock{%
\begin{tikzpicture}[ch1fig, node distance=7mm and 9mm]
  \node[ch1box] (shell) {$\Gamma_R$};
  \node[ch1box, right=12mm of shell] (dens) {Eq.~\eqref{eq:ch3.shell-power-density-normalized}};
  \node[ch1box, right=12mm of dens] (tot) {Eq.~\eqref{eq:ch3.shell-power-normalized-mode}};
  \node[ch1box, below=10mm of dens] (probe) {local $\mathbf{S}\cdot\hat{\mathbf{r}}$};
  \draw[ch1flow] (probe) -- (dens) -- (tot);
  \draw[ch1flow] (shell) -- (dens);
\end{tikzpicture}%
}
\caption{Subsection~\ref{sec:ch351} shell geometry: local Poynting density integrates to the unit-flux shell functional.}
\label{fig:ch3.351-shell-density}
\end{figure}

\paragraph{Problem (local probe versus shell integral).}
\textbf{Problem.} A probe on \(\Gamma_R\) reads \(\langle\mathbf{S}\cdot\hat{\mathbf{r}}\rangle=1/(4\pi R^{2})\,\mathrm{W\,m^{-2}}\)
for single-mode \(\widehat{\mathbf{u}}_{00}^{(\mathbf{N})}\). Verify total shell power is \(1\,\mathrm{W}\)
\cite{poynting1884,stratton1941,harrington2001}.

\textbf{Step 1.} Integrate constant density over \(4\pi R^{2}\) \cite{jackson1999}.

\textbf{Step 2.} Apply Eq.~\eqref{eq:ch3.shell-power-normalized-mode} \cite{harrington2001}.

\textbf{Step 3.} Recover \(1\,\mathrm{W}\) \cite{stratton1941}.

\textbf{Physical meaning.} Local and global flux meters must agree after normalization \cite{poynting1884}.

\paragraph{Radial decay audit.}
Normalized far fields carry the same \(1/r\) amplitude decay as geometric eigenfields; only overall scale changes under
\(\mathbf{u}\to\widehat{\mathbf{u}}\) \cite{stratton1941,hansen1988}. Watt calibration therefore does not alter
Sommerfeld asymptotics locked in Section~\ref{sec:ch24} \cite{stratton1941,harrington2001}.

Subsection~\ref{sec:ch351} has fixed flux and energy normalization for discrete VSH modes on
\(\mathcal{F}_{\mathrm{rad}}\) \cite{stratton1941,poynting1884,harrington2001}. Subsection~\ref{sec:ch352} extends
measure conventions to the continuous spectrum imported from Subsection~\ref{sec:ch313} \cite{hansen1988}.


\subsection{Continuous-spectrum normalization}
\label{sec:ch352}
% TOC tag: [HOME][USE SS3.1.3]

Discrete VSH normalization Subsection~\ref{sec:ch351} does not cover direction continua on \(S^{2}\). Subsection~\ref{sec:ch352}
is the HOME completion of the measure program previewed in Eqs.~\eqref{eq:ch3.continuous-spectral-measure}--\eqref{eq:ch3.continuous-mode-orthogonality}
\cite{stratton1941,jackson1999,harrington2001}.

\paragraph{Standing imports.}
Spectral partition Eq.~\eqref{eq:ch3.spectral-index-partition}, continuous measure preview
Eq.~\eqref{eq:ch3.continuous-spectral-measure}, \(\delta\)-orthogonality
Eq.~\eqref{eq:ch3.continuous-mode-orthogonality}, coefficient density
Eq.~\eqref{eq:ch3.continuous-coefficient-density}, plane-wave bridge
Eq.~\eqref{eq:ch3.plane-wave-continuum-bridge}, and rigged embeddings
Eqs.~\eqref{eq:ch3.gelfand-triple-rad}--\eqref{eq:ch3.delta-directional-mode} enter from Subsections~\ref{sec:ch313}--\ref{sec:ch314}
by reference \cite{stratton1941,hansen1988}. Subsection~\ref{sec:ch352} fixes \(d\mu\) absolutely, not up to scaling
\cite{harrington2001}.

\paragraph{Flux-fixed measure (HOME).}
On \(I_{\mathrm{c}}\cong S^{2}\times\{\sigma\}\), declare
\begin{equation}
\label{eq:ch3.continuous-measure-flux-fix}
d\mu(\beta)
=
\frac{1}{P_{0}}\,
\mathcal{P}_{\mathrm{rad}}\!\bigl(r;\mathbf{f}(\beta)\bigr)\,d\Omega(\hat{\mathbf{k}})\,d\rho(\sigma),
\qquad
\beta=(\hat{\mathbf{k}},\sigma),
\end{equation}
with \(d\Omega\) the solid-angle measure on \(S^{2}\), \(P_{0}\) a reference watt scale, and \(\mathbf{f}(\beta)\) outgoing
generalized modes Eq.~\eqref{eq:ch3.generalized-radiative-eigenmode} \cite{jackson1999,stratton1941}. Equation
\eqref{eq:ch3.continuous-measure-flux-fix} ties \(d\mu\) to the same flux meters as
Eq.~\eqref{eq:ch3.VSH-unit-flux-normalization} \cite{poynting1884,harrington2001}.

\paragraph{Plane-wave \(\delta\) normalization (HOME).}
For plane-wave continuum modes Eq.~\eqref{eq:ch3.plane-wave-continuum-bridge},
\begin{equation}
\label{eq:ch3.plane-wave-delta-normalization}
\big\langle \mathbf{f}(\hat{\mathbf{k}},\sigma),\,\mathbf{f}(\hat{\mathbf{k}}',\sigma')\big\rangle_{\mathrm{rad}}
=
\frac{\delta(\hat{\mathbf{k}}-\hat{\mathbf{k}}')}{\eta_0}\,\delta_{\sigma\sigma'}
\quad\text{under }d\mu,
\end{equation}
matching Eq.~\eqref{eq:ch3.continuous-mode-orthogonality} with explicit flux units \cite{jackson1999,stratton1941,hansen1988}.
Equation~\eqref{eq:ch3.plane-wave-delta-normalization} is the Subsection~\ref{sec:ch352} continuous analogue of unit
energy normalization Eq.~\eqref{eq:ch3.VSH-unit-energy-norm} \cite{hansen1988}.

\paragraph{Continuous Parseval with fixed measure.}
\begin{equation}
\label{eq:ch3.continuous-parseval-flux-fixed}
\mathcal{P}_{\mathrm{rad}}(r;\mathbf{v}_\omega)
=
\int_{I_{\mathrm{c}}}
\big|c(\beta)\big|^{2}\,d\mu(\beta)
\;+\;
\sum_{\ell m\sigma}\big|\widehat{c}_{\ell m}^{(\sigma)}\big|^{2}
\;+\;\mathcal{O}(r^{-1}),
\end{equation}
for hybrid fields combining discrete VSH and continuous direction content via
Eq.~\eqref{eq:ch3.hybrid-spectral-decomposition} \cite{stratton1941,harrington2001,jackson1999}. Equation
\eqref{eq:ch3.continuous-parseval-flux-fixed} extends Eq.~\eqref{eq:ch3.mode-power-from-coefficient} to
\(I_{\mathrm{c}}\) \cite{hansen1988}.

\begin{figure}[t]
\centering
\ChOneFigBlock{%
\begin{tikzpicture}[ch1fig, node distance=7mm and 9mm]
  \node[ch1box] (id) {$I_{\mathrm{d}}$\\Eq.~\eqref{eq:ch3.normalized-coefficient-resolution}};
  \node[ch1box, right=13mm of id] (ic) {$I_{\mathrm{c}}$\\Eq.~\eqref{eq:ch3.continuous-measure-flux-fix}};
  \node[ch1box, below=11mm of id] (w) {watts};
  \node[ch1box, below=11mm of ic] (pw) {Eq.~\eqref{eq:ch3.plane-wave-delta-normalization}};
  \draw[ch1flow] (id) -- (w);
  \draw[ch1flow] (ic) -- (pw) -- (w);
\end{tikzpicture}%
}
\caption{Subsection~\ref{sec:ch352} unifies discrete and continuous normalization under a common flux measure.}
\label{fig:ch3.352-hybrid-measure}
\end{figure}

\paragraph{Worked illustration (narrow beam as localized \(d\mu\)).}
A Gaussian pencil beam concentrated near \(\hat{\mathbf{k}}_{0}\) has coefficient density \(c(\beta)\) peaked on
\(S^{2}\); integrated power \(\int|c|^{2}d\mu\) returns watts under Eq.~\eqref{eq:ch3.continuous-parseval-flux-fixed}
\cite{jackson1999,harrington2001}. Replacing the integral by a single VSH label without \(d\mu\) misstates sidelobe
continuum \cite{hansen1988,stratton1941}.

\paragraph{Problem (check dimensions of \(c(\beta)\)).}
\textbf{Problem.} With \(d\mu\) in Eq.~\eqref{eq:ch3.continuous-measure-flux-fix}, determine the units of \(c(\beta)\)
so that \(|c|^{2}d\mu\) is watts \cite{jackson1999,harrington2001,stratton1941}.

\textbf{Step 1.} Use Eq.~\eqref{eq:ch3.continuous-parseval-flux-fixed} \cite{stratton1941}.

\textbf{Step 2.} Read \(d\mu\) dimensions from flux \(\times\) solid angle \cite{jackson1999}.

\textbf{Step 3.} Conclude \(c\) carries \(\mathrm{W}^{1/2}\) per measure unit \cite{hansen1988}.

\textbf{Physical meaning.} Continuous amplitudes are densities against \(d\mu\), not pointwise watts
\cite{harrington2001}.

\paragraph{Rigged completion pointer.}
Generalized modes \(\mathbf{f}(\beta)\notin\mathcal{F}_{\mathrm{rad}}\) are paired through
Eq.~\eqref{eq:ch3.distributional-pairing-extension} \cite{stratton1941,harrington2001}. Subsection~\ref{sec:ch352}
fixes the measure on \(I_{\mathrm{c}}\); Subsection~\ref{sec:ch314} supplied the topological embedding
\cite{hansen1988,kato1995}. Together they make Eq.~\eqref{eq:ch3.continuous-coefficient-density} measurable
\cite{harrington2001}.

\paragraph{Quadrature on \(S^{2}\) (engineering form).}
Discrete quadrature weights \(\{w_{q}\}\) approximate Eq.~\eqref{eq:ch3.continuous-parseval-flux-fixed} via
Eq.~\eqref{eq:ch3.continuous-parseval-quadrature} with weights scaled to match
Eq.~\eqref{eq:ch3.continuous-measure-flux-fix} \cite{harrington2001,jackson1999}. Quadrature error is separate from
normalization choice but must respect the same \(d\mu\) \cite{hansen1988}.

\paragraph{Worked illustration (hybrid discrete + continuum budget).}
A source with one dominant \(\ell=1\) VSH component plus a broad sidelobe continuum has discrete watts from
\(|\widehat{c}_{1}|^{2}\) and continuous watts from \(\int|c(\beta)|^{2}d\mu\) in
Eq.~\eqref{eq:ch3.continuous-parseval-flux-fixed} \cite{stratton1941,harrington2001}. Omitting the integral
underestimates total export \cite{hansen1988}.

\paragraph{Problem (reject inconsistent \(d\mu\) scaling).}
\textbf{Problem.} Two codes differ by a factor \(a\) in \(d\mu\). Identify the matching rescaling of \(c(\beta)\) that
preserves watts \cite{harrington2001,jackson1999}.

\textbf{Step 1.} Impose invariance of \(\int|c|^{2}d\mu\) \cite{stratton1941}.

\textbf{Step 2.} Replace \(d\mu\to a\,d\mu\), \(c\to a^{-1/2}c\) \cite{hansen1988}.

\textbf{Physical meaning.} Measure and amplitude conventions must be paired \cite{harrington2001}.

\begin{table}[t]
\centering
\caption{Subsection~\ref{sec:ch352} continuous normalization contracts.}
\label{tab:ch3.352-continuous-contracts}
\small
\begin{tabular}{@{}p{0.24\linewidth}p{0.40\linewidth}p{0.28\linewidth}@{}}
\toprule
Object & Equation & Role \\
\midrule
Flux measure & Eq.~\eqref{eq:ch3.continuous-measure-flux-fix} & fixes \(d\mu\) \\
Plane-wave \(\delta\) & Eq.~\eqref{eq:ch3.plane-wave-delta-normalization} & angular orthogonality \\
Hybrid Parseval & Eq.~\eqref{eq:ch3.continuous-parseval-flux-fixed} & watt budget \\
\bottomrule
\end{tabular}
\end{table}

\paragraph{Absolute measure fixing (no residual scaling).}
Equation~\eqref{eq:ch3.continuous-measure-flux-fix} removes the ``up to scaling'' freedom noted below
Eq.~\eqref{eq:ch3.continuous-spectral-measure} \cite{stratton1941,harrington2001}. The reference watt \(P_{0}\) ties
\(d\mu\) to laboratory power meters, not to abstract spectral theory alone \cite{poynting1884,hansen1988}.

\paragraph{Directional delta mode linkage.}
Eq.~\eqref{eq:ch3.delta-directional-mode} from Subsection~\ref{sec:ch314} pairs with
Eq.~\eqref{eq:ch3.plane-wave-delta-normalization}: a sharp direction label carries watt density only after integration
against \(d\mu\) \cite{jackson1999,stratton1941}. Subsection~\ref{sec:ch352} makes that pairing flux-consistent
\cite{hansen1988,harrington2001}.

\paragraph{Worked illustration (finite cone replaces delta).}
Replace \(\delta(\hat{\mathbf{k}}-\hat{\mathbf{k}}_{0})\) by a cone of solid angle \(\Delta\Omega\) with uniform
\(|c|^{2}\). Integrated power \(\int|c|^{2}d\mu\approx |c_{0}|^{2}\Delta\Omega/(4\pi)\) under
Eq.~\eqref{eq:ch3.angular-DOS-continuum-limit} \cite{jackson1999,harrington2001}. Narrow beams approach the delta limit
as \(\Delta\Omega\to0\) \cite{hansen1988}.

\paragraph{Problem (convert W/sr to coefficient density).}
\textbf{Problem.} A pattern carries \(0.1\,\mathrm{W/sr}\) at \(\hat{\mathbf{k}}_{0}\). Estimate \(|c(\hat{\mathbf{k}}_{0})|^{2}\)
under Eqs.~\eqref{eq:ch3.continuous-measure-flux-fix}--\eqref{eq:ch3.angular-DOS-continuum-limit}
\cite{jackson1999,harrington2001,stratton1941}.

\textbf{Step 1.} Use \(\rho_{\Omega}=1/(4\pi)\) \cite{hansen1988}.

\textbf{Step 2.} Match \(0.1=|c|^{2}\rho_{\Omega}\) locally \cite{jackson1999}.

\textbf{Step 3.} Obtain \(|c|^{2}=0.1\times4\pi\approx1.26\) in the chosen gauge \cite{harrington2001}.

\textbf{Physical meaning.} Coefficient density is not the same as radiance; \(d\mu\) converts between them
\cite{stratton1941}.

\paragraph{Bandwidth coupling to \(d\mu\).}
For temporal bandwidth \(\Delta\omega\), hybrid measure products scale with mode count
Eq.~\eqref{eq:ch3.hybrid-DOS-count} \cite{jackson1999,harrington2001}. Subsection~\ref{sec:ch352} fixes angular \(d\mu\);
frequency normalization is inherited from Chapter~\ref{ch:02} spectral density previews without re-derivation here
\cite{stratton1941,hansen1988}.

\begin{figure}[t]
\centering
\ChOneFigBlock{%
\begin{tikzpicture}[ch1fig, node distance=7mm and 9mm]
  \node[ch1box] (pw) {plane-wave family};
  \node[ch1box, right=12mm of pw] (delta) {Eq.~\eqref{eq:ch3.plane-wave-delta-normalization}};
  \node[ch1box, right=12mm of delta] (int) {Eq.~\eqref{eq:ch3.continuous-parseval-flux-fixed}};
  \node[ch1box, below=10mm of delta] (rig) {Eq.~\eqref{eq:ch3.distributional-pairing-extension}};
  \draw[ch1flow] (pw) -- (delta) -- (int);
  \draw[ch1flow] (rig) -- (delta);
\end{tikzpicture}%
}
\caption{Subsection~\ref{sec:ch352} continuous chain: plane-wave modes, \(\delta\) orthogonality, flux Parseval integral.}
\label{fig:ch3.352-plane-wave-chain}
\end{figure}

\paragraph{Quadrature weight calibration.}
Given quadrature nodes \(\{\hat{\mathbf{k}}_{q}\}\) with raw weights \(\tilde{w}_{q}\), calibrate
\(w_{q}=\tilde{w}_{q}\times P_{0}/\mathcal{P}_{\mathrm{ref}}\) so a known reference pattern satisfies
Eq.~\eqref{eq:ch3.continuous-parseval-flux-fixed} \cite{harrington2001,jackson1999}. Miscalibrated \(w_{q}\) mimics
normalization drift in Eq.~\eqref{eq:ch3.coupling-normalization-propagation} \cite{harrington2001}.

\paragraph{Worked illustration (discrete VSH embedded in continuum).}
A very high \(\ell_{\max}\) VSH sum approximates a continuum integral when \(d\mu\) is chosen consistent with
Eq.~\eqref{eq:ch3.continuous-measure-flux-fix} \cite{hansen1988,stratton1941}. The hybrid Parseval law
Eq.~\eqref{eq:ch3.continuous-parseval-flux-fixed} prevents double counting when both sums and integrals appear
\cite{harrington2001}.

\paragraph{Problem (hybrid budget sanity check).}
\textbf{Problem.} Discrete watts \(W_{d}=0.7\) and continuum integral \(W_{c}=0.3\) are reported. Verify total
\(W=1\,\mathrm{W}\) under Eq.~\eqref{eq:ch3.continuous-parseval-flux-fixed} \cite{stratton1941,harrington2001}.

\textbf{Step 1.} Sum components \cite{poynting1884}.

\textbf{Step 2.} Compare to independent calorimetry \cite{harrington2001}.

\textbf{Step 3.} If mismatch, audit \(d\mu\) and \(\Pi_{\mathrm{rad}}\) \cite{stratton1941}.

\textbf{Physical meaning.} Hybrid budgets must partition watts without overlap \cite{hansen1988}.

\paragraph{Solid-angle Jacobian reminder.}
When \(\beta=(\theta,\phi)\) coordinates replace \(\hat{\mathbf{k}}\), \(d\mu\) includes \(\sin\theta\,d\theta\,d\phi\) consistently with
Eq.~\eqref{eq:ch3.continuous-measure-flux-fix} \cite{jackson1999,hansen1988}. Jacobian omissions mimic DOS errors in
Subsection~\ref{sec:ch353} \cite{harrington2001}.

\paragraph{Worked illustration (hemisphere continuum).}
Restricting \(I_{\mathrm{c}}\) to a hemisphere halves the total solid angle; hybrid watts integrate only over the visible
patch with the same \(d\mu\) restriction \cite{jackson1999,harrington2001,stratton1941}. Full-sphere totals require
extending the integral, not rescaling \(d\mu\) alone \cite{hansen1988}.

\paragraph{Problem (verify plane-wave orthogonality dimensions).}
\textbf{Problem.} Show Eq.~\eqref{eq:ch3.plane-wave-delta-normalization} is dimensionally consistent with
Eq.~\eqref{eq:ch3.continuous-parseval-flux-fixed} \cite{jackson1999,harrington2001,stratton1941}.

\textbf{Step 1.} Integrate \(\delta(\hat{\mathbf{k}}-\hat{\mathbf{k}}')\) over \(S^{2}\) \cite{jackson1999}.

\textbf{Step 2.} Pair with \(|c|^{2}d\mu\) \cite{hansen1988}.

\textbf{Step 3.} Read watts \cite{poynting1884}.

\textbf{Physical meaning.} \(\delta\) orthogonality is a measure statement, not a pointwise equality of watts
\cite{harrington2001}.

\begin{table}[t]
\centering
\caption{Subsection~\ref{sec:ch352} continuous normalization audit.}
\label{tab:ch3.352-measure-audit}
\small
\begin{tabular}{@{}p{0.30\linewidth}p{0.36\linewidth}p{0.28\linewidth}@{}}
\toprule
Check & Object & Failure mode \\
\midrule
Measure & Eq.~\eqref{eq:ch3.continuous-measure-flux-fix} & wrong \(P_{0}\) \\
Orthogonality & Eq.~\eqref{eq:ch3.plane-wave-delta-normalization} & missing \(\eta_0\) \\
Budget & Eq.~\eqref{eq:ch3.continuous-parseval-flux-fixed} & double counting \\
\bottomrule
\end{tabular}
\end{table}

\paragraph{Rigged triple consistency.}
Eqs.~\eqref{eq:ch3.gelfand-triple-rad}--\eqref{eq:ch3.distributional-pairing-extension} require that \(d\mu\) in
Eq.~\eqref{eq:ch3.continuous-measure-flux-fix} be the same measure appearing in
Eq.~\eqref{eq:ch3.continuous-mode-orthogonality} \cite{stratton1941,harrington2001,kato1995}. Subsection~\ref{sec:ch352}
closes that consistency for flux-calibrated continua \cite{hansen1988}.

\paragraph{Worked illustration (imported plane-wave continuum).}
Eq.~\eqref{eq:ch3.plane-wave-continuum-bridge} connects \(\beta=\hat{\mathbf{k}}\) to physical plane waves
\cite{jackson1999,stratton1941}. After Eq.~\eqref{eq:ch3.plane-wave-delta-normalization}, coefficient density \(c(\hat{\mathbf{k}})\)
square-integrates to watts against \(d\mu\) \cite{harrington2001,poynting1884}.

\paragraph{Problem (discretize \(d\mu\) on a \(4\times4\) direction grid).}
\textbf{Problem.} Replace \(\int|c|^{2}d\mu\) by \(\sum_{q}w_{q}|c_{q}|^{2}\) with uniform solid-angle cells. Choose \(w_{q}\)
so a unit plane-wave reference yields \(1\,\mathrm{W}\) \cite{harrington2001,jackson1999,stratton1941}.

\textbf{Step 1.} Cell area \(\Delta\Omega\approx4\pi/16\) \cite{jackson1999}.

\textbf{Step 2.} Set \(w_{q}=\Delta\Omega/P_{0}\) per Eq.~\eqref{eq:ch3.continuous-measure-flux-fix} \cite{hansen1988}.

\textbf{Step 3.} Calibrate \(P_{0}\) from reference flux \cite{poynting1884}.

\textbf{Physical meaning.} Quadrature is a discrete measure approximation, not a new normalization \cite{harrington2001}.

\paragraph{Hybrid discrete--continuous split (worked).}
When \(W_{d}+W_{c}=W_{\mathrm{tot}}\), verify \(W_{\mathrm{tot}}\) matches calorimetry before interpreting \(W_{c}\) as
sidelobe continuum \cite{stratton1941,harrington2001}. Subsection~\ref{sec:ch352} treats mismatch as measure or outgoing-class
failure, not as ``missing physics'' \cite{hansen1988}.

\paragraph{Continuum limit of discrete sums.}
As \(\ell_{\max}\to\infty\) with appropriate \(d\mu\), discrete Parseval sums Eq.~\eqref{eq:ch3.mode-power-from-coefficient}
approach continuum integrals Eq.~\eqref{eq:ch3.continuous-parseval-flux-fixed} \cite{hansen1988,stratton1941,jackson1999}.
Subsection~\ref{sec:ch352} supplies the continuum measure needed for that limit \cite{harrington2001}.

\paragraph{Solid-angle Jacobian (HOME).}
Directional continua use
\begin{equation}
\label{eq:ch3.solid-angle-jacobian-measure}
d\mu(\hat{\mathbf{k}})
=
\frac{P_{0}}{4\pi}\,d\Omega(\hat{\mathbf{k}}),
\qquad
d\Omega=\sin\theta\,d\theta\,d\phi,
\end{equation}
with \(P_{0}\) the reference power in Eq.~\eqref{eq:ch3.continuous-measure-flux-fix} \cite{jackson1999,stratton1941}.
Equation~\eqref{eq:ch3.solid-angle-jacobian-measure} ties Lebesgue measure on \(S^{2}\) to watt-normalized plane-wave
coefficients in Eq.~\eqref{eq:ch3.plane-wave-delta-normalization} \cite{harrington2001,hansen1988}.

\paragraph{Worked illustration (hemisphere continuum budget).}
Restricting \(\hat{\mathbf{k}}\) to a hemisphere with uniform \(|c(\hat{\mathbf{k}})|^{2}=1\) gives
\(\int|c|^{2}d\mu=2\pi P_{0}/(4\pi)=P_{0}/2\) watts \cite{jackson1999,harrington2001}. Half-space illumination therefore
occupies half the unit continuum reference of Eq.~\eqref{eq:ch3.continuous-parseval-flux-fixed} \cite{stratton1941}.

\begin{figure}[t]
\centering
\ChOneFigBlock{%
\begin{tikzpicture}[ch1fig, node distance=7mm and 9mm]
  \node[ch1box] (s2) {$S^{2}$};
  \node[ch1box, right=11mm of s2] (jac) {Eq.~\eqref{eq:ch3.solid-angle-jacobian-measure}};
  \node[ch1box, right=11mm of jac] (pw) {Eq.~\eqref{eq:ch3.plane-wave-delta-normalization}};
  \node[ch1box, below=10mm of jac] (par) {Eq.~\eqref{eq:ch3.continuous-parseval-flux-fixed}};
  \draw[ch1flow] (s2) -- (jac) -- (pw) -- (par);
\end{tikzpicture}%
}
\caption{Subsection~\ref{sec:ch352} Jacobian chain: solid-angle measure on \(S^{2}\) feeds flux-fixed Parseval integrals.}
\label{fig:ch3.352-jacobian-chain}
\end{figure}

\paragraph{Problem (convert \(\int|c|^{2}d\Omega\) to watts).}
\textbf{Problem.} A coefficient density satisfies \(\int_{S^{2}}|c|^{2}\,d\Omega=4\pi\) in geometric solid-angle units.
Find watts under Eq.~\eqref{eq:ch3.solid-angle-jacobian-measure} with \(P_{0}=1\,\mathrm{W}\)
\cite{harrington2001,jackson1999,stratton1941}.

\textbf{Step 1.} Insert \(d\mu=P_{0}\,d\Omega/(4\pi)\) \cite{hansen1988}.

\textbf{Step 2.} Evaluate \(\int|c|^{2}d\mu=P_{0}\) \cite{poynting1884}.

\textbf{Step 3.} Obtain \(1\,\mathrm{W}\) \cite{harrington2001}.

\textbf{Physical meaning.} Solid-angle integrals acquire watts only through \(P_{0}\) \cite{stratton1941}.

\begin{table}[t]
\centering
\caption{Subsection~\ref{sec:ch352} measure comparison for direction continua.}
\label{tab:ch3.352-measure-comparison}
\small
\begin{tabular}{@{}p{0.26\linewidth}p{0.34\linewidth}p{0.32\linewidth}@{}}
\toprule
Measure & Definition & Watt link \\
\midrule
\(d\Omega\) & geometric steradian & requires \(P_{0}\) \\
\(d\mu\) & Eq.~\eqref{eq:ch3.solid-angle-jacobian-measure} & flux-fixed \\
Discrete \(w_{q}\) & quadrature weights & calibrate to \(P_{0}\) \\
\bottomrule
\end{tabular}
\end{table}

\paragraph{Plane-wave flux closure (expanded).}
For a unit reference plane wave \(\mathbf{E}_{\mathrm{pw}}(\mathbf{r})=\mathbf{E}_{0}\,e^{-\jj k_{0}\hat{\mathbf{k}}\cdot\mathbf{r}}\) with
\(\hat{\mathbf{k}}\) fixed, Eq.~\eqref{eq:ch3.plane-wave-delta-normalization} fixes \(\mathbf{E}_{0}\) so that
\(\int|c(\hat{\mathbf{k}})|^{2}d\mu=|\mathbf{E}_{0}|^{2}/(2\eta_{0})\) equals one watt when \(|c|=1\) on that direction
\cite{jackson1999,stratton1941,harrington2001}. Subsection~\ref{sec:ch352} therefore closes the \(\eta_{0}\) factor audit flagged in
Table~\ref{tab:ch3.352-measure-audit} \cite{hansen1988,poynting1884}.

\paragraph{Worked illustration (Gaussian beam continuum slice).}
A narrow Gaussian \(|c(\hat{\mathbf{k}})|^{2}\propto\exp[-(\theta/\theta_{0})^{2}]\) integrates to finite watts when
\(\theta_{0}\gg(\Delta\Omega_{\mathrm{cell}})^{1/2}\) on the quadrature grid \cite{harrington2001,jackson1999}. Violating the
resolution condition underestimates sidelobe continuum power in hybrid budgets Eq.~\eqref{eq:ch3.hybrid-DOS-count}
\cite{stratton1941,hansen1988}.

\paragraph{Problem (finite cone solid angle).}
\textbf{Problem.} A conical sector of half-angle \(\theta_{0}=\pi/6\) carries uniform \(|c|^{2}=1\). Find watts under
Eq.~\eqref{eq:ch3.solid-angle-jacobian-measure} with \(P_{0}=1\,\mathrm{W}\) \cite{jackson1999,harrington2001,stratton1941}.

\textbf{Step 1.} Compute \(\Omega_{\mathrm{cone}}=2\pi(1-\cos\theta_{0})\) \cite{jackson1999}.

\textbf{Step 2.} Insert into \(\int|c|^{2}d\mu=\Omega_{\mathrm{cone}}P_{0}/(4\pi)\) \cite{hansen1988}.

\textbf{Step 3.} With \(\cos(\pi/6)=\sqrt{3}/2\), obtain \(\Omega_{\mathrm{cone}}=\pi(2-\sqrt{3})\) and
\(P=\pi(2-\sqrt{3})/(4\pi)\approx0.067\,\mathrm{W}\) \cite{harrington2001}.

\textbf{Physical meaning.} Directional continua assign watts proportional to subtended solid angle \cite{stratton1941}.

\paragraph{Outgoing-class requirement on continua.}
Only \(\Pi_{\mathrm{rad}}\)-qualified plane-wave components enter Eq.~\eqref{eq:ch3.continuous-parseval-flux-fixed}
\cite{harrington2001,stratton1941}. Evanescent or standing-wave continua belong to different measures and must not be
folded into \(d\mu\) without a separate gauge declaration \cite{hansen1988,harrington2001}.

Subsection~\ref{sec:ch352} has fixed continuous-spectrum normalization against flux meters
\cite{stratton1941,jackson1999,harrington2001}. Subsection~\ref{sec:ch353} counts modes under the same conventions
\cite{hansen1988}.


\subsection{Density of states and multiplicity}
\label{sec:ch353}
% TOC tag: [HOME]

Normalization Eqs.~\eqref{eq:ch3.VSH-unit-flux-normalization}--\eqref{eq:ch3.continuous-measure-flux-fix} convert
amplitudes to watts; density of states (DOS) answers how many orthogonal channels occupy a spectral interval
\cite{jackson1999,stratton1941,harrington2001}. Subsection~\ref{sec:ch353} is the HOME mode-counting layer for
angular radiation \cite{hansen1988}.

\paragraph{Discrete multiplet multiplicity (HOME).}
For each \((\ell,\sigma)\), Eq.~\eqref{eq:ch3.J-multiplet-degeneracy-count} gives
\begin{equation}
\label{eq:ch3.multiplet-multiplicity-def}
g_{\ell}^{(\sigma)}=2\ell+1,
\end{equation}
magnetic sublevels at fixed \(\ell\) \cite{jackson1999,hansen1988}. Equation~\eqref{eq:ch3.multiplet-multiplicity-def}
is the discrete DOS at fixed \(\ell\): each sublevel can carry one watt independently after
Eq.~\eqref{eq:ch3.mode-power-from-coefficient} \cite{harrington2001,stratton1941}.

\paragraph{Angular DOS on \(S^{2}\) (HOME).}
Define the exterior angular density of states by
\begin{equation}
\label{eq:ch3.angular-DOS-def}
\rho_{\Omega}(\hat{\mathbf{k}})
=
\frac{dN}{d\Omega},
\qquad
\int_{S^{2}}\rho_{\Omega}\,d\Omega
=
N_{\mathrm{tot}},
\end{equation}
with \(N_{\mathrm{tot}}\) the number of orthogonal direction channels in a bandwidth--aperture class
\cite{jackson1999,stratton1941}. In the ideal full-sphere continuum limit,
\begin{equation}
\label{eq:ch3.angular-DOS-continuum-limit}
\rho_{\Omega}(\hat{\mathbf{k}})=\frac{1}{4\pi},
\qquad
\int_{S^{2}}\rho_{\Omega}\,d\Omega=1,
\end{equation}
consistent with Eq.~\eqref{eq:ch3.continuous-measure-flux-fix} when one direction mode carries unit flux reference
\cite{hansen1988,harrington2001}. Equations~\eqref{eq:ch3.angular-DOS-def}--\eqref{eq:ch3.angular-DOS-continuum-limit}
count \emph{channels per steradian}, not watts; watts follow after multiplying by \(|c|^{2}\) in
Eq.~\eqref{eq:ch3.continuous-parseval-flux-fixed} \cite{poynting1884}.

\paragraph{Hybrid DOS bookkeeping.}
For bandwidth \(\Delta\omega\) and maximum \(\ell_{\max}\),
\begin{equation}
\label{eq:ch3.hybrid-DOS-count}
N_{\mathrm{disc}}
=
\sum_{\ell=0}^{\ell_{\max}}\sum_{\sigma\in\{\mathbf{M},\mathbf{N}\}}(2\ell+1),
\qquad
N_{\mathrm{cont}}
\sim
\frac{\Delta\omega}{\Delta\omega_{\min}}\times N_{\mathrm{aperture}},
\end{equation}
where \(N_{\mathrm{aperture}}\) is the number of resolved direction bins on \(S^{2}\) set by the measurement map
\cite{harrington2001,jackson1999,stratton1941}. Equation~\eqref{eq:ch3.hybrid-DOS-count} links representation
multiplicity Eq.~\eqref{eq:ch3.multiplet-multiplicity-def} to laboratory resolution without confusing \(\ell\) bins
with steradians \cite{hansen1988}.

\begin{figure}[t]
\centering
\ChOneFigBlock{%
\begin{tikzpicture}[ch1fig, node distance=7mm and 10mm]
  \node[ch1box] (ell) {fixed $\ell$};
  \node[ch1box, right=12mm of ell] (g) {Eq.~\eqref{eq:ch3.multiplet-multiplicity-def}};
  \node[ch1box, right=12mm of g] (dos) {Eq.~\eqref{eq:ch3.angular-DOS-def}};
  \node[ch1box, below=10mm of dos] (mu) {Eq.~\eqref{eq:ch3.continuous-measure-flux-fix}};
  \draw[ch1flow] (ell) -- (g);
  \draw[ch1flow] (dos) -- (mu);
\end{tikzpicture}%
}
\caption{Subsection~\ref{sec:ch353} separates discrete multiplet multiplicity from continuous direction DOS on \(S^{2}\).}
\label{fig:ch3.353-dos-split}
\end{figure}

\paragraph{Worked illustration (\(\ell=2\) channel count).}
At \(\ell=2\), Eq.~\eqref{eq:ch3.multiplet-multiplicity-def} gives \(g_{2}^{(\sigma)}=5\) per \(\sigma\), hence ten
orthogonal channels for both \(\sigma\) sectors \cite{jackson1999,hansen1988}. A five-term fit claiming \(\ell=2\) but
using only three independent coefficients violates DOS bookkeeping \cite{harrington2001}.

\paragraph{Problem (count channels to \(\ell_{\max}=3\)).}
\textbf{Problem.} Compute \(N_{\mathrm{disc}}\) from Eq.~\eqref{eq:ch3.hybrid-DOS-count} for \(\ell_{\max}=3\), both
\(\sigma\) \cite{jackson1999,hansen1988}.

\textbf{Step 1.} Sum \(2\ell+1\) for \(\ell=0,1,2,3\) \cite{jackson1999}.

\textbf{Step 2.} Multiply by two \(\sigma\) sectors \cite{belinfante1940}.

\textbf{Step 3.} Obtain \(N_{\mathrm{disc}}=32\) \cite{hansen1988}.

\textbf{Physical meaning.} Mode catalogs must list DOF count, not only peak multipole index \cite{harrington2001}.

\paragraph{Wigner block dimension reminder.}
Rotation blocks Eq.~\eqref{eq:ch3.eigenfield-Wigner-mixing} have dimension \(g_{\ell}^{(\sigma)}\); DOS counts
independent coefficients within each block \cite{jackson1999,hansen1988}. Subsection~\ref{sec:ch353} supplies the
integer \(g_{\ell}^{(\sigma)}\); Subsection~\ref{sec:ch345} supplied the mixing law \cite{stratton1941}.

\paragraph{Measurement-limited DOS.}
Finite apertures replace Eq.~\eqref{eq:ch3.angular-DOS-continuum-limit} with a truncated
\(\rho_{\Omega}\) supported on visible \(\hat{\mathbf{k}}\) \cite{harrington2001,stratton1941}. Accessible angular
budget Eq.~\eqref{eq:ch2.accessible-angular-budget} caps \(N_{\mathrm{tot}}\) independently of mathematical
completeness \cite{harrington2001}. DOS is therefore a \emph{measurement-aware} count \cite{stratton1941}.

\paragraph{Worked illustration (aperture-limited direction bins).}
A \(30^{\circ}\) cone about bore-sight covers \(\Delta\Omega\approx2\pi(1-\cos15^{\circ})\) steradian
\cite{jackson1999,harrington2001}. With one watt-normalized mode per resolved bin under
Eq.~\eqref{eq:ch3.continuous-measure-flux-fix}, total measurable continuum power integrates
\(\int_{\mathrm{cone}}|c|^{2}d\mu\) \cite{hansen1988}.

\begin{table}[t]
\centering
\caption{Subsection~\ref{sec:ch353} DOS summary.}
\label{tab:ch3.353-dos-summary}
\small
\begin{tabular}{@{}p{0.26\linewidth}p{0.34\linewidth}p{0.32\linewidth}@{}}
\toprule
Spectrum & Count law & Normalization partner \\
\midrule
Discrete \(\ell\) & Eq.~\eqref{eq:ch3.multiplet-multiplicity-def} & Eq.~\eqref{eq:ch3.mode-power-from-coefficient} \\
Continuous \(\hat{\mathbf{k}}\) & Eq.~\eqref{eq:ch3.angular-DOS-def} & Eq.~\eqref{eq:ch3.continuous-measure-flux-fix} \\
Hybrid & Eq.~\eqref{eq:ch3.hybrid-DOS-count} & Eq.~\eqref{eq:ch3.continuous-parseval-flux-fixed} \\
\bottomrule
\end{tabular}
\end{table}

\paragraph{Per-\(\ell\) DOS vector.}
Define \(\mathbf{g}_{\ell}=(g_{\ell}^{(\mathbf{M})},g_{\ell}^{(\mathbf{N})})\) with components
Eq.~\eqref{eq:ch3.multiplet-multiplicity-def} \cite{jackson1999,hansen1988}. Total discrete channels up to
\(\ell_{\max}\) are \(\sum_{\ell\le\ell_{\max}}(g_{\ell}^{(\mathbf{M})}+g_{\ell}^{(\mathbf{N})})\), matching
\(N_{\mathrm{disc}}\) in Eq.~\eqref{eq:ch3.hybrid-DOS-count} \cite{harrington2001}.

\paragraph{Local DOS on \(S^{2}\) (engineering).}
When beams are narrow, define local \(\rho_{\Omega}(\hat{\mathbf{k}}_{0})\approx(\Delta\Omega)^{-1}\) on the resolved
patch \cite{jackson1999,harrington2001}. This departs from the uniform limit
Eq.~\eqref{eq:ch3.angular-DOS-continuum-limit} but matches array resolution cells \cite{harrington2001,hansen1988}.

\paragraph{Worked illustration (array element as DOS bin).}
An \(N\)-element uniform circular array resolves \(N\) independent direction bins under narrowband assumptions
\cite{harrington2001,stratton1941}. \(N_{\mathrm{aperture}}\) in Eq.~\eqref{eq:ch3.hybrid-DOS-count} is then
\(\mathcal{O}(N)\), not \(\infty\) \cite{harrington2001}. DOS is finite because the measurement map has finite rank
\cite{stratton1941}.

\paragraph{Problem (DOS versus accessible budget).}
\textbf{Problem.} \(N_{\mathrm{disc}}=32\) but Eq.~\eqref{eq:ch2.accessible-angular-budget} reports \(N_{\mathrm{acc}}=12\).
Reconcile \cite{harrington2001,jackson1999,stratton1941}.

\textbf{Step 1.} \(N_{\mathrm{disc}}\) counts mathematical modes \cite{hansen1988}.

\textbf{Step 2.} \(N_{\mathrm{acc}}\) counts measurement-accessible modes \cite{harrington2001}.

\textbf{Step 3.} Expect \(N_{\mathrm{acc}}\le N_{\mathrm{disc}}\) \cite{stratton1941}.

\textbf{Physical meaning.} Full DOS is not fully observable \cite{harrington2001}.

\paragraph{Multiplicity versus normalization.}
Each of the \(g_{\ell}^{(\sigma)}\) channels can carry \(1\,\mathrm{W}\) independently only after
Eq.~\eqref{eq:ch3.VSH-unit-flux-normalization} \cite{poynting1884,harrington2001}. DOS counts channels; normalization
assigns watts per channel \cite{stratton1941,hansen1988}.

\begin{figure}[t]
\centering
\ChOneFigBlock{%
\begin{tikzpicture}[ch1fig, node distance=7mm and 9mm]
  \node[ch1box] (math) {$N_{\mathrm{disc}}$};
  \node[ch1box, right=12mm of math] (meas) {$N_{\mathrm{acc}}$};
  \node[ch1box, right=12mm of meas] (w) {watts via Eq.~\eqref{eq:ch3.mode-power-from-coefficient}};
  \draw[ch1flow] (math) -- (meas) -- (w);
\end{tikzpicture}%
}
\caption{Subsection~\ref{sec:ch353} separates mathematical DOS from accessible DOF before watt assignment.}
\label{fig:ch3.353-accessible-dos}
\end{figure}

\paragraph{Worked illustration (\(\ell=0\) monopole channel).}
\(\ell=0\) contributes \(g_{0}^{(\sigma)}=1\) per \(\sigma\) sector \cite{jackson1999}. Even a ``monopole'' term
consumes one watt slot in Eq.~\eqref{eq:ch3.mode-power-from-coefficient} if excited \cite{harrington2001,stratton1941}.

\paragraph{Problem (estimate continuum bins at \(\lambda/2\) spacing).}
\textbf{Problem.} A spherical scan at radius \(R\) uses \(\lambda_0/2\) sample spacing on \(S^{2}\). Estimate
\(N_{\mathrm{aperture}}\) scaling \cite{harrington2001,jackson1999}.

\textbf{Step 1.} Count samples \(\propto(4\pi R^{2})/(\lambda_0/2)^{2}\) \cite{jackson1999}.

\textbf{Step 2.} Insert into Eq.~\eqref{eq:ch3.hybrid-DOS-count} \cite{hansen1988}.

\textbf{Physical meaning.} Spatial sampling caps continuous DOS \cite{harrington2001}.

\paragraph{Cavity versus exterior DOS (pointer).}
Bounded cavities quantize \(I_{\mathrm{d}}\) with additional density laws deferred to Chapter~\ref{ch:04} previews
\cite{jackson1999,stratton1941}. Subsection~\ref{sec:ch353} counts exterior angular channels on \(S^{2}\) and discrete
\(\ell\) multiplets without reopening cavity derivations \cite{hansen1988,harrington2001}.

\paragraph{Worked illustration (duplicate \(\ell\) labels).}
Reporting ``two \(\ell=1\) modes'' without \((m,\sigma)\) splits undercounts DOS: there are six unit-flux slots at
\(\ell=1\) after Eq.~\eqref{eq:ch3.VSH-unit-flux-normalization} \cite{belinfante1940,harrington2001,hansen1988}.
Subsection~\ref{sec:ch353} requires full label tuples in mode catalogs \cite{stratton1941}.

\paragraph{Problem (DOS-weighted average power).}
\textbf{Problem.} Powers \(\{P_{n}\}\) are listed on \(N_{\mathrm{acc}}=12\) accessible modes. Relate \(\sum P_{n}\) to total
\(\mathcal{P}_{\mathrm{rad}}\) when \(N_{\mathrm{disc}}=32\) \cite{harrington2001,stratton1941,jackson1999}.

\textbf{Step 1.} Sum accessible powers \cite{poynting1884}.

\textbf{Step 2.} Unobserved modes may still carry power in \(\mathcal{R}_\omega[\mathfrak{D}]\) \cite{harrington2001}.

\textbf{Step 3.} Expect \(\sum_{n\in\mathrm{acc}}P_{n}\le\mathcal{P}_{\mathrm{rad}}\) \cite{stratton1941}.

\textbf{Physical meaning.} Accessible DOS is a measurement cap, not a physical power cap \cite{harrington2001}.

\begin{table}[t]
\centering
\caption{Subsection~\ref{sec:ch353} label-resolved channel count at low \(\ell\).}
\label{tab:ch3.353-lowell-channel-count}
\small
\begin{tabular}{@{}p{0.12\linewidth}p{0.22\linewidth}p{0.22\linewidth}p{0.22\linewidth}@{}}
\toprule
\(\ell\) & \(g_{\ell}^{(\mathbf{M})}\) & \(g_{\ell}^{(\mathbf{N})}\) & total \\
\midrule
0 & 1 & 1 & 2 \\
1 & 3 & 3 & 6 \\
2 & 5 & 5 & 10 \\
\bottomrule
\end{tabular}
\end{table}

\paragraph{Wigner block dimension check.}
Eq.~\eqref{eq:ch3.eigenfield-Wigner-mixing} acts on spaces of dimension \(g_{\ell}^{(\sigma)}\); DOS tables must list
\(g_{\ell}^{(\sigma)}\), not only \(\ell\) \cite{jackson1999,hansen1988}. Subsection~\ref{sec:ch353} supplies integers for
that dimension audit \cite{stratton1941,harrington2001}.

\paragraph{Problem (watts if all \(\ell\le1\) channels excited equally).}
\textbf{Problem.} Excite all six \(\ell\le1\) normalized channels with \(|\widehat{c}|^{2}=0.1\,\mathrm{W}\) each. Total power?
\cite{harrington2001,stratton1941,poynting1884}.

\textbf{Step 1.} Count six channels from Table~\ref{tab:ch3.353-lowell-channel-count} \cite{hansen1988}.

\textbf{Step 2.} Sum \(6\times0.1\) \cite{poynting1884}.

\textbf{Step 3.} Obtain \(0.6\,\mathrm{W}\) \cite{harrington2001}.

\textbf{Physical meaning.} Filling all DOS slots at equal power is a structured budget test \cite{stratton1941}.

\paragraph{Angular Nyquist heuristic.}
Resolve \(\ell_{\max}\) on \(S^{2}\) only when sample spacing satisfies \(\Delta\theta\lesssim\pi/\ell_{\max}\) heuristically
\cite{harrington2001,jackson1999}. Violating the heuristic undercounts \(N_{\mathrm{aperture}}\) and inflates truncation error
analyzed in Section~\ref{sec:ch36} \cite{hansen1988,stratton1941}.

\paragraph{Worked illustration (rank cap from Eq.~\eqref{eq:ch2.reconstruction-rank-cap}).}
If the observation map has rank \(r\), only \(r\) normalized coefficients can be recovered stably regardless of
\(N_{\mathrm{disc}}\) \cite{harrington2001,stratton1941}. DOS must be intersected with rank caps before watt assignment
\cite{harrington2001}.

\paragraph{DOS density per steradian (continuum).}
When \(\rho_{\Omega}=1/(4\pi)\), one watt-normalized mode occupies \(4\pi\) steradians of direction space in the uniform
continuum idealization \cite{jackson1999,hansen1988}. Laboratory beams with \(\Delta\Omega\ll4\pi\) therefore occupy a
fraction \(\Delta\Omega/(4\pi)\) of the unit continuum reference \cite{harrington2001,stratton1941}.

\paragraph{Problem (accessible DOF from aperture rank).}
\textbf{Problem.} A measurement map has rank \(r=8\) but \(N_{\mathrm{disc}}=32\). How many watt coefficients can be
independently identified? \cite{harrington2001,stratton1941,jackson1999}.

\textbf{Step 1.} Independent recoverable coefficients \(\le r\) \cite{harrington2001}.

\textbf{Step 2.} Remaining modes are null under the observation functional \cite{stratton1941}.

\textbf{Physical meaning.} DOS exceeds accessible DOF when apertures are small \cite{harrington2001}.

\paragraph{Cumulative discrete DOS (HOME).}
Define the running channel count
\begin{equation}
\label{eq:ch3.cumulative-disc-DOS}
N_{\mathrm{disc}}(\ell_{\max})
=
\sum_{\ell=0}^{\ell_{\max}}\sum_{\sigma\in\{\mathbf{M},\mathbf{N}\}} g_{\ell}^{(\sigma)},
\qquad
g_{\ell}^{(\sigma)}\ \text{from Eq.~\eqref{eq:ch3.multiplet-multiplicity-def}},
\end{equation}
so \(N_{\mathrm{disc}}(2)=22\) by Table~\ref{tab:ch3.353-lowell-channel-count} extended to \(\ell=2\)
\cite{jackson1999,hansen1988}. Equation~\eqref{eq:ch3.cumulative-disc-DOS} is the discrete DOS staircase used when
assigning watt budgets to truncated catalogs \cite{harrington2001,stratton1941}.

\paragraph{High-\(\ell\) asymptotics.}
For large \(\ell\), \(g_{\ell}^{(\mathbf{M})}+g_{\ell}^{(\mathbf{N})}=4\ell+2\) per Eq.~\eqref{eq:ch3.J-multiplet-degeneracy-count},
so \(N_{\mathrm{disc}}(\ell_{\max})\sim 2\ell_{\max}^{2}\) \cite{jackson1999,hansen1988}. Angular Nyquist heuristics in
Subsection~\ref{sec:ch353} therefore scale aperture sampling with \(\ell_{\max}^{2}\), not linearly in \(\ell_{\max}\)
\cite{harrington2001,stratton1941}.

\begin{figure}[t]
\centering
\ChOneFigBlock{%
\begin{tikzpicture}[ch1fig, node distance=7mm and 9mm]
  \node[ch1box] (ell) {$\ell_{\max}$};
  \node[ch1box, right=11mm of ell] (cum) {Eq.~\eqref{eq:ch3.cumulative-disc-DOS}};
  \node[ch1box, right=11mm of cum] (hyb) {Eq.~\eqref{eq:ch3.hybrid-DOS-count}};
  \node[ch1box, below=10mm of cum] (acc) {Eq.~\eqref{eq:ch2.accessible-angular-budget}};
  \draw[ch1flow] (ell) -- (cum) -- (hyb);
  \draw[ch1flow] (acc) -- (hyb);
\end{tikzpicture}%
}
\caption{Subsection~\ref{sec:ch353} cumulative DOS: discrete staircase meets accessible angular budget before watt assignment.}
\label{fig:ch3.353-cumulative-dos}
\end{figure}

\paragraph{Problem (hybrid count at \(\ell_{\max}=1\)).}
\textbf{Problem.} Using Eqs.~\eqref{eq:ch3.cumulative-disc-DOS} and \eqref{eq:ch3.angular-DOS-def}, report discrete and
continuum contributions to \(N_{\mathrm{hyb}}\) at \(\ell_{\max}=1\) \cite{harrington2001,hansen1988,stratton1941}.

\textbf{Step 1.} Read \(N_{\mathrm{disc}}(1)=8\) from Table~\ref{tab:ch3.353-lowell-channel-count} \cite{hansen1988}.

\textbf{Step 2.} Continuum DOS on \(S^{2}\) adds direction degrees beyond the multipole cut \cite{jackson1999}.

\textbf{Step 3.} Assemble Eq.~\eqref{eq:ch3.hybrid-DOS-count} without double counting shared labels \cite{harrington2001}.

\textbf{Physical meaning.} Hybrid counts separate multiplet slots from direction continua \cite{stratton1941}.

\paragraph{Spectral occupancy and watt floors.}
Exciting all \(N_{\mathrm{disc}}(\ell_{\max})\) normalized channels at \(|\widehat{c}|^{2}=P_{0}/N_{\mathrm{disc}}\) yields total
\(P_{0}\) watts \cite{harrington2001,poynting1884}. Uniform occupancy is a structured test of
Eq.~\eqref{eq:ch3.mode-power-from-coefficient} against DOS tables \cite{stratton1941,hansen1988}.

\paragraph{Label-resolved DOS ledger (expanded).}
At fixed \(\ell\), electric and magnetic multiplets contribute \(g_{\ell}^{(\mathbf{N})}\) and \(g_{\ell}^{(\mathbf{M})}\) channels
separately in Eq.~\eqref{eq:ch3.cumulative-disc-DOS} \cite{belinfante1940,hansen1988}. Laboratory sources that excite only
\(\mathbf{N}\)-type sectors therefore fill at most \(\sum_{\ell\le\ell_{\max}}g_{\ell}^{(\mathbf{N})}\) normalized slots, not the
full \(N_{\mathrm{disc}}(\ell_{\max})\) staircase \cite{stratton1941,harrington2001}.

\paragraph{Worked illustration (\(\ell_{\max}=2\) cumulative count).}
From Eq.~\eqref{eq:ch3.cumulative-disc-DOS} and Table~\ref{tab:ch3.353-lowell-channel-count},
\(N_{\mathrm{disc}}(2)=2+6+10=18\) normalized channels \cite{hansen1988,jackson1999}. Assigning \(0.05\,\mathrm{W}\) to each occupied
channel yields \(0.9\,\mathrm{W}\) total under Eq.~\eqref{eq:ch3.mode-power-from-coefficient} \cite{poynting1884,harrington2001}.

\paragraph{Problem (compare discrete and continuum DOS at fixed \(\ell_{\max}\)).}
\textbf{Problem.} At \(\ell_{\max}=1\), report \(N_{\mathrm{disc}}(1)\) and the continuum direction measure of a full \(4\pi\) sphere
under Eq.~\eqref{eq:ch3.angular-DOS-def} \cite{jackson1999,harrington2001,stratton1941}.

\textbf{Step 1.} Read \(N_{\mathrm{disc}}(1)=8\) \cite{hansen1988}.

\textbf{Step 2.} Continuum \(\rho_{\Omega}=1/(4\pi)\) integrates to one direction unit on \(S^{2}\) \cite{jackson1999}.

\textbf{Step 3.} Assemble Eq.~\eqref{eq:ch3.hybrid-DOS-count} without identifying multipoles with isolated directions
\cite{harrington2001}.

\textbf{Physical meaning.} Discrete multiplets and direction continua answer different counting questions \cite{stratton1941}.

\begin{table}[t]
\centering
\caption{Subsection~\ref{sec:ch353} cumulative discrete DOS through \(\ell_{\max}=2\).}
\label{tab:ch3.353-cumulative-ledger}
\small
\begin{tabular}{@{}p{0.18\linewidth}p{0.22\linewidth}p{0.22\linewidth}p{0.28\linewidth}@{}}
\toprule
\(\ell_{\max}\) & \(N_{\mathrm{disc}}\) & per-channel \(|\widehat{c}|^{2}\) & total \(P\) \\
\midrule
0 & 2 & $0.5\,\mathrm{W}$ & $1\,\mathrm{W}$ \\
1 & 8 & $0.125\,\mathrm{W}$ & $1\,\mathrm{W}$ \\
2 & 18 & $0.0556\,\mathrm{W}$ & $1\,\mathrm{W}$ \\
\bottomrule
\end{tabular}
\end{table}

Subsection~\ref{sec:ch353} has defined discrete multiplicity and angular DOS consistent with Sections~\ref{sec:ch34}--\ref{sec:ch352}
\cite{jackson1999,hansen1988,harrington2001}. Subsection~\ref{sec:ch354} records how normalization choices propagate into
coupling and power reports \cite{stratton1941}.


\subsection{Consequences of normalization choice}
\label{sec:ch354}
% TOC tag: [HOME]

Once Eqs.~\eqref{eq:ch3.VSH-unit-flux-normalization}--\eqref{eq:ch3.continuous-measure-flux-fix} are fixed, every
downstream coupling matrix, power report, and reciprocity statement inherits the same gauge \cite{stratton1941,harrington2001}.
Subsection~\ref{sec:ch354} is the HOME consequence calculus for that inheritance \cite{hansen1988,poynting1884}.

\paragraph{Coupling normalization propagation (HOME).}
Let \(K_{mn}=\langle \widehat{\mathbf{u}}_{m},\mathcal{O}\widehat{\mathbf{u}}_{n}\rangle_{\mathrm{rad}}\) for a linear
observation operator \(\mathcal{O}\) on \(\mathcal{F}_{\mathrm{rad}}\) \cite{harrington2001,stratton1941}. Changing
normalization \(\widehat{\mathbf{u}}\to a_{m}\widehat{\mathbf{u}}_{m}\) scales rows and columns as
\begin{equation}
\label{eq:ch3.coupling-normalization-propagation}
K'_{mn}
=
\frac{K_{mn}}{a_{m}a_{n}},
\qquad
\widehat{c}'_{m}
=
\frac{\widehat{c}_{m}}{a_{m}}.
\end{equation}
Equation~\eqref{eq:ch3.coupling-normalization-propagation} is the Subsection~\ref{sec:ch354} warning: comparing coupling
matrices across codes requires matching Eqs.~\eqref{eq:ch3.VSH-unit-flux-normalization}--\eqref{eq:ch3.VSH-unit-energy-norm}
\cite{harrington2001,hansen1988}.

\paragraph{Power-budget invariance (HOME).}
Physical power is invariant under paired rescalings:
\begin{equation}
\label{eq:ch3.power-budget-invariance}
\mathcal{P}_{\mathrm{rad}}
=
\sum_{m,n}\widehat{c}_{m}^{*}K_{mn}\widehat{c}_{n}
\quad\Rightarrow\quad
\mathcal{P}'_{\mathrm{rad}}=\mathcal{P}_{\mathrm{rad}}
\end{equation}
when \(K,\widehat{c}\) transform together under Eq.~\eqref{eq:ch3.coupling-normalization-propagation} on reciprocal
media \cite{stratton1941,poynting1884,harrington2001}. Equation~\eqref{eq:ch3.power-budget-invariance} separates
\emph{gauge} (normalization choice) from \emph{physics} (watts exported by \(\mathcal{R}_\omega\)) \cite{harrington2001}.

\paragraph{Reciprocity constraint (HOME).}
On lossless reciprocal media, reciprocity pairing Eq.~\eqref{eq:ch3.reciprocity-symmetric-inner-product} implies
\begin{equation}
\label{eq:ch3.reciprocity-coupling-symmetry}
K_{mn}=K_{nm},
\qquad
\mathcal{O}\ \text{self-adjoint on }\mathcal{F}_{\mathrm{rad}},
\end{equation}
for passive observation operators \cite{stratton1941,harrington2001}. Subsection~\ref{sec:ch354} records that reciprocity
is preserved only if normalization is applied identically to transmit and receive modes
\cite{harrington2001,hansen1988}.

\begin{figure}[t]
\centering
\ChOneFigBlock{%
\begin{tikzpicture}[ch1fig, node distance=7mm and 9mm]
  \node[ch1box] (norm) {Eqs.~\eqref{eq:ch3.VSH-unit-flux-normalization}--\eqref{eq:ch3.continuous-measure-flux-fix}};
  \node[ch1box, right=11mm of norm] (coup) {Eq.~\eqref{eq:ch3.coupling-normalization-propagation}};
  \node[ch1box, right=11mm of coup] (pow) {Eq.~\eqref{eq:ch3.power-budget-invariance}};
  \node[ch1box, below=10mm of coup] (rec) {Eq.~\eqref{eq:ch3.reciprocity-coupling-symmetry}};
  \draw[ch1flow] (norm) -- (coup) -- (pow);
  \draw[ch1flow] (rec) -- (coup);
\end{tikzpicture}%
}
\caption{Subsection~\ref{sec:ch354} consequence chain: normalization gauge fixes coupling reports while physical watts remain invariant.}
\label{fig:ch3.354-consequence-chain}
\end{figure}

\paragraph{Measurement composition restatement.}
Eq.~\eqref{eq:ch3.measurement-on-coefficients} becomes
\begin{equation}
\label{eq:ch3.measurement-normalized-coefficients}
\mathcal{M}[\mathfrak{F}]
=
\mathcal{N}\circ\mathcal{C}\circ\mathcal{R}_\omega[\mathfrak{D}],
\qquad
\mathcal{N}\ \text{implements Eqs.~\eqref{eq:ch3.VSH-unit-flux-normalization}--\eqref{eq:ch3.continuous-measure-flux-fix}},
\end{equation}
importing the Chapter~\ref{ch:03} measurement preview with normalization explicit \cite{harrington2001,stratton1941}.
Subsection~\ref{sec:ch354} therefore closes the loop opened in Subsection~\ref{sec:ch301}: coefficients are now
\(\mathcal{N}\)-fixed amplitudes, not raw projections \cite{hansen1988}.

\paragraph{Worked illustration (mismatched library normalization).}
If transmit modes use Eq.~\eqref{eq:ch3.VSH-unit-flux-normalization} but receive modes omit energy normalization
Eq.~\eqref{eq:ch3.VSH-unit-energy-norm}, reported coupling \(K_{mn}\) is off by known factors \(a_{m}a_{n}\) from
Eq.~\eqref{eq:ch3.coupling-normalization-propagation} \cite{harrington2001}. Link budgets remain correct if power is
assembled through Eq.~\eqref{eq:ch3.mode-power-from-coefficient} with paired norms \cite{stratton1941,poynting1884}.

\paragraph{Problem (rescale coupling after normalization change).}
\textbf{Problem.} A code switches from geometric to flux normalization with \(a_{\ell m}^{(\sigma)}=\sqrt{\mathcal{P}_{\mathrm{rad}}(\mathbf{u})}\).
Update stored coupling \(K\) \cite{harrington2001,hansen1988,stratton1941}.

\textbf{Step 1.} Record old \(K_{mn}\) in geometric gauge \cite{harrington2001}.

\textbf{Step 2.} Apply Eq.~\eqref{eq:ch3.coupling-normalization-propagation} with \(a_{m},a_{n}\) \cite{hansen1988}.

\textbf{Step 3.} Verify \(\mathcal{P}_{\mathrm{rad}}\) unchanged via Eq.~\eqref{eq:ch3.power-budget-invariance}
\cite{poynting1884}.

\textbf{Physical meaning.} Coupling reports are gauge-dependent; watts are not \cite{stratton1941}.

\paragraph{Far-zone pattern amplitude linkage.}
Pattern functions \(|\E_{0}(\theta,\phi)|\) from Eq.~\eqref{eq:ch2.far-zone-power-from-pattern} inherit overall scale from
\(\widehat{c}_{\ell m}^{(\sigma)}\) through Eq.~\eqref{eq:ch3.normalized-coefficient-resolution}
\cite{stratton1941,jackson1999}. Changing Condon--Shortley phase conventions does not alter watts under
Eq.~\eqref{eq:ch3.mode-power-from-coefficient} but can rotate coupling phases in
Eq.~\eqref{eq:ch3.reciprocity-coupling-symmetry} \cite{hansen1988,harrington2001}.

\paragraph{Truncation consequence.}
Finite \(\ell_{\max}\) truncations leave a remainder orthogonal to the truncated subspace only when inner products use
the same \(\mathcal{N}\) \cite{hansen1988,stratton1941}. Mixing truncated geometric projections with normalized power
sums violates Eq.~\eqref{eq:ch3.power-budget-invariance} \cite{harrington2001}. Section~\ref{sec:ch36} analyzes
truncation error; Subsection~\ref{sec:ch354} records the normalization prerequisite \cite{stratton1941}.

\paragraph{Worked illustration (reciprocal two-port).}
In a lossless two-port measurement, transmit coefficient \(\widehat{c}_{\mathrm{tx}}\) and receive coefficient
\(\widehat{c}_{\mathrm{rx}}\) enter \(|\widehat{c}_{\mathrm{tx}}|^{2}\) and \(|\widehat{c}_{\mathrm{rx}}|^{2}\) watt
terms separately; reciprocity equates certain coupling entries via Eq.~\eqref{eq:ch3.reciprocity-coupling-symmetry}
\cite{stratton1941,harrington2001}. Normalizing only one port breaks symmetry without changing true exported power
\cite{harrington2001}.

\begin{table}[t]
\centering
\caption{Subsection~\ref{sec:ch354} normalization consequence map.}
\label{tab:ch3.354-consequence-map}
\small
\begin{tabular}{@{}p{0.28\linewidth}p{0.36\linewidth}p{0.30\linewidth}@{}}
\toprule
Report object & Transforms as & Invariant \\
\midrule
Coupling \(K_{mn}\) & Eq.~\eqref{eq:ch3.coupling-normalization-propagation} & --- \\
Coefficients \(\widehat{c}\) & \(1/a_{m}\) & --- \\
Power \(\mathcal{P}_{\mathrm{rad}}\) & --- & Eq.~\eqref{eq:ch3.power-budget-invariance} \\
Reciprocity & Eq.~\eqref{eq:ch3.reciprocity-coupling-symmetry} & paired norms \\
\bottomrule
\end{tabular}
\end{table}

\paragraph{Link-budget form (HOME corollary).}
For normalized coefficients,
\begin{equation}
\label{eq:ch3.link-budget-normalized}
P_{\mathrm{rx}}
=
\Big|\sum_{n}\widehat{c}_{n}^{\mathrm{rx}}\,K_{n,\mathrm{tx}}\,\widehat{c}_{\mathrm{tx}}\Big|^{2}
\quad\text{(single-frequency, matched normalization)},
\end{equation}
with \(K\) the coupling matrix under Eq.~\eqref{eq:ch3.coupling-normalization-propagation} and ports using the same
\(\mathcal{N}\) \cite{harrington2001,stratton1941}. Equation~\eqref{eq:ch3.link-budget-normalized} is the Subsection~\ref{sec:ch354}
template for communication-style budgets once Eq.~\eqref{eq:ch3.mode-power-from-coefficient} fixes transmit watts
\cite{harrington2001,poynting1884}.

\paragraph{Phase convention propagation.}
Condon--Shortley phases in Eq.~\eqref{eq:ch3.Ylm-Condon-Shortley-def} rotate coupling phases in \(K_{mn}\) but leave
\(|\widehat{c}|^{2}\) sums invariant \cite{hansen1988,harrington2001}. Comparing coupling phases across software
requires matching both \(\mathcal{N}\) and spherical phase conventions \cite{hansen1988,stratton1941}.

\paragraph{Worked illustration (one-port recalibration).}
If receive modes are renormalized by \(a_{\mathrm{rx}}\) only, \(K_{mn}\) picks up \(1/a_{\mathrm{rx}}\) while transmit
power remains correct when computed from transmit coefficients alone \cite{harrington2001}. Reciprocity symmetry
Eq.~\eqref{eq:ch3.reciprocity-coupling-symmetry} fails unless both ports rescale \cite{stratton1941,harrington2001}.

\paragraph{Problem (trace identity for power).}
\textbf{Problem.} For diagonal \(K\) in the normalized VSH basis, show \(\mathcal{P}=\sum_{n}|\widehat{c}_{n}|^{2}\)
\cite{hansen1988,stratton1941,harrington2001}.

\textbf{Step 1.} Use Eq.~\eqref{eq:ch3.mode-power-from-coefficient} \cite{poynting1884}.

\textbf{Step 2.} Orthogonality kills cross terms \cite{hansen1988}.

\textbf{Physical meaning.} Diagonal coupling is the watt-sum gauge of Subsection~\ref{sec:ch351} \cite{harrington2001}.

\paragraph{Inverse-problem caution.}
Reconstructing \(\widehat{c}\) from far-field samples without declaring \(\mathcal{N}\) in
Eq.~\eqref{eq:ch3.measurement-normalized-coefficients} yields gauge-ambiguous amplitudes \cite{devaney2004,harrington2001}.
Subsection~\ref{sec:ch354} requires normalization metadata alongside any inversion report \cite{harrington2001,stratton1941}.

\begin{figure}[t]
\centering
\ChOneFigBlock{%
\begin{tikzpicture}[ch1fig, node distance=7mm and 9mm]
  \node[ch1box] (tx) {TX $\widehat{c}$};
  \node[ch1box, right=11mm of tx] (k) {$K_{mn}$};
  \node[ch1box, right=11mm of k] (rx) {RX $\widehat{c}$};
  \node[ch1box, below=10mm of k] (n) {$\mathcal{N}$};
  \draw[ch1flow] (tx) -- (k) -- (rx);
  \draw[ch1flow] (n) -- (k);
\end{tikzpicture}%
}
\caption{Subsection~\ref{sec:ch354} two-port reporting: coupling inherits normalization \(\mathcal{N}\) on both sides.}
\label{fig:ch3.354-two-port-normalization}
\end{figure}

\paragraph{Software interchange contract.}
Archive files should record Eqs.~\eqref{eq:ch3.VSH-unit-flux-normalization}, \eqref{eq:ch3.VSH-unit-energy-norm},
and \eqref{eq:ch3.continuous-measure-flux-fix} scalings as metadata \cite{harrington2001,hansen1988}. Without that
metadata, Eq.~\eqref{eq:ch3.coupling-normalization-propagation} cannot reconcile libraries \cite{harrington2001}.

\paragraph{Worked illustration (pattern vs watts).}
Radiation patterns are often plotted normalized to peak unity; watt budgets require
Eq.~\eqref{eq:ch3.normalized-coefficient-resolution} \cite{stratton1941,jackson1999}. Subsection~\ref{sec:ch354}
separates shape plots from calibrated amplitudes \cite{harrington2001,poynting1884}.

\paragraph{Problem (audit a published coupling matrix).}
\textbf{Problem.} A paper reports \(K\) in geometric gauge. Convert to flux gauge given per-mode powers \(\{P_{n}\}\)
\cite{harrington2001,hansen1988,stratton1941}.

\textbf{Step 1.} Set \(a_{n}=\sqrt{P_{n}}\) \cite{poynting1884}.

\textbf{Step 2.} Apply Eq.~\eqref{eq:ch3.coupling-normalization-propagation} \cite{hansen1988}.

\textbf{Step 3.} Recompute \(\mathcal{P}\) with Eq.~\eqref{eq:ch3.power-budget-invariance} \cite{harrington2001}.

\textbf{Physical meaning.} Literature values are not comparable without \(\mathcal{N}\) metadata \cite{stratton1941}.

\paragraph{Far-field pattern scaling law.}
Let \(\E_{0}(\theta,\phi)=\sum_{\ell m\sigma}\widehat{c}_{\ell m}^{(\sigma)}\mathbf{E}_{\ell m\sigma}^{(0)}(\theta,\phi)\) be the
normalized far amplitude expansion \cite{stratton1941,jackson1999}. Changing \(\widehat{c}\to a\widehat{c}\) scales
\(|\E_{0}|^{2}\) by \(|a|^{2}\) and \(\mathcal{P}_{\mathrm{rad}}\) likewise under Eq.~\eqref{eq:ch3.power-budget-invariance}
\cite{poynting1884,harrington2001}. Pattern \emph{shape} plots often hide that scaling \cite{harrington2001}.

\paragraph{Worked illustration (coupling phase for alignment).}
Two ports aligned to the same normalized mode have \(|K_{nn}|=1\) when both use identical \(\mathcal{N}\) and lossless
reciprocity holds \cite{stratton1941,harrington2001}. A \(\pi/2\) phase in \(\widehat{c}\) rotates link phase in
Eq.~\eqref{eq:ch3.link-budget-normalized} without changing \(|K_{nn}|\) \cite{hansen1988,harrington2001}.

\paragraph{Problem (verify power invariance under dual rescaling).}
\textbf{Problem.} Rescale \(\widehat{\mathbf{u}}_{m}\to a_{m}\widehat{\mathbf{u}}_{m}\) and \(\widehat{c}_{m}\to\widehat{c}_{m}/a_{m}\).
Show \(\mathcal{P}_{\mathrm{rad}}\) unchanged \cite{harrington2001,stratton1941,hansen1988}.

\textbf{Step 1.} Insert into Eq.~\eqref{eq:ch3.mode-power-from-coefficient} \cite{poynting1884}.

\textbf{Step 2.} Track \(| \widehat{c}_{m}/a_{m}|^{2}\times|a_{m}|^{2}\) \cite{hansen1988}.

\textbf{Step 3.} Conclude invariance \cite{stratton1941}.

\textbf{Physical meaning.} Physics is invariant; reports are not unless paired \cite{harrington2001}.

\begin{table}[t]
\centering
\caption{Subsection~\ref{sec:ch354} reporting metadata required for interchange.}
\label{tab:ch3.354-metadata-contract}
\small
\begin{tabular}{@{}p{0.34\linewidth}p{0.56\linewidth}@{}}
\toprule
Metadata field & Equation reference \\
\midrule
Discrete flux norm & Eq.~\eqref{eq:ch3.VSH-unit-flux-normalization} \\
Discrete energy norm & Eq.~\eqref{eq:ch3.VSH-unit-energy-norm} \\
Continuous measure & Eq.~\eqref{eq:ch3.continuous-measure-flux-fix} \\
Phase convention & Eq.~\eqref{eq:ch3.Ylm-Condon-Shortley-def} \\
\bottomrule
\end{tabular}
\end{table}

\paragraph{Operational measurement chain (restated).}
Eq.~\eqref{eq:ch3.measurement-normalized-coefficients} is the Subsection~\ref{sec:ch354} cap on
Eq.~\eqref{eq:ch3.measurement-on-coefficients}: every laboratory report of \(\widehat{c}\) must specify \(\mathcal{N}\)
\cite{harrington2001,stratton1941}. Otherwise coupling comparisons across experiments are undefined even when patterns match
\cite{hansen1988,harrington2001}.

\paragraph{Worked illustration (S-parameter analogy).}
Normalized coupling magnitudes behave like scattering amplitudes only after \(\mathcal{N}\) is shared; otherwise ratios
resemble \(S_{21}\) referenced to mismatched impedances \cite{harrington2001,stratton1941}. Eq.~\eqref{eq:ch3.reciprocity-coupling-symmetry}
is the lossless symmetry check analogous to reciprocal \(S\) matrices \cite{stratton1941,harrington2001}.

\paragraph{Truncation consequence (preview).}
Section~\ref{sec:ch36} analyzes incomplete bases; Subsection~\ref{sec:ch354} records that watt sums over truncated
\(\ell_{\max}\) are only lower bounds on full \(\mathcal{P}_{\mathrm{rad}}\) unless remainder modes are shown orthogonal
and negligible \cite{harrington2001,stratton1941,hansen1988}. Normalization must be declared on the same truncated set used
in coupling reports \cite{harrington2001}.

\paragraph{Problem (compare geometric and flux \(K\) for one port).}
\textbf{Problem.} Geometric \(K_{10}=0.5\); flux normalization uses \(a_{10}=2\). Find flux-gauge \(K'_{10}\)
\cite{harrington2001,hansen1988,stratton1941}.

\textbf{Step 1.} Apply Eq.~\eqref{eq:ch3.coupling-normalization-propagation} with \(a_{10}=2\) \cite{hansen1988}.

\textbf{Step 2.} Obtain \(K'_{10}=0.5/4=0.125\) \cite{harrington2001}.

\textbf{Physical meaning.} Coupling magnitudes shrink when modes are flux-normalized to \(1\,\mathrm{W}\) \cite{stratton1941}.

\paragraph{Chapter~\ref{ch:02} flux-meter composition (imported).}
Eq.~\eqref{eq:ch2.flux-meter-composition} chains observation with \(\mathcal{R}_\omega\); Subsection~\ref{sec:ch354}
inserts \(\mathcal{N}\) between \(\mathcal{C}\) and reported coefficients so every meter in the chain shares one watt gauge
\cite{harrington2001,stratton1941,poynting1884}. Skipping \(\mathcal{N}\) breaks power-budget invariance
Eq.~\eqref{eq:ch3.power-budget-invariance} in software pipelines \cite{harrington2001}.

\paragraph{Multi-port normalization product (HOME).}
For \(P\) ports indexed by \(p\), let \(a_{p,n}\) be the flux-gauge scale of mode \(n\) at port \(p\). Coupling transforms as
\begin{equation}
\label{eq:ch3.multiport-normalization-product}
K'_{mn}
=
\frac{K_{mn}}{a_{\mathrm{tx},m}\,a_{\mathrm{rx},n}},
\end{equation}
generalizing Eq.~\eqref{eq:ch3.coupling-normalization-propagation} to distinct transmit and receive libraries
\cite{harrington2001,stratton1941}. Equation~\eqref{eq:ch3.multiport-normalization-product} is the interchange rule for
multi-antenna campaigns: each port must publish its \(\mathcal{N}\) metadata \cite{hansen1988,harrington2001}.

\paragraph{Worked illustration (numeric link budget).}
Take \(\widehat{c}_{\mathrm{tx}}=1\) (one watt transmit), \(|K_{n,\mathrm{tx}}|^{2}=0.1\) for a single receive mode \(n\), and
\(\widehat{c}_{n}^{\mathrm{rx}}=1\). Eq.~\eqref{eq:ch3.link-budget-normalized} gives \(P_{\mathrm{rx}}=0.1\,\mathrm{W}\)
\cite{harrington2001,poynting1884,stratton1941}. Halving transmit normalization scale doubles reported \(\widehat{c}_{\mathrm{tx}}\)
but leaves \(P_{\mathrm{rx}}\) unchanged when \(K\) is updated through Eq.~\eqref{eq:ch3.multiport-normalization-product}
\cite{harrington2001}.

\begin{figure}[t]
\centering
\ChOneFigBlock{%
\begin{tikzpicture}[ch1fig, node distance=7mm and 9mm]
  \node[ch1box] (p1) {Port 1 $\mathcal{N}_{1}$};
  \node[ch1box, right=10mm of p1] (k) {Eq.~\eqref{eq:ch3.multiport-normalization-product}};
  \node[ch1box, right=10mm of k] (p2) {Port 2 $\mathcal{N}_{2}$};
  \node[ch1box, below=10mm of k] (lb) {Eq.~\eqref{eq:ch3.link-budget-normalized}};
  \draw[ch1flow] (p1) -- (k) -- (p2);
  \draw[ch1flow] (lb) -- (k);
\end{tikzpicture}%
}
\caption{Subsection~\ref{sec:ch354} multi-port gauge: each port carries its own \(\mathcal{N}\) before coupling assembly.}
\label{fig:ch3.354-multiport-gauge}
\end{figure}

\paragraph{Impedance-analogy caution (no re-derivation).}
Transmission-line \(S\) parameters require common reference impedances \cite{harrington2001,collin2001}. Modal \(K_{mn}\) under
Eq.~\eqref{eq:ch3.coupling-normalization-propagation} obeys the same logical rule: ratios are meaningless when ports use
different \(\mathcal{N}\) \cite{stratton1941,harrington2001}. Subsection~\ref{sec:ch354} treats \(\mathcal{N}\) as the modal
reference impedance for watt-conserving algebra \cite{harrington2001}.

\paragraph{Problem (three-port audit).}
\textbf{Problem.} Ports A, B, C publish \(a_{A}=1\), \(a_{B}=2\), \(a_{C}=\sqrt{2}\). A stored coupling \(K_{BC}=0.4\) is in
geometric gauge. Find flux-gauge \(K'_{BC}\) \cite{harrington2001,hansen1988,stratton1941}.

\textbf{Step 1.} Identify \(a_{\mathrm{tx},B}=2\), \(a_{\mathrm{rx},C}=\sqrt{2}\) \cite{hansen1988}.

\textbf{Step 2.} Apply Eq.~\eqref{eq:ch3.multiport-normalization-product} \cite{harrington2001}.

\textbf{Step 3.} Obtain \(K'_{BC}=0.4/(2\sqrt{2})\approx0.141\) \cite{stratton1941}.

\textbf{Physical meaning.} Each port scale divides once in the coupling denominator \cite{poynting1884}.

\begin{table}[t]
\centering
\caption{Subsection~\ref{sec:ch354} numeric link-budget template (matched \(\mathcal{N}\)).}
\label{tab:ch3.354-link-budget-numeric}
\small
\begin{tabular}{@{}p{0.22\linewidth}p{0.28\linewidth}p{0.22\linewidth}p{0.20\linewidth}@{}}
\toprule
Quantity & Symbol & Value & Unit \\
\midrule
TX coefficient & $\widehat{c}_{\mathrm{tx}}$ & $1$ & $\sqrt{\mathrm{W}}$ \\
Coupling power & $|K|^{2}$ & $0.1$ & dimensionless \\
RX coefficient & $\widehat{c}_{\mathrm{rx}}$ & $1$ & $\sqrt{\mathrm{W}}$ \\
Received power & $P_{\mathrm{rx}}$ & $0.1$ & W \\
\bottomrule
\end{tabular}
\end{table}

\paragraph{Continuous--discrete consequence split.}
When \(K\) mixes discrete \(\widehat{\mathbf{u}}_{\ell m}^{(\sigma)}\) with continuum \(c(\hat{\mathbf{k}})\), both sides must
declare Eqs.~\eqref{eq:ch3.VSH-unit-flux-normalization} and \eqref{eq:ch3.continuous-measure-flux-fix} separately
\cite{harrington2001,stratton1941}. Eq.~\eqref{eq:ch3.power-budget-invariance} holds only after the hybrid \(\mathcal{N}\) is
specified on each factor \cite{hansen1988,harrington2001}.

\paragraph{Receive-only inversion consequence.}
Far-field sampling that estimates \(\widehat{c}\) without transmitting a calibrated watt reference can fix pattern shape but leaves
overall scale ambiguous by a factor \(|a|\) until Eq.~\eqref{eq:ch3.shell-power-normalized-mode} is enforced on a known excitation
\cite{devaney2004,harrington2001,stratton1941}. Subsection~\ref{sec:ch354} therefore pairs inversion outputs with at least one
flux-anchored normalization measurement \cite{harrington2001,poynting1884}.

\paragraph{Problem (hybrid link budget with mismatched \(P_{0}\)).}
\textbf{Problem.} Transmit side uses Eq.~\eqref{eq:ch3.VSH-unit-flux-normalization}; receive continuum uses \(P_{0}=2\,\mathrm{W}\) in
Eq.~\eqref{eq:ch3.continuous-measure-flux-fix}. A coupling entry \(K_{d,c}=0.2\) is reported. Rescale \(K\) to a common
\(P_{0}=1\,\mathrm{W}\) convention \cite{harrington2001,hansen1988,stratton1941}.

\textbf{Step 1.} Continuum coefficients scale as \(1/\sqrt{P_{0}}\) \cite{hansen1988}.

\textbf{Step 2.} Ratio \(\sqrt{2/1}\) adjusts the continuum leg \cite{harrington2001}.

\textbf{Step 3.} Obtain \(K'_{d,c}=0.2/\sqrt{2}\approx0.141\) \cite{stratton1941}.

\textbf{Physical meaning.} Hybrid campaigns require explicit \(P_{0}\) on every continuum factor \cite{poynting1884}.

Subsection~\ref{sec:ch354} has recorded how flux and measure normalizations propagate through coupling, power, and
reciprocity reports \cite{stratton1941,harrington2001,poynting1884}. Section~\ref{sec:ch35} is complete: discrete and
continuous amplitudes are watt-calibrated, DOS is counted, and downstream consequences are fixed for completeness and
truncation analysis in Section~\ref{sec:ch36} \cite{hansen1988,jackson1999}.

The four subsections of Section~\ref{sec:ch35} close the amplitude layer of Chapter~\ref{ch:03}: flux normalization
Eqs.~\eqref{eq:ch3.VSH-unit-flux-normalization}--\eqref{eq:ch3.mode-power-from-coefficient} in Subsection~\ref{sec:ch351};
continuous measure Eqs.~\eqref{eq:ch3.continuous-measure-flux-fix}--\eqref{eq:ch3.continuous-parseval-flux-fixed} in
Subsection~\ref{sec:ch352}; DOS Eqs.~\eqref{eq:ch3.multiplet-multiplicity-def}--\eqref{eq:ch3.hybrid-DOS-count} in
Subsection~\ref{sec:ch353}; and coupling consequences Eqs.~\eqref{eq:ch3.coupling-normalization-propagation}--\eqref{eq:ch3.reciprocity-coupling-symmetry}
in Subsection~\ref{sec:ch354} \cite{stratton1941,poynting1884,harrington2001}. Section~\ref{sec:ch36} now asks when the
normalized basis is complete and how completeness fails under truncation and symmetry breaking \cite{harrington2001}.

\section{Completeness, Truncation, and Symmetry Breaking}
\label{sec:ch36}

Complete bases are convenient; finite laboratories never enjoy them. Section~\ref{sec:ch36} explains when
angular-momentum eigenfields form complete sets and how completeness fails under truncation, boundaries, and
symmetry breaking---the representation-theoretic prelude to realizability in Chapter~\ref{ch:04}
\cite{harrington2001,stratton1941}. Subsection~\ref{sec:ch361} states completeness on unbounded domains using
Subsection~\ref{sec:ch343}. Subsection~\ref{sec:ch362} analyzes truncation and loss of completeness when
sources or observation supports are compact. Subsection~\ref{sec:ch363} treats boundary-induced mode mixing.
Subsection~\ref{sec:ch364} classifies accessible structure when symmetry is reduced and links the taxonomy to
null-space and accessibility screens from Subsection~\ref{sec:ch274} \cite{devaney2004,harrington2001}.


\subsection{Completeness in unbounded domains}
\label{sec:ch361}
% TOC tag: [HOME][USE SS3.4.3]

\emph{Completeness} means every admissible radiative state is a convergent superposition of basis modes;
\emph{unbounded domains} mean the exterior problem on \(\mathbb{R}^{3}\) (or an exterior region) without
enclosing conductors that quantize the spectrum \cite{stratton1941,jackson1999,hansen1988}. Subsection~\ref{sec:ch361}
is the HOME certificate layer that upgrades Theorem~\ref{thm:ch3.vsh-completeness} from Subsection~\ref{sec:ch343}
to the normalized watt-calibrated basis of Section~\ref{sec:ch35} \cite{harrington2001,kato1995}.

\paragraph{Standing imports (no re-proof).}
Hilbert closure Eq.~\eqref{eq:ch3.F-rad-hilbert-closure}, radiative completeness characterization
Eq.~\eqref{eq:ch3.completeness-radiative-characterization}, VSH resolution
Eq.~\eqref{eq:ch3.VSH-resolution-identity}, orthogonality Theorem~\ref{thm:ch3.vsh-orthogonality}, and
Theorem~\ref{thm:ch3.vsh-completeness} are cited from Subsection~\ref{sec:ch343} \cite{hansen1988,stratton1941}.
Subsection~\ref{sec:ch361} does not re-derive spherical harmonic completeness or Hankel asymptotics
\cite{kato1995,hansen1988}.

\paragraph{Assumptions.}
Work on outgoing \(\mathcal{F}_{\mathrm{rad}}\) with inheritance
Eq.~\eqref{eq:ch3.outgoing-inheritance-criterion}, reciprocal lossless exterior media, and single-frequency
phasors \(\exp(-\jj\omega t)\) \cite{stratton1941,harrington2001}. Unbounded completeness is an
\emph{infinite}-rank statement: every \(\mathbf{v}_\omega\in\mathcal{F}_{\mathrm{rad}}\) is a limit of finite
VSH sums \cite{kato1995,hansen1988}.

\paragraph{Unbounded-domain completeness statement (HOME).}
Let \(\{\widehat{\mathbf{u}}_{\ell m}^{(\sigma)}\}\) be flux-normalized outgoing pairs from
Eqs.~\eqref{eq:ch3.VSH-unit-flux-normalization}--\eqref{eq:ch3.VSH-unit-energy-norm}. Then
\begin{equation}
\label{eq:ch3.unbounded-completeness-statement}
\overline{\operatorname{span}\{\widehat{\mathbf{u}}_{\ell m}^{(\sigma)}\}}
=
\mathcal{F}_{\mathrm{rad}},
\qquad
\text{closure in }\|\cdot\|_{\mathrm{rad}},
\end{equation}
importing Theorem~\ref{thm:ch3.vsh-completeness} and Eq.~\eqref{eq:ch3.completeness-radiative-characterization}
\cite{hansen1988,kato1995,stratton1941}. Equation~\eqref{eq:ch3.unbounded-completeness-statement} is the
Subsection~\ref{sec:ch361} physical headline: on free-space exterior problems, no radiative degree of freedom lies
\emph{outside} the angular-momentum catalog once outgoing class and reciprocity are fixed \cite{harrington2001}.

\paragraph{\(\mathcal{L}^{2}\) convergence of VSH sums (HOME).}
For \(\mathbf{v}_\omega\in\mathcal{F}_{\mathrm{rad}}\), the normalized expansion
Eq.~\eqref{eq:ch3.normalized-coefficient-resolution} converges as
\begin{equation}
\label{eq:ch3.L2-convergence-VSH}
\lim_{\ell_{\max}\to\infty}
\Big\|
\mathbf{v}_\omega
-
\sum_{\ell\le\ell_{\max}}\sum_{m,\sigma}
\widehat{c}_{\ell m}^{(\sigma)}\,\widehat{\mathbf{u}}_{\ell m}^{(\sigma)}
\Big\|_{\mathrm{rad}}
=
0,
\end{equation}
with coefficients \(\widehat{c}_{\ell m}^{(\sigma)}=\langle\widehat{\mathbf{u}}_{\ell m}^{(\sigma)},\mathbf{v}_\omega\rangle_{\mathrm{rad}}\)
\cite{hansen1988,kato1995,stratton1941}. Equation~\eqref{eq:ch3.L2-convergence-VSH} separates \emph{mathematical}
completeness (the limit exists) from \emph{numerical} truncation (finite \(\ell_{\max}\) in Subsection~\ref{sec:ch362})
\cite{harrington2001}.

\begin{theorem}[Unbounded radiative completeness (HOME restatement)]
\label{thm:ch3.unbounded-domain-completeness}
Under Eqs.~\eqref{eq:ch3.outgoing-inheritance-criterion}--\eqref{eq:ch3.radiative-inner-product-def} on exterior
domains, the normalized VSH family is complete in \(\mathcal{F}_{\mathrm{rad}}\) and every
\(\mathcal{R}_\omega[\mathfrak{D}]\) admits expansion Eq.~\eqref{eq:ch3.normalized-coefficient-resolution} with
convergence Eq.~\eqref{eq:ch3.L2-convergence-VSH} \cite{hansen1988,stratton1941,kato1995}.
\end{theorem}
\begin{proof}[Proof sketch]
Invoke Theorem~\ref{thm:ch3.vsh-completeness} and normalization isometry on diagonal norms from
Section~\ref{sec:ch35} \cite{hansen1988,harrington2001}.
% PROOF: AUTHOR REVIEW
\end{proof}

\begin{figure}[t]
\centering
\ChOneFigBlock{%
\begin{tikzpicture}[ch1fig, node distance=7mm and 9mm]
  \node[ch1box] (frad) {$\mathcal{F}_{\mathrm{rad}}$};
  \node[ch1box, right=11mm of frad] (span) {Eq.~\eqref{eq:ch3.unbounded-completeness-statement}};
  \node[ch1box, right=11mm of span] (conv) {Eq.~\eqref{eq:ch3.L2-convergence-VSH}};
  \node[ch1box, below=10mm of span] (res) {Eq.~\eqref{eq:ch3.VSH-resolution-identity}};
  \draw[ch1flow] (res) -- (span) -- (conv);
  \draw[ch1flow] (frad) -- (span);
\end{tikzpicture}%
}
\caption{Subsection~\ref{sec:ch361} completeness chain on unbounded domains: resolution identity, dense span, and
\(\|\cdot\|_{\mathrm{rad}}\) convergence.}
\label{fig:ch3.361-completeness-chain}
\end{figure}

\paragraph{Physical meaning of completeness.}
Completeness guarantees that \emph{shape} residuals in far-zone patterns are not mysterious ``missing physics'' on
exterior domains: any square-integrable radiative pattern is synthesizable by summing \(\ell,m,\sigma\) channels
\cite{stratton1941,jackson1999}. The certificate is operator-theoretic---it concerns
\(\operatorname{ran}\mathcal{R}_\omega\) inside \(\mathcal{F}_{\mathrm{rad}}\)---not a claim that every laboratory
antenna physically excites all \(\ell\) \cite{harrington2001}.

\paragraph{Worked illustration (dipole-plus-quadrupole synthesis).}
A target \(\mathbf{v}_\omega\) with dominant \(\widehat{c}_{10}^{(\mathbf{N})}\) and small
\(\widehat{c}_{20}^{(\mathbf{N})}\) is \emph{exactly} representable as \(\widehat{c}_{10}\widehat{\mathbf{u}}_{10}^{(\mathbf{N})}
+\widehat{c}_{20}\widehat{\mathbf{u}}_{20}^{(\mathbf{N})}+\cdots\) with convergence Eq.~\eqref{eq:ch3.L2-convergence-VSH}
\cite{hansen1988,stratton1941}. Finite arrays that keep only \(\ell\le1\) incur truncation error analyzed in
Subsection~\ref{sec:ch362}, not incompleteness of the full catalog \cite{harrington2001}.

\paragraph{Problem (verify Parseval on complete basis).}
\textbf{Problem.} For \(\mathbf{v}_\omega=\sum_{\ell m\sigma}\widehat{c}_{\ell m}^{(\sigma)}\widehat{\mathbf{u}}_{\ell m}^{(\sigma)}\)
with convergence Eq.~\eqref{eq:ch3.L2-convergence-VSH}, show
\(\|\mathbf{v}_\omega\|_{\mathrm{rad}}^{2}=\sum_{\ell m\sigma}|\widehat{c}_{\ell m}^{(\sigma)}|^{2}\)
\cite{hansen1988,kato1995,harrington2001}.

\textbf{Step 1.} Apply Theorem~\ref{thm:ch3.vsh-orthogonality} termwise \cite{hansen1988}.

\textbf{Step 2.} Pass to the limit \(\ell_{\max}\to\infty\) using Eq.~\eqref{eq:ch3.L2-convergence-VSH} \cite{kato1995}.

\textbf{Step 3.} Read normalized energy norm Eq.~\eqref{eq:ch3.VSH-unit-energy-norm} \cite{stratton1941}.

\textbf{Physical meaning.} Complete bases permit infinite Parseval identities; watt sums are finite truncations
\cite{harrington2001,poynting1884}.

\paragraph{Continuous spectrum on open domains.}
When discrete \(\ell\) labels do not close the space---previewed in Eq.~\eqref{eq:ch3.continuous-eigenfield-expansion}---the
continuous direction integral supplements discrete multiplets on unbounded domains \cite{stratton1941,jackson1999}.
Subsection~\ref{sec:ch361} records that discrete VSH completeness and continuous direction measures coexist in the
rigged frame of Subsection~\ref{sec:ch313} \cite{harrington2001,kato1995}.

\paragraph{Separation from operator nullity.}
Completeness addresses \emph{representation} of \(\mathcal{F}_{\mathrm{rad}}\); \(\ker\mathcal{R}_\omega\) from
Eq.~\eqref{eq:ch2.ker-R-omega-preview} addresses \emph{export} \cite{harrington2001,stratton1941}. A non-radiating
source lies outside the completeness question because it never enters \(\mathcal{F}_{\mathrm{rad}}\) through
\(\mathcal{R}_\omega\) \cite{harrington2001}.

\begin{table}[t]
\centering
\caption{Subsection~\ref{sec:ch361} completeness objects on unbounded domains.}
\label{tab:ch3.361-completeness-objects}
\small
\begin{tabular}{@{}p{0.24\linewidth}p{0.38\linewidth}p{0.30\linewidth}@{}}
\toprule
Concept & Equation & Meaning \\
\midrule
Dense span & Eq.~\eqref{eq:ch3.unbounded-completeness-statement} & no missing radiative modes \\
Convergence & Eq.~\eqref{eq:ch3.L2-convergence-VSH} & \(\ell_{\max}\to\infty\) limit \\
Coefficients & Eq.~\eqref{eq:ch3.normalized-coefficient-resolution} & projections, not guesses \\
\bottomrule
\end{tabular}
\end{table}

\paragraph{Far-zone equivalence.}
Completeness in \(\|\cdot\|_{\mathrm{rad}}\) implies completeness of far amplitudes
\(\mathbf{E}_{0}(\theta,\phi)\) in the pattern norm induced by Eq.~\eqref{eq:ch2.far-zone-power-from-pattern}
\cite{stratton1941,jackson1999}. Subsection~\ref{sec:ch361} therefore certifies pattern synthesis on \(S^{2}\) for
outgoing transverse fields \cite{hansen1988,harrington2001}.

\paragraph{Worked illustration (monopole forbidden on exterior).}
Pure \(\ell=0\) transverse radiation is absent in \(\mathcal{F}_{\mathrm{rad}}\); completeness is among
\(\ell\ge1\) sectors certified in Subsection~\ref{sec:ch343} \cite{stratton1941,jackson1999}. One must not
interpret Eq.~\eqref{eq:ch3.unbounded-completeness-statement} as ``every \(\ell\) is physically excited,'' but as
``every \emph{admissible} radiative harmonic is representable'' \cite{harrington2001}.

\paragraph{Problem (diagnose non-convergent partial sum).}
\textbf{Problem.} A partial sum to \(\ell_{\max}=5\) oscillates as \(\ell_{\max}\) increases on \(\Gamma_R\).
Distinguish incompleteness from wrong branch \cite{stratton1941,hansen1988,harrington2001}.

\textbf{Step 1.} Test outgoing criterion Eq.~\eqref{eq:ch3.outgoing-inheritance-criterion} \cite{stratton1941}.

\textbf{Step 2.} Verify coefficients arise from projections Eq.~\eqref{eq:ch3.normalized-coefficient-resolution}
\cite{hansen1988}.

\textbf{Step 3.} If branch is correct, increase \(\ell_{\max}\) toward Eq.~\eqref{eq:ch3.L2-convergence-VSH}
\cite{kato1995}.

\textbf{Physical meaning.} Non-convergence usually signals reactive admixture or wrong Hankel sheet, not basis defect
\cite{harrington2001,stratton1941}.

\paragraph{Reciprocity compatibility.}
Completeness is stated with reciprocity-symmetric inner product
Eq.~\eqref{eq:ch3.reciprocity-symmetric-inner-product}; lossy media require biorthogonal pairings deferred beyond
Volume~I \cite{stratton1941,harrington2001}. Subsection~\ref{sec:ch361} fixes the lossless exterior class used in
Chapters~\ref{ch:02}--\ref{ch:03} \cite{hansen1988}.

\paragraph{Rigged triple alignment.}
Completeness on \(\mathcal{F}_{\mathrm{rad}}\) is compatible with Gelfand triple embeddings
Eqs.~\eqref{eq:ch3.gelfand-triple-rad}--\eqref{eq:ch3.distributional-pairing-extension} because smooth test functions
are dense in the radiative energy norm \cite{kato1995,stratton1941}. Subsection~\ref{sec:ch361} therefore aligns
discrete VSH completeness with continuous measures of Section~\ref{sec:ch35} \cite{harrington2001}.

\paragraph{Worked illustration (plane-wave as completeness witness).}
A single plane-wave mode lies in the span of VSH sums by Eq.~\eqref{eq:ch3.plane-wave-continuum-bridge} and
Eq.~\eqref{eq:ch3.L2-convergence-VSH} \cite{jackson1999,hansen1988}. Plane-wave beams are not ``outside'' the
angular-momentum basis; they are different coordinates on the same complete space \cite{stratton1941,harrington2001}.

\paragraph{Validity boundary.}
Ground planes, cavities, and inhomogeneous exteriors break the unbounded symmetry assumptions underlying
Theorem~\ref{thm:ch3.vsh-completeness}; those regimes are analyzed in Subsections~\ref{sec:ch362}--\ref{sec:ch364}
\cite{harrington2001,stratton1941}. Eq.~\eqref{eq:ch3.unbounded-completeness-statement} is the \emph{reference}
complete world against which laboratory incompleteness is measured \cite{hansen1988}.

\paragraph{Shell-wise completeness witness (HOME corollary).}
On any far shell \(\Gamma_R\), completeness implies that spherical harmonic expansions of
\(\mathbf{E}_{t}(\theta,\phi)\) converge in mean square:
\begin{equation}
\label{eq:ch3.shell-harmonic-completeness}
\lim_{\ell_{\max}\to\infty}
\int_{\Gamma_R}
\big|\mathbf{E}_{t}-\mathbf{E}_{t}^{(\ell_{\max})}\big|^{2}\,dS
=
0,
\end{equation}
with \(\mathbf{E}_{t}^{(\ell_{\max})}\) the truncated VSH synthesis \cite{hansen1988,stratton1941,jackson1999}.
Equation~\eqref{eq:ch3.shell-harmonic-completeness} is the measurable form of
Eq.~\eqref{eq:ch3.L2-convergence-VSH}: anechoic chambers test completeness through \(\Gamma_R\) integrals, not abstract
norms \cite{harrington2001,poynting1884}.

\paragraph{Worked illustration (synthesize a skew beam).}
A target beam tilted \(15^{\circ}\) from broadside requires admixture of multiple \(m\) sectors at fixed \(\ell\)
through Eq.~\eqref{eq:ch3.eigenfield-Wigner-mixing} \cite{jackson1999,hansen1988}. Completeness guarantees a convergent
sum exists; it does not guarantee a \emph{short} sum---large \(\ell_{\max}\) may be required for tight beams
\cite{harrington2001,stratton1941}.

\begin{figure}[t]
\centering
\ChOneFigBlock{%
\begin{tikzpicture}[ch1fig, node distance=7mm and 9mm]
  \node[ch1box] (r3) {$\mathbb{R}^{3}$ exterior};
  \node[ch1box, right=10mm of r3] (frad) {$\mathcal{F}_{\mathrm{rad}}$};
  \node[ch1box, right=10mm of frad] (span) {Eq.~\eqref{eq:ch3.unbounded-completeness-statement}};
  \node[ch1box, below=10mm of frad] (shell) {Eq.~\eqref{eq:ch3.shell-harmonic-completeness}};
  \draw[ch1flow] (r3) -- (frad) -- (span);
  \draw[ch1flow] (shell) -- (span);
\end{tikzpicture}%
}
\caption{Subsection~\ref{sec:ch361} domain picture: unbounded exterior domains host a complete outgoing span tested on
far shells.}
\label{fig:ch3.361-domain-complete}
\end{figure}

\paragraph{Problem (pointwise versus mean-square convergence).}
\textbf{Problem.} One expects \(\mathbf{E}_{t}^{(\ell_{\max})}\to\mathbf{E}_{t}\) pointwise on \(\Gamma_R\). State the
correct convergence mode \cite{hansen1988,kato1995,harrington2001}.

\textbf{Step 1.} Cite Eq.~\eqref{eq:ch3.shell-harmonic-completeness} \cite{hansen1988}.

\textbf{Step 2.} Note Gibbs-like overshoot can occur near sharp nulls even when mean-square error vanishes \cite{jackson1999}.

\textbf{Step 3.} Recommend mean-square flux metrics for pass/fail \cite{poynting1884,harrington2001}.

\textbf{Physical meaning.} Completeness is an \(\mathcal{L}^{2}\) statement on radiative energy, not uniform pointwise
convergence \cite{kato1995,stratton1941}.

\paragraph{Hansen--Stratton closure (imported).}
Classical vector spherical harmonic treatments establish completeness of outgoing sets on exterior domains
\cite{hansen1988,stratton1941}. Subsection~\ref{sec:ch361} restates that closure in the \(\mathcal{F}_{\mathrm{rad}}\) inner
product and watt normalization of Section~\ref{sec:ch35} \cite{harrington2001,poynting1884}.

\begin{table}[t]
\centering
\caption{Subsection~\ref{sec:ch361} domain completeness comparison.}
\label{tab:ch3.361-domain-comparison}
\small
\begin{tabular}{@{}p{0.26\linewidth}p{0.34\linewidth}p{0.32\linewidth}@{}}
\toprule
Domain & Completeness & Section \\
\midrule
\(\mathbb{R}^{3}\) exterior & Eq.~\eqref{eq:ch3.unbounded-completeness-statement} & \ref{sec:ch361} \\
Half-space PEC & broken & \ref{sec:ch363} \\
Finite \(\ell_{\max}\) & lost & \ref{sec:ch362} \\
\bottomrule
\end{tabular}
\end{table}

\paragraph{Coefficient uniqueness under completeness.}
When Eq.~\eqref{eq:ch3.unbounded-completeness-statement} holds, normalized coefficients in
Eq.~\eqref{eq:ch3.normalized-coefficient-resolution} are unique for each \(\mathbf{v}_\omega\in\mathcal{F}_{\mathrm{rad}}\)
\cite{kato1995,hansen1988}. Non-uniqueness at finite \(\ell_{\max}\) is truncation ambiguity, not gauge freedom
\cite{harrington2001,stratton1941}.

\paragraph{Worked illustration (continuous spectrum coexistence).}
Completeness of discrete VSH does not eliminate the continuous direction integral of
Eq.~\eqref{eq:ch3.continuous-eigenfield-expansion}; both descriptions can represent the same \(\mathbf{v}_\omega\)
\cite{stratton1941,jackson1999}. Subsection~\ref{sec:ch361} certifies the discrete span is dense; Section~\ref{sec:ch37}
compares discrete and continuous coordinates \cite{hansen1988,harrington2001}.

\paragraph{Export pipeline closure.}
Every \(\mathcal{R}_\omega[\mathfrak{D}]\in\mathcal{F}_{\mathrm{rad}}\) admits coefficients \(\widehat{c}_{\ell m}^{(\sigma)}\)
through Eqs.~\eqref{eq:ch3.normalized-coefficient-resolution} and \eqref{eq:ch3.L2-convergence-VSH} on unbounded domains
\cite{hansen1988,kato1995,stratton1941}. Subsection~\ref{sec:ch361} is therefore the last \emph{existence} certificate before
truncation and boundary caveats in Subsections~\ref{sec:ch362}--\ref{sec:ch364} \cite{harrington2001}.

\paragraph{Problem (finite shell versus unbounded limit).}
\textbf{Problem.} Show Eq.~\eqref{eq:ch3.shell-harmonic-completeness} approaches Eq.~\eqref{eq:ch3.L2-convergence-VSH} as
\(R\to\infty\) with fixed \(\ell_{\max}\) increasing \cite{hansen1988,kato1995,harrington2001}.

\textbf{Step 1.} Use far-zone asymptotics Eq.~\eqref{eq:ch3.hankel-far-zone-asymptotic} on \(\Gamma_R\) \cite{stratton1941}.

\textbf{Step 2.} Let shell integrals approximate \(\|\cdot\|_{\mathrm{rad}}\) as \(k_{0}R\gg\ell_{\max}\) \cite{hansen1988}.

\textbf{Step 3.} Conclude shell tests are practical proxies for completeness when \(R\) is large enough \cite{harrington2001}.

\textbf{Physical meaning.} Laboratories verify completeness on finite spheres, not abstract limits \cite{poynting1884,stratton1941}.

\paragraph{Dense span operational test.}
Given \(\epsilon>0\), choose \(\ell_{\max}\) such that \(\|\mathbf{v}_\omega-\mathcal{T}_{\ell_{\max}}\mathbf{v}_\omega\|_{\mathrm{rad}}<\epsilon\)
for every target pattern in a campaign library \cite{kato1995,harrington2001}. Subsection~\ref{sec:ch361} makes that test the
acceptance criterion for ``complete representation'' in software pipelines \cite{hansen1988,stratton1941}.

\begin{figure}[t]
\centering
\ChOneFigBlock{%
\begin{tikzpicture}[ch1fig, node distance=7mm and 9mm]
  \node[ch1box] (lib) {pattern library};
  \node[ch1box, right=10mm of lib] (conv) {Eq.~\eqref{eq:ch3.L2-convergence-VSH}};
  \node[ch1box, right=10mm of conv] (pass) {$\epsilon$ pass};
  \draw[ch1flow] (lib) -- (conv) -- (pass);
\end{tikzpicture}%
}
\caption{Subsection~\ref{sec:ch361} acceptance flow: library patterns must satisfy \(\mathcal{L}^{2}\) convergence on
\(\mathcal{F}_{\mathrm{rad}}\).}
\label{fig:ch3.361-acceptance-flow}
\end{figure}

\paragraph{Worked illustration (catalog coverage test).}
A calibration library of \(50\) far patterns with distinct sidelobe textures can be expanded in VSH until each satisfies
\(\epsilon=10^{-3}\) in \(\|\cdot\|_{\mathrm{rad}}\) \cite{kato1995,harrington2001}. Failure on a single library member
signals incomplete outgoing classification or reactive contamination, not basis defect \cite{stratton1941,hansen1988}.

\paragraph{Completeness and uniqueness of angular power.}
When Eq.~\eqref{eq:ch3.unbounded-completeness-statement} holds, angular power partitions from
Eq.~\eqref{eq:ch3.mode-power-from-coefficient} are unique in the normalized gauge \cite{poynting1884,harrington2001}.
Subsection~\ref{sec:ch361} therefore underwrites watt bookkeeping in Sections~\ref{sec:ch35}--\ref{sec:ch36} on exterior
domains \cite{stratton1941,hansen1988}.

\begin{table}[t]
\centering
\caption{Subsection~\ref{sec:ch361} completeness acceptance checklist on \(\Gamma_R\).}
\label{tab:ch3.361-acceptance-checklist}
\small
\begin{tabular}{@{}p{0.08\linewidth}p{0.44\linewidth}p{0.36\linewidth}@{}}
\toprule
Step & Test & Pass criterion \\
\midrule
1 & Outgoing class & Eq.~\eqref{eq:ch3.outgoing-inheritance-criterion} \\
2 & Shell convergence & Eq.~\eqref{eq:ch3.shell-harmonic-completeness} \\
3 & Energy norm & Eq.~\eqref{eq:ch3.VSH-unit-energy-norm} \\
4 & Library \(\epsilon\) & Eq.~\eqref{eq:ch3.L2-convergence-VSH} \\
\bottomrule
\end{tabular}
\end{table}

\paragraph{Problem (library pass/fail).}
\textbf{Problem.} One library pattern fails the \(\epsilon\) test while others pass. Diagnose \cite{harrington2001,hansen1988,stratton1941}.

\textbf{Step 1.} Inspect reactive admixture on that pattern \cite{harrington2001}.

\textbf{Step 2.} Raise \(\ell_{\max}\) locally for the failing member \cite{kato1995}.

\textbf{Step 3.} If failure persists, classify outgoing violation \cite{stratton1941}.

\textbf{Physical meaning.} Completeness is global; failures are usually local classification errors \cite{hansen1988}.

Subsection~\ref{sec:ch361} has certified completeness on unbounded exterior domains via
Eqs.~\eqref{eq:ch3.unbounded-completeness-statement}--\eqref{eq:ch3.L2-convergence-VSH}
\cite{hansen1988,kato1995,stratton1941}. Subsection~\ref{sec:ch362} quantifies how finite \(\ell_{\max}\) and compact
supports destroy that completeness in practice \cite{harrington2001}.


\subsection{Truncation and loss of completeness}
\label{sec:ch362}
% TOC tag: [HOME]

\emph{Truncation} replaces an infinite index set by a finite cut \(\ell\le\ell_{\max}\);
\emph{loss of completeness} means the finite sum span is a proper subspace of \(\mathcal{F}_{\mathrm{rad}}\), so
pattern residuals cannot be driven to zero by increasing local sources within the cut \cite{harrington2001,stratton1941}.
Subsection~\ref{sec:ch362} is the HOME error calculus for that loss, importing remainder preview
Eq.~\eqref{eq:ch3.eigenfield-truncation-remainder} and compression bounds from Chapter~\ref{ch:02}
\cite{hansen1988,harrington2001}.

\paragraph{Standing imports.}
Remainder definition Eq.~\eqref{eq:ch3.eigenfield-truncation-remainder}, compression error
Eq.~\eqref{eq:ch2.compression-error-bound}, reconstruction rank cap Eq.~\eqref{eq:ch2.reconstruction-rank-cap},
and accessible angular budget Eq.~\eqref{eq:ch2.accessible-angular-budget} enter by reference
\cite{harrington2001,stratton1941}. Subsection~\ref{sec:ch362} does not re-derive singular-value tail bounds
\cite{harrington2001}.

\paragraph{Assumptions.}
Fix outgoing \(\mathcal{F}_{\mathrm{rad}}\), normalized modes
Eqs.~\eqref{eq:ch3.VSH-unit-flux-normalization}--\eqref{eq:ch3.VSH-unit-energy-norm}, and a truncation level
\(\ell_{\max}\in\mathbb{N}_{0}\) \cite{hansen1988,stratton1941}. Compact source support \(\Omega_{s}\) and finite
observation apertures may further cap recoverable coefficients \cite{harrington2001,jackson1999}.

\paragraph{Finite-\(\ell\) projection operator (HOME).}
Define the truncating projector on coefficient space by
\begin{equation}
\label{eq:ch3.ellmax-truncation-projector}
(\mathcal{T}_{\ell_{\max}}\mathbf{v})_\omega
=
\sum_{\ell\le\ell_{\max}}\sum_{m,\sigma}
\langle\widehat{\mathbf{u}}_{\ell m}^{(\sigma)},\mathbf{v}_\omega\rangle_{\mathrm{rad}}\,
\widehat{\mathbf{u}}_{\ell m}^{(\sigma)}.
\end{equation}
Equation~\eqref{eq:ch3.ellmax-truncation-projector} is the Subsection~\ref{sec:ch362} mathematical model of ``keep
multipoles through \(\ell_{\max}\) only'' \cite{hansen1988,harrington2001}. For \(\ell_{\max}<\infty\),
\(\mathcal{T}_{\ell_{\max}}\mathcal{F}_{\mathrm{rad}}\subsetneq\mathcal{F}_{\mathrm{rad}}\): completeness is lost
\cite{kato1995,stratton1941}.

\paragraph{Truncation remainder energy (HOME).}
The radiative remainder field is
\begin{equation}
\label{eq:ch3.truncation-remainder-energy}
\mathbf{R}_{\ell_{\max}}
=
\mathbf{v}_\omega-\mathcal{T}_{\ell_{\max}}\mathbf{v}_\omega,
\qquad
\|\mathbf{R}_{\ell_{\max}}\|_{\mathrm{rad}}^{2}
=
\sum_{\ell>\ell_{\max}}\sum_{m,\sigma}|\widehat{c}_{\ell m}^{(\sigma)}|^{2},
\end{equation}
by Theorem~\ref{thm:ch3.vsh-orthogonality} and normalized Parseval
Eq.~\eqref{eq:ch3.mode-power-from-coefficient} \cite{hansen1988,harrington2001,poynting1884}. Equation
\eqref{eq:ch3.truncation-remainder-energy} turns watt residuals into tail coefficient energy: raising \(\ell_{\max}\)
reduces \(\|\mathbf{R}_{\ell_{\max}}\|_{\mathrm{rad}}\) monotonically when coefficients decay \cite{stratton1941}.

\begin{figure}[t]
\centering
\ChOneFigBlock{%
\begin{tikzpicture}[ch1fig, node distance=7mm and 9mm]
  \node[ch1box] (v) {$\mathbf{v}_\omega$};
  \node[ch1box, right=11mm of v] (t) {$\mathcal{T}_{\ell_{\max}}$};
  \node[ch1box, right=11mm of t] (tail) {$\mathbf{R}_{\ell_{\max}}$};
  \node[ch1box, below=10mm of t] (pow) {Eq.~\eqref{eq:ch3.truncation-remainder-energy}};
  \draw[ch1flow] (v) -- (t);
  \draw[ch1flow] (v) -- (tail);
  \draw[ch1flow] (pow) -- (tail);
\end{tikzpicture}%
}
\caption{Subsection~\ref{sec:ch362} truncation split: retained subspace under \(\mathcal{T}_{\ell_{\max}}\) and
orthogonal remainder \(\mathbf{R}_{\ell_{\max}}\).}
\label{fig:ch3.362-truncation-split}
\end{figure}

\paragraph{Compact support and spectral tail (HOME).}
For sources supported in a ball of radius \(a\), far coefficients exhibit accelerated decay for
\(\ell\gg k_{0}a\) \cite{stratton1941,jackson1999,hansen1988}. Subsection~\ref{sec:ch362} records the engineering
rule: choose \(\ell_{\max}\gtrsim k_{0}a\) before attributing residuals to noise \cite{harrington2001}. When
\(k_{0}a\ll1\), low-\(\ell\) truncation often suffices; electrically large sources demand higher cuts
\cite{stratton1941,hansen1988}.

\paragraph{Compression bound import.}
When \(\{\widehat{\mathbf{u}}_{\ell m}^{(\sigma)}\}\) orders decreasing singular values of a compact source operator,
Eq.~\eqref{eq:ch2.compression-error-bound} bounds \(\|\mathbf{R}_{\ell_{\max}}\|_{\mathrm{rad}}\) by a tail singular
value \(\sigma_{\ell_{\max}+1}\) \cite{harrington2001}. Subsection~\ref{sec:ch362} uses that bound as the bridge to
Chapter~\ref{ch:04} realizability screens \cite{harrington2001,stratton1941}.

\paragraph{Worked illustration (quadrupole tail after \(\ell_{\max}=1\)).}
If \(\widehat{c}_{20}^{(\mathbf{N})}\neq0\) but \(\ell_{\max}=1\), remainder norm equals
\(|\widehat{c}_{20}^{(\mathbf{N})}|^{2}\) watts plus higher-\(\ell\) terms from
Eq.~\eqref{eq:ch3.truncation-remainder-energy} \cite{poynting1884,harrington2001}. A dipole-only fit cannot absorb
quadrupole power; the tail is representation error, not measurement noise \cite{stratton1941,hansen1988}.

\paragraph{Problem (watt residual from truncation).}
\textbf{Problem.} Coefficients \(|\widehat{c}_{10}^{(\mathbf{N})}|^{2}=0.8\,\mathrm{W}\),
\(|\widehat{c}_{20}^{(\mathbf{N})}|^{2}=0.15\,\mathrm{W}\), \(|\widehat{c}_{30}^{(\mathbf{N})}|^{2}=0.05\,\mathrm{W}\).
Find remainder power for \(\ell_{\max}=1\) \cite{harrington2001,poynting1884,stratton1941}.

\textbf{Step 1.} Apply Eq.~\eqref{eq:ch3.truncation-remainder-energy} \cite{hansen1988}.

\textbf{Step 2.} Sum \(\ell>1\) terms: \(0.15+0.05=0.2\,\mathrm{W}\) \cite{poynting1884}.

\textbf{Step 3.} Interpret \(0.2\,\mathrm{W}\) as unmodeled tail under dipole truncation \cite{harrington2001}.

\textbf{Physical meaning.} Truncation lowers reported total power unless tail is declared \cite{stratton1941}.

\begin{table}[t]
\centering
\caption{Subsection~\ref{sec:ch362} truncation loss taxonomy.}
\label{tab:ch3.362-truncation-taxonomy}
\small
\begin{tabular}{@{}p{0.28\linewidth}p{0.36\linewidth}p{0.30\linewidth}@{}}
\toprule
Mechanism & Object & Symptom \\
\midrule
Finite \(\ell_{\max}\) & Eq.~\eqref{eq:ch3.ellmax-truncation-projector} & tail watts \\
Compact source & \(k_{0}a\) cutoff & slow tail decay \\
Finite aperture & Eq.~\eqref{eq:ch2.reconstruction-rank-cap} & rank-limited fit \\
\bottomrule
\end{tabular}
\end{table}

\paragraph{Observation rank versus truncation.}
Even when \(\ell_{\max}\) is large, a finite receiver basis implements only \(r\) independent measurements
\cite{harrington2001,stratton1941}. Eq.~\eqref{eq:ch2.reconstruction-rank-cap} caps recoverable coefficients below
\(N_{\mathrm{disc}}(\ell_{\max})\) from Eq.~\eqref{eq:ch3.cumulative-disc-DOS} \cite{harrington2001}. Truncation and
rank limits multiply: first cut the basis, then cut what the laboratory resolves \cite{stratton1941}.

\paragraph{Worked illustration (SV tail controls \(\ell_{\max}\)).}
If singular values \(\sigma_{n}\) decay as \(n^{-3/2}\), Eq.~\eqref{eq:ch2.compression-error-bound} suggests
\(\ell_{\max}\approx50\) for \(1\%\) tail error when the 51st mode carries \(\sigma_{51}\approx10^{-2}\sigma_{1}\)
\cite{harrington2001}. Subsection~\ref{sec:ch362} recommends reporting \(\ell_{\max}\) alongside \(\sigma\)-tail
estimates in inversion reports \cite{harrington2001,devaney2004}.

\paragraph{Problem (choose \(\ell_{\max}\) from \(k_{0}a\)).}
\textbf{Problem.} A source has \(k_{0}a=5\). Propose \(\ell_{\max}\) using \(\ell_{\max}\gtrsim k_{0}a\)
\cite{stratton1941,jackson1999,hansen1988}.

\textbf{Step 1.} Set baseline \(\ell_{\max}=5\) \cite{jackson1999}.

\textbf{Step 2.} Add margin for \(\sigma\)-type sectors: \(\ell_{\max}=7\) \cite{hansen1988}.

\textbf{Step 3.} Verify tail watts with Eq.~\eqref{eq:ch3.truncation-remainder-energy} \cite{harrington2001}.

\textbf{Physical meaning.} Electrical size sets first-pass multipole cut \cite{stratton1941}.

\paragraph{Compact observation support.}
When measurements are taken only on a subsector of \(\Gamma_R\), extrapolated coefficients beyond the aperture are
ill-conditioned even if \(\ell_{\max}\) is large \cite{harrington2001,devaney2004}. Subsection~\ref{sec:ch362}
separates \emph{basis truncation} from \emph{aperture truncation}; both shrink effective completeness
\cite{stratton1941,harrington2001}.

\paragraph{Near-field contamination under finite cuts.}
Reactive fields near compact sources can project onto high-\(\ell\) modes that decay slowly in the far zone if
\(\Pi_{\mathrm{rad}}\) is not applied before truncation \cite{harrington2001,stratton1941}. Outgoing projection
Eq.~\eqref{eq:ch3.outgoing-inheritance-criterion} is prerequisite to interpreting
Eq.~\eqref{eq:ch3.truncation-remainder-energy} as far-zone watts \cite{poynting1884,hansen1988}.

\begin{figure}[t]
\centering
\ChOneFigBlock{%
\begin{tikzpicture}[ch1fig, node distance=7mm and 9mm]
  \node[ch1box] (ell) {$\ell_{\max}$};
  \node[ch1box, right=10mm of ell] (tail) {Eq.~\eqref{eq:ch3.truncation-remainder-energy}};
  \node[ch1box, right=10mm of tail] (sv) {Eq.~\eqref{eq:ch2.compression-error-bound}};
  \node[ch1box, below=10mm of tail] (rank) {Eq.~\eqref{eq:ch2.reconstruction-rank-cap}};
  \draw[ch1flow] (ell) -- (tail) -- (sv);
  \draw[ch1flow] (rank) -- (tail);
\end{tikzpicture}%
}
\caption{Subsection~\ref{sec:ch362} error chain: multipole cut, tail energy, compression bound, and observation rank.}
\label{fig:ch3.362-error-chain}
\end{figure}

\paragraph{Monotone watt lower bound.}
Truncated power sums \(\sum_{\ell\le\ell_{\max}}|\widehat{c}|^{2}\) are lower bounds on full
\(\mathcal{P}_{\mathrm{rad}}\) only when remainder \(\mathbf{R}_{\ell_{\max}}\) is orthogonal to the retained span
\cite{harrington2001,poynting1884}. Eq.~\eqref{eq:ch3.truncation-remainder-energy} certifies that orthogonality in the
normalized VSH gauge \cite{hansen1988,stratton1941}.

\paragraph{Worked illustration (array element count versus \(\ell_{\max}\)).}
An \(N\)-element array with fixed spacing may synthesize only \(\ell_{\max}\sim\sqrt{N}\) independent far patterns
before grating lobes alias high-\(\ell\) content \cite{harrington2001,stratton1941}. Subsection~\ref{sec:ch362} warns
against setting \(\ell_{\max}\) from mathematics alone without element budget \cite{hansen1988,harrington2001}.

\paragraph{Problem (rank cap dominates truncation).}
\textbf{Problem.} \(\ell_{\max}=10\) but observation rank \(r=6\). How many coefficients are stably recoverable?
\cite{harrington2001,stratton1941,devaney2004}.

\textbf{Step 1.} Recoverable count \(\le r\) \cite{harrington2001}.

\textbf{Step 2.} Remaining labels lie in measurement null space \cite{stratton1941}.

\textbf{Physical meaning.} Instrument rank can truncate before multipole cut \cite{harrington2001}.

\paragraph{Link to Chapter~\ref{ch:04}.}
Truncation loss Eq.~\eqref{eq:ch3.truncation-remainder-energy} is the angular-momentum form of finite-rank
realizability: not every coefficient vector with small tail is source-realizable on a given \(\Omega_{s}\)
\cite{harrington2001,stratton1941}. Subsection~\ref{sec:ch362} supplies the harmonic tail language for
Section~\ref{sec:ch46} \cite{harrington2001}.

\paragraph{Tail decay bound for compact sources (HOME).}
When \(|\widehat{c}_{\ell m}^{(\sigma)}|\le C(\ell_{0}/\ell)^{p}\) for \(\ell>\ell_{0}\), remainder energy satisfies
\begin{equation}
\label{eq:ch3.truncation-tail-decay-bound}
\|\mathbf{R}_{\ell_{\max}}\|_{\mathrm{rad}}^{2}
\le
C'\sum_{\ell>\ell_{\max}}(2\ell+1)\Big(\frac{\ell_{0}}{\ell}\Big)^{2p},
\end{equation}
importing multiplet count Eq.~\eqref{eq:ch3.multiplet-multiplicity-def} \cite{hansen1988,harrington2001,stratton1941}.
Equation~\eqref{eq:ch3.truncation-tail-decay-bound} turns algebraic \(\ell\)-decay into a watt tail budget for campaign
planning \cite{poynting1884,harrington2001}.

\paragraph{Worked illustration (super-algebraic decay).}
Smooth compact sources often exhibit faster-than-polynomial coefficient decay, making \(\ell_{\max}\approx k_{0}a+5\)
sufficient for sub-percent tails \cite{stratton1941,jackson1999}. Subsection~\ref{sec:ch362} recommends measuring tail
watts directly on \(\Gamma_R\) rather than trusting decay laws alone \cite{harrington2001,poynting1884}.

\paragraph{Problem (estimate tail from decay exponent).}
\textbf{Problem.} \(p=2\), \(\ell_{0}=3\), \(\ell_{\max}=8\). Bound relative tail power using
Eq.~\eqref{eq:ch3.truncation-tail-decay-bound} qualitatively \cite{harrington2001,hansen1988,stratton1941}.

\textbf{Step 1.} Tail begins at \(\ell=9\) \cite{hansen1988}.

\textbf{Step 2.} Each skipped \(\ell\) contributes \(\sim(3/\ell)^{4}\) weighted by \(2\ell+1\) \cite{jackson1999}.

\textbf{Step 3.} Conclude rapid decay when \(\ell\gg\ell_{0}\) \cite{harrington2001}.

\textbf{Physical meaning.} Smooth sources shrink truncation error faster than minimum Nyquist cuts suggest \cite{stratton1941}.

\paragraph{Aliasing under spatial undersampling.}
Finite element spacing \(\Delta\) on an array can alias high-\(\ell\) content into low-\(\ell\) apparent coefficients
\cite{harrington2001,stratton1941}. Subsection~\ref{sec:ch362} distinguishes aliasing (sampling artifact) from
Eq.~\eqref{eq:ch3.truncation-remainder-energy} (basis cut) \cite{hansen1988,harrington2001}.

\begin{figure}[t]
\centering
\ChOneFigBlock{%
\begin{tikzpicture}[ch1fig, node distance=7mm and 9mm]
  \node[ch1box] (true) {true $\widehat{c}_{\ell}$};
  \node[ch1box, right=10mm of true] (cut) {$\mathcal{T}_{\ell_{\max}}$};
  \node[ch1box, right=10mm of cut] (alias) {array alias};
  \node[ch1box, below=10mm of cut] (tail) {Eq.~\eqref{eq:ch3.truncation-tail-decay-bound}};
  \draw[ch1flow] (true) -- (cut) -- (tail);
  \draw[ch1flow] (alias) -- (cut);
\end{tikzpicture}%
}
\caption{Subsection~\ref{sec:ch362} separates multipole truncation tails from spatial aliasing on finite arrays.}
\label{fig:ch3.362-alias-truncation}
\end{figure}

\paragraph{Accessible budget intersection.}
Eq.~\eqref{eq:ch2.accessible-angular-budget} may cap observable \(\ell\) below the mathematically retained
\(\ell_{\max}\) \cite{harrington2001,stratton1941}. Subsection~\ref{sec:ch362} records that intersection as an
effective truncation tighter than Eq.~\eqref{eq:ch3.ellmax-truncation-projector} alone \cite{hansen1988,harrington2001}.

\begin{table}[t]
\centering
\caption{Subsection~\ref{sec:ch362} truncation error budget worksheet.}
\label{tab:ch3.362-error-worksheet}
\small
\begin{tabular}{@{}p{0.28\linewidth}p{0.32\linewidth}p{0.32\linewidth}@{}}
\toprule
Term & Formula & Typical source \\
\midrule
Basis tail & Eq.~\eqref{eq:ch3.truncation-remainder-energy} & low \(\ell_{\max}\) \\
Decay bound & Eq.~\eqref{eq:ch3.truncation-tail-decay-bound} & smooth \(\Jimp\) \\
Rank cap & Eq.~\eqref{eq:ch2.reconstruction-rank-cap} & small aperture \\
\bottomrule
\end{tabular}
\end{table}

\paragraph{Worked illustration (two-term tail budget).}
If \(\|\mathbf{R}_{8}\|_{\mathrm{rad}}^{2}=0.05\,\mathrm{W}\) and rank cap loses another \(0.02\,\mathrm{W}\) to
\(\mathcal{L}_{\mathrm{meas}}\), total unreported export is \(0.07\,\mathrm{W}\) against a \(1\,\mathrm{W}\) reference
\cite{poynting1884,harrington2001}. Campaign reports should list both terms separately \cite{stratton1941,hansen1988}.

\paragraph{Problem (combine tail and rank).}
\textbf{Problem.} Tail \(0.04\,\mathrm{W}\), rank-lost \(0.01\,\mathrm{W}\), measured total \(0.90\,\mathrm{W}\). Estimate
full export \cite{harrington2001,poynting1884,stratton1941}.

\textbf{Step 1.} Sum unmodeled terms \(0.05\,\mathrm{W}\) \cite{poynting1884}.

\textbf{Step 2.} Add measured \(0.90\,\mathrm{W}\) \cite{harrington2001}.

\textbf{Step 3.} Obtain \(0.95\,\mathrm{W}\) lower bound on true export \cite{stratton1941}.

\textbf{Physical meaning.} Truncation and rank deficits accumulate in watt budgets \cite{harrington2001}.

\paragraph{Remainder orthogonality audit.}
Eq.~\eqref{eq:ch3.truncation-remainder-energy} holds because \(\mathbf{R}_{\ell_{\max}}\) lies in the orthogonal complement
of \(\mathcal{T}_{\ell_{\max}}\mathcal{F}_{\mathrm{rad}}\) under Theorem~\ref{thm:ch3.vsh-orthogonality}
\cite{hansen1988,kato1995}. Violating normalized gauge Eq.~\eqref{eq:ch3.VSH-unit-energy-norm} can introduce spurious
cross-terms that masquerade as truncation tails \cite{harrington2001,stratton1941}.

\paragraph{Worked illustration (SVD reordering versus \(\ell\) cut).}
Singular-value truncation in Chapter~\ref{ch:02} may retain non-contiguous \(\ell\) sectors when source operators are
non-spherical \cite{harrington2001,stratton1941}. Eq.~\eqref{eq:ch3.ellmax-truncation-projector} is the spherical harmonic
\emph{ordered} cut; SVD cuts are operator-optimal but label-mixed \cite{harrington2001,devaney2004}. Campaigns should declare
which truncation philosophy is used \cite{hansen1988,harrington2001}.

\paragraph{Problem (set \(\ell_{\max}\) from tail watt target).}
\textbf{Problem.} Require \(\|\mathbf{R}_{\ell_{\max}}\|_{\mathrm{rad}}^{2}\le0.01\,\mathrm{W}\) on a \(1\,\mathrm{W}\) pattern.
Choose \(\ell_{\max}\) given \(|\widehat{c}_{\ell m}^{(\sigma)}|\) data \cite{harrington2001,poynting1884,stratton1941}.

\textbf{Step 1.} Sum tail coefficients for candidate cuts \cite{hansen1988}.

\textbf{Step 2.} Increase \(\ell_{\max}\) until tail falls below \(0.01\,\mathrm{W}\) \cite{poynting1884}.

\textbf{Step 3.} Record chosen \(\ell_{\max}\) in metadata \cite{harrington2001}.

\textbf{Physical meaning.} Watt targets translate abstract completeness loss into acceptance thresholds \cite{stratton1941}.

\begin{figure}[t]
\centering
\ChOneFigBlock{%
\begin{tikzpicture}[ch1fig, node distance=7mm and 9mm]
  \node[ch1box] (target) {1\% tail};
  \node[ch1box, right=10mm of target] (scan) {scan $\ell_{\max}$};
  \node[ch1box, right=10mm of scan] (ok) {Eq.~\eqref{eq:ch3.truncation-remainder-energy}};
  \draw[ch1flow] (target) -- (scan) -- (ok);
\end{tikzpicture}%
}
\caption{Subsection~\ref{sec:ch362} acceptance scan: raise \(\ell_{\max}\) until tail watts meet tolerance.}
\label{fig:ch3.362-tail-scan}
\end{figure}

\paragraph{Truncation versus eigenfield remainder preview.}
Eq.~\eqref{eq:ch3.eigenfield-truncation-remainder} is the abstract index form; Eqs.~\eqref{eq:ch3.ellmax-truncation-projector}--\eqref{eq:ch3.truncation-remainder-energy}
specialize it to VSH labels with explicit watt interpretation \cite{harrington2001,stratton1941}. Subsection~\ref{sec:ch362}
therefore closes the preview opened in Subsection~\ref{sec:ch301} \cite{hansen1988,harrington2001}.

\paragraph{Worked illustration (store \(\ell_{\max}\) in inversion archives).}
Inverse solvers should archive \(\ell_{\max}\), tail watts, and rank \(r\) with every coefficient vector so
Eq.~\eqref{eq:ch3.truncation-remainder-energy} can be recomputed in peer review \cite{devaney2004,harrington2001,stratton1941}.

\begin{table}[t]
\centering
\caption{Subsection~\ref{sec:ch362} recommended truncation metadata.}
\label{tab:ch3.362-truncation-metadata}
\small
\begin{tabular}{@{}p{0.30\linewidth}p{0.58\linewidth}@{}}
\toprule
Field & Equation \\
\midrule
Cut level & \(\ell_{\max}\) in Eq.~\eqref{eq:ch3.ellmax-truncation-projector} \\
Tail watts & Eq.~\eqref{eq:ch3.truncation-remainder-energy} \\
Decay bound & Eq.~\eqref{eq:ch3.truncation-tail-decay-bound} \\
Rank cap & Eq.~\eqref{eq:ch2.reconstruction-rank-cap} \\
\bottomrule
\end{tabular}
\end{table}

\paragraph{Problem (full truncation audit worksheet).}
\textbf{Problem.} For a measured pattern, execute the Subsection~\ref{sec:ch362} audit: (i) choose \(\ell_{\max}\),
(ii) compute tail watts, (iii) apply rank cap, (iv) report residual \cite{harrington2001,poynting1884,stratton1941,devaney2004}.

\textbf{Step 1.} Estimate \(k_{0}a\) and set initial \(\ell_{\max}\) \cite{jackson1999,stratton1941}.

\textbf{Step 2.} Evaluate Eq.~\eqref{eq:ch3.truncation-remainder-energy} on \(\Gamma_R\) \cite{poynting1884}.

\textbf{Step 3.} Intersect with Eq.~\eqref{eq:ch2.reconstruction-rank-cap} \cite{harrington2001}.

\textbf{Step 4.} Archive Table~\ref{tab:ch3.362-truncation-metadata} fields \cite{devaney2004}.

\textbf{Physical meaning.} Truncation audits are mandatory for any claim of ``complete'' modal fitting \cite{harrington2001}.

\begin{figure}[t]
\centering
\ChOneFigBlock{%
\begin{tikzpicture}[ch1fig, node distance=7mm and 9mm]
  \node[ch1box] (data) {measured $\widehat{c}$};
  \node[ch1box, right=9mm of data] (meta) {Table~\ref{tab:ch3.362-truncation-metadata}};
  \node[ch1box, right=9mm of meta] (ok) {peer review};
  \draw[ch1flow] (data) -- (meta) -- (ok);
\end{tikzpicture}%
}
\caption{Subsection~\ref{sec:ch362} archival chain: coefficients without truncation metadata are not peer-reviewable.}
\label{fig:ch3.362-archive-chain}
\end{figure}

Subsection~\ref{sec:ch362} has quantified truncation-induced incompleteness through
Eqs.~\eqref{eq:ch3.ellmax-truncation-projector}--\eqref{eq:ch3.truncation-remainder-energy}
\cite{hansen1988,harrington2001,stratton1941}. Subsection~\ref{sec:ch363} treats a distinct failure mode: boundaries
that \emph{mix} previously orthogonal labels \cite{stratton1941}.


\subsection{Boundary-induced mode mixing}
\label{sec:ch363}
% TOC tag: [HOME]

\emph{Boundary-induced mode mixing} occurs when conductors, ground planes, or inhomogeneous interfaces break the
symmetry that diagonalized VSH labels in Subsection~\ref{sec:ch343}, coupling \((\ell,m,\sigma)\) sectors that were
orthogonal on unbounded domains \cite{stratton1941,harrington2001,jackson1999}. Subsection~\ref{sec:ch363} is the
HOME coupling calculus for that mixing---distinct from truncation tails of Subsection~\ref{sec:ch362}
\cite{hansen1988,harrington2001}.

\paragraph{Standing imports.}
Rotation commutation Eq.~\eqref{eq:ch2.radiation-rotation-commutation}, geometric symmetry operator
Eq.~\eqref{eq:ch2.geometric-symmetry-operator-def}, and boundary image constraints from
Theorem~\ref{thm:ch2.geometric-spectral-constraints} enter by reference \cite{stratton1941,harrington2001}.
Subsection~\ref{sec:ch363} does not re-derive image theory or Fresnel reflection \cite{jackson1999,stratton1941}.

\paragraph{Assumptions.}
Consider exterior problems with additional boundaries \(\Gamma_{b}\) (ground plane, cavity aperture, radome shell)
that reduce the symmetry group from \(\mathrm{SO}(3)\) to a subgroup \(G_{\mathrm{red}}\subset\mathrm{SO}(3)\)
\cite{harrington2001,stratton1941}. Fields remain outgoing on the exterior radiation sphere; mixing is diagnosed on
\(\mathcal{F}_{\mathrm{rad}}\) after \(\Pi_{\mathrm{rad}}\) \cite{harrington2001,kato1995}.

\paragraph{Boundary scattering matrix on labels (HOME).}
Let \(\mathbf{c}\) be the normalized coefficient vector on the unbounded VSH basis. Boundary presence induces
\begin{equation}
\label{eq:ch3.boundary-scattering-mixing-matrix}
\mathbf{c}^{\mathrm{obs}}
=
\mathbf{S}_{b}\,\mathbf{c},
\qquad
[\mathbf{S}_{b}]_{(\ell m\sigma),(\ell'm'\sigma')}
=
\langle\widehat{\mathbf{u}}_{\ell m}^{(\sigma)},\,\mathcal{B}\widehat{\mathbf{u}}_{\ell'm'}^{(\sigma')}\rangle_{\mathrm{rad}},
\end{equation}
where \(\mathcal{B}\) is the linear boundary map (reflection plus transmission) on radiative states
\cite{stratton1941,harrington2001}. Equation~\eqref{eq:ch3.boundary-scattering-mixing-matrix} is the Subsection~\ref{sec:ch363}
definition of mode mixing: off-diagonal entries of \(\mathbf{S}_{b}\) couple labels that were orthogonal before
\(\Gamma_{b}\) \cite{hansen1988,harrington2001}.

\paragraph{Reduced-symmetry block form (HOME).}
When \(G_{\mathrm{red}}\) preserves one axis (e.g., ground plane at \(z=0\)), the unblocked VSH basis block-diagonalizes
into sectors labeled by conserved quantum numbers (often \(m\) or parity) only partially:
\begin{equation}
\label{eq:ch3.reduced-symmetry-block-form}
\mathbf{S}_{b}
=
\mathrm{diag}(\mathbf{S}^{(m_{1})},\mathbf{S}^{(m_{2})},\ldots)
\;+\;
\mathbf{S}_{\mathrm{cross}},
\end{equation}
with \(\mathbf{S}_{\mathrm{cross}}\neq0\) when symmetry is only approximate \cite{jackson1999,harrington2001,stratton1941}.
Equation~\eqref{eq:ch3.reduced-symmetry-block-form} explains why half-space problems retain some \(m\) structure yet
still mix \(\ell\) within a sector \cite{stratton1941,hansen1988}.

\begin{figure}[t]
\centering
\ChOneFigBlock{%
\begin{tikzpicture}[ch1fig, node distance=7mm and 9mm]
  \node[ch1box] (free) {free-space VSH};
  \node[ch1box, right=11mm of free] (bnd) {$\Gamma_{b}$};
  \node[ch1box, right=11mm of bnd] (mix) {Eq.~\eqref{eq:ch3.boundary-scattering-mixing-matrix}};
  \node[ch1box, below=10mm of bnd] (blk) {Eq.~\eqref{eq:ch3.reduced-symmetry-block-form}};
  \draw[ch1flow] (free) -- (bnd) -- (mix);
  \draw[ch1flow] (blk) -- (mix);
\end{tikzpicture}%
}
\caption{Subsection~\ref{sec:ch363} boundary chain: symmetry reduction on \(\Gamma_{b}\) induces label coupling
\(\mathbf{S}_{b}\).}
\label{fig:ch3.363-boundary-mixing-chain}
\end{figure}

\paragraph{Image-mode admixture (HOME).}
For a perfect electric conductor ground plane, image sources admix sectors related by reflection:
\begin{equation}
\label{eq:ch3.image-mode-admixture}
\widehat{\mathbf{u}}_{\ell m}^{(\sigma)\,\mathrm{PEC}}
\propto
\widehat{\mathbf{u}}_{\ell m}^{(\sigma)}+(-1)^{\ell}\,\mathcal{M}_{z}\widehat{\mathbf{u}}_{\ell m}^{(\sigma)},
\end{equation}
with \(\mathcal{M}_{z}\) the mirror operator on fields \cite{jackson1999,stratton1941}. Equation
\eqref{eq:ch3.image-mode-admixture} is schematic: the physical content is that parity under \(z\to -z\) replaces
pure \(\mathrm{SO}(3)\) multiplets with standing-plus-traveling hybrids in the half-space \cite{harrington2001,stratton1941}.

\paragraph{Physical meaning of mixing versus truncation.}
Truncation omits high-\(\ell\) modes but leaves retained labels orthogonal
(Subsection~\ref{sec:ch362}) \cite{hansen1988}. Boundary mixing rotates coefficient energy across labels at fixed
\(\ell_{\max}\): a dipole template can acquire quadrupole coupling through \(\mathbf{S}_{b}\) even when both modes are
in the truncated set \cite{harrington2001,stratton1941}. Confusing the two mechanisms misattributes fit residuals
\cite{harrington2001}.

\paragraph{Worked illustration (ground plane breaks \(m\) degeneracy).}
In free space, \(\pm m\) partners are degenerate under rotation about \(z\)
Eq.~\eqref{eq:ch3.eigenfield-Wigner-mixing} \cite{jackson1999,hansen1988}. A ground plane breaks full rotation to
\(\mathrm{C}_{\infty\nu}\); degeneracy splits and \(\mathbf{S}_{b}\) develops nonzero off-diagonal coupling between
sectors that shared energy on \(\mathbb{R}^{3}\) \cite{stratton1941,harrington2001}.

\paragraph{Problem (detect mixing from Gram matrix).}
\textbf{Problem.} After installing a radome, off-diagonal Gram entries
\(G_{(\ell m\sigma),(\ell'm'\sigma')}\neq0\) appear. Diagnose \cite{harrington2001,hansen1988,stratton1941}.

\textbf{Step 1.} Form \(\mathbf{S}_{b}\) via Eq.~\eqref{eq:ch3.boundary-scattering-mixing-matrix} \cite{hansen1988}.

\textbf{Step 2.} Compare to truncation tail Eq.~\eqref{eq:ch3.truncation-remainder-energy} \cite{harrington2001}.

\textbf{Step 3.} If tail is small but \(G\) off-diagonal persists, attribute to boundary mixing \cite{stratton1941}.

\textbf{Physical meaning.} Non-diagonal Gram signals symmetry breaking, not necessarily insufficient \(\ell_{\max}\)
\cite{harrington2001}.

\begin{table}[t]
\centering
\caption{Subsection~\ref{sec:ch363} boundary mixing versus truncation.}
\label{tab:ch3.363-mixing-vs-truncation}
\small
\begin{tabular}{@{}p{0.22\linewidth}p{0.34\linewidth}p{0.36\linewidth}@{}}
\toprule
Effect & Matrix signature & Remedy \\
\midrule
Truncation & tail in Eq.~\eqref{eq:ch3.truncation-remainder-energy} & raise \(\ell_{\max}\) \\
Boundary mixing & \(\mathbf{S}_{b}\) off-diagonal & change basis / model \(\mathcal{B}\) \\
Rank limit & Eq.~\eqref{eq:ch2.reconstruction-rank-cap} & enlarge aperture \\
\bottomrule
\end{tabular}
\end{table}

\paragraph{Cavity and guide quantization (pointer).}
Enclosed volumes quantize spectra discretely; Subsection~\ref{sec:ch372} develops angular spectra under boundaries
\cite{stratton1941,harrington2001}. Subsection~\ref{sec:ch363} supplies the \emph{mixing} vocabulary: boundary
conditions are linear constraints that replace continuous \(\ell\) with mixed normal modes \cite{hansen1988,jackson1999}.

\paragraph{Worked illustration (two-port radome coupling).}
A transmit mode \(\widehat{\mathbf{u}}_{10}^{(\mathbf{N})}\) scattering from \(\Gamma_{b}\) excites receive projection on
\(\widehat{\mathbf{u}}_{20}^{(\mathbf{N})}\) when \([\mathbf{S}_{b}]_{(20),(10)}\neq0\) \cite{harrington2001,stratton1941}.
Reciprocity Eq.~\eqref{eq:ch3.reciprocity-coupling-symmetry} constrains \(\mathbf{S}_{b}\) on lossless exteriors
\cite{stratton1941,harrington2001}.

\paragraph{Problem (estimate cross-coupling magnitude).}
\textbf{Problem.} \(|[\mathbf{S}_{b}]_{(11),(10)}|=0.15\) with unit-flux modes. What fraction of dipole power appears in
adjacent sector after one boundary bounce? \cite{harrington2001,hansen1988,stratton1941}.

\textbf{Step 1.} Read coupling magnitude squared as watt cross-term if coherent \cite{poynting1884}.

\textbf{Step 2.} For small coupling, \(0.15^{2}=2.25\%\) secondary-sector leakage \cite{harrington2001}.

\textbf{Physical meaning.} Mixing redistributes labeled power without changing total watts in lossless media
\cite{stratton1941}.

\paragraph{Anisotropic media admixture.}
Bianisotropic or tensor \(\varepsilon,\mu\) break reciprocity-symmetric diagonalization of VSH sectors even without
conductors \cite{stratton1941,harrington2001}. \(\mathbf{S}_{b}\) then includes material rotation of
\(\widehat{\mathbf{u}}_{\ell m}^{(\sigma)}\) into non-degenerate hybrid polarizations \cite{hansen1988,harrington2001}.

\begin{figure}[t]
\centering
\ChOneFigBlock{%
\begin{tikzpicture}[ch1fig, node distance=7mm and 9mm]
  \node[ch1box] (diag) {diagonal VSH};
  \node[ch1box, right=11mm of diag] (pec) {PEC / radome};
  \node[ch1box, right=11mm of pec] (gram) {off-diagonal $G$};
  \node[ch1box, below=10mm of pec] (img) {Eq.~\eqref{eq:ch3.image-mode-admixture}};
  \draw[ch1flow] (diag) -- (pec) -- (gram);
  \draw[ch1flow] (img) -- (pec);
\end{tikzpicture}%
}
\caption{Subsection~\ref{sec:ch363} diagnostic: boundaries turn diagonal VSH Gram matrices into coupled label networks.}
\label{fig:ch3.363-gram-offdiagonal}
\end{figure}

\paragraph{Wigner mixing reinterpretation.}
Eq.~\eqref{eq:ch3.eigenfield-Wigner-mixing} describes rotation within degenerate multiplets; boundary mixing replaces
degeneracy with forced coupling between \emph{distinct} multipoles \cite{jackson1999,hansen1988}. Subsection~\ref{sec:ch363}
extends representation language from group rotations to boundary operators \(\mathcal{B}\) \cite{stratton1941,harrington2001}.

\paragraph{Worked illustration (slot through ground plane).}
A slot antenna breaks translational symmetry; \(\mathbf{S}_{b}\) couples many \(m\) sectors simultaneously
\cite{harrington2001,stratton1941}. Pattern nulls on the ground-plane side are geometric accessibility features
Eq.~\eqref{eq:ch2.angular-accessibility-sector}, not truncation artifacts \cite{stratton1941,harrington2001}.

\paragraph{Problem (parity selection on PEC).}
\textbf{Problem.} Which \(\ell\) sectors dominate over a PEC plane under image law Eq.~\eqref{eq:ch3.image-mode-admixture}?
\cite{jackson1999,stratton1941,hansen1988}.

\textbf{Step 1.} Vertical dipole images add constructively for odd parity sectors \cite{jackson1999}.

\textbf{Step 2.} Even-\(\ell\) magnetic sectors may cancel on broadside \cite{stratton1941}.

\textbf{Physical meaning.} Parity under reflection replaces free-space degeneracy rules \cite{harrington2001}.

\paragraph{Measurement consequence.}
Harmonic receivers calibrated in free space mis-estimate \(\widehat{c}_{\ell m}^{(\sigma)}\) when \(\mathbf{S}_{b}\neq\mathbf{I}\)
\cite{harrington2001,hansen1988}. Subsection~\ref{sec:ch363} requires boundary metadata in
Eq.~\eqref{eq:ch3.measurement-normalized-coefficients} pipelines \cite{harrington2001,stratton1941}.

\paragraph{Mutual coupling Gram matrix (HOME).}
Boundary mixing induces a non-diagonal receive Gram matrix
\begin{equation}
\label{eq:ch3.mutual-coupling-gram}
G_{mn}
=
\langle\widehat{\mathbf{u}}_{m},\,\mathbf{S}_{b}^{\dagger}\mathbf{S}_{b}\,\widehat{\mathbf{u}}_{n}\rangle_{\mathrm{rad}},
\end{equation}
whose off-diagonal terms quantify how \(\Gamma_{b}\) couples normalized ports \(m,n\) \cite{harrington2001,stratton1941}.
Equation~\eqref{eq:ch3.mutual-coupling-gram} reduces to identity on unbounded domains where \(\mathbf{S}_{b}=\mathbf{I}\)
\cite{hansen1988,harrington2001}.

\paragraph{Worked illustration (radome bistatic leakage).}
Transmit \(\widehat{\mathbf{u}}_{10}^{(\mathbf{N})}\) scattering from a radome can project onto
\(\widehat{\mathbf{u}}_{11}^{(\mathbf{M})}\) with \(|G_{(11),(10)}|\neq0\) even when free-space coupling was null
\cite{harrington2001,stratton1941}. Bistatic coupling is boundary mixing, not a change of \(\ell\) label alone
\cite{hansen1988,harrington2001}.

\paragraph{Problem (diagonalize \(\mathbf{S}_{b}\) in a symmetry block).}
\textbf{Problem.} In a \(G_{\mathrm{red}}\)-invariant block, find eigenmodes of \(\mathbf{S}_{b}\) that replace raw VSH
labels \cite{jackson1999,harrington2001,stratton1941}.

\textbf{Step 1.} Restrict \(\mathbf{S}_{b}\) to block Eq.~\eqref{eq:ch3.reduced-symmetry-block-form} \cite{hansen1988}.

\textbf{Step 2.} Diagonalize within the block \cite{harrington2001}.

\textbf{Step 3.} Express watts in eigenmode coordinates to remove off-diagonal leakage \cite{stratton1941}.

\textbf{Physical meaning.} Symmetry-adapted bases can re-diagonalize mixing caused by boundaries \cite{harrington2001}.

\begin{figure}[t]
\centering
\ChOneFigBlock{%
\begin{tikzpicture}[ch1fig, node distance=7mm and 9mm]
  \node[ch1box] (tx) {TX mode};
  \node[ch1box, right=10mm of tx] (b) {$\mathbf{S}_{b}$};
  \node[ch1box, right=10mm of b] (gram) {Eq.~\eqref{eq:ch3.mutual-coupling-gram}};
  \node[ch1box, below=10mm of b] (rx) {RX modes};
  \draw[ch1flow] (tx) -- (b) -- (gram);
  \draw[ch1flow] (rx) -- (gram);
\end{tikzpicture}%
}
\caption{Subsection~\ref{sec:ch363} bistatic chain: boundaries map transmit modes into coupled receive Gram entries.}
\label{fig:ch3.363-bistatic-gram}
\end{figure}

\paragraph{Dielectric interface mixing.}
Fresnel reflection admixes \(\widehat{\mathbf{u}}_{\ell m}^{(\mathbf{N})}\) with \(\widehat{\mathbf{u}}_{\ell m}^{(\mathbf{M})}\)
at lossless dielectric boundaries through tangential continuity \cite{jackson1999,stratton1941}. Subsection~\ref{sec:ch363}
treats that admixture as \(\mathbf{S}_{b}\) rather than as a failure of VSH orthogonality on \(\mathbb{R}^{3}\)
\cite{hansen1988,harrington2001}.

\begin{table}[t]
\centering
\caption{Subsection~\ref{sec:ch363} boundary types and mixing signatures.}
\label{tab:ch3.363-boundary-types}
\small
\begin{tabular}{@{}p{0.24\linewidth}p{0.36\linewidth}p{0.32\linewidth}@{}}
\toprule
Boundary & Symmetry left & Mixing object \\
\midrule
PEC plane & \(m\) partial & Eq.~\eqref{eq:ch3.image-mode-admixture} \\
Radome & none & Eq.~\eqref{eq:ch3.mutual-coupling-gram} \\
Cavity & discrete & Section~\ref{sec:ch372} \\
\bottomrule
\end{tabular}
\end{table}

\paragraph{Worked illustration (corrugated ground plane).}
Corrugations restore approximate \(\mathrm{SO}(3)\) symmetry locally while global \(\mathbf{S}_{b}\) remains nontrivial
\cite{harrington2001,stratton1941}. Mixing matrices should be measured per installation, not assumed from free-space
calibration \cite{hansen1988,harrington2001}.

\paragraph{Problem (compare mixing versus truncation Gram).}
\textbf{Problem.} Off-diagonal Gram persists at \(\ell_{\max}=12\). Decide mixing versus insufficient cut
\cite{harrington2001,hansen1988,stratton1941}.

\textbf{Step 1.} Evaluate tail Eq.~\eqref{eq:ch3.truncation-remainder-energy} \cite{harrington2001}.

\textbf{Step 2.} If tail small, inspect \(\mathbf{S}_{b}\) via Eq.~\eqref{eq:ch3.boundary-scattering-mixing-matrix}
\cite{hansen1988}.

\textbf{Physical meaning.} Persistent off-diagonals at high \(\ell_{\max}\) implicate boundaries \cite{stratton1941}.

\paragraph{Mode mixing versus Wigner rotation (distinction).}
Rotating the laboratory relabels degenerate partners through Eq.~\eqref{eq:ch3.eigenfield-Wigner-mixing} without creating
new off-diagonal \(G_{mn}\) in a fixed basis \cite{jackson1999,hansen1988}. Boundary mixing is \emph{not} a rotation: it is
a non-unitary \(\mathbf{S}_{b}\) in the fixed VSH catalog \cite{harrington2001,stratton1941}. Confusing the two misattributes
pattern changes to source rotation \cite{harrington2001}.

\paragraph{Worked illustration (chamber wall scattering).}
Compact range walls produce \(\mathbf{S}_{b}\neq\mathbf{I}\) even when the intended test zone is free-space-like
\cite{harrington2001,stratton1941}. Wall treatments reduce mixing but rarely restore full
Eq.~\eqref{eq:ch3.unbounded-completeness-statement} orthogonality on finite \(\Gamma_R\) \cite{hansen1988,harrington2001}.

\paragraph{Problem (estimate mixing from single reflection).}
\textbf{Problem.} A PEC wall at distance \(d\) produces one dominant image contribution to \(\mathbf{S}_{b}\). Qualitatively
rank the coupled sectors \cite{jackson1999,stratton1941,harrington2001}.

\textbf{Step 1.} Identify image parity rule Eq.~\eqref{eq:ch3.image-mode-admixture} \cite{jackson1999}.

\textbf{Step 2.} Couple sectors sharing reflected tangential structure \cite{stratton1941}.

\textbf{Step 3.} Expect largest \(G_{mn}\) among parity-aligned partners \cite{harrington2001}.

\textbf{Physical meaning.} Single-bounce images are the leading mixing mechanism in half-space problems \cite{stratton1941}.

\paragraph{Lossy boundary caveat.}
Conductive losses break strict reciprocity symmetry Eq.~\eqref{eq:ch3.reciprocity-coupling-symmetry} for \(\mathbf{S}_{b}\);
passive lossless exteriors are the Subsection~\ref{sec:ch363} default \cite{stratton1941,harrington2001}. Lossy radomes require
biorthogonal pairings deferred beyond Volume~I \cite{harrington2001,hansen1988}.

\paragraph{Coupling to accessibility.}
Mixed sectors may lie in \(\mathcal{L}_{\mathrm{geo}}\) if \(\mathbf{S}_{b}\) maps energy into
\(\mathcal{K}_{\mathrm{sup}}\) Eq.~\eqref{eq:ch2.suppressed-angular-sheet-def} \cite{stratton1941,harrington2001}.
Subsection~\ref{sec:ch363} supplies \(\mathbf{S}_{b}\); Subsection~\ref{sec:ch364} decides which mixed amplitudes are
reportable \cite{harrington2001}.

\paragraph{Near-boundary reactive overlap (caution).}
Reactive fields near \(\Gamma_{b}\) can produce large local inner products that are not far-zone mixing of
Eq.~\eqref{eq:ch3.boundary-scattering-mixing-matrix} \cite{harrington2001,stratton1941}. Apply \(\Pi_{\mathrm{rad}}\) before
forming \(\mathbf{S}_{b}\) on \(\Gamma_R\) in the far zone of the radiator \cite{hansen1988,harrington2001,kato1995}.

\paragraph{Worked illustration (two-boundary round-trip).}
Transmit mode reflected twice from parallel plates produces higher-order \(\mathbf{S}_{b}^{2}\) couplings across \(\ell\)
sectors \cite{harrington2001,stratton1941}. Multi-bounce chambers require matrix products of per-bounce maps, not a single
free-space calibration \cite{hansen1988,harrington2001}.

\paragraph{Problem (symmetrize \(\mathbf{S}_{b}\) for reciprocity audit).}
\textbf{Problem.} Verify lossless reciprocity by checking \(\mathbf{S}_{b}=\mathbf{S}_{b}^{\mathsf{T}}\) in the normalized
VSH gauge \cite{stratton1941,harrington2001,hansen1988}.

\textbf{Step 1.} Measure forward and reverse couplings on \(\Gamma_R\) \cite{poynting1884}.

\textbf{Step 2.} Compare matrix entries within calibration tolerance \cite{harrington2001}.

\textbf{Physical meaning.} Reciprocity failure flags loss, measurement error, or wrong \(\mathcal{N}\) \cite{stratton1941}.

Subsection~\ref{sec:ch363} has defined boundary-induced mixing through
Eqs.~\eqref{eq:ch3.boundary-scattering-mixing-matrix}--\eqref{eq:ch3.image-mode-admixture}
\cite{stratton1941,harrington2001,jackson1999}. Subsection~\ref{sec:ch364} classifies which mixed sectors remain
\emph{accessible} under broken symmetry \cite{harrington2001}.


\subsection{Accessible structure under broken symmetry}
\label{sec:ch364}
% TOC tag: [HOME]

\emph{Accessible structure} is the subset of modal labels that survive geometric, operator, and measurement screens
with nonzero export and resolvable amplitude \cite{harrington2001,stratton1941}.
\emph{Broken symmetry} means the laboratory configuration realizes only a subgroup \(G_{\mathrm{red}}\) of the
\(\mathrm{SO}(3)\) structure used in Subsection~\ref{sec:ch343}, activating mixing of Subsection~\ref{sec:ch363}
while shrinking the reportable index set \cite{jackson1999,harrington2001}. Subsection~\ref{sec:ch364} is the HOME
taxonomy that composes Chapter~\ref{ch:02} accessibility with VSH labels \cite{stratton1941,harrington2001}.

\paragraph{Standing imports.}
Accessible modal sector Eq.~\eqref{eq:ch2.accessible-modal-sector}, angular accessibility sector
Eq.~\eqref{eq:ch2.angular-accessibility-sector}, constraint composition
Eq.~\eqref{eq:ch2.constraint-composition-law}, reconstructability budget
Eq.~\eqref{eq:ch2.reconstructability-budget}, rank cap Eq.~\eqref{eq:ch2.reconstruction-rank-cap}, and
Subsection~\ref{sec:ch274} null-space laws Eq.~\eqref{eq:ch2.angular-spectrum-nullspace-law} are cited without
re-proof \cite{harrington2001,stratton1941}. Subsection~\ref{sec:ch364} does not re-derive Fredholm screens
\cite{harrington2001}.

\paragraph{Assumptions.}
Fix normalized VSH coefficients \(\widehat{c}_{\ell m}^{(\sigma)}\), boundary map \(\mathbf{S}_{b}\) from
Eq.~\eqref{eq:ch3.boundary-scattering-mixing-matrix}, and observation operator \(\mathcal{O}_{\mathrm{meas}}\) with
finite rank \cite{harrington2001,stratton1941}. Accessibility is evaluated on outgoing \(\mathcal{F}_{\mathrm{rad}}\)
after \(\Pi_{\mathrm{rad}}\) \cite{harrington2001,kato1995}.

\paragraph{Accessible VSH label set (HOME).}
Define
\begin{equation}
\label{eq:ch3.accessible-VSH-label-set}
\mathcal{L}_{\mathrm{acc}}
=
\{(\ell,m,\sigma):\ 
\widehat{\mathbf{u}}_{\ell m}^{(\sigma)}\ \text{survives Eqs.~\eqref{eq:ch2.accessible-modal-sector}--\eqref{eq:ch2.angular-accessibility-sector}},\
[\mathbf{S}_{b}\widehat{\mathbf{c}}]_{\ell m}^{(\sigma)}\ \text{not in }\mathcal{K}_{\mathrm{sup}}
\},
\end{equation}
importing suppressed sheet Eq.~\eqref{eq:ch2.suppressed-angular-sheet-def} \cite{stratton1941,harrington2001}.
Equation~\eqref{eq:ch3.accessible-VSH-label-set} is the Subsection~\ref{sec:ch364} discrete label form of
reportable structure: a mode is accessible only if geometry permits export \emph{and} measurement can resolve it
\cite{harrington2001,devaney2004}.

\paragraph{Broken-symmetry sector partition (HOME).}
Decompose coefficient space into
\begin{equation}
\label{eq:ch3.broken-symmetry-sector-partition}
\mathcal{L}
=
\mathcal{L}_{\mathrm{acc}}
\;\uplus\;
\mathcal{L}_{\mathrm{geo}}
\;\uplus\;
\mathcal{L}_{\mathrm{null}}
\;\uplus\;
\mathcal{L}_{\mathrm{meas}},
\end{equation}
where \(\mathcal{L}_{\mathrm{geo}}\) is geometrically forbidden by
Eq.~\eqref{eq:ch2.geometric-modal-annihilation}, \(\mathcal{L}_{\mathrm{null}}\) lies in
\(\ker\mathcal{R}_\omega\), and \(\mathcal{L}_{\mathrm{meas}}\) is exportable but unresolved under
Eq.~\eqref{eq:ch2.reconstruction-rank-cap} \cite{harrington2001,stratton1941}. Equation
\eqref{eq:ch3.broken-symmetry-sector-partition} is the Subsection~\ref{sec:ch364} master partition for broken
symmetry: each label has exactly one exclusion reason when classification is exhaustive \cite{harrington2001}.

\begin{figure}[t]
\centering
\ChOneFigBlock{%
\begin{tikzpicture}[ch1fig, node distance=7mm and 9mm]
  \node[ch1box] (all) {all labels $\mathcal{L}$};
  \node[ch1box, right=10mm of all] (part) {Eq.~\eqref{eq:ch3.broken-symmetry-sector-partition}};
  \node[ch1box, right=10mm of part] (acc) {$\mathcal{L}_{\mathrm{acc}}$};
  \node[ch1box, below=10mm of part] (scr) {Eq.~\eqref{eq:ch2.constraint-composition-law}};
  \draw[ch1flow] (all) -- (part) -- (acc);
  \draw[ch1flow] (scr) -- (part);
\end{tikzpicture}%
}
\caption{Subsection~\ref{sec:ch364} partition: broken symmetry splits full VSH labels into accessible and excluded sectors.}
\label{fig:ch3.364-sector-partition}
\end{figure}

\paragraph{Reportable coefficient screen (HOME).}
Let
\begin{equation}
\label{eq:ch3.reportable-coefficient-screen}
\widehat{c}_{\ell m}^{(\sigma)\,\mathrm{rep}}
=
\begin{cases}
[\mathbf{S}_{b}\widehat{\mathbf{c}}]_{\ell m}^{(\sigma)}, & (\ell,m,\sigma)\in\mathcal{L}_{\mathrm{acc}},\\
0, & \text{otherwise},
\end{cases}
\end{equation}
with observation rank applied afterward \cite{harrington2001,stratton1941}. Equation
\eqref{eq:ch3.reportable-coefficient-screen} is the coefficient law laboratories should publish: accessible,
normalized, and boundary-corrected \cite{hansen1988,harrington2001}.

\paragraph{Physical meaning of accessibility.}
A label can be mathematically present in Eq.~\eqref{eq:ch3.normalized-coefficient-resolution} yet inaccessible if
\(\mathcal{C}_{\mathrm{geo}}\) annihilates it Eq.~\eqref{eq:ch2.geometric-modal-annihilation} or if the receiver
rank cannot resolve it \cite{harrington2001,stratton1941}. Accessibility is therefore \emph{three-layered}: operator
export, geometric transport, and measurement resolution \cite{harrington2001}.

\paragraph{Worked illustration (half-space suppresses backward sheet).}
Over a PEC ground, backward half-space \(\mathcal{K}_{\mathrm{sup}}\) from Eq.~\eqref{eq:ch2.suppressed-angular-sheet-def}
removes directions regardless of \(\ell_{\max}\) \cite{stratton1941,jackson1999}. \(\mathcal{L}_{\mathrm{geo}}\) therefore
includes labels whose far patterns would peak in the suppressed sheet \cite{harrington2001}.

\paragraph{Problem (classify a silent coefficient).}
\textbf{Problem.} \(\widehat{c}_{3,-2}^{(\mathbf{M})}\) is below noise. Decide whether label is in
\(\mathcal{L}_{\mathrm{null}}\), \(\mathcal{L}_{\mathrm{geo}}\), \(\mathcal{L}_{\mathrm{meas}}\), or
\(\mathcal{L}_{\mathrm{acc}}\) \cite{harrington2001,stratton1941,devaney2004}.

\textbf{Step 1.} Test export: is \(\mathcal{R}_\omega[\Jimp]\neq\mathbf{0}\)? \cite{harrington2001}.

\textbf{Step 2.} Test geometry: is label in \(\mathcal{K}_{\mathrm{sup}}\)? \cite{stratton1941}.

\textbf{Step 3.} Test rank: is label resolved by \(\mathcal{O}_{\mathrm{meas}}\)? \cite{devaney2004}.

\textbf{Physical meaning.} Silence has distinct mechanisms; do not default to truncation \cite{harrington2001}.

\begin{table}[t]
\centering
\caption{Subsection~\ref{sec:ch364} broken-symmetry sector ledger.}
\label{tab:ch3.364-sector-ledger}
\small
\begin{tabular}{@{}p{0.18\linewidth}p{0.32\linewidth}p{0.40\linewidth}@{}}
\toprule
Sector & Mechanism & Chapter anchor \\
\midrule
\(\mathcal{L}_{\mathrm{acc}}\) & export + resolve & Eq.~\eqref{eq:ch3.accessible-VSH-label-set} \\
\(\mathcal{L}_{\mathrm{geo}}\) & symmetry forbid & Eq.~\eqref{eq:ch2.geometric-modal-annihilation} \\
\(\mathcal{L}_{\mathrm{null}}\) & NR source & Eq.~\eqref{eq:ch2.ker-R-omega-preview} \\
\(\mathcal{L}_{\mathrm{meas}}\) & finite rank & Eq.~\eqref{eq:ch2.reconstruction-rank-cap} \\
\bottomrule
\end{tabular}
\end{table}

\paragraph{Composition with truncation.}
Even accessible labels outside \(\ell\le\ell_{\max}\) reside in truncation remainder
Eq.~\eqref{eq:ch3.truncation-remainder-energy} \cite{hansen1988,harrington2001}. Subsection~\ref{sec:ch364} therefore
intersects \(\mathcal{L}_{\mathrm{acc}}\) with the truncated index set before watt assignment
\cite{harrington2001,poynting1884}.

\paragraph{Worked illustration (array symmetry reduces \(\mathcal{L}_{\mathrm{acc}}\)).}
A four-element square array has discrete rotational symmetry of order four; only certain \(m\) harmonics of the
array factor survive \cite{harrington2001,stratton1941}. \(\mathcal{L}_{\mathrm{geo}}\) grows compared to a single
dipole in free space \cite{hansen1988,harrington2001}.

\paragraph{Link to Subsection~\ref{sec:ch274}.}
Angular spectrum nullity law Eq.~\eqref{eq:ch2.angular-spectrum-nullspace-law} annihilates \emph{all} directions for
NR sources, whereas \(\mathcal{L}_{\mathrm{geo}}\) annihilates selected directions only \cite{harrington2001,stratton1941}.
Subsection~\ref{sec:ch364} keeps those mechanisms disjoint in Eq.~\eqref{eq:ch3.broken-symmetry-sector-partition}
\cite{harrington2001}.

\paragraph{Problem (count accessible labels at \(\ell\le2\)).}
\textbf{Problem.} Half-space PEC removes backward directions. Estimate \(|\mathcal{L}_{\mathrm{acc}}|\) among
\(N_{\mathrm{disc}}(2)=18\) labels \cite{harrington2001,jackson1999,stratton1941}.

\textbf{Step 1.} List geometric survivors using \(\mathcal{K}_{\mathrm{acc}}\) \cite{stratton1941}.

\textbf{Step 2.} Intersect with rank cap \(r\) \cite{harrington2001}.

\textbf{Step 3.} Report \(|\mathcal{L}_{\mathrm{acc}}|\le\min(18,r,|\mathcal{K}_{\mathrm{acc}}|)\) \cite{devaney2004}.

\textbf{Physical meaning.} Accessibility is an upper bound on identifiable multipoles \cite{harrington2001}.

\begin{figure}[t]
\centering
\ChOneFigBlock{%
\begin{tikzpicture}[ch1fig, node distance=7mm and 9mm]
  \node[ch1box] (op) {$\mathcal{R}_\omega$};
  \node[ch1box, right=9mm of op] (geo) {$\mathcal{C}_{\mathrm{geo}}$};
  \node[ch1box, right=9mm of geo] (meas) {$\mathcal{O}_{\mathrm{meas}}$};
  \node[ch1box, right=9mm of meas] (rep) {Eq.~\eqref{eq:ch3.reportable-coefficient-screen}};
  \draw[ch1flow] (op) -- (geo) -- (meas) -- (rep);
\end{tikzpicture}%
}
\caption{Subsection~\ref{sec:ch364} reporting pipeline: operator export, geometric screen, measurement rank, then
reportable coefficients.}
\label{fig:ch3.364-reporting-pipeline}
\end{figure}

\paragraph{Watt budget on accessible sectors only.}
Normalized power sums should use \(\mathcal{L}_{\mathrm{acc}}\) intersected with the truncated set:
\(\mathcal{P}_{\mathrm{rep}}=\sum_{(\ell,m,\sigma)\in\mathcal{L}_{\mathrm{acc}}\cap\mathcal{L}_{\le\ell_{\max}}}
|\widehat{c}_{\ell m}^{(\sigma)\,\mathrm{rep}}|^{2}\) \cite{poynting1884,harrington2001}. Including
\(\mathcal{L}_{\mathrm{geo}}\) or \(\mathcal{L}_{\mathrm{meas}}\) labels inflates or deflates budgets without physical
interpretation \cite{stratton1941,harrington2001}.

\paragraph{Reciprocity under broken symmetry.}
Eq.~\eqref{eq:ch3.reciprocity-coupling-symmetry} applies on the lossless exterior; \(\mathbf{S}_{b}\) from passive
boundaries preserves reciprocity when media are reciprocal \cite{stratton1941,harrington2001}. Subsection~\ref{sec:ch364}
requires paired \(\mathbf{S}_{b}\) on transmit and receive ports in link budgets Eq.~\eqref{eq:ch3.link-budget-normalized}
\cite{harrington2001}.

\paragraph{Worked illustration (NR source versus geometric null).}
\(\Jimp\in\ker\mathcal{R}_\omega\) yields \(\widehat{c}_{\ell m}^{(\sigma)}=0\) for \emph{all} labels
Eq.~\eqref{eq:ch2.angular-spectrum-nullspace-law} \cite{harrington2001,stratton1941}. A radiating source over a ground
plane may still excite \(\mathcal{L}_{\mathrm{acc}}\) while \(\mathcal{L}_{\mathrm{geo}}\) forbids backward sectors---distinct
from NR nullity \cite{harrington2001}.

\paragraph{Forward pointer to Chapter~\ref{ch:04}.}
Realizability asks which \(\widehat{\mathbf{c}}\) in \(\mathcal{L}_{\mathrm{acc}}\) are source-achievable on a given
\(\Omega_{s}\) \cite{harrington2001,stratton1941}. Subsection~\ref{sec:ch364} supplies the label taxonomy;
Chapter~\ref{ch:04} supplies the source inverse map \cite{harrington2001,devaney2004}.

\paragraph{Problem (verify partition exhaustiveness).}
\textbf{Problem.} Show every \((\ell,m,\sigma)\) lies in exactly one sector of
Eq.~\eqref{eq:ch3.broken-symmetry-sector-partition} when classification rules are applied in order null \(\to\) geo
\(\to\) meas \(\to\) acc \cite{harrington2001,stratton1941}.

\textbf{Step 1.} NR membership is tested first via \(\ker\mathcal{R}_\omega\) \cite{harrington2001}.

\textbf{Step 2.} Geometric forbidding is tested on exported states \cite{stratton1941}.

\textbf{Step 3.} Remaining unresolved labels occupy \(\mathcal{L}_{\mathrm{meas}}\); survivors are
\(\mathcal{L}_{\mathrm{acc}}\) \cite{devaney2004}.

\textbf{Physical meaning.} Ordered screening prevents double counting exclusion reasons \cite{harrington2001}.

\paragraph{DOS intersection.}
Eq.~\eqref{eq:ch3.cumulative-disc-DOS} counts mathematical channels; \(|\mathcal{L}_{\mathrm{acc}}|\) counts
reportable channels under broken symmetry \cite{harrington2001,hansen1988}. Subsection~\ref{sec:ch364} recommends
publishing both integers in campaign metadata \cite{harrington2001,stratton1941}.

\paragraph{Accessible watt budget (HOME).}
Reportable power on broken-symmetry campaigns is
\begin{equation}
\label{eq:ch3.accessible-watt-budget}
\mathcal{P}_{\mathrm{rep}}
=
\sum_{(\ell,m,\sigma)\in\mathcal{L}_{\mathrm{acc}}\cap\mathcal{L}_{\le\ell_{\max}}}
\big|\widehat{c}_{\ell m}^{(\sigma)\,\mathrm{rep}}\big|^{2},
\end{equation}
importing Eq.~\eqref{eq:ch3.reportable-coefficient-screen} and normalized flux gauge
Eq.~\eqref{eq:ch3.mode-power-from-coefficient} \cite{poynting1884,harrington2001,stratton1941}. Equation
\eqref{eq:ch3.accessible-watt-budget} is the Subsection~\ref{sec:ch364} watt law under broken symmetry: excluded
sectors do not enter sums even if mathematically present in Eq.~\eqref{eq:ch3.normalized-coefficient-resolution}
\cite{harrington2001}.

\paragraph{Worked illustration (half-space campaign metadata).}
Publish \(|\mathcal{L}_{\mathrm{acc}}|\), \(\ell_{\max}\), \(\mathbf{S}_{b}\) identifier, and
\(\mathcal{P}_{\mathrm{rep}}\) alongside raw pattern files \cite{harrington2001,stratton1941}. Without that tuple,
cross-laboratory comparisons confuse geometric suppression with source weakness \cite{hansen1988,harrington2001}.

\begin{figure}[t]
\centering
\ChOneFigBlock{%
\begin{tikzpicture}[ch1fig, node distance=6mm and 8mm]
  \node[ch1box, minimum width=22mm] (acc) {$\mathcal{L}_{\mathrm{acc}}$};
  \node[ch1box, right=8mm of acc, minimum width=22mm] (geo) {$\mathcal{L}_{\mathrm{geo}}$};
  \node[ch1box, right=8mm of geo, minimum width=22mm] (null) {$\mathcal{L}_{\mathrm{null}}$};
  \node[ch1box, right=8mm of null, minimum width=22mm] (meas) {$\mathcal{L}_{\mathrm{meas}}$};
  \node[ch1box, below=10mm of geo, minimum width=30mm] (part) {Eq.~\eqref{eq:ch3.broken-symmetry-sector-partition}};
  \draw[ch1flow] (part) -- (acc);
  \draw[ch1flow] (part) -- (geo);
  \draw[ch1flow] (part) -- (null);
  \draw[ch1flow] (part) -- (meas);
\end{tikzpicture}%
}
\caption{Subsection~\ref{sec:ch364} four-way partition of VSH labels under broken symmetry and finite observation.}
\label{fig:ch3.364-four-sector}
\end{figure}

\paragraph{Problem (compute \(\mathcal{P}_{\mathrm{rep}}\) with exclusions).}
\textbf{Problem.} Three accessible modes carry \(0.4\), \(0.3\), and \(0.2\,\mathrm{W}\); two geometrically forbidden modes
would carry \(0.05\,\mathrm{W}\) each if included. Find \(\mathcal{P}_{\mathrm{rep}}\) \cite{harrington2001,poynting1884,stratton1941}.

\textbf{Step 1.} Exclude \(\mathcal{L}_{\mathrm{geo}}\) from the sum \cite{stratton1941}.

\textbf{Step 2.} Sum accessible watts only \cite{poynting1884}.

\textbf{Step 3.} Obtain \(\mathcal{P}_{\mathrm{rep}}=0.9\,\mathrm{W}\) \cite{harrington2001}.

\textbf{Physical meaning.} Forbidden sectors must not inflate or deflate reported budgets \cite{harrington2001}.

\begin{table}[t]
\centering
\caption{Subsection~\ref{sec:ch364} campaign metadata for broken-symmetry reports.}
\label{tab:ch3.364-campaign-metadata}
\small
\begin{tabular}{@{}p{0.30\linewidth}p{0.58\linewidth}@{}}
\toprule
Field & Content \\
\midrule
Label partition & Eq.~\eqref{eq:ch3.broken-symmetry-sector-partition} \\
Boundary map & \(\mathbf{S}_{b}\) or \(\Gamma_{b}\) class \\
Reportable power & Eq.~\eqref{eq:ch3.accessible-watt-budget} \\
DOS versus accessible & \(N_{\mathrm{disc}}\) vs.\ \(|\mathcal{L}_{\mathrm{acc}}|\) \\
\bottomrule
\end{tabular}
\end{table}

\paragraph{Constraint composition restatement.}
Eq.~\eqref{eq:ch2.constraint-composition-law} orders null, geometric, and measurement deficits before accessibility is
declared \cite{harrington2001,stratton1941}. Subsection~\ref{sec:ch364} is the VSH label realization of that composition
law imported from Subsection~\ref{sec:ch274} \cite{harrington2001}.

\paragraph{Worked illustration (mobile platform symmetry).}
A moving vehicle breaks time-translation symmetry and reduces rotational symmetry to a subgroup; \(\mathcal{L}_{\mathrm{geo}}\)
grows while \(\mathcal{L}_{\mathrm{null}}\) is unchanged for fixed \(\Jimp\) \cite{harrington2001,stratton1941}. Accessibility
reports must name the platform frame used for \(m\) labels \cite{hansen1988,harrington2001}.

\paragraph{Problem (map \(\mathcal{L}_{\mathrm{meas}}\) to SNR threshold).}
\textbf{Problem.} Coefficients below receiver noise belong to \(\mathcal{L}_{\mathrm{meas}}\) even if geometry permits export
\cite{harrington2001,devaney2004,stratton1941}.

\textbf{Step 1.} Compare \(|\widehat{c}|\) to noise floor in \(\sqrt{\mathrm{W}}\) \cite{harrington2001}.

\textbf{Step 2.} Classify unresolved labels as \(\mathcal{L}_{\mathrm{meas}}\) \cite{devaney2004}.

\textbf{Physical meaning.} Accessibility requires resolvability, not existence alone \cite{harrington2001}.

\paragraph{Angular spectrum restatement.}
On \(\mathcal{K}_{\mathrm{acc}}\), continuous coefficients \(a(\widehat{\mathbf{k}})\) from
Eq.~\eqref{eq:ch2.exterior-continuous-spectral-integral} inherit the same partition: directions in
\(\mathcal{K}_{\mathrm{sup}}\) remain null in reports even when discrete \(\widehat{c}_{\ell m}^{(\sigma)}\) are used
\cite{stratton1941,harrington2001}. Subsection~\ref{sec:ch364} unifies discrete and continuous accessibility languages
imported from Subsection~\ref{sec:ch274} \cite{harrington2001}.

\paragraph{Worked illustration (link budget on \(\mathcal{L}_{\mathrm{acc}}\) only).}
Eq.~\eqref{eq:ch3.link-budget-normalized} should use coupling blocks restricted to
\(\mathcal{L}_{\mathrm{acc}}\times\mathcal{L}_{\mathrm{acc}}\) after \(\mathbf{S}_{b}\) correction
\cite{harrington2001,stratton1941}. Feeding geometrically forbidden ports inflates apparent coupling without physical
transport \cite{hansen1988,harrington2001}.

\paragraph{Reconstructability budget import.}
Eq.~\eqref{eq:ch2.reconstructability-budget} caps how many \(\mathcal{L}_{\mathrm{acc}}\) labels can be inverted jointly
\cite{harrington2001,stratton1941}. Accessibility is necessary but not sufficient for stable reconstruction: a label can
lie in \(\mathcal{L}_{\mathrm{acc}}\) yet remain ill-conditioned under finite SNR \cite{devaney2004,harrington2001}.
Subsection~\ref{sec:ch364} therefore pairs Eq.~\eqref{eq:ch3.accessible-VSH-label-set} with observation rank
Eq.~\eqref{eq:ch2.reconstruction-rank-cap} in every inversion report \cite{harrington2001}.

\paragraph{Problem (build minimal reportable set).}
\textbf{Problem.} Given \(\mathcal{L}_{\mathrm{geo}}\) and rank \(r\), construct a maximal subset
\(\mathcal{L}_{\mathrm{acc}}\subset\mathcal{L}_{\le\ell_{\max}}\) with \(|\mathcal{L}_{\mathrm{acc}}|\le r\)
\cite{harrington2001,devaney2004,stratton1941}.

\textbf{Step 1.} Remove \(\mathcal{L}_{\mathrm{geo}}\) and \(\mathcal{L}_{\mathrm{null}}\) \cite{stratton1941}.

\textbf{Step 2.} Sort remaining labels by SNR \cite{devaney2004}.

\textbf{Step 3.} Keep top \(r\) labels as \(\mathcal{L}_{\mathrm{acc}}\) \cite{harrington2001}.

\textbf{Physical meaning.} Reportable sets are rank-limited even after geometry is fixed \cite{harrington2001}.

\begin{figure}[t]
\centering
\ChOneFigBlock{%
\begin{tikzpicture}[ch1fig, node distance=7mm and 9mm]
  \node[ch1box] (pool) {$\mathcal{L}_{\le\ell_{\max}}$};
  \node[ch1box, right=10mm of pool] (scr) {screens};
  \node[ch1box, right=10mm of scr] (acc) {$\mathcal{L}_{\mathrm{acc}}$};
  \node[ch1box, below=10mm of scr] (rank) {rank $\le r$};
  \draw[ch1flow] (pool) -- (scr) -- (acc);
  \draw[ch1flow] (rank) -- (scr);
\end{tikzpicture}%
}
\caption{Subsection~\ref{sec:ch364} selection flow: geometric and null screens, then rank cap, yield reportable labels.}
\label{fig:ch3.364-selection-flow}
\end{figure}

\paragraph{Worked illustration (publish partition fractions).}
Report \(|\mathcal{L}_{\mathrm{acc}}|/|\mathcal{L}_{\le\ell_{\max}}|\) alongside
Eq.~\eqref{eq:ch3.accessible-watt-budget} so readers see what fraction of mathematical DOS was actually resolved
\cite{harrington2001,hansen1988}. A high DOS with low accessible fraction signals broken symmetry or rank limits, not
weak sources \cite{stratton1941,harrington2001}.

\begin{table}[t]
\centering
\caption{Subsection~\ref{sec:ch364} example partition counts at \(\ell_{\max}=2\), half-space PEC.}
\label{tab:ch3.364-example-counts}
\small
\begin{tabular}{@{}p{0.22\linewidth}p{0.18\linewidth}p{0.48\linewidth}@{}}
\toprule
Sector & count (illustrative) & note \\
\midrule
\(\mathcal{L}_{\mathrm{acc}}\) & 10 & forward half + rank \\
\(\mathcal{L}_{\mathrm{geo}}\) & 6 & backward sheet \\
\(\mathcal{L}_{\mathrm{meas}}\) & 2 & below SNR \\
\bottomrule
\end{tabular}
\end{table}

Subsection~\ref{sec:ch364} has classified accessible structure under broken symmetry through
Eqs.~\eqref{eq:ch3.accessible-VSH-label-set}--\eqref{eq:ch3.reportable-coefficient-screen}
\cite{harrington2001,stratton1941,jackson1999}. Section~\ref{sec:ch36} is complete: unbounded completeness,
truncation loss, boundary mixing, and accessibility partition close the representation-theoretic arc before
Section~\ref{sec:ch37} unifies discrete and continuous angular spectra \cite{hansen1988,harrington2001}.

The four subsections of Section~\ref{sec:ch36} establish when VSH bases are complete
Eqs.~\eqref{eq:ch3.unbounded-completeness-statement}--\eqref{eq:ch3.L2-convergence-VSH} in Subsection~\ref{sec:ch361};
how finite \(\ell_{\max}\) and compact supports destroy completeness
Eqs.~\eqref{eq:ch3.ellmax-truncation-projector}--\eqref{eq:ch3.truncation-remainder-energy} in Subsection~\ref{sec:ch362};
how boundaries mix orthogonal labels Eqs.~\eqref{eq:ch3.boundary-scattering-mixing-matrix}--\eqref{eq:ch3.image-mode-admixture}
in Subsection~\ref{sec:ch363}; and which labels remain accessible under reduced symmetry
Eqs.~\eqref{eq:ch3.accessible-VSH-label-set}--\eqref{eq:ch3.broken-symmetry-sector-partition} in
Subsection~\ref{sec:ch364} \cite{stratton1941,harrington2001,jackson1999}. Section~\ref{sec:ch37} now compares
discrete multipole indices with continuous direction spectra on the same \(\mathcal{R}_\omega\) image
\cite{hansen1988,jackson1999}.

\section{Discrete and Continuous Angular Spectra}
\label{sec:ch37}

Laboratory antennas often report discrete modal indices; free-space scattering theory prefers continuous
angular spectra. Section~\ref{sec:ch37} unifies those pictures as representation choices on the same
\(\mathcal{R}_\omega\) image \cite{jackson1999,stratton1941}. Subsection~\ref{sec:ch371} develops angular
spectra in unbounded space, importing plane-wave radiation operators and accessibility budgets from
Section~\ref{sec:ch29}. Subsection~\ref{sec:ch372} explains quantization induced by boundaries, using
mode-mixing results from Subsection~\ref{sec:ch363}. Subsection~\ref{sec:ch373} treats integer and
non-integer angular indices without collapsing them into informal beam orders. Subsection~\ref{sec:ch374}
connects phase singularities and topological limits to the phase-topology material of
Subsection~\ref{sec:ch01.2.2} \cite{allen1992,stratton1941}.


\subsection{Angular spectra in unbounded space}
\label{sec:ch371}
% TOC tag: [HOME][USE SS2.9.1][USE SS2.9.4]

An \emph{angular spectrum} is a direction-resolved amplitude density on outgoing wave vectors; \emph{unbounded
space} means no enclosing boundary quantizes those directions into discrete normal modes
\cite{jackson1999,stratton1941,harrington2001}. Subsection~\ref{sec:ch371} is the HOME coordinate bridge that places
Chapter~\ref{ch:02} plane-wave spectra Subsection~\ref{sec:ch291} and far-zone accessibility Subsection~\ref{sec:ch294}
on the same \(\mathcal{R}_\omega\) image as VSH coefficients Eq.~\eqref{eq:ch3.normalized-coefficient-resolution}
\cite{hansen1988,harrington2001}.

\paragraph{Standing imports (no re-derivation).}
Plane-wave angular spectrum Eqs.~\eqref{eq:ch2.plane-wave-angular-spectrum-def}--\eqref{eq:ch2.plane-wave-z-propagation},
radiation operator Eq.~\eqref{eq:ch2.plane-wave-radiation-operator}, exterior continuous integral
Eq.~\eqref{eq:ch2.exterior-continuous-spectral-integral}, far angular accessibility
Eq.~\eqref{eq:ch2.far-angular-accessibility-def}, accessible angular budget
Eq.~\eqref{eq:ch2.accessible-angular-budget}, continuous measure
Eq.~\eqref{eq:ch3.continuous-measure-flux-fix}, and plane-wave bridge
Eq.~\eqref{eq:ch3.plane-wave-continuum-bridge} enter by reference \cite{stratton1941,jackson1999,harrington2001}.
Subsection~\ref{sec:ch371} does not re-derive Fourier conventions or outgoing \(k_{z}\) gates
\cite{stratton1941,harrington2001}.

\paragraph{Assumptions.}
Fix \(\omega>0\), homogeneous reciprocal exterior media, outgoing \(\mathcal{F}_{\mathrm{rad}}\), and phasor time
dependence \(\exp(-\jj\omega t)\) \cite{stratton1941,harrington2001}. Angular spectra are defined on the propagating
outgoing sheet \(k_{t}\le k_{0}\) unless evanescent sheets are explicitly labeled
\cite{stratton1941,jackson1999}.

\paragraph{Unbounded angular spectrum density (HOME).}
Let \(a(\widehat{\mathbf{k}})\) denote the scalar outgoing angular spectrum on \(S^{2}\) imported from
Eq.~\eqref{eq:ch2.exterior-continuous-spectral-integral}. Subsection~\ref{sec:ch371} fixes the watt-calibrated density
\begin{equation}
\label{eq:ch3.unbounded-angular-spectrum-def}
\mathcal{A}(\widehat{\mathbf{k}})
=
|a(\widehat{\mathbf{k}})|^{2}\,d\mu(\widehat{\mathbf{k}}),
\qquad
d\mu=\frac{P_{0}}{4\pi}\,d\Omega,
\end{equation}
with \(P_{0}\) from Eq.~\eqref{eq:ch3.continuous-measure-flux-fix} \cite{jackson1999,harrington2001,stratton1941}.
Equation~\eqref{eq:ch3.unbounded-angular-spectrum-def} states the physical meaning of angular spectra in Volume~I:
\emph{power per solid-angle increment} on the unit sphere, not informal ``beam shape'' plots
\cite{poynting1884,harrington2001}.

\paragraph{Discrete--continuous bridge (HOME).}
VSH coefficients and angular spectra describe the same \(\mathbf{v}_\omega\in\mathcal{F}_{\mathrm{rad}}\) when
\begin{equation}
\label{eq:ch3.discrete-continuous-spectrum-bridge}
\mathbf{v}_\omega
=
\sum_{\ell m\sigma}\widehat{c}_{\ell m}^{(\sigma)}\widehat{\mathbf{u}}_{\ell m}^{(\sigma)}
\;\Longleftrightarrow\;
\int_{S^{2}} a(\widehat{\mathbf{k}})\,\mathbf{e}(\widehat{\mathbf{k}})\,d\mu(\widehat{\mathbf{k}}),
\end{equation}
with \(\mathbf{e}(\widehat{\mathbf{k}})\) the unit outgoing plane-wave polarization bundle from
Eq.~\eqref{eq:ch3.plane-wave-continuum-bridge} \cite{hansen1988,jackson1999,stratton1941}. Equation
\eqref{eq:ch3.discrete-continuous-spectrum-bridge} is the Subsection~\ref{sec:ch371} unification claim: discrete
\(\ell\) and continuous \(\hat{\mathbf{k}}\) are coordinates, not different physics \cite{harrington2001}.

\begin{figure}[t]
\centering
\ChOneFigBlock{%
\begin{tikzpicture}[ch1fig, node distance=7mm and 9mm]
  \node[ch1box] (vsh) {VSH $\widehat{c}_{\ell m}^{(\sigma)}$};
  \node[ch1box, right=10mm of vsh] (bridge) {Eq.~\eqref{eq:ch3.discrete-continuous-spectrum-bridge}};
  \node[ch1box, right=10mm of bridge] (ang) {$a(\widehat{\mathbf{k}})$};
  \node[ch1box, below=10mm of bridge] (mu) {Eq.~\eqref{eq:ch3.unbounded-angular-spectrum-def}};
  \draw[ch1flow] (vsh) -- (bridge) -- (ang);
  \draw[ch1flow] (mu) -- (ang);
\end{tikzpicture}%
}
\caption{Subsection~\ref{sec:ch371} bridge: multipole coefficients and direction spectra are dual coordinates on
\(\mathcal{F}_{\mathrm{rad}}\) in unbounded space.}
\label{fig:ch3.371-spectrum-bridge}
\end{figure}

\begin{theorem}[Unbounded angular representation (HOME)]
\label{thm:ch3.unbounded-angular-spectrum}
On exterior domains, every \(\mathbf{v}_\omega\in\mathcal{F}_{\mathrm{rad}}\) admits both VSH expansion
Eq.~\eqref{eq:ch3.L2-convergence-VSH} and angular spectrum Eq.~\eqref{eq:ch3.unbounded-angular-spectrum-def} with
Parseval closure Eq.~\eqref{eq:ch3.continuous-parseval-flux-fixed} \cite{hansen1988,jackson1999,stratton1941,kato1995}.
\end{theorem}
\begin{proof}[Proof sketch]
Combine Theorem~\ref{thm:ch3.unbounded-domain-completeness} with plane-wave continuum bridge
Eq.~\eqref{eq:ch3.plane-wave-continuum-bridge} and Section~\ref{sec:ch35} measure \cite{hansen1988,harrington2001}.
% PROOF: AUTHOR REVIEW
\end{proof}

\paragraph{Lateral spectrum linkage.}
Eq.~\eqref{eq:ch2.radiative-lateral-spectrum-def} on a plane \(z=z_{0}\) is the Fourier chart of
\(a(\widehat{\mathbf{k}})\) near the forward hemisphere \cite{stratton1941,jackson1999}. Subsection~\ref{sec:ch371}
therefore connects compact-range plane scans to \(S^{2}\) spectra through outgoing \(k_{z}\) selection
Eq.~\eqref{eq:ch2.plane-wave-z-propagation} \cite{harrington2001,stratton1941}.

\paragraph{Worked illustration (Gaussian beam spectrum).}
A Gaussian \(\widetilde{\E}(k_{x},k_{y})\propto\exp[-(k_{x}^{2}+k_{y}^{2})/k_{0}^{2}w_{0}^{2}]\) concentrates
\(\mathcal{A}(\widehat{\mathbf{k}})\) near broadside \cite{jackson1999,harrington2001}. Watt integral
\(\int\mathcal{A}\) equals total export when \(d\mu\) is flux-fixed \cite{poynting1884,stratton1941}.

\paragraph{Problem (convert VSH dipole to angular sheet).}
\textbf{Problem.} A normalized \(\widehat{\mathbf{u}}_{10}^{(\mathbf{N})}\) mode exports \(1\,\mathrm{W}\). Sketch
\(\mathcal{A}(\widehat{\mathbf{k}})\) qualitatively on \(S^{2}\) \cite{hansen1988,jackson1999,harrington2001}.

\textbf{Step 1.} Dipole radiation peaks along \(\theta=\pi/2\) \cite{jackson1999}.

\textbf{Step 2.} Map to \(d\mu\) via Eq.~\eqref{eq:ch3.unbounded-angular-spectrum-def} \cite{stratton1941}.

\textbf{Step 3.} Integrate to \(1\,\mathrm{W}\) \cite{poynting1884}.

\textbf{Physical meaning.} Multipole labels induce anisotropic angular spectra, not isotropic noise \cite{harrington2001}.

\paragraph{Accessibility import.}
Eq.~\eqref{eq:ch2.accessible-angular-budget} restricts which \(\widehat{\mathbf{k}}\) sheets enter reports after
geometry and measurement \cite{harrington2001,stratton1941}. Subsection~\ref{sec:ch371} applies that budget to
continuous spectra on unbounded domains before boundary quantization in Subsection~\ref{sec:ch372}
\cite{hansen1988,harrington2001}.

\begin{table}[t]
\centering
\caption{Subsection~\ref{sec:ch371} angular spectrum objects on unbounded domains.}
\label{tab:ch3.371-spectrum-objects}
\small
\begin{tabular}{@{}p{0.24\linewidth}p{0.38\linewidth}p{0.30\linewidth}@{}}
\toprule
Object & Equation & Unit \\
\midrule
Spectrum density & Eq.~\eqref{eq:ch3.unbounded-angular-spectrum-def} & W per \(d\mu\) \\
Bridge & Eq.~\eqref{eq:ch3.discrete-continuous-spectrum-bridge} & coordinate map \\
Parseval & Eq.~\eqref{eq:ch3.continuous-parseval-flux-fixed} & watts \\
\bottomrule
\end{tabular}
\end{table}

\paragraph{Outgoing filter discipline.}
Only \(\Pi_{\mathrm{rad}}\)-qualified content enters Eq.~\eqref{eq:ch2.plane-wave-radiation-operator}
\cite{harrington2001,stratton1941}. Evanescent lateral modes may be large near sources but lie outside
\(\mathcal{F}_{\mathrm{rad}}\) export \cite{stratton1941,jackson1999}. Subsection~\ref{sec:ch371} treats them as
reactive sheets, not angular spectrum watts \cite{harrington2001}.

\paragraph{Worked illustration (hemisphere power budget).}
Restricting \(\mathcal{A}\) to a forward hemisphere integrates to half the full-sphere watts when
\(|a|^{2}\) is uniform on that hemisphere \cite{jackson1999,harrington2001}. Laboratory main beams are often
hemisphere-dominant even in unbounded theory \cite{stratton1941,hansen1988}.

\paragraph{Problem (audit lateral FFT without \(\Pi_{\mathrm{rad}}\)).}
\textbf{Problem.} A raw FFT of \(\E\) near a source overestimates watts by \(20\%\). Diagnose
\cite{harrington2001,stratton1941,jackson1999}.

\textbf{Step 1.} Compare Eq.~\eqref{eq:ch2.radiative-lateral-spectrum-def} to raw FFT \cite{stratton1941}.

\textbf{Step 2.} Identify evanescent/reactive admixture \cite{harrington2001}.

\textbf{Step 3.} Apply \(\Pi_{\mathrm{rad}}\) before spectrum formation \cite{stratton1941}.

\textbf{Physical meaning.} Angular spectra are radiative objects, not near-field FFTs \cite{harrington2001}.

\paragraph{Far-zone amplitude map.}
Eq.~\eqref{eq:ch2.far-zone-power-from-pattern} ties \(|\E_{0}(\theta,\phi)|^{2}\) to \(a(\widehat{\mathbf{k}})\) on
\(S^{2}\) \cite{stratton1941,jackson1999}. Subsection~\ref{sec:ch371} identifies pattern plots as polar charts of
\(|a|^{2}\), while Eq.~\eqref{eq:ch3.unbounded-angular-spectrum-def} supplies watt measure \cite{poynting1884,harrington2001}.

\begin{figure}[t]
\centering
\ChOneFigBlock{%
\begin{tikzpicture}[ch1fig, node distance=7mm and 9mm]
  \node[ch1box] (lat) {Eq.~\eqref{eq:ch2.radiative-lateral-spectrum-def}};
  \node[ch1box, right=10mm of lat] (s2) {$S^{2}$ spectrum};
  \node[ch1box, right=10mm of s2] (acc) {Eq.~\eqref{eq:ch2.accessible-angular-budget}};
  \draw[ch1flow] (lat) -- (s2) -- (acc);
\end{tikzpicture}%
}
\caption{Subsection~\ref{sec:ch371} pipeline: lateral FFT on outgoing fields, lift to \(S^{2}\), then apply accessibility
budget.}
\label{fig:ch3.371-lateral-to-sphere}
\end{figure}

\paragraph{Reciprocity on angular sheets.}
Lossless reciprocity pairs transmit and receive spectra on the same \(d\mu\) \cite{stratton1941,harrington2001}.
Subsection~\ref{sec:ch371} requires identical \(P_{0}\) in Eq.~\eqref{eq:ch3.unbounded-angular-spectrum-def} on both
ports before comparing \(|a(\widehat{\mathbf{k}})|\) \cite{hansen1988,harrington2001}.

\paragraph{Worked illustration (dual VSH and plane-wave fits).}
Fitting the same far pattern in VSH and angular spectrum coordinates yields the same \(\mathcal{P}_{\mathrm{rad}}\) when
both expansions converge and share \(\mathcal{N}\) \cite{hansen1988,harrington2001,stratton1941}. Residual differences
indicate truncation or windowing, not coordinate failure \cite{harrington2001}.

\paragraph{Problem (Parseval on \(S^{2}\)).}
\textbf{Problem.} Show \(\int_{S^{2}}|a|^{2}d\mu=\mathcal{P}_{\mathrm{rad}}\) for normalized outgoing states
\cite{jackson1999,harrington2001,stratton1941}.

\textbf{Step 1.} Use Eq.~\eqref{eq:ch3.continuous-parseval-flux-fixed} \cite{hansen1988}.

\textbf{Step 2.} Identify \(d\mu\) with Eq.~\eqref{eq:ch3.solid-angle-jacobian-measure} \cite{jackson1999}.

\textbf{Physical meaning.} Angular spectra participate in the same watt calculus as VSH sums \cite{poynting1884}.

\paragraph{Transport gate reminder.}
Eq.~\eqref{eq:ch2.far-angular-transport-gate} governs whether reported spectral content is transport-qualified
\cite{stratton1941,harrington2001}. Subsection~\ref{sec:ch371} cites the gate without reopening Chapter~\ref{ch:01}
topology proofs \cite{harrington2001}.

\paragraph{Spectral index alignment.}
Eq.~\eqref{eq:ch3.spectral-index-partition} prevents double-counting when hybrid sums mix discrete and continuous
indices \cite{stratton1941,harrington2001}. Subsection~\ref{sec:ch371} uses the partition before assembling hybrid
budgets in Eq.~\eqref{eq:ch3.hybrid-DOS-count} \cite{hansen1988,harrington2001}.

\paragraph{Angular spectrum operator closure (HOME).}
Define the spectral analysis operator on \(\mathcal{F}_{\mathrm{rad}}\) by
\begin{equation}
\label{eq:ch3.angular-spectrum-operator}
(\mathcal{A}a)(\widehat{\mathbf{k}})
=
\big\langle \mathbf{e}(\widehat{\mathbf{k}}),\,\mathbf{v}_\omega\big\rangle_{\mathrm{rad}},
\qquad
\mathbf{v}_\omega=\int a(\widehat{\mathbf{k}}')\,\mathbf{e}(\widehat{\mathbf{k}}')\,d\mu(\widehat{\mathbf{k}}'),
\end{equation}
importing Eq.~\eqref{eq:ch3.plane-wave-continuum-bridge} polarization factors \(\mathbf{e}\)
\cite{jackson1999,hansen1988,stratton1941}. Equation~\eqref{eq:ch3.angular-spectrum-operator} states that
\(a(\widehat{\mathbf{k}})\) is a Riesz coefficient on direction modes, dual to VSH projections
Eq.~\eqref{eq:ch3.normalized-coefficient-resolution} \cite{harrington2001,kato1995}.

\paragraph{Worked illustration (main beam solid angle).}
Integrating \(\mathcal{A}\) over a main-lobe cone \(\Omega_{\mathrm{main}}\) yields partial power
\(P_{\mathrm{main}}=\int_{\Omega_{\mathrm{main}}}|a|^{2}d\mu\) \cite{jackson1999,poynting1884,harrington2001}. Sidelobe
continuum outside \(\Omega_{\mathrm{main}}\) carries the remainder without changing total export
\cite{stratton1941,hansen1988}.

\paragraph{Problem (discretize \(\mathcal{A}\) on a \(6\times6\) direction grid).}
\textbf{Problem.} Replace \(\int|a|^{2}d\mu\) by \(\sum_{q}w_{q}|a_{q}|^{2}\) on a uniform \(S^{2}\) grid. Choose \(w_{q}\)
for \(P_{0}=1\,\mathrm{W}\) \cite{harrington2001,jackson1999,stratton1941}.

\textbf{Step 1.} Cell solid angle \(\Delta\Omega_{q}\approx4\pi/36\) \cite{jackson1999}.

\textbf{Step 2.} Set \(w_{q}=\Delta\Omega_{q}/(4\pi)\) in Eq.~\eqref{eq:ch3.solid-angle-jacobian-measure} \cite{hansen1988}.

\textbf{Step 3.} Calibrate against shell flux \cite{poynting1884}.

\textbf{Physical meaning.} Direction grids approximate \(d\mu\), not a new physics model \cite{harrington2001}.

\begin{figure}[t]
\centering
\ChOneFigBlock{%
\begin{tikzpicture}[ch1fig, node distance=7mm and 9mm]
  \node[ch1box] (op) {Eq.~\eqref{eq:ch3.angular-spectrum-operator}};
  \node[ch1box, right=10mm of op] (dens) {Eq.~\eqref{eq:ch3.unbounded-angular-spectrum-def}};
  \node[ch1box, right=10mm of dens] (par) {Eq.~\eqref{eq:ch3.continuous-parseval-flux-fixed}};
  \draw[ch1flow] (op) -- (dens) -- (par);
\end{tikzpicture}%
}
\caption{Subsection~\ref{sec:ch371} operator chain: direction coefficients, watt density, Parseval integral.}
\label{fig:ch3.371-operator-chain}
\end{figure}

\paragraph{Retained and suppressed sheets (imported).}
Eqs.~\eqref{eq:ch2.retained-angular-amplitude-def} and \eqref{eq:ch2.suppressed-angular-amplitude-def} split \(a(\widehat{\mathbf{k}})\)
before accessibility screens \cite{stratton1941,harrington2001}. Subsection~\ref{sec:ch371} applies those splits on unbounded
domains prior to boundary quantization \cite{hansen1988,harrington2001}.

\paragraph{Worked illustration (plane-wave filter).}
Eq.~\eqref{eq:ch2.plane-wave-accessibility-filter} zeroes directions in \(\mathcal{K}_{\mathrm{sup}}\) while leaving
\(\mathcal{K}_{\mathrm{acc}}\) intact \cite{harrington2001,stratton1941}. Angular spectrum reports must declare the filter
or risk interpreting geometric nulls as source weakness \cite{hansen1988,harrington2001}.

\begin{table}[t]
\centering
\caption{Subsection~\ref{sec:ch371} unbounded angular spectrum checklist.}
\label{tab:ch3.371-checklist}
\small
\begin{tabular}{@{}p{0.08\linewidth}p{0.44\linewidth}p{0.36\linewidth}@{}}
\toprule
Step & Test & Pass \\
\midrule
1 & Outgoing & Eq.~\eqref{eq:ch3.outgoing-inheritance-criterion} \\
2 & Measure & Eq.~\eqref{eq:ch3.unbounded-angular-spectrum-def} \\
3 & Parseval & Eq.~\eqref{eq:ch3.continuous-parseval-flux-fixed} \\
4 & Access & Eq.~\eqref{eq:ch2.accessible-angular-budget} \\
\bottomrule
\end{tabular}
\end{table}

\paragraph{Problem (compare lateral and spherical peaks).}
\textbf{Problem.} A lateral FFT peak at \((k_{x},k_{y})=(k_{0}\sin\theta_{0},0)\) maps to which \(\widehat{\mathbf{k}}\) on \(S^{2}\)?
\cite{jackson1999,stratton1941,harrington2001}.

\textbf{Step 1.} Identify \(\theta=\theta_{0}\), \(\phi=0\) \cite{jackson1999}.

\textbf{Step 2.} Read \(\mathcal{A}(\widehat{\mathbf{k}})\) peak location \cite{hansen1988}.

\textbf{Physical meaning.} Lateral spectra are coordinate charts of \(S^{2}\) spectra \cite{stratton1941,harrington2001}.

\paragraph{Hybrid coordinate reporting.}
Campaigns may publish both \(\widehat{c}_{\ell m}^{(\sigma)}\) and \(a(\widehat{\mathbf{k}})\) tied by
Theorem~\ref{thm:ch3.unbounded-angular-spectrum} \cite{hansen1988,harrington2001}. Subsection~\ref{sec:ch371} requires
shared \(\mathcal{N}\) and outgoing class on both sides \cite{stratton1941,poynting1884}.

\paragraph{Far-field Jacobian reminder.}
Mapping \(\widetilde{\E}(k_{x},k_{y})\) to \(a(\widehat{\mathbf{k}})\) requires the outgoing branch and Jacobian factor on
\(S^{2}\) \cite{jackson1999,stratton1941}. Subsection~\ref{sec:ch371} treats Jacobian errors as measure errors in
Eq.~\eqref{eq:ch3.unbounded-angular-spectrum-def}, not as new source physics \cite{harrington2001,hansen1988}.

\paragraph{Worked illustration (broadside versus endfire).}
Endfire peaks map to \(\widehat{\mathbf{k}}\) near \(\theta=0,\pi\); broadside peaks near \(\theta=\pi/2\)
\cite{jackson1999,harrington2001}. Accessibility Eq.~\eqref{eq:ch2.accessible-angular-budget} may remove one hemisphere from
reports even when \(\mathcal{A}\) is mathematically nonzero there \cite{stratton1941,hansen1988}.

\paragraph{Problem (windowed spectrum energy).}
\textbf{Problem.} A Hann window on \(\widetilde{\E}(k_{x},k_{y})\) reduces total reconstructed watts by \(15\%\). Explain
\cite{harrington2001,jackson1999,stratton1941}.

\textbf{Step 1.} Windowing is a measurement taper, not \(\mathcal{K}_{\mathrm{sup}}\) \cite{harrington2001}.

\textbf{Step 2.} Compare windowed integral to Eq.~\eqref{eq:ch3.continuous-parseval-flux-fixed} \cite{poynting1884}.

\textbf{Physical meaning.} Tapers redistribute spectral energy without violating Parseval when applied consistently
\cite{stratton1941,hansen1988}.

\paragraph{Subsection~\ref{sec:ch294} import closure.}
Angular accessibility screens from Subsection~\ref{sec:ch294} act on \(a(\widehat{\mathbf{k}})\) after formation of
Eq.~\eqref{eq:ch3.angular-spectrum-operator} \cite{harrington2001,stratton1941}. Subsection~\ref{sec:ch371} positions
unbounded spectra \emph{before} those screens in the Chapter~\ref{ch:02} pipeline \cite{hansen1988,harrington2001}.

\begin{figure}[t]
\centering
\ChOneFigBlock{%
\begin{tikzpicture}[ch1fig, node distance=7mm and 9mm]
  \node[ch1box] (s2) {$S^{2}$};
  \node[ch1box, right=10mm of s2] (acc) {$\mathcal{K}_{\mathrm{acc}}$};
  \node[ch1box, right=10mm of acc] (rep) {reported $\mathcal{A}$};
  \draw[ch1flow] (s2) -- (acc) -- (rep);
\end{tikzpicture}%
}
\caption{Subsection~\ref{sec:ch371} accessibility screen on the unbounded direction sphere before laboratory reporting.}
\label{fig:ch3.371-accessibility-screen}
\end{figure}

\paragraph{Energy density on \(S^{2}\) (expanded).}
The scalar density \(|a(\widehat{\mathbf{k}})|^{2}\) is not power per steradian until multiplied by \(d\mu\) in
Eq.~\eqref{eq:ch3.unbounded-angular-spectrum-def} \cite{jackson1999,poynting1884}. One often plots \(|a|^{2}\) versus
\(\theta\) without \( \sin\theta\,d\theta\,d\phi\) weighting, underestimating equatorial contribution relative to poles
\cite{harrington2001,stratton1941,hansen1988}. Subsection~\ref{sec:ch371} therefore distinguishes \emph{amplitude density plots}
from \emph{watt-integrated} spectra \cite{poynting1884,harrington2001}.

\paragraph{Worked illustration (equatorial weighting).}
Uniform \(|a|^{2}\) on \(\theta\in[\pi/4,3\pi/4]\) integrates to \(\int_{\pi/4}^{3\pi/4}\int_{0}^{2\pi}|a|^{2}\sin\theta\,d\theta\,d\phi
= |a|^{2}\cdot\pi\sqrt{2}\) steradians-weighted measure before \(P_{0}\) scaling \cite{jackson1999,harrington2001}. Pole-dominated
patterns carry less equatorial weight even when \(|a|^{2}\) peaks at \(\theta=0\) \cite{stratton1941,hansen1988}.

\paragraph{Problem (integrate \(\mathcal{A}\) over a cone).}
\textbf{Problem.} For constant \(|a|^{2}=A_{0}\) on a cone \(\theta\le\theta_{0}\), find \(\int\mathcal{A}\) with \(P_{0}=1\,\mathrm{W}\)
\cite{jackson1999,harrington2001,stratton1941}.

\textbf{Step 1.} Integrate \(\sin\theta\) from \(0\) to \(\theta_{0}\) \cite{jackson1999}.

\textbf{Step 2.} Multiply by \(2\pi\) and \(P_{0}/(4\pi)\) \cite{poynting1884}.

\textbf{Step 3.} Obtain \(P=(1-\cos\theta_{0})/2\) watts for unit \(A_{0}\) normalization \cite{harrington2001}.

\textbf{Physical meaning.} Cone subtense sets partial power budgets on \(S^{2}\) \cite{stratton1941,hansen1988}.

\begin{figure}[t]
\centering
\ChOneFigBlock{%
\begin{tikzpicture}[ch1fig, node distance=7mm and 9mm]
  \node[ch1box] (plot) {$|a|^{2}$ plot};
  \node[ch1box, right=10mm of plot] (weight) {$\sin\theta\,d\Omega$};
  \node[ch1box, right=10mm of weight] (watts) {Eq.~\eqref{eq:ch3.unbounded-angular-spectrum-def}};
  \draw[ch1flow] (plot) -- (weight) -- (watts);
\end{tikzpicture}%
}
\caption{Subsection~\ref{sec:ch371} measure correction: amplitude plots require \(\sin\theta\) weighting before watt integration.}
\label{fig:ch3.371-sintheta-weight}
\end{figure}

\paragraph{Campaign archive fields (Subsection~\ref{sec:ch371}).}
Archive \(P_{0}\), outgoing certificate, lateral FFT grid, \(S^{2}\) quadrature weights, and
Eq.~\eqref{eq:ch2.accessible-angular-budget} mask with every \(a(\widehat{\mathbf{k}})\) file
\cite{harrington2001,stratton1941,hansen1988}. Reproducibility of angular spectra requires the measure, not only amplitudes
\cite{harrington2001,jackson1999}.

Subsection~\ref{sec:ch371} has unified angular spectra on unbounded domains through
Eqs.~\eqref{eq:ch3.unbounded-angular-spectrum-def}--\eqref{eq:ch3.discrete-continuous-spectrum-bridge}
\cite{jackson1999,stratton1941,harrington2001}. Subsection~\ref{sec:ch372} explains how boundaries quantize those continua
into discrete normal modes \cite{hansen1988}.


\subsection{Quantization from boundary constraints}
\label{sec:ch372}
% TOC tag: [HOME][USE SS3.6.3]

\emph{Quantization} means a continuous direction or multipole parameter becomes a discrete normal-mode index when
boundary conditions compactify or break the symmetry group \cite{stratton1941,jackson1999,harrington2001}.
\emph{Boundary constraints} are the linear relations imposed by conductors, interfaces, and cavity walls on tangential
fields \cite{stratton1941,hansen1988}. Subsection~\ref{sec:ch372} is the HOME spectral-discretization layer importing
mode mixing Eq.~\eqref{eq:ch3.boundary-scattering-mixing-matrix} from Subsection~\ref{sec:ch363}
\cite{harrington2001,stratton1941}.

\paragraph{Standing imports.}
Boundary mixing matrix Eq.~\eqref{eq:ch3.boundary-scattering-mixing-matrix}, image admixture
Eq.~\eqref{eq:ch3.image-mode-admixture}, cavity discrete spectrum Eq.~\eqref{eq:ch2.cavity-discrete-radiation-spectrum},
and geometric modal annihilation Eq.~\eqref{eq:ch2.geometric-modal-annihilation} enter by reference
\cite{stratton1941,harrington2001,jackson1999}. Subsection~\ref{sec:ch372} does not re-derive cavity Green functions
\cite{hansen1988,stratton1941}.

\paragraph{Assumptions.}
Consider configurations where \(\Gamma_{b}\) reduces \(\mathrm{SO}(3)\) symmetry and replaces continuous direction
degeneracy with discrete normal modes \(\{q\}\) \cite{stratton1941,harrington2001}. Quantization is analyzed on
\(\mathcal{F}_{\mathrm{rad}}\) after outgoing selection \cite{harrington2001,kato1995}.

\paragraph{Boundary-quantized angular spectrum (HOME).}
Let \(\{a_{q}\}_{q\in\mathcal{Q}}\) denote discrete angular amplitudes induced by \(\Gamma_{b}\). Subsection~\ref{sec:ch372}
defines
\begin{equation}
\label{eq:ch3.boundary-quantized-spectrum}
\mathbf{v}_\omega
=
\sum_{q\in\mathcal{Q}}
a_{q}\,\mathbf{e}_{q},
\qquad
\mathbf{e}_{q}\ \text{orthonormal on }\mathcal{F}_{\mathrm{rad}}\ \text{under }\Gamma_{b},
\end{equation}
replacing continuous integral Eq.~\eqref{eq:ch3.discrete-continuous-spectrum-bridge} when boundaries compactify the
problem \cite{stratton1941,jackson1999,harrington2001}. Equation~\eqref{eq:ch3.boundary-quantized-spectrum} is the
Subsection~\ref{sec:ch372} mathematical statement of quantization: the continuous \(S^{2}\) chart is replaced by a
countable mode set compatible with \(\Gamma_{b}\) \cite{hansen1988,harrington2001}.

\paragraph{Cavity angular index (HOME).}
For enclosed volumes, import Eq.~\eqref{eq:ch2.cavity-discrete-radiation-spectrum} and define the Chapter~\ref{ch:03}
index map
\begin{equation}
\label{eq:ch3.cavity-angular-index}
q\;\longleftrightarrow\;(\ell,m,n_{r}),
\end{equation}
with radial quantum number \(n_{r}\) from cavity separation and \((\ell,m)\) the surviving angular labels under
\(\Gamma_{b}\) \cite{stratton1941,jackson1999}. Equation~\eqref{eq:ch3.cavity-angular-index} warns that open-region
\(\ell\) alone is insufficient inside cavities: radial indices are part of the angular spectrum once boundaries quantize
all three coordinates \cite{harrington2001,hansen1988}.

\begin{figure}[t]
\centering
\ChOneFigBlock{%
\begin{tikzpicture}[ch1fig, node distance=7mm and 9mm]
  \node[ch1box] (cont) {continuous $S^{2}$};
  \node[ch1box, right=10mm of cont] (bnd) {$\Gamma_{b}$};
  \node[ch1box, right=10mm of bnd] (disc) {Eq.~\eqref{eq:ch3.boundary-quantized-spectrum}};
  \node[ch1box, below=10mm of bnd] (mix) {Eq.~\eqref{eq:ch3.boundary-scattering-mixing-matrix}};
  \draw[ch1flow] (cont) -- (bnd) -- (disc);
  \draw[ch1flow] (mix) -- (bnd);
\end{tikzpicture}%
}
\caption{Subsection~\ref{sec:ch372} quantization chain: boundaries turn continuous direction spectra into discrete
normal modes.}
\label{fig:ch3.372-quantization-chain}
\end{figure}

\paragraph{Physical meaning of quantization.}
Quantization is not a numerical grid choice: it is imposed by boundary eigenvalue problems \cite{stratton1941,jackson1999}.
A open-region VSH mode \(\widehat{\mathbf{u}}_{\ell m}^{(\sigma)}\) may not satisfy \(\Gamma_{b}\); the admissible
superposition is Eq.~\eqref{eq:ch3.boundary-quantized-spectrum}, not the continuous integral
\cite{harrington2001,hansen1988}.

\paragraph{Worked illustration (waveguide TE\(_{10}\) analogy).}
A rectangular guide quantizes transverse wavenumbers to discrete \(k_{x},k_{y}\) \cite{jackson1999,harrington2001}.
Subsection~\ref{sec:ch372} treats that as boundary quantization of the lateral spectrum
Eq.~\eqref{eq:ch2.radiative-lateral-spectrum-def}: continuous Fourier parameters become countable modes
\cite{stratton1941,hansen1988}.

\paragraph{Problem (count quantized modes below cutoff).}
\textbf{Problem.} A cavity supports only modes with \(k_{0,\mathrm{eff}}<k_{0}\). How many \(q\) labels are radiatively
active? \cite{stratton1941,jackson1999,harrington2001}.

\textbf{Step 1.} Solve boundary eigenvalue problem for \(k_{0,\mathrm{eff},q}\) \cite{jackson1999}.

\textbf{Step 2.} Keep \(q\) with \(k_{0,\mathrm{eff},q}\le k_{0}\) on outgoing sheet \cite{stratton1941}.

\textbf{Step 3.} Sum watts only over active \(q\) \cite{poynting1884}.

\textbf{Physical meaning.} Quantized modes can be evanescent or radiative; only outgoing ones export \cite{harrington2001}.

\begin{table}[t]
\centering
\caption{Subsection~\ref{sec:ch372} quantization taxonomy.}
\label{tab:ch3.372-quantization-taxonomy}
\small
\begin{tabular}{@{}p{0.26\linewidth}p{0.34\linewidth}p{0.32\linewidth}@{}}
\toprule
Geometry & Spectrum & Equation \\
\midrule
Open exterior & continuous \(S^{2}\) & Eq.~\eqref{eq:ch3.unbounded-angular-spectrum-def} \\
PEC half-space & mixed & Eq.~\eqref{eq:ch3.image-mode-admixture} \\
Cavity & discrete \(q\) & Eq.~\eqref{eq:ch3.boundary-quantized-spectrum} \\
\bottomrule
\end{tabular}
\end{table}

\paragraph{Mixing before quantization.}
Eq.~\eqref{eq:ch3.boundary-scattering-mixing-matrix} acts on open-region labels before discrete \(q\) are identified
\cite{harrington2001,stratton1941}. Subsection~\ref{sec:ch372} order of operations: (i) impose \(\Gamma_{b}\), (ii) solve
for quantized \(\mathbf{e}_{q}\), (iii) normalize with Section~\ref{sec:ch35} meters \cite{hansen1988,harrington2001}.

\paragraph{Worked illustration (half-space not a cavity).}
A PEC plane quantizes parity but does not produce Eq.~\eqref{eq:ch2.cavity-discrete-radiation-spectrum} radial indices
\cite{jackson1999,stratton1941}. Subsection~\ref{sec:ch372} classifies half-space problems as \emph{mixed} continua, not
fully discrete cavities \cite{harrington2001,hansen1988}.

\paragraph{Problem (map open \(\ell\) to cavity \(q\)).}
\textbf{Problem.} Explain why a single open-region \(\ell\) label does not map one-to-one to cavity \(q\)
\cite{stratton1941,jackson1999,harrington2001}.

\textbf{Step 1.} Cite Eq.~\eqref{eq:ch3.cavity-angular-index} \cite{jackson1999}.

\textbf{Step 2.} Note radial \(n_{r}\) dependence \cite{stratton1941}.

\textbf{Physical meaning.} Quantization is three-dimensional, not purely angular \cite{harrington2001}.

\paragraph{Watt sums on discrete \(q\).}
After normalization, \(\mathcal{P}_{\mathrm{rad}}=\sum_{q}|a_{q}|^{2}\) when \(\{\mathbf{e}_{q}\}\) is orthonormal on
\(\mathcal{F}_{\mathrm{rad}}\) under \(\Gamma_{b}\) \cite{poynting1884,harrington2001}. Mixing from
Subsection~\ref{sec:ch363} may require non-diagonal Gram matrices before diagonalization \cite{hansen1988,stratton1941}.

\begin{figure}[t]
\centering
\ChOneFigBlock{%
\begin{tikzpicture}[ch1fig, node distance=7mm and 9mm]
  \node[ch1box] (open) {open VSH};
  \node[ch1box, right=10mm of open] (mix) {$\mathbf{S}_{b}$};
  \node[ch1box, right=10mm of mix] (eig) {eigenmodes $\mathbf{e}_{q}$};
  \draw[ch1flow] (open) -- (mix) -- (eig);
\end{tikzpicture}%
}
\caption{Subsection~\ref{sec:ch372} diagonalization path: boundary mixing then quantized eigenmode extraction.}
\label{fig:ch3.372-diagonalization}
\end{figure}

\paragraph{Accessibility under quantization.}
\(\mathcal{L}_{\mathrm{acc}}\) from Eq.~\eqref{eq:ch3.accessible-VSH-label-set} must be reindexed in \(q\) after
quantization \cite{harrington2001,stratton1941}. Forbidden \(q\) modes are geometrically annihilated like continuous
directions in \(\mathcal{K}_{\mathrm{sup}}\) \cite{stratton1941,harrington2001}.

\paragraph{Worked illustration (aperture into cavity).}
A small aperture couples exterior continuous spectra to discrete cavity \(q\) \cite{harrington2001,stratton1941,jackson1999}.
Subsection~\ref{sec:ch372} models the aperture as a coupling matrix between Eq.~\eqref{eq:ch3.unbounded-angular-spectrum-def}
and Eq.~\eqref{eq:ch3.boundary-quantized-spectrum} \cite{hansen1988,harrington2001}.

\paragraph{Problem (avoid open-region DOS in cavity).}
\textbf{Problem.} Show why Eq.~\eqref{eq:ch3.cumulative-disc-DOS} overcounts modes inside a closed cavity
\cite{harrington2001,stratton1941,jackson1999}.

\textbf{Step 1.} Count \(q\) from Eq.~\eqref{eq:ch3.boundary-quantized-spectrum} \cite{jackson1999}.

\textbf{Step 2.} Compare to open \(N_{\mathrm{disc}}\) \cite{hansen1988}.

\textbf{Physical meaning.} DOS is geometry-specific after quantization \cite{harrington2001}.

\paragraph{Quantized Parseval law (HOME).}
For orthonormal \(\{\mathbf{e}_{q}\}\) under \(\Gamma_{b}\),
\begin{equation}
\label{eq:ch3.quantized-parseval-law}
\mathcal{P}_{\mathrm{rad}}
=
\sum_{q\in\mathcal{Q}_{\mathrm{rad}}}
|a_{q}|^{2},
\qquad
\mathcal{Q}_{\mathrm{rad}}=\{q:\mathbf{e}_{q}\in\mathcal{F}_{\mathrm{rad}}\},
\end{equation}
importing Section~\ref{sec:ch35} normalization on each \(q\) \cite{poynting1884,harrington2001,stratton1941}. Equation
\eqref{eq:ch3.quantized-parseval-law} is the watt closure for discrete boundary spectra \cite{hansen1988,harrington2001}.

\paragraph{Worked illustration (radial quantum jump).}
Increasing cavity height shifts \(n_{r}\) ladders in Eq.~\eqref{eq:ch3.cavity-angular-index} without changing open-region
\(\ell\) counting \cite{jackson1999,stratton1941}. Quantized spectra therefore require geometry metadata in every mode table
\cite{harrington2001,hansen1988}.

\paragraph{Problem (couple exterior \(a\) to cavity \(a_{q}\)).}
\textbf{Problem.} An aperture coupling matrix \(C_{q}(\widehat{\mathbf{k}})\) maps exterior spectrum to cavity coefficients. Write
\(a_{q}=\int C_{q}(\widehat{\mathbf{k}})a(\widehat{\mathbf{k}})\,d\mu\) schematically \cite{harrington2001,stratton1941,jackson1999}.

\textbf{Step 1.} Identify aperture as boundary operator \cite{stratton1941}.

\textbf{Step 2.} Apply Eq.~\eqref{eq:ch3.boundary-scattering-mixing-matrix} \cite{hansen1988}.

\textbf{Physical meaning.} Quantization is a coupling map from continua, not a rename \cite{harrington2001}.

\begin{figure}[t]
\centering
\ChOneFigBlock{%
\begin{tikzpicture}[ch1fig, node distance=7mm and 9mm]
  \node[ch1box] (ext) {$a(\widehat{\mathbf{k}})$};
  \node[ch1box, right=10mm of ext] (ap) {aperture};
  \node[ch1box, right=10mm of ap] (cav) {$a_{q}$};
  \node[ch1box, below=10mm of ap] (par) {Eq.~\eqref{eq:ch3.quantized-parseval-law}};
  \draw[ch1flow] (ext) -- (ap) -- (cav);
  \draw[ch1flow] (par) -- (cav);
\end{tikzpicture}%
}
\caption{Subsection~\ref{sec:ch372} aperture coupling: exterior angular spectrum feeds quantized cavity modes.}
\label{fig:ch3.372-aperture-coupling}
\end{figure}

\begin{table}[t]
\centering
\caption{Subsection~\ref{sec:ch372} quantization metadata.}
\label{tab:ch3.372-quantization-metadata}
\small
\begin{tabular}{@{}p{0.30\linewidth}p{0.58\linewidth}@{}}
\toprule
Field & Content \\
\midrule
Boundary class & \(\Gamma_{b}\), PEC/Radome/cavity \\
Index map & Eq.~\eqref{eq:ch3.cavity-angular-index} \\
Active set & \(\mathcal{Q}_{\mathrm{rad}}\) in Eq.~\eqref{eq:ch3.quantized-parseval-law} \\
Mixing & Eq.~\eqref{eq:ch3.boundary-scattering-mixing-matrix} \\
\bottomrule
\end{tabular}
\end{table}

\paragraph{Standing-wave versus propagating quanta.}
Quantized modes with \(k_{0,\mathrm{eff},q}>k_{0}\) remain evanescent and do not contribute to
Eq.~\eqref{eq:ch3.quantized-parseval-law} export sums \cite{stratton1941,jackson1999}. Subsection~\ref{sec:ch372} separates
stored cavity energy from radiative \(q\) labels \cite{harrington2001,hansen1988}.

\paragraph{Worked illustration (openings in shielded enclosures).}
A slot in a cavity selects a subset \(\mathcal{Q}_{\mathrm{slot}}\subset\mathcal{Q}\) coupled to the exterior
\cite{harrington2001,stratton1941}. Reported radiation uses \(\mathcal{Q}_{\mathrm{slot}}\), not the full cavity book
Eq.~\eqref{eq:ch2.cavity-discrete-radiation-spectrum} \cite{jackson1999,hansen1988}.

\paragraph{Problem (reindex \(\mathcal{L}_{\mathrm{acc}}\) after quantization).}
\textbf{Problem.} Map open-region \(\mathcal{L}_{\mathrm{acc}}\) to \(q\) indices when a ground plane is added
\cite{harrington2001,stratton1941,jackson1999}.

\textbf{Step 1.} Diagonalize \(\mathbf{S}_{b}\) \cite{hansen1988}.

\textbf{Step 2.} Intersect with \(\mathcal{K}_{\mathrm{acc}}\) \cite{stratton1941}.

\textbf{Physical meaning.} Accessibility is reindexed, not preserved label-for-label \cite{harrington2001}.

\paragraph{PEC half-space reclassification.}
Image law Eq.~\eqref{eq:ch3.image-mode-admixture} produces mixed continua, not Eq.~\eqref{eq:ch3.boundary-quantized-spectrum}
discrete ladders \cite{jackson1999,stratton1941}. Subsection~\ref{sec:ch372} labels half-space problems as
\emph{mixed}, reserving discrete \(q\) for fully enclosed cavities \cite{harrington2001,hansen1988}.

\paragraph{Worked illustration (frequency sweep in cavity).}
As \(\omega\) increases, additional \(q\) cross the radiative threshold into \(\mathcal{Q}_{\mathrm{rad}}\)
\cite{jackson1999,stratton1941}. Quantized spectra are frequency lists, not static integer cuts \cite{harrington2001,hansen1988}.

\paragraph{Problem (compare \(\mathcal{Q}\) and \(N_{\mathrm{disc}}\)).}
\textbf{Problem.} At the same \(\ell_{\max}\), show \(|\mathcal{Q}_{\mathrm{rad}}|\) for a cavity can be less than
\(N_{\mathrm{disc}}(\ell_{\max})\) \cite{harrington2001,jackson1999,stratton1941}.

\textbf{Step 1.} Count radiative \(q\) only \cite{jackson1999}.

\textbf{Step 2.} Compare to Eq.~\eqref{eq:ch3.cumulative-disc-DOS} \cite{hansen1988}.

\textbf{Physical meaning.} Boundaries reduce active modes below open DOS \cite{harrington2001}.

\begin{table}[t]
\centering
\caption{Subsection~\ref{sec:ch372} radiative versus stored \(q\) labels.}
\label{tab:ch3.372-radiative-stored}
\small
\begin{tabular}{@{}p{0.22\linewidth}p{0.36\linewidth}p{0.34\linewidth}@{}}
\toprule
Class & Criterion & Power law \\
\midrule
\(\mathcal{Q}_{\mathrm{rad}}\) & outgoing & Eq.~\eqref{eq:ch3.quantized-parseval-law} \\
Stored & evanescent & not in \(\mathcal{P}_{\mathrm{rad}}\) \\
\bottomrule
\end{tabular}
\end{table}

\paragraph{Quantization workflow (laboratory).}
Subsection~\ref{sec:ch372} recommends the ordered workflow: (i) measure exterior \(a(\widehat{\mathbf{k}})\) if coupling exists;
(ii) identify \(\mathbf{S}_{b}\); (iii) solve for \(\mathbf{e}_{q}\); (iv) normalize each \(q\); (v) report
Eq.~\eqref{eq:ch3.quantized-parseval-law} \cite{harrington2001,stratton1941,jackson1999}. Skipping (ii) misattributes mixing to
source spectra \cite{hansen1988,harrington2001}.

\paragraph{Worked illustration (resonance hop).}
When a cavity mode crosses from below-cutoff to radiative as \(\omega\) sweeps, a new \(q\) enters \(\mathcal{Q}_{\mathrm{rad}}\)
discontinuously in the ideal lossless model \cite{jackson1999,stratton1941}. Measured watts can jump even when source drive is smooth
because a stored mode becomes exporting \cite{harrington2001,hansen1988}.

\paragraph{Problem (estimate stored fraction).}
\textbf{Problem.} Total cavity energy partitions into \(\sum_{q\in\mathcal{Q}_{\mathrm{rad}}}|a_{q}|^{2}\) plus stored terms. If
stored sum is \(0.3\,\mathrm{W}\) and radiative sum is \(0.7\,\mathrm{W}\), report export fraction \cite{poynting1884,harrington2001,stratton1941}.

\textbf{Step 1.} Sum radiative watts \cite{poynting1884}.

\textbf{Step 2.} Divide by total \(1\,\mathrm{W}\) \cite{harrington2001}.

\textbf{Physical meaning.} Quantized cavities can store non-exporting \(q\) \cite{stratton1941,jackson1999}.

\begin{figure}[t]
\centering
\ChOneFigBlock{%
\begin{tikzpicture}[ch1fig, node distance=7mm and 9mm]
  \node[ch1box] (omega) {$\omega$ sweep};
  \node[ch1box, right=10mm of omega] (hop) {$q$ hop};
  \node[ch1box, right=10mm of hop] (rad) {$\mathcal{Q}_{\mathrm{rad}}$};
  \draw[ch1flow] (omega) -- (hop) -- (rad);
\end{tikzpicture}%
}
\caption{Subsection~\ref{sec:ch372} resonance hop: modes enter \(\mathcal{Q}_{\mathrm{rad}}\) as frequency crosses cutoff.}
\label{fig:ch3.372-resonance-hop}
\end{figure}

\paragraph{Peer-review metadata for quantized campaigns.}
Publish \(\Gamma_{b}\) class, \(\mathcal{Q}_{\mathrm{rad}}\) list, Eq.~\eqref{eq:ch3.cavity-angular-index} map, and aperture coupling matrix
alongside \(a_{q}\) \cite{harrington2001,stratton1941,jackson1999}. Without that tuple, open-region spectra cannot be distinguished
from cavity spectra in post-processing \cite{hansen1988,harrington2001}.

\paragraph{Boundary constraint vocabulary (expanded).}
\emph{Tangential continuity} forces \(\hat{\mathbf{n}}\times\E\) and \(\hat{\mathbf{n}}\times\H\) to match partner fields across
\(\Gamma_{b}\) \cite{stratton1941,jackson1999}. \emph{Normal conductor conditions} set \(\hat{\mathbf{n}}\times\E=\mathbf{0}\) on PEC
surfaces \cite{stratton1941}. Subsection~\ref{sec:ch372} uses these constraints as the physical source of Eq.~\eqref{eq:ch3.boundary-quantized-spectrum}:
only eigenmodes satisfying the constraints simultaneously are admissible \(q\) labels \cite{harrington2001,hansen1988}.

\paragraph{Worked illustration (dielectric slab modes).}
A dielectric slab quantizes lateral wavenumbers while leaving exterior regions continuous \cite{jackson1999,harrington2001}. The
slab therefore contributes discrete \(q\) to Eq.~\eqref{eq:ch3.boundary-quantized-spectrum} while the exterior retains
Eq.~\eqref{eq:ch3.unbounded-angular-spectrum-def} \cite{stratton1941,hansen1988}.

\paragraph{Problem (classify a shielded box).}
\textbf{Problem.} A metal box with a coax feed is (i) cavity, (ii) half-space, or (iii) hybrid? \cite{harrington2001,stratton1941,jackson1999}.

\textbf{Step 1.} Identify closed \(\Gamma_{b}\) \cite{stratton1941}.

\textbf{Step 2.} Use Eq.~\eqref{eq:ch3.boundary-quantized-spectrum} inside \cite{jackson1999}.

\textbf{Step 3.} Model feed aperture as coupling to exterior \(a(\widehat{\mathbf{k}})\) \cite{harrington2001}.

\textbf{Physical meaning.} Real hardware is usually hybrid \cite{hansen1988,harrington2001}.

\begin{figure}[t]
\centering
\ChOneFigBlock{%
\begin{tikzpicture}[ch1fig, node distance=7mm and 9mm]
  \node[ch1box] (cav) {cavity $q$};
  \node[ch1box, right=10mm of cav] (ap) {aperture};
  \node[ch1box, right=10mm of ap] (ext) {exterior $a(\widehat{\mathbf{k}})$};
  \draw[ch1flow] (cav) -- (ap) -- (ext);
\end{tikzpicture}%
}
\caption{Subsection~\ref{sec:ch372} hybrid hardware: quantized interior coupled to continuous exterior spectra.}
\label{fig:ch3.372-hybrid-hardware}
\end{figure}

\begin{table}[t]
\centering
\caption{Subsection~\ref{sec:ch372} geometry versus spectrum type.}
\label{tab:ch3.372-geometry-spectrum}
\small
\begin{tabular}{@{}p{0.28\linewidth}p{0.62\linewidth}@{}}
\toprule
Geometry & Spectrum model \\
\midrule
Open exterior & Eq.~\eqref{eq:ch3.unbounded-angular-spectrum-def} \\
PEC half-space & mixed Eq.~\eqref{eq:ch3.image-mode-admixture} \\
Closed cavity & Eq.~\eqref{eq:ch3.boundary-quantized-spectrum} \\
Fed cavity & hybrid coupling \\
\bottomrule
\end{tabular}
\end{table}

\paragraph{Boundary eigenvalue problem (HOME).}
Quantized labels \(q\) are not chosen by the analyst; they are eigenvalues of a boundary-constrained Helmholtz problem
\cite{stratton1941,jackson1999,harrington2001}. Let \(\mathcal{L}_{\Gamma_{b}}\) denote the self-adjoint operator obtained
by restricting Maxwell's separated angular--radial system to fields satisfying \(\Gamma_{b}\). Subsection~\ref{sec:ch372}
fixes the spectral statement
\begin{equation}
\label{eq:ch3.boundary-eigenvalue-quantization}
\mathcal{L}_{\Gamma_{b}}\,\mathbf{e}_{q}
=
k_{0,\mathrm{eff},q}^{2}\,\mathbf{e}_{q},
\qquad
q\in\mathcal{Q},
\qquad
\mathbf{e}_{q}\perp\mathbf{e}_{q'}\ \text{under }\Gamma_{b},
\end{equation}
with \(k_{0,\mathrm{eff},q}\) the effective wavenumber set by cavity geometry \cite{jackson1999,stratton1941}. Equation
\eqref{eq:ch3.boundary-eigenvalue-quantization} is the operator backbone of Eq.~\eqref{eq:ch3.boundary-quantized-spectrum}:
discrete \(q\) are eigenmode indices, and orthogonality is imposed by the boundary inner product, not by open-region VSH orthogonality
alone \cite{hansen1988,harrington2001,kato1995}. Physically, a conducting enclosure selects standing patterns whose tangential fields
vanish on \(\Gamma_{b}\); only those patterns become admissible normal modes, and their eigenfrequencies replace the continuous
direction parameter of Eq.~\eqref{eq:ch3.unbounded-angular-spectrum-def} \cite{stratton1941,poynting1884}.

\paragraph{Cutoff and propagating quanta (expanded).}
\emph{Cutoff} means \(k_{0,\mathrm{eff},q}<k_{0}\) so the mode is evanescent at the operating frequency; \emph{propagating quanta}
means \(k_{0,\mathrm{eff},q}\le k_{0}\) on the outgoing sheet and the mode belongs to \(\mathcal{Q}_{\mathrm{rad}}\) in
Eq.~\eqref{eq:ch3.quantized-parseval-law} \cite{jackson1999,harrington2001,stratton1941}. Subsection~\ref{sec:ch372} treats cutoff as a
\emph{spectral threshold} in Eq.~\eqref{eq:ch3.boundary-eigenvalue-quantization}, not as a numerical truncation order: a mode below cutoff
can store reactive energy without contributing export watts \cite{harrington2001,hansen1988}. Raising \(\omega\) slides the operating point
along the eigenvalue ladder and can activate new \(q\) labels discontinuously in the lossless idealization of
Figure~\ref{fig:ch3.372-resonance-hop} \cite{jackson1999,stratton1941}.

\paragraph{Worked illustration (rectangular cavity index triple).}
In a rectangular PEC cavity, transverse indices \((n_{x},n_{y})\) and axial index \(n_{z}\) produce
\(k_{0,\mathrm{eff},q}^{2}=(n_{x}\pi/a)^{2}+(n_{y}\pi/b)^{2}+(n_{z}\pi/d)^{2}\) \cite{jackson1999,harrington2001}. Mapping to
Eq.~\eqref{eq:ch3.cavity-angular-index} requires a spherical-harmonic analysis of the rectangular eigenfields; the point for
Subsection~\ref{sec:ch372} is that three integers \((n_{x},n_{y},n_{z})\) quantize the spectrum before any open-region \(\ell\) label is
meaningful inside the box \cite{stratton1941,hansen1988}. Exterior reports that ignore \((n_{x},n_{y},n_{z})\) confuse feed coupling with
angular-spectrum physics \cite{harrington2001,stratton1941}.

\paragraph{Problem (eigenvalue threshold for a new \(q\)).}
\textbf{Problem.} A cubic cavity of side \(d\) has \(k_{0,\mathrm{eff},q}=(\pi/d)\sqrt{n_{x}^{2}+n_{y}^{2}+n_{z}^{2}}\). For
\((n_{x},n_{y},n_{z})=(1,1,1)\), find the lowest \(\omega\) at which this \(q\) enters \(\mathcal{Q}_{\mathrm{rad}}\)
\cite{jackson1999,stratton1941,harrington2001}.

\textbf{Step 1.} Set \(k_{0,\mathrm{eff},q}=k_{0}=\omega\sqrt{\varepsilon\mu}\) \cite{jackson1999}.

\textbf{Step 2.} Solve \(\omega_{\min}=(\pi\sqrt{3}/d)\,c\) in free-space filling \cite{stratton1941}.

\textbf{Step 3.} Below \(\omega_{\min}\), the triplet stores but does not export on the outgoing sheet \cite{harrington2001}.

\textbf{Physical meaning.} Quantization thresholds are geometric eigenvalues, not fit tolerances \cite{hansen1988,harrington2001}.

\begin{figure}[t]
\centering
\ChOneFigBlock{%
\begin{tikzpicture}[ch1fig, node distance=7mm and 9mm]
  \node[ch1box] (ladder) {Eq.~\eqref{eq:ch3.boundary-eigenvalue-quantization}};
  \node[ch1box, right=10mm of ladder] (cut) {cutoff $k_{0,\mathrm{eff},q}=k_{0}$};
  \node[ch1box, right=10mm of cut] (q) {$\mathcal{Q}_{\mathrm{rad}}$};
  \node[ch1box, below=10mm of cut] (par) {Eq.~\eqref{eq:ch3.quantized-parseval-law}};
  \draw[ch1flow] (ladder) -- (cut) -- (q);
  \draw[ch1flow] (par) -- (q);
\end{tikzpicture}%
}
\caption{Subsection~\ref{sec:ch372} eigenvalue ladder: boundary eigenmodes cross into \(\mathcal{Q}_{\mathrm{rad}}\) only above cutoff.}
\label{fig:ch3.372-eigenvalue-ladder}
\end{figure}

\begin{table}[t]
\centering
\caption{Subsection~\ref{sec:ch372} boundary condition classes and quantization outcome.}
\label{tab:ch3.372-bc-quantization}
\small
\begin{tabular}{@{}p{0.26\linewidth}p{0.34\linewidth}p{0.32\linewidth}@{}}
\toprule
\(\Gamma_{b}\) class & Operator effect & Spectrum \\
\midrule
PEC enclosure & discrete eigenvalues Eq.~\eqref{eq:ch3.boundary-eigenvalue-quantization} & Eq.~\eqref{eq:ch3.boundary-quantized-spectrum} \\
PEC half-space & image parity & Eq.~\eqref{eq:ch3.image-mode-admixture} \\
Open exterior & no compactification & Eq.~\eqref{eq:ch3.unbounded-angular-spectrum-def} \\
\bottomrule
\end{tabular}
\end{table}

\paragraph{Degeneracy lifting by \(\Gamma_{b}\).}
Open-region degeneracy in \((\ell,m)\) is lifted when \(\Gamma_{b}\) breaks \(\mathrm{SO}(3)\): distinct \(q\) labels can share
similar \(k_{0,\mathrm{eff},q}\) but differ in boundary parity or feed coupling \cite{stratton1941,harrington2001,jackson1999}.
Subsection~\ref{sec:ch372} therefore pairs Eq.~\eqref{eq:ch3.boundary-eigenvalue-quantization} with mixing matrix
Eq.~\eqref{eq:ch3.boundary-scattering-mixing-matrix} before reporting sorted \(q\) lists: eigenvalues alone do not specify the
radiating content until outgoing projection and aperture coupling are applied \cite{hansen1988,harrington2001}.

\paragraph{Worked illustration (loss broadens eigensteps).}
Finite conductivity and dielectric loss broaden eigensteps in Eq.~\eqref{eq:ch3.boundary-eigenvalue-quantization} into resonances with
finite \(Q\) \cite{jackson1999,harrington2001,stratton1941}. Measured \(\mathcal{Q}_{\mathrm{rad}}\) membership then becomes frequency
fuzzy near threshold, even though the ideal lossless model in Eq.~\eqref{eq:ch3.quantized-parseval-law} uses sharp hops
\cite{harrington2001,hansen1988}. Laboratory quantization reports should attach \(Q\) or linewidth metadata beside each active \(q\)
\cite{stratton1941,jackson1999}.

\paragraph{Problem (diagonalize before watt sum).}
\textbf{Problem.} Boundary mixing yields non-orthogonal raw modes. Outline the diagonalization needed before
Eq.~\eqref{eq:ch3.quantized-parseval-law} \cite{harrington2001,stratton1941,hansen1988}.

\textbf{Step 1.} Form Gram matrix from Eq.~\eqref{eq:ch3.boundary-scattering-mixing-matrix} \cite{hansen1988}.

\textbf{Step 2.} Extract orthonormal \(\mathbf{e}_{q}\) satisfying Eq.~\eqref{eq:ch3.boundary-eigenvalue-quantization} \cite{kato1995}.

\textbf{Step 3.} Normalize each \(q\) with Section~\ref{sec:ch35} meters \cite{poynting1884,harrington2001}.

\textbf{Physical meaning.} Watt sums require orthonormal eigenmodes under \(\Gamma_{b}\), not raw image superpositions \cite{stratton1941}.

\paragraph{Quantization archive tuple (expanded).}
Peer archives for quantized campaigns should store \((\Gamma_{b},\mathcal{L}_{\Gamma_{b}},\mathcal{Q},\mathcal{Q}_{\mathrm{rad}})\),
the eigenvalue list \(\{k_{0,\mathrm{eff},q}\}\), aperture coupling if hybrid, and Eq.~\eqref{eq:ch3.cavity-angular-index} maps
\cite{harrington2001,stratton1941,jackson1999,hansen1988}. Without the eigenvalue list, post-processors cannot distinguish a stored mode
from a radiating one when \(\omega\) sits near cutoff \cite{harrington2001,jackson1999}.

Subsection~\ref{sec:ch372} has explained boundary-induced quantization through
Eqs.~\eqref{eq:ch3.boundary-quantized-spectrum}--\eqref{eq:ch3.cavity-angular-index}
\cite{stratton1941,jackson1999,harrington2001}. Subsection~\ref{sec:ch373} separates integer and non-integer angular
indices without informal beam-order collapse \cite{hansen1988}.


\subsection{Integer and non-integer angular indices}
\label{sec:ch373}
% TOC tag: [HOME]

\emph{Integer angular indices} are the quantum numbers \((\ell,m)\) aligned with \(\mathrm{SO}(3)\) irreducible
representations on unbounded domains \cite{jackson1999,hansen1988,stratton1941}.
\emph{Non-integer indices} arise when symmetry is reduced, conical geometries appear, or continuation parameters label
non-physical sheets---they are not informal ``fractional beam orders'' \cite{harrington2001,stratton1941}. Subsection~\ref{sec:ch373}
is the HOME index discipline subsection for Volume~I \cite{hansen1988,harrington2001}.

\paragraph{Standing imports.}
Multiplet multiplicity Eq.~\eqref{eq:ch3.multiplet-multiplicity-def}, Wigner mixing
Eq.~\eqref{eq:ch3.eigenfield-Wigner-mixing}, and cumulative DOS Eq.~\eqref{eq:ch3.cumulative-disc-DOS} enter by
reference \cite{jackson1999,hansen1988}. Subsection~\ref{sec:ch373} does not re-derive addition of angular momentum
\cite{jackson1999,stratton1941}.

\paragraph{Assumptions.}
Work in the normalized VSH gauge unless continuation is explicitly declared \cite{hansen1988,harrington2001}. Integer
\(\ell\in\mathbb{N}_{0}\) on open exteriors; non-integer labels appear only under declared symmetry breaking or
continuation paths \cite{stratton1941,harrington2001}.

\paragraph{Integer multipole index (HOME).}
On unbounded domains, the integer label is
\begin{equation}
\label{eq:ch3.integer-ell-def}
\ell\in\{0,1,2,\ldots\},
\qquad
m\in\{-\ell,\ldots,\ell\},
\qquad
\sigma\in\{\mathbf{M},\mathbf{N}\},
\end{equation}
with physical content fixed by Eqs.~\eqref{eq:ch3.Mlm-def}--\eqref{eq:ch3.Nlm-def} and eigenvalue
Eq.~\eqref{eq:ch3.eigenfield-Jsq-eigenvalue} \cite{hansen1988,stratton1941,jackson1999}. Equation
\eqref{eq:ch3.integer-ell-def} is the Subsection~\ref{sec:ch373} integer contract: these are representation indices,
not adjustable fit orders \cite{harrington2001}.

\paragraph{Non-integer continuation (HOME).}
When symmetry is reduced, define continuation parameters \(\nu\) only through analytic continuation of separated
solutions with boundary matching:
\begin{equation}
\label{eq:ch3.noninteger-ell-continuation}
\mathbf{u}_{\nu m}^{(\sigma)}
=
\text{continuation of separated angular solutions with declared boundary sheet},
\qquad
\nu\notin\mathbb{N}_{0}\ \text{allowed only with boundary certificate},
\end{equation}
not through fitting heuristics \cite{stratton1941,harrington2001,jackson1999}. Equation
\eqref{eq:ch3.noninteger-ell-continuation} blocks popular misuse of ``fractional \(\ell\)'' in OAM discourse without
geometry \cite{allen1992,harrington2001}.

\paragraph{Beam-order warning (HOME).}
\begin{equation}
\label{eq:ch3.beam-order-vs-ell-warning}
\ell_{\mathrm{fit}}\ \text{(algorithmic truncation order)}
\;\neq\;
\ell\ \text{(representation index in Eq.~\eqref{eq:ch3.integer-ell-def})}.
\end{equation}
Equation~\eqref{eq:ch3.beam-order-vs-ell-warning} is the Subsection~\ref{sec:ch373} anti-confusion law: numerical beam
orders count fitted terms; \(\ell\) labels irreducible sectors certified by harmonics and operators
\cite{harrington2001,hansen1988,stratton1941}.

\begin{figure}[t]
\centering
\ChOneFigBlock{%
\begin{tikzpicture}[ch1fig, node distance=7mm and 9mm]
  \node[ch1box] (int) {integer $\ell$};
  \node[ch1box, right=10mm of int] (break) {symmetry break};
  \node[ch1box, right=10mm of break] (nu) {non-integer $\nu$};
  \node[ch1box, below=10mm of int] (warn) {Eq.~\eqref{eq:ch3.beam-order-vs-ell-warning}};
  \draw[ch1flow] (int) -- (break) -- (nu);
  \draw[ch1flow] (warn) -- (int);
\end{tikzpicture}%
}
\caption{Subsection~\ref{sec:ch373} index discipline: integer \(\ell\) on open domains; non-integer \(\nu\) only with
boundary continuation certificates.}
\label{fig:ch3.373-index-discipline}
\end{figure}

\paragraph{Physical meaning of integer \(\ell\).}
Integer \(\ell\) indexes the angular momentum content of radiative harmonics on \(S^{2}\) \cite{jackson1999,hansen1988}.
Measuring a pattern does not automatically reveal a single \(\ell\); superpositions carry multiple integers with
coefficients \(\widehat{c}_{\ell m}^{(\sigma)}\) \cite{harrington2001,stratton1941}. Subsection~\ref{sec:ch373} therefore
rejects single-number ``beam \(\ell\)'' summaries without expansion data \cite{harrington2001}.

\paragraph{Worked illustration (pure \(\ell=1\) versus mixed).}
\(|\widehat{c}_{10}^{(\mathbf{N})}|^{2}=1\,\mathrm{W}\) with all other coefficients zero is pure dipole
\cite{hansen1988,poynting1884}. Adding \(|\widehat{c}_{20}^{(\mathbf{N})}|^{2}=0.1\,\mathrm{W}\) makes the state
\(\ell=1\) plus \(\ell=2\) superposition---not a ``1.2-order'' beam \cite{harrington2001,stratton1941}.

\paragraph{Problem (decode a fit order list).}
\textbf{Problem.} An optimizer reports orders \((0,1,1,2)\). Map to Eq.~\eqref{eq:ch3.integer-ell-def} labels
\cite{harrington2001,hansen1988,stratton1941}.

\textbf{Step 1.} Identify each term as separate \((\ell,m,\sigma)\) \cite{hansen1988}.

\textbf{Step 2.} Reject averaged ``effective \(\ell\)'' \cite{harrington2001}.

\textbf{Physical meaning.} Fit orders are algorithmic; \(\ell\) is representation-theoretic \cite{stratton1941}.

\begin{table}[t]
\centering
\caption{Subsection~\ref{sec:ch373} integer versus non-integer index comparison.}
\label{tab:ch3.373-index-comparison}
\small
\begin{tabular}{@{}p{0.22\linewidth}p{0.36\linewidth}p{0.34\linewidth}@{}}
\toprule
Label & Domain & Certificate \\
\midrule
\(\ell\in\mathbb{N}_{0}\) & open exterior & Eq.~\eqref{eq:ch3.integer-ell-def} \\
\(\nu\notin\mathbb{N}_{0}\) & reduced symmetry & Eq.~\eqref{eq:ch3.noninteger-ell-continuation} \\
\(\ell_{\mathrm{fit}}\) & numerical & Eq.~\eqref{eq:ch3.beam-order-vs-ell-warning} \\
\bottomrule
\end{tabular}
\end{table}

\paragraph{Conical geometry pointer.}
Conical wave problems introduce non-integer \(\nu\) through separated angular equations with wedge boundary conditions
\cite{jackson1999,stratton1941}. Subsection~\ref{sec:ch373} records the label without constructing conical functions:
Volume~I treats them as symmetry-broken continuations Eq.~\eqref{eq:ch3.noninteger-ell-continuation}
\cite{hansen1988,harrington2001}.

\paragraph{Worked illustration (cylindrical versus spherical labels).}
Cylindrical modes use \(m\) and axial wavenumber; spherical uses \(\ell,m\) \cite{stratton1941,jackson1999}. Section~\ref{sec:ch38}
transforms between bases; Subsection~\ref{sec:ch373} warns against identifying cylindrical \(m\) with spherical \(\ell\)
\cite{hansen1988,harrington2001}.

\paragraph{Problem (non-integer claim audit).}
\textbf{Problem.} A paper claims ``\(\ell=3/2\) OAM beam'' in free space. Audit using Eq.~\eqref{eq:ch3.noninteger-ell-continuation}
\cite{allen1992,harrington2001,stratton1941}.

\textbf{Step 1.} Demand boundary certificate \cite{stratton1941}.

\textbf{Step 2.} If none, re-express as integer superposition \cite{hansen1988}.

\textbf{Physical meaning.} Free-space \(\ell\) is integer in the VSH catalog \cite{jackson1999,harrington2001}.

\begin{figure}[t]
\centering
\ChOneFigBlock{%
\begin{tikzpicture}[ch1fig, node distance=7mm and 9mm]
  \node[ch1box] (catalog) {Eq.~\eqref{eq:ch3.integer-ell-def}};
  \node[ch1box, right=10mm of catalog] (fit) {$\ell_{\mathrm{fit}}$};
  \node[ch1box, right=10mm of fit] (reject) {Eq.~\eqref{eq:ch3.beam-order-vs-ell-warning}};
  \draw[ch1flow] (catalog) -- (reject);
  \draw[ch1flow] (fit) -- (reject);
\end{tikzpicture}%
}
\caption{Subsection~\ref{sec:ch373} rejects identifying numerical fit orders with representation indices \(\ell\).}
\label{fig:ch3.373-fit-vs-rep}
\end{figure}

\paragraph{DOS and integer ladders.}
Eq.~\eqref{eq:ch3.cumulative-disc-DOS} counts integer slots only \cite{hansen1988,harrington2001}. Non-integer
continuations do not increment DOS until a discrete boundary spectrum is certified \cite{stratton1941,jackson1999}.

\paragraph{Worked illustration (superposition uniqueness).}
Two distinct integer superpositions can produce similar patterns but different \(\widehat{c}_{\ell m}^{(\sigma)}\) vectors;
inversion must recover the vector, not a single \(\ell\) slogan \cite{devaney2004,harrington2001,stratton1941}.

\paragraph{Problem (watt partition across integers).}
\textbf{Problem.} Given \(|\widehat{c}_{10}^{(\mathbf{N})}|^{2}=0.7\,\mathrm{W}\) and
\(|\widehat{c}_{20}^{(\mathbf{N})}|^{2}=0.3\,\mathrm{W}\), report integer power partition \cite{poynting1884,harrington2001}.

\textbf{Step 1.} Read \(\ell=1\) sector \(0.7\,\mathrm{W}\) \cite{hansen1988}.

\textbf{Step 2.} Read \(\ell=2\) sector \(0.3\,\mathrm{W}\) \cite{poynting1884}.

\textbf{Physical meaning.} Integer partitions are orthogonal watt sums \cite{stratton1941,harrington2001}.

\paragraph{Integer ladder degeneracy (HOME).}
At each integer \(\ell\), degeneracy \(2\ell+1\) in \(m\) and dual \(\sigma\) sectors follows
Eq.~\eqref{eq:ch3.J-multiplet-degeneracy-count} \cite{jackson1999,hansen1988}. Subsection~\ref{sec:ch373} treats that
degeneracy as representation multiplicity, not as independent fit parameters unless Wigner mixing
Eq.~\eqref{eq:ch3.eigenfield-Wigner-mixing} is resolved \cite{harrington2001,stratton1941}.

\paragraph{Worked illustration (rotate integer labels).}
Rotating the laboratory mixes \(m\) within fixed \(\ell\) via Eq.~\eqref{eq:ch3.eigenfield-Wigner-mixing} without changing
integer \(\ell\) \cite{jackson1999,hansen1988}. Non-integer claims cannot be generated by rotation alone on open domains
\cite{harrington2001,stratton1941}.

\paragraph{Problem (effective \(\ell\) marketing audit).}
\textbf{Problem.} A vendor reports ``effective \(\ell=2.3\).'' Re-express using Eqs.~\eqref{eq:ch3.beam-order-vs-ell-warning}
and \eqref{eq:ch3.integer-ell-def} \cite{harrington2001,hansen1988,stratton1941}.

\textbf{Step 1.} Request full coefficient vector \cite{devaney2004}.

\textbf{Step 2.} Replace scalar slogan with integer partition \cite{harrington2001}.

\textbf{Physical meaning.} Effective indices are not representation labels \cite{stratton1941}.

\begin{figure}[t]
\centering
\ChOneFigBlock{%
\begin{tikzpicture}[ch1fig, node distance=7mm and 9mm]
  \node[ch1box] (ell) {integer $\ell$};
  \node[ch1box, right=10mm of ell] (deg) {Eq.~\eqref{eq:ch3.J-multiplet-degeneracy-count}};
  \node[ch1box, right=10mm of deg] (w) {Eq.~\eqref{eq:ch3.eigenfield-Wigner-mixing}};
  \draw[ch1flow] (ell) -- (deg) -- (w);
\end{tikzpicture}%
}
\caption{Subsection~\ref{sec:ch373} integer ladder: multiplet degeneracy and rotation mixing within fixed \(\ell\).}
\label{fig:ch3.373-integer-ladder}
\end{figure}

\begin{table}[t]
\centering
\caption{Subsection~\ref{sec:ch373} index reporting contract.}
\label{tab:ch3.373-reporting-contract}
\small
\begin{tabular}{@{}p{0.34\linewidth}p{0.56\linewidth}@{}}
\toprule
Report field & Required content \\
\midrule
Integer expansion & full \(\widehat{c}_{\ell m}^{(\sigma)}\) vector \\
Non-integer claim & Eq.~\eqref{eq:ch3.noninteger-ell-continuation} certificate \\
Fit orders & mapped via Eq.~\eqref{eq:ch3.beam-order-vs-ell-warning} \\
\bottomrule
\end{tabular}
\end{table}

\paragraph{Continuation sheet discipline.}
Non-integer \(\nu\) in Eq.~\eqref{eq:ch3.noninteger-ell-continuation} must name the analytic sheet and boundary matching
data \cite{stratton1941,harrington2001}. Without that data, fits must remain in integer span
Eq.~\eqref{eq:ch3.integer-ell-def} \cite{hansen1988,jackson1999}.

\paragraph{Worked illustration (sparse versus dense integer spectra).}
Sparse \(\widehat{c}\) vectors with only \(\ell\le3\) active are still integer-indexed; density of active integers is not
non-integer \(\ell\) \cite{harrington2001,stratton1941}. Subsection~\ref{sec:ch373} separates sparsity from index continuity
\cite{hansen1988,harrington2001}.

\paragraph{Problem (integer count at \(\ell_{\max}=3\)).}
\textbf{Problem.} Count distinct integer labels in Eq.~\eqref{eq:ch3.integer-ell-def} with \(\ell\le3\), both \(\sigma\)
\cite{hansen1988,jackson1999,harrington2001}.

\textbf{Step 1.} Sum \(2(2\ell+1)\) for \(\ell=0,1,2,3\) \cite{jackson1999}.

\textbf{Step 2.} Obtain \(32\) channels \cite{hansen1988}.

\textbf{Physical meaning.} Integer catalogs grow quadratically in \(\ell_{\max}\) \cite{harrington2001}.

\paragraph{Algorithmic order versus physics (expanded).}
Nonlinear optimizers may adaptively increase \(\ell_{\mathrm{fit}}\) until a cost threshold is met \cite{harrington2001,devaney2004}.
Eq.~\eqref{eq:ch3.beam-order-vs-ell-warning} requires publishing the final integer coefficient vector even when the optimizer
used \(\ell_{\mathrm{fit}}\) internally \cite{stratton1941,hansen1988}. Subsection~\ref{sec:ch373} separates algorithmic stopping
rules from representation content \cite{harrington2001}.

\paragraph{Worked illustration (pencil beam not a single \(\ell\)).}
Narrow beams generally require many integer \(\ell\) sectors with structured \(\widehat{c}_{\ell m}^{(\sigma)}\) interference
\cite{hansen1988,jackson1999,harrington2001}. Calling such beams ``\(\ell=0\)'' because the peak is broadside misreads
Eq.~\eqref{eq:ch3.integer-ell-def} \cite{stratton1941,harrington2001}.

\paragraph{Problem (integer support from inversion).}
\textbf{Problem.} An inversion returns nonzero coefficients only for \(\ell\le4\). State the integer support set
\cite{devaney2004,harrington2001,hansen1988}.

\textbf{Step 1.} List all \((\ell,m,\sigma)\) with \(|\widehat{c}|>\epsilon\) \cite{harrington2001}.

\textbf{Step 2.} Do not collapse to one \(\ell\) \cite{stratton1941}.

\textbf{Physical meaning.} Support sets are finite integer lists, not fractional orders \cite{hansen1988}.

\begin{figure}[t]
\centering
\ChOneFigBlock{%
\begin{tikzpicture}[ch1fig, node distance=7mm and 9mm]
  \node[ch1box] (supp) {support set};
  \node[ch1box, right=10mm of supp] (int) {Eq.~\eqref{eq:ch3.integer-ell-def}};
  \node[ch1box, right=10mm of int] (w) {watt partition};
  \draw[ch1flow] (supp) -- (int) -- (w);
\end{tikzpicture}%
}
\caption{Subsection~\ref{sec:ch373} inversion output: declare integer support sets before any scalar beam-order summary.}
\label{fig:ch3.373-support-set}
\end{figure}

\paragraph{Magnetic versus electric integer tracks.}
Integer \(\ell\) on \(\mathbf{N}\)-type and \(\mathbf{M}\)-type sectors are independent ladders in
Eq.~\eqref{eq:ch3.integer-ell-def} \cite{belinfante1940,hansen1988}. Reporting ``\(\ell=2\)'' without \(\sigma\) is ambiguous
\cite{harrington2001,stratton1941}. Subsection~\ref{sec:ch373} requires \(\sigma\) in every integer label tuple
\cite{hansen1988,harrington2001}.

\paragraph{Integer projection operators (HOME sketch).}
Define sector projectors \(\mathcal{P}_{\ell\sigma}\) that extract the integer block at fixed \((\ell,\sigma)\) from any
expansion \cite{hansen1988,kato1995}. Watt power in that block is \(\|\mathcal{P}_{\ell\sigma}\mathbf{v}_\omega\|_{\mathrm{rad}}^{2}\)
under normalized VSH \cite{poynting1884,harrington2001}. Subsection~\ref{sec:ch373} uses block powers as the integer
\emph{observable} associated with labels, not a single fitted scalar \cite{stratton1941,hansen1988}.

\paragraph{Worked illustration (dominant \(\ell\) sector).}
If \(\|\mathcal{P}_{2\mathbf{N}}\mathbf{v}\|^{2}=0.6\,\mathrm{W}\) and all other blocks sum to \(0.4\,\mathrm{W}\), the dominant
integer sector is \(\ell=2\), \(\mathbf{N}\)-type, but the state is still a multi-sector superposition \cite{harrington2001,poynting1884}.
Calling the beam ``pure \(\ell=2\)'' omits \(0.4\,\mathrm{W}\) elsewhere \cite{stratton1941,hansen1988}.

\paragraph{Problem (rank integer blocks by power).}
\textbf{Problem.} Given block powers \(\{0.5,0.3,0.15,0.05\}\,\mathrm{W}\) on four integer labels, rank sectors without merging labels
\cite{harrington2001,hansen1988,stratton1941}.

\textbf{Step 1.} Sort descending \cite{harrington2001}.

\textbf{Step 2.} Keep full label list in the report \cite{stratton1941}.

\textbf{Physical meaning.} Dominance is partial, not exclusive \cite{hansen1988,harrington2001}.

\begin{figure}[t]
\centering
\ChOneFigBlock{%
\begin{tikzpicture}[ch1fig, node distance=7mm and 9mm]
  \node[ch1box] (blocks) {integer blocks};
  \node[ch1box, right=10mm of blocks] (pow) {block watts};
  \node[ch1box, right=10mm of pow] (rank) {ranked report};
  \draw[ch1flow] (blocks) -- (pow) -- (rank);
\end{tikzpicture}%
}
\caption{Subsection~\ref{sec:ch373} integer block powers: rank sectors by watts without collapsing to one effective \(\ell\).}
\label{fig:ch3.373-block-ranking}
\end{figure}

\begin{table}[t]
\centering
\caption{Subsection~\ref{sec:ch373} example integer block watt ledger (\(\ell\le2\)).}
\label{tab:ch3.373-block-ledger}
\small
\begin{tabular}{@{}p{0.20\linewidth}p{0.22\linewidth}p{0.22\linewidth}p{0.26\linewidth}@{}}
\toprule
Label & block power & cumulative & note \\
\midrule
$(1,0,\mathbf{N})$ & $0.45\,\mathrm{W}$ & $0.45$ & dominant \\
$(2,0,\mathbf{N})$ & $0.30\,\mathrm{W}$ & $0.75$ & secondary \\
$(1,1,\mathbf{M})$ & $0.15\,\mathrm{W}$ & $0.90$ & cross-type \\
other & $0.10\,\mathrm{W}$ & $1.00$ & tail \\
\bottomrule
\end{tabular}
\end{table}

\paragraph{Peer-review rule for integer claims.}
Any public ``beam \(\ell\)'' statement must attach Table~\ref{tab:ch3.373-block-ledger} style watt partitions or the full
\(\widehat{c}_{\ell m}^{(\sigma)}\) vector \cite{harrington2001,devaney2004,stratton1941}. Subsection~\ref{sec:ch373} treats scalar
\(\ell\) slogans as non-compliant unless backed by integer block data \cite{hansen1988,harrington2001}.

\paragraph{Non-integer vocabulary (expanded).}
\emph{Non-integer angular indices} in Subsection~\ref{sec:ch373} mean continuation parameters \(\nu\) tied to reduced symmetry
equations, not curve-fit orders \cite{stratton1941,harrington2001}. \emph{Integer indices} mean \(\ell\in\mathbb{N}_{0}\) in the
\(\mathrm{SO}(3)\) catalog Eq.~\eqref{eq:ch3.integer-ell-def} \cite{jackson1999,hansen1988}. Collapsing the two vocabularies
produces the common error flagged by Eq.~\eqref{eq:ch3.beam-order-vs-ell-warning} \cite{harrington2001,stratton1941}.

\paragraph{Worked illustration (continuation sheet naming).}
A wedge guide may use \(\nu\) as separation constant; reports must print \(\nu\) and boundary angles, not a rounded ``\(\ell\)'' scalar
\cite{jackson1999,stratton1941,harrington2001}. Subsection~\ref{sec:ch373} treats missing sheet data as an incomplete index declaration
\cite{hansen1988,harrington2001}.

\paragraph{Problem (translate vendor table to integer vector).}
\textbf{Problem.} A vendor table lists ``mode orders'' \(\{0,1,1,2\}\). Convert to labeled pairs \((\ell,m,\sigma)\) with watt column
\cite{harrington2001,hansen1988,stratton1941}.

\textbf{Step 1.} Map each row to a distinct label \cite{hansen1988}.

\textbf{Step 2.} Attach measured watts per row \cite{poynting1884}.

\textbf{Step 3.} Reject merged effective \(\ell\) \cite{harrington2001}.

\textbf{Physical meaning.} Vendor tables are not physics labels until mapped \cite{stratton1941,devaney2004}.

Subsection~\ref{sec:ch373} has fixed integer and non-integer angular index discipline through
Eqs.~\eqref{eq:ch3.integer-ell-def}--\eqref{eq:ch3.beam-order-vs-ell-warning}
\cite{jackson1999,hansen1988,harrington2001}. Subsection~\ref{sec:ch374} connects phase singularities and topological
limits to angular spectra \cite{allen1992,stratton1941}.


\subsection{Phase singularities and topological limits}
\label{sec:ch374}
% TOC tag: [HOME][USE SS1.2.2]

A \emph{phase singularity} is a point (or curve) where electromagnetic phase is undefined while magnitude may remain
bounded \cite{allen1992,stratton1941,jackson1999}. A \emph{topological limit} is the constraint that links phase winding
around such singularities to transport and radiation diagnostics certified in Chapter~\ref{ch:01}
\cite{stratton1941,harrington2001}. Subsection~\ref{sec:ch374} is the HOME bridge between angular spectra and
phase-topology material Subsection~\ref{sec:ch01.2.2} \cite{allen1992,belinfante1940,harrington2001}.

\paragraph{Standing imports.}
Phase singularity condition Eq.~\eqref{eq:ch1.phase-singularity-condition}, winding number
Eq.~\eqref{eq:ch1.phase-winding-number}, and topology-qualified gate Eq.~\eqref{eq:ch1.topology-qualified-gate} are
cited from Subsection~\ref{sec:ch01.2.2} without re-proof \cite{stratton1941,allen1992}. Angular spectrum density
Eq.~\eqref{eq:ch3.unbounded-angular-spectrum-def} and integer index Eq.~\eqref{eq:ch3.integer-ell-def} are presupposed
\cite{hansen1988,harrington2001}.

\paragraph{Assumptions.}
Consider transverse far-zone fields where phase maps \(\arg\E_{t}\) are defined away from singular loci
\cite{stratton1941,jackson1999}. Topological claims require transport gates from Chapter~\ref{ch:01}; Subsection~\ref{sec:ch374}
does not promote phase plots to observables without that gate \cite{harrington2001,stratton1941}.

\paragraph{Phase singularity link to angular harmonics (HOME).}
Near a phase singularity on \(S^{2}\), azimuthal phase windings coincide with integer \(m\) sectors in
Eq.~\eqref{eq:ch3.integer-ell-def} for pure harmonic templates \cite{allen1992,jackson1999,hansen1988}:
\begin{equation}
\label{eq:ch3.phase-singularity-angular-link}
\arg\E_{\ell m}^{(t)}\sim m\,\phi,
\qquad
\text{singularity when }\ell=|m|\ \text{and }m\neq0,
\end{equation}
importing Eq.~\eqref{eq:ch1.phase-winding-number} on circuits enclosing the singularity \cite{stratton1941,allen1992}.
Equation~\eqref{eq:ch3.phase-singularity-angular-link} explains why ``OAM'' discussions often reference integer \(m\):
winding is a phase-topological index, not a separate energy quantum \cite{harrington2001,belinfante1940}.

\paragraph{Topological angular limit (HOME).}
\begin{equation}
\label{eq:ch3.topological-angular-limit}
q
=
\frac{1}{2\pi}\oint_{\mathcal{C}}\nabla_{\perp}\arg\E\cdot d\boldsymbol{\ell}
\;\Longrightarrow\;
m\ \text{in Eq.~\eqref{eq:ch3.integer-ell-def} for pure harmonics},
\end{equation}
with \(\mathcal{C}\) a small loop on a far-zone sphere and transport qualified by
Eq.~\eqref{eq:ch1.topology-qualified-gate} \cite{stratton1941,allen1992,harrington2001}. Equation
\eqref{eq:ch3.topological-angular-limit} is the Subsection~\ref{sec:ch374} closure: topological winding counts enter the
same integer catalog as VSH labels when harmonics are pure \cite{hansen1988,harrington2001}.

\begin{figure}[t]
\centering
\ChOneFigBlock{%
\begin{tikzpicture}[ch1fig, node distance=7mm and 9mm]
  \node[ch1box] (sing) {phase singularity};
  \node[ch1box, right=10mm of sing] (wind) {Eq.~\eqref{eq:ch1.phase-winding-number}};
  \node[ch1box, right=10mm of wind] (m) {$m$ index};
  \node[ch1box, below=10mm of wind] (gate) {Eq.~\eqref{eq:ch1.topology-qualified-gate}};
  \draw[ch1flow] (sing) -- (wind) -- (m);
  \draw[ch1flow] (gate) -- (wind);
\end{tikzpicture}%
}
\caption{Subsection~\ref{sec:ch374} topology chain: singular phase, winding number, integer \(m\) label, gated by transport
qualification.}
\label{fig:ch3.374-topology-chain}
\end{figure}

\paragraph{Physical meaning versus representation hype.}
Phase singularities are features of \emph{phase maps}, not automatic proof of independent angular momentum flux
\cite{belinfante1940,harrington2001,allen1992}. Subsection~\ref{sec:ch374} requires gauge discipline from Section~\ref{sec:ch39}
before elevating \(m\) winds to observable claims \cite{belinfante1940,stratton1941,harrington2001}.

\paragraph{Worked illustration (Laguerre--Gaussian template).}
Paraxial LG modes exhibit \(m\) phase winds on transverse planes \cite{allen1992,jackson1999}. Far-zone projection to
Eq.~\eqref{eq:ch3.unbounded-angular-spectrum-def} spreads power across multiple \(\ell\) sectors unless the beam is pure
harmonic \cite{harrington2001,hansen1988,stratton1941}.

\paragraph{Problem (winding on a small loop).}
\textbf{Problem.} Phase advances by \(2\pi\) around a singularity. Find \(q\) in Eq.~\eqref{eq:ch3.topological-angular-limit}
\cite{allen1992,stratton1941,harrington2001}.

\textbf{Step 1.} Apply Eq.~\eqref{eq:ch1.phase-winding-number} \cite{stratton1941}.

\textbf{Step 2.} Obtain \(q=1\) \cite{allen1992}.

\textbf{Physical meaning.} One wind equals unit topological charge \cite{jackson1999,harrington2001}.

\begin{table}[t]
\centering
\caption{Subsection~\ref{sec:ch374} phase-topology objects.}
\label{tab:ch3.374-topology-objects}
\small
\begin{tabular}{@{}p{0.26\linewidth}p{0.38\linewidth}p{0.30\linewidth}@{}}
\toprule
Object & Equation & Role \\
\midrule
Singularity locus & Eq.~\eqref{eq:ch1.phase-singularity-condition} & phase undefined \\
Winding & Eq.~\eqref{eq:ch1.phase-winding-number} & integer \(q\) \\
Angular link & Eq.~\eqref{eq:ch3.phase-singularity-angular-link} & maps to \(m\) \\
Transport gate & Eq.~\eqref{eq:ch1.topology-qualified-gate} & observable filter \\
\bottomrule
\end{tabular}
\end{table}

\paragraph{Superposition spoils pure winding.}
Superpositions of multiple \(\widehat{c}_{\ell m}^{(\sigma)}\) can destroy a clean singularity even when individual terms
carry winds \cite{harrington2001,hansen1988,stratton1941}. Subsection~\ref{sec:ch374} therefore analyzes singularities on
\emph{decomposed} harmonics, not on arbitrary fitted patterns \cite{harrington2001,devaney2004}.

\paragraph{Worked illustration (two-mode interference null).}
Equal superposition of opposite \(m\) modes can move or cancel singularities \cite{allen1992,jackson1999,harrington2001}.
Topological charge is not additive in watt space; it is a feature of complex field maps \cite{stratton1941,harrington2001}.

\paragraph{Problem (topology gate before OAM claim).}
\textbf{Problem.} A measured vortex pattern is reported as ``OAM beam.'' List gates required
\cite{belinfante1940,harrington2001,stratton1941,allen1992}.

\textbf{Step 1.} Apply Eq.~\eqref{eq:ch1.topology-qualified-gate} \cite{stratton1941}.

\textbf{Step 2.} Decompose into Eq.~\eqref{eq:ch3.integer-ell-def} coefficients \cite{hansen1988}.

\textbf{Step 3.} Report \(\mathcal{P}_{\mathrm{rep}}\) from Eq.~\eqref{eq:ch3.accessible-watt-budget} \cite{harrington2001}.

\textbf{Physical meaning.} Topology diagnoses require transport and gauge discipline \cite{belinfante1940,harrington2001}.

\begin{figure}[t]
\centering
\ChOneFigBlock{%
\begin{tikzpicture}[ch1fig, node distance=7mm and 9mm]
  \node[ch1box] (pat) {pattern plot};
  \node[ch1box, right=10mm of pat] (decomp) {VSH coefficients};
  \node[ch1box, right=10mm of decomp] (topo) {Eq.~\eqref{eq:ch3.topological-angular-limit}};
  \draw[ch1flow] (pat) -- (decomp) -- (topo);
\end{tikzpicture}%
}
\caption{Subsection~\ref{sec:ch374} analysis order: decompose into harmonics before topological winding claims.}
\label{fig:ch3.374-decompose-first}
\end{figure}

\paragraph{Limits at high \(\ell\).}
As \(\ell\) grows, phase maps oscillate rapidly on \(S^{2}\); singularity density can increase without increasing export
watts \cite{hansen1988,harrington2001,stratton1941}. Subsection~\ref{sec:ch374} separates topological complexity from
\(\mathcal{P}_{\mathrm{rad}}\) \cite{poynting1884,harrington2001}.

\paragraph{Gauge caveat (pointer).}
Gauge transformations can move phase singularities in non-transverse potentials; Subsection~\ref{sec:ch374} works on
transverse far-zone \(\E_{t}\) only \cite{belinfante1940,stratton1941}. Full gauge-invariant discussion is Section~\ref{sec:ch39}
\cite{belinfante1940,harrington2001}.

\paragraph{Worked illustration (spiral phase plate).}
A plate imprinting \(\exp(\jj m\phi)\) on a Gaussian beam creates \(m\) winds in the near field; far-zone content still
must be decomposed into Eq.~\eqref{eq:ch3.discrete-continuous-spectrum-bridge} \cite{allen1992,harrington2001,stratton1941}.

\paragraph{Problem (link winding to \(m\) degeneracy).}
\textbf{Problem.} For pure \(\mathbf{u}_{\ell m}^{(\sigma)}\) with \(\ell=|m|\), relate \(q\) to \(m\)
\cite{allen1992,jackson1999,hansen1988}.

\textbf{Step 1.} Use Eq.~\eqref{eq:ch3.phase-singularity-angular-link} \cite{hansen1988}.

\textbf{Step 2.} Identify \(q=m\) for azimuthal circuits \cite{jackson1999}.

\textbf{Physical meaning.} Topological charge aligns with azimuthal index for pure harmonics \cite{allen1992,stratton1941}.

\paragraph{Forward pointer.}
Section~\ref{sec:ch384} analyzes representation dependence of OAM modes; Subsection~\ref{sec:ch374} supplies topological
limits that Section~\ref{sec:ch39} elevates to gauge-invariant language \cite{belinfante1940,allen1992,harrington2001}.

\paragraph{Topological transport budget (HOME).}
Let \(q_{\mathrm{top}}\) be the winding count from Eq.~\eqref{eq:ch3.topological-angular-limit}. Subsection~\ref{sec:ch374}
records the diagnostic inequality
\begin{equation}
\label{eq:ch3.topological-transport-budget}
|q_{\mathrm{top}}|
\le
\max\{|m|:\ \widehat{c}_{\ell m}^{(\sigma)}\neq0,\ \ell=|m|\},
\end{equation}
for pure harmonic dominance in decompositions \cite{allen1992,hansen1988,harrington2001}. Equation
\eqref{eq:ch3.topological-transport-budget} prevents reporting \(|q_{\mathrm{top}}|> \ell_{\max}\) without corresponding
integer support in Eq.~\eqref{eq:ch3.integer-ell-def} \cite{stratton1941,harrington2001}.

\paragraph{Worked illustration (off-axis singularity motion).}
Tilting a vortex beam moves phase singularities in transverse planes without changing far-zone integer spectrum if the same
\(\widehat{c}_{\ell m}^{(\sigma)}\) vector is preserved \cite{allen1992,harrington2001,stratton1941}. Topology of phase maps is
not invariant under arbitrary pattern plots \cite{harrington2001,hansen1988}.

\paragraph{Problem (topology gate checklist).}
\textbf{Problem.} Before claiming OAM transport, verify Eqs.~\eqref{eq:ch1.topology-qualified-gate}, \eqref{eq:ch3.topological-angular-limit},
and \eqref{eq:ch3.accessible-watt-budget} \cite{belinfante1940,harrington2001,stratton1941,allen1992}.

\textbf{Step 1.} Decompose far field into VSH coefficients \cite{hansen1988}.

\textbf{Step 2.} Measure winding on small loops away from nulls of magnitude \cite{allen1992}.

\textbf{Step 3.} Report only \(\mathcal{L}_{\mathrm{acc}}\) watts \cite{harrington2001}.

\textbf{Physical meaning.} Topological claims require harmonic decomposition plus transport gates \cite{belinfante1940}.

\begin{figure}[t]
\centering
\ChOneFigBlock{%
\begin{tikzpicture}[ch1fig, node distance=7mm and 9mm]
  \node[ch1box] (gate) {Eq.~\eqref{eq:ch1.topology-qualified-gate}};
  \node[ch1box, right=10mm of gate] (wind) {Eq.~\eqref{eq:ch3.topological-angular-limit}};
  \node[ch1box, right=10mm of wind] (bud) {Eq.~\eqref{eq:ch3.topological-transport-budget}};
  \draw[ch1flow] (gate) -- (wind) -- (bud);
\end{tikzpicture}%
}
\caption{Subsection~\ref{sec:ch374} gated topology: transport qualification precedes winding and budget inequalities.}
\label{fig:ch3.374-gated-topology}
\end{figure}

\begin{table}[t]
\centering
\caption{Subsection~\ref{sec:ch374} phase-topology audit steps.}
\label{tab:ch3.374-audit-steps}
\small
\begin{tabular}{@{}p{0.08\linewidth}p{0.44\linewidth}p{0.36\linewidth}@{}}
\toprule
Step & Action & Equation \\
\midrule
1 & Gate transport & Eq.~\eqref{eq:ch1.topology-qualified-gate} \\
2 & Decompose & Eq.~\eqref{eq:ch3.integer-ell-def} \\
3 & Count wind & Eq.~\eqref{eq:ch1.phase-winding-number} \\
4 & Budget & Eq.~\eqref{eq:ch3.topological-transport-budget} \\
\bottomrule
\end{tabular}
\end{table}

\paragraph{Paraxial limit pointer.}
Paraxial LG vortices approximate integer \(m\) winds only within paraxial validity \cite{allen1992,jackson1999}. Subsection~\ref{sec:ch393}
records scalar/paraxial reductions; Subsection~\ref{sec:ch374} does not extend paraxial winds to full vector spectra without
decomposition \cite{harrington2001,stratton1941}.

\paragraph{Worked illustration (multiple singularities).}
Superpositions can carry several singularities with net \(q\) summing only when branch cuts are tracked \cite{allen1992,jackson1999}.
Far-zone \(\mathcal{P}_{\mathrm{rad}}\) remains quadratic in coefficients Eq.~\eqref{eq:ch3.mode-power-from-coefficient}
regardless of singularity count \cite{poynting1884,harrington2001,stratton1941}.

\paragraph{Problem (compare \(q\) and \(m\) after tilt).}
\textbf{Problem.} A tilted beam shows \(q=1\) on a loop but dominant \(\ell=2\) coefficients. Explain
\cite{allen1992,harrington2001,hansen1988,stratton1941}.

\textbf{Step 1.} Recognize superposition of multiple \(m\) sectors \cite{hansen1988}.

\textbf{Step 2.} Apply Eq.~\eqref{eq:ch3.topological-transport-budget} as diagnostic bound \cite{harrington2001}.

\textbf{Physical meaning.} Local winding does not equal a single multipole label \cite{stratton1941,allen1992}.

\paragraph{Branch cuts and measurement domains.}
Phase windings are defined on loops that avoid magnitude nulls of \(\E_{t}\) \cite{allen1992,jackson1999}. Subsection~\ref{sec:ch374}
requires declaring loop locations on \(\Gamma_R\) when reporting \(q_{\mathrm{top}}\) \cite{harrington2001,stratton1941,hansen1988}.

\paragraph{Worked illustration (circular polarization versus phase vortex).}
Circular polarization can produce rotating phase patterns without a transverse phase singularity on \(\Gamma_R\)
\cite{jackson1999,belinfante1940,harrington2001}. Subsection~\ref{sec:ch374} distinguishes polarization phase from azimuthal
vortex winds Eq.~\eqref{eq:ch3.phase-singularity-angular-link} \cite{stratton1941,allen1992}.

\paragraph{Problem (null of magnitude versus phase singularity).}
\textbf{Problem.} At a radiation null, \(|\E_{t}|=0\) but phase is undefined. Classify using
Eq.~\eqref{eq:ch1.phase-singularity-condition} \cite{allen1992,stratton1941,harrington2001}.

\textbf{Step 1.} Magnitude nulls need not be phase singularities \cite{jackson1999}.

\textbf{Step 2.} Test winding on a loop enclosing the null \cite{allen1992}.

\textbf{Physical meaning.} Nulls and vortices are distinct features \cite{harrington2001,stratton1941}.

\begin{figure}[t]
\centering
\ChOneFigBlock{%
\begin{tikzpicture}[ch1fig, node distance=7mm and 9mm]
  \node[ch1box] (null) {magnitude null};
  \node[ch1box, right=10mm of null] (vort) {phase vortex};
  \node[ch1box, right=10mm of vort] (eqs) {Eqs.~\eqref{eq:ch1.phase-singularity-condition}--\eqref{eq:ch3.topological-angular-limit}};
  \draw[ch1flow] (null) -- (eqs);
  \draw[ch1flow] (vort) -- (eqs);
\end{tikzpicture}%
}
\caption{Subsection~\ref{sec:ch374} separates magnitude nulls from phase vortices before topological counting.}
\label{fig:ch3.374-null-vs-vortex}
\end{figure}

\paragraph{Topological limits on measurable \(m\) winds (expanded).}
Eq.~\eqref{eq:ch3.topological-transport-budget} is a consistency check, not a substitute for helicity or gauge-invariant
observables Section~\ref{sec:ch39} \cite{belinfante1940,harrington2001,allen1992}. Subsection~\ref{sec:ch374} therefore positions
\(q_{\mathrm{top}}\) as a \emph{diagnostic} on decomposed harmonics, not as a standalone OAM quantum \cite{harrington2001,stratton1941}.

\paragraph{Worked illustration (array of vortices).}
Phased arrays can imprint multiple vortices across \(\Gamma_R\); each may satisfy Eq.~\eqref{eq:ch1.phase-winding-number} locally
while global \(\widehat{c}_{\ell m}^{(\sigma)}\) vectors remain sparse in \(\ell\) \cite{allen1992,harrington2001,stratton1941}.
Topological maps and multipole coefficients are complementary data products \cite{hansen1988,harrington2001}.

\paragraph{Problem (gate failure interpretation).}
\textbf{Problem.} A winding count is measured but Eq.~\eqref{eq:ch1.topology-qualified-gate} fails. What can be claimed?
\cite{belinfante1940,harrington2001,stratton1941,allen1992}.

\textbf{Step 1.} Report winding as phase-map feature only \cite{allen1992}.

\textbf{Step 2.} Withhold transport/OAM claims \cite{belinfante1940,harrington2001}.

\textbf{Physical meaning.} Topology without transport gate is not an export claim \cite{stratton1941,harrington2001}.

\begin{table}[t]
\centering
\caption{Subsection~\ref{sec:ch374} claim strength ladder.}
\label{tab:ch3.374-claim-ladder}
\small
\begin{tabular}{@{}p{0.22\linewidth}p{0.68\linewidth}@{}}
\toprule
Level & Allowed claim \\
\midrule
Phase map & winding on declared loop \\
Decomposed & Eq.~\eqref{eq:ch3.topological-transport-budget} \\
Gated & transport per Eq.~\eqref{eq:ch1.topology-qualified-gate} \\
\bottomrule
\end{tabular}
\end{table}

\paragraph{Peer-review rule for OAM claims.}
Publications must state which row of Table~\ref{tab:ch3.374-claim-ladder} is claimed and include VSH coefficients for level~2 or higher
\cite{belinfante1940,harrington2001,allen1992,stratton1941}. Subsection~\ref{sec:ch374} rejects level~3 claims without
Eq.~\eqref{eq:ch1.topology-qualified-gate} documentation \cite{harrington2001,belinfante1940}.

Subsection~\ref{sec:ch374} has connected phase singularities and topological limits to angular spectra through
Eqs.~\eqref{eq:ch3.phase-singularity-angular-link}--\eqref{eq:ch3.topological-angular-limit}
\cite{allen1992,stratton1941,harrington2001}. Section~\ref{sec:ch37} is complete: unbounded angular spectra,
boundary quantization, integer index discipline, and topological limits now share one \(\mathcal{R}_\omega\) narrative
\cite{hansen1988,jackson1999}.

The four subsections of Section~\ref{sec:ch37} unify discrete and continuous angular coordinates:
Eqs.~\eqref{eq:ch3.unbounded-angular-spectrum-def}--\eqref{eq:ch3.discrete-continuous-spectrum-bridge} in
Subsection~\ref{sec:ch371}; Eqs.~\eqref{eq:ch3.boundary-quantized-spectrum}--\eqref{eq:ch3.cavity-angular-index} in
Subsection~\ref{sec:ch372}; Eqs.~\eqref{eq:ch3.integer-ell-def}--\eqref{eq:ch3.beam-order-vs-ell-warning} in
Subsection~\ref{sec:ch373}; and Eqs.~\eqref{eq:ch3.phase-singularity-angular-link}--\eqref{eq:ch3.topological-angular-limit}
in Subsection~\ref{sec:ch374} \cite{stratton1941,jackson1999,harrington2001,allen1992}. Section~\ref{sec:ch38} now makes
basis changes among plane-wave, cylindrical, parabolic, and spherical representations explicit
\cite{hansen1988,stratton1941}.

\section{Change of Representation Basis}
\label{sec:ch38}

The same radiative state can be written in plane-wave, spherical, cylindrical, or parabolic coordinates;
only the coefficients change. Section~\ref{sec:ch38} makes those changes explicit matrices rather than
hand-waving equivalences \cite{stratton1941,hansen1988}. Subsection~\ref{sec:ch381} develops plane-wave and
Fourier representations, extending the operator view of Subsection~\ref{sec:ch291}. Subsection~\ref{sec:ch382}
introduces cylindrical and parabolic bases used in feeds, reflectors, and laser cavities.
Subsection~\ref{sec:ch383} constructs equivalence transformations among bases, importing the VSH definitions
of Subsection~\ref{sec:ch341} as the spherical anchor. Subsection~\ref{sec:ch384} analyzes representation
dependence of OAM modes, using spin-orbit limits from Subsection~\ref{sec:ch01.4.5} \cite{allen1992,belinfante1940}.


\subsection{Plane-wave and Fourier representations}
\label{sec:ch381}
% TOC tag: [HOME][USE SS2.9.1]

A \emph{plane-wave representation} expands outgoing radiative content as a superposition of transverse
eigenmodes indexed by wave vector \(\mathbf{k}\) on the propagating sheet \(k_{t}\le k_{0}\)
\cite{stratton1941,jackson1999,harrington2001}. A \emph{Fourier representation} is the lateral chart of that
expansion on a measurement plane \(z=z_{0}\), where \((k_{x},k_{y})\) play the role of spatial frequencies
\cite{stratton1941,jackson1999}. Subsection~\ref{sec:ch381} is the HOME basis layer that upgrades
Subsection~\ref{sec:ch291} plane-wave spectra to the normalized watt-calibrated coordinates of
Section~\ref{sec:ch35} and the angular-spectrum bridge of Subsection~\ref{sec:ch371}
\cite{hansen1988,harrington2001}.

\paragraph{Standing imports (no re-derivation).}
Plane-wave angular spectrum Eqs.~\eqref{eq:ch2.plane-wave-angular-spectrum-def}--\eqref{eq:ch2.plane-wave-z-propagation},
radiative lateral spectrum Eq.~\eqref{eq:ch2.radiative-lateral-spectrum-def}, radiation operator
Eq.~\eqref{eq:ch2.plane-wave-radiation-operator}, plane-wave accessibility filter
Eq.~\eqref{eq:ch2.plane-wave-accessibility-filter}, continuum bridge Eq.~\eqref{eq:ch3.plane-wave-continuum-bridge},
flux-fixed Parseval Eq.~\eqref{eq:ch3.continuous-parseval-flux-fixed}, and plane-wave delta normalization
Eq.~\eqref{eq:ch3.plane-wave-delta-normalization} enter by reference \cite{stratton1941,jackson1999,harrington2001}.
Subsection~\ref{sec:ch381} does not re-derive Fourier sign conventions or Maxwell's equations
\cite{stratton1941,harrington2001}.

\paragraph{Assumptions.}
Fix \(\omega>0\), homogeneous reciprocal exterior media, outgoing \(\mathcal{F}_{\mathrm{rad}}\), phasor time
\(\exp(-\jj\omega t)\), and a declared measurement plane \(z=z_{0}\) when lateral Fourier charts are used
\cite{stratton1941,harrington2001}. Evanescent sheets \(k_{t}>k_{0}\) may appear in near-field data but lie outside
export sums unless explicitly labeled reactive \cite{stratton1941,jackson1999}.

\paragraph{Plane-wave basis resolution (HOME).}
Let \(\widetilde{\mathbf{E}}_{\mathrm{rad}}(k_{x},k_{y})\) denote the radiative lateral spectrum from
Eq.~\eqref{eq:ch2.radiative-lateral-spectrum-def}. Subsection~\ref{sec:ch381} fixes the watt-ready plane-wave expansion
\begin{equation}
\label{eq:ch3.plane-wave-basis-resolution}
\mathbf{v}_\omega(\mathbf{r})
=
\int_{\mathbb{R}^{2}}
\widetilde{\mathbf{A}}(k_{x},k_{y})\,
\boldsymbol{\mathcal{E}}(k_{x},k_{y};\mathbf{r})\,dk_{x}\,dk_{y},
\qquad
\boldsymbol{\mathcal{E}}\in\mathcal{F}_{\mathrm{rad}},
\end{equation}
with \(\widetilde{\mathbf{A}}\) the flux-normalized coefficient density tied to
Eq.~\eqref{eq:ch3.plane-wave-delta-normalization} \cite{stratton1941,harrington2001,jackson1999}. Equation
\eqref{eq:ch3.plane-wave-basis-resolution} states the physical meaning of plane-wave representations in Volume~I:
\(\widetilde{\mathbf{A}}\) is not a raw FFT of \(\E\); it is the spectral coordinate of
\(\mathcal{R}_\omega[\mathfrak{D}]\) after outgoing selection \cite{harrington2001,stratton1941}. Each
\((k_{x},k_{y})\) pair labels a transverse propagation direction on the \(z=z_{0}\) chart; only admissible outgoing
\(k_{z}\) from Eq.~\eqref{eq:ch2.plane-wave-z-propagation} contributes to far-zone export \cite{jackson1999,stratton1941}.

\paragraph{Fourier--radiation filter chain (HOME).}
Subsection~\ref{sec:ch381} packages the laboratory pipeline as
\begin{equation}
\label{eq:ch3.fourier-radiation-filter-chain}
\widetilde{\mathbf{A}}(k_{x},k_{y})
=
\mathcal{F}_{z_{0}}\bigl[\Pi_{\mathrm{rad}}\,\mathbf{E}(\mathbf{r};z_{0})\bigr](k_{x},k_{y}),
\qquad
\mathcal{F}_{z_{0}}:\ \text{lateral Fourier on plane }z=z_{0},
\end{equation}
importing Eq.~\eqref{eq:ch2.plane-wave-radiation-operator} as the operator that precedes \(\mathcal{F}_{z_{0}}\) in source-driven
problems \cite{stratton1941,harrington2001}. Equation~\eqref{eq:ch3.fourier-radiation-filter-chain} is the Subsection~\ref{sec:ch381}
discipline law: raw FFTs of near-field \(\E\) omit \(\Pi_{\mathrm{rad}}\) and therefore misstate watts
\cite{harrington2001,stratton1941}. Mathematically, \(\mathcal{F}_{z_{0}}\) is a unitary map only after the outgoing
measure and normalization of Section~\ref{sec:ch35} are fixed \cite{hansen1988,harrington2001}.

\begin{figure}[t]
\centering
\ChOneFigBlock{%
\begin{tikzpicture}[ch1fig, node distance=7mm and 9mm]
  \node[ch1box] (field) {$\mathbf{E}$ on $z_{0}$};
  \node[ch1box, right=10mm of field] (pi) {$\Pi_{\mathrm{rad}}$};
  \node[ch1box, right=10mm of pi] (fft) {$\mathcal{F}_{z_{0}}$};
  \node[ch1box, right=10mm of fft] (A) {Eq.~\eqref{eq:ch3.plane-wave-basis-resolution}};
  \node[ch1box, below=10mm of pi] (op) {Eq.~\eqref{eq:ch2.plane-wave-radiation-operator}};
  \draw[ch1flow] (field) -- (pi) -- (fft) -- (A);
  \draw[ch1flow] (op) -- (pi);
\end{tikzpicture}%
}
\caption{Subsection~\ref{sec:ch381} Fourier chain: outgoing filter precedes lateral FFT before plane-wave coefficients are reported.}
\label{fig:ch3.381-fourier-chain}
\end{figure}

\begin{theorem}[Plane-wave/Fourier equivalence on outgoing domains (HOME)]
\label{thm:ch3.plane-wave-Fourier-equivalence}
Every \(\mathbf{v}_\omega\in\mathcal{F}_{\mathrm{rad}}\) admits Eq.~\eqref{eq:ch3.plane-wave-basis-resolution} with
Parseval closure Eq.~\eqref{eq:ch3.continuous-parseval-flux-fixed} and identification to the \(S^{2}\) chart
Eq.~\eqref{eq:ch3.plane-wave-continuum-bridge} \cite{stratton1941,jackson1999,hansen1988,kato1995}.
\end{theorem}
\begin{proof}[Proof sketch]
Combine Theorem~\ref{thm:ch3.unbounded-angular-spectrum} with lateral spectrum definition
Eq.~\eqref{eq:ch2.radiative-lateral-spectrum-def} and Section~\ref{sec:ch35} normalization \cite{stratton1941,harrington2001}.
% PROOF: AUTHOR REVIEW
\end{proof}

\paragraph{Physical meaning of lateral Fourier coefficients.}
\(\widetilde{\mathbf{A}}(k_{x},k_{y})\) measures how much radiative content propagates with transverse wavenumber
\((k_{x},k_{y})\) through the \(z=z_{0}\) chart \cite{stratton1941,jackson1999}. Large \(|k_{x}|,|k_{y}|\) correspond to
steep propagation angles; \(k_{t}\approx0\) corresponds to broadside-dominated beams \cite{harrington2001,jackson1999}.
Subsection~\ref{sec:ch381} therefore reads compact-range scans as samples of Eq.~\eqref{eq:ch3.plane-wave-basis-resolution},
not as direct images of far-zone \(S^{2}\) spectra until the lift of Subsection~\ref{sec:ch371} is applied
\cite{hansen1988,stratton1941}.

\paragraph{Worked illustration (Gaussian lateral spectrum).}
A Gaussian \(\widetilde{\mathbf{A}}\propto\exp[-(k_{x}^{2}+k_{y}^{2})/k_{0}^{2}w_{0}^{2}]\) concentrates export on
small transverse wavenumbers \cite{jackson1999,harrington2001}. Integrating \(| \widetilde{\mathbf{A}}|^{2}\) with the
measure of Eq.~\eqref{eq:ch3.plane-wave-delta-normalization} recovers total watts on \(\mathcal{F}_{\mathrm{rad}}\)
\cite{poynting1884,stratton1941}. The Gaussian is a coordinate shape, not a separate physical mode \cite{harrington2001}.

\paragraph{Problem (diagnose raw FFT overcount).}
\textbf{Problem.} A near-field FFT reports \(1.2\,\mathrm{W}\) while the far-zone meter reads \(1.0\,\mathrm{W}\). Reconcile using
Eq.~\eqref{eq:ch3.fourier-radiation-filter-chain} \cite{harrington2001,stratton1941,jackson1999}.

\textbf{Step 1.} Identify reactive/evanescent admixture in raw FFT \cite{stratton1941}.

\textbf{Step 2.} Apply \(\Pi_{\mathrm{rad}}\) before \(\mathcal{F}_{z_{0}}\) \cite{harrington2001}.

\textbf{Step 3.} Renormalize with Eq.~\eqref{eq:ch3.plane-wave-delta-normalization} \cite{hansen1988}.

\textbf{Physical meaning.} Fourier coefficients count radiative content only after outgoing selection \cite{harrington2001}.

\begin{table}[t]
\centering
\caption{Subsection~\ref{sec:ch381} plane-wave/Fourier objects.}
\label{tab:ch3.381-fourier-objects}
\small
\begin{tabular}{@{}p{0.24\linewidth}p{0.38\linewidth}p{0.30\linewidth}@{}}
\toprule
Object & Equation & Role \\
\midrule
Lateral spectrum & Eq.~\eqref{eq:ch2.radiative-lateral-spectrum-def} & measurement chart \\
Basis resolution & Eq.~\eqref{eq:ch3.plane-wave-basis-resolution} & expansion \\
Filter chain & Eq.~\eqref{eq:ch3.fourier-radiation-filter-chain} & laboratory pipeline \\
Parseval & Eq.~\eqref{eq:ch3.continuous-parseval-flux-fixed} & watt closure \\
\bottomrule
\end{tabular}
\end{table}

\paragraph{Propagating versus evanescent Fourier sheets.}
Eq.~\eqref{eq:ch2.plane-wave-z-propagation} splits the \(k\)-plane into propagating \(k_{t}\le k_{0}\) and evanescent
\(k_{t}>k_{0}\) sheets \cite{jackson1999,stratton1941}. Subsection~\ref{sec:ch381} assigns export watts only to the
propagating sheet in Eq.~\eqref{eq:ch3.plane-wave-basis-resolution}; evanescent amplitudes may be large near sources but
do not survive \(\Pi_{\mathrm{rad}}\) \cite{harrington2001,stratton1941}. Laboratory reports must label which sheet is plotted
\cite{hansen1988,harrington2001}.

\paragraph{Worked illustration (tilted beam in \(k\)-space).}
A broadside beam peaks near \(k_{x}=k_{y}=0\); a tilted beam shifts the peak to finite \((k_{x},k_{y})\) while preserving
total watts under Parseval \cite{jackson1999,harrington2001}. Rotation of the laboratory rotates the \(k\)-plane chart; it does not
change \(\mathcal{P}_{\mathrm{rad}}\) \cite{stratton1941,hansen1988}.

\paragraph{Problem (map \(k_{x},k_{y}\) to polar angle).}
\textbf{Problem.} For propagating modes, relate \((k_{x},k_{y})\) to \((\theta,\phi)\) on \(S^{2}\)
\cite{jackson1999,stratton1941,harrington2001}.

\textbf{Step 1.} Use \(k_{t}=\sqrt{k_{x}^{2}+k_{y}^{2}}\) and \(k_{z}=\sqrt{k_{0}^{2}-k_{t}^{2}}\) \cite{jackson1999}.

\textbf{Step 2.} Identify \(\sin\theta=k_{t}/k_{0}\), \(\phi=\operatorname{atan2}(k_{y},k_{x})\) \cite{stratton1941}.

\textbf{Physical meaning.} Lateral Fourier charts and direction spheres are linked by outgoing kinematics \cite{hansen1988}.

\paragraph{Accessibility filter on Fourier data.}
Eq.~\eqref{eq:ch2.plane-wave-accessibility-filter} restricts which \((k_{x},k_{y})\) cells enter reports after geometry and
measurement windows \cite{harrington2001,stratton1941}. Subsection~\ref{sec:ch381} applies that filter before any comparison with
VSH coefficients Eq.~\eqref{eq:ch3.normalized-coefficient-resolution} \cite{hansen1988,harrington2001}.

\begin{figure}[t]
\centering
\ChOneFigBlock{%
\begin{tikzpicture}[ch1fig, node distance=7mm and 9mm]
  \node[ch1box] (k) {$k_{x},k_{y}$ plane};
  \node[ch1box, right=10mm of k] (prop) {$k_{t}\le k_{0}$};
  \node[ch1box, right=10mm of prop] (s2) {$S^{2}$ lift};
  \node[ch1box, below=10mm of prop] (eva) {evanescent sheet};
  \draw[ch1flow] (k) -- (prop) -- (s2);
  \draw[ch1flow] (eva) -- (k);
\end{tikzpicture}%
}
\caption{Subsection~\ref{sec:ch381} propagating sheet: only \(k_{t}\le k_{0}\) contributes to outgoing plane-wave watts.}
\label{fig:ch3.381-propagating-sheet}
\end{figure}

\paragraph{Near-to-far linkage without re-derivation.}
Eq.~\eqref{eq:ch2.plane-wave-radiation-operator} already maps sources to \(\widetilde{\mathbf{A}}\); Subsection~\ref{sec:ch381}
uses that map as given \cite{stratton1941,harrington2001}. The present subsection supplies the \emph{basis} interpretation of the
output: \(\widetilde{\mathbf{A}}\) is a coordinate on \(\mathcal{F}_{\mathrm{rad}}\), dual to VSH coefficients
\cite{hansen1988,jackson1999}.

\paragraph{Worked illustration (discrete FFT sampling).}
Digital FFTs sample \(\widetilde{\mathbf{A}}\) on a grid; quadrature weights must match Eq.~\eqref{eq:ch3.plane-wave-delta-normalization}
to preserve watts \cite{harrington2001,stratton1941}. Uneven spacing near \(k_{t}\approx k_{0}\) requires explicit Jacobian factors before
integration \cite{hansen1988,jackson1999}.

\paragraph{Problem (verify Parseval on a grid).}
\textbf{Problem.} Given sampled \(|\widetilde{\mathbf{A}}|^{2}\) on a uniform \(k_{x},k_{y}\) grid, outline watt summation
\cite{harrington2001,stratton1941,hansen1988}.

\textbf{Step 1.} Multiply by \(\Delta k_{x}\Delta k_{y}\) and normalization factor \cite{stratton1941}.

\textbf{Step 2.} Restrict sum to \(k_{t}\le k_{0}\) \cite{jackson1999}.

\textbf{Step 3.} Compare to Eq.~\eqref{eq:ch3.continuous-parseval-flux-fixed} continuum limit \cite{hansen1988}.

\textbf{Physical meaning.} Discrete sums approximate continuous Parseval only with correct measure \cite{harrington2001}.

\paragraph{Polarization tracks in plane-wave bases.}
Each \((k_{x},k_{y})\) carries two transverse polarizations \(\sigma\in\{1,2\}\) in
Eq.~\eqref{eq:ch3.plane-wave-continuum-bridge} \cite{stratton1941,jackson1999}. Subsection~\ref{sec:ch381} requires \(\sigma\) in every
coefficient tuple; collapsing to scalar amplitude erases helicity information used in Subsection~\ref{sec:ch384}
\cite{belinfante1940,harrington2001}.

\begin{figure}[t]
\centering
\ChOneFigBlock{%
\begin{tikzpicture}[ch1fig, node distance=7mm and 9mm]
  \node[ch1box] (pw) {Eq.~\eqref{eq:ch3.plane-wave-basis-resolution}};
  \node[ch1box, right=10mm of pw] (vsh) {Eq.~\eqref{eq:ch3.normalized-coefficient-resolution}};
  \node[ch1box, right=10mm of vsh] (bridge) {Subsection~\ref{sec:ch383}};
  \draw[ch1flow] (pw) -- (bridge) -- (vsh);
\end{tikzpicture}%
}
\caption{Subsection~\ref{sec:ch381} plane-wave chart feeds explicit basis-change matrices in Subsection~\ref{sec:ch383}.}
\label{fig:ch3.381-to-transform}
\end{figure}

\begin{table}[t]
\centering
\caption{Subsection~\ref{sec:ch381} laboratory Fourier checklist.}
\label{tab:ch3.381-checklist}
\small
\begin{tabular}{@{}p{0.08\linewidth}p{0.52\linewidth}p{0.32\linewidth}@{}}
\toprule
Step & Action & Equation \\
\midrule
1 & Outgoing filter & Eq.~\eqref{eq:ch3.fourier-radiation-filter-chain} \\
2 & Lateral FFT & Eq.~\eqref{eq:ch2.radiative-lateral-spectrum-def} \\
3 & Normalize & Eq.~\eqref{eq:ch3.plane-wave-delta-normalization} \\
4 & Parseval & Eq.~\eqref{eq:ch3.continuous-parseval-flux-fixed} \\
\bottomrule
\end{tabular}
\end{table}

\paragraph{Propagating-sheet Jacobian (measure transport).}
On the outgoing sheet \(k_{t}=\sqrt{k_{x}^{2}+k_{y}^{2}}\le k_{0}\), directions \((\theta,\phi)\) and lateral coordinates are linked by
\(k_{x}=k_{0}\sin\theta\cos\phi\), \(k_{y}=k_{0}\sin\theta\sin\phi\) \cite{jackson1999,stratton1941}. Watt closure under the change of
variables requires the Jacobian factor that converts \(dk_{x}\,dk_{y}\) into solid-angle measure \(d\Omega\):
\begin{equation}
\label{eq:ch3.plane-wave-direction-jacobian}
d\Omega
=
\frac{k_{z}}{k_{0}^{2}}\,dk_{x}\,dk_{y},
\qquad
k_{z}=\sqrt{k_{0}^{2}-k_{x}^{2}-k_{y}^{2}},
\end{equation}
with \(k_{z}>0\) on the propagating branch \cite{stratton1941,jackson1999,hansen1988}. Equation~\eqref{eq:ch3.plane-wave-direction-jacobian}
is the geometric content behind Eq.~\eqref{eq:ch3.plane-wave-continuum-bridge}: plane-wave amplitudes are densities, and comparing a lateral FFT
to an \(S^{2}\) plot without this factor misstates peak directivity and integrated watts \cite{harrington2001,jackson1999}. Near \(k_{t}\approx k_{0}\)
the Jacobian amplifies grazing directions; compact-range grids must refine sampling there before claiming angular resolution
\cite{hansen1988,stratton1941}.

\paragraph{Helicity-resolved lateral spectrum.}
Each propagating \((k_{x},k_{y})\) cell carries two transverse polarizations \(\sigma\in\{+,-\}\) in
Eq.~\eqref{eq:ch3.plane-wave-continuum-bridge} \cite{belinfante1940,stratton1941}. Subsection~\ref{sec:ch381} therefore decomposes
\(\widetilde{\mathbf{A}}\) into helicity-sector sheets
\begin{equation}
\label{eq:ch3.plane-wave-helicity-spectrum}
\mathcal{P}_{\mathrm{rad}}^{(\sigma)}
=
\int_{k_{t}\le k_{0}}
\bigl|\widetilde{\mathbf{A}}_{\sigma}(k_{x},k_{y})\bigr|^{2}\,
\frac{k_{z}}{k_{0}^{2}}\,dk_{x}\,dk_{y},
\qquad
\sum_{\sigma}\mathcal{P}_{\mathrm{rad}}^{(\sigma)}=\mathcal{P}_{\mathrm{rad}},
\end{equation}
importing helicity density Eq.~\eqref{eq:ch1.helicity-density} only after full vector reconstruction on \(\mathcal{R}_\omega\)
\cite{belinfante1940,harrington2001,jackson1999}. Equation~\eqref{eq:ch3.plane-wave-helicity-spectrum} prevents a common laboratory error:
collapsing a dual-polarization scan to scalar magnitude erases sector watts that Subsection~\ref{sec:ch392} treats as conserved when gauge gates pass
\cite{belinfante1940,hansen1988}. Circular-feed designs that imprint opposite helicity lobes therefore require \(\sigma\)-resolved archives, not a
single ``angular spectrum'' image \cite{harrington2001,stratton1941}.

\paragraph{Finite-aperture aliasing and angular Nyquist bound.}
A scan window of lateral extent \(L_{x},L_{y}\) on \(z=z_{0}\) imposes a minimum resolvable transverse wavenumber
\(\Delta k\sim\pi/L\) \cite{jackson1999,harrington2001}. Subsection~\ref{sec:ch381} records the associated angular Nyquist limit
\begin{equation}
\label{eq:ch3.plane-wave-angular-nyquist}
\sin\theta_{\min}
\approx
\frac{\lambda}{2L_{\mathrm{eff}}},
\qquad
L_{\mathrm{eff}}=\min(L_{x},L_{y}),
\end{equation}
with \(\lambda=2\pi/k_{0}\) \cite{stratton1941,hansen1988}. Equation~\eqref{eq:ch3.plane-wave-angular-nyquist} states when a discrete FFT can
resolve two plane-wave components: attempting to read finer structure is spectral aliasing, not sub-wavelength physics
\cite{harrington2001,devaney2004}. Array designers should compare \(\theta_{\min}\) to the beamwidth implied by
Eq.~\eqref{eq:ch3.plane-wave-basis-resolution} before promoting FFT peaks to propagation angles \cite{jackson1999,stratton1941}.

\begin{figure}[t]
\centering
\ChOneFigBlock{%
\begin{tikzpicture}[ch1fig, node distance=7mm and 9mm]
  \node[ch1box] (fft) {lateral FFT};
  \node[ch1box, right=10mm of fft] (jac) {Eq.~\eqref{eq:ch3.plane-wave-direction-jacobian}};
  \node[ch1box, right=10mm of jac] (hel) {Eq.~\eqref{eq:ch3.plane-wave-helicity-spectrum}};
  \node[ch1box, below=10mm of fft] (nyq) {Eq.~\eqref{eq:ch3.plane-wave-angular-nyquist}};
  \draw[ch1flow] (fft) -- (jac) -- (hel);
  \draw[ch1flow] (nyq) -- (fft);
\end{tikzpicture}%
}
\caption{Subsection~\ref{sec:ch381} measure chain: Jacobian converts FFT cells to solid-angle watts; Nyquist bound limits resolvable directions.}
\label{fig:ch3.381-measure-chain}
\end{figure}

\paragraph{Campaign archive fields (Subsection~\ref{sec:ch381}).}
Archive \(z_{0}\), FFT grid, \(\Pi_{\mathrm{rad}}\) certificate, normalization constants, propagating mask, and accessibility filter
Eq.~\eqref{eq:ch2.plane-wave-accessibility-filter} with every \(\widetilde{\mathbf{A}}\) file \cite{harrington2001,stratton1941,hansen1988}.
Reproducible plane-wave reports require the measure and filter chain, not amplitudes alone \cite{harrington2001,jackson1999}.

\paragraph{Vocabulary and implicit chart philosophy.}
The scientific vocabulary of Subsection~\ref{sec:ch381} is deliberately dual: \emph{plane-wave} names a mode indexed by
propagation direction on the unit sphere of directions, while \emph{Fourier} names the same information after a lateral
chart \((k_{x},k_{y})\) is chosen on a measurement plane \cite{jackson1999,stratton1941}. Neither word, by itself,
specifies the outgoing radiation filter of Eq.~\eqref{eq:ch2.plane-wave-accessibility-filter}; the implicit concept is
\emph{spatial frequency as a coordinate on propagating content}, not as a label for evanescent or reactive sheets
\cite{harrington2001,jackson1999}. Equation~\eqref{eq:ch3.plane-wave-basis-resolution} therefore embeds three commitments
at once: a transverse polarization gauge on each \(\mathbf{k}\), a watt normalization inherited from
Section~\ref{sec:ch35}, and a far-zone radiation certificate compatible with \(\mathcal{F}_{\mathrm{rad}}\)
\cite{stratton1941,hansen1988}. The Fourier chain
Eq.~\eqref{eq:ch3.fourier-radiation-filter-chain} makes the filter explicit: lateral spectra are meaningful only after
the propagating hemisphere is declared and the evanescent null is enforced on \(\mathcal{R}_\omega\)
\cite{harrington2001,jackson1999}. Theorem~\ref{thm:ch3.plane-wave-Fourier-equivalence} is not a convenience identity;
it states that two charts describe one radiative object when---and only when---the same measure, outgoing phase
convention \(\exp(-\jj\omega t)\), and accessibility screen are shared \cite{stratton1941,harrington2001}. Readers who
collapse ``Fourier transform'' into ``angular spectrum'' without naming the plane \(z=z_{0}\) and the filter order
confuse \emph{representation} with \emph{observation}: the embedded mathematics always assumes a radiative equivalence
class, not an arbitrary near-field snapshot \cite{devaney2004,harrington2001}.

\paragraph{Steerable-beam synthesis in lateral \(k\)-space.}
Phased arrays and metasurfaces steer beams by imposing a linear phase ramp across aperture samples, which translates to a shift
\((k_{x0},k_{y0})\) of the spectral peak in Eq.~\eqref{eq:ch3.plane-wave-basis-resolution} \cite{jackson1999,harrington2001,stratton1941}.
Subsection~\ref{sec:ch381} records the synthesis law
\begin{equation}
\label{eq:ch3.plane-wave-steering-phase}
\widetilde{\mathbf{A}}_{\mathrm{steer}}(k_{x},k_{y})
=
\widetilde{\mathbf{A}}_{0}(k_{x}-k_{x0},\,k_{y}-k_{y0})\,
\exp\!\bigl(-\jj\,(k_{x0}x+k_{y0}y)\bigr),
\end{equation}
valid on the propagating sheet with Jacobian weight Eq.~\eqref{eq:ch3.plane-wave-direction-jacobian} included in watt integrals
\cite{hansen1988,jackson1999}. Mathematically, Eq.~\eqref{eq:ch3.plane-wave-steering-phase} is a convolution in spatial frequency; physically,
it is a rigid rotation of the radiative lobes on \(S^{2}\) when \(|k_{x0}^{2}+k_{y0}^{2}|\le k_{0}^{2}\) \cite{stratton1941,harrington2001}.
Grating lobes appear when the shift pushes content toward \(k_{t}\approx k_{0}\); Eq.~\eqref{eq:ch3.plane-wave-angular-nyquist} then limits
whether the steered peak is resolvable or aliased \cite{hansen1988,devaney2004}. Campaigns must archive \((k_{x0},k_{y0})\), aperture weights, and
\(\sigma\) tracks because steering without helicity metadata can rotate one sector while appearing to steer ``the beam'' \cite{belinfante1940,harrington2001}.

\paragraph{Broadside-to-grazing watt distortion (illustration).}
When a narrow \(\widetilde{\mathbf{A}}_{0}\) peak is shifted toward grazing, the Jacobian factor \(k_{z}/k_{0}^{2}\) shrinks and the same spectral
amplitude can integrate to fewer watts on \(\Gamma_{R}\) \cite{jackson1999,poynting1884,stratton1941}. Subsection~\ref{sec:ch381} treats that as a
\emph{measure} effect, not as dissipation: total \(\mathcal{P}_{\mathrm{rad}}\) is preserved only when the full propagating sheet integral is
performed with Eq.~\eqref{eq:ch3.continuous-parseval-flux-fixed} \cite{harrington2001,hansen1988}. Comparing FFT peak heights across steer
angles without Jacobian weighting routinely mis-rank sidelobe levels \cite{jackson1999,stratton1941}.

Subsection~\ref{sec:ch381} has fixed plane-wave and Fourier representations through
Eqs.~\eqref{eq:ch3.plane-wave-basis-resolution}--\eqref{eq:ch3.fourier-radiation-filter-chain}
\cite{stratton1941,jackson1999,harrington2001}. Subsection~\ref{sec:ch382} introduces cylindrical and parabolic coordinates used in
feeds and focused optics \cite{hansen1988}.


\subsection{Cylindrical and parabolic bases}
\label{sec:ch382}
% TOC tag: [HOME]

\emph{Cylindrical bases} index modes by axial wavenumber \(k_{z}\) and azimuthal integer \(m\) in
\((\rho,\phi,z)\) coordinates---the natural language of waveguides, horns, and line sources
\cite{stratton1941,jackson1999,harrington2001}. \emph{Parabolic bases} use confocal parabolic coordinates suited to
focused reflectors, Gaussian-beam cavities, and long-path optical links \cite{stratton1941,jackson1999}. Subsection~\ref{sec:ch382}
is the HOME feed-and-focus coordinate layer: it does not replace VSH or plane-wave charts but supplies the representations
laboratory hardware often prefers before export to \(\mathcal{F}_{\mathrm{rad}}\) \cite{hansen1988,harrington2001}.

\paragraph{Standing imports.}
Vector Helmholtz separation Subsection~\ref{sec:ch332}, normalized VSH resolution
Eq.~\eqref{eq:ch3.normalized-coefficient-resolution}, and plane-wave chart Eq.~\eqref{eq:ch3.plane-wave-basis-resolution} enter by
reference \cite{stratton1941,hansen1988}. Subsection~\ref{sec:ch382} does not re-derive cylindrical Bessel functions or parabolic
wave equations \cite{jackson1999,stratton1941}.

\paragraph{Assumptions.}
Work on simply connected exterior regions or uniform cylindrical waveguides with declared boundary \(\Gamma_{b}\) when
discrete \(m\) quantization is invoked \cite{stratton1941,harrington2001}. Parabolic modes are used in the paraxial
validity regime unless full vector continuation is explicitly certified \cite{jackson1999,harrington2001}.

\paragraph{Cylindrical mode expansion (HOME).}
Let \(\mathbf{e}_{m}^{(c)}(\rho,\phi,z)\) denote cylindrical vector eigenmodes with azimuthal index \(m\in\mathbb{Z}\) and axial
wavenumber \(k_{z}\). Subsection~\ref{sec:ch382} writes
\begin{equation}
\label{eq:ch3.cylindrical-mode-expansion}
\mathbf{v}_\omega
=
\sum_{m\in\mathbb{Z}}
\int_{\mathbb{R}}
\widetilde{c}_{m}^{(c)}(k_{z})\,
\mathbf{e}_{m}^{(c)}(\rho,\phi,z;k_{z})\,dk_{z},
\end{equation}
with \(\widetilde{c}_{m}^{(c)}\) flux-normalized when restricted to \(\mathcal{F}_{\mathrm{rad}}\)
\cite{stratton1941,jackson1999,harrington2001}. Equation~\eqref{eq:ch3.cylindrical-mode-expansion} states the physical meaning of
cylindrical bases: \(m\) counts azimuthal phase winds around the \(z\)-axis, while \(k_{z}\) tracks axial propagation
\cite{jackson1999,harrington2001}. The same exported state generally requires many \(m\) values; a single \(m\) label is a
harmonic template, not a universal beam descriptor \cite{hansen1988,stratton1941}.

\paragraph{Parabolic beam basis (HOME).}
In parabolic coordinates \((\xi,\eta,\phi)\), Subsection~\ref{sec:ch382} adopts
\begin{equation}
\label{eq:ch3.parabolic-beam-basis}
\mathbf{v}_\omega
=
\sum_{\nu}
\int
\widetilde{c}_{\nu}^{(p)}(\beta)\,
\mathbf{e}_{\nu}^{(p)}(\xi,\eta,\phi;\beta)\,d\mu_{p}(\beta),
\end{equation}
with \(\mathbf{e}_{\nu}^{(p)}\) parabolic vector modes and \(\nu\) a separation index tied to focal geometry
\cite{stratton1941,jackson1999}. Equation~\eqref{eq:ch3.parabolic-beam-basis} captures focused-beam physics: energy is organized
around a focal axis and waist parameter, not around the origin-centered spheres of Eq.~\eqref{eq:ch3.normalized-coefficient-resolution}
\cite{harrington2001,jackson1999}. Parabolic coefficients are therefore the natural coordinates for reflector feeds and laser
resonators before far-zone projection \cite{stratton1941,hansen1988}.

\begin{figure}[t]
\centering
\ChOneFigBlock{%
\begin{tikzpicture}[ch1fig, node distance=7mm and 9mm]
  \node[ch1box] (cyl) {Eq.~\eqref{eq:ch3.cylindrical-mode-expansion}};
  \node[ch1box, right=10mm of cyl] (par) {Eq.~\eqref{eq:ch3.parabolic-beam-basis}};
  \node[ch1box, right=10mm of par] (rad) {$\mathcal{F}_{\mathrm{rad}}$};
  \draw[ch1flow] (cyl) -- (rad);
  \draw[ch1flow] (par) -- (rad);
\end{tikzpicture}%
}
\caption{Subsection~\ref{sec:ch382} cylindrical and parabolic charts both describe the same \(\mathcal{F}_{\mathrm{rad}}\) image with different
coordinate singularities and symmetries.}
\label{fig:ch3.382-dual-charts}
\end{figure}

\paragraph{Cylindrical--parabolic bridge (HOME).}
There exists a linear change of coordinates (not re-derived here) such that
\begin{equation}
\label{eq:ch3.cylindrical-parabolic-bridge}
\{\widetilde{c}_{m}^{(c)}(k_{z})\}
\;\longleftrightarrow\;
\{\widetilde{c}_{\nu}^{(p)}(\beta)\}
\;\longleftrightarrow\;
\{\widehat{c}_{\ell m}^{(\sigma)}\},
\end{equation}
with explicit transformation kernels deferred to Subsection~\ref{sec:ch383} \cite{hansen1988,stratton1941,jackson1999}. Equation
\eqref{eq:ch3.cylindrical-parabolic-bridge} is the Subsection~\ref{sec:ch382} unification claim: feed coordinates and focus coordinates are
representations, not distinct radiation mechanisms \cite{harrington2001,stratton1941}.

\paragraph{Physical meaning of azimuthal index \(m\) in cylindrical charts.}
Integer \(m\) in Eq.~\eqref{eq:ch3.cylindrical-mode-expansion} counts \(2\pi m\) phase advance around a cylindrical axis
\cite{jackson1999,stratton1941}. This \(m\) is not identical to the spherical \(m\) in Eq.~\eqref{eq:ch3.normalized-coefficient-resolution}
until Subsection~\ref{sec:ch383} supplies the transformation law \cite{hansen1988,harrington2001}. One must not equate cylindrical and
spherical \(m\) labels without the explicit matrix \cite{stratton1941,harrington2001}.

\paragraph{Worked illustration (circular waveguide TE\(_{m n}\)).}
A circular guide quantizes \(m\) and radial index \(n\) while leaving \(k_{z}\) discrete or continuous depending on length
\cite{jackson1999,harrington2001}. Subsection~\ref{sec:ch382} reads that hardware as a partial implementation of
Eq.~\eqref{eq:ch3.cylindrical-mode-expansion} before radiation export \cite{stratton1941,hansen1988}.

\paragraph{Problem (why cylindrical \(m\) is not spherical \(\ell\)).}
\textbf{Problem.} A horn datasheet lists ``\(m=1\)''. Explain why this is not automatically \(\ell=1\) in VSH
\cite{harrington2001,stratton1941,hansen1988}.

\textbf{Step 1.} Identify coordinate system used in the datasheet \cite{stratton1941}.

\textbf{Step 2.} Cite Eq.~\eqref{eq:ch3.cylindrical-parabolic-bridge} \cite{hansen1988}.

\textbf{Step 3.} Require Subsection~\ref{sec:ch383} transformation \cite{harrington2001}.

\textbf{Physical meaning.} Index labels are representation-dependent \cite{jackson1999,stratton1941}.

\begin{table}[t]
\centering
\caption{Subsection~\ref{sec:ch382} cylindrical versus parabolic comparison.}
\label{tab:ch3.382-chart-comparison}
\small
\begin{tabular}{@{}p{0.22\linewidth}p{0.34\linewidth}p{0.36\linewidth}@{}}
\toprule
Chart & Natural hardware & Dominant index \\
\midrule
Cylindrical & waveguides, horns & \(m,k_{z}\) \\
Parabolic & reflectors, cavities & \(\nu\), waist \\
Spherical & far-zone patterns & \(\ell,m,\sigma\) \\
\bottomrule
\end{tabular}
\end{table}

\paragraph{Parabolic waist and focal length vocabulary.}
\emph{Beam waist} \(w_{0}\) is the radius where paraxial intensity falls to \(1/e^{2}\) of axial value; \emph{focal length} \(f\)
locates the phase center of a focused beam \cite{jackson1999,harrington2001}. Subsection~\ref{sec:ch382} uses these as \emph{geometry parameters}
in Eq.~\eqref{eq:ch3.parabolic-beam-basis}, not as substitutes for \(\ell\) or \(m\) indices \cite{stratton1941,hansen1988}. Far-zone export
still decomposes into Eq.~\eqref{eq:ch3.normalized-coefficient-resolution} regardless of waist-based fitting \cite{harrington2001}.

\paragraph{Worked illustration (reflector feed to far zone).}
A parabolic reflector transforms a feed pattern in Eq.~\eqref{eq:ch3.parabolic-beam-basis} into a narrow far-zone cone
\cite{harrington2001,stratton1941,jackson1999}. The narrow cone may require large \(\ell\) in spherical coordinates even when the feed
used low parabolic indices \(\nu\) \cite{hansen1988,harrington2001}.

\paragraph{Problem (sparse parabolic, dense spherical).}
\textbf{Problem.} A Gaussian parabolic mode appears sparse in \(\nu\) but requires many \(\ell\) in VSH. Explain
\cite{hansen1988,jackson1999,harrington2001,stratton1941}.

\textbf{Step 1.} Recognize different coordinate singularities \cite{stratton1941}.

\textbf{Step 2.} Invoke Eq.~\eqref{eq:ch3.cylindrical-parabolic-bridge} \cite{hansen1988}.

\textbf{Physical meaning.} Sparsity is representation-dependent \cite{harrington2001,jackson1999}.

\begin{figure}[t]
\centering
\ChOneFigBlock{%
\begin{tikzpicture}[ch1fig, node distance=7mm and 9mm]
  \node[ch1box] (horn) {horn $m$};
  \node[ch1box, right=10mm of horn] (refl) {reflector $\nu$};
  \node[ch1box, right=10mm of refl] (far) {far-zone $\ell$};
  \draw[ch1flow] (horn) -- (far);
  \draw[ch1flow] (refl) -- (far);
\end{tikzpicture}%
}
\caption{Subsection~\ref{sec:ch382} feed charts compress description near sources; far-zone spherical labels generally remain dense.}
\label{fig:ch3.382-feed-to-far}
\end{figure}

\paragraph{Boundary quantization in cylindrical guides.}
When \(\Gamma_{b}\) is a conducting cylinder, \(m\) and radial \(n\) quantize as in Subsection~\ref{sec:ch372}
\cite{stratton1941,jackson1999,harrington2001}. Subsection~\ref{sec:ch382} treats that as boundary-induced discreteness in the
cylindrical chart, not as a change in \(\mathcal{F}_{\mathrm{rad}}\) definition \cite{hansen1988,harrington2001}.

\paragraph{Worked illustration (line source radiation).}
An infinite line current excites cylindrical modes with preferred \(m\) and continuous \(k_{z}\) spectrum
\cite{jackson1999,stratton1941}. Exported power integrates over \(k_{z}\) with the cylindrical measure before any spherical display
\cite{harrington2001,hansen1988}.

\paragraph{Problem (select chart for a horn datasheet).}
\textbf{Problem.} When should a report use Eq.~\eqref{eq:ch3.cylindrical-mode-expansion} instead of
Eq.~\eqref{eq:ch3.plane-wave-basis-resolution}? \cite{harrington2001,stratton1941,jackson1999}.

\textbf{Step 1.} Match chart to feed symmetry \cite{stratton1941}.

\textbf{Step 2.} Document transformation to far-zone VSH \cite{hansen1988}.

\textbf{Physical meaning.} Coordinates follow hardware symmetry \cite{harrington2001}.

\paragraph{Parabolic versus plane-wave near focus.}
Near a focus, plane-wave expansions Eq.~\eqref{eq:ch3.plane-wave-basis-resolution} are ill-conditioned because many
\((k_{x},k_{y})\) are required; parabolic modes Eq.~\eqref{eq:ch3.parabolic-beam-basis} remain economical
\cite{jackson1999,harrington2001,stratton1941}. Subsection~\ref{sec:ch382} recommends parabolic charts in focal regions and plane-wave
charts only after the far-zone condition is met \cite{hansen1988,harrington2001}.

\begin{figure}[t]
\centering
\ChOneFigBlock{%
\begin{tikzpicture}[ch1fig, node distance=7mm and 9mm]
  \node[ch1box] (near) {focal region};
  \node[ch1box, right=10mm of near] (par) {Eq.~\eqref{eq:ch3.parabolic-beam-basis}};
  \node[ch1box, right=10mm of par] (far) {Eq.~\eqref{eq:ch3.plane-wave-basis-resolution}};
  \draw[ch1flow] (near) -- (par) -- (far);
\end{tikzpicture}%
}
\caption{Subsection~\ref{sec:ch382} parabolic coordinates are economical near focus; plane-wave charts dominate in the far zone.}
\label{fig:ch3.382-focus-handoff}
\end{figure}

\begin{table}[t]
\centering
\caption{Subsection~\ref{sec:ch382} coordinate selection guide.}
\label{tab:ch3.382-selection-guide}
\small
\begin{tabular}{@{}p{0.30\linewidth}p{0.58\linewidth}@{}}
\toprule
Hardware context & Preferred chart \\
\midrule
Circular waveguide / horn & Eq.~\eqref{eq:ch3.cylindrical-mode-expansion} \\
Parabolic reflector / cavity & Eq.~\eqref{eq:ch3.parabolic-beam-basis} \\
Compact range far zone & Eq.~\eqref{eq:ch3.plane-wave-basis-resolution} \\
\bottomrule
\end{tabular}
\end{table}

\paragraph{Hankel--Fourier azimuthal spectrum.}
For axisymmetric feeds, the azimuthal dependence in Eq.~\eqref{eq:ch3.cylindrical-mode-expansion} reduces to a single integer \(m\), while the
radial profile enters through a Hankel transform at fixed \(k_{z}\) \cite{jackson1999,stratton1941,hansen1988}:
\begin{equation}
\label{eq:ch3.cylindrical-hankel-spectrum}
\widetilde{c}_{m}^{(c)}(k_{z})
=
\int_{0}^{\infty}
\widetilde{E}_{m}(\rho;k_{z})\,J_{m}(k_{t}\rho)\,k_{t}\,dk_{t},
\qquad
k_{t}=\sqrt{k_{0}^{2}-k_{z}^{2}},
\end{equation}
with \(\widetilde{E}_{m}\) the \(m\)-indexed transverse spectrum on the aperture plane and \(J_{m}\) the Bessel function of the first kind
\cite{stratton1941,jackson1999}. Equation~\eqref{eq:ch3.cylindrical-hankel-spectrum} is the cylindrical analogue of
Eq.~\eqref{eq:ch3.fourier-radiation-filter-chain}: it is meaningful only after \(\Pi_{\mathrm{rad}}\) and outgoing \(k_{z}\) branches are fixed
\cite{harrington2001,stratton1941}. Horn datasheets that quote a single ``mode \(m\)'' without the Hankel spectrum therefore omit the radial
coupling that generally admixes neighboring \(m\) values after truncation \cite{hansen1988,harrington2001}.

\paragraph{Gouy phase and focal-axis bookkeeping.}
Parabolic-beam charts carry a longitudinal phase correction absent from spherical VSH phases \cite{allen1992,jackson1999,harrington2001}.
Subsection~\ref{sec:ch382} records
\begin{equation}
\label{eq:ch3.parabolic-gouy-phase}
\Psi_{\mathrm{Gouy}}(z)
=
(m_{\mathrm{LG}}+1)\,
\arctan\!\left(\frac{z}{f}\right),
\qquad
f=\frac{\pi w_{0}^{2}}{\lambda},
\end{equation}
for dominant LG-like indices \(m_{\mathrm{LG}}\) near the waist \cite{allen1992,jackson1999}. Equation~\eqref{eq:ch3.parabolic-gouy-phase} is not a
gauge artifact: it tracks accumulated phase along the focal axis and must be archived when comparing near-focus interferograms to far-zone VSH phases
\cite{harrington2001,stratton1941}. Omitting \(\Psi_{\mathrm{Gouy}}\) when stitching parabolic coefficients to
Eq.~\eqref{eq:ch3.normalized-coefficient-resolution} produces false vortex counts in phase-unwrapping pipelines \cite{allen1992,hansen1988}.

\paragraph{Edge diffraction and azimuthal-index contamination.}
Finite apertures break the ideal \(e^{\jj m\phi}\) dependence assumed in Eq.~\eqref{eq:ch3.cylindrical-mode-expansion} \cite{stratton1941,harrington2001}.
Subsection~\ref{sec:ch382} bounds the leakage of power into neighboring \(m\) by the edge-to-wavelength ratio
\begin{equation}
\label{eq:ch3.cylindrical-edge-mixing}
\eta_{m\rightarrow m'}
\lesssim
\mathcal{O}\!\left(\left(\frac{\lambda}{D_{\mathrm{ap}}}\right)^{|m-m'|}\right),
\qquad
D_{\mathrm{ap}}>0,
\end{equation}
with \(D_{\mathrm{ap}}\) the physical aperture diameter \cite{stratton1941,harrington2001}.
importing the same symmetry-breaking taxonomy as Eq.~\eqref{eq:ch3.boundary-scattering-mixing-matrix} without reopening its proof
\cite{harrington2001,stratton1941,hansen1988}. Equation~\eqref{eq:ch3.cylindrical-edge-mixing} explains why a horn marketed as ``pure \(m=1\)'' can
still radiate measurable \(m=0\) and \(m=2\) content: the contaminant is geometry, not gauge freedom \cite{jackson1999,harrington2001}.

\paragraph{Archive metadata for cylindrical/parabolic campaigns.}
Store axis orientation, \(w_{0}\), \(f\), guide radius, boundary class \(\Gamma_{b}\), and the declared chart alongside coefficients
\cite{harrington2001,stratton1941,jackson1999}. Without that tuple, Subsection~\ref{sec:ch383} cannot reconstruct transformation matrices
\cite{hansen1988,harrington2001}.

\paragraph{Feed geometry and separable-chart vocabulary.}
Cylindrical and parabolic vocabularies arise from \emph{engineering separability}: feeds, reflectors, and cavities
often possess a preferred axis or focal line that makes one transverse coordinate natural in aperture or waist charts
\cite{stratton1941,harrington2001}. A \emph{cylindrical mode} is not a different physical wave; it is an eigenfunction
of the transverse Laplacian in \((\rho,\phi)\) charts linked to line-source and reflector symmetries
\cite{hansen1988,jackson1999}. A \emph{parabolic beam} names Gauss--Laguerre or parabolic-cylinder reductions valid in
paraxial domains; the implicit concept is \emph{controlled asymptotics about a focal axis}, not global spherical
completeness \cite{allen1992,harrington2001}. Equation~\eqref{eq:ch3.cylindrical-mode-expansion} embeds the same
normalization discipline as VSH coefficients: modal amplitudes are flux meters only after the cylindrical measure and
outgoing branch are fixed \cite{stratton1941,hansen1988}. Equation~\eqref{eq:ch3.parabolic-beam-basis} carries the
waist parameter \(w_{0}\), Rayleigh range, and paraxial indicator as part of its meaning; omitting them leaves a
formally similar expansion that is not the same radiative report \cite{allen1992,harrington2001}. The bridge
Eq.~\eqref{eq:ch3.cylindrical-parabolic-bridge} states the philosophical point directly: separable charts are
\emph{local dictionaries} glued to the spherical anchor of Subsection~\ref{sec:ch341}, and singularities on focal axes
are coordinate artifacts to be classified, not physical sources of power \cite{jackson1999,stratton1941,hansen1988}.

\paragraph{Waveguide hybrid-mode coupling on cylindrical charts.}
Finite feed transitions admix cylindrical modes that Eq.~\eqref{eq:ch3.cylindrical-mode-expansion} lists separately only in an ideal separable limit \cite{stratton1941,harrington2001,jackson1999}. Subsection~\ref{sec:ch382} models the admixture by
\begin{equation}
\label{eq:ch3.cylindrical-hybrid-coupling}
\begin{pmatrix}\widetilde{c}_{m}^{(c)}\\\widetilde{c}_{m'}^{(c)}\end{pmatrix}
=
\mathbf{H}_{m m'}(z)
\begin{pmatrix}\widetilde{c}_{m}^{(c)}(0)\\\widetilde{c}_{m'}^{(c)}(0)\end{pmatrix},
\end{equation}
with \(\mathbf{H}_{m m'}\) unitary on flux-normalized \(\mathcal{F}_{\mathrm{rad}}\) when both modes are outgoing-qualified \cite{poynting1884,harrington2001,stratton1941}. Mathematically, \(\mathbf{H}_{m m'}\) is a feed transfer matrix; physically, it explains why single-\(m\) horn headlines fail on measured patterns \cite{jackson1999,hansen1988}. Archive \(\mathbf{H}_{m m'}\) with Gouy phase Eq.~\eqref{eq:ch3.parabolic-gouy-phase} before parabolic--cylindrical comparisons \cite{allen1992,stratton1941}.

\paragraph{Reflector focal-line curvature and chart handoff.}
Parabolic charts compress focal regions; plane-wave charts dominate the far zone \cite{jackson1999,harrington2001}. Publish a handoff distance \(z_{\mathrm{handoff}}\): use Eq.~\eqref{eq:ch3.parabolic-beam-basis} for \(z<z_{\mathrm{handoff}}\) and Eq.~\eqref{eq:ch3.plane-wave-basis-resolution} beyond, linked by \(\mathbf{T}\) from Subsection~\ref{sec:ch383} \cite{hansen1988,stratton1941}. Omitting \(z_{\mathrm{handoff}}\) double-counts watts when both charts appear on one ledger \cite{harrington2001,devaney2004}.

\paragraph{Cylindrical Hankel spectrum and edge contamination (expanded).}
Eq.~\eqref{eq:ch3.cylindrical-hankel-spectrum} expresses azimuthal content as a Hankel/Fourier superposition in \(\rho\) \cite{stratton1941,jackson1999,hansen1988}.
Edge diffraction Eq.~\eqref{eq:ch3.cylindrical-edge-mixing} contaminates integer \(m\) labels when feed apertures truncate separable modes \cite{harrington2001,stratton1941}.
Subsection~\ref{sec:ch382} treats that contamination as a \emph{chart effect}: publish aperture diameter, guide radius, and \(\mathbf{H}_{m m'}\) from
Eq.~\eqref{eq:ch3.cylindrical-hybrid-coupling} before claiming single-\(m\) horn patterns \cite{jackson1999,hansen1988,devaney2004}. Gouy phase
Eq.~\eqref{eq:ch3.parabolic-gouy-phase} must accompany parabolic campaigns so phase-sensitive comparisons to cylindrical feeds are not misaligned
\cite{allen1992,jackson1999,harrington2001}.

\paragraph{Parabolic waist metrics and failure boundaries.}
Rayleigh range \(f=\pi w_{0}^{2}/\lambda\) and waist \(w_{0}\) enter Eq.~\eqref{eq:ch3.parabolic-beam-basis} as part of the coordinate definition, not as optional metadata
\cite{jackson1999,allen1992,harrington2001}. When \(f\sim w_{0}\), paraxial reductions fail and vector continuation to
Eq.~\eqref{eq:ch3.normalized-coefficient-resolution} is mandatory on \(\Gamma_R\) \cite{allen1992,hansen1988,stratton1941}. Feed designers should
archive \(w_{0}\), \(f\), and declared paraxial parameter \(\epsilon\) whenever parabolic charts are used in compact-range reports
\cite{harrington2001,jackson1999}.

\paragraph{Horn-to-far-zone transfer matrix (worked composition).}
A circular horn campaign often exports cylindrical coefficients \(\widetilde{c}_{m}^{(c)}(k_{z})\) on a finite \(k_{z}\) grid before any VSH report
\cite{stratton1941,jackson1999,harrington2001}. Subsection~\ref{sec:ch382} records the two-stage transfer
\begin{equation}
\label{eq:ch3.cylindrical-far-transfer}
\widehat{c}_{\ell m}^{(\sigma)}
=
\sum_{m'}\int
\bigl[\mathbf{T}_{\mathrm{VSH}\leftarrow\mathrm{cyl}}\bigr]_{\ell m\sigma;\,m'k_{z}}\,
\widetilde{c}_{m'}^{(c)}(k_{z})\,dk_{z},
\end{equation}
with \(\mathbf{T}_{\mathrm{VSH}\leftarrow\mathrm{cyl}}\) the certified block from Subsection~\ref{sec:ch383} rather than an informal fit
\cite{hansen1988,harrington2001,stratton1941}. Mathematically, Eq.~\eqref{eq:ch3.cylindrical-far-transfer} is a change of coordinates on
\(\mathcal{F}_{\mathrm{rad}}\); physically, it is how a feed datasheet becomes a spherical harmonic table on \(\Gamma_R\)
\cite{jackson1999,devaney2004}. Watt closure requires \(\mathcal{P}_{\mathrm{rad}}\) from Eq.~\eqref{eq:ch3.mode-power-from-coefficient} to match
before and after the integral when flux normalization is active \cite{poynting1884,harrington2001}.

\paragraph{Parabolic focus chart and equiphase surfaces.}
Confocal parabolic coordinates align with equiphase surfaces near a focal line \cite{stratton1941,jackson1999}. Eq.~\eqref{eq:ch3.parabolic-gouy-phase}
tracks the \(\pi/2\) phase slip between waist and far zone that cylindrical feeds do not display in the same variable
\cite{allen1992,jackson1999,harrington2001}. Subsection~\ref{sec:ch382} therefore pairs every parabolic coefficient file with Gouy metadata and
\(z_{\mathrm{handoff}}\) to the plane-wave chart Eq.~\eqref{eq:ch3.plane-wave-basis-resolution} \cite{hansen1988,stratton1941}. Omitting Gouy
phase when comparing horn and reflector campaigns rotates interference nulls without changing declared \(\mathcal{P}_{\mathrm{rad}}\)
\cite{harrington2001,jackson1999}.

\begin{figure}[t]
\centering
\ChOneFigBlock{%
\begin{tikzpicture}[ch1fig, node distance=7mm and 9mm]
  \node[ch1box] (horn) {cylindrical $\widetilde{c}_{m}^{(c)}$};
  \node[ch1box, right=10mm of horn] (T) {Eq.~\eqref{eq:ch3.cylindrical-far-transfer}};
  \node[ch1box, right=10mm of T] (vsh) {$\widehat{c}_{\ell m}^{(\sigma)}$};
  \node[ch1box, below=10mm of horn] (para) {Eq.~\eqref{eq:ch3.parabolic-beam-basis}};
  \draw[ch1flow] (horn) -- (T) -- (vsh);
  \draw[ch1flow] (para) -- (T);
\end{tikzpicture}%
}
\caption{Subsection~\ref{sec:ch382} feed transfer: cylindrical and parabolic charts converge on VSH export only through declared $\mathbf{T}$ blocks.}
\label{fig:ch3.382-feed-transfer}
\end{figure}

Subsection~\ref{sec:ch382} has introduced cylindrical and parabolic bases through
Eqs.~\eqref{eq:ch3.cylindrical-mode-expansion}--\eqref{eq:ch3.cylindrical-parabolic-bridge}
\cite{stratton1941,jackson1999,harrington2001}. Subsection~\ref{sec:ch383} constructs the explicit equivalence transformations
\cite{hansen1988}.


\subsection{Equivalence transformations}
\label{sec:ch383}
% TOC tag: [HOME][USE SS3.4][USE SS3.8.1][USE SS3.8.2]

An \emph{equivalence transformation} is a linear isomorphism between coordinate expansions of the same
\(\mathbf{v}_\omega\in\mathcal{F}_{\mathrm{rad}}\) that preserves the radiative inner product and watt budgets
\cite{stratton1941,hansen1988,harrington2001}. Subsection~\ref{sec:ch383} is the HOME matrix layer importing VSH definitions
Subsection~\ref{sec:ch341}, plane-wave resolution Eq.~\eqref{eq:ch3.plane-wave-basis-resolution}, and cylindrical/parabolic charts
Eqs.~\eqref{eq:ch3.cylindrical-mode-expansion}--\eqref{eq:ch3.parabolic-beam-basis} \cite{jackson1999,stratton1941}.

\paragraph{Standing imports.}
Normalized VSH resolution Eq.~\eqref{eq:ch3.normalized-coefficient-resolution}, mode power law
Eq.~\eqref{eq:ch3.mode-power-from-coefficient}, coefficient rotation law Eq.~\eqref{eq:ch3.coefficient-rotation-law},
plane-wave continuum bridge Eq.~\eqref{eq:ch3.plane-wave-continuum-bridge}, discrete--continuous bridge
Eq.~\eqref{eq:ch3.discrete-continuous-spectrum-bridge}, and VSH orthogonality Theorem~\ref{thm:ch3.vsh-orthogonality} enter by
reference \cite{hansen1988,stratton1941,harrington2001}. Subsection~\ref{sec:ch383} does not re-prove addition of angular momentum
\cite{jackson1999}.

\paragraph{Assumptions.}
All bases are expressed in the same outgoing gauge and Section~\ref{sec:ch35} normalization meters \cite{hansen1988,harrington2001}.
Transformation matrices are finite when truncated to \(\ell\le\ell_{\max}\) and band-limited in \((k_{x},k_{y})\) or \(k_{z}\)
\cite{stratton1941,harrington2001,kato1995}.

\paragraph{Basis-change matrix (HOME).}
Let \(\mathbf{c}^{(\mathrm{a})}\) and \(\mathbf{c}^{(\mathrm{b})}\) be coefficient vectors in bases \(\mathcal{B}_{a}\) and
\(\mathcal{B}_{b}\). Subsection~\ref{sec:ch383} defines
\begin{equation}
\label{eq:ch3.basis-change-matrix}
\mathbf{c}^{(\mathrm{b})}
=
\mathbf{T}_{b\leftarrow a}\,\mathbf{c}^{(\mathrm{a})},
\qquad
\mathbf{T}_{b\leftarrow a}\ \text{block-sparse under symmetry},
\end{equation}
with \(\mathbf{T}\) constructed from overlap integrals of normalized modes \cite{hansen1988,stratton1941,harrington2001}. Equation
\eqref{eq:ch3.basis-change-matrix} is the Subsection~\ref{sec:ch383} computational law: changing coordinates is matrix multiplication,
not a change in \(\mathcal{P}_{\mathrm{rad}}\) \cite{poynting1884,harrington2001}. Sparsity patterns of \(\mathbf{T}\) depend on which
symmetries the two charts share \cite{jackson1999,stratton1941}.

\paragraph{Power invariance under basis change (HOME).}
\begin{equation}
\label{eq:ch3.power-invariance-basis-change}
\mathcal{P}_{\mathrm{rad}}
=
\|\mathbf{c}^{(\mathrm{a})}\|_{2}^{2}
=
\|\mathbf{c}^{(\mathrm{b})}\|_{2}^{2}
\quad\text{when both bases are flux-normalized,}
\end{equation}
importing Eq.~\eqref{eq:ch3.mode-power-from-coefficient} and Theorem~\ref{thm:ch3.vsh-orthogonality} \cite{poynting1884,harrington2001,stratton1941}.
Equation~\eqref{eq:ch3.power-invariance-basis-change} is the watt closure test for any claimed transformation: if a matrix changes
coefficients but alters \(\mathcal{P}_{\mathrm{rad}}\), normalization or outgoing selection was violated \cite{hansen1988,harrington2001}.

\begin{theorem}[Basis-equivalence invariance (HOME)]
\label{thm:ch3.basis-equivalence-invariance}
For admissible \(\mathbf{v}_\omega\in\mathcal{F}_{\mathrm{rad}}\), expansions in VSH, plane-wave, cylindrical, and parabolic bases
describe the same state if and only if their coefficient vectors satisfy Eq.~\eqref{eq:ch3.basis-change-matrix} with
Eq.~\eqref{eq:ch3.power-invariance-basis-change} \cite{hansen1988,stratton1941,harrington2001,kato1995}.
\end{theorem}
\begin{proof}[Proof sketch]
Uniqueness of Riesz coefficients Eq.~\eqref{eq:ch3.normalized-coefficient-resolution} on each orthonormal basis and
Proposition~\ref{prop:ch3.continuous-parseval} on continuous charts \cite{stratton1941,hansen1988}.
% PROOF: AUTHOR REVIEW
\end{proof}

\begin{figure}[t]
\centering
\ChOneFigBlock{%
\begin{tikzpicture}[ch1fig, node distance=7mm and 9mm]
  \node[ch1box] (a) {$\mathcal{B}_{a}$};
  \node[ch1box, right=10mm of a] (T) {Eq.~\eqref{eq:ch3.basis-change-matrix}};
  \node[ch1box, right=10mm of T] (b) {$\mathcal{B}_{b}$};
  \node[ch1box, below=10mm of T] (P) {Eq.~\eqref{eq:ch3.power-invariance-basis-change}};
  \draw[ch1flow] (a) -- (T) -- (b);
  \draw[ch1flow] (P) -- (T);
\end{tikzpicture}%
}
\caption{Subsection~\ref{sec:ch383} equivalence: transformation matrices preserve watt budgets on normalized outgoing bases.}
\label{fig:ch3.383-matrix-chain}
\end{figure}

\paragraph{Physical meaning of transformation sparsity.}
When a horn is axisymmetric, \(\mathbf{T}\) from cylindrical to spherical blocks by \(m\) conservation modulo mixing induced by
finite \(\ell_{\max}\) \cite{harrington2001,stratton1941,hansen1988}. Breaking axisymmetry fills formerly zero blocks; that filling is
physics from Subsection~\ref{sec:ch363}, not an artifact of ``bad'' coordinates \cite{harrington2001,jackson1999}. Subsection~\ref{sec:ch383}
therefore reports block structure as diagnostic of symmetry, not merely as a computational convenience \cite{stratton1941,hansen1988}.

\paragraph{Worked illustration (plane-wave to VSH on broadside).}
A narrow \(k_{x},k_{y}\) patch around origin maps predominantly to low \(\ell\) in
Eq.~\eqref{eq:ch3.normalized-coefficient-resolution} \cite{hansen1988,jackson1999}. The \(\mathbf{T}\) matrix is dense in general but
may be truncated for error bounds from Subsection~\ref{sec:ch362} \cite{harrington2001,stratton1941}.

\paragraph{Problem (verify power invariance numerically).}
\textbf{Problem.} Given \(\mathbf{c}^{(\mathrm{pw})}\) and \(\mathbf{T}_{\mathrm{sph}\leftarrow\mathrm{pw}}\), verify
Eq.~\eqref{eq:ch3.power-invariance-basis-change} \cite{harrington2001,hansen1988,stratton1941}.

\textbf{Step 1.} Compute \(\mathbf{c}^{(\mathrm{sph})}=\mathbf{T}\mathbf{c}^{(\mathrm{pw})}\) \cite{hansen1988}.

\textbf{Step 2.} Sum squared magnitudes in each basis \cite{poynting1884}.

\textbf{Step 3.} Tolerance-compare totals \cite{harrington2001}.

\textbf{Physical meaning.} Correct \(\mathbf{T}\) preserves quadratic power functionals \cite{stratton1941}.

\begin{table}[t]
\centering
\caption{Subsection~\ref{sec:ch383} transformation objects.}
\label{tab:ch3.383-transform-objects}
\small
\begin{tabular}{@{}p{0.26\linewidth}p{0.38\linewidth}p{0.28\linewidth}@{}}
\toprule
Map & Equation & Preserves \\
\midrule
General & Eq.~\eqref{eq:ch3.basis-change-matrix} & state \\
Power & Eq.~\eqref{eq:ch3.power-invariance-basis-change} & watts \\
Rotation & Eq.~\eqref{eq:ch3.coefficient-rotation-law} & \(\mathcal{P}_{\mathrm{rad}}\) \\
\bottomrule
\end{tabular}
\end{table}

\paragraph{Composition law for chained transforms.}
If \(\mathcal{B}_{a}\to\mathcal{B}_{b}\) and \(\mathcal{B}_{b}\to\mathcal{B}_{c}\), then
\(\mathbf{T}_{c\leftarrow a}=\mathbf{T}_{c\leftarrow b}\mathbf{T}_{b\leftarrow a}\) \cite{hansen1988,stratton1941}. Subsection~\ref{sec:ch383}
uses composition to pipeline horn cylindrical modes \(\to\) parabolic focus modes \(\to\) far-zone VSH without ad hoc rescalings
\cite{harrington2001,jackson1999}. Each stage must preserve Eq.~\eqref{eq:ch3.power-invariance-basis-change} \cite{poynting1884,harrington2001}.

\paragraph{Worked illustration (rotation as unitary block).}
Laboratory rotation acts as Wigner blocks on VSH coefficients Eq.~\eqref{eq:ch3.coefficient-rotation-law} while mixing plane-wave
phases linearly \cite{jackson1999,hansen1988}. Subsection~\ref{sec:ch383} treats rotation as a symmetry-certified \(\mathbf{T}\), not as a
new physical excitation \cite{stratton1941,harrington2001}.

\paragraph{Problem (compose horn-to-far pipeline).}
\textbf{Problem.} Write \(\mathbf{T}_{\mathrm{sph}\leftarrow\mathrm{cyl}}\) as composition through parabolic focus chart
\cite{harrington2001,stratton1941,hansen1988}.

\textbf{Step 1.} \(\mathbf{T}_{\mathrm{par}\leftarrow\mathrm{cyl}}\) from Eq.~\eqref{eq:ch3.cylindrical-parabolic-bridge} \cite{hansen1988}.

\textbf{Step 2.} \(\mathbf{T}_{\mathrm{sph}\leftarrow\mathrm{par}}\) from far-zone asymptotics \cite{stratton1941}.

\textbf{Physical meaning.} Hardware pipelines are matrix products \cite{harrington2001}.

\begin{figure}[t]
\centering
\ChOneFigBlock{%
\begin{tikzpicture}[ch1fig, node distance=7mm and 9mm]
  \node[ch1box] (cyl) {cyl};
  \node[ch1box, right=8mm of cyl] (par) {par};
  \node[ch1box, right=8mm of par] (pw) {pw};
  \node[ch1box, right=8mm of pw] (sph) {VSH};
  \draw[ch1flow] (cyl) -- (par) -- (pw) -- (sph);
\end{tikzpicture}%
}
\caption{Subsection~\ref{sec:ch383} composition pipeline: chained \(\mathbf{T}\) matrices from feed charts to spherical far-zone reports.}
\label{fig:ch3.383-composition-pipeline}
\end{figure}

\paragraph{Truncation error under \(\mathbf{T}\).}
Finite \(\ell_{\max}\) and band-limited \((k_{x},k_{y})\) induce truncation remainder Eq.~\eqref{eq:ch3.truncation-remainder-energy}
after transformation \cite{harrington2001,hansen1988}. Subsection~\ref{sec:ch383} bounds reporting error by the larger of source truncation and
\(\mathbf{T}\) tail norms imported from Subsection~\ref{sec:ch362} \cite{stratton1941,harrington2001}.

\paragraph{Worked illustration (ill-conditioned \(\mathbf{T}\)).}
Near caustics, parabolic-to-plane-wave maps become ill-conditioned: many plane waves are needed \cite{jackson1999,harrington2001}. The
condition number of \(\mathbf{T}\) signals when to remain in parabolic coordinates per Subsection~\ref{sec:ch382} \cite{stratton1941,hansen1988}.

\paragraph{Problem (audit a published equivalence claim).}
\textbf{Problem.} A paper claims ``plane-wave spectrum equals VSH spectrum'' without matrix. Audit
\cite{harrington2001,stratton1941,hansen1988}.

\textbf{Step 1.} Demand explicit \(\mathbf{T}\) \cite{hansen1988}.

\textbf{Step 2.} Test Eq.~\eqref{eq:ch3.power-invariance-basis-change} \cite{poynting1884}.

\textbf{Physical meaning.} Verbal equivalence is insufficient \cite{stratton1941,harrington2001}.

\paragraph{Adjoint transforms and reciprocity.}
When \(\mathbf{T}\) is unitary on \(\mathcal{F}_{\mathrm{rad}}\), the adjoint equals the inverse for power purposes
\cite{stratton1941,harrington2001}. Reciprocity relations from Subsection~\ref{sec:ch282} pair source and observation coefficients through
\(\mathbf{T}^{\mathsf{H}}\) in dual charts \cite{stratton1941,hansen1988}. Subsection~\ref{sec:ch383} records that pairing without re-deriving
reciprocity \cite{harrington2001}.

\begin{figure}[t]
\centering
\ChOneFigBlock{%
\begin{tikzpicture}[ch1fig, node distance=7mm and 9mm]
  \node[ch1box] (T) {$\mathbf{T}$};
  \node[ch1box, right=10mm of T] (uni) {unitary on $\mathcal{F}_{\mathrm{rad}}$};
  \node[ch1box, right=10mm of uni] (adj) {$\mathbf{T}^{\mathsf{H}}=\mathbf{T}^{-1}$};
  \draw[ch1flow] (T) -- (uni) -- (adj);
\end{tikzpicture}%
}
\caption{Subsection~\ref{sec:ch383} unitary \(\mathbf{T}\): adjoint maps implement reciprocal observation charts.}
\label{fig:ch3.383-unitary-adjoint}
\end{figure}

\begin{table}[t]
\centering
\caption{Subsection~\ref{sec:ch383} transformation audit checklist.}
\label{tab:ch3.383-audit-checklist}
\small
\begin{tabular}{@{}p{0.08\linewidth}p{0.50\linewidth}p{0.34\linewidth}@{}}
\toprule
Step & Test & Pass criterion \\
\midrule
1 & Normalization & Section~\ref{sec:ch35} meters \\
2 & Matrix & Eq.~\eqref{eq:ch3.basis-change-matrix} \\
3 & Power & Eq.~\eqref{eq:ch3.power-invariance-basis-change} \\
4 & Truncation & Subsection~\ref{sec:ch362} bound \\
\bottomrule
\end{tabular}
\end{table}

\paragraph{Truncation error propagation through \(\mathbf{T}\).}
Let \(\mathbf{c}^{(\mathrm{a})}_{N}\) denote a truncated coefficient vector and \(\delta\mathbf{c}^{(\mathrm{a})}\) the omitted tail from
Subsection~\ref{sec:ch362} \cite{harrington2001,stratton1941}. Subsection~\ref{sec:ch383} bounds the exported error in chart \(\mathcal{B}_{b}\) by
\begin{equation}
\label{eq:ch3.basis-change-error-propagation}
\|\delta\mathbf{c}^{(\mathrm{b})}\|_{2}
\le
\|\mathbf{T}_{b\leftarrow a}\|_{2}\,
\|\delta\mathbf{c}^{(\mathrm{a})}\|_{2},
\qquad
\delta\mathbf{c}^{(\mathrm{b})}=\mathbf{T}_{b\leftarrow a}\,\delta\mathbf{c}^{(\mathrm{a})},
\end{equation}
with \(\|\cdot\|_{2}\) the flux-normalized coefficient norm \cite{hansen1988,kato1995,harrington2001}. Equation~\eqref{eq:ch3.basis-change-error-propagation}
states that an ill-conditioned \(\mathbf{T}\) \emph{amplifies} truncation noise: a sparse VSH list can become a dense plane-wave spectrum not because
physics changed, but because the transformation matrix spreads omitted modes \cite{jackson1999,harrington2001}. Campaigns should therefore publish
\(\kappa(\mathbf{T})=\|\mathbf{T}\|_{2}\|\mathbf{T}^{-1}\|_{2}\) alongside \(\ell_{\max}\) when comparing charts \cite{stratton1941,hansen1988}.

\paragraph{Wigner-block composition on rotated pipelines.}
When a laboratory rotation \(R(g)\) precedes a chart change, Subsection~\ref{sec:ch383} implements
\begin{equation}
\label{eq:ch3.basis-change-rotation-compose}
\mathbf{T}_{b\leftarrow a}^{\mathrm{(rot)}}
=
\mathbf{T}_{b\leftarrow a}\,
\mathbf{D}^{(\ell)}(g),
\end{equation}
with \(\mathbf{D}^{(\ell)}(g)\) from Eq.~\eqref{eq:ch3.coefficient-rotation-law}, block-diagonal in \(\ell\) under \(\mathrm{SO}(3)\)
\cite{jackson1999,hansen1988,stratton1941}. Equation~\eqref{eq:ch3.basis-change-rotation-compose}
separates symmetry action from dictionary transport: rotations rephase sectors within a fixed basis; \(\mathbf{T}\) changes the basis itself
\cite{harrington2001,jackson1999}. Confusing the two produces sign errors in reciprocity checks that compare rotated feeds to static transformation tables
\cite{stratton1941,hansen1988}.

\begin{figure}[t]
\centering
\ChOneFigBlock{%
\begin{tikzpicture}[ch1fig, node distance=7mm and 9mm]
  \node[ch1box] (tail) {truncation tail $\delta\mathbf{c}^{(a)}$};
  \node[ch1box, right=10mm of tail] (T) {Eq.~\eqref{eq:ch3.basis-change-error-propagation}};
  \node[ch1box, right=10mm of T] (b) {$\delta\mathbf{c}^{(b)}$};
  \node[ch1box, below=10mm of T] (D) {Eq.~\eqref{eq:ch3.basis-change-rotation-compose}};
  \draw[ch1flow] (tail) -- (T) -- (b);
  \draw[ch1flow] (D) -- (T);
\end{tikzpicture}%
}
\caption{Subsection~\ref{sec:ch383} error and symmetry pipeline: \(\mathbf{T}\) propagates truncation; \(\mathbf{D}^{(\ell)}\) handles rotation within a chart.}
\label{fig:ch3.383-error-rotation}
\end{figure}

\paragraph{Peer-review metadata for transformation claims.}
Publish \(\mathbf{T}\) sparsity pattern, normalization convention, \(\ell_{\max}\), band limits, and power test residuals alongside
coefficient vectors \cite{harrington2001,stratton1941,hansen1988}. Subsection~\ref{sec:ch383} treats missing \(\mathbf{T}\) as a non-reproducible
basis claim \cite{harrington2001,jackson1999}.

\paragraph{Representation relativity and power invariance.}
A \emph{basis change} is a change of coordinates on coefficient space that leaves the radiative equivalence class
invariant when executed with a certified transformation matrix \cite{hansen1988,harrington2001,stratton1941}. The
implicit concept behind Eq.~\eqref{eq:ch3.basis-change-matrix} is \emph{dictionary transport}: each row specifies how
one complete coordinate system expresses a single physical direction or mode label in another \cite{jackson1999,hansen1988}.
Power invariance Eq.~\eqref{eq:ch3.power-invariance-basis-change} embeds unitary (or certified isometric) structure as
a physical statement: watts computed from spherical, plane-wave, or cylindrical coefficients agree because the inner
product on \(\mathcal{F}_{\mathrm{rad}}\) is the invariant object, not any particular sparse chart
\cite{poynting1884,harrington2001,stratton1941}. Theorem~\ref{thm:ch3.basis-equivalence-invariance} therefore belongs
to the same epistemic class as gauge-invariant energy bookkeeping: a mismatch in reported power after a declared
basis change signals a matrix error, a missing evanescent null, or an incompatible normalization tuple---not a new
degree of freedom \cite{harrington2001,jackson1999}. Scientific discourse that treats ``transform to plane waves'' as a
mere plotting step forgets that \(\mathbf{T}\) is part of the measurement protocol; without it, coefficients are
non-reproducible symbols \cite{harrington2001,devaney2004}.

\paragraph{Reciprocity-adjoint pairing of transformation matrices.}
Reciprocity imports from Subsection~\ref{sec:ch282} constrain certified basis changes \cite{stratton1941,harrington2001}. For unitary \(\mathbf{T}_{b\leftarrow a}\),
\begin{equation}
\label{eq:ch3.basis-change-reciprocity-adjoint}
\mathbf{T}_{b\leftarrow a}^{-1}
=
\mathbf{T}_{a\leftarrow b}
=
\mathbf{T}_{b\leftarrow a}^{\mathsf{H}},
\end{equation}
and inner products agree on both charts when flux normalization is active \cite{hansen1988,poynting1884,stratton1941}. Equation~\eqref{eq:ch3.basis-change-reciprocity-adjoint} is the transmit/receive audit after chart conversion: reciprocity failures often trace to \(\kappa(\mathbf{T})\gg1\) in Eq.~\eqref{eq:ch3.basis-change-error-propagation}, not to new physics \cite{harrington2001,jackson1999,kato1995}.

\paragraph{Sparse-block structure under shared symmetry axes.}
When \(\mathcal{B}_{a}\) and \(\mathcal{B}_{b}\) share a rotation axis, \(\mathbf{T}_{b\leftarrow a}\) block-diagonalizes in \(\ell\) up to truncation \cite{jackson1999,hansen1988}. Publish the sparsity mask with \(\mathbf{T}\); unexpected density signals mislabeled axes or missing evanescent nulls \cite{stratton1941,harrington2001}.

\paragraph{Condition-number audit on published $\mathbf{T}$ blocks.}
Ill-conditioned basis changes amplify coefficient noise without changing \(\mathcal{P}_{\mathrm{rad}}\) when Eq.~\eqref{eq:ch3.power-invariance-basis-change} holds exactly
\cite{hansen1988,harrington2001,kato1995}. Subsection~\ref{sec:ch383} recommends archiving
\begin{equation}
\label{eq:ch3.basis-change-condition-audit}
\kappa(\mathbf{T}_{b\leftarrow a})
=
\|\mathbf{T}_{b\leftarrow a}\|\,\|\mathbf{T}_{b\leftarrow a}^{-1}\|,
\end{equation}
with declared \(\kappa_{\max}\) in campaign metadata, alongside error propagation Eq.~\eqref{eq:ch3.basis-change-error-propagation} \cite{harrington2001,devaney2004}. Mathematically, large \(\kappa\) explains why sparse cylindrical headlines disagree with dense VSH tables for the same \(\mathbf{v}_\omega\)
\cite{jackson1999,stratton1941}; physically, it is an inversion warning before Chapter~\ref{ch:04} realizability arguments begin
\cite{devaney2004,harrington2001}. Figure~\ref{fig:ch3.383-error-rotation} remains the geometric reminder that rotations are unitary while generic feed transforms need not be
\cite{hansen1988,jackson1999}.

\paragraph{Chained pipelines (horn $\to$ plane-wave $\to$ VSH).}
Composition Eq.~\eqref{eq:ch3.basis-change-rotation-compose} extends to feed pipelines: the product
\(\mathbf{T}_{\mathrm{VSH}\leftarrow\mathrm{pw}}\,\mathbf{T}_{\mathrm{pw}\leftarrow\mathrm{cyl}}\) must be published as one certified map, not as two undocumented plotting steps
\cite{harrington2001,stratton1941,hansen1988}. Subsection~\ref{sec:ch383} treats missing intermediate normalization as a power-invariance failure even when endpoint patterns look similar
\cite{poynting1884,harrington2001,jackson1999}.

\paragraph{Unitary versus isometric transforms in laboratory charts.}
Not every practical \(\mathbf{T}_{b\leftarrow a}\) is unitary on truncated spaces: evanescent nulls and band-limits yield isometric embeddings with nontrivial \(\kappa(\mathbf{T})\)
\cite{kato1995,hansen1988,harrington2001}. Subsection~\ref{sec:ch383} therefore distinguishes \emph{power invariance} (Eq.~\eqref{eq:ch3.power-invariance-basis-change}) from \emph{coefficient stability}: watts may close while individual \(\widehat{c}\) entries amplify by \(\kappa\)
\cite{devaney2004,stratton1941}. Campaigns must publish both tests when comparing sparse horn charts to dense VSH tables on the same \(\mathcal{P}_{\mathrm{rad}}\)
\cite{harrington2001,jackson1999}.

\paragraph{Adjoint verification on reciprocal campaigns.}
Reciprocal antennas require Eq.~\eqref{eq:ch3.basis-change-reciprocity-adjoint} on transmit and receive paths \cite{stratton1941,harrington2001}. Subsection~\ref{sec:ch383} recommends logging \(\|\mathbf{T}\mathbf{T}^{\mathsf{H}}-\mathbf{I}\|\) and \(\|\mathbf{T}^{\mathsf{H}}\mathbf{T}-\mathbf{I}\|\) on active blocks as a pre-publication screen alongside Eq.~\eqref{eq:ch3.basis-change-condition-audit}
\cite{hansen1988,poynting1884,devaney2004}. Failure implicates normalization tuple mismatch, not new coupling physics
\cite{harrington2001,stratton1941}.

\paragraph{Broadside calibration as a $\mathbf{T}$ acceptance test.}
A broadside-dominated \(\mathbf{v}_\omega\) concentrates flux near \(k_{x}=k_{y}=0\) in Eq.~\eqref{eq:ch3.plane-wave-basis-resolution} and near \(\ell=0,1\) in Eq.~\eqref{eq:ch3.normalized-coefficient-resolution}
\cite{jackson1999,hansen1988,harrington2001}. Subsection~\ref{sec:ch383} uses that limit as a calibration instance: \(\mathbf{T}_{\mathrm{VSH}\leftarrow\mathrm{pw}}\) should map a single dominant plane-wave cell to low-\(\ell\) VSH weights with closure error below the campaign threshold in Eq.~\eqref{eq:ch3.basis-change-error-propagation}
\cite{stratton1941,devaney2004}. Failure at broadside signals Jacobian omission Eq.~\eqref{eq:ch3.plane-wave-direction-jacobian} or evanescent leakage, not genuine high-\(\ell\) physics
\cite{harrington2001,jackson1999}.

\paragraph{Error budget composition for chained transforms.}
When \(\mathbf{T}=\mathbf{T}_{2}\mathbf{T}_{1}\), error propagation satisfies \(\|\delta\mathbf{c}_{3}\|\le\|\mathbf{T}_{2}\|\,\|\delta\mathbf{c}_{2}\|+\|\delta\mathbf{T}_{2}\|\,\|\mathbf{c}_{2}\|\) in the usual operator norm \cite{kato1995,harrington2001}. Subsection~\ref{sec:ch383} archives each factor's \(\kappa\) from Eq.~\eqref{eq:ch3.basis-change-condition-audit} because chained feed pipelines multiply condition numbers
\cite{devaney2004,hansen1988,stratton1941}. Watt closure can still pass while coefficient vectors become uninterpretable---a state Chapter~\ref{ch:04} inversion must not inherit uncertified
\cite{harrington2001,jackson1999}.

\paragraph{Evanescent null as a $\mathbf{T}$ precondition.}
Basis-change matrices are defined on outgoing \(\mathcal{F}_{\mathrm{rad}}\); evanescent sheets must be annihilated before \(\mathbf{T}_{b\leftarrow a}\) is applied
\cite{stratton1941,harrington2001,jackson1999}. Subsection~\ref{sec:ch383} treats reactive FFT content as a precondition violation: watts may appear to close while \(\kappa(\mathbf{T})\) explodes because near-field energy was folded into propagating coefficients
\cite{hansen1988,devaney2004,kato1995}. The corrective order repeats Eq.~\eqref{eq:ch3.fourier-radiation-filter-chain}: \(\Pi_{\mathrm{rad}}\) first, normalization second, \(\mathbf{T}\) third
\cite{harrington2001,stratton1941}.

\paragraph{Rotation as the unitary reference transform.}
Wigner blocks Eq.~\eqref{eq:ch3.basis-change-rotation-compose} are the unitary reference against which feed transforms should be compared: rotations preserve \(\kappa=1\) on active subspaces
\cite{jackson1999,hansen1988}. Subsection~\ref{sec:ch383} uses rotation as a calibration sanity check---if a declared \(\mathbf{T}\) deviates strongly from block-Wigner form without documented symmetry breaking, the matrix is suspect
\cite{harrington2001,stratton1941,devaney2004}.

Subsection~\ref{sec:ch383} has constructed equivalence transformations through
Eqs.~\eqref{eq:ch3.basis-change-matrix}--\eqref{eq:ch3.power-invariance-basis-change} and
Theorem~\ref{thm:ch3.basis-equivalence-invariance} \cite{hansen1988,stratton1941,harrington2001}. Subsection~\ref{sec:ch384} analyzes how OAM
labels depend on representation choices \cite{allen1992,belinfante1940}.


\subsection{Representation dependence of OAM modes}
\label{sec:ch384}
% TOC tag: [HOME][USE SS1.4.5]

\emph{Representation dependence} means a single radiative state carries different coefficient sparsity patterns, phase winds, and informal
``OAM'' labels in plane-wave, cylindrical, parabolic, and spherical charts \cite{allen1992,belinfante1940,harrington2001}.
\emph{OAM modes} in popular discourse usually refer to helical phase fronts with integer azimuthal index; in this monograph they are
\emph{coordinates} on \(\mathcal{F}_{\mathrm{rad}}\) subject to gauge and spin-orbit limits from Subsection~\ref{sec:ch01.4.5}
\cite{belinfante1940,stratton1941,allen1992}. Principle~\ref{prin:ch3.osr-separation} and
Principle~\ref{prin:ch1.rep-physics} forbid substituting coordinate labels for operator export or measurement
reports. Subsection~\ref{sec:ch384} is the HOME caution layer connecting basis transformations of
Subsection~\ref{sec:ch383} to topology gates of Subsection~\ref{sec:ch374} \cite{harrington2001,hansen1988}.

\paragraph{Standing imports.}
Local spin-orbit split Eq.~\eqref{eq:ch1.local-spin-orbit-split}, equivalence class
Eq.~\eqref{eq:ch1.spin-orbit-equivalence-class}, gauge redistribution Eq.~\eqref{eq:ch1.spin-orbit-gauge-redistribution},
phase-topology link Eq.~\eqref{eq:ch3.phase-singularity-angular-link}, integer index discipline
Eq.~\eqref{eq:ch3.integer-ell-def}, basis-change matrix Eq.~\eqref{eq:ch3.basis-change-matrix}, and beam-order warning
Eq.~\eqref{eq:ch3.beam-order-vs-ell-warning} enter by reference \cite{belinfante1940,allen1992,harrington2001,stratton1941}. Subsection~\ref{sec:ch384}
does not re-derive Belinfante symmetrization \cite{belinfante1940}.

\paragraph{Assumptions.}
OAM claims are evaluated on transverse far-zone fields after \(\Pi_{\mathrm{rad}}\) selection and decomposition into declared bases
\cite{harrington2001,stratton1941}. Paraxial LG templates are cited only within paraxial validity; full-vector claims require VSH coefficients
\cite{allen1992,jackson1999,harrington2001}.

\paragraph{OAM representation-dependence law (HOME).}
Let \(\mathcal{B}\) denote any of the charts in Section~\ref{sec:ch38}. Subsection~\ref{sec:ch384} records
\begin{equation}
\label{eq:ch3.OAM-representation-dependence}
\text{``OAM label'' in }\mathcal{B}
\;\neq\;
\text{``OAM label'' in }\mathcal{B}'
\quad\text{unless }\mathbf{T}_{\mathcal{B}'\leftarrow\mathcal{B}}\ \text{preserves the labeled sector},
\end{equation}
importing Eq.~\eqref{eq:ch3.basis-change-matrix} and spin-orbit equivalence
Eq.~\eqref{eq:ch1.spin-orbit-equivalence-class} \cite{belinfante1940,allen1992,harrington2001}. Equation
\eqref{eq:ch3.OAM-representation-dependence} is the Subsection~\ref{sec:ch384} anti-essentialism law: OAM is not a representation-independent
quantum number for arbitrary beams \cite{harrington2001,stratton1941,belinfante1940}.

\paragraph{Paraxial versus spherical OAM (HOME).}
\begin{equation}
\label{eq:ch3.paraxial-vs-spherical-OAM}
m_{\mathrm{LG}}
\;\Longleftrightarrow\;
\{m\ \text{sectors in Eq.~\eqref{eq:ch3.integer-ell-def}}\}
\quad\text{only after }\mathbf{T}_{\mathrm{sph}\leftarrow\mathrm{LG}}\ \text{and gauge screen},
\end{equation}
with \(m_{\mathrm{LG}}\) the Laguerre--Gaussian azimuthal index \cite{allen1992,jackson1999}. Equation
\eqref{eq:ch3.paraxial-vs-spherical-OAM} warns that paraxial vortex labels do not transfer verbatim to VSH labels: full-vector
decomposition generally activates multiple \(\ell\) sectors \cite{harrington2001,hansen1988,stratton1941}.

\paragraph{Gauge screen for OAM claims (HOME).}
\begin{equation}
\label{eq:ch3.gauge-screen-OAM-claim}
\text{OAM transport claim allowed}
\;\Longrightarrow\;
\text{Eq.~\eqref{eq:ch1.topology-qualified-gate}}
\ \wedge\
\text{Eq.~\eqref{eq:ch3.power-invariance-basis-change}}
\ \wedge\
\text{declared }\mathcal{B},
\end{equation}
importing Subsection~\ref{sec:ch01.4.5} spin-orbit limits \cite{belinfante1940,harrington2001,stratton1941,allen1992}. Equation
\eqref{eq:ch3.gauge-screen-OAM-claim} is the Subsection~\ref{sec:ch384} publication gate: topology, power invariance, and basis declaration
must accompany any OAM headline \cite{harrington2001,belinfante1940}.

\begin{figure}[t]
\centering
\ChOneFigBlock{%
\begin{tikzpicture}[ch1fig, node distance=7mm and 9mm]
  \node[ch1box] (B) {chart $\mathcal{B}$};
  \node[ch1box, right=10mm of B] (T) {Eq.~\eqref{eq:ch3.basis-change-matrix}};
  \node[ch1box, right=10mm of T] (Bp) {chart $\mathcal{B}'$};
  \node[ch1box, below=10mm of T] (dep) {Eq.~\eqref{eq:ch3.OAM-representation-dependence}};
  \draw[ch1flow] (B) -- (T) -- (Bp);
  \draw[ch1flow] (dep) -- (T);
\end{tikzpicture}%
}
\caption{Subsection~\ref{sec:ch384} OAM labels change under basis change unless symmetry preserves the labeled sector.}
\label{fig:ch3.384-OAM-basis-shift}
\end{figure}

\paragraph{Physical meaning: OAM is not a single number.}
A beam may display \(m_{\mathrm{LG}}=1\) phase wind in a transverse plane yet decompose into \(\ell\ge1\) VSH sectors with multiple
\(m\) values \cite{allen1992,hansen1988,harrington2001}. Subsection~\ref{sec:ch384} therefore rejects scalar OAM slogans and requires integer
block partitions Table~\ref{tab:ch3.373-block-ledger} style data from Subsection~\ref{sec:ch373} \cite{harrington2001,stratton1941}.

\paragraph{Worked illustration (LG beam far-zone spread).}
Paraxial LG\(_{0}^{1}\) concentrates near axis; its far-zone VSH spectrum is not single-\(\ell\) \cite{allen1992,jackson1999,harrington2001}.
Eq.~\eqref{eq:ch3.paraxial-vs-spherical-OAM} documents the spread before any transport claim \cite{hansen1988,stratton1941}.

\paragraph{Problem (audit ``pure OAM beam'').}
\textbf{Problem.} A paper claims a ``pure \(\ell=1\) OAM beam'' from a spiral phase plate. Audit with
Eqs.~\eqref{eq:ch3.gauge-screen-OAM-claim} and \eqref{eq:ch3.beam-order-vs-ell-warning}
\cite{allen1992,belinfante1940,harrington2001,stratton1941}.

\textbf{Step 1.} Decompose far field into Eq.~\eqref{eq:ch3.normalized-coefficient-resolution} \cite{hansen1988}.

\textbf{Step 2.} Report full coefficient vector, not a slogan \cite{harrington2001}.

\textbf{Step 3.} Apply topology gate Eq.~\eqref{eq:ch1.topology-qualified-gate} \cite{belinfante1940}.

\textbf{Physical meaning.} Phase plates do not guarantee pure multipole sectors \cite{stratton1941,allen1992}.

\begin{table}[t]
\centering
\caption{Subsection~\ref{sec:ch384} OAM label sources and certificates.}
\label{tab:ch3.384-OAM-sources}
\small
\begin{tabular}{@{}p{0.22\linewidth}p{0.36\linewidth}p{0.34\linewidth}@{}}
\toprule
Label source & Index & Certificate \\
\midrule
LG paraxial & \(m_{\mathrm{LG}}\) & Eq.~\eqref{eq:ch3.paraxial-vs-spherical-OAM} \\
VSH spherical & \(m\) in Eq.~\eqref{eq:ch3.integer-ell-def} & decomposition \\
Cylindrical & \(m\) in Eq.~\eqref{eq:ch3.cylindrical-mode-expansion} & Eq.~\eqref{eq:ch3.basis-change-matrix} \\
Topology & \(q_{\mathrm{top}}\) & Subsection~\ref{sec:ch374} \\
\bottomrule
\end{tabular}
\end{table}

\paragraph{Spin-orbit redistribution under gauge (import).}
Eq.~\eqref{eq:ch1.spin-orbit-gauge-redistribution} shows local spin and orbit densities can move between labels without changing
\(\mathbf{j}_{\mathrm{em}}\) \cite{belinfante1940,stratton1941}. Subsection~\ref{sec:ch384} uses that import to block pixel-level OAM
narratives on near-field maps unless Belinfante-class splits are declared \cite{belinfante1940,harrington2001}.

\paragraph{Worked illustration (cylindrical \(m\) vs topological \(q\)).}
A cylindrical mode with \(m=1\) produces azimuthal phase winds; topological charge \(q_{\mathrm{top}}\) from
Eq.~\eqref{eq:ch3.topological-angular-limit} matches \(m\) only for pure harmonics on declared loops \cite{allen1992,hansen1988,stratton1941}.
Superpositions break the equality even when cylindrical \(m\) is fixed \cite{harrington2001,hansen1988}.

\paragraph{Problem (compare three \(m\) notions).}
\textbf{Problem.} Distinguish \(m_{\mathrm{LG}}\), cylindrical \(m\), and VSH \(m\) for the same hardware output
\cite{allen1992,harrington2001,stratton1941,hansen1988}.

\textbf{Step 1.} Declare measurement chart \cite{stratton1941}.

\textbf{Step 2.} Apply Eq.~\eqref{eq:ch3.basis-change-matrix} between charts \cite{hansen1988}.

\textbf{Step 3.} Tabulate Table~\ref{tab:ch3.384-OAM-sources} \cite{harrington2001}.

\textbf{Physical meaning.} \(m\) is chart-specific until transformed \cite{jackson1999,belinfante1940}.

\begin{figure}[t]
\centering
\ChOneFigBlock{%
\begin{tikzpicture}[ch1fig, node distance=7mm and 9mm]
  \node[ch1box] (lg) {$m_{\mathrm{LG}}$};
  \node[ch1box, right=8mm of lg] (cyl) {cyl $m$};
  \node[ch1box, right=8mm of cyl] (vsh) {VSH $m$};
  \node[ch1box, right=8mm of vsh] (topo) {$q_{\mathrm{top}}$};
  \node[ch1box, below=10mm of cyl] (gate) {Eq.~\eqref{eq:ch3.gauge-screen-OAM-claim}};
  \draw[ch1flow] (lg) -- (gate);
  \draw[ch1flow] (cyl) -- (gate);
  \draw[ch1flow] (vsh) -- (gate);
  \draw[ch1flow] (topo) -- (gate);
\end{tikzpicture}%
}
\caption{Subsection~\ref{sec:ch384} multiple OAM-like indices converge only after basis declaration and gauge screen.}
\label{fig:ch3.384-multi-index-gate}
\end{figure}

\paragraph{Helicity coupling pointer.}
Helicity sectors from Subsection~\ref{sec:ch325} mix under \(\mathbf{T}\) unless symmetry blocks them \cite{belinfante1940,hansen1988}.
Subsection~\ref{sec:ch384} records that OAM discussions cannot ignore polarization/helicity tracks \(\sigma\) in
Eq.~\eqref{eq:ch3.integer-ell-def} \cite{harrington2001,stratton1941}.

\paragraph{Worked illustration (array-imprinted vortices).}
Phased arrays imprint local phase winds that need not align with a single VSH \(m\) sector after decomposition
\cite{allen1992,harrington2001,stratton1941}. Array synthesis should target \(\widehat{c}_{\ell m}^{(\sigma)}\) vectors, not informal OAM
headlines \cite{hansen1988,harrington2001}.

\paragraph{Problem (representation change alters sparsity).}
\textbf{Problem.} A beam sparse in plane-wave chart appears dense in VSH. Does OAM content change?
\cite{harrington2001,hansen1988,stratton1941,jackson1999}.

\textbf{Step 1.} Recognize \(\mathcal{P}_{\mathrm{rad}}\) is invariant Eq.~\eqref{eq:ch3.power-invariance-basis-change} \cite{poynting1884}.

\textbf{Step 2.} Accept sparsity changes Eq.~\eqref{eq:ch3.OAM-representation-dependence} \cite{hansen1988}.

\textbf{Physical meaning.} Physics is single-valued; labels are not \cite{stratton1941,harrington2001}.

\paragraph{Accessibility screen on OAM reports.}
Eq.~\eqref{eq:ch3.accessible-watt-budget} restricts which \(\widehat{c}_{\ell m}^{(\sigma)}\) sectors are reportable after geometry
\cite{harrington2001,stratton1941}. Subsection~\ref{sec:ch384} applies that screen before equating measured vortices with exported OAM
transport \cite{hansen1988,harrington2001}.

\begin{figure}[t]
\centering
\ChOneFigBlock{%
\begin{tikzpicture}[ch1fig, node distance=7mm and 9mm]
  \node[ch1box] (claim) {OAM claim};
  \node[ch1box, right=10mm of claim] (scr) {Eq.~\eqref{eq:ch3.gauge-screen-OAM-claim}};
  \node[ch1box, right=10mm of scr] (acc) {Eq.~\eqref{eq:ch3.accessible-watt-budget}};
  \draw[ch1flow] (claim) -- (scr) -- (acc);
\end{tikzpicture}%
}
\caption{Subsection~\ref{sec:ch384} gated OAM reporting: gauge screen then accessibility budget.}
\label{fig:ch3.384-gated-OAM-report}
\end{figure}

\begin{table}[t]
\centering
\caption{Subsection~\ref{sec:ch384} OAM claim strength ladder.}
\label{tab:ch3.384-claim-ladder}
\small
\begin{tabular}{@{}p{0.22\linewidth}p{0.68\linewidth}@{}}
\toprule
Level & Required evidence \\
\midrule
Phase image & winding on declared loop \\
Chart-specific & Table~\ref{tab:ch3.384-OAM-sources} row \\
Decomposed & full \(\widehat{c}_{\ell m}^{(\sigma)}\) vector \\
Transport & Eqs.~\eqref{eq:ch3.gauge-screen-OAM-claim}--\eqref{eq:ch1.topology-qualified-gate} \\
\bottomrule
\end{tabular}
\end{table}

\paragraph{Volume test import (pointer).}
Eq.~\eqref{eq:ch1.spin-orbit-volume-test} from Subsection~\ref{sec:ch01.4.5} supplies integrated closure tests for spin-orbit narratives
\cite{belinfante1940,stratton1941}. Subsection~\ref{sec:ch384} cites that test as a promotion criterion from phase images to transport claims
\cite{harrington2001,belinfante1940}; full gauge-invariant treatment continues in Section~\ref{sec:ch39} \cite{belinfante1940,harrington2001}.

\paragraph{Worked illustration (same state, four headlines).}
The same \(\mathbf{v}_\omega\) might be called ``LG \(m=1\),'' ``cylindrical \(m=1\),'' ``\(\ell=1\) VSH,'' and ``\(q_{\mathrm{top}}=1\)''
in different charts \cite{allen1992,hansen1988,harrington2001,stratton1941}. Subsection~\ref{sec:ch384} requires publishing the chart and
transformation steps, not the shortest slogan \cite{harrington2001,belinfante1940}.

\paragraph{Problem (design review for OAM feed).}
\textbf{Problem.} An OAM feed proposal cites only \(m_{\mathrm{LG}}\). List additional Subsection~\ref{sec:ch384} deliverables
\cite{allen1992,harrington2001,belinfante1940,stratton1941}.

\textbf{Step 1.} Far-zone VSH vector \cite{hansen1988}.

\textbf{Step 2.} \(\mathbf{T}\) from Eq.~\eqref{eq:ch3.paraxial-vs-spherical-OAM} \cite{harrington2001}.

\textbf{Step 3.} Table~\ref{tab:ch3.384-claim-ladder} level declaration \cite{belinfante1940}.

\textbf{Physical meaning.} Feed design reviews need representation certificates \cite{stratton1941,harrington2001}.

\paragraph{Peer-review rule.}
OAM manuscripts must state chart \(\mathcal{B}\), supply \(\mathbf{T}\) when comparing charts, and cite Eq.~\eqref{eq:ch3.gauge-screen-OAM-claim}
\cite{belinfante1940,harrington2001,allen1992,stratton1941}. Subsection~\ref{sec:ch384} treats missing chart declaration as a representation
artifact, not a completed physical claim \cite{harrington2001,belinfante1940}.

\paragraph{OAM vocabulary and representation dependence.}
\emph{Orbital angular momentum} in paraxial optics names a winding of phase in transverse charts; in spherical radiation
theory the same colloquial phrase must be split into \emph{magnetic quantum number} labels tied to
Eq.~\eqref{eq:ch3.integer-ell-def} and chart-dependent vortex structure \cite{allen1992,belinfante1940,harrington2001}.
Equation~\eqref{eq:ch3.OAM-representation-dependence} makes the philosophical warning explicit: \(\ell\) or OAM
headlines computed in LG charts need not coincide with spherical VSH labels on \(\mathcal{R}_\omega\) without a
declared bridge \cite{allen1992,hansen1988,stratton1941}. Equation~\eqref{eq:ch3.paraxial-vs-spherical-OAM} embeds the
validity boundary: paraxial winding is a controlled approximation coefficient, not an automatic substitute for
Belinfante--Maxwell totals \cite{belinfante1940,harrington2001}. The gauge screen
Eq.~\eqref{eq:ch3.gauge-screen-OAM-claim} closes the vocabulary loop: spin, orbit, and total angular momentum may not
appear in one sentence unless the gauge class, chart, and outgoing certificate match the definitions imported from
Chapter~\ref{ch:01} \cite{belinfante1940,stratton1941}. Implicit in every OAM discussion is \emph{representation
dependence of vortex counting}: phase singularities follow coordinate choices and beam envelopes, while conserved
fluxes follow the normalized eigenfield meters of Section~\ref{sec:ch35} \cite{harrington2001,allen1992,belinfante1940}.

\paragraph{Topological charge versus magnetic quantum number.}
Phase winds in LG charts carry topological index \(q_{\mathrm{top}}\); spherical labels use integer \(\ell,m\) from Eq.~\eqref{eq:ch3.integer-ell-def} \cite{allen1992,hansen1988,stratton1941}. Subsection~\ref{sec:ch384} records
\begin{equation}
\label{eq:ch3.topological-vs-magnetic-index}
\mathcal{S}_{\mathrm{OAM}}=1
\iff
\left(q_{\mathrm{top}}=m_{\mathrm{mag}}\ \vee\ \mathcal{T}_{\mathrm{decl}}=1\right),
\end{equation}
with \(\mathcal{T}_{\mathrm{decl}}=1\) denoting a declared basis-change matrix Eq.~\eqref{eq:ch3.basis-change-matrix}, importing Eq.~\eqref{eq:ch3.gauge-screen-OAM-claim} \cite{belinfante1940,harrington2001}. Vortex imagery is not a substitute for \(\widehat{c}_{\ell m}^{(\sigma)}\) on \(\mathcal{R}_\omega\) \cite{hansen1988,stratton1941,allen1992}.

\paragraph{Array-imprinted vortex lattices.}
Periodic arrays imprint multiple near-field vortices while far-zone VSH spectra remain dense \cite{jackson1999,harrington2001}. Publish lattice pitch and \(\Pi_{\mathrm{rad}}\)-filtered coefficients together \cite{hansen1988,belinfante1940,stratton1941}.

\paragraph{Four-chart headline equivalence test.}
Subsection~\ref{sec:ch384} closes with an explicit falsification protocol: the same \(\mathbf{v}_\omega\) must yield identical \(\mathcal{P}_{\mathrm{rad}}\) and certified fluxes while permitting different sparse headlines in plane-wave, cylindrical, LG, and VSH charts
\cite{allen1992,hansen1988,harrington2001,belinfante1940}. The test record is
\begin{equation}
\label{eq:ch3.oam-four-chart-flux-lock}
\mathcal{P}_{\mathrm{rad}}^{(\mathrm{pw})}
=
\mathcal{P}_{\mathrm{rad}}^{(\mathrm{cyl})}
=
\mathcal{P}_{\mathrm{rad}}^{(\mathrm{VSH})}
\quad\Longrightarrow\quad
\mathcal{H}_{m}=0,
\end{equation}
with \(\mathcal{H}_{m}=0\) denoting that headline \(m\) labels may differ across charts while fluxes agree
with each chart normalized by its Section~\ref{sec:ch35} meter and linked by Eq.~\eqref{eq:ch3.basis-change-matrix} \cite{poynting1884,stratton1941,harrington2001}. Failure of Eq.~\eqref{eq:ch3.oam-four-chart-flux-lock} is a representation bug; agreement with different \(m\) headlines is expected representation dependence, not physics paradox
\cite{belinfante1940,allen1992,hansen1988}.

\paragraph{Spin--orbit wording gate (expanded).}
Eq.~\eqref{eq:ch3.gauge-screen-OAM-claim} forbids mixing spin, orbital, and total angular-momentum phrases unless gauge class, chart, and outgoing certificate match Chapter~\ref{ch:01} definitions
\cite{belinfante1940,stratton1941,harrington2001}. Subsection~\ref{sec:ch384} applies that gate to marketing language: ``pure OAM beam'' is admissible only with \(\mathcal{S}_{\mathrm{OAM}}=1\) in Eq.~\eqref{eq:ch3.topological-vs-magnetic-index} and attached \(\mathbf{T}\) metadata
\cite{allen1992,hansen1988,devaney2004}.

\paragraph{Measurement-plane versus far-zone vortex counting.}
Near-field vortex counts on a fixed plane can exceed far-zone topological indices declared on \(\Gamma_R\) after \(\Pi_{\mathrm{rad}}\)
\cite{allen1992,harrington2001,belinfante1940}. Subsection~\ref{sec:ch384} requires that any \(q_{\mathrm{top}}\) headline specify the integration loop and filter order: near-plane winds are pre-export unless Eq.~\eqref{eq:ch3.gauge-screen-OAM-claim} and Eq.~\eqref{eq:ch3.debye-vsh-promotion-screen} pass
\cite{stratton1941,hansen1988,jackson1999}. Far-zone \(\widehat{c}_{\ell m}^{(\sigma)}\) remain the export standard for OAM transport claims in Volume~I
\cite{harrington2001,belinfante1940}.

\paragraph{Array lattice imprint on sparse OAM headlines.}
Periodic arrays imprint multiple vortex cores in near fields while far-zone VSH spectra remain dense in \(\ell\)
\cite{jackson1999,harrington2001,devaney2004}. Subsection~\ref{sec:ch384} pairs lattice pitch with Eq.~\eqref{eq:ch3.oam-four-chart-flux-lock}: flux agreement is mandatory; sparse OAM headlines are optional chart-dependent summaries
\cite{hansen1988,belinfante1940,stratton1941}.

\paragraph{Design-review checklist for OAM feeds.}
Subsection~\ref{sec:ch384} closes with a design-review sequence: declare chart \(\mathcal{B}\); publish \(\mathbf{T}\) to VSH; verify Eq.~\eqref{eq:ch3.oam-four-chart-flux-lock}; apply Eq.~\eqref{eq:ch3.gauge-screen-OAM-claim}; attach \(\widehat{c}_{\ell m}^{(\sigma)}\) on \(\Gamma_R\)
\cite{allen1992,belinfante1940,harrington2001,hansen1988}. Skipping any step leaves a representation artifact masquerading as a transport claim
\cite{stratton1941,devaney2004,jackson1999}. The checklist is not administrative; it is the scientific condition for comparing beams across laboratories
\cite{harrington2001,belinfante1940}.

Subsection~\ref{sec:ch384} has analyzed representation dependence of OAM modes through
Eqs.~\eqref{eq:ch3.OAM-representation-dependence}--\eqref{eq:ch3.gauge-screen-OAM-claim}
\cite{allen1992,belinfante1940,harrington2001}. Section~\ref{sec:ch38} is complete: plane-wave/Fourier, cylindrical/parabolic, equivalence
transformations, and OAM representation limits now share one matrix discipline on \(\mathcal{R}_\omega\) \cite{hansen1988,stratton1941}.

The four subsections of Section~\ref{sec:ch38} make basis changes explicit:
Eqs.~\eqref{eq:ch3.plane-wave-basis-resolution}--\eqref{eq:ch3.fourier-radiation-filter-chain} in Subsection~\ref{sec:ch381};
Eqs.~\eqref{eq:ch3.cylindrical-mode-expansion}--\eqref{eq:ch3.cylindrical-parabolic-bridge} in Subsection~\ref{sec:ch382};
Eqs.~\eqref{eq:ch3.basis-change-matrix}--\eqref{eq:ch3.power-invariance-basis-change} in Subsection~\ref{sec:ch383}; and
Eqs.~\eqref{eq:ch3.OAM-representation-dependence}--\eqref{eq:ch3.gauge-screen-OAM-claim} in Subsection~\ref{sec:ch384}
\cite{stratton1941,jackson1999,harrington2001,allen1992,belinfante1940}. Section~\ref{sec:ch39} now closes the gauge thread on the same
eigenfield language \cite{belinfante1940,stratton1941}.


\section{Gauge-Invariant Interpretation and Approximation Limits}
\label{sec:ch39}

Representations can be mathematically complete yet physically misleading if gauge-dependent quantities are
read as observables. Section~\ref{sec:ch39} closes the gauge thread opened in Chapter~\ref{ch:01} on the
eigenfield language developed in Sections~\ref{sec:ch32}--\ref{sec:ch38} \cite{belinfante1940,stratton1941}.
Subsection~\ref{sec:ch391} states gauge-invariant meaning of eigenfields using Section~\ref{sec:ch01.4} by
reference. Subsection~\ref{sec:ch392} separates conserved from representation-dependent quantities, importing
helicity discipline from Subsection~\ref{sec:ch01.4.3}. Subsection~\ref{sec:ch393} records scalar and paraxial
limits where vector eigenfield theory reduces consistently. Subsection~\ref{sec:ch394} contrasts observables
with representation artifacts, connecting to measurement inference rules from
Subsection~\ref{sec:ch01.5.4} \cite{harrington2001}.


\subsection{Gauge-invariant meaning of eigenfields}
\label{sec:ch391}
% TOC tag: [USE SS1.4]

\emph{Gauge-invariant meaning} asks which statements about angular-momentum eigenfields remain true when potentials are
reparameterized by Eq.~\eqref{eq:ch1.gauge-transform} while \(\E\) and \(\H\) are held fixed
\cite{belinfante1940,stratton1941,harrington2001}. \emph{Eigenfields} in Chapter~\ref{ch:03} are vector modes
\(\widehat{\mathbf{u}}_{\ell m}^{(\sigma)}\in\mathcal{F}_{\mathrm{rad}}\) from Eq.~\eqref{eq:ch3.normalized-coefficient-resolution},
not potential eigenfunctions \cite{hansen1988,stratton1941}. By Principle~\ref{prin:ch3.osr-separation} and
Principle~\ref{prin:ch1.rep-physics}, gauge-invariant meaning attaches to field observables and flux functionals,
not to coordinate relabelings alone. Subsection~\ref{sec:ch391} is the USE bridge that applies
Section~\ref{sec:ch01.4} gauge screening to the representation block of Sections~\ref{sec:ch32}--\ref{sec:ch38}
\cite{belinfante1940,harrington2001}.

\paragraph{Standing imports (no re-derivation).}
Gauge transformation Eq.~\eqref{eq:ch1.gauge-transform}, field invariance check Eq.~\eqref{eq:ch1.gauge-field-invariance-check},
gauge-observable invariance Eq.~\eqref{eq:ch1.gauge-observable-invariance}, gauge-qualified decision
Eq.~\eqref{eq:ch1.gauge-qualified-decision}, potential-field relations Eq.~\eqref{eq:ch1.potential-field-relations}, and
normalized eigenfield resolution Eq.~\eqref{eq:ch3.normalized-coefficient-resolution} enter by reference
\cite{stratton1941,belinfante1940,hansen1988}. Subsection~\ref{sec:ch391} does not re-derive Lorenz gauge or retarded potentials
\cite{stratton1941,harrington2001}.

\paragraph{Assumptions.}
Work on transverse far-zone \(\mathbf{v}_\omega\in\mathcal{F}_{\mathrm{rad}}\) with outgoing selection already applied
\cite{harrington2001,stratton1941}. Eigenfield statements are evaluated on \(\E,\H\) and quadratic flux functionals, not on
\((\mathbf{A},\phi)\) unless a gauge certificate is attached \cite{belinfante1940,stratton1941}.

\paragraph{Eigenfield gauge-invariance class (HOME).}
Let \(\mathfrak{G}\) denote the gauge orbit of potentials generating the same \((\E,\H)\). Subsection~\ref{sec:ch391} defines the
admissible statement class
\begin{equation}
\label{eq:ch3.eigenfield-gauge-invariance-class}
\mathcal{S}_{\mathrm{gi}}
=
\bigl\{\text{claims about }\mathbf{v}_\omega\text{ expressible solely through }
\E,\H,\mathcal{P}_{\mathrm{rad}},\widehat{c}_{\ell m}^{(\sigma)}\bigr\},
\end{equation}
with coefficients from Eq.~\eqref{eq:ch3.normalized-coefficient-resolution} \cite{stratton1941,hansen1988,harrington2001}. Equation
\eqref{eq:ch3.eigenfield-gauge-invariance-class} is the Subsection~\ref{sec:ch391} meaning law: eigenfield expansions are
\emph{gauge-invariant coordinates} on \(\mathcal{F}_{\mathrm{rad}}\) because projections use the radiative inner product on
\(\E,\H\), not on \(\mathbf{A}\) \cite{belinfante1940,harrington2001,kato1995}. Potential-based OAM formulas that move under
Eq.~\eqref{eq:ch1.gauge-transform} lie outside \(\mathcal{S}_{\mathrm{gi}}\) unless promoted by
Eq.~\eqref{eq:ch1.gauge-qualified-decision} \cite{belinfante1940,stratton1941}.

\paragraph{Gauge-qualified eigenfield report (HOME).}
\begin{equation}
\label{eq:ch3.gauge-qualified-eigenfield-report}
\mathcal{G}_{\mathrm{rep}}=1
\iff
\left(
\mathcal{G}_{\mathrm{gauge}}=1\
\wedge\
\mathcal{P}_{\mathrm{inv}}=1
\right),
\end{equation}
with \(\mathcal{G}_{\mathrm{gauge}}=1\) the gauge-qualified decision of Section~\ref{sec:ch01.4} and \(\mathcal{P}_{\mathrm{inv}}=1\) the watt invariance screen Eq.~\eqref{eq:ch3.power-invariance-basis-change} from Section~\ref{sec:ch38}
\cite{poynting1884,harrington2001,stratton1941,belinfante1940}. Equation
\eqref{eq:ch3.gauge-qualified-eigenfield-report} is the publication gate for eigenfield tables: coefficients are admissible
observables only when gauge screening and normalization meters are documented \cite{harrington2001,belinfante1940,hansen1988}.

\begin{theorem}[Gauge-invariant eigenfield meaning (HOME)]
\label{thm:ch3.gauge-invariant-eigenfield-meaning}
For \(\mathbf{v}_\omega\in\mathcal{F}_{\mathrm{rad}}\), the normalized expansion
Eq.~\eqref{eq:ch3.normalized-coefficient-resolution} and power law Eq.~\eqref{eq:ch3.mode-power-from-coefficient} define
gauge-invariant content in \(\mathcal{S}_{\mathrm{gi}}\); potential-level azimuthal densities are gauge-orbit coordinates unless
Eq.~\eqref{eq:ch1.gauge-observable-invariance} certifies otherwise \cite{belinfante1940,stratton1941,hansen1988,harrington2001}.
\end{theorem}
\begin{proof}[Proof sketch]
Coefficients are Riesz projections on \(\mathcal{F}_{\mathrm{rad}}\) built from \(\E,\H\); gauge transforms cancel in
Eq.~\eqref{eq:ch1.gauge-field-invariance-check} before projection \cite{stratton1941,hansen1988,kato1995}.
% PROOF: AUTHOR REVIEW
\end{proof}

\begin{figure}[t]
\centering
\ChOneFigBlock{%
\begin{tikzpicture}[ch1fig, node distance=7mm and 9mm]
  \node[ch1box] (pot) {$(\mathbf{A},\phi)$ orbit};
  \node[ch1box, right=10mm of pot] (EH) {$\E,\H$ fixed};
  \node[ch1box, right=10mm of EH] (vsh) {Eq.~\eqref{eq:ch3.normalized-coefficient-resolution}};
  \node[ch1box, below=10mm of EH] (gate) {Eq.~\eqref{eq:ch3.gauge-qualified-eigenfield-report}};
  \draw[ch1flow] (pot) -- (EH) -- (vsh);
  \draw[ch1flow] (gate) -- (vsh);
\end{tikzpicture}%
}
\caption{Subsection~\ref{sec:ch391} eigenfield meaning: projections on \(\E,\H\) survive gauge orbits; potential formulas require separate gates.}
\label{fig:ch3.391-gauge-orbit}
\end{figure}

\paragraph{Physical meaning of eigenfield gauge invariance.}
One must not treat \(\widehat{\mathbf{u}}_{\ell m}^{(\sigma)}\) as potential eigenfunctions; in Volume~I they are
\emph{field} eigenmodes of the export operator image \cite{hansen1988,harrington2001,stratton1941}. Changing Lorenz gauge does not
rotate \(\widehat{c}_{\ell m}^{(\sigma)}\) if \(\mathbf{v}_\omega\) is unchanged, because coefficients measure watt content of the
same \(\E,\H\) pattern \cite{poynting1884,belinfante1940,harrington2001}. Subsection~\ref{sec:ch391} therefore elevates coefficient
vectors and \(\mathcal{P}_{\mathrm{rad}}\) to the primary gauge-invariant data products of the representation block
\cite{stratton1941,hansen1988}.

\paragraph{Worked illustration (same fields, different potentials).}
Two analysts choose different \(\chi\) in Eq.~\eqref{eq:ch1.gauge-transform} yet extract identical \(\widehat{c}_{\ell m}^{(\sigma)}\)
from the same \(\Gamma_R\) samples when projections are correctly normalized \cite{stratton1941,harrington2001,hansen1988}. Their
potential plots differ; their eigenfield tables agree \cite{belinfante1940,harrington2001}.

\paragraph{Problem (audit a potential-based OAM formula).}
\textbf{Problem.} A formula uses \(\mathbf{A}\) alone to define ``orbital'' density. Classify using
Eqs.~\eqref{eq:ch3.eigenfield-gauge-invariance-class} and \eqref{eq:ch1.gauge-qualified-decision}
\cite{belinfante1940,stratton1941,harrington2001}.

\textbf{Step 1.} Test Eq.~\eqref{eq:ch1.gauge-transform} invariance \cite{stratton1941}.

\textbf{Step 2.} If it fails, re-express in \(\mathcal{S}_{\mathrm{gi}}\) via Eq.~\eqref{eq:ch3.normalized-coefficient-resolution} \cite{hansen1988}.

\textbf{Physical meaning.} Potential formulas are not eigenfield observables by default \cite{belinfante1940,harrington2001}.

\begin{table}[t]
\centering
\caption{Subsection~\ref{sec:ch391} gauge-invariant eigenfield objects.}
\label{tab:ch3.391-gi-objects}
\small
\begin{tabular}{@{}p{0.26\linewidth}p{0.38\linewidth}p{0.30\linewidth}@{}}
\toprule
Object & Equation & Gauge status \\
\midrule
Coefficients & Eq.~\eqref{eq:ch3.normalized-coefficient-resolution} & invariant \\
Power & Eq.~\eqref{eq:ch3.mode-power-from-coefficient} & invariant \\
Potential OAM & Eq.~\eqref{eq:ch1.gauge-transform} orbit & conditional \\
Report gate & Eq.~\eqref{eq:ch3.gauge-qualified-eigenfield-report} & required \\
\bottomrule
\end{tabular}
\end{table}

\paragraph{Transverse far-zone reduction.}
Subsection~\ref{sec:ch391} restricts gauge discussion to transverse \(\E_{t},\H_{t}\) on \(\Gamma_R\) where
Eq.~\eqref{eq:ch1.radiation-trace-condition} applies \cite{stratton1941,harrington2001}. Longitudinal gauge artifacts in near-field
potentials do not invalidate far-zone \(\widehat{c}_{\ell m}^{(\sigma)}\) if extraction uses \(\Pi_{\mathrm{rad}}\) first
\cite{harrington2001,belinfante1940,jackson1999}.

\paragraph{Worked illustration (Debye potentials versus VSH).}
Debye potentials Subsection~\ref{sec:ch334} are construction tools; exported eigenfields are the gauge-invariant end product
Eq.~\eqref{eq:ch3.normalized-coefficient-resolution} \cite{hansen1988,stratton1941}. Subsection~\ref{sec:ch391} forbids identifying
Debye gauge choices with physical eigenfield labels \cite{harrington2001,belinfante1940}.

\paragraph{Problem (gauge residual on coefficient table).}
\textbf{Problem.} Define when a published \(\widehat{c}\) table fails Eq.~\eqref{eq:ch3.gauge-qualified-eigenfield-report}
\cite{harrington2001,stratton1941,belinfante1940,hansen1988}.

\textbf{Step 1.} Check normalization meters Section~\ref{sec:ch35} \cite{hansen1988}.

\textbf{Step 2.} Verify outgoing certificate \cite{harrington2001}.

\textbf{Step 3.} Apply Eq.~\eqref{eq:ch1.gauge-qualified-decision} \cite{belinfante1940}.

\textbf{Physical meaning.} Incomplete metadata voids gauge claims \cite{stratton1941,harrington2001}.

\begin{figure}[t]
\centering
\ChOneFigBlock{%
\begin{tikzpicture}[ch1fig, node distance=7mm and 9mm]
  \node[ch1box] (deb) {Debye construction};
  \node[ch1box, right=10mm of deb] (v) {$\mathbf{v}_\omega$};
  \node[ch1box, right=10mm of v] (c) {$\widehat{c}_{\ell m}^{(\sigma)}$};
  \node[ch1box, below=10mm of v] (gi) {Eq.~\eqref{eq:ch3.eigenfield-gauge-invariance-class}};
  \draw[ch1flow] (deb) -- (v) -- (c);
  \draw[ch1flow] (gi) -- (c);
\end{tikzpicture}%
}
\caption{Subsection~\ref{sec:ch391} construction path: intermediate potentials are not the final gauge-invariant report.}
\label{fig:ch3.391-construction-path}
\end{figure}

\paragraph{Gauge orbit, rotation, and field-first observability.}
An \emph{eigenfield} in Chapter~\ref{ch:03} is a simultaneous solution of the vector Helmholtz operator and the angular-momentum operators on
\(\mathcal{F}_{\mathrm{rad}}\); a \emph{gauge orbit} is the family of potentials generating the same \((\E,\H)\)
\cite{belinfante1940,stratton1941,hansen1988}. Laboratory rotations act on \(\mathbf{v}_\omega\) through
Eq.~\eqref{eq:ch3.coefficient-rotation-law}, whereas gauge transformations move potentials along orbits without altering \(\mathbf{v}_\omega\)
\cite{jackson1999,hansen1988,belinfante1940}. Subsection~\ref{sec:ch391} encodes that separation on fields by
\begin{equation}
\label{eq:ch3.gauge-rotation-commute-fields}
\mathcal{J}_{g}\,\mathbf{v}_\omega
=
\mathbf{v}_\omega^{(g)},
\qquad
\mathfrak{G}:\ \mathbf{A}\mapsto\mathbf{A}+\nabla\chi
\ \Longrightarrow\
(\E,\H)\ \text{fixed on }\mathcal{R}_\omega,
\end{equation}
importing rotation covariance from Section~\ref{sec:ch01.3} and gauge orbits from Section~\ref{sec:ch01.4} by reference only
\cite{belinfante1940,jackson1999,stratton1941}. Equation~\eqref{eq:ch3.gauge-rotation-commute-fields} is the operational rule: rotated eigenfield
tables are symmetry certificates; uncited potential relabelings are not \cite{harrington2001,hansen1988}. Equations~\eqref{eq:ch3.eigenfield-gauge-invariance-class}
and~\eqref{eq:ch3.gauge-qualified-eigenfield-report} then supply the admissible reporting tuple---outgoing certificate, transverse radiation gauge, and
normalization meters---without which coefficient statements about \(\mathbf{M}_{\ell m},\mathbf{N}_{\ell m}\) collapse into informal mode naming
\cite{stratton1941,hansen1988}. Theorem~\ref{thm:ch3.gauge-invariant-eigenfield-meaning} elevates \emph{field-first observability}: fluxes,
orthogonality inner products, and operator eigenvalues on \(\mathcal{F}_{\mathrm{rad}}\) are gauge-invariant reports; longitudinal potentials and
uncited phase winds are not \cite{belinfante1940,harrington2001,jackson1999}.

\paragraph{Far-zone transverse projector on eigenfield reports.}
On \(\mathcal{R}_\omega\), only transverse components carry the outward flux budget of Eq.~\eqref{eq:ch3.mode-power-from-coefficient}
\cite{poynting1884,stratton1941,harrington2001}. Subsection~\ref{sec:ch391} therefore applies
\begin{equation}
\label{eq:ch3.far-zone-transverse-projector}
\mathbf{v}_\omega^{\perp}
=
\hat{\mathbf{r}}\times(\hat{\mathbf{r}}\times\mathbf{v}_\omega),
\qquad
\mathcal{P}_{\mathrm{rad}}
=
\langle \mathbf{v}_\omega^{\perp},\mathbf{v}_\omega^{\perp}\rangle_{\mathcal{F}},
\end{equation}
before any eigenfield coefficient table is promoted to an observable \cite{hansen1988,jackson1999,stratton1941}. Equation~\eqref{eq:ch3.far-zone-transverse-projector}
removes longitudinal contamination that gauge choices can introduce in near-zone reconstructions but that must not enter far-zone watt ledgers
\cite{belinfante1940,harrington2001}. Lorenz and Coulomb gauges differ in near-field potentials yet yield identical
Eq.~\eqref{eq:ch3.far-zone-transverse-projector} outputs when evaluated on the same \(\mathbf{v}_\omega\) \cite{stratton1941,jackson1999,harrington2001}.

\paragraph{Gauge residual on projected coefficients.}
Even after Eq.~\eqref{eq:ch3.far-zone-transverse-projector}, potential-based fitting can leave a gauge-dependent phase on individual
\(\widehat{c}_{\ell m}\) that cancels only in quadratic observables \cite{belinfante1940,hansen1988,harrington2001}. Subsection~\ref{sec:ch391}
bounds the promotion defect by
\begin{equation}
\label{eq:ch3.gauge-coefficient-residual}
\Delta_{\mathrm{gauge}}
=
\min_{\chi}\,
\bigl\|\widehat{\mathbf{c}}-\widehat{\mathbf{c}}^{(\chi)}\bigr\|_{2}.
\end{equation}
A gauge-qualified report requires \(\Delta_{\mathrm{gauge}}=0\) within numerical tolerance; \(\widehat{\mathbf{c}}^{(\chi)}\) denotes the coefficient vector after an admissible gauge transform in the generating-potential chart
\cite{belinfante1940,stratton1941}. Equation~\eqref{eq:ch3.gauge-coefficient-residual} is the numerical face of
Theorem~\ref{thm:ch3.gauge-invariant-eigenfield-meaning}: flux and orthogonality tests may pass while \(\Delta_{\mathrm{gauge}}\) remains large,
signalling that only quadratic summaries are safe observables \cite{harrington2001,hansen1988,jackson1999}.

\paragraph{Worked illustration (Lorenz versus Coulomb at export).}
Lorenz and Coulomb gauges differ in near-field potentials yet yield the same \(\mathcal{P}_{\mathrm{rad}}\) when
Eq.~\eqref{eq:ch3.mode-power-from-coefficient} is evaluated on identical \(\mathbf{v}_\omega\) \cite{stratton1941,jackson1999,harrington2001}.
Subsection~\ref{sec:ch391} uses that invariance as a laboratory consistency check \cite{belinfante1940,hansen1988}.

\paragraph{Problem (promote near-field potential map).}
\textbf{Problem.} A near-field \(\mathbf{A}\) plot is proposed as eigenfield evidence. List promotion requirements
\cite{belinfante1940,stratton1941,harrington2001,hansen1988}.

\textbf{Step 1.} Map to \(\mathbf{v}_\omega\) via \(\Pi_{\mathrm{rad}}\) \cite{harrington2001}.

\textbf{Step 2.} Project Eq.~\eqref{eq:ch3.normalized-coefficient-resolution} \cite{hansen1988}.

\textbf{Step 3.} Satisfy Eq.~\eqref{eq:ch3.gauge-qualified-eigenfield-report} \cite{belinfante1940}.

\textbf{Physical meaning.} Near-field potentials are pre-observable unless gated \cite{stratton1941,harrington2001}.

\paragraph{Peer-review metadata.}
Eigenfield publications must archive gauge class, outgoing certificate, normalization tuple, and
Eq.~\eqref{eq:ch3.gauge-qualified-eigenfield-report} checklist \cite{harrington2001,stratton1941,belinfante1940,hansen1988}.
Subsection~\ref{sec:ch391} treats missing gauge metadata as a representation risk, not a cosmetic omission \cite{belinfante1940,harrington2001}.

\paragraph{Debye-to-VSH promotion protocol.}
Debye potentials from Subsection~\ref{sec:ch334} are construction coordinates; VSH coefficients are export coordinates \cite{tai1994,hansen1988,stratton1941}. Subsection~\ref{sec:ch391} requires
\begin{equation}
\label{eq:ch3.debye-vsh-promotion-screen}
\mathcal{P}_{\mathrm{Debye}\to\mathrm{VSH}}=1
\iff
\left(
\Pi_{\mathrm{rad}}\ \wedge\
\Delta_{\mathrm{gauge}}=0\
\wedge\
\mathcal{T}_{\perp}=1
\right),
\end{equation}
with \(\mathcal{T}_{\perp}=1\) denoting passage of Eq.~\eqref{eq:ch3.far-zone-transverse-projector}, before Debye-derived amplitudes enter Eq.~\eqref{eq:ch3.gauge-qualified-eigenfield-report} \cite{belinfante1940,harrington2001,tai1994}. Mathematically, the screen is a composition of projection operators; physically, it prevents near-field potential fits from masquerading as eigenfield observables \cite{stratton1941,hansen1988,jackson1999}.

\paragraph{Lorenz--Coulomb coefficient drift (quantitative).}
Even when \(\mathcal{P}_{\mathrm{rad}}\) matches, intermediate gauges can drift \(\arg\widehat{c}_{\ell m}\) while \(|\widehat{c}_{\ell m}|^{2}\) remains fixed \cite{belinfante1940,jackson1999,stratton1941}. Subsection~\ref{sec:ch391} therefore promotes only magnitude-squared observables and orthogonality tests unless \(\Delta_{\mathrm{gauge}}=0\) in Eq.~\eqref{eq:ch3.gauge-coefficient-residual} \cite{harrington2001,hansen1988}.

\paragraph{Belinfante symmetrization and eigenfield observability.}
Total angular momentum on Maxwell fields is a Belinfante-symmetrized construction imported from Section~\ref{sec:ch01.4} \cite{belinfante1940,stratton1941}.
Subsection~\ref{sec:ch391} does not re-derive that symmetrization; it records the consequence for eigenfield reporting: \(\widehat{c}_{\ell m}^{(\sigma)}\)
are projections of \((\E,\H)\) pairs, not of canonical orbital currents in potentials \cite{hansen1988,harrington2001}. When a manuscript cites
``orbital'' density from \(\mathbf{A}\) alone, the correct classification is \emph{pre-observable} until Eq.~\eqref{eq:ch3.gauge-qualified-eigenfield-report}
passes \cite{belinfante1940,stratton1941}. Field-first observability is therefore a publication discipline, not a optional philosophical preference
\cite{harrington2001,jackson1999}.

\paragraph{Transverse projector composition with \(\Pi_{\mathrm{rad}}\).}
The export pipeline applies \(\Pi_{\mathrm{rad}}\) before far-zone projection in Eq.~\eqref{eq:ch3.far-zone-transverse-projector}
\cite{harrington2001,stratton1941}. Mathematically, \(\mathbf{v}_\omega^{\perp}\) is a composition of outgoing selection and transverse extraction;
physically, it is the object that watt meters integrate on \(\Gamma_R\) \cite{poynting1884,hansen1988}. Near-field gauges that inflate longitudinal
potential content must not enter \(\mathcal{S}_{\mathrm{gi}}\) in Eq.~\eqref{eq:ch3.eigenfield-gauge-invariance-class} until that composition is executed
\cite{belinfante1940,harrington2001}. Figure~\ref{fig:ch3.391-gauge-orbit} and Figure~\ref{fig:ch3.391-construction-path} together show that
potentials are construction orbits while \(\widehat{c}_{\ell m}^{(\sigma)}\) are export coordinates \cite{tai1994,stratton1941,hansen1988}.

\paragraph{Worked audit (gauge residual on a three-sector table).}
Suppose \(\widehat{c}_{10}^{(\mathbf{N})}\), \(\widehat{c}_{20}^{(\mathbf{N})}\), and \(\widehat{c}_{11}^{(\mathbf{M})}\) are extracted from the same
\(\Gamma_R\) samples under Lorenz and Coulomb gauges \cite{stratton1941,jackson1999}. Magnitude-squared rows should match within numerical tolerance;
phase rows may differ unless \(\Delta_{\mathrm{gauge}}=0\) in Eq.~\eqref{eq:ch3.gauge-coefficient-residual} \cite{belinfante1940,harrington2001}.
Subsection~\ref{sec:ch391} recommends archiving both gauge classes when phase-sensitive nulls are reported, because interferometric claims require
\(\Delta_{\mathrm{gauge}}\) certification while watt-only claims do not \cite{hansen1988,harrington2001,stratton1941}.

\paragraph{Publication checklist for eigenfield tables (consolidated).}
Subsection~\ref{sec:ch391} consolidates the admissible report tuple as
\begin{equation}
\label{eq:ch3.eigenfield-publication-checklist}
\mathcal{G}_{\mathrm{rep}}=1
\iff
\left(
\mathcal{G}_{\mathrm{gauge}}=1\
\wedge\
\mathcal{T}_{\perp}=1\
\wedge\
\Delta_{\mathrm{gauge}}=0\
\wedge\
\mathcal{N}_{\mathrm{mtr}}=1
\right),
\end{equation}
importing Eq.~\eqref{eq:ch3.gauge-qualified-eigenfield-report}, Eq.~\eqref{eq:ch3.far-zone-transverse-projector}, Eq.~\eqref{eq:ch3.gauge-coefficient-residual}, and Section~\ref{sec:ch35} meters by reference
\cite{belinfante1940,harrington2001,stratton1941,hansen1988}. Equation~\eqref{eq:ch3.eigenfield-publication-checklist} is the reviewer-facing distillation of Theorem~\ref{thm:ch3.gauge-invariant-eigenfield-meaning}: field-first coefficients are observables only when the full tuple passes
\cite{poynting1884,jackson1999,harrington2001}.

\paragraph{Near-field potential maps and export promotion workflow.}
Near-field \(\mathbf{A}\) maps are common in OAM feeds yet remain pre-observable until the promotion chain \(\Pi_{\mathrm{rad}}\to\) Eq.~\eqref{eq:ch3.far-zone-transverse-projector} \(\to\) Eq.~\eqref{eq:ch3.normalized-coefficient-resolution} \(\to\) Eq.~\eqref{eq:ch3.eigenfield-publication-checklist} completes
\cite{belinfante1940,harrington2001,stratton1941,tai1994}. Subsection~\ref{sec:ch391} orders that chain explicitly because potential heatmaps visually suggest eigenfield content before \(\mathcal{T}_{\perp}=1\) is certified
\cite{hansen1988,jackson1999}. Debye-derived amplitudes additionally require Eq.~\eqref{eq:ch3.debye-vsh-promotion-screen} before entering gauge-qualified tables
\cite{tai1994,stratton1941,harrington2001}.

\paragraph{Coefficient phase book versus magnitude book.}
Subsection~\ref{sec:ch391} separates the \emph{magnitude book} \(|\widehat{c}_{\ell m}^{(\sigma)}|^{2}\) certified by Eq.~\eqref{eq:ch3.mode-power-from-coefficient} from the \emph{phase book} \(\arg\widehat{c}_{\ell m}\) certified only when \(\Delta_{\mathrm{gauge}}=0\)
\cite{belinfante1940,harrington2001,jackson1999}. Watt-only campaigns may publish the magnitude book alone; interferometric campaigns must publish both books with gauge class metadata
\cite{stratton1941,hansen1988,devaney2004}. Confusing the two books is the dominant source of false gauge alarms in eigenfield tables
\cite{harrington2001,belinfante1940}.

Subsection~\ref{sec:ch391} has stated gauge-invariant eigenfield meaning through
Eqs.~\eqref{eq:ch3.eigenfield-gauge-invariance-class}--\eqref{eq:ch3.gauge-qualified-eigenfield-report} and
Theorem~\ref{thm:ch3.gauge-invariant-eigenfield-meaning} \cite{belinfante1940,stratton1941,harrington2001}. Subsection~\ref{sec:ch392}
separates conserved from representation-dependent quantities using helicity discipline \cite{stratton1941,belinfante1940}.


\subsection{Conserved versus representation-dependent quantities}
\label{sec:ch392}
% TOC tag: [USE SS1.4.3]

\emph{Conserved quantities} here mean flux-integrable totals whose closure tests survive gauge relabeling and rigid-frame motion
\cite{belinfante1940,stratton1941,poynting1884}. \emph{Representation-dependent quantities} change numerical values under legitimate
basis changes Eq.~\eqref{eq:ch3.basis-change-matrix} even when \(\mathbf{v}_\omega\) is fixed \cite{hansen1988,harrington2001}.
Subsection~\ref{sec:ch392} is the USE layer importing helicity and dual-symmetry observables from Subsection~\ref{sec:ch01.4.3} and
contrasting them with chart-specific labels from Section~\ref{sec:ch38} \cite{belinfante1940,stratton1941,harrington2001}.

\paragraph{Standing imports.}
Observable admissibility triad Eq.~\eqref{eq:ch1.observable-admissibility-triad}, helicity density
Eq.~\eqref{eq:ch1.helicity-density}, helicity flux integral Eq.~\eqref{eq:ch1.helicity-flux-integral}, helicity balance
Eq.~\eqref{eq:ch1.helicity-balance}, helicity observable decision Eq.~\eqref{eq:ch1.helicity-observable-decision}, power
invariance Eq.~\eqref{eq:ch3.power-invariance-basis-change}, and mode power law Eq.~\eqref{eq:ch3.mode-power-from-coefficient} enter
by reference \cite{belinfante1940,stratton1941,hansen1988}. Subsection~\ref{sec:ch392} does not re-derive Belinfante symmetrization
\cite{belinfante1940}.

\paragraph{Assumptions.}
Totals are evaluated on \(\mathcal{F}_{\mathrm{rad}}\) with Section~\ref{sec:ch35} normalization \cite{harrington2001,stratton1941}.
Local densities may be reported only when Eq.~\eqref{eq:ch1.observable-admissibility-triad} passes \cite{belinfante1940,stratton1941}.

\paragraph{Conserved versus representation split (HOME).}
Subsection~\ref{sec:ch392} partitions reported quantities as
\begin{equation}
\label{eq:ch3.conserved-vs-representation-split}
\underbrace{\mathcal{P}_{\mathrm{rad}},\ \oint h\,\mathrm{d}V}_{\text{conserved class}}
\;\oplus\;
\underbrace{\widehat{c}^{(\mathrm{pw})},\ \widehat{c}^{(\mathrm{cyl})},\ m_{\mathrm{LG}}}_{\text{representation-dependent class}},
\end{equation}
with the first group invariant under Eq.~\eqref{eq:ch1.gauge-transform} and the second covariant under
Eq.~\eqref{eq:ch3.basis-change-matrix} \cite{poynting1884,belinfante1940,harrington2001,hansen1988}. Equation
\eqref{eq:ch3.conserved-vs-representation-split} is the Subsection~\ref{sec:ch392} ledger law: watts and certified helicity fluxes are
\emph{physics totals}; raw chart coefficients are \emph{coordinates} \cite{stratton1941,harrington2001}.

\paragraph{Helicity sector conservation law (HOME).}
When Eq.~\eqref{eq:ch1.helicity-observable-decision} passes,
\begin{equation}
\label{eq:ch3.helicity-sector-conservation-law}
\oint_{\partial V} \mathbf{F}_{h}\cdot\hat{\mathbf{n}}\,\mathrm{d}A
=
\int_{V} s_{h}\,\mathrm{d}V
\;+\;
\text{source terms},
\end{equation}
importing Eq.~\eqref{eq:ch1.helicity-balance} with flux \(\mathbf{F}_{h}\) and source density \(s_{h}\) from Subsection~\ref{sec:ch01.4.3}
\cite{belinfante1940,stratton1941}. Equation~\eqref{eq:ch3.helicity-sector-conservation-law} states the physical meaning of helicity in the
eigenfield block: signed rotational overlap between \(\E\) and \(\H\) sectors, certified only when the Chapter~\ref{ch:01} gates pass
\cite{harrington2001,belinfante1940}. Helicity is not interchangeable with azimuthal index \(m\) in
Eq.~\eqref{eq:ch3.integer-ell-def} \cite{hansen1988,stratton1941}.

\paragraph{Representation-dependent ledger (HOME).}
\begin{equation}
\label{eq:ch3.representation-dependent-ledger}
\|\widehat{c}^{(\mathrm{a})}\|_{2}^{2}
=
\|\widehat{c}^{(\mathrm{b})}\|_{2}^{2}
\quad\text{but}\quad
\widehat{c}^{(\mathrm{a})}\neq \widehat{c}^{(\mathrm{b})}
\ \Longrightarrow\
\text{chart labels are not invariants},
\end{equation}
with \(\mathbf{T}_{b\leftarrow a}\) from Eq.~\eqref{eq:ch3.basis-change-matrix} \cite{hansen1988,harrington2001,stratton1941}. Equation
\eqref{eq:ch3.representation-dependent-ledger} prevents confusing Parseval invariance with label invariance: power is conserved, coefficient
\emph{vectors} are not \cite{poynting1884,harrington2001,belinfante1940}.

\begin{figure}[t]
\centering
\ChOneFigBlock{%
\begin{tikzpicture}[ch1fig, node distance=7mm and 9mm]
  \node[ch1box] (cons) {conserved totals};
  \node[ch1box, right=10mm of cons] (rep) {chart coefficients};
  \node[ch1box, right=10mm of rep] (T) {Eq.~\eqref{eq:ch3.basis-change-matrix}};
  \node[ch1box, below=10mm of cons] (hel) {Eq.~\eqref{eq:ch3.helicity-sector-conservation-law}};
  \draw[ch1flow] (cons) -- (rep) -- (T);
  \draw[ch1flow] (hel) -- (cons);
\end{tikzpicture}%
}
\caption{Subsection~\ref{sec:ch392} split: conserved flux totals versus representation-dependent coefficient vectors.}
\label{fig:ch3.392-conserved-split}
\end{figure}

\paragraph{Physical meaning of dual symmetry versus gauge.}
Dual rotations mix \(\E\) and \(\H\) while leaving source-free Maxwell dynamics invariant \cite{belinfante1940,stratton1941}; gauge
transforms relabel potentials without changing \(\E,\H\) \cite{stratton1941,belinfante1940}. Subsection~\ref{sec:ch392} assigns dual-symmetric
helicity reports to the conserved class only after Eq.~\eqref{eq:ch1.helicity-observable-decision}; gauge moves without dual symmetry stay in
the representation-dependent class \cite{harrington2001,belinfante1940}.

\paragraph{Worked illustration (same watts, different sparsity).}
A state may have sparse VSH coefficients yet dense plane-wave coefficients while \(\mathcal{P}_{\mathrm{rad}}\) is unchanged
Eq.~\eqref{eq:ch3.power-invariance-basis-change} \cite{hansen1988,harrington2001,stratton1941}. Subsection~\ref{sec:ch392} classifies sparsity
as a chart property, not a conserved quantum number \cite{harrington2001,jackson1999}.

\paragraph{Problem (classify helicity versus \(m\) label).}
\textbf{Problem.} A report lists helicity flux and ``\(m=1\)'' on the same line. Separate conserved from representation-dependent content
\cite{belinfante1940,harrington2001,stratton1941,hansen1988}.

\textbf{Step 1.} Test helicity with Eq.~\eqref{eq:ch1.helicity-observable-decision} \cite{belinfante1940}.

\textbf{Step 2.} Map \(m\) through Eq.~\eqref{eq:ch3.representation-dependent-ledger} \cite{hansen1988}.

\textbf{Physical meaning.} Only certified fluxes are conserved-class \cite{stratton1941,harrington2001}.

\begin{table}[t]
\centering
\caption{Subsection~\ref{sec:ch392} conserved versus representation-dependent taxonomy.}
\label{tab:ch3.392-taxonomy}
\small
\begin{tabular}{@{}p{0.28\linewidth}p{0.34\linewidth}p{0.32\linewidth}@{}}
\toprule
Quantity & Class & Test \\
\midrule
\(\mathcal{P}_{\mathrm{rad}}\) & conserved & Eq.~\eqref{eq:ch3.mode-power-from-coefficient} \\
Helicity flux & conserved & Eq.~\eqref{eq:ch1.helicity-observable-decision} \\
\(\widehat{c}_{\ell m}^{(\sigma)}\) & representation & Eq.~\eqref{eq:ch3.basis-change-matrix} \\
\(m_{\mathrm{LG}}\) & representation & Eq.~\eqref{eq:ch3.paraxial-vs-spherical-OAM} \\
\bottomrule
\end{tabular}
\end{table}

\paragraph{Belinfante total versus local splits.}
Integrated angular momentum totals from Eq.~\eqref{eq:ch1.belinfante-total-am-density} are conserved under the symmetric stress framework
\cite{belinfante1940,stratton1941}; local spin-orbit splits Eq.~\eqref{eq:ch1.local-spin-orbit-split} are representation-dependent unless
certified \cite{belinfante1940,harrington2001}. Subsection~\ref{sec:ch392} promotes totals to the conserved ledger and demotes local splits to
interpretive coordinates \cite{stratton1941,belinfante1940}.

\paragraph{Worked illustration (helicity sign flip under duality).}
Duality rotation can flip helicity sign while preserving \(\mathcal{P}_{\mathrm{rad}}\) \cite{belinfante1940,jackson1999,harrington2001}.
Subsection~\ref{sec:ch392} records helicity as a pseudoscalar conserved report, not as a second watt budget \cite{stratton1941,belinfante1940}.

\paragraph{Problem (audit a ``conserved OAM'' claim).}
\textbf{Problem.} A paper claims ``conserved OAM'' from a phase map. Apply Eq.~\eqref{eq:ch3.conserved-vs-representation-split}
\cite{belinfante1940,harrington2001,allen1992,stratton1941}.

\textbf{Step 1.} Demand flux closure or reject conserved class \cite{poynting1884}.

\textbf{Step 2.} Relabel phase maps as representation-dependent \cite{allen1992,harrington2001}.

\textbf{Physical meaning.} Phase winds are not automatic conserved charges \cite{belinfante1940,stratton1941}.

\begin{figure}[t]
\centering
\ChOneFigBlock{%
\begin{tikzpicture}[ch1fig, node distance=7mm and 9mm]
  \node[ch1box] (tot) {Belinfante totals};
  \node[ch1box, right=10mm of tot] (loc) {local splits};
  \node[ch1box, right=10mm of loc] (chart) {chart labels};
  \draw[ch1flow] (tot) -- (loc) -- (chart);
\end{tikzpicture}%
}
\caption{Subsection~\ref{sec:ch392} demotion ladder: integrated totals are conserved; local and chart labels are conditional.}
\label{fig:ch3.392-demotion-ladder}
\end{figure}

\paragraph{Helicity on eigenfield sectors.}
Helicity density Eq.~\eqref{eq:ch1.helicity-density} can be decomposed across VSH sectors only after full vector reconstruction
Eq.~\eqref{eq:ch3.normalized-coefficient-resolution} \cite{hansen1988,belinfante1940,harrington2001}. Subsection~\ref{sec:ch392} rejects
sector helicity slogans without coefficient vectors \cite{stratton1941,harrington2001}.

\paragraph{Worked illustration (circular polarization eigenfields).}
Circularly polarized \(\widehat{\mathbf{u}}_{\ell m}^{(\sigma)}\) templates carry definite helicity sign in the conserved class when
Eq.~\eqref{eq:ch1.helicity-density-wave-reduction} applies \cite{belinfante1940,jackson1999,hansen1988}. Changing \(\sigma\) label
notation without changing \(\mathbf{v}_\omega\) is representation bookkeeping, not physics \cite{harrington2001,stratton1941}.

\paragraph{Problem (compute conserved versus chart watts).}
\textbf{Problem.} Given \(\mathcal{P}_{\mathrm{rad}}=1\,\mathrm{W}\) and sparse \(\widehat{c}\) in VSH, show chart sparsity is not conserved
\cite{poynting1884,harrington2001,hansen1988,stratton1941}.

\textbf{Step 1.} Transform to plane-wave chart \cite{hansen1988}.

\textbf{Step 2.} Observe coefficient density change \cite{harrington2001}.

\textbf{Step 3.} Confirm \(\mathcal{P}_{\mathrm{rad}}\) unchanged \cite{poynting1884}.

\textbf{Physical meaning.} Conserved class is quadratic functionals, not sparsity \cite{stratton1941,harrington2001}.

\paragraph{Accessibility and conserved totals.}
Eq.~\eqref{eq:ch3.accessible-watt-budget} restricts which sectors enter conserved sums even when mathematically present in expansions
\cite{harrington2001,stratton1941}. Subsection~\ref{sec:ch392} applies that screen before promoting partial sector watts to conserved
narratives \cite{hansen1988,harrington2001}.

\begin{figure}[t]
\centering
\ChOneFigBlock{%
\begin{tikzpicture}[ch1fig, node distance=7mm and 9mm]
  \node[ch1box] (P) {$\mathcal{P}_{\mathrm{rad}}$};
  \node[ch1box, right=10mm of P] (acc) {Eq.~\eqref{eq:ch3.accessible-watt-budget}};
  \node[ch1box, right=10mm of acc] (hel) {Eq.~\eqref{eq:ch3.helicity-sector-conservation-law}};
  \draw[ch1flow] (P) -- (acc) -- (hel);
\end{tikzpicture}%
}
\caption{Subsection~\ref{sec:ch392} accessibility screen precedes conserved helicity promotion.}
\label{fig:ch3.392-accessibility-hel}
\end{figure}

\begin{table}[t]
\centering
\caption{Subsection~\ref{sec:ch392} promotion checklist for conserved reports.}
\label{tab:ch3.392-promotion-checklist}
\small
\begin{tabular}{@{}p{0.08\linewidth}p{0.50\linewidth}p{0.34\linewidth}@{}}
\toprule
Step & Requirement & Equation \\
\midrule
1 & Gauge screen & Eq.~\eqref{eq:ch1.observable-admissibility-triad} \\
2 & Watt closure & Eq.~\eqref{eq:ch3.mode-power-from-coefficient} \\
3 & Chart declared & Eq.~\eqref{eq:ch3.basis-change-matrix} \\
4 & Helicity gate & Eq.~\eqref{eq:ch1.helicity-observable-decision} \\
\bottomrule
\end{tabular}
\end{table}

\paragraph{Belinfante shell flux ledger.}
Conserved angular-momentum reports on an outgoing sphere \(S_{R}\) must use Belinfante--Maxwell totals imported from
Section~\ref{sec:ch01.4} rather than raw mode labels \cite{belinfante1940,stratton1941,harrington2001}. Subsection~\ref{sec:ch392} records
\begin{equation}
\label{eq:ch3.belinfante-shell-flux-ledger}
\mathbf{J}_{\mathrm{BF}}
=
\mathbf{J}_{\mathrm{orb}}
+
\mathbf{J}_{\mathrm{spin}},
\qquad
\Phi_{J}
=
\oint_{S_{R}}
\hat{\mathbf{r}}\cdot\mathbf{J}_{\mathrm{BF}}\,dS,
\end{equation}
with \(\mathbf{J}_{\mathrm{BF}}\) the symmetrized angular-momentum current and \(\Phi_{J}\) the outward flux through \(S_{R}\)
\cite{belinfante1940,jackson1999,stratton1941}. Equation~\eqref{eq:ch3.belinfante-shell-flux-ledger} is the conserved-class anchor: sector sums built
from Eq.~\eqref{eq:ch3.helicity-sector-conservation-law} must reconcile with \(\Phi_{J}\) within measurement tolerance, whereas chart-sparse headlines
such as ``\(m=1\) only'' need not \cite{harrington2001,belinfante1940}. Dual-symmetry sectors partition the watt budget as
\begin{equation}
\label{eq:ch3.dual-sector-watt-split}
\mathcal{P}_{\mathrm{rad}}
=
\mathcal{P}_{+}+\mathcal{P}_{-},
\qquad
\mathcal{P}_{\sigma}
=
\sum_{(\ell m)\in\mathcal{L}_{\mathrm{acc}}}
\bigl|\widehat{c}_{\ell m}^{(\sigma)}\bigr|^{2},
\end{equation}
importing helicity indices \(\sigma\) from Eq.~\eqref{eq:ch3.plane-wave-helicity-spectrum} after VSH transport
\cite{belinfante1940,hansen1988,harrington2001}. Equation~\eqref{eq:ch3.dual-sector-watt-split} shows that conserved totals are quadratic functionals
on \(\mathcal{F}_{\mathrm{rad}}\): transforming charts with Eq.~\eqref{eq:ch3.basis-change-matrix} may densify coefficients while leaving
\(\mathcal{P}_{\sigma}\) invariant \cite{poynting1884,stratton1941,harrington2001}.

\paragraph{Quadratic invariant transport across charts.}
Let \(Q(\mathbf{v}_\omega)\) denote any conserved quadratic functional admitted by Table~\ref{tab:ch3.392-taxonomy}. Subsection~\ref{sec:ch392}
requires
\begin{equation}
\label{eq:ch3.conserved-quadratic-transport}
Q(\mathbf{v}_\omega)
=
Q\!\left(\mathbf{T}_{b\leftarrow a}\,\mathbf{c}^{(\mathrm{a})}\right)
=
Q\!\left(\mathbf{c}^{(\mathrm{b})}\right),
\end{equation}
for unitary \(\mathbf{T}_{b\leftarrow a}\) on flux-normalized charts \cite{hansen1988,stratton1941,harrington2001}. Equation~\eqref{eq:ch3.conserved-quadratic-transport}
is the mathematical reason Belinfante totals and helicity-sector sums survive basis change while raw LG or cylindrical indices need not:
\(Q\) is defined on the field equivalence class, not on sparse coordinate displays \cite{belinfante1940,jackson1999}. Campaigns that export only
\(\mathcal{P}_{\mathrm{rad}}\) and \(\Phi_{J}\) therefore remain chart-robust; campaigns that export headline \(m\) or \(\ell\) without
Eq.~\eqref{eq:ch3.basis-change-matrix} do not \cite{harrington2001,allen1992}.

\paragraph{Peer-review rule.}
Conserved-class claims must cite a row of Table~\ref{tab:ch3.392-taxonomy} with passing equation tests; representation-dependent labels must
not appear in the same sentence without chart declaration \cite{belinfante1940,harrington2001,stratton1941,hansen1988}. Subsection~\ref{sec:ch392}
treats mixed wording as a classification error \cite{harrington2001,belinfante1940}.

\paragraph{Conservation laws versus chart sparsity.}
\emph{Conserved quantities} in Maxwell radiation belong to four-momentum and angular-momentum budgets certified on
outgoing shells; \emph{representation-dependent quantities} are chart-sparse encodings that change under basis change
\cite{belinfante1940,stratton1941,harrington2001}. Equation~\eqref{eq:ch3.conserved-vs-representation-split} encodes that
distinction as an admissibility rule: helicity-sector powers and Belinfante totals may be compared across charts only
after Eq.~\eqref{eq:ch3.basis-change-matrix} transport; raw coefficient magnitudes without \(\mathbf{T}\) are ledger
entries, not invariants \cite{belinfante1940,hansen1988,harrington2001}. Equation~\eqref{eq:ch3.helicity-sector-conservation-law}
embeds the physical content of sector conservation on \(\mathcal{R}_\omega\): when sources respect the assumptions of
Subsection~\ref{sec:ch01.4.3}, sector fluxes are observables in the same sense as total radiated power
\cite{belinfante1940,jackson1999}. Equation~\eqref{eq:ch3.representation-dependent-ledger} lists the complementary
vocabulary---OAM headlines, LG indices, paraxial vortex counts---as artifacts unless the bridge maps of
Section~\ref{sec:ch38} are attached \cite{allen1992,harrington2001,belinfante1940}. The implicit philosophy is
\emph{sparse charts do not create new conservation laws}; they reveal or hide coupling among modes already fixed by
Maxwell structure \cite{stratton1941,harrington2001,belinfante1940}.

\paragraph{Belinfante shell flux ledger (expanded arithmetic).}
When helicity gates pass, Subsection~\ref{sec:ch392} records sector fluxes through Eq.~\eqref{eq:ch3.belinfante-shell-flux-ledger} on a closed surface \(\partial V\) enclosing sources
\cite{belinfante1940,stratton1941,poynting1884}. The ledger is conserved under Eq.~\eqref{eq:ch3.basis-change-matrix} when \(Q\) in Eq.~\eqref{eq:ch3.conserved-quadratic-transport} is chosen from the conserved class of Table~\ref{tab:ch3.392-taxonomy}
\cite{hansen1988,harrington2001}. Representation-dependent headlines that disagree across charts therefore cannot overturn a certified Belinfante total obtained on \(\Gamma_R\)
\cite{belinfante1940,jackson1999,stratton1941}.

\paragraph{Dual-sector watt split under circular polarization.}
Circularly polarized eigenfields mix \(\mathbf{M}\) and \(\mathbf{N}\) tracks within fixed \(\ell\) \cite{belinfante1940,hansen1988}. Subsection~\ref{sec:ch392} imports Eq.~\eqref{eq:ch3.dual-sector-watt-split} as the executable split between dual templates before helicity-sector conservation is invoked
\cite{harrington2001,stratton1941,jackson1999}. Campaigns must publish \(\sigma\) resolution whenever helicity or dual-symmetry language appears in the same table as \(m\) labels
\cite{belinfante1940,harrington2001}.

\paragraph{Worked watt ledger (same state, two charts).}
Consider \(\mathcal{P}_{\mathrm{rad}}=1\,\mathrm{W}\) fixed on \(\Gamma_R\) with plane-wave coefficients \(\widehat{c}^{(\mathrm{pw})}\) sparse at broadside and VSH coefficients \(\widehat{c}^{(\mathrm{VSH})}\) dense across \(\ell\le3\)
\cite{hansen1988,harrington2001,stratton1941}. Eq.~\eqref{eq:ch3.representation-dependent-ledger} permits \(\widehat{c}^{(\mathrm{pw})}\neq\widehat{c}^{(\mathrm{VSH})}\) while \(\|\widehat{c}^{(\mathrm{pw})}\|_{2}^{2}=\|\widehat{c}^{(\mathrm{VSH})}\|_{2}^{2}\)
\cite{poynting1884,harrington2001}. Subsection~\ref{sec:ch392} uses that instance to train readers: conserved totals survive chart change; headline sparsity does not
\cite{belinfante1940,jackson1999,stratton1941}.

\paragraph{Helicity-sector audit versus $m$ labels.}
Helicity-sector flux Eq.~\eqref{eq:ch3.helicity-sector-conservation-law} is not equivalent to azimuthal index \(m\) in Eq.~\eqref{eq:ch3.integer-ell-def}
\cite{belinfante1940,hansen1988,stratton1941}. Subsection~\ref{sec:ch392} rejects manuscripts that infer conserved OAM from helicity totals without Eq.~\eqref{eq:ch1.helicity-observable-decision} and explicit \(\sigma\) resolution
\cite{harrington2001,belinfante1940,jackson1999}.

\paragraph{Flux transport under rigid rotations.}
Rigid rotations commute with conserved totals in Table~\ref{tab:ch3.392-taxonomy} when Eq.~\eqref{eq:ch3.conserved-quadratic-transport} is evaluated on \(\mathcal{F}_{\mathrm{rad}}\)
\cite{belinfante1940,jackson1999,harrington2001}. Subsection~\ref{sec:ch392} contrasts that covariance with representation-dependent sparsity: rotating the laboratory Wigner-mixes \(m\) within \(\ell\) while leaving \(\mathcal{P}_{\mathrm{rad}}\) fixed
\cite{hansen1988,stratton1941}. One must not infer conserved OAM from rotation-invariant watts without helicity gates
\cite{belinfante1940,harrington2001}.

\paragraph{Chart watt ledger (numeric sketch).}
With \(\mathcal{P}_{\mathrm{rad}}=1\,\mathrm{W}\) on \(\Gamma_R\), plane-wave chart coefficients may show \(0.85\,\mathrm{W}\) in a single \((k_{x},k_{y})\) cell while VSH coefficients spread the same watt across \(\ell\le3\)
\cite{hansen1988,jackson1999,harrington2001}. Eq.~\eqref{eq:ch3.representation-dependent-ledger} permits that sparsity difference; Eq.~\eqref{eq:ch3.conserved-vs-representation-split} forbids interpreting sparsity as a new conservation law
\cite{poynting1884,belinfante1940,stratton1941}.

\paragraph{Transport of Belinfante totals across $\mathbf{T}$.}
When Eq.~\eqref{eq:ch1.helicity-observable-decision} passes, Belinfante-sector totals in Eq.~\eqref{eq:ch3.helicity-sector-conservation-law} transport under unitary \(\mathbf{T}\) via Eq.~\eqref{eq:ch3.conserved-quadratic-transport}
\cite{belinfante1940,hansen1988,harrington2001}. Subsection~\ref{sec:ch392} uses that transport law to reject mixed sentences that equate helicity flux with azimuthal index \(m\) or LG label \(m_{\mathrm{LG}}\)
\cite{allen1992,stratton1941,jackson1999}. Conserved language must cite a row of Table~\ref{tab:ch3.392-taxonomy}; representation language must cite Eq.~\eqref{eq:ch3.basis-change-matrix}
\cite{harrington2001,belinfante1940}.

Subsection~\ref{sec:ch392} has separated conserved from representation-dependent quantities through
Eqs.~\eqref{eq:ch3.conserved-vs-representation-split}--\eqref{eq:ch3.representation-dependent-ledger}
\cite{belinfante1940,stratton1941,harrington2001}. Subsection~\ref{sec:ch393} records scalar and paraxial limits where vector theory reduces
\cite{jackson1999,allen1992}.


\subsection{Scalar and paraxial limits}
\label{sec:ch393}
% TOC tag: [HOME]

\emph{Scalar limits} mean consistent reduction of vector eigenfield theory to a single transverse component obeying a scalar Helmholtz
equation \cite{jackson1999,stratton1941,harrington2001}. \emph{Paraxial limits} mean slowly varying envelope approximations where
longitudinal field components and depolarization are controlled by a small parameter \cite{allen1992,jackson1999}. Subsection~\ref{sec:ch393}
is the HOME validity layer that states when Laguerre--Gaussian and scalar diffraction formulas from popular OAM discourse agree with full
vector VSH reports \cite{allen1992,harrington2001,stratton1941,hansen1988}.

\paragraph{Standing imports.}
Paraxial versus spherical OAM link Eq.~\eqref{eq:ch3.paraxial-vs-spherical-OAM}, integer index discipline
Eq.~\eqref{eq:ch3.integer-ell-def}, normalized VSH resolution Eq.~\eqref{eq:ch3.normalized-coefficient-resolution}, and gauge screen
Eq.~\eqref{eq:ch3.gauge-screen-OAM-claim} enter by reference \cite{allen1992,hansen1988,harrington2001}. Subsection~\ref{sec:ch393} does not
re-derive paraxial wave equations from Maxwell's system \cite{stratton1941,jackson1999}.

\paragraph{Assumptions.}
Introduce paraxial parameter \(\epsilon\ll1\) measuring beam divergence \(k_{t}/k_{z}\) \cite{jackson1999,allen1992}. Scalar reduction
requires polarization track fixed and weak depolarization on the propagation path \cite{stratton1941,harrington2001}.

\paragraph{Paraxial reduction condition (HOME).}
Subsection~\ref{sec:ch393} fixes
\begin{equation}
\label{eq:ch3.paraxial-reduction-condition}
\epsilon
=
\frac{k_{t}}{k_{z}}
\ll 1,
\qquad
\bigl|\E_{z}\bigr|\ll \bigl|\E_{t}\bigr|,
\qquad
\bigl|\partial_{z}\mathbf{u}\bigr|\ll k_{0}\,\bigl|\mathbf{u}\bigr|,
\end{equation}
with \(\mathbf{u}\) a transverse field component and \(k_{t}=\sqrt{k_{x}^{2}+k_{y}^{2}}\) \cite{jackson1999,allen1992,harrington2001}. Equation
\eqref{eq:ch3.paraxial-reduction-condition} is the Subsection~\ref{sec:ch393} validity certificate: paraxial LG formulas apply only inside this
inequality chain \cite{allen1992,stratton1941}. The reduction is a \emph{controlled contraction} of the vector Helmholtz problem---not an alternate
physics---and every coefficient report must carry \(\epsilon\) or \(f/w_{0}\) alongside LG indices so that Theorem~\ref{thm:ch3.paraxial-consistent-reduction}
can be invoked without overclaiming \cite{allen1992,harrington2001,belinfante1940}. Outside Eq.~\eqref{eq:ch3.paraxial-reduction-condition}, vector VSH
reconstruction Eq.~\eqref{eq:ch3.normalized-coefficient-resolution} is mandatory for watt reports \cite{hansen1988,harrington2001}.

\paragraph{Scalar Helmholtz limit (HOME).}
When Eq.~\eqref{eq:ch3.paraxial-reduction-condition} holds and polarization is uniform,
\begin{equation}
\label{eq:ch3.scalar-helmholtz-limit}
(\nabla_{t}^{2}+k_{z}^{2})\,u(\mathbf{r}_{t},z)
=
0,
\qquad
\mathbf{E}_{t}\approx u\,\hat{\mathbf{e}}_{0},
\end{equation}
with transverse Laplacian \(\nabla_{t}^{2}\) and fixed polarization \(\hat{\mathbf{e}}_{0}\) \cite{jackson1999,stratton1941,allen1992}. Equation
\eqref{eq:ch3.scalar-helmholtz-limit} states the physical meaning of scalar limits: transverse dynamics decouple from longitudinal Maxwell
coupling to leading order in \(\epsilon\) \cite{harrington2001,jackson1999}. The scalar mode is a \emph{single-component proxy}, not a
gauge-invariant replacement for \(\mathbf{v}_\omega\) \cite{belinfante1940,harrington2001}.

\paragraph{Paraxial validity boundary (HOME).}
\begin{equation}
\label{eq:ch3.paraxial-validity-boundary}
f
=
\frac{\pi w_{0}^{2}}{\lambda}
\gg w_{0},
\qquad
\mathcal{N}_{\mathrm{par}}
=
\frac{1}{\epsilon^{2}},
\end{equation}
with Rayleigh range \(f\), waist \(w_{0}\), wavelength \(\lambda\), and paraxial figure \(\mathcal{N}_{\mathrm{par}}\) \cite{jackson1999,allen1992}.
Equation~\eqref{eq:ch3.paraxial-validity-boundary} supplies the Subsection~\ref{sec:ch393} focus metric: when \(f\sim w_{0}\), paraxial
expansions fail and vector corrections grow as \(\mathcal{O}(\epsilon)\) \cite{allen1992,harrington2001,stratton1941}. Laboratory claims
using LG indices must report \(\epsilon\) or \(f/w_{0}\) alongside \(m_{\mathrm{LG}}\) \cite{harrington2001,jackson1999}.

\begin{theorem}[Paraxial-consistent reduction (HOME)]
\label{thm:ch3.paraxial-consistent-reduction}
If Eqs.~\eqref{eq:ch3.paraxial-reduction-condition}--\eqref{eq:ch3.scalar-helmholtz-limit} hold, LG modes approximate solutions of
Eq.~\eqref{eq:ch3.scalar-helmholtz-limit} and map to VSH coefficients through Eq.~\eqref{eq:ch3.paraxial-vs-spherical-OAM} with remainder
\(\mathcal{O}(\epsilon)\) in \(\mathcal{P}_{\mathrm{rad}}\) \cite{allen1992,jackson1999,hansen1988,harrington2001}.
\end{theorem}
\begin{proof}[Proof sketch]
Standard paraxial expansion \cite{allen1992,jackson1999} plus Subsection~\ref{sec:ch383} transformation bounds \cite{hansen1988,harrington2001}.
% PROOF: AUTHOR REVIEW
\end{proof}

\begin{figure}[t]
\centering
\ChOneFigBlock{%
\begin{tikzpicture}[ch1fig, node distance=7mm and 9mm]
  \node[ch1box] (max) {Maxwell vector};
  \node[ch1box, right=10mm of max] (par) {Eq.~\eqref{eq:ch3.paraxial-reduction-condition}};
  \node[ch1box, right=10mm of par] (sc) {Eq.~\eqref{eq:ch3.scalar-helmholtz-limit}};
  \node[ch1box, below=10mm of par] (vsh) {Eq.~\eqref{eq:ch3.normalized-coefficient-resolution}};
  \draw[ch1flow] (max) -- (par) -- (sc);
  \draw[ch1flow] (par) -- (vsh);
\end{tikzpicture}%
}
\caption{Subsection~\ref{sec:ch393} reduction chain: paraxial scalar limits are approximate; vector VSH remains the far-zone anchor.}
\label{fig:ch3.393-reduction-chain}
\end{figure}

\paragraph{Physical meaning of paraxial LG modes.}
Laguerre--Gaussian modes are convenient \emph{scalar} eigensolutions of Eq.~\eqref{eq:ch3.scalar-helmholtz-limit} with azimuthal index
\(m_{\mathrm{LG}}\) \cite{allen1992,jackson1999}. Their vortex phase winds are faithful only where
Eq.~\eqref{eq:ch3.paraxial-reduction-condition} holds \cite{allen1992,harrington2001}. Far-zone vector content generally spreads across
multiple \(\ell\) in Eq.~\eqref{eq:ch3.integer-ell-def} per Eq.~\eqref{eq:ch3.paraxial-vs-spherical-OAM} \cite{hansen1988,stratton1941}.

\paragraph{Worked illustration (Gaussian beam far zone).}
A fundamental Gaussian satisfies Eq.~\eqref{eq:ch3.paraxial-reduction-condition} near waist with small \(\epsilon\) \cite{jackson1999,allen1992}.
Its far-zone VSH spectrum is dominated by low \(\ell\) but not strictly single-mode \cite{hansen1988,harrington2001}. Subsection~\ref{sec:ch393}
uses that spread as a scalar-limit warning \cite{stratton1941,jackson1999}.

\paragraph{Problem (estimate \(\epsilon\) for a beam).}
\textbf{Problem.} Given half-angle \(\theta_{1/2}\), estimate \(\epsilon\approx\sin\theta_{1/2}\) and state whether
Eq.~\eqref{eq:ch3.paraxial-reduction-condition} applies \cite{jackson1999,allen1992,harrington2001,stratton1941}.

\textbf{Step 1.} Map \(\theta_{1/2}\) to \(k_{t}/k_{0}\) \cite{jackson1999}.

\textbf{Step 2.} Compare to protocol threshold \(\epsilon_{\max}\) \cite{harrington2001}.

\textbf{Physical meaning.} Wide-angle beams exit paraxial validity \cite{allen1992,stratton1941}.

\begin{table}[t]
\centering
\caption{Subsection~\ref{sec:ch393} scalar and paraxial limit objects.}
\label{tab:ch3.393-limit-objects}
\small
\begin{tabular}{@{}p{0.26\linewidth}p{0.38\linewidth}p{0.30\linewidth}@{}}
\toprule
Object & Equation & Role \\
\midrule
Paraxial parameter & Eq.~\eqref{eq:ch3.paraxial-reduction-condition} & validity \\
Scalar wave & Eq.~\eqref{eq:ch3.scalar-helmholtz-limit} & reduced dynamics \\
Focus metric & Eq.~\eqref{eq:ch3.paraxial-validity-boundary} & failure boundary \\
Vector anchor & Eq.~\eqref{eq:ch3.normalized-coefficient-resolution} & far-zone truth \\
\bottomrule
\end{tabular}
\end{table}

\paragraph{Vector correction beyond paraxial.}
When \(\epsilon\) is not small, longitudinal \(\E_{z}\) and depolarization feed \(\mathcal{O}(\epsilon)\) corrections to
\(\mathcal{P}_{\mathrm{rad}}\) \cite{jackson1999,harrington2001,stratton1941}. Subsection~\ref{sec:ch393} requires reporting those corrections
or switching fully to vector VSH before watt claims \cite{hansen1988,harrington2001,belinfante1940}.

\paragraph{Worked illustration (tight focus failure).}
Near \(f\sim w_{0}\), Eq.~\eqref{eq:ch3.paraxial-validity-boundary} fails; scalar LG winds mis-predict far-zone \(\widehat{c}_{\ell m}^{(\sigma)}\)
\cite{allen1992,jackson1999,harrington2001}. High-NA optics must use vector fields throughout \cite{stratton1941,hansen1988}.

\paragraph{Problem (scalar-to-vector uplift).}
\textbf{Problem.} Given scalar \(u\) satisfying Eq.~\eqref{eq:ch3.scalar-helmholtz-limit}, outline uplift to
Eq.~\eqref{eq:ch3.normalized-coefficient-resolution} \cite{hansen1988,stratton1941,harrington2001,allen1992}.

\textbf{Step 1.} Fix \(\hat{\mathbf{e}}_{0}\) and reconstruct \(\mathbf{E}_{t}\) \cite{stratton1941}.

\textbf{Step 2.} Apply \(\Pi_{\mathrm{rad}}\) \cite{harrington2001}.

\textbf{Step 3.} Project VSH coefficients \cite{hansen1988}.

\textbf{Physical meaning.} Scalar solutions are uplift templates, not completed vector states \cite{belinfante1940,harrington2001}.

\begin{figure}[t]
\centering
\ChOneFigBlock{%
\begin{tikzpicture}[ch1fig, node distance=7mm and 9mm]
  \node[ch1box] (eps) {$\epsilon$};
  \node[ch1box, right=10mm of eps] (ok) {paraxial OK};
  \node[ch1box, right=10mm of ok] (fail) {vector required};
  \draw[ch1flow] (eps) -- (ok);
  \draw[ch1flow] (eps) -- (fail);
\end{tikzpicture}%
}
\caption{Subsection~\ref{sec:ch393} validity fork: small \(\epsilon\) permits scalar LG; large \(\epsilon\) requires full vector reconstruction.}
\label{fig:ch3.393-validity-fork}
\end{figure}

\paragraph{Scalar limit and gauge.}
Scalar reductions fix polarization direction \(\hat{\mathbf{e}}_{0}\), removing part of the gauge-dual structure tracked in
Subsection~\ref{sec:ch392} \cite{belinfante1940,harrington2001}. Subsection~\ref{sec:ch393} therefore tags scalar limits as
\emph{approximation coordinates}, not as gauge-invariant subtheories \cite{stratton1941,belinfante1940}.

\paragraph{Worked illustration (LG\(_{0}^{1}\) winding).}
In valid paraxial domains, LG\(_{0}^{1}\) exhibits \(m_{\mathrm{LG}}=1\) phase wind \cite{allen1992,jackson1999}. Topological diagnostic
\(q_{\mathrm{top}}\) from Subsection~\ref{sec:ch374} may agree on declared loops yet VSH vector still shows multi-\(\ell\) content
\cite{hansen1988,harrington2001,stratton1941}.

\paragraph{Problem (Rayleigh range audit).}
\textbf{Problem.} For \(w_{0}=1\,\mathrm{mm}\), \(\lambda=30\,\mathrm{mm}\), compute \(f\) and test
Eq.~\eqref{eq:ch3.paraxial-validity-boundary} \cite{jackson1999,allen1992,harrington2001}.

\textbf{Step 1.} \(f=\pi w_{0}^{2}/\lambda\approx0.105\,\mathrm{mm}\) \cite{jackson1999}.

\textbf{Step 2.} \(f\ll w_{0}\) fails the paraxial focus metric \cite{allen1992}.

\textbf{Physical meaning.} Long-wavelength beams can leave paraxial validity quickly \cite{harrington2001,stratton1941}.

\paragraph{Handoff to far-zone VSH.}
Subsection~\ref{sec:ch393} ends paraxial validity at the shell where Eq.~\eqref{eq:ch3.paraxial-reduction-condition} breaks; beyond that,
Eq.~\eqref{eq:ch3.normalized-coefficient-resolution} is the reporting standard \cite{hansen1988,harrington2001,stratton1941}. Campaign archives
should record the handoff radius or angle \cite{harrington2001,jackson1999}.

\begin{figure}[t]
\centering
\ChOneFigBlock{%
\begin{tikzpicture}[ch1fig, node distance=7mm and 9mm]
  \node[ch1box] (near) {near paraxial};
  \node[ch1box, right=10mm of near] (hand) {handoff};
  \node[ch1box, right=10mm of hand] (far) {VSH far zone};
  \draw[ch1flow] (near) -- (hand) -- (far);
\end{tikzpicture}%
}
\caption{Subsection~\ref{sec:ch393} handoff: paraxial scalar charts yield to vector VSH on the far observation shell.}
\label{fig:ch3.393-handoff}
\end{figure}

\begin{table}[t]
\centering
\caption{Subsection~\ref{sec:ch393} paraxial campaign metadata.}
\label{tab:ch3.393-campaign-metadata}
\small
\begin{tabular}{@{}p{0.34\linewidth}p{0.56\linewidth}@{}}
\toprule
Field & Content \\
\midrule
\(\epsilon,\ f/w_{0}\) & Eq.~\eqref{eq:ch3.paraxial-validity-boundary} \\
\(w_{0},\lambda\) & beam geometry \\
\(m_{\mathrm{LG}}\) & scalar index only \\
VSH uplift & Eq.~\eqref{eq:ch3.normalized-coefficient-resolution} \\
\bottomrule
\end{tabular}
\end{table}

\paragraph{Paraxial-to-VSH coupling coefficients (HOME).}
When Eq.~\eqref{eq:ch3.paraxial-reduction-condition} holds, LG modes map to VSH coefficients through the bridge already recorded in
Eq.~\eqref{eq:ch3.paraxial-vs-spherical-OAM} \cite{allen1992,hansen1988}. Subsection~\ref{sec:ch393} states the coupling remainder law
\begin{equation}
\label{eq:ch3.paraxial-vsh-coupling-remainder}
\widehat{c}_{\ell m}^{(\sigma)}
=
\mathcal{T}_{\mathrm{par}\to\mathrm{VSH}}\bigl[m_{\mathrm{LG}},p_{\mathrm{LG}}\bigr]_{\ell m}^{(\sigma)}
+
\mathcal{O}(\epsilon),
\qquad
\sum_{\ell m\sigma}\bigl|\widehat{c}_{\ell m}^{(\sigma)}\bigr|^{2}
=
\mathcal{P}_{\mathrm{rad}}+\mathcal{O}(\epsilon^{2}),
\end{equation}
with \(\mathcal{T}_{\mathrm{par}\to\mathrm{VSH}}\) the fixed transformation induced by Eq.~\eqref{eq:ch3.paraxial-vs-spherical-OAM},
Laguerre index \(p_{\mathrm{LG}}\), and azimuthal index \(m_{\mathrm{LG}}\) \cite{allen1992,jackson1999,harrington2001}. Equation
\eqref{eq:ch3.paraxial-vsh-coupling-remainder} is the quantitative face of Theorem~\ref{thm:ch3.paraxial-consistent-reduction}: paraxial LG
labels are single-chart coordinates, while far-zone watts distribute across multiple integer \(\ell\) sectors at \(\mathcal{O}(\epsilon)\) beyond
the paraxial inequality chain \cite{hansen1988,stratton1941}. A beam with \(m_{\mathrm{LG}}=1\) may therefore show \(q_{\mathrm{top}}=1\) on a
declared loop while \(\widehat{c}_{\ell m}^{(\sigma)}\) still exhibits \(\ell>1\) admixture once vector uplift is performed
\cite{allen1992,harrington2001,belinfante1940}.

\paragraph{Depolarization and longitudinal feed at order \(\epsilon\).}
When \(\bigl|\E_{z}\bigr|\) is no longer negligible, uniform polarization \(\hat{\mathbf{e}}_{0}\) in Eq.~\eqref{eq:ch3.scalar-helmholtz-limit} fails
and \(\mathcal{O}(\epsilon)\) corrections enter both \(\mathcal{P}_{\mathrm{rad}}\) and helicity sectors \cite{jackson1999,harrington2001,stratton1941}.
Subsection~\ref{sec:ch393} treats that failure as a \emph{vector uplift trigger}: once \(\epsilon\) exceeds the campaign threshold, scalar LG fits must
be abandoned for Eq.~\eqref{eq:ch3.normalized-coefficient-resolution} on \(\Gamma_R\) \cite{hansen1988,harrington2001}. High-NA focusing and compact
resonant feeds are the common laboratory paths to that trigger \cite{stratton1941,jackson1999}.

\paragraph{Paraxial LG mode family in scalar language.}
Laguerre--Gaussian modes are eigenfunctions of Eq.~\eqref{eq:ch3.scalar-helmholtz-limit} with azimuthal index \(m_{\mathrm{LG}}\) and radial index
\(p_{\mathrm{LG}}\) \cite{allen1992,jackson1999}. Their vortex structure is faithful only when Eqs.~\eqref{eq:ch3.paraxial-reduction-condition} and
\eqref{eq:ch3.paraxial-validity-boundary} hold simultaneously \cite{allen1992,harrington2001}. Subsection~\ref{sec:ch393} therefore treats LG modes as
\emph{controlled proxies}: they accelerate near-field intuition and laser-cavity design, but far-zone watt and angular-momentum reports require uplift to
Eq.~\eqref{eq:ch3.normalized-coefficient-resolution} on \(\Gamma_R\) \cite{hansen1988,stratton1941,belinfante1940}.

\paragraph{Gaussian-beam uplift workflow (stepwise).}
Begin from scalar \(u\) satisfying Eq.~\eqref{eq:ch3.scalar-helmholtz-limit} with fixed \(\hat{\mathbf{e}}_{0}\) \cite{stratton1941,jackson1999}.
Reconstruct \(\mathbf{E}_{t}\approx u\hat{\mathbf{e}}_{0}\), apply \(\Pi_{\mathrm{rad}}\), and project VSH coefficients on \(\Gamma_R\)
\cite{harrington2001,hansen1988}. Compare uplifted \(\widehat{c}_{\ell m}^{(\sigma)}\) to paraxial map
Eq.~\eqref{eq:ch3.paraxial-vsh-coupling-remainder}; remainder \(\mathcal{O}(\epsilon)\) certifies Theorem~\ref{thm:ch3.paraxial-consistent-reduction}
\cite{allen1992,harrington2001}. Archive \(\epsilon\), \(w_{0}\), \(f\), and uplift residuals with every far-zone claim originating in LG fits
\cite{jackson1999,stratton1941}.

\paragraph{High-NA and tight-focus exit ramp.}
When \(\epsilon\) is not small, longitudinal \(\E_{z}\) and depolarization invalidate Eq.~\eqref{eq:ch3.scalar-helmholtz-limit}
\cite{jackson1999,harrington2001,stratton1941}. Subsection~\ref{sec:ch393} mandates full vector reconstruction before any \(\mathcal{P}_{\mathrm{rad}}\) or
helicity statement \cite{belinfante1940,hansen1988}. Figure~\ref{fig:ch3.393-validity-fork} and Figure~\ref{fig:ch3.393-handoff} together define the exit ramp:
inside paraxial domain use scalar certificates; on the far shell use VSH meters \cite{allen1992,harrington2001,jackson1999}.

\paragraph{Paraxial campaign acceptance inequality chain.}
Subsection~\ref{sec:ch393} packages the acceptance test as
\begin{equation}
\label{eq:ch3.paraxial-campaign-acceptance}
\mathcal{A}_{\mathrm{par}}=1
\iff
\left(
\epsilon\le\epsilon_{\max}\
\wedge\
f/w_{0}\gg 1\
\wedge\
\|\mathbf{u}_{\mathrm{uplift}}-\mathbf{u}_{\mathrm{LG}}\|_{\mathrm{rad}}\le \tau_{\mathrm{uplift}}
\right),
\end{equation}
with \(\epsilon\) from Eq.~\eqref{eq:ch3.paraxial-reduction-condition}, \(f/w_{0}\) from Eq.~\eqref{eq:ch3.paraxial-validity-boundary}, and uplift residual \(\tau_{\mathrm{uplift}}\) comparing vector VSH reconstruction to the LG fit on \(\Gamma_R\)
\cite{allen1992,hansen1988,harrington2001,stratton1941}. Equation~\eqref{eq:ch3.paraxial-campaign-acceptance} prevents paraxial vocabulary from outrunning vector validity in far-zone reports
\cite{belinfante1940,jackson1999,harrington2001}.

\paragraph{Numeric instance (\(\epsilon\) audit).}
For half-angle \(\theta_{1/2}=8^{\circ}\), estimate \(\epsilon\approx\sin\theta_{1/2}\approx0.14\) \cite{jackson1999,allen1992}. With \(\epsilon_{\max}=0.1\) in a campaign protocol, Eq.~\eqref{eq:ch3.paraxial-campaign-acceptance} fails and vector uplift to Eq.~\eqref{eq:ch3.normalized-coefficient-resolution} is mandatory on \(\Gamma_R\)
\cite{harrington2001,hansen1988,stratton1941}. Subsection~\ref{sec:ch393} records such audits as numeric companions to Theorem~\ref{thm:ch3.paraxial-consistent-reduction}, not as optional remarks
\cite{allen1992,belinfante1940,jackson1999}.

\paragraph{Rayleigh-range failure instance (numeric).}
With \(w_{0}=1\,\mathrm{mm}\) and \(\lambda=30\,\mathrm{mm}\), \(f=\pi w_{0}^{2}/\lambda\approx0.105\,\mathrm{mm}\) so \(f/w_{0}\ll1\) and Eq.~\eqref{eq:ch3.paraxial-validity-boundary} fails \cite{jackson1999,allen1992}. Subsection~\ref{sec:ch393} uses that instance to show long-wavelength beams can exit paraxial validity quickly even when near-field LG plots look smooth
\cite{harrington2001,stratton1941,belinfante1940}. Vector uplift Eq.~\eqref{eq:ch3.normalized-coefficient-resolution} is then mandatory on \(\Gamma_R\)
\cite{hansen1988,harrington2001}.

\paragraph{Peer-review rule for paraxial OAM.}
Paraxial OAM claims must publish \(\epsilon\) or \(f/w_{0}\), cite Eq.~\eqref{eq:ch3.paraxial-reduction-condition}, and attach vector uplift
coefficients for far-zone statements \cite{allen1992,harrington2001,stratton1941,hansen1988}. Subsection~\ref{sec:ch393} rejects far-zone watt
claims based solely on LG fits \cite{harrington2001,belinfante1940}.

Subsection~\ref{sec:ch393} has recorded scalar and paraxial limits through
Eqs.~\eqref{eq:ch3.paraxial-reduction-condition}--\eqref{eq:ch3.paraxial-validity-boundary} and
Theorem~\ref{thm:ch3.paraxial-consistent-reduction} \cite{allen1992,jackson1999,harrington2001}. Subsection~\ref{sec:ch394} contrasts
observables with representation artifacts using Chapter~\ref{ch:01} inference rules \cite{harrington2001}.


\subsection{Observables versus representation artifacts}
\label{sec:ch394}
% TOC tag: [HOME][USE SS1.5.4]

An \emph{observable} is a quantity admitted by calibrated channel maps and upstream gates with documented uncertainty
\cite{harrington2001,stratton1941,belinfante1940}. A \emph{representation artifact} is a visually compelling or algebraically convenient
label that changes under basis change Eq.~\eqref{eq:ch3.basis-change-matrix} or fails inference tests
Eq.~\eqref{eq:ch1.observable-decision-rule} \cite{hansen1988,harrington2001,allen1992}. Subsection~\ref{sec:ch394} is the HOME closure that
imports Subsection~\ref{sec:ch01.5.4} decision rules onto the eigenfield language of Sections~\ref{sec:ch34}--\ref{sec:ch39}
\cite{harrington2001,stratton1941,belinfante1940}.

\paragraph{Standing imports.}
Direct versus inferred split Eq.~\eqref{eq:ch1.direct-inferred-role-split}, observable inference criterion
Eq.~\eqref{eq:ch1.observable-inference-criterion}, model mismatch index Eq.~\eqref{eq:ch1.inference-model-mismatch-index},
observable decision rule Eq.~\eqref{eq:ch1.observable-decision-rule}, observable report contract
Eq.~\eqref{eq:ch1.observable-report-contract}, representation dependence Eq.~\eqref{eq:ch3.OAM-representation-dependence}, and gauge screen
Eq.~\eqref{eq:ch3.gauge-screen-OAM-claim} enter by reference \cite{harrington2001,stratton1941,belinfante1940,allen1992}. Subsection~\ref{sec:ch394}
does not re-derive channel covariance models \cite{harrington2001}.

\paragraph{Assumptions.}
Channel noise floor \(\sigma_d\), thresholds \(\tau_A,M_{\max}\), and symmetry penalty \(\chi_{\mathrm{sym}}\) are fixed before inversion
\cite{harrington2001,stratton1941}. Eigenfield coefficients are inferred quantities unless measured by direct flux calibration
\cite{poynting1884,harrington2001,hansen1988}.

\paragraph{Observable versus artifact screen (HOME).}
Subsection~\ref{sec:ch394} defines
\begin{equation}
\label{eq:ch3.observable-vs-artifact-screen}
\mathfrak{S}_{\mathrm{obs}}
=
\bigl\{\mathcal{P}_{\mathrm{rad}},\ \text{certified fluxes},\ \mathcal{O}_d=\mathcal{M}[\mathcal{F}]\bigr\},
\qquad
\mathfrak{S}_{\mathrm{art}}
=
\bigl\{\text{un-gated }m_{\mathrm{LG}},\ \text{raw FFT peaks},\ \text{scalar }\ell_{\mathrm{fit}}\bigr\},
\end{equation}
importing Eq.~\eqref{eq:ch1.direct-inferred-role-split} \cite{harrington2001,stratton1941,belinfante1940}. Equation
\eqref{eq:ch3.observable-vs-artifact-screen} is the Subsection~\ref{sec:ch394} classifier: \(\mathfrak{S}_{\mathrm{obs}}\) items survive
measurement and gate chains; \(\mathfrak{S}_{\mathrm{art}}\) items require promotion or must be labeled inference-only
\cite{harrington2001,allen1992,hansen1988}.

\paragraph{Inference artifact rejection (HOME).}
\begin{equation}
\label{eq:ch3.inference-artifact-rejection}
\mathfrak{A}(\mathcal{Q}_i)>\tau_A
\ \vee\
\mathcal{M}_{\mathrm{mis}}>M_{\max}
\;\Longrightarrow\;
\mathcal{Q}_i\in\mathfrak{S}_{\mathrm{art}},
\end{equation}
with \(\mathfrak{A},\mathcal{M}_{\mathrm{mis}}\) from Eqs.~\eqref{eq:ch1.observable-inference-criterion}--\eqref{eq:ch1.inference-model-mismatch-index}
\cite{harrington2001,stratton1941}. Equation~\eqref{eq:ch3.inference-artifact-rejection} operationalizes artifact rejection: high inversion
sensitivity or model mismatch demotes a descriptor even if plots look structured \cite{harrington2001,devaney2004,hansen1988}.

\paragraph{Reportable eigenfield observable set (HOME).}
\begin{equation}
\label{eq:ch3.reportable-eigenfield-observable-set}
\mathcal{O}_{\mathrm{rep}}
=
\bigl\{\widehat{c}_{\ell m}^{(\sigma)}:\ \mathcal{G}_{\mathrm{eig}}=1\bigr\}
\ \cup\
\{\mathcal{P}_{\mathrm{rad}}\},
\end{equation}
with \(\mathcal{G}_{\mathrm{eig}}=1\) denoting passage of Eqs.~\eqref{eq:ch3.gauge-qualified-eigenfield-report} and~\eqref{eq:ch1.observable-decision-rule},
and accessibility Eq.~\eqref{eq:ch3.accessible-watt-budget} enforced on accessible sectors \cite{harrington2001,stratton1941,hansen1988,belinfante1940}.
Equation~\eqref{eq:ch3.reportable-eigenfield-observable-set} is the Subsection~\ref{sec:ch394} export list for Chapter~\ref{ch:03}: only
coefficients passing gauge, inference, and accessibility gates are observables; others remain artifacts or constrained surrogates
\cite{harrington2001,belinfante1940,stratton1941}.

\begin{figure}[t]
\centering
\ChOneFigBlock{%
\begin{tikzpicture}[ch1fig, node distance=7mm and 9mm]
  \node[ch1box] (data) {channel data};
  \node[ch1box, right=10mm of data] (inf) {inversion $\mathcal{Q}_i$};
  \node[ch1box, right=10mm of inf] (rule) {Eq.~\eqref{eq:ch1.observable-decision-rule}};
  \node[ch1box, right=10mm of rule] (obs) {$\mathcal{O}_{\mathrm{rep}}$};
  \draw[ch1flow] (data) -- (inf) -- (rule) -- (obs);
\end{tikzpicture}%
}
\caption{Subsection~\ref{sec:ch394} inference gate: reconstructed eigenfield labels become observables only after decision-rule passage.}
\label{fig:ch3.394-inference-gate}
\end{figure}

\paragraph{Physical meaning of representation artifacts.}
Artifacts are not ``wrong data''; they are \emph{pre-observable} descriptors whose numerical values depend on chart, gauge, regularization, or
plot choice \cite{hansen1988,harrington2001,allen1992}. A beautiful phase vortex image is an artifact until decomposed into
Eq.~\eqref{eq:ch3.reportable-eigenfield-observable-set} with passing gates \cite{belinfante1940,harrington2001,stratton1941}. Subsection~\ref{sec:ch394}
therefore trains readers to ask \emph{which screen passed}, not \emph{which picture appeared} \cite{harrington2001,devaney2004}.

\paragraph{Worked illustration (NFF to VSH inversion).}
Near-field scans inverted to \(\widehat{c}_{\ell m}^{(\sigma)}\) produce \(\mathcal{Q}_i\) with \(\kappa_{\mathrm{eff}}\) sensitivity
\cite{devaney2004,harrington2001,hansen1988}. Subsection~\ref{sec:ch394} reports coefficients as constrained observables only when
Eq.~\eqref{eq:ch3.inference-artifact-rejection} does not fire \cite{harrington2001,stratton1941}.

\paragraph{Problem (classify a vendor OAM metric).}
\textbf{Problem.} A vendor metric ``OAM purity'' lacks \(\mathbf{T}\) and \(\mathfrak{A}\). Classify using
Eq.~\eqref{eq:ch3.observable-vs-artifact-screen} \cite{harrington2001,allen1992,belinfante1940,stratton1941}.

\textbf{Step 1.} Request chart and gates \cite{hansen1988}.

\textbf{Step 2.} Compute \(\mathfrak{A}\) if data allow \cite{harrington2001}.

\textbf{Step 3.} Default to \(\mathfrak{S}_{\mathrm{art}}\) \cite{belinfante1940}.

\textbf{Physical meaning.} Proprietary metrics are artifacts until certified \cite{stratton1941,harrington2001}.

\begin{table}[t]
\centering
\caption{Subsection~\ref{sec:ch394} observable versus artifact examples.}
\label{tab:ch3.394-examples}
\small
\begin{tabular}{@{}p{0.30\linewidth}p{0.30\linewidth}p{0.34\linewidth}@{}}
\toprule
Report & Class & Screen \\
\midrule
\(\mathcal{P}_{\mathrm{rad}}\) & observable & Eq.~\eqref{eq:ch3.mode-power-from-coefficient} \\
Certified \(\widehat{c}\) & observable & Eq.~\eqref{eq:ch3.reportable-eigenfield-observable-set} \\
Raw FFT peak & artifact & Eq.~\eqref{eq:ch3.fourier-radiation-filter-chain} \\
``OAM purity'' & artifact & Eq.~\eqref{eq:ch3.inference-artifact-rejection} \\
\bottomrule
\end{tabular}
\end{table}

\paragraph{Symmetry penalty on inferred labels.}
\(\chi_{\mathrm{sym}}\) in Eq.~\eqref{eq:ch1.observable-inference-criterion} rises when inferred labels violate declared symmetry class, as in
broken-axis horn reports using cylindrical \(m\) without mixing matrix \cite{harrington2001,stratton1941,hansen1988}. Subsection~\ref{sec:ch394}
uses \(\chi_{\mathrm{sym}}\) to demote misaligned descriptors to artifacts \cite{harrington2001,devaney2004}.

\paragraph{Worked illustration (phase plate versus observable set).}
A spiral phase plate produces striking vortex images in \(\mathfrak{S}_{\mathrm{art}}\) until far-zone VSH coefficients pass
Eq.~\eqref{eq:ch3.reportable-eigenfield-observable-set} \cite{allen1992,harrington2001,stratton1941,belinfante1940}. Marketing images are not
observables \cite{harrington2001,belinfante1940}.

\paragraph{Problem (promote artifact to constrained surrogate).}
\textbf{Problem.} An inverted \(\widehat{c}\) vector has \(\mathfrak{A}=1.2\,\tau_A\) but \(\mathcal{M}_{\mathrm{mis}}<M_{\max}\). Classify
\cite{harrington2001,stratton1941,devaney2004,hansen1988}.

\textbf{Step 1.} Apply Eq.~\eqref{eq:ch1.observable-decision-rule} \cite{harrington2001}.

\textbf{Step 2.} Report as inference-only descriptor \cite{stratton1941}.

\textbf{Physical meaning.} Partial gates yield constrained surrogates, not full observables \cite{harrington2001,belinfante1940}.

\begin{figure}[t]
\centering
\ChOneFigBlock{%
\begin{tikzpicture}[ch1fig, node distance=7mm and 9mm]
  \node[ch1box] (art) {$\mathfrak{S}_{\mathrm{art}}$};
  \node[ch1box, right=10mm of art] (sur) {constrained surrogate};
  \node[ch1box, right=10mm of sur] (obs) {$\mathcal{O}_{\mathrm{rep}}$};
  \draw[ch1flow] (art) -- (sur) -- (obs);
\end{tikzpicture}%
}
\caption{Subsection~\ref{sec:ch394} promotion ladder: artifacts may become constrained surrogates before full observables.}
\label{fig:ch3.394-promotion-ladder}
\end{figure}

\paragraph{Report contract import.}
Eq.~\eqref{eq:ch1.observable-report-contract} requires tagging each quantity with class, uncertainty, and gate list
\cite{harrington2001,stratton1941}. Subsection~\ref{sec:ch394} extends that contract to eigenfield tables exported to Chapter~\ref{ch:04}
\cite{harrington2001,devaney2004,hansen1988}.

\paragraph{Worked illustration (two labs, same artifact).}
Two labs plot different ``effective \(\ell\)'' from the same data because they use different regularizers; both values are artifacts under
Eq.~\eqref{eq:ch3.beam-order-vs-ell-warning} until full \(\widehat{c}\) vectors are published \cite{harrington2001,devaney2004,stratton1941}.
Subsection~\ref{sec:ch394} resolves disputes by demanding \(\mathcal{O}_{\mathrm{rep}}\) certificates, not scalar slogans \cite{hansen1988,harrington2001}.

\paragraph{Problem (design acceptance test).}
\textbf{Problem.} An acceptance test lists only phase vortex count. Replace with observable checklist
\cite{harrington2001,belinfante1940,allen1992,stratton1941,hansen1988}.

\textbf{Step 1.} Require \(\mathcal{P}_{\mathrm{rad}}\) tolerance \cite{poynting1884}.

\textbf{Step 2.} Require Eq.~\eqref{eq:ch3.reportable-eigenfield-observable-set} on \(\mathcal{L}_{\mathrm{acc}}\) \cite{harrington2001}.

\textbf{Step 3.} Forbid sole reliance on \(\mathfrak{S}_{\mathrm{art}}\) \cite{belinfante1940}.

\textbf{Physical meaning.} Acceptance tests must target observables, not images \cite{stratton1941,harrington2001}.

\paragraph{Chapter closure on interpretation.}
Subsection~\ref{sec:ch394} completes the gauge-and-interpretation arc begun in Chapter~\ref{ch:01}: eigenfields are coordinates,
coefficients are inferred unless gated, paraxial limits are approximate, and popular OAM language is artifact-prone without
Eq.~\eqref{eq:ch3.reportable-eigenfield-observable-set} \cite{belinfante1940,harrington2001,allen1992,stratton1941,hansen1988}.

\begin{figure}[t]
\centering
\ChOneFigBlock{%
\begin{tikzpicture}[ch1fig, node distance=7mm and 9mm]
  \node[ch1box] (ch1) {Ch.~\ref{ch:01} gates};
  \node[ch1box, right=10mm of ch1] (ch3) {eigenfield coords};
  \node[ch1box, right=10mm of ch3] (out) {Eq.~\eqref{eq:ch3.reportable-eigenfield-observable-set}};
  \draw[ch1flow] (ch1) -- (ch3) -- (out);
\end{tikzpicture}%
}
\caption{Subsection~\ref{sec:ch394} exports Chapter~\ref{ch:01} measurement discipline onto Chapter~\ref{ch:03} eigenfield reports.}
\label{fig:ch3.394-chapter-closure}
\end{figure}

\begin{table}[t]
\centering
\caption{Subsection~\ref{sec:ch394} final report tuple for eigenfield campaigns.}
\label{tab:ch3.394-final-tuple}
\small
\begin{tabular}{@{}p{0.32\linewidth}p{0.60\linewidth}@{}}
\toprule
Field & Content \\
\midrule
Observable set & Eq.~\eqref{eq:ch3.reportable-eigenfield-observable-set} \\
Inference indices & Eqs.~\eqref{eq:ch1.observable-inference-criterion}--\eqref{eq:ch1.inference-model-mismatch-index} \\
Artifact screen & Eq.~\eqref{eq:ch3.observable-vs-artifact-screen} \\
Report contract & Eq.~\eqref{eq:ch1.observable-report-contract} \\
\bottomrule
\end{tabular}
\end{table}

\paragraph{Modal identifiability and finite-channel bounds.}
Finite apertures and noise turn eigenfield inversion into a restricted linear map; not every label in a complete basis is identifiable from data
\cite{harrington2001,devaney2004,stratton1941}. Subsection~\ref{sec:ch394} couples the inference screen
Eq.~\eqref{eq:ch1.observable-inference-criterion} to a resolution inequality on certified coefficients:
\begin{equation}
\label{eq:ch3.modal-identifiability-bound}
\operatorname{Var}\!\bigl(\widehat{c}_{\ell m}\bigr)
\ge
\bigl[\mathbf{F}^{-1}\bigr]_{(\ell m),(\ell m)},
\qquad
\mathbf{F}=\mathbf{G}^{\mathsf{H}}\mathbf{C}_{\eta}^{-1}\mathbf{G},
\end{equation}
with \(\mathbf{G}\) the forward channel matrix from Eq.~\eqref{eq:ch1.modal-projection-matrix}, \(\mathbf{C}_{\eta}\) the noise covariance, and
\(\mathbf{F}\) the Fisher information on \(\widehat{c}\) \cite{harrington2001,devaney2004}. Equation~\eqref{eq:ch3.modal-identifiability-bound}
states the physical ceiling: when \(\bigl[\mathbf{F}^{-1}\bigr]_{(\ell m),(\ell m)}\) exceeds the promotion tolerance
\(\tau_{A}\) in Eq.~\eqref{eq:ch1.observable-decision-rule}, the corresponding index belongs to \(\mathfrak{S}_{\mathrm{art}}\) even if a fitting
routine returns a numeric value \cite{harrington2001,stratton1941}. Two laboratories may therefore report different ``effective \(\ell\)'' from identical
raw scans without contradiction until \(\mathbf{G}\), \(\mathbf{C}_{\eta}\), and \(\mathcal{L}_{\mathrm{acc}}\) are matched \cite{devaney2004,hansen1988}.

\paragraph{Promotion ladder with explicit mismatch indices.}
Subsection~\ref{sec:ch394} orders inference outcomes by the mismatch functional
Eq.~\eqref{eq:ch1.inference-model-mismatch-index}:
\begin{equation}
\label{eq:ch3.observable-promotion-ladder}
\mathfrak{S}_{\mathrm{art}}
\ \xrightarrow{\ \mathcal{M}_{\mathrm{mis}}<M_{\max}\ }
\mathcal{S}_{\mathrm{sur}}
\ \xrightarrow{\ \mathcal{G}_{\mathrm{obs}}=1\ }
\mathcal{O}_{\mathrm{rep}},
\end{equation}
with \(\mathcal{S}_{\mathrm{sur}}\) a constrained surrogate and \(\mathcal{G}_{\mathrm{obs}}=1\) denoting satisfaction of Eq.~\eqref{eq:ch1.global-observable-admissibility}
\cite{harrington2001,stratton1941,belinfante1940}. Equation~\eqref{eq:ch3.observable-promotion-ladder}
makes the epistemic ladder executable: phase-vortex images and LG fits remain at the artifact stage until
Eqs.~\eqref{eq:ch3.gauge-qualified-eigenfield-report} and~\eqref{eq:ch3.power-invariance-basis-change} pass, after which
Eq.~\eqref{eq:ch3.reportable-eigenfield-observable-set} lists the admissible observables for Chapter~\ref{ch:04} transport claims
\cite{hansen1988,harrington2001,devaney2004}.

\begin{figure}[t]
\centering
\ChOneFigBlock{%
\begin{tikzpicture}[ch1fig, node distance=7mm and 9mm]
  \node[ch1box] (F) {Eq.~\eqref{eq:ch3.modal-identifiability-bound}};
  \node[ch1box, right=10mm of F] (lad) {Eq.~\eqref{eq:ch3.observable-promotion-ladder}};
  \node[ch1box, right=10mm of lad] (rep) {Eq.~\eqref{eq:ch3.reportable-eigenfield-observable-set}};
  \draw[ch1flow] (F) -- (lad) -- (rep);
\end{tikzpicture}%
}
\caption{Subsection~\ref{sec:ch394} identifiability bound precedes the promotion ladder to reportable observables.}
\label{fig:ch3.394-identifiability-ladder}
\end{figure}

\paragraph{Peer-review rule.}
Manuscripts must tag every reported eigenfield quantity as \(\mathfrak{S}_{\mathrm{obs}}\), constrained surrogate, or
\(\mathfrak{S}_{\mathrm{art}}\) using Table~\ref{tab:ch3.394-examples} \cite{harrington2001,stratton1941,belinfante1940,devaney2004}. Subsection~\ref{sec:ch394}
rejects untagged OAM or \(\ell\) headlines \cite{allen1992,harrington2001,belinfante1940}.

\paragraph{Observables, artifacts, and inference hygiene.}
An \emph{observable} in Subsection~\ref{sec:ch394} is a gauge-invariant functional of \(\mathbf{E},\mathbf{H}\) on
\(\mathcal{R}_\omega\) with a declared normalization and outgoing certificate; a \emph{representation artifact} is any
mode label that changes under admissible basis change or gauge relabeling without a certified bridge
\cite{belinfante1940,harrington2001,devaney2004}. Equation~\eqref{eq:ch3.observable-vs-artifact-screen} embeds the
screening logic imported from Subsection~\ref{sec:ch01.5.4}: inference may promote a quantity to observable status only
after it survives Eqs.~\eqref{eq:ch3.gauge-qualified-eigenfield-report} and \eqref{eq:ch3.power-invariance-basis-change}
\cite{harrington2001,stratton1941,belinfante1940}. Equation~\eqref{eq:ch3.inference-artifact-rejection} is the
operational rejection rule for headlines that mix paraxial OAM, spherical \(\ell\), and helicity sectors without chart
metadata \cite{allen1992,harrington2001,belinfante1940}. Equation~\eqref{eq:ch3.reportable-eigenfield-observable-set}
closes the chapter interpretation layer with an explicit positive list---fluxes, orthogonality inner products,
sector powers, and certified spectrum coefficients---so that embedded equations elsewhere in Chapter~\ref{ch:03} share
one reportable vocabulary \cite{hansen1988,harrington2001,stratton1941,belinfante1940}. The implicit concept is
\emph{epistemic hygiene in harmonic coordinates}: mathematics supplies many complete bases; measurement supplies one
declared protocol \cite{devaney2004,harrington2001}.

\paragraph{Vendor metric audit template.}
Subsection~\ref{sec:ch394} recommends a three-line audit for any vendor OAM or \(\ell\) metric: classify the quantity against Eq.~\eqref{eq:ch3.observable-vs-artifact-screen}; compute \(\mathcal{M}_{\mathrm{mis}}\) from Eq.~\eqref{eq:ch1.inference-model-mismatch-index}; verify promotion against Eq.~\eqref{eq:ch3.observable-promotion-ladder}
\cite{harrington2001,devaney2004,stratton1941}. Metrics that fail the ladder remain \(\mathfrak{S}_{\mathrm{art}}\) even when plots are visually compelling
\cite{allen1992,belinfante1940,hansen1988}. Figure~\ref{fig:ch3.394-identifiability-ladder} orders that audit: Fisher bound first, promotion second, reportable set last
\cite{harrington2001,stratton1941}.

\paragraph{Chapter interpretation closure.}
Eq.~\eqref{eq:ch3.reportable-eigenfield-observable-set} is the positive vocabulary Chapter~\ref{ch:04} must inherit: fluxes, orthogonality inner products, sector powers, and certified spectrum coefficients
\cite{belinfante1940,hansen1988,harrington2001}. Any transport or realizability headline built from labels outside that set must carry an explicit surrogate certificate or be rejected by Eq.~\eqref{eq:ch3.inference-artifact-rejection}
\cite{stratton1941,devaney2004,allen1992}.

\paragraph{Inference acceptance record (template).}
Subsection~\ref{sec:ch394} recommends archiving the tuple \((\mathcal{M}_{\mathrm{mis}},\tau_{A},M_{\max},\chi_{\mathrm{sym}},\mathcal{G}_{\mathrm{obs}})\) imported from Section~\ref{sec:ch01.5.4} beside every inferred \(\widehat{c}_{\ell m}^{(\sigma)}\) table
\cite{harrington2001,stratton1941,devaney2004}. Without that tuple, Eq.~\eqref{eq:ch3.observable-promotion-ladder} cannot be evaluated and the table defaults to \(\mathfrak{S}_{\mathrm{art}}\)
\cite{belinfante1940,hansen1988,allen1992}. Vendor metrics that omit \(\mathbf{G}\) and \(\mathbf{C}_{\eta}\) fail Eq.~\eqref{eq:ch3.modal-identifiability-bound} by construction
\cite{harrington2001,devaney2004}.

\paragraph{Two-laboratory reconciliation protocol.}
When two laboratories report different \(\ell\) or OAM headlines for the same \(\mathcal{P}_{\mathrm{rad}}\), Subsection~\ref{sec:ch394} orders reconciliation: match \(\mathbf{G}\) and \(\mathbf{C}_{\eta}\); match \(\mathcal{N}_{\mathrm{mtr}}\); match \(\mathbf{T}\); re-evaluate Eq.~\eqref{eq:ch3.modal-identifiability-bound}; only then compare sparse labels
\cite{harrington2001,devaney2004,stratton1941,hansen1988}. Disagreement after those steps is a representation artifact, not a physics disagreement
\cite{belinfante1940,allen1992,jackson1999}.

Subsection~\ref{sec:ch394} has contrasted observables with representation artifacts through
Eqs.~\eqref{eq:ch3.observable-vs-artifact-screen}--\eqref{eq:ch3.reportable-eigenfield-observable-set}
\cite{harrington2001,stratton1941,belinfante1940}. Section~\ref{sec:ch39} is complete: gauge-invariant eigenfield meaning, conserved ledger,
paraxial limits, and observable classification close the interpretation thread on \(\mathcal{R}_\omega\) \cite{belinfante1940,hansen1988,harrington2001}.

The four subsections of Section~\ref{sec:ch39} close gauge and approximation discipline:
Eqs.~\eqref{eq:ch3.eigenfield-gauge-invariance-class}--\eqref{eq:ch3.gauge-qualified-eigenfield-report} in Subsection~\ref{sec:ch391};
Eqs.~\eqref{eq:ch3.conserved-vs-representation-split}--\eqref{eq:ch3.representation-dependent-ledger} in Subsection~\ref{sec:ch392};
Eqs.~\eqref{eq:ch3.paraxial-reduction-condition}--\eqref{eq:ch3.paraxial-validity-boundary} in Subsection~\ref{sec:ch393}; and
Eqs.~\eqref{eq:ch3.observable-vs-artifact-screen}--\eqref{eq:ch3.reportable-eigenfield-observable-set} in Subsection~\ref{sec:ch394}
\cite{stratton1941,belinfante1940,harrington2001,allen1992,jackson1999}. Section~\ref{sec:ch310} now supplies worked examples that arithmetic-check
the representation laws \cite{hansen1988,stratton1941}.


\section{Examples}
\label{sec:ch310}

Representation claims earn credibility only when arithmetic on concrete geometries reproduces the
orthogonality, normalization, and coupling laws proved earlier. Section~\ref{sec:ch310} supplies four
worked problems in that diagnostic spirit \cite{stratton1941,hansen1988}. Each subsection is a FIG+PROB instance:
a stated problem, schematic figure, equation steps with physical interpretation, and acceptance thresholds tied to
Chapter~\ref{ch:03} anchors \cite{harrington2001,stratton1941}. Subsection~\ref{sec:ch3101}
constructs free-space angular-momentum eigenfields using Subsection~\ref{sec:ch344}. Subsection~\ref{sec:ch3102}
verifies orthogonality and completeness of vector spherical harmonics via Subsection~\ref{sec:ch343}.
Subsection~\ref{sec:ch3103} exercises continuous-spectrum normalization and mode counting from
Section~\ref{sec:ch35}. Subsection~\ref{sec:ch3104} demonstrates symmetry breaking and mode coupling using
Subsection~\ref{sec:ch363} \cite{harrington2001,jackson1999}. Readers should treat the section as a coordinated audit:
watt partitions in Subsection~\ref{sec:ch3101} presuppose orthogonality proved in Subsection~\ref{sec:ch3102}; discrete
Parseval in Subsection~\ref{sec:ch3103} presupposes normalization meters of Section~\ref{sec:ch35}; boundary coupling in
Subsection~\ref{sec:ch3104} presupposes mixing taxonomy of Subsection~\ref{sec:ch363} \cite{hansen1988,poynting1884,devaney2004}.


\subsection{Free-space angular-momentum eigenfields}
\label{sec:ch3101}
% TOC tag: [FIG+PROB][USE SS3.4.4]

\emph{Free-space angular-momentum eigenfields} are outgoing vector modes \(\mathbf{u}_{\ell m}^{(\sigma)}\) that simultaneously
diagonalize \(\mathbf{J}^{2}\) and \(\mathbf{J}_{z}\) on \(\mathcal{F}_{\mathrm{rad}}\) \cite{jackson1999,hansen1988,stratton1941}.
\emph{Construction} means building those modes from VSH pairs and verifying their quantum numbers against the separation algebra of
Subsection~\ref{sec:ch344}, not merely plotting far-zone patterns \cite{hansen1988,tai1994}. Subsection~\ref{sec:ch3101} is a worked
FIG+PROB instance that arithmetic-checks eigenfield definitions before Chapter~\ref{ch:04} realizability arguments consume them
\cite{harrington2001,stratton1941}.

\paragraph{Standing imports (no re-derivation).}
Eigenfield definition Eq.~\eqref{eq:ch3.eigenfield-VSH-def}, \(\mathbf{M}\)/\(\mathbf{N}\) constructions
Eqs.~\eqref{eq:ch3.Mlm-def}--\eqref{eq:ch3.Nlm-def}, simultaneous system Eq.~\eqref{eq:ch3.J-simultaneous-eigenfield-system},
eigenvalue Eq.~\eqref{eq:ch3.eigenfield-Jsq-eigenvalue}, normalized resolution Eq.~\eqref{eq:ch3.normalized-coefficient-resolution},
mode power law Eq.~\eqref{eq:ch3.mode-power-from-coefficient}, and Theorem~\ref{thm:ch3.vsh-J-eigenfields} enter by reference
\cite{hansen1988,stratton1941,jackson1999}. Subsection~\ref{sec:ch3101} does not re-derive Debye potentials
\cite{tai1994,harrington2001}.

\paragraph{Problem.}
A lossless reciprocal antenna test at \(\omega\) exports \(\mathcal{P}_{\mathrm{rad}}=1\,\mathrm{W}\) on \(\Gamma_R\) with normalized
superposition
\begin{equation}
\label{eq:ch3.wp3101-eigenfield-superposition}
\mathbf{v}_\omega
=
\widehat{c}_{10}^{(\mathbf{N})}\,\widehat{\mathbf{u}}_{10}^{(\mathbf{N})}
+
\widehat{c}_{20}^{(\mathbf{N})}\,\widehat{\mathbf{u}}_{20}^{(\mathbf{N})},
\qquad
\bigl|\widehat{c}_{10}^{(\mathbf{N})}\bigr|^{2}=0.64\,\mathrm{W},
\quad
\bigl|\widehat{c}_{20}^{(\mathbf{N})}\bigr|^{2}=0.36\,\mathrm{W},
\end{equation}
using flux-normalized eigenfields from Eqs.~\eqref{eq:ch3.VSH-unit-flux-normalization}--\eqref{eq:ch3.VSH-unit-energy-norm}
\cite{poynting1884,hansen1988,harrington2001}. Tasks:
\begin{enumerate}
\item[(a)] verify \((\ell,m)=(1,0)\) and \((2,0)\) labels satisfy Eq.~\eqref{eq:ch3.eigenfield-Jsq-eigenvalue};
\item[(b)] confirm \(\mathbf{J}_{z}\) eigenvalues on each term;
\item[(c)] compute far-zone angular anisotropy qualitatively from sector weights;
\item[(d)] state why a single ``\(\ell=1.36\)'' effective index is forbidden under Eq.~\eqref{eq:ch3.integer-ell-def}.
\end{enumerate}

\begin{figure}[t]
\centering
\ChOneFigBlock{%
\begin{tikzpicture}[ch1fig, node distance=7mm and 10mm]
  \node[ch1box] (u10) {$\widehat{\mathbf{u}}_{10}^{(\mathbf{N})}$\\$0.64\,\mathrm{W}$};
  \node[ch1box, right=12mm of u10] (u20) {$\widehat{\mathbf{u}}_{20}^{(\mathbf{N})}$\\$0.36\,\mathrm{W}$};
  \node[ch1box, right=12mm of u20] (sum) {Eq.~\eqref{eq:ch3.wp3101-eigenfield-superposition}};
  \node[ch1box, below=11mm of u10] (j1) {$\ell=1$, $m=0$};
  \node[ch1box, below=11mm of u20] (j2) {$\ell=2$, $m=0$};
  \draw[ch1flow] (u10) -- (sum);
  \draw[ch1flow] (u20) -- (sum);
  \draw[ch1flow] (j1) -- (u10);
  \draw[ch1flow] (j2) -- (u20);
\end{tikzpicture}%
}
\caption{Worked problem~\ref{sec:ch3101}: two integer eigenfield sectors superpose on \(\mathcal{F}_{\mathrm{rad}}\) with watt weights
\(0.64\) and \(0.36\). Each block carries its own \((\ell,m,\sigma)\) label from Eq.~\eqref{eq:ch3.eigenfield-VSH-def}.}
\label{fig:ch3.3101-dipole-quadrupole}
\end{figure}

Figure~\ref{fig:ch3.3101-dipole-quadrupole} fixes the geometry: two integer ladders interfere without merging labels
\cite{hansen1988,harrington2001,stratton1941}.

\paragraph{Step 1 ---\(\mathbf{J}^{2}\) eigenvalues (USE §3.4.4).}
Theorem~\ref{thm:ch3.vsh-J-eigenfields} certifies
\(\mathbf{J}^{2}\mathbf{u}_{\ell m}^{(\sigma)}=\ell(\ell+1)\mathbf{u}_{\ell m}^{(\sigma)}\) through
Eq.~\eqref{eq:ch3.eigenfield-Jsq-eigenvalue} \cite{jackson1999,hansen1988}. For \((1,0)\),
\(\ell(\ell+1)=2\); for \((2,0)\), \(\ell(\ell+1)=6\) \cite{jackson1999,stratton1941}. Mathematically, each superposition term
lives in a distinct eigenspace of \(\mathbf{J}^{2}\); physically, dipole and quadrupole angular momentum content coexist without
assigning a single fractional \(\ell\) \cite{harrington2001,hansen1988}. Equation~\eqref{eq:ch3.wp3101-eigenfield-superposition} is
therefore a direct sum of eigenspaces, not a single eigenmode \cite{stratton1941,jackson1999}.

\paragraph{Step 2 ---\(\mathbf{J}_{z}\) labels.}
Azimuthal index \(m=0\) on both terms implies \(\mathbf{J}_{z}\mathbf{u}_{\ell 0}^{(\mathbf{N})}=0\) \cite{jackson1999,hansen1988}.
The pattern is axisymmetric about the laboratory \(z\)-axis; no azimuthal phase wind appears despite quadrupole admixture
\cite{stratton1941,harrington2001}. One must not confuse \(m=0\) with ``no angular structure'': \(\ell=2\) still produces
four-lobe far-zone anisotropy \cite{jackson1999,hansen1988}.

\paragraph{Step 3 ---Watt partition (HOME check).}
Normalized orthogonality gives
\begin{equation}
\label{eq:ch3.wp3101-power-partition}
\mathcal{P}_{\mathrm{rad}}
=
\bigl|\widehat{c}_{10}^{(\mathbf{N})}\bigr|^{2}
+
\bigl|\widehat{c}_{20}^{(\mathbf{N})}\bigr|^{2}
=
0.64+0.36
=
1\,\mathrm{W},
\end{equation}
by Eq.~\eqref{eq:ch3.mode-power-from-coefficient} \cite{poynting1884,harrington2001,stratton1941}. Mathematically, cross terms vanish
because distinct \((\ell,m,\sigma)\) labels are orthogonal on \(\mathcal{F}_{\mathrm{rad}}\); physically, watt meters integrate
incoherent sector powers when eigenfield projectors are resolved \cite{hansen1988,harrington2001}.

\paragraph{Step 4 ---Pattern anisotropy (qualitative).}
\(\widehat{\mathbf{u}}_{10}^{(\mathbf{N})}\) peaks at \(\theta=\pi/2\); \(\widehat{\mathbf{u}}_{20}^{(\mathbf{N})}\) introduces
deeper nulls along \(\theta=0,\pi\) \cite{jackson1999,hansen1988}. With \(64\%\) dipole and \(36\%\) quadrupole watts, the composite
pattern is dipole-dominated but not pure \(\ell=1\) \cite{harrington2001,stratton1941}. Subsection~\ref{sec:ch3101} uses that split to
reject single-number beam summaries \cite{harrington2001,devaney2004}.

\paragraph{Step 5 ---Integer index audit.}
Eq.~\eqref{eq:ch3.integer-ell-def} requires integer \(\ell\) labels; \(0.64\) and \(0.36\) are \emph{power weights}, not \(\ell\) values
\cite{hansen1988,harrington2001}. An ``effective \(\ell=1.36\)'' collapses orthogonal sectors and violates
Eq.~\eqref{eq:ch3.beam-order-vs-ell-warning} \cite{stratton1941,harrington2001}.

\paragraph{Physical meaning (closure).}
Free-space eigenfields are integer-labeled building blocks; superpositions carry multiple labels simultaneously
\cite{jackson1999,hansen1988,stratton1941}. Subsection~\ref{sec:ch3101} certifies that laboratory watt splits map to projector
coefficients, not to informal fit orders \cite{harrington2001,poynting1884}.

\begin{table}[t]
\centering
\caption{Subsection~\ref{sec:ch3101} eigenfield quantum-number ledger.}
\label{tab:ch3.3101-quantum-ledger}
\small
\begin{tabular}{@{}p{0.20\linewidth}p{0.22\linewidth}p{0.22\linewidth}p{0.26\linewidth}@{}}
\toprule
Sector & \(\mathbf{J}^{2}\) & \(\mathbf{J}_{z}\) & power \\
\midrule
$(1,0,\mathbf{N})$ & $2$ & $0$ & $0.64\,\mathrm{W}$ \\
$(2,0,\mathbf{N})$ & $6$ & $0$ & $0.36\,\mathrm{W}$ \\
\bottomrule
\end{tabular}
\end{table}

\paragraph{Step 6 --- \(\sigma\) track and dual types.}
Eq.~\eqref{eq:ch3.eigenfield-VSH-def} uses \(\sigma=\mathbf{N}\) on both sectors; an \(\mathbf{M}\)-type admixture
\(\widehat{c}_{11}^{(\mathbf{M})}\neq0\) would add \(\ell=1\), \(m=1\) content without changing \(\ell=2\) block
\cite{belinfante1940,hansen1988,stratton1941}. Subsection~\ref{sec:ch3101} keeps \(\sigma\) explicit because polarization/helicity
tracks are not determined by \(m=0\) axisymmetry alone \cite{harrington2001,belinfante1940}.

\paragraph{Step 7 --- Source projection preview (USE §3.4.4).}
Eq.~\eqref{eq:ch3.eigenfield-source-projection} implies short electric dipoles excite primarily \((1,0,\mathbf{N})\) with tail in
\(\ell\ge2\) from finite source support \cite{stratton1941,jackson1999,harrington2001}. The worked split \(0.64/0.36\) can be read as a
\emph{measurement-derived} projection, not as a universal source law \cite{harrington2001,devaney2004}.

\paragraph{Step 8 --- Far-zone phase and eigenfield meaning.}
Eigenfield phases are defined by outgoing Hankel factors and Condon--Shortley conventions
Eq.~\eqref{eq:ch3.Ylm-Condon-Shortley-def} \cite{hansen1988,jackson1999}. Relative phase between \((1,0)\) and \((2,0)\) sectors affects
interference nulls but not Eq.~\eqref{eq:ch3.wp3101-power-partition} when orthogonality holds \cite{stratton1941,harrington2001}.

\begin{figure}[t]
\centering
\ChOneFigBlock{%
\begin{tikzpicture}[ch1fig, node distance=7mm and 10mm]
  \node[ch1box] (src) {finite source};
  \node[ch1box, right=12mm of src] (proj) {Eq.~\eqref{eq:ch3.eigenfield-source-projection}};
  \node[ch1box, right=12mm of proj] (vec) {$\widehat{c}_{\ell m}^{(\sigma)}$};
  \node[ch1box, below=11mm of vec] (pow) {Eq.~\eqref{eq:ch3.wp3101-power-partition}};
  \draw[ch1flow] (src) -- (proj) -- (vec) -- (pow);
\end{tikzpicture}%
}
\caption{Subsection~\ref{sec:ch3101} construction chain: sources project onto eigenfields; watt sums follow only after normalization and outgoing selection.}
\label{fig:ch3.3101-source-projection}
\end{figure}

\paragraph{Step 9 --- Accessibility screen (pointer).}
If \(\mathcal{L}_{\mathrm{acc}}\) forbids \(\ell=2\) sectors for a given geometry, only \(0.64\,\mathrm{W}\) is reportable even when
\(0.36\,\mathrm{W}\) is mathematically present in Eq.~\eqref{eq:ch3.wp3101-eigenfield-superposition} \cite{harrington2001,stratton1941}.
Subsection~\ref{sec:ch3101} separates \emph{fitted} from \emph{reportable} coefficients \cite{hansen1988,harrington2001}.

\paragraph{Step 10 --- Numerical acceptance band.}
Require \(|\mathcal{P}_{\mathrm{rad}}-1\,\mathrm{W}|<10^{-3}\,\mathrm{W}\) after projection and \(|\widehat{c}_{10}|^{2}=0.64\pm10^{-3}\,\mathrm{W}\)
in reciprocal normalization \cite{harrington2001,hansen1988,stratton1941}. Failure indicates flux meter mismatch or reactive admixture before
\(\Pi_{\mathrm{rad}}\) \cite{harrington2001,stratton1941}.

\paragraph{Eigenvalue ledger in operator form.}
The simultaneous system Eq.~\eqref{eq:ch3.J-simultaneous-eigenfield-system} can be read as a commuting-family diagonalization on
\(\mathcal{F}_{\mathrm{rad}}\): each normalized mode \(\widehat{\mathbf{u}}_{\ell m}^{(\sigma)}\) carries eigenvalues
\(\hbar\ell(\ell+1)\) for \(\mathbf{J}^{2}\) and \(\hbar m\) for \(\mathbf{J}_{z}\) in the angular-momentum convention of
Section~\ref{sec:ch32} \cite{jackson1999,hansen1988,stratton1941}. For the worked pair \((1,0)\) and \((2,0)\), the operator labels
are orthogonal eigenspaces, so the superposition Eq.~\eqref{eq:ch3.wp3101-eigenfield-superposition} is not an eigenvector of
\(\mathbf{J}^{2}\) unless one coefficient vanishes \cite{harrington2001,stratton1941}. Mathematically, that non-eigenvector status is
the reason a single effective \(\ell\) is undefined; physically, it is the difference between a dipole--quadrupole blend and a pure
multipole source \cite{jackson1999,hansen1988}. It is essential to report the full coefficient vector, not a fitted scalar order.

\paragraph{Far-zone template shapes and angular momentum.}
The \(\ell=1\), \(m=0\), \(\mathbf{N}\)-type template peaks in the equatorial plane because \(\mathbf{J}_{z}\) eigenvalue zero suppresses
azimuthal phase variation while \(\ell=1\) fixes the lowest nontrivial angular anisotropy \cite{jackson1999,hansen1988}. The \(\ell=2\),
\(m=0\) template deepens polar nulls: quadrupole angular momentum couples opposite lobes along \(\theta=0,\pi\) without introducing
\(m=\pm1\) winds \cite{stratton1941,harrington2001}. Equation~\eqref{eq:ch3.wp3101-power-partition} weights those templates by watts, so a
\(36\%\) quadrupole admixture can sharpen nulls even when the pattern remains dipole-dominated in peak direction \cite{poynting1884,hansen1988}.
Figure~\ref{fig:ch3.3101-dipole-quadrupole} is the schematic reminder that quantum numbers attach to \emph{sectors}, not to composite peaks.

\paragraph{Outgoing Hankel phase and relative interference.}
On \(\Gamma_R\), each sector carries outgoing radial dependence through Hankel functions selected in Subsection~\ref{sec:ch343}
\cite{hansen1988,stratton1941}. Relative phase between \(\widehat{c}_{10}^{(\mathbf{N})}\) and \(\widehat{c}_{20}^{(\mathbf{N})}\) enters
far-zone \(\E\) linearly and can move null locations while leaving Eq.~\eqref{eq:ch3.wp3101-power-partition} unchanged when orthogonality holds
\cite{jackson1999,harrington2001}. Subsection~\ref{sec:ch3101} therefore distinguishes \emph{power observables} from \emph{field-phase
observables}: watt meters integrate quadratic forms; interferometric channels require explicit relative phase metadata
\cite{stratton1941,hansen1988}.

\paragraph{Source projection and near-field admixture.}
Eq.~\eqref{eq:ch3.eigenfield-source-projection} states that impressed currents project onto eigenfields through the radiative inner product
\cite{stratton1941,harrington2001}. A short dipole excites primarily \((1,0,\mathbf{N})\) but finite support and higher-order current moments
feed \(\ell\ge2\) tails \cite{jackson1999,hansen1988}. The split \(0.64/0.36\) in Eq.~\eqref{eq:ch3.wp3101-eigenfield-superposition} is
interpreted here as a \emph{measurement-derived} projection on \(\Gamma_R\), not as a universal source theorem \cite{harrington2001,devaney2004}.
If reactive near-field content is included before \(\Pi_{\mathrm{rad}}\), coefficients absorb non-export energy and violate
Eq.~\eqref{eq:ch3.wp3101-power-partition}; Figure~\ref{fig:ch3.3101-source-projection} orders the correct pipeline as source \(\to\) outgoing
filter \(\to\) projection \(\to\) watt sum \cite{harrington2001,stratton1941}.

\paragraph{Integer-index audit and beam-order confusion.}
Eq.~\eqref{eq:ch3.integer-ell-def} requires integer \(\ell\) for pure VSH sectors \cite{hansen1988,harrington2001}. Popular ``beam order''
fits that return non-integer \(\ell_{\mathrm{eff}}\) collapse distinct eigenspaces and violate Eq.~\eqref{eq:ch3.beam-order-vs-ell-warning}
\cite{stratton1941,harrington2001}. The worked weights \(0.64\) and \(0.36\) are Parseval coefficients, not angular indices; conflating them
is the most common post-processing error in OAM-themed antenna papers \cite{allen1992,harrington2001,belinfante1940}. Subsection~\ref{sec:ch3101}
uses the arithmetic instance to train readers to publish full \((\ell,m,\sigma)\) tables with Eq.~\eqref{eq:ch3.wp3101-power-partition} closure.

\paragraph{Dual-sector decomposition at fixed $(\ell,m)$.}
Even with \(m=0\) axisymmetry, \(\sigma\in\{\mathbf{M},\mathbf{N}\}\) tracks distinct dual templates from Eq.~\eqref{eq:ch3.eigenfield-VSH-def}
\cite{belinfante1940,hansen1988,stratton1941}. Subsection~\ref{sec:ch3101} records the dual-sector watt law
\begin{equation}
\label{eq:ch3.wp3101-dual-sector-split}
\mathcal{P}_{\mathrm{rad}}^{(\mathbf{N})}
=
\bigl|\widehat{c}_{10}^{(\mathbf{N})}\bigr|^{2}
+
\bigl|\widehat{c}_{20}^{(\mathbf{N})}\bigr|^{2},
\qquad
\mathcal{P}_{\mathrm{rad}}^{(\mathbf{M})}
=
\sum_{\ell m}\bigl|\widehat{c}_{\ell m}^{(\mathbf{M})}\bigr|^{2},
\end{equation}
for the worked axisymmetric case with \(\mathcal{P}_{\mathrm{rad}}^{(\mathbf{M})}=0\) \cite{poynting1884,harrington2001}. Mathematically, Eq.~\eqref{eq:ch3.wp3101-dual-sector-split} is a block Parseval sum; physically, it is the reminder that axisymmetric patterns can still carry \(\mathbf{M}\)-type admixture when feeds break dual symmetry
\cite{belinfante1940,jackson1999,harrington2001}.

\paragraph{Riesz projection interpretation of $\widehat{c}_{\ell m}^{(\sigma)}$.}
Coefficients in Eq.~\eqref{eq:ch3.normalized-coefficient-resolution} are Riesz projections of \(\mathbf{v}_\omega\) onto normalized eigenmodes
\cite{kato1995,hansen1988,stratton1941}. Subsection~\ref{sec:ch3101} makes that interpretation explicit for the \(0.64/0.36\) split: each weight is a projector norm squared, not a nonlinear fit parameter
\cite{harrington2001,devaney2004}. Changing \(\Gamma_R\) radius without updating flux normalization shifts \(\widehat{c}\) while leaving \(\mathbf{v}_\omega\) fixed---a meter effect, not a source change
\cite{stratton1941,poynting1884,hansen1988}.

\paragraph{Worked sensitivity (phase rotation of quadrupole sector).}
Rotating the relative phase \(\delta\) between \(\widehat{c}_{10}^{(\mathbf{N})}\) and \(\widehat{c}_{20}^{(\mathbf{N})}\) moves far-zone nulls on \(\Gamma_R\) while preserving Eq.~\eqref{eq:ch3.wp3101-power-partition}
\cite{jackson1999,harrington2001,stratton1941}. Subsection~\ref{sec:ch3101} uses that sensitivity to separate interferometric observables from watt observables: publish \(\delta\) when null locations matter; publish only magnitude squares when watt meters alone are authoritative
\cite{hansen1988,harrington2001}.

\begin{table}[t]
\centering
\caption{Subsection~\ref{sec:ch3101} observable class for dipole--quadrupole blend.}
\label{tab:ch3.3101-observable-class}
\small
\begin{tabular}{@{}p{0.32\linewidth}p{0.30\linewidth}p{0.30\linewidth}@{}}
\toprule
Quantity & Invariant under & Requires metadata \\
\midrule
$|\widehat{c}_{\ell m}|^{2}$ & gauge (field-first) & flux norm \\
Null locations & -- & relative phase $\delta$ \\
$\mathcal{P}_{\mathrm{rad}}$ & basis change & Eq.~\eqref{eq:ch3.wp3101-power-partition} \\
\bottomrule
\end{tabular}
\end{table}

\paragraph{Campaign note.}
Archive \(\ell,m,\sigma\), normalization meters, and Eq.~\eqref{eq:ch3.wp3101-power-partition} residuals with every eigenfield table
\cite{harrington2001,hansen1988,stratton1941}. Subsection~\ref{sec:ch3102} verifies orthogonality that makes
Eq.~\eqref{eq:ch3.wp3101-power-partition} possible \cite{stratton1941,hansen1988}.


\subsection{Orthogonality and completeness of VSH}
\label{sec:ch3102}
% TOC tag: [FIG+PROB][USE SS3.4.3]

\emph{Orthogonality} means distinct normalized VSH eigenmodes have zero radiative inner product on \(\mathcal{F}_{\mathrm{rad}}\)
\cite{hansen1988,stratton1941,harrington2001}. \emph{Completeness} means partial sums converge to arbitrary outgoing states in the
\(\|\cdot\|_{\mathrm{rad}}\) norm \cite{kato1995,hansen1988}. Subsection~\ref{sec:ch3102} is a FIG+PROB arithmetic instance importing
Theorem~\ref{thm:ch3.vsh-orthogonality} and Theorem~\ref{thm:ch3.vsh-completeness} from Subsection~\ref{sec:ch343} without re-proof
\cite{stratton1941,harrington2001}.

\paragraph{Standing imports.}
Radiative inner product Eq.~\eqref{eq:ch3.radiative-inner-product-def}, orthogonality Theorem~\ref{thm:ch3.vsh-orthogonality},
completeness Theorem~\ref{thm:ch3.vsh-completeness}, resolution identity Eq.~\eqref{eq:ch3.VSH-resolution-identity}, and
\(L^{2}\) convergence Eq.~\eqref{eq:ch3.L2-convergence-VSH} enter by reference \cite{hansen1988,stratton1941,kato1995}. Subsection~\ref{sec:ch3102}
does not re-derive spherical harmonic addition theorems \cite{jackson1999,hansen1988}.

\paragraph{Problem.}
Consider normalized pair \(\widehat{\mathbf{u}}_{10}^{(\mathbf{N})}\) and \(\widehat{\mathbf{u}}_{20}^{(\mathbf{N})}\) on \(\Gamma_R\).
\begin{enumerate}
\item[(a)] evaluate orthogonality integral Eq.~\eqref{eq:ch3.wp3102-orthogonality-integral};
\item[(b)] truncate an arbitrary \(\mathbf{v}_\omega\) at \(\ell_{\max}=2\) and bound remainder energy;
\item[(c)] explain why orthogonality failure in a numerical code usually signals normalization or outgoing-gauge mismatch.
\end{enumerate}

\begin{figure}[t]
\centering
\ChOneFigBlock{%
\begin{tikzpicture}[ch1fig, node distance=7mm and 10mm]
  \node[ch1box] (u) {$\widehat{\mathbf{u}}_{10}^{(\mathbf{N})}$};
  \node[ch1box, right=12mm of u] (v) {$\widehat{\mathbf{u}}_{20}^{(\mathbf{N})}$};
  \node[ch1box, right=12mm of v] (ip) {$\langle\cdot,\cdot\rangle_{\mathrm{rad}}$};
  \node[ch1box, below=11mm of ip] (zero) {$=0$};
  \node[ch1box, below=11mm of u] (shell) {$\Gamma_R$};
  \draw[ch1flow] (u) -- (ip);
  \draw[ch1flow] (v) -- (ip);
  \draw[ch1flow] (zero) -- (ip);
  \draw[ch1flow] (shell) -- (ip);
\end{tikzpicture}%
}
\caption{Worked problem~\ref{sec:ch3102}: orthogonality is tested on the far observation shell \(\Gamma_R\) with flux-normalized VSH pairs.}
\label{fig:ch3.3102-orthogonality-geometry}
\end{figure}

\paragraph{Step 1 ---Orthogonality integral (HOME).}
By Theorem~\ref{thm:ch3.vsh-orthogonality},
\begin{equation}
\label{eq:ch3.wp3102-orthogonality-integral}
\bigl\langle \widehat{\mathbf{u}}_{10}^{(\mathbf{N})},\,\widehat{\mathbf{u}}_{20}^{(\mathbf{N})}\bigr\rangle_{\mathrm{rad}}
=
0,
\end{equation}
because \(\ell=1\neq2\) sectors are orthogonal under Eq.~\eqref{eq:ch3.radiative-inner-product-def} \cite{hansen1988,stratton1941,harrington2001}.
Mathematically, distinct \(\mathbf{J}^{2}\) eigenspaces decouple the inner product; physically, watt coupling between dipole and quadrupole
templates vanishes when both are properly normalized \cite{poynting1884,hansen1988}. Figure~\ref{fig:ch3.3102-orthogonality-geometry} is the
geometry behind Eq.~\eqref{eq:ch3.wp3101-power-partition} \cite{stratton1941,harrington2001}.

\paragraph{Step 2 ---Same-\(\ell\) non-orthogonality caveat.}
Modes sharing \(\ell\) but differing in \(m\) or \(\sigma\) need not be orthogonal in degenerate subspaces until \(m\) is fixed
\cite{jackson1999,hansen1988}. Subsection~\ref{sec:ch3102} tests \emph{distinct} \(\ell\) first; Wigner mixing within fixed \(\ell\) is
imported from Eq.~\eqref{eq:ch3.eigenfield-Wigner-mixing} without re-derivation \cite{harrington2001,stratton1941}.

\paragraph{Step 3 ---Truncation remainder (HOME).}
Let \(\mathbf{v}_\omega\) have coefficients \(\widehat{c}_{\ell m}^{(\sigma)}\) with \(\sum_{\ell>2}|\widehat{c}|^{2}=0.08\,\mathrm{W}\).
Truncation projector Eq.~\eqref{eq:ch3.ellmax-truncation-projector} at \(\ell_{\max}=2\) leaves
\begin{equation}
\label{eq:ch3.wp3102-partial-sum-error}
\|\mathbf{R}_{\ell_{\max}}\mathbf{v}_\omega\|_{\mathrm{rad}}^{2}
=
0.08\,\mathrm{W},
\qquad
\|\mathbf{v}_\omega-\mathbf{R}_{\ell_{\max}}\mathbf{v}_\omega\|_{\mathrm{rad}}^{2}
=
0.92\,\mathrm{W},
\end{equation}
importing Eq.~\eqref{eq:ch3.truncation-remainder-energy} \cite{harrington2001,hansen1988,stratton1941}. Mathematically, the remainder is
the tail beyond \(\ell_{\max}\); physically, an \(\ell_{\max}=2\) instrument misses \(8\%\) of export watts \cite{harrington2001,poynting1884}.

\paragraph{Step 4 ---Completeness interpretation (USE §3.4.3).}
Theorem~\ref{thm:ch3.vsh-completeness} promises that increasing \(\ell_{\max}\) drives
\(\|\mathbf{v}_\omega-\mathbf{R}_{\ell_{\max}}\mathbf{v}_\omega\|_{\mathrm{rad}}\to0\) \cite{kato1995,hansen1988}. Subsection~\ref{sec:ch3102}
interprets that as a \emph{convergence rate} problem: smooth patterns decay quickly in \(\ell\); sharp features require large \(\ell_{\max}\)
\cite{harrington2001,devaney2004,stratton1941}.

\paragraph{Step 5 ---Numerical failure modes.}
Orthogonality residuals at \(10^{-3}\) usually indicate mismatched flux normalization
Eqs.~\eqref{eq:ch3.VSH-unit-flux-normalization}--\eqref{eq:ch3.VSH-unit-energy-norm} or reactive near-field admixture before \(\Pi_{\mathrm{rad}}\)
\cite{harrington2001,stratton1941,hansen1988}. Subsection~\ref{sec:ch3102} treats such residuals as implementation bugs, not as physics coupling
\cite{harrington2001,kato1995}.

\paragraph{Physical meaning (closure).}
Orthogonality makes watt sums diagonal; completeness makes truncation auditable \cite{hansen1988,stratton1941,harrington2001}. Together they
justify Eq.~\eqref{eq:ch3.normalized-coefficient-resolution} as a stable coordinate system on \(\mathcal{F}_{\mathrm{rad}}\)
\cite{kato1995,harrington2001}.

\begin{table}[t]
\centering
\caption{Subsection~\ref{sec:ch3102} orthogonality and truncation audit.}
\label{tab:ch3.3102-audit}
\small
\begin{tabular}{@{}p{0.30\linewidth}p{0.34\linewidth}p{0.30\linewidth}@{}}
\toprule
Test & Equation & Pass \\
\midrule
Cross sector & Eq.~\eqref{eq:ch3.wp3102-orthogonality-integral} & $0$ \\
Tail energy & Eq.~\eqref{eq:ch3.wp3102-partial-sum-error} & declared \\
Completeness & Theorem~\ref{thm:ch3.vsh-completeness} & $\ell_{\max}\uparrow$ \\
\bottomrule
\end{tabular}
\end{table}

\paragraph{Worked follow-up (rotation within \(\ell\)).}
Rotating the laboratory mixes \(m\) within \(\ell=2\) via Eq.~\eqref{eq:ch3.eigenfield-Wigner-mixing} without exciting \(\ell=1\)
\cite{jackson1999,hansen1988}. Subsection~\ref{sec:ch3102} records that orthogonality across \(\ell\) is preserved under rotation even when
\(m\) blocks mix \cite{harrington2001,stratton1941}.

\paragraph{Step 6 --- Inner product definition audit.}
Eq.~\eqref{eq:ch3.radiative-inner-product-def} integrates over \(\Gamma_R\) with flux weighting \cite{stratton1941,harrington2001}.
Subsection~\ref{sec:ch3102} treats missing \(\sin\theta\) weights as the dominant cause of false non-orthogonality in pattern codes
\cite{hansen1988,jackson1999,harrington2001}.

\paragraph{Step 7 --- Resolution identity spot check.}
Eq.~\eqref{eq:ch3.VSH-resolution-identity} implies \(\sum_{\ell m\sigma}\mathbf{P}_{\ell m\sigma}=\mathbf{I}\) on \(\mathcal{F}_{\mathrm{rad}}\)
\cite{hansen1988,kato1995}. Truncating at \(\ell_{\max}=2\) violates identity by the \(0.08\,\mathrm{W}\) tail in
Eq.~\eqref{eq:ch3.wp3102-partial-sum-error} \cite{harrington2001,stratton1941}.

\begin{figure}[t]
\centering
\ChOneFigBlock{%
\begin{tikzpicture}[ch1fig, node distance=7mm and 9mm]
  \node[ch1box] (R2) {$\mathbf{R}_{2}$};
  \node[ch1box, right=10mm of R2] (tail) {tail $0.08\,\mathrm{W}$};
  \node[ch1box, right=10mm of tail] (conv) {Theorem~\ref{thm:ch3.vsh-completeness}};
  \draw[ch1flow] (R2) -- (tail) -- (conv);
\end{tikzpicture}%
}
\caption{Subsection~\ref{sec:ch3102} truncation tail: completeness raises $\ell_{\max}$ until tail falls below measurement noise.}
\label{fig:ch3.3102-truncation-tail}
\end{figure}

\paragraph{Step 8 --- Cross-power diagnostic.}
Define cross-power \(P_{12}=\Re\{\widehat{c}_{10}^{*}c_{20}\}\); orthogonality forces \(P_{12}=0\) in the eigenbasis
\cite{hansen1988,harrington2001,stratton1941}. Nonzero \(P_{12}\) in a vendor report signals basis mismatch or normalization error, not
physical coupling between \(\ell\) sectors \cite{harrington2001,devaney2004}.

\paragraph{Step 9 --- Quadrature on $\Gamma_R$.}
Discrete quadrature on \(\Gamma_R\) must reproduce Eq.~\eqref{eq:ch3.wp3102-orthogonality-integral} to numerical tolerance
\cite{hansen1988,stratton1941}. Subsection~\ref{sec:ch3102} recommends doubling angular samples until cross-integral falls below
\(10^{-4}\,\mathrm{W}\) \cite{harrington2001,kato1995}.

\begin{table}[t]
\centering
\caption{Subsection~\ref{sec:ch3102} numerical orthogonality protocol.}
\label{tab:ch3.3102-numerical-protocol}
\small
\begin{tabular}{@{}p{0.10\linewidth}p{0.52\linewidth}p{0.30\linewidth}@{}}
\toprule
Iter & Action & Target \\
\midrule
1 & flux weights & Eq.~\eqref{eq:ch3.radiative-inner-product-def} \\
2 & cross integral & Eq.~\eqref{eq:ch3.wp3102-orthogonality-integral} \\
3 & refine grid & $|P_{12}|<10^{-4}\,\mathrm{W}$ \\
\bottomrule
\end{tabular}
\end{table}

\paragraph{Orthogonality as a watt-diagonalization principle.}
Theorem~\ref{thm:ch3.vsh-orthogonality} is the mathematical license for Eq.~\eqref{eq:ch3.wp3101-power-partition}: distinct
\((\ell,m,\sigma)\) labels decouple the radiative inner product so cross terms vanish in quadratic watt sums \cite{hansen1988,kato1995}.
Subsection~\ref{sec:ch3102} makes that abstraction executable on \(\Gamma_R\) through Eq.~\eqref{eq:ch3.wp3102-orthogonality-integral}.
Physically, orthogonality means a dipole template and a quadrupole template do not exchange power under proper flux weighting
\cite{poynting1884,stratton1941}; numerically, it means cross-integrals should fall below \(10^{-4}\,\mathrm{W}\) once
Eq.~\eqref{eq:ch3.radiative-inner-product-def} quadrature is converged \cite{harrington2001,hansen1988}.

\paragraph{Truncation as controlled incompleteness.}
Theorem~\ref{thm:ch3.vsh-completeness} promises exact representation only as \(\ell_{\max}\to\infty\) \cite{kato1995,hansen1988}.
The \(0.08\,\mathrm{W}\) tail in Eq.~\eqref{eq:ch3.wp3102-partial-sum-error} is therefore not a fitting error but a \emph{band-limit}
declaration: an instrument with \(\ell_{\max}=2\) cannot resolve finer angular structure \cite{harrington2001,devaney2004}. Figure~\ref{fig:ch3.3102-truncation-tail}
orders the remediation---raise \(\ell_{\max}\) until the tail falls below noise---mirroring the angular grid refinement protocol of
Table~\ref{tab:ch3.3102-numerical-protocol} \cite{stratton1941,hansen1988}.

\paragraph{Rotation covariance within fixed \(\ell\).}
Eq.~\eqref{eq:ch3.eigenfield-Wigner-mixing} mixes \(m\) within a fixed \(\ell\) block under laboratory rotation while preserving
orthogonality across distinct \(\ell\) \cite{jackson1999,hansen1988}. Subsection~\ref{sec:ch3102} records that split because rotation
residuals are often misread as coupling between \(\ell\) sectors when the code failed to rotate the inner-product weights
\cite{harrington2001,stratton1941}. Wigner matrices are symmetry certificates; they are not substitutes for boundary mixing matrices
Eq.~\eqref{eq:ch3.boundary-scattering-mixing-matrix} treated in Subsection~\ref{sec:ch3104} \cite{harrington2001,jackson1999}.

\paragraph{Resolution identity and projector audit.}
Eq.~\eqref{eq:ch3.VSH-resolution-identity} states that partial projectors sum to the identity on \(\mathcal{F}_{\mathrm{rad}}\) only when
all sectors are retained \cite{hansen1988,kato1995}. Truncation at \(\ell_{\max}=2\) produces a proper subprojector whose defect equals the
tail energy in Eq.~\eqref{eq:ch3.wp3102-partial-sum-error} \cite{harrington2001,stratton1941}. Campaigns should archive both the declared
\(\ell_{\max}\) and the measured tail so completeness claims remain falsifiable \cite{devaney2004,harrington2001}.

\paragraph{Shell quadrature as an inner-product experiment.}
Eq.~\eqref{eq:ch3.radiative-inner-product-def} is an integral over \(\Gamma_R\) with flux weighting \(\sin\theta\) and outgoing phasor conventions
\cite{stratton1941,harrington2001}. Subsection~\ref{sec:ch3102} treats numerical quadrature as an \emph{experiment on the inner product}: refine
\((\theta,\phi)\) samples until Eq.~\eqref{eq:ch3.wp3102-orthogonality-integral} stabilizes below \(10^{-4}\,\mathrm{W}\)
\cite{hansen1988,kato1995}. Mathematically, quadrature error is a discretization of the continuous orthogonality theorem; physically, it is the
difference between a pattern code that respects flux meters and one that reports spurious cross-coupling \cite{poynting1884,harrington2001}.
Table~\ref{tab:ch3.3102-numerical-protocol} orders the refinement loop: fix flux weights first, then cross integrals, then grid density
\cite{stratton1941,hansen1988}.

\paragraph{Cross-power \(P_{12}\) and basis diagnostics.}
Define the cross-power \(P_{12}=\Re\{\widehat{c}_{10}^{*}c_{20}\}\) in the eigenbasis of Theorem~\ref{thm:ch3.vsh-orthogonality}
\cite{hansen1988,harrington2001}. Orthogonality across distinct \(\ell\) forces \(P_{12}=0\) for the normalized pair in the worked problem
\cite{stratton1941,poynting1884}. A vendor report that shows \(|P_{12}|\sim10^{-2}\,\mathrm{W}\) while claiming orthonormal VSH modes is
almost always suffering normalization mismatch Eqs.~\eqref{eq:ch3.VSH-unit-flux-normalization}--\eqref{eq:ch3.VSH-unit-energy-norm} or reactive
admixture before \(\Pi_{\mathrm{rad}}\) \cite{harrington2001,stratton1941}. Subsection~\ref{sec:ch3102} elevates \(P_{12}\) as a first-line
diagnostic because it is cheaper than full coupling-matrix analysis yet sensitive to the same implementation bugs \cite{devaney2004,hansen1988}.

\paragraph{Completeness rate and smooth versus sharp patterns.}
Theorem~\ref{thm:ch3.vsh-completeness} guarantees convergence of partial sums in \(\|\cdot\|_{\mathrm{rad}}\) \cite{kato1995,hansen1988}.
The \(0.08\,\mathrm{W}\) tail at \(\ell_{\max}=2\) is therefore a \emph{rate} question: smooth broadside-dominated patterns decay rapidly in
\(\ell\), while sharp features, grating sidelobes, and compact vortex cores require large \(\ell_{\max}\) \cite{harrington2001,devaney2004,jackson1999}.
Subsection~\ref{sec:ch3102} uses the tail as a falsifiable completeness certificate: publish \(\ell_{\max}\), measured tail watts, and the
convergence slope as \(\ell_{\max}\) increases \cite{stratton1941,harrington2001}. Figure~\ref{fig:ch3.3102-truncation-tail} is the schematic
reminder that completeness is an asymptotic statement, not a single finite truncation \cite{kato1995,hansen1988}.

\paragraph{Worked extension (three-sector orthogonality).}
Extend the worked pair to include \(\widehat{\mathbf{u}}_{11}^{(\mathbf{M})}\) at \(\ell=1\), \(m=1\): orthogonality to \((2,0)\) still holds by
distinct \(\ell\), while orthogonality to \((1,0)\) within \(\ell=1\) requires fixed \(m\) and \(\sigma\) quantum numbers
\cite{jackson1999,hansen1988}. Subsection~\ref{sec:ch3102} uses that extension to train readers that \emph{same-\(\ell\)} blocks need explicit
\(m,\sigma\) resolution before watt sums are interpreted \cite{harrington2001,stratton1941,belinfante1940}.

\paragraph{Reactive contamination and false coupling.}
Near-field reactive content violates Eq.~\eqref{eq:ch3.wp3102-orthogonality-integral} when \(\Pi_{\mathrm{rad}}\) is omitted: evanescent admixture
creates non-zero cross-integrals that mimic physical coupling \cite{stratton1941,harrington2001}. The corrective order is outgoing filter, then flux
normalization, then orthogonality test \cite{harrington2001,jackson1999}. Subsection~\ref{sec:ch3102} records that order as part of the FIG+PROB
contract because orthogonality failures are implementation bugs until proven otherwise \cite{hansen1988,kato1995}.

\paragraph{Numerical acceptance protocol (consolidated).}
Subsection~\ref{sec:ch3102} closes with the executable protocol
\begin{equation}
\label{eq:ch3.wp3102-acceptance-protocol}
\mathcal{A}_{\mathrm{ortho}}=1
\iff
\left(
|P_{12}|<10^{-4}\,\mathrm{W}\
\wedge\
\|\mathbf{R}_{\ell_{\max}}\mathbf{v}_\omega\|_{\mathrm{rad}}^{2}\le\tau_{\mathrm{tail}}\
\wedge\
\mathcal{Q}_{\Gamma}=1
\right),
\end{equation}
with \(\mathcal{Q}_{\Gamma}=1\) denoting converged \(\Gamma_R\) quadrature for Eq.~\eqref{eq:ch3.radiative-inner-product-def}, importing Eqs.~\eqref{eq:ch3.wp3102-orthogonality-integral} and~\eqref{eq:ch3.wp3102-partial-sum-error}
\cite{hansen1988,kato1995,harrington2001,stratton1941}. Equation~\eqref{eq:ch3.wp3102-acceptance-protocol} is the FIG+PROB acceptance certificate that Subsection~\ref{sec:ch3101} watt partitions presuppose
\cite{poynting1884,harrington2001,devaney2004}.

\paragraph{Gram matrix spot check on three-mode truncation.}
Extend the worked pair to \(\ell_{\max}=2\) with normalized Gram \(\mathbf{G}_{ij}=\langle\widehat{\mathbf{u}}_{i},\widehat{\mathbf{u}}_{j}\rangle_{\mathrm{rad}}\)
\cite{hansen1988,kato1995,harrington2001}. Theorem~\ref{thm:ch3.vsh-orthogonality} forces \(\mathbf{G}\) diagonal in the eigenbasis; off-diagonal entries measure implementation error, not physics coupling across \(\ell\)
\cite{stratton1941,poynting1884}. Subsection~\ref{sec:ch3102} recommends publishing \(\max_{i\neq j}|G_{ij}|\) alongside Eq.~\eqref{eq:ch3.wp3102-acceptance-protocol} as a compact quality metric for pattern codes
\cite{harrington2001,devaney2004}.

\paragraph{Tail convergence rate estimate.}
Increasing \(\ell_{\max}\) from \(2\) to \(4\) and halving the tail in Eq.~\eqref{eq:ch3.wp3102-partial-sum-error} suggests rapid VSH decay for smooth patterns \cite{kato1995,hansen1988,harrington2001}. Subsection~\ref{sec:ch3102} records that slope as a \emph{completeness rate} certificate: publish tail values at two successive \(\ell_{\max}\) values so readers can judge whether truncation is adequate for the claimed angular resolution
\cite{devaney2004,stratton1941,jackson1999}.

\paragraph{Projection operator algebra on $\Gamma_R$.}
Let \(\mathbf{P}_{\ell m\sigma}\) denote the Riesz projector onto \(\widehat{\mathbf{u}}_{\ell m}^{(\sigma)}\). Theorem~\ref{thm:ch3.vsh-orthogonality} states \(\mathbf{P}_{\ell m\sigma}\mathbf{P}_{\ell'm'\sigma'}=\delta_{\ell\ell'}\delta_{mm'}\delta_{\sigma\sigma'}\mathbf{P}_{\ell m\sigma}\)
\cite{kato1995,hansen1988}. Subsection~\ref{sec:ch3102} interprets Eq.~\eqref{eq:ch3.wp3102-orthogonality-integral} as \(\langle\mathbf{P}_{10}\mathbf{v},\mathbf{P}_{20}\mathbf{v}\rangle=0\) on \(\Gamma_R\)
\cite{stratton1941,harrington2001}. Truncation projector \(\mathbf{R}_{\ell_{\max}}\) is a partial sum of \(\mathbf{P}_{\ell m\sigma}\); its defect is exactly the tail in Eq.~\eqref{eq:ch3.wp3102-partial-sum-error}
\cite{kato1995,devaney2004,poynting1884}.

\paragraph{Vendor code regression suite.}
Subsection~\ref{sec:ch3102} recommends a three-test regression on any new pattern code: reproduce Eq.~\eqref{eq:ch3.wp3102-orthogonality-integral} on a known pair; reproduce Eq.~\eqref{eq:ch3.wp3102-partial-sum-error} tail scaling; satisfy Eq.~\eqref{eq:ch3.wp3102-acceptance-protocol}
\cite{hansen1988,harrington2001,stratton1941}. Failure at test one is almost always flux-weight or Condon--Shortley convention mismatch
\cite{jackson1999,hansen1988,devaney2004}.

Subsection~\ref{sec:ch3102} has verified orthogonality and completeness numerically through
Eqs.~\eqref{eq:ch3.wp3102-orthogonality-integral}--\eqref{eq:ch3.wp3102-partial-sum-error}
\cite{hansen1988,stratton1941,harrington2001}. Subsection~\ref{sec:ch3103} exercises continuous-spectrum normalization and mode counting
\cite{jackson1999}.


\subsection{Continuous-spectrum normalization and mode counting}
\label{sec:ch3103}
% TOC tag: [FIG+PROB][USE SS3.5]

\emph{Continuous-spectrum normalization} fixes watt measure on plane-wave and direction charts so Parseval integrals close
\cite{stratton1941,harrington2001,jackson1999}. \emph{Mode counting} tracks how many discrete VSH slots exist at finite \(\ell_{\max}\) versus
open-region density-of-states estimates \cite{hansen1988,harrington2001}. Subsection~\ref{sec:ch3103} is a FIG+PROB ledger exercise importing
Section~\ref{sec:ch35} normalization and Eq.~\eqref{eq:ch3.cumulative-disc-DOS} without re-derivation \cite{stratton1941,hansen1988}.

\paragraph{Standing imports.}
Flux-fixed Parseval Eq.~\eqref{eq:ch3.continuous-parseval-flux-fixed}, plane-wave delta normalization
Eq.~\eqref{eq:ch3.plane-wave-delta-normalization}, shell power density Eq.~\eqref{eq:ch3.shell-power-density-normalized},
cumulative discrete DOS Eq.~\eqref{eq:ch3.cumulative-disc-DOS}, and accessible watt budget Eq.~\eqref{eq:ch3.accessible-watt-budget} enter by
reference \cite{harrington2001,stratton1941,hansen1988}. Subsection~\ref{sec:ch3103} does not re-derive measure Jacobians
\cite{stratton1941,jackson1999}.

\paragraph{Problem.}
A broadside-dominated pattern has discrete coefficients \(\widehat{c}_{\ell m}^{(\sigma)}\) with \(\ell_{\max}=3\) and total
\(\mathcal{P}_{\mathrm{rad}}=1\,\mathrm{W}\). A parallel plane-wave chart gives \(\widetilde{\mathbf{A}}(k_{x},k_{y})\) on a propagating
patch \(\mathcal{K}_{\mathrm{prop}}\).
\begin{enumerate}
\item[(a)] verify discrete Parseval Eq.~\eqref{eq:ch3.wp3103-parseval-discrete};
\item[(b)] count \(N_{\mathrm{disc}}(\ell_{\max}=3)\) from Eq.~\eqref{eq:ch3.cumulative-disc-DOS};
\item[(c)] compare to a continuous direction integral and explain the \(\ell_{\max}\) cutoff as a laboratory band-limit, not a physics cutoff.
\end{enumerate}

\begin{figure}[t]
\centering
\ChOneFigBlock{%
\begin{tikzpicture}[ch1fig, node distance=7mm and 10mm]
  \node[ch1box] (disc) {discrete $\widehat{c}_{\ell m}^{(\sigma)}$};
  \node[ch1box, right=12mm of disc] (cont) {continuous $\widetilde{\mathbf{A}}$};
  \node[ch1box, right=12mm of cont] (par) {Parseval};
  \node[ch1box, below=11mm of disc] (dos) {Eq.~\eqref{eq:ch3.cumulative-disc-DOS}};
  \node[ch1box, below=11mm of cont] (mu) {Eq.~\eqref{eq:ch3.continuous-parseval-flux-fixed}};
  \draw[ch1flow] (disc) -- (par);
  \draw[ch1flow] (cont) -- (par);
  \draw[ch1flow] (dos) -- (disc);
  \draw[ch1flow] (mu) -- (cont);
\end{tikzpicture}%
}
\caption{Worked problem~\ref{sec:ch3103}: discrete VSH and continuous plane-wave ledgers must agree on \(\mathcal{P}_{\mathrm{rad}}\) under Section~\ref{sec:ch35} meters.}
\label{fig:ch3.3103-spectrum-ledger}
\end{figure}

\paragraph{Step 1 ---Discrete Parseval (HOME).}
With normalized coefficients,
\begin{equation}
\label{eq:ch3.wp3103-parseval-discrete}
\mathcal{P}_{\mathrm{rad}}
=
\sum_{\ell=0}^{3}\sum_{m=-\ell}^{\ell}\sum_{\sigma}
\bigl|\widehat{c}_{\ell m}^{(\sigma)}\bigr|^{2}
=
1\,\mathrm{W},
\end{equation}
by Eq.~\eqref{eq:ch3.mode-power-from-coefficient} \cite{poynting1884,harrington2001,hansen1988}. Mathematically, the sum is finite because
\(\ell_{\max}=3\); physically, the watt meter reads the full discrete ledger, not a fitted subset \cite{stratton1941,harrington2001}.

\paragraph{Step 2 ---Mode count (HOME).}
Eq.~\eqref{eq:ch3.cumulative-disc-DOS} gives
\begin{equation}
\label{eq:ch3.wp3103-DOS-count}
N_{\mathrm{disc}}(3)
=
\sum_{\ell=0}^{3}2(2\ell+1)
=
32
\end{equation}
independent \(\sigma\)-doubled channels \cite{hansen1988,jackson1999,harrington2001}. Each channel is one normalized eigenfield slot at
\(\ell_{\max}=3\); accessibility may forbid some slots via Eq.~\eqref{eq:ch3.accessible-watt-budget} even when mathematically present
\cite{harrington2001,stratton1941}.

\paragraph{Step 3 ---Continuous chart closure.}
On \(\mathcal{K}_{\mathrm{prop}}\),
\(\int_{\mathcal{K}_{\mathrm{prop}}}|\widetilde{\mathbf{A}}|^{2}\,dk_{x}\,dk_{y}\) must equal \(1\,\mathrm{W}\) only after
Eq.~\eqref{eq:ch3.plane-wave-delta-normalization} and outgoing mask \cite{stratton1941,jackson1999,harrington2001}. Subsection~\ref{sec:ch3103}
uses Figure~\ref{fig:ch3.3103-spectrum-ledger} to stress that continuous integrals are limit cases of dense discrete sums, not separate power
laws \cite{hansen1988,harrington2001}.

\paragraph{Step 4 ---Band-limit versus physics.}
\(\ell_{\max}=3\) is an instrument band-limit; Theorem~\ref{thm:ch3.vsh-completeness} requires \(\ell_{\max}\to\infty\) for exact representation
\cite{kato1995,hansen1988}. Subsection~\ref{sec:ch3103} records \(N_{\mathrm{disc}}(3)=32\) as \emph{accessible chart size}, not as the number of
propagating plane waves in free space \cite{jackson1999,stratton1941,harrington2001}.

\paragraph{Step 5 ---Shell density check.}
Eq.~\eqref{eq:ch3.shell-power-density-normalized} ties \(|\widehat{c}|^{2}\) to \(S^{2}\) plots; mis-weighted \(|\E|^{2}\) without
\(\sin\theta\) violates flux closure \cite{stratton1941,harrington2001,jackson1999}. Subsection~\ref{sec:ch3103} flags that as a normalization
error, not a DOS ambiguity \cite{hansen1988,harrington2001}.

\paragraph{Physical meaning (closure).}
Discrete mode counting and continuous Parseval are dual bookkeeping on \(\mathcal{R}_\omega\) \cite{stratton1941,hansen1988,harrington2001}.
Laboratories must publish \(\ell_{\max}\), measure, and accessibility mask with every watt table \cite{harrington2001,devaney2004}.

\begin{table}[t]
\centering
\caption{Subsection~\ref{sec:ch3103} discrete versus continuous ledger at \(\ell_{\max}=3\).}
\label{tab:ch3.3103-ledger}
\small
\begin{tabular}{@{}p{0.28\linewidth}p{0.34\linewidth}p{0.30\linewidth}@{}}
\toprule
Chart & Count / closure & Value \\
\midrule
VSH discrete & Eq.~\eqref{eq:ch3.wp3103-DOS-count} & $32$ channels \\
VSH power & Eq.~\eqref{eq:ch3.wp3103-parseval-discrete} & $1\,\mathrm{W}$ \\
Plane-wave & Eq.~\eqref{eq:ch3.continuous-parseval-flux-fixed} & matched \\
\bottomrule
\end{tabular}
\end{table}

\paragraph{Worked follow-up (accessibility).}
If geometry annihilates half the \(m\) sectors, accessible watts sum to less than \(1\,\mathrm{W}\) even when
Eq.~\eqref{eq:ch3.wp3103-parseval-discrete} holds mathematically \cite{harrington2001,stratton1941}. Subsection~\ref{sec:ch3103} separates
\emph{mathematical} channel count from \emph{reportable} channel count \cite{hansen1988,harrington2001}.

\paragraph{Step 6 --- Plane-wave patch Jacobian.}
On \(\mathcal{K}_{\mathrm{prop}}\), each cell area \(\Delta k_{x}\Delta k_{y}\) must include propagating mask and normalization factor from
Eq.~\eqref{eq:ch3.plane-wave-delta-normalization} \cite{stratton1941,jackson1999,harrington2001}. Subsection~\ref{sec:ch3103} treats missing
Jacobian as a continuous-spectrum normalization bug, not as physics \cite{hansen1988,harrington2001}.

\paragraph{Step 7 --- Shell versus lateral ledgers.}
Eq.~\eqref{eq:ch3.shell-power-density-normalized} on \(S^{2}\) and lateral FFT ledger Eq.~\eqref{eq:ch3.plane-wave-basis-resolution} must
agree after \(\mathbf{T}\) from Subsection~\ref{sec:ch383} \cite{hansen1988,stratton1941}. A \(3\%\) mismatch in the worked campaign signals
band-limit or accessibility inconsistency \cite{harrington2001,devaney2004}.

\begin{figure}[t]
\centering
\ChOneFigBlock{%
\begin{tikzpicture}[ch1fig, node distance=7mm and 9mm]
  \node[ch1box] (lat) {lateral FFT};
  \node[ch1box, right=10mm of lat] (s2) {$S^{2}$ density};
  \node[ch1box, right=10mm of s2] (vsh) {VSH $\widehat{c}$};
  \node[ch1box, below=10mm of s2] (par) {Eq.~\eqref{eq:ch3.continuous-parseval-flux-fixed}};
  \draw[ch1flow] (lat) -- (s2) -- (vsh);
  \draw[ch1flow] (par) -- (vsh);
\end{tikzpicture}%
}
\caption{Subsection~\ref{sec:ch3103} triple ledger: lateral, shell, and VSH charts must close Parseval on the same $\mathcal{P}_{\mathrm{rad}}$.}
\label{fig:ch3.3103-triple-ledger}
\end{figure}

\paragraph{Step 8 --- DOS versus accessible count.}
\(N_{\mathrm{disc}}(3)=32\) counts mathematical channels; accessible count may be smaller after
Eq.~\eqref{eq:ch3.accessible-watt-budget} \cite{harrington2001,stratton1941}. Subsection~\ref{sec:ch3103} publishes both integers in
campaign metadata \cite{hansen1988,harrington2001}.

\paragraph{Step 9 --- Truncation of continuous integral.}
Replacing the direction integral by a finite \(k\)-grid introduces discretization error analogous to \(\ell_{\max}\) truncation
\cite{jackson1999,harrington2001,stratton1941}. Refining the grid until watt error falls below \(10^{-3}\,\mathrm{W}\) mirrors
Subsection~\ref{sec:ch3102} angular refinement \cite{hansen1988,kato1995}.

\begin{table}[t]
\centering
\caption{Subsection~\ref{sec:ch3103} normalization acceptance thresholds.}
\label{tab:ch3.3103-acceptance}
\small
\begin{tabular}{@{}p{0.34\linewidth}p{0.56\linewidth}@{}}
\toprule
Test & Threshold \\
\midrule
Discrete Parseval & Eq.~\eqref{eq:ch3.wp3103-parseval-discrete} $\pm10^{-3}\,\mathrm{W}$ \\
Chart cross-check & $3\%$ max mismatch \\
DOS audit & Eq.~\eqref{eq:ch3.wp3103-DOS-count} $=32$ \\
\bottomrule
\end{tabular}
\end{table}

\paragraph{Discrete DOS and polarization doubling.}
Equation~\eqref{eq:ch3.wp3103-DOS-count} counts \(N_{\mathrm{disc}}(3)=32\) channels because each \(\ell\) carries \(2\ell+1\) magnetic
quantum numbers and two polarization tracks \(\sigma\) in the normalized VSH chart \cite{hansen1988,jackson1999}. Mathematically, the count is
the dimension of a truncated coefficient vector; physically, it is the size of a square Parseval ledger at \(\ell_{\max}=3\)
\cite{harrington2001,stratton1941}. Accessibility may forbid some rows of that ledger even when the mathematical count is \(32\)
\cite{harrington2001,stratton1941}; Subsection~\ref{sec:ch3103} therefore publishes both \(N_{\mathrm{disc}}\) and the accessible subset.

\paragraph{Plane-wave Jacobian and triple-ledger closure.}
Figure~\ref{fig:ch3.3103-triple-ledger} orders lateral FFT, shell density, and VSH coefficients as three views of one \(\mathcal{P}_{\mathrm{rad}}\)
\cite{jackson1999,hansen1988}. Eq.~\eqref{eq:ch3.plane-wave-direction-jacobian} is the geometric bridge between FFT cells and \(S^{2}\) plots;
omitting it produces the \(3\%\) chart mismatch flagged in the acceptance table \cite{stratton1941,harrington2001}. Continuous Parseval
Eq.~\eqref{eq:ch3.continuous-parseval-flux-fixed} is the continuum limit of Eq.~\eqref{eq:ch3.wp3103-parseval-discrete} as \(\ell_{\max}\) and
FFT grids refine jointly \cite{kato1995,hansen1988}.

\paragraph{Band-limit semantics in campaigns.}
\(\ell_{\max}=3\) is an instrument band-limit, not a propagation cutoff in free space \cite{jackson1999,stratton1941}. Theorem~\ref{thm:ch3.vsh-completeness}
requires \(\ell_{\max}\to\infty\) for exact fields; the worked ledger declares what a finite instrument can resolve
\cite{kato1995,harrington2001}. Peer reviewers should treat undeclared \(\ell_{\max}\) as an incomplete normalization statement comparable to
missing flux-meter calibration \cite{poynting1884,harrington2001,devaney2004}.

\paragraph{Explicit DOS arithmetic at \(\ell_{\max}=3\).}
Summing \(2(2\ell+1)\) for \(\ell=0,1,2,3\) yields \(1+6+10+14=31\) magnetic quantum numbers per polarization track; doubling for \(\sigma\) gives
\(N_{\mathrm{disc}}(3)=32\) in Eq.~\eqref{eq:ch3.wp3103-DOS-count} \cite{hansen1988,jackson1999}. Each channel is one slot in the Parseval sum
Eq.~\eqref{eq:ch3.wp3103-parseval-discrete}; accessibility may zero rows without changing the mathematical count \cite{harrington2001,stratton1941}.
Subsection~\ref{sec:ch3103} uses that arithmetic to prevent ``mode counting by plotting peaks'' errors in vendor reports \cite{devaney2004,hansen1988}.

\paragraph{Continuous integral discretization mirror.}
Replacing the direction integral by a finite \(k\)-grid introduces error analogous to \(\ell_{\max}\) truncation \cite{jackson1999,harrington2001,stratton1941}.
Refine \(\Delta k_{x},\Delta k_{y}\) until watt error falls below \(10^{-3}\,\mathrm{W}\), mirroring angular grid refinement in
Subsection~\ref{sec:ch3102} \cite{hansen1988,kato1995}. Figure~\ref{fig:ch3.3103-triple-ledger} is the audit diagram: lateral, shell, and VSH ledgers must
close on the same \(\mathcal{P}_{\mathrm{rad}}\) within Table~\ref{tab:ch3.3103-acceptance} thresholds \cite{poynting1884,harrington2001}.

\paragraph{Shell density versus FFT magnitude plots.}
Eq.~\eqref{eq:ch3.shell-power-density-normalized} weights \(|\E|^{2}\) by \(\sin\theta\) on \(S^{2}\); raw FFT magnitudes omit that weight
\cite{stratton1941,harrington2001,jackson1999}. Subsection~\ref{sec:ch3103} treats \(3\%\) cross-chart mismatch as a normalization failure, not as
physics ambiguity \cite{hansen1988,harrington2001}. One must compute shell density and discrete Parseval on the same campaign file before
claiming agreement between compact-range FFTs and spherical harmonic fits \cite{devaney2004,stratton1941}.

\paragraph{Worked ledger closure (numeric instance).}
Suppose the \(\ell_{\max}=3\) discrete sum in Eq.~\eqref{eq:ch3.wp3103-parseval-discrete} closes to \(1.000\,\mathrm{W}\) while the lateral FFT integral on \(\mathcal{K}_{\mathrm{prop}}\), weighted by Eq.~\eqref{eq:ch3.plane-wave-direction-jacobian}, closes to \(1.028\,\mathrm{W}\) before refinement
\cite{jackson1999,harrington2001,stratton1941}. Subsection~\ref{sec:ch3103} diagnoses the \(2.8\%\) gap as discretization or evanescent leakage, not as new physics: refine the \(k\)-grid, enforce \(\Pi_{\mathrm{rad}}\), and re-test Table~\ref{tab:ch3.3103-acceptance} until the mismatch falls below \(3\%\)
\cite{hansen1988,kato1995,devaney2004}. The DOS count \(N_{\mathrm{disc}}(3)=32\) is unchanged throughout because channel cardinality is a property of the declared \(\ell_{\max}\), not of the FFT grid density
\cite{hansen1988,jackson1999,harrington2001}.

\paragraph{Accessibility split on the discrete ledger.}
When geometry annihilates half the \(m\) sectors, accessible watts may sum to \(0.50\,\mathrm{W}\) even when Eq.~\eqref{eq:ch3.wp3103-parseval-discrete} holds at \(1\,\mathrm{W}\) mathematically
\cite{harrington2001,stratton1941}. Subsection~\ref{sec:ch3103} requires publishing both totals: mathematical Parseval on the full label set and reportable Parseval on \(\mathcal{L}_{\mathrm{acc}}\) from Eq.~\eqref{eq:ch3.accessible-watt-budget}
\cite{hansen1988,devaney2004,jackson1999}.

Subsection~\ref{sec:ch3103} has exercised continuous-spectrum normalization and mode counting through
Eqs.~\eqref{eq:ch3.wp3103-parseval-discrete}--\eqref{eq:ch3.wp3103-DOS-count}
\cite{stratton1941,hansen1988,harrington2001}. Subsection~\ref{sec:ch3104} demonstrates symmetry breaking and mode coupling at a conducting
boundary \cite{jackson1999}.


\subsection{Symmetry breaking and mode coupling}
\label{sec:ch3104}
% TOC tag: [FIG+PROB][USE SS3.6.3]

\emph{Symmetry breaking} means a boundary \(\Gamma_{b}\) reduces \(\mathrm{SO}(3)\) invariance so pure VSH labels no longer diagonalize the
scattering problem \cite{stratton1941,jackson1999,harrington2001}. \emph{Mode coupling} is the off-diagonal mixing induced by
Eq.~\eqref{eq:ch3.boundary-scattering-mixing-matrix} \cite{hansen1988,harrington2001}. Subsection~\ref{sec:ch3104} is a FIG+PROB half-space
instance importing Subsection~\ref{sec:ch363} mixing and image laws without re-derivation \cite{stratton1941,jackson1999}.

\paragraph{Standing imports.}
Boundary mixing matrix Eq.~\eqref{eq:ch3.boundary-scattering-mixing-matrix}, image admixture Eq.~\eqref{eq:ch3.image-mode-admixture},
mutual coupling Gram Eq.~\eqref{eq:ch3.mutual-coupling-gram}, broken-symmetry partition Eq.~\eqref{eq:ch3.broken-symmetry-sector-partition},
and accessible label set Eq.~\eqref{eq:ch3.accessible-VSH-label-set} enter by reference \cite{harrington2001,stratton1941,hansen1988}. Subsection~\ref{sec:ch3104}
does not re-derive image theory \cite{jackson1999,stratton1941}.

\paragraph{Problem.}
An electric dipole \(\widehat{\mathbf{u}}_{10}^{(\mathbf{N})}\) radiates above a PEC half-space at height \(h=\lambda/4\). Open-region
coefficients \(\widehat{c}_{10}^{(\mathbf{N})}=1\) (normalized to \(1\,\mathrm{W}\) in free space) are mixed by \(\mathbf{S}_{b}\).
\begin{enumerate}
\item[(a)] form \(2\times2\) coupling block on \(\{\ell=1,\,m=0,\,\mathbf{N}\}\) and its image partner;
\item[(b)] diagonalize to find coupled eigenmodes and rotated watt partition;
\item[(c)] identify which labels remain accessible under parity partition Eq.~\eqref{eq:ch3.broken-symmetry-sector-partition}.
\end{enumerate}

\begin{figure}[t]
\centering
\ChOneFigBlock{%
\begin{tikzpicture}[ch1fig, node distance=7mm and 10mm]
  \node[ch1box] (src) {dipole $\widehat{\mathbf{u}}_{10}^{(\mathbf{N})}$};
  \node[ch1box, right=12mm of src] (Sb) {$\mathbf{S}_{b}$};
  \node[ch1box, right=12mm of Sb] (mix) {coupled modes};
  \node[ch1box, below=11mm of Sb] (pec) {PEC half-space};
  \node[ch1box, below=11mm of mix] (gram) {Eq.~\eqref{eq:ch3.mutual-coupling-gram}};
  \draw[ch1flow] (src) -- (Sb) -- (mix);
  \draw[ch1flow] (pec) -- (Sb);
  \draw[ch1flow] (gram) -- (mix);
\end{tikzpicture}%
}
\caption{Worked problem~\ref{sec:ch3104}: a PEC plane induces mode coupling through $\mathbf{S}_{b}$ before watts are summed on accessible sectors.}
\label{fig:ch3.3104-half-space-mixing}
\end{figure}

\paragraph{Step 1 ---Image admixture (USE §3.6.3).}
Eq.~\eqref{eq:ch3.image-mode-admixture} introduces an image coefficient with phase \(\exp(-2\jj k_{0}h)\) at \(h=\lambda/4\)
\cite{jackson1999,stratton1941}. Mathematically, the open-region label pairs with its image; physically, the ground plane creates
constructive/destructive interference in exported power \cite{harrington2001,jackson1999}. Subsection~\ref{sec:ch3104} uses that phase to
populate off-diagonal entries of \(\mathbf{S}_{b}\) \cite{hansen1988,stratton1941}.

\paragraph{Step 2 ---Gram matrix (HOME).}
On the two-mode block,
\begin{equation}
\label{eq:ch3.wp3104-mixing-gram}
\mathbf{G}
=
\begin{pmatrix}
1 & g_{12}\\[2pt]
g_{12}^{*} & 1
\end{pmatrix},
\qquad
g_{12}=\bigl[\mathbf{S}_{b}\bigr]_{(10),(10')},
\end{equation}
importing Eq.~\eqref{eq:ch3.mutual-coupling-gram} \cite{harrington2001,hansen1988,stratton1941}. With \(h=\lambda/4\),
\(g_{12}=-\jj\) schematically (phase factor from image term) \cite{jackson1999,stratton1941}. Mathematically, \(\mathbf{G}\) is the
energy inner product matrix before diagonalization; physically, mutual coupling redistributes watts between source and image sectors
\cite{harrington2001,poynting1884}.

\paragraph{Step 3 ---Diagonalization and coupled powers (HOME).}
Unitary \(\mathbf{U}\) diagonalizes \(\mathbf{G}\) to eigenvalues \(\lambda_{\pm}=1\pm|g_{12}|\) \cite{hansen1988,kato1995}. Coupled powers
\begin{equation}
\label{eq:ch3.wp3104-coupled-power}
\mathcal{P}_{\pm}
=
|\alpha_{\pm}|^{2}\,\lambda_{\pm},
\qquad
\mathcal{P}_{+}+\mathcal{P}_{-}
=
\mathcal{P}_{\mathrm{rad}},
\end{equation}
with \(\alpha_{\pm}\) expansion coefficients in the diagonal basis \cite{poynting1884,harrington2001,stratton1941}. Equation
\eqref{eq:ch3.wp3104-coupled-power} shows symmetry breaking: \emph{eigenmodes of coupling}, not original \(\ell m\) labels, diagonalize watts
\cite{harrington2001,hansen1988}. One must report \(\mathbf{U}\) when publishing ``dominant mode'' after a ground plane
\cite{stratton1941,devaney2004}.

\paragraph{Step 4 ---Accessibility partition.}
Eq.~\eqref{eq:ch3.broken-symmetry-sector-partition} may remove image-sector labels from \(\mathcal{L}_{\mathrm{acc}}\) when probes face only
the upper half-space \cite{harrington2001,stratton1941}. Subsection~\ref{sec:ch3104} counts reportable watts only on accessible rows of
Eq.~\eqref{eq:ch3.wp3104-coupled-power} \cite{hansen1988,harrington2001}.

\paragraph{Step 5 ---Comparison to free space.}
Without \(\Gamma_{b}\), Eq.~\eqref{eq:ch3.wp3101-power-partition} would apply with a single \(\widehat{c}_{10}^{(\mathbf{N})}\)
\cite{hansen1988,stratton1941}. Boundary presence forces \(\mathbf{S}_{b}\) \emph{before} any Parseval sum \cite{harrington2001,jackson1999}.
Subsection~\ref{sec:ch3104} quantifies that workflow difference as the Chapter~\ref{ch:03} symmetry-breaking diagnostic
\cite{stratton1941,harrington2001}.

\paragraph{Physical meaning (closure).}
Half-space problems are mixed continua, not new eigenvalue ladders \cite{jackson1999,stratton1941}. Coupling matrices are the executable form
of symmetry breaking; ignoring \(\mathbf{S}_{b}\) misattributes distorted patterns to source spectra \cite{harrington2001,hansen1988}.

\begin{table}[t]
\centering
\caption{Subsection~\ref{sec:ch3104} half-space coupling ledger (schematic).}
\label{tab:ch3.3104-coupling-ledger}
\small
\begin{tabular}{@{}p{0.24\linewidth}p{0.30\linewidth}p{0.34\linewidth}@{}}
\toprule
Basis & Object & Role \\
\midrule
Open VSH & $\mathbf{S}_{b}$ & Eq.~\eqref{eq:ch3.boundary-scattering-mixing-matrix} \\
Coupled & $\mathbf{G}$ & Eq.~\eqref{eq:ch3.wp3104-mixing-gram} \\
Diagonal & $\mathcal{P}_{\pm}$ & Eq.~\eqref{eq:ch3.wp3104-coupled-power} \\
\bottomrule
\end{tabular}
\end{table}

\paragraph{Peer-review note.}
Publish \(h/\lambda\), \(\mathbf{S}_{b}\) block, \(\mathbf{G}\) eigenvalues, and accessible partition alongside far-zone plots
\cite{harrington2001,stratton1941,hansen1988,devaney2004}. Subsection~\ref{sec:ch3104} treats missing coupling data as a symmetry-breaking
artifact, not a completed boundary problem \cite{harrington2001,jackson1999}.

\paragraph{Step 6 --- Height sweep sensitivity.}
Varying \(h/\lambda\) rotates phase of \(g_{12}\) in Eq.~\eqref{eq:ch3.wp3104-mixing-gram}, modulating \(\mathcal{P}_{\pm}\) while preserving
total watts only if \(\mathbf{G}\) is recomputed \cite{jackson1999,harrington2001,stratton1941}. Subsection~\ref{sec:ch3104} warns against
using free-space \(\widehat{c}\) tables at multiple heights without re-diagonalizing \(\mathbf{G}\) \cite{harrington2001,devaney2004}.

\paragraph{Step 7 --- Mutual coupling versus mixing.}
Eq.~\eqref{eq:ch3.mutual-coupling-gram} generalizes the \(2\times2\) block to multi-element arrays near \(\Gamma_{b}\)
\cite{harrington2001,stratton1941,hansen1988}. The half-space lesson scales: report eigenmodes of \(\mathbf{G}\), not raw VSH labels
\cite{harrington2001,jackson1999}.

\begin{figure}[t]
\centering
\ChOneFigBlock{%
\begin{tikzpicture}[ch1fig, node distance=7mm and 9mm]
  \node[ch1box] (h) {$h/\lambda$};
  \node[ch1box, right=10mm of h] (g) {$g_{12}(h)$};
  \node[ch1box, right=10mm of g] (P) {Eq.~\eqref{eq:ch3.wp3104-coupled-power}};
  \draw[ch1flow] (h) -- (g) -- (P);
\end{tikzpicture}%
}
\caption{Subsection~\ref{sec:ch3104} height sensitivity: image phase drives coupled watt eigenvalues at fixed $\mathcal{P}_{\mathrm{rad}}$.}
\label{fig:ch3.3104-height-sweep}
\end{figure}

\paragraph{Step 8 --- Broken parity labels.}
Eq.~\eqref{eq:ch3.broken-symmetry-sector-partition} may classify image sectors as distinct parity classes even when \(\ell,m\) labels match
open-region book entries \cite{harrington2001,stratton1941}. Subsection~\ref{sec:ch3104} documents parity in metadata alongside \(\mathbf{G}\)
\cite{hansen1988,harrington2001}.

\paragraph{Step 9 --- Reportable watt sum.}
After accessibility, suppose only the symmetric coupled mode carries \(0.92\,\mathrm{W}\) and the antisymmetric mode is probe-invisible
\cite{harrington2001,stratton1941}. Reportable power can fall below mathematical \(\mathcal{P}_{\mathrm{rad}}\) even when coupling is strong
\cite{hansen1988,harrington2001,jackson1999}.

\begin{table}[t]
\centering
\caption{Subsection~\ref{sec:ch3104} boundary-coupling acceptance checklist.}
\label{tab:ch3.3104-acceptance}
\small
\begin{tabular}{@{}p{0.08\linewidth}p{0.50\linewidth}p{0.34\linewidth}@{}}
\toprule
Step & Deliverable & Equation \\
\midrule
1 & $\mathbf{S}_{b}$ block & Eq.~\eqref{eq:ch3.boundary-scattering-mixing-matrix} \\
2 & Gram & Eq.~\eqref{eq:ch3.wp3104-mixing-gram} \\
3 & Coupled watts & Eq.~\eqref{eq:ch3.wp3104-coupled-power} \\
4 & Accessibility & Eq.~\eqref{eq:ch3.accessible-VSH-label-set} \\
\bottomrule
\end{tabular}
\end{table}

\paragraph{Image phase and height sensitivity in operator language.}
At \(h=\lambda/4\), the image contribution enters Eq.~\eqref{eq:ch3.image-mode-admixture} with phase \(\exp(-2\jj k_{0}h)\), producing a
schematic off-diagonal \(g_{12}=-\jj\) in Eq.~\eqref{eq:ch3.wp3104-mixing-gram} \cite{jackson1999,stratton1941}. Mathematically, \(\mathbf{G}\) is
the energy Gram matrix of mixed open-region labels; physically, it is the matrix that must be diagonalized before watt sums acquire meaning above
a ground plane \cite{harrington2001,poynting1884}. Figure~\ref{fig:ch3.3104-height-sweep} warns that height sweeps rotate \(g_{12}\) while
preserving total watts only if \(\mathbf{G}\) is recomputed at each \(h/\lambda\) \cite{harrington2001,devaney2004}.

\paragraph{Coupled eigenmodes versus raw VSH labels.}
Equation~\eqref{eq:ch3.wp3104-coupled-power} expresses watts in the eigenbasis of \(\mathbf{G}\), not in the open-region VSH basis
\cite{hansen1988,kato1995}. Symmetry breaking therefore changes the \emph{diagonalization basis} for power, analogous to normal modes in
coupled circuits \cite{harrington2001,stratton1941}. Reporting \(\widehat{c}_{10}^{(\mathbf{N})}\) from a free-space table without \(\mathbf{S}_{b}\)
misattributes pattern distortion to source spectra \cite{jackson1999,harrington2001}. Table~\ref{tab:ch3.3104-acceptance} is the minimum publishable
bundle: \(\mathbf{S}_{b}\) block, Gram eigenvalues, coupled powers, and accessibility partition.

\paragraph{Accessibility and probe-visible watts.}
Eq.~\eqref{eq:ch3.broken-symmetry-sector-partition} may classify image sectors as probe-invisible when measurements face only the upper half-space
\cite{harrington2001,stratton1941}. Reportable power can fall below mathematical \(\mathcal{P}_{\mathrm{rad}}\) even when coupling is strong
\cite{hansen1988,harrington2001}. Subsection~\ref{sec:ch3104} separates \emph{mathematical} coupling from \emph{reportable} coupling so
Chapter~\ref{ch:04} realizability arguments inherit the correct watt ledger \cite{devaney2004,harrington2001}.

\paragraph{Half-space eigenmode diagonalization (expanded).}
Let \(\mathbf{G}\) in Eq.~\eqref{eq:ch3.wp3104-mixing-gram} have eigenvectors \(\mathbf{u}_{\pm}\) with eigenvalues \(\lambda_{\pm}=1\pm|g_{12}|\)
\cite{hansen1988,kato1995,harrington2001}. Coupled coefficients \(\alpha_{\pm}=\mathbf{u}_{\pm}^{\mathsf{H}}\widehat{\mathbf{c}}\) diagonalize the watt form
Eq.~\eqref{eq:ch3.wp3104-coupled-power} so \(\mathcal{P}_{\pm}=|\alpha_{\pm}|^{2}\lambda_{\pm}\) \cite{poynting1884,stratton1941}. Subsection~\ref{sec:ch3104} uses that diagonalization to teach that \emph{dominant mode} headlines after a ground plane refer to \(\mathbf{u}_{\pm}\), not to raw \(\ell,m\) labels
\cite{harrington2001,jackson1999,devaney2004}. Height sweeps rotate \(g_{12}\) and therefore exchange \(\mathcal{P}_{+}\) and \(\mathcal{P}_{-}\) while preserving \(\mathcal{P}_{\mathrm{rad}}\) only when \(\mathbf{G}\) is recomputed at each \(h/\lambda\)
\cite{jackson1999,harrington2001}.

\paragraph{Mutual coupling generalization pointer.}
Eq.~\eqref{eq:ch3.mutual-coupling-gram} lifts the schematic \(2\times2\) block to multi-port arrays near \(\Gamma_{b}\)
\cite{harrington2001,stratton1941,hansen1988}. Subsection~\ref{sec:ch3104} is the minimal FIG+PROB instance; Chapter~\ref{ch:04} owns full coupling ranks on finite supports
\cite{devaney2004,harrington2001}. The publishable bundle remains Table~\ref{tab:ch3.3104-acceptance}: \(\mathbf{S}_{b}\), \(\mathbf{G}\) spectrum, coupled powers, accessibility partition
\cite{stratton1941,jackson1999}.

\paragraph{PEC half-space limiting cases (pedagogical).}
As \(h/\lambda\to\infty\), image admixture vanishes and Eq.~\eqref{eq:ch3.wp3104-mixing-gram} approaches the identity on open-region labels, recovering Subsection~\ref{sec:ch3101} free-space ledger
\cite{jackson1999,harrington2001,stratton1941}. As \(h/\lambda\to0\), coupling strengthens and \(\mathcal{P}_{\pm}\) split diverges unless accessibility restricts reportable sectors
\cite{harrington2001,hansen1988}. Subsection~\ref{sec:ch3104} uses those limits to show symmetry breaking is a \emph{continuous} deformation of the watt diagonalization basis, not a separate modal alphabet
\cite{kato1995,devaney2004,poynting1884}.

\paragraph{Array-scale coupling preview.}
Multi-element arrays near \(\Gamma_{b}\) lift Eq.~\eqref{eq:ch3.wp3104-mixing-gram} to full \(\mathbf{G}\) from Eq.~\eqref{eq:ch3.mutual-coupling-gram} with dimension equal to the number of retained open-region labels
\cite{harrington2001,stratton1941,hansen1988}. Subsection~\ref{sec:ch3104} is the \(2\times2\) pedagogical core: diagonalize \(\mathbf{G}\) before watt sums; publish eigenvectors with eigenvalues \(\lambda_{\pm}\) from Eq.~\eqref{eq:ch3.wp3104-coupled-power}
\cite{kato1995,poynting1884,jackson1999}. Chapter~\ref{ch:04} owns quantitative realizability on those coupled eigenmodes \cite{devaney2004,harrington2001}.

\paragraph{Reportable versus mathematical watt reconciliation.}
When accessibility hides a coupled mode, reportable power can fall below mathematical \(\mathcal{P}_{\mathrm{rad}}\) even though Eq.~\eqref{eq:ch3.wp3104-coupled-power} closes in the full space
\cite{harrington2001,stratton1941}. Subsection~\ref{sec:ch3104} requires both totals in metadata so inheritance into Chapter~\ref{ch:04} does not confuse probe-visible watts with the full coupling ledger
\cite{hansen1988,devaney2004,jackson1999}.

\paragraph{Image height catalog for regression tests.}
Subsection~\ref{sec:ch3104} recommends archiving a small height catalog \(h/\lambda\in\{0.25,0.5,1,2\}\) with recomputed \(\mathbf{G}\) and \(\mathcal{P}_{\pm}\) for each entry
\cite{jackson1999,harrington2001,stratton1941}. Regression against free-space Subsection~\ref{sec:ch3101} at large \(h/\lambda\) validates implementation of Eq.~\eqref{eq:ch3.image-mode-admixture}
\cite{hansen1988,devaney2004,kato1995}. Coupled eigenvectors \(\mathbf{u}_{\pm}\) should be published alongside raw \(\widehat{c}\) labels for every height in the catalog
\cite{harrington2001,stratton1941}.

Subsection~\ref{sec:ch3104} has demonstrated symmetry breaking and mode coupling through
Eqs.~\eqref{eq:ch3.wp3104-mixing-gram}--\eqref{eq:ch3.wp3104-coupled-power}
\cite{stratton1941,jackson1999,harrington2001}. Section~\ref{sec:ch310} is complete: four worked examples now arithmetic-check eigenfield
construction, orthogonality, spectrum normalization, and boundary coupling on \(\mathcal{R}_\omega\) \cite{hansen1988,stratton1941}.

The four subsections of Section~\ref{sec:ch310} supply FIG+PROB verification:
Eqs.~\eqref{eq:ch3.wp3101-eigenfield-superposition}--\eqref{eq:ch3.wp3101-power-partition} in Subsection~\ref{sec:ch3101};
Eqs.~\eqref{eq:ch3.wp3102-orthogonality-integral}--\eqref{eq:ch3.wp3102-partial-sum-error} in Subsection~\ref{sec:ch3102};
Eqs.~\eqref{eq:ch3.wp3103-parseval-discrete}--\eqref{eq:ch3.wp3103-DOS-count} in Subsection~\ref{sec:ch3103}; and
Eqs.~\eqref{eq:ch3.wp3104-mixing-gram}--\eqref{eq:ch3.wp3104-coupled-power} in Subsection~\ref{sec:ch3104}
\cite{hansen1988,stratton1941,harrington2001,jackson1999}. Section~\ref{sec:ch3inherit} now records the chapter forward map to
Chapter~\ref{ch:04} \cite{devaney2004,harrington2001}.


\section{Chapter Summary and Forward Map}
\label{sec:ch3inherit}

Chapter~\ref{ch:03} is the sole inheritance section of the representation block: it records what was proved,
what remains deferred, and which exports Chapter~\ref{ch:04} must cite without proof. Subsection~\ref{sec:ch3111}
lists results established in the chapter. Subsection~\ref{sec:ch3112} catalogs concepts deferred to
Chapter~\ref{ch:04} and later volumes. Subsection~\ref{sec:ch3113} states the bridge map to realizability theory,
supplying VSH bases, normalization, truncation taxonomy, angular spectra, and basis-transformation rules by
reference only \cite{stratton1941,harrington2001,devaney2004}.


\subsection{Results established in Chapter 3}
\label{sec:ch3111}
% TOC tag: [DEFER]

Results are listed in Table~\ref{tab:ch3.results-ledger}; inheritance gates are fixed in Subsection~\ref{sec:ch303}.
Define the Chapter~\ref{ch:03} inheritance criterion
\begin{equation}
\label{eq:ch3.inheritance-admissibility}
\mathcal{D}_{\mathrm{ch3}}=1
\iff
\left(
\mathcal{D}_{\mathrm{ch2}}=1\
\wedge\
\mathbf{v}_\omega\in\mathcal{V}_{\mathrm{qual}}\
\wedge\
\mathcal{N}_{\mathrm{mtr}}=1
\right),
\end{equation}
with \(\mathcal{V}_{\mathrm{qual}}\subset\mathcal{F}_{\mathrm{rad}}\) the VSH- or spectrum-qualified export class and \(\mathcal{N}_{\mathrm{mtr}}=1\) the normalization-meter declaration of Section~\ref{sec:ch35},
importing \(\mathcal{D}_{\mathrm{ch2}}\) from Eq.~\eqref{eq:ch2.inheritance-admissibility} \cite{stratton1941,harrington2001}.
Let \(\mathcal{R}_{\mathrm{ch3}}\) denote results established in Sections~\ref{sec:ch30}--\ref{sec:ch310} under
\(\mathcal{D}_{\mathrm{ch3}}=1\). Table~\ref{tab:ch3.results-ledger} records the ledger without re-displaying proofs;
forward chapters cite anchors only \cite{hansen1988,stratton1941,harrington2001}.

\begin{table}[t]
\centering
\caption{Chapter~\ref{ch:03} results ledger: section ownership and primary anchors (no new mathematics).}
\label{tab:ch3.results-ledger}
\small
\begin{tabular}{@{}p{0.22\linewidth}p{0.14\linewidth}p{0.38\linewidth}p{0.18\linewidth}@{}}
\toprule
Result bundle & Section & Primary anchor(s) & Report stage \\
\midrule
Representation bridge, Hilbert closure &
§3.0--§3.1 &
Eqs.~\eqref{eq:ch3.F-rad-hilbert-closure}, \eqref{eq:ch3.radiative-inner-product-def} &
domain setup \\
Angular momentum, SO(3), Wigner &
§3.2 &
Eq.~\eqref{eq:ch3.J-total-maxwell-def}; Subsection~\ref{sec:ch324} &
symmetry labels \\
Vector Helmholtz, Debye &
§3.3 &
Eqs.~\eqref{eq:ch3.debye-E-representation}--\eqref{eq:ch3.debye-separated-expansion} &
construction \\
VSH definitions, eigenfields &
§3.4 &
Thms.~\ref{thm:ch3.vsh-orthogonality}, \ref{thm:ch3.vsh-completeness}, \ref{thm:ch3.vsh-J-eigenfields};
Eq.~\eqref{eq:ch3.normalized-coefficient-resolution} &
harmonic coords \\
Normalization, DOS &
§3.5 &
Eqs.~\eqref{eq:ch3.VSH-unit-flux-normalization}, \eqref{eq:ch3.mode-power-from-coefficient},
\eqref{eq:ch3.cumulative-disc-DOS} &
watt meters \\
Completeness, truncation, boundaries &
§3.6 &
Eqs.~\eqref{eq:ch3.unbounded-completeness-statement}, \eqref{eq:ch3.boundary-scattering-mixing-matrix},
\eqref{eq:ch3.accessible-VSH-label-set} &
symmetry break \\
Angular spectra, quantization &
§3.7 &
Eqs.~\eqref{eq:ch3.unbounded-angular-spectrum-def}, \eqref{eq:ch3.boundary-quantized-spectrum},
\eqref{eq:ch3.integer-ell-def} &
continuous/discrete \\
Basis transformations, OAM screens &
§3.8--§3.9 &
Eqs.~\eqref{eq:ch3.basis-change-matrix}, \eqref{eq:ch3.power-invariance-basis-change},
\eqref{eq:ch3.gauge-screen-OAM-claim}, \eqref{eq:ch3.reportable-eigenfield-observable-set} &
interpretation \\
Worked verification suite &
§3.10 &
Eqs.~\eqref{eq:ch3.wp3101-power-partition}--\eqref{eq:ch3.wp3104-coupled-power} &
FIG+PROB audit \\
\bottomrule
\end{tabular}
\end{table}

Every quantity reported from \(\mathcal{R}_{\mathrm{ch3}}\) must carry normalization tuple, outgoing certificate, and
accessible label set when geometry restricts \(\mathcal{L}_{\mathrm{acc}}\) \cite{harrington2001,stratton1941,hansen1988}.
Chapter~\ref{ch:03} does not re-open Chapter~\ref{ch:02} radiation proofs or Chapter~\ref{ch:01} gauge derivations
\cite{belinfante1940,stratton1941}.

\paragraph{Representation bridge and Hilbert closure (§3.0--§3.1).}
Sections~\ref{sec:ch30}--\ref{sec:ch31} established \(\mathcal{F}_{\mathrm{rad}}\) as the finite-energy outgoing carrier and fixed the radiative
inner product Eq.~\eqref{eq:ch3.radiative-inner-product-def} that all later projections presume \cite{kato1995,stratton1941,harrington2001}.
Downstream chapters cite Eq.~\eqref{eq:ch3.F-rad-hilbert-closure} as the domain certificate: fields live in a Hilbert space only after reactive
content is removed by \(\Pi_{\mathrm{rad}}\) \cite{harrington2001,jackson1999}. The inheritance screen Eq.~\eqref{eq:ch3.inheritance-admissibility}
requires \(\mathbf{v}_\omega\in\mathcal{V}_{\mathrm{qual}}\subset\mathcal{F}_{\mathrm{rad}}\), so realizability work in Chapter~\ref{ch:04} never
reopens the domain construction \cite{stratton1941,hansen1988}.

\paragraph{Angular momentum and SO(3) labels (§3.2).}
Section~\ref{sec:ch32} defined \(\mathbf{J}\) on Maxwell fields and recorded Wigner coupling laws used whenever rotated laboratories mix
\(m\) within fixed \(\ell\) \cite{jackson1999,hansen1988,belinfante1940}. Eq.~\eqref{eq:ch3.J-total-maxwell-def} is the operator anchor;
Subsection~\ref{sec:ch324} supplies the executable Wigner algebra \cite{harrington2001,stratton1941}. Integer labels in later sections inherit
this symmetry structure; fractional effective indices remain outside \(\mathcal{R}_{\mathrm{ch3}}\) \cite{hansen1988,harrington2001}.

\paragraph{Vector Helmholtz and Debye construction (§3.3).}
Section~\ref{sec:ch33} separated the vector Helmholtz equation, spherical separation, and Debye potential route to \(\E,\H\)
\cite{stratton1941,tai1994,hansen1988}. Eqs.~\eqref{eq:ch3.debye-E-representation}--\eqref{eq:ch3.debye-separated-expansion} are construction
coordinates; export coordinates remain VSH coefficients Eq.~\eqref{eq:ch3.normalized-coefficient-resolution} screened by
Eq.~\eqref{eq:ch3.debye-vsh-promotion-screen} \cite{belinfante1940,harrington2001}. Chapter~\ref{ch:04} may use Debye models for sources but cites
Section~\ref{sec:ch33} rather than reproving separation \cite{tai1994,stratton1941}.

\paragraph{VSH eigenfields and harmonic coordinates (§3.4--§3.5).}
Sections~\ref{sec:ch34}--\ref{sec:ch35} are the harmonic core: Theorems~\ref{thm:ch3.vsh-orthogonality}, \ref{thm:ch3.vsh-completeness}, and
\ref{thm:ch3.vsh-J-eigenfields} with flux normalization Eqs.~\eqref{eq:ch3.VSH-unit-flux-normalization}--\eqref{eq:ch3.mode-power-from-coefficient}
\cite{hansen1988,stratton1941,poynting1884}. Mode counting Eq.~\eqref{eq:ch3.cumulative-disc-DOS} and continuous Parseval
Eq.~\eqref{eq:ch3.continuous-parseval-flux-fixed} are dual ledgers for discrete and plane-wave charts \cite{jackson1999,harrington2001}.
Subsections~\ref{sec:ch3101}--\ref{sec:ch3103} arithmetic-check this bundle on concrete watt partitions \cite{hansen1988,stratton1941}.

\paragraph{Completeness, boundaries, and spectra (§3.6--§3.9).}
Sections~\ref{sec:ch36}--\ref{sec:ch39} extend the core to truncation remainders, boundary mixing, angular spectra, basis transformations, and
gauge/inference screens \cite{harrington2001,stratton1941,belinfante1940,allen1992}. Eqs.~\eqref{eq:ch3.boundary-scattering-mixing-matrix},
\eqref{eq:ch3.basis-change-matrix}, \eqref{eq:ch3.gauge-qualified-eigenfield-report}, and \eqref{eq:ch3.reportable-eigenfield-observable-set}
close the interpretation layer that Chapter~\ref{ch:04} must inherit intact \cite{hansen1988,devaney2004}. Section~\ref{sec:ch310} supplies
FIG+PROB verification Eqs.~\eqref{eq:ch3.wp3101-power-partition}--\eqref{eq:ch3.wp3104-coupled-power} as auditable templates
\cite{harrington2001,stratton1941}.

\paragraph{Worked inheritance check (admissibility screen).}
A Chapter~\ref{ch:04} manuscript satisfies Eq.~\eqref{eq:ch3.inheritance-admissibility} only when \(\mathcal{D}_{\mathrm{ch2}}=1\),
\(\mathbf{v}_\omega\in\mathcal{V}_{\mathrm{qual}}\), and \(\mathcal{N}_{\mathrm{mtr}}=1\) are documented on the first page of coefficient data
\cite{stratton1941,harrington2001,hansen1988}. Table~\ref{tab:ch3.results-ledger} is the positive list of bundles that may be cited without proof;
any equation outside that list re-opened in Chapter~\ref{ch:04} violates Eq.~\eqref{eq:ch3.bridge-inheritance-map}
\cite{harrington2001,devaney2004}. Reviewers should treat missing normalization tuples as hard stops before realizability claims are evaluated
\cite{poynting1884,harrington2001}.

\paragraph{FIG+PROB verification as inheritance evidence.}
Section~\ref{sec:ch310} is part of \(\mathcal{R}_{\mathrm{ch3}}\): Eqs.~\eqref{eq:ch3.wp3101-power-partition}--\eqref{eq:ch3.wp3104-coupled-power}
are auditable instances of the abstract theorems in Sections~\ref{sec:ch34}--\ref{sec:ch36} \cite{hansen1988,stratton1941,harrington2001}.
Downstream chapters should cite those worked equations when explaining how watt partitions, orthogonality, DOS counts, and boundary Gram matrices
appear in laboratory pipelines \cite{harrington2001,devaney2004}. The inheritance subsection does not add new mathematics; it certifies that the
chapter's exports are both proved and arithmetic-checked \cite{stratton1941,hansen1988}.

\paragraph{Normalization tuple as mandatory metadata.}
Every row of Table~\ref{tab:ch3.results-ledger} presumes Section~\ref{sec:ch35} meters: flux normalization, outgoing certificate, and declared
\(\ell_{\max}\) or continuous measure \cite{poynting1884,harrington2001,hansen1988}. \(\mathcal{N}_{\mathrm{mtr}}=1\) in
Eq.~\eqref{eq:ch3.inheritance-admissibility} is not bureaucratic; it is the statement that watt sums in Chapter~\ref{ch:04} will use the same
normalization as Eq.~\eqref{eq:ch3.mode-power-from-coefficient} \cite{stratton1941,harrington2001}. Changing meters between chapters breaks
Eq.~\eqref{eq:ch3.bridge-inheritance-map} even if intermediate algebra appears locally consistent \cite{devaney2004,hansen1988}.

\paragraph{Gauge and inference rows in the ledger.}
The §3.8--§3.9 bundle in Table~\ref{tab:ch3.results-ledger} exports basis-change matrices, power invariance, gauge screens, and reportable observable sets
\cite{belinfante1940,harrington2001,allen1992}. Chapter~\ref{ch:04} realizability arguments must not strip those screens when promoting coefficients to
source claims \cite{stratton1941,harrington2001}. Subsection~\ref{sec:ch394} positive list Eq.~\eqref{eq:ch3.reportable-eigenfield-observable-set}
is the interpretation capstone cited by inheritance \cite{hansen1988,belinfante1940}.

\paragraph{Section-by-section anchor map for reviewers.}
Subsection~\ref{sec:ch3111} closes the positive inheritance list with explicit anchors so Chapter~\ref{ch:04} citations are auditable without opening proofs
\cite{stratton1941,harrington2001}. Domain closure uses Eq.~\eqref{eq:ch3.F-rad-hilbert-closure}; projections use Eq.~\eqref{eq:ch3.radiative-inner-product-def}; harmonic construction uses
Eqs.~\eqref{eq:ch3.eigenfield-VSH-def}--\eqref{eq:ch3.normalized-coefficient-resolution}; watt meters use Eqs.~\eqref{eq:ch3.VSH-unit-flux-normalization}--\eqref{eq:ch3.mode-power-from-coefficient}; symmetry breaking uses
Eq.~\eqref{eq:ch3.boundary-scattering-mixing-matrix}; chart transport uses Eq.~\eqref{eq:ch3.basis-change-matrix}; interpretation uses Eq.~\eqref{eq:ch3.gauge-qualified-eigenfield-report}
\cite{hansen1988,poynting1884,belinfante1940,jackson1999}. Any Chapter~\ref{ch:04} manuscript that replaces these anchors with prose summaries fails Eq.~\eqref{eq:ch3.inheritance-admissibility}
\cite{harrington2001,devaney2004}.

\paragraph{Arithmetic templates bound to theorems.}
Subsection~\ref{sec:ch3101} certifies Eq.~\eqref{eq:ch3.wp3101-power-partition} against Theorem~\ref{thm:ch3.vsh-J-eigenfields}; Subsection~\ref{sec:ch3102} certifies Eq.~\eqref{eq:ch3.wp3102-orthogonality-integral} against Theorem~\ref{thm:ch3.vsh-orthogonality}; Subsection~\ref{sec:ch3103} certifies Eq.~\eqref{eq:ch3.wp3103-parseval-discrete} against Section~\ref{sec:ch35} meters; Subsection~\ref{sec:ch3104} certifies Eq.~\eqref{eq:ch3.wp3104-coupled-power} against Subsection~\ref{sec:ch363} mixing taxonomy
\cite{hansen1988,stratton1941,harrington2001,kato1995}. Inheritance therefore carries both \emph{proof ownership} and \emph{numeric acceptance bands} into Chapter~\ref{ch:04}
\cite{devaney2004,harrington2001}.

\begin{figure}[t]
\centering
\ChOneFigBlock{%
\begin{tikzpicture}[ch1fig, node distance=7mm and 8mm]
  \node[ch1box] (thm) {Chapter~\ref{ch:03} theorems};
  \node[ch1box, right=9mm of thm] (wp) {Section~\ref{sec:ch310} FIG+PROB};
  \node[ch1box, right=9mm of wp] (ch4) {Chapter~\ref{ch:04} cite};
  \node[ch1box, below=9mm of wp] (adm) {Eq.~\eqref{eq:ch3.inheritance-admissibility}};
  \draw[ch1flow] (thm) -- (wp) -- (ch4);
  \draw[ch1flow] (adm) -- (wp);
\end{tikzpicture}%
}
\caption{Subsection~\ref{sec:ch3111} inheritance: theorems, worked audits, and forward citation form a single admissible bundle.}
\label{fig:ch3.3111-inheritance-chain}
\end{figure}

\paragraph{Reviewer-facing admissibility sentences.}
Subsection~\ref{sec:ch3111} supplies three admissible citation templates for Chapter~\ref{ch:04} authors: cite Eq.~\eqref{eq:ch3.normalized-coefficient-resolution} for harmonic coordinates; cite Eq.~\eqref{eq:ch3.mode-power-from-coefficient} for watt closure; cite Eq.~\eqref{eq:ch3.inheritance-admissibility} for domain qualification
\cite{stratton1941,harrington2001,hansen1988,poynting1884}. Any sentence that re-derives those objects inline should be rewritten as a forward reference; any sentence that uses coefficients without \(\mathcal{N}_{\mathrm{mtr}}=1\) should be rejected at review
\cite{devaney2004,harrington2001,belinfante1940}. Table~\ref{tab:ch3.results-ledger} is the authoritative index of which template applies to which result bundle
\cite{hansen1988,stratton1941}.

\paragraph{Numeric templates as inheritance evidence (recap).}
Eqs.~\eqref{eq:ch3.wp3101-power-partition}, \eqref{eq:ch3.wp3102-acceptance-protocol}, \eqref{eq:ch3.wp3103-parseval-discrete}, and \eqref{eq:ch3.wp3104-coupled-power} are not optional illustrations; they are the numeric face of \(\mathcal{R}_{\mathrm{ch3}}\) cited by inheritance
\cite{harrington2001,stratton1941,kato1995}. Chapter~\ref{ch:04} may reproduce those arithmetic checks on new geometries but must not relabel them as new theorems
\cite{devaney2004,hansen1988,jackson1999}.

\paragraph{Chapter~\ref{ch:03} export certificate (closing statement).}
Subsection~\ref{sec:ch3111} certifies \(\mathcal{R}_{\mathrm{ch3}}\) as the set of results whose proofs live in Sections~\ref{sec:ch30}--\ref{sec:ch310} and whose numeric acceptance is demonstrated by Section~\ref{sec:ch310} FIG+PROB equations
\cite{stratton1941,hansen1988,harrington2001,kato1995}. Chapter~\ref{ch:04} may treat that certificate as frozen: cite, do not reopen, do not relabel as realizability
\cite{devaney2004,belinfante1940,poynting1884}. Eq.~\eqref{eq:ch3.inheritance-admissibility} is the formal gate; Table~\ref{tab:ch3.results-ledger} is the human-readable index
\cite{harrington2001,stratton1941}.

\paragraph{Inheritance admissibility walk-through.}
A compliant Chapter~\ref{ch:04} opening section should read, in order: \(\mathcal{D}_{\mathrm{ch2}}=1\) from Chapter~\ref{ch:02}; \(\mathbf{v}_\omega\in\mathcal{V}_{\mathrm{qual}}\) from Section~\ref{sec:ch31}; \(\mathcal{N}_{\mathrm{mtr}}=1\) from Section~\ref{sec:ch35}; harmonic expansion Eq.~\eqref{eq:ch3.normalized-coefficient-resolution} from Section~\ref{sec:ch34}
\cite{stratton1941,harrington2001,hansen1988,poynting1884}. Subsection~\ref{sec:ch3111} treats that ordering as part of the inheritance proof obligation: skipping a rung invalidates downstream realizability claims even if local algebra is correct
\cite{devaney2004,belinfante1940,jackson1999}. Figure~\ref{fig:ch3.3111-inheritance-chain} is the schematic of that obligation
\cite{harrington2001,stratton1941}.

Subsection~\ref{sec:ch3112} catalogs deferred realizability, inverse, and superresolution topics for
Chapter~\ref{ch:04} and later volumes \cite{devaney2004,harrington2001}.


\subsection{Concepts deferred to Chapter 4 and Volume II}
\label{sec:ch3112}
% TOC tag: [DEFER]

Chapter~\ref{ch:03} deliberately stops before source realizability proofs, equivalence aperture constructions, and
singular-value realizability chains that constitute Chapter~\ref{ch:04} \cite{devaney2004,harrington2001,stratton1941}.
\emph{Deferred} means a topic is named, ownership is assigned forward, and no proof is attempted in the present chapter
\cite{stratton1941,harrington2001}. Table~\ref{tab:ch3.deferred-ledger} lists the principal deferrals; Volume~II items are
named only as pointers \cite{hansen1988,devaney2004}.

\begin{table}[t]
\centering
\caption{Chapter~\ref{ch:03} deferred-concept ledger (ownership forward).}
\label{tab:ch3.deferred-ledger}
\small
\begin{tabular}{@{}p{0.28\linewidth}p{0.18\linewidth}p{0.46\linewidth}@{}}
\toprule
Concept & Owner & Reason deferred from Ch.~3 \\
\midrule
Realizable source spaces & Ch.~4 §4.1 &
requires impressed-current passivity/support proofs \\
Love/Huygens equivalence & Ch.~4 §4.2 &
aperture realizability, not representation \\
Operator SVD chains & Ch.~4 §4.3 &
rank of \(\mathcal{R}_\omega\) on finite support \\
Fourier--angular bandlimits & Ch.~4 §4.4 &
inherits §3.8.1; needs source geometry \\
Eigenfield coupling matrices & Ch.~4 §4.5 &
imports §3.4--§3.6 truncation taxonomy \\
Truncation scaling laws & Ch.~4 §4.6--§4.7 &
imports §3.6.2 remainder bounds \\
Full inverse uniqueness & Vol.~II &
Fredholm inversion beyond Volume~I scope \\
\bottomrule
\end{tabular}
\end{table}

Truncation remainders Eq.~\eqref{eq:ch3.truncation-remainder-energy} and boundary mixing
Eq.~\eqref{eq:ch3.boundary-scattering-mixing-matrix} supply \emph{taxonomy} only; Chapter~\ref{ch:04} owns quantitative
realizability and measurement inversion \cite{harrington2001,devaney2004,stratton1941}. Paraxial reductions
Eq.~\eqref{eq:ch3.paraxial-reduction-condition} remain approximation certificates; full vector inversion is deferred
\cite{allen1992,harrington2001}.

\paragraph{Realizable source spaces (Chapter~\ref{ch:04}, §4.1).}
Finite impressed-current supports, passivity, and realizability norms require source-domain proofs that Chapter~\ref{ch:03} deliberately omits
\cite{stratton1941,devaney2004,harrington2001}. Chapter~\ref{ch:03} supplies harmonic coordinates on \(\mathcal{F}_{\mathrm{rad}}\); Chapter~\ref{ch:04}
will ask which coefficient vectors arise from physically realizable \(\mathbf{J}^{\mathrm{imp}}\) \cite{harrington2001,stratton1941}. The deferral
is ownership, not ignorance: Eq.~\eqref{eq:ch2.near-to-far-operator-composition} already composed the export operator, but rank and accessibility on
finite supports belong to the realizability chapter \cite{devaney2004,harrington2001}.

\paragraph{Love/Huygens equivalence (Chapter~\ref{ch:04}, §4.2).}
Equivalence apertures replace sources by equivalent surface currents while preserving exterior fields \cite{stratton1941,harrington2001}. Chapter~\ref{ch:03}
provides the VSH charts in which equivalent currents would be expanded, yet the uniqueness and realizability of those currents on finite apertures
is deferred \cite{harrington2001,devaney2004}. Boundary mixing Eq.~\eqref{eq:ch3.boundary-scattering-mixing-matrix} supplies symmetry-breaking taxonomy
only \cite{jackson1999,harrington2001}.

\paragraph{Operator SVD and coupling chains (Chapter~\ref{ch:04}, §4.3--§4.5).}
Singular-value chains on \(\mathcal{R}_\omega\) restricted to finite supports determine which angular spectra are measurable versus invisible
\cite{devaney2004,harrington2001}. Chapter~\ref{ch:03} normalizes and labels modes; Chapter~\ref{ch:04} will compute rank-deficiency and coupling
matrices between eigenfield sectors induced by source geometry \cite{harrington2001,stratton1941}. Subsection~\ref{sec:ch3104} is the schematic
\(2\times2\) instance; full arrays defer to §4.5 \cite{harrington2001,jackson1999}.

\paragraph{Bandlimits and truncation scaling (Chapter~\ref{ch:04}, §4.4--§4.7).}
Fourier--angular bandlimits inherit Subsection~\ref{sec:ch381} plane-wave charts but require source-supported geometry to state which \(k_{t}\) or
\(\ell\) sectors are excitable \cite{jackson1999,harrington2001,devaney2004}. Truncation scaling laws connect \(\ell_{\max}\) to measurement noise
and superresolution claims; Eq.~\eqref{eq:ch3.truncation-remainder-energy} supplies the remainder definition, not the scaling proof
\cite{harrington2001,kato1995}. Volume~II owns full Fredholm inversion and uniqueness beyond Volume~I scope \cite{devaney2004,harrington2001}.

\paragraph{Why realizability is deferred (not forgotten).}
Chapter~\ref{ch:03} ends at harmonic coordinates on \(\mathcal{F}_{\mathrm{rad}}\) because realizability is a \emph{source-domain} question:
which coefficient vectors arise from passive, finite-support impressed currents, and which are mere mathematical fits to exterior patterns
\cite{stratton1941,devaney2004,harrington2001}? Eq.~\eqref{eq:ch2.near-to-far-operator-composition} already maps sources to exports, but rank,
null spaces, and equivalence apertures require Chapter~\ref{ch:04} proofs \cite{harrington2001,stratton1941}. Table~\ref{tab:ch3.deferred-ledger}
assigns ownership so Volume~I readers know exactly where each proof lives \cite{devaney2004,hansen1988}.

\paragraph{Truncation and inversion without recursion.}
Truncation remainder Eq.~\eqref{eq:ch3.truncation-remainder-energy} quantifies band-limit error on \(\mathcal{F}_{\mathrm{rad}}\); it does not
prove scaling laws for inverse problems \cite{kato1995,harrington2001,devaney2004}. Boundary mixing Eq.~\eqref{eq:ch3.boundary-scattering-mixing-matrix}
classifies symmetry breaking; it does not construct Love equivalents \cite{jackson1999,stratton1941}. Paraxial reduction
Eq.~\eqref{eq:ch3.paraxial-reduction-condition} certifies approximations; it does not replace vector inversion on \(\Gamma_R\)
\cite{allen1992,harrington2001,belinfante1940}. These distinctions prevent Chapter~\ref{ch:04} from smuggling deferred proofs back into
Chapter~\ref{ch:03} citations \cite{stratton1941,harrington2001}.

\paragraph{Volume~II pointer (Fredholm and uniqueness).}
Full inverse uniqueness and Fredholm theory for angular spectra exceed Volume~I scope \cite{devaney2004,harrington2001}. Chapter~\ref{ch:03} supplies
normalized charts and gauge screens so Volume~II can state inversion theorems on a fixed harmonic infrastructure \cite{hansen1988,stratton1941}.
Readers pursuing superresolution claims should treat Table~\ref{tab:ch3.deferred-ledger} as the ownership map before citing Chapter~\ref{ch:03}
equations as if they were inversion theorems \cite{harrington2001,devaney2004}.

\paragraph{Deferred realizability in source language.}
Realizable source spaces (Chapter~\ref{ch:04} §4.1) ask which \(\mathbf{J}^{\mathrm{imp}}\) with finite support and passivity constraints generate a
prescribed \(\mathbf{v}_\omega\) on \(\mathcal{R}_\omega\) \cite{stratton1941,devaney2004,harrington2001}. Chapter~\ref{ch:03} supplies the target
field expansion Eq.~\eqref{eq:ch3.normalized-coefficient-resolution} but not the inverse map from coefficients to currents \cite{harrington2001,stratton1941}.
Love/Huygens equivalence (§4.2) replaces sources by surface currents on apertures while preserving exterior fields \cite{stratton1941,harrington2001};
Chapter~\ref{ch:03} provides the harmonic basis for those surface expansions yet defers uniqueness and realizability proofs \cite{devaney2004,hansen1988}.

\paragraph{Deferred operator chains and coupling matrices.}
Operator SVD chains (§4.3) determine measurable subspaces of \(\mathcal{R}_\omega\) on finite supports \cite{devaney2004,harrington2001}. Eigenfield coupling
matrices (§4.5) import truncation taxonomy from Sections~\ref{sec:ch34}--\ref{sec:ch36} but compute source-induced ranks not present in Chapter~\ref{ch:03}
\cite{harrington2001,stratton1941}. Subsection~\ref{sec:ch3104} schematically diagonalizes a \(2\times2\) Gram matrix; array-scale coupling defers to
Chapter~\ref{ch:04} \cite{jackson1999,harrington2001,devaney2004}.

\paragraph{Deferred bandlimits and scaling laws.}
Fourier--angular bandlimits (§4.4) inherit Subsection~\ref{sec:ch381} plane-wave charts but require source geometry to state excitable \(k_{t}\) and \(\ell\)
sectors \cite{jackson1999,harrington2001,devaney2004}. Truncation scaling laws (§4.6--§4.7) connect \(\ell_{\max}\) to measurement noise; Eq.~\eqref{eq:ch3.truncation-remainder-energy}
defines remainders only \cite{kato1995,harrington2001}. Paraxial reductions Eq.~\eqref{eq:ch3.paraxial-reduction-condition} remain approximation certificates;
vector uplift Eq.~\eqref{eq:ch3.normalized-coefficient-resolution} is still the far-zone anchor \cite{allen1992,harrington2001,belinfante1940}.

\paragraph{How deferrals interact with inheritance.}
Deferred topics are not optional reading: they are \emph{proof ownership} boundaries \cite{stratton1941,harrington2001}. A Chapter~\ref{ch:04} argument that
re-proves VSH completeness while claiming a new realizability theorem violates Eq.~\eqref{eq:ch3.bridge-inheritance-map}
\cite{harrington2001,devaney2004}. Conversely, a realizability proof that cites Eq.~\eqref{eq:ch3.normalized-coefficient-resolution} and
Eq.~\eqref{eq:ch3.mode-power-from-coefficient} without altering them satisfies the bridge screen \cite{poynting1884,hansen1988,stratton1941}.

\paragraph{Deferred inverse problems and measurement rank.}
Full inverse uniqueness and Fredholm stability for angular spectra are Volume~II topics \cite{devaney2004,harrington2001}. Chapter~\ref{ch:03} nevertheless fixes the forward harmonic coordinates on which those inversions will be posed: normalized VSH coefficients, accessible label sets, and gauge screens
\cite{hansen1988,stratton1941,belinfante1940}. Chapter~\ref{ch:04} may bound measurement rank using operator SVD chains, but may not redefine Eq.~\eqref{eq:ch3.normalized-coefficient-resolution} while doing so
\cite{harrington2001,devaney2004}.

\paragraph{Superresolution claims and truncation ownership.}
Truncation scaling laws connect \(\ell_{\max}\) to recoverable detail; Eq.~\eqref{eq:ch3.truncation-remainder-energy} supplies the remainder norm, not the scaling exponent
\cite{kato1995,harrington2001,devaney2004}. Campaigns that claim superresolved OAM modes without declaring \(\ell_{\max}\), noise floor, and accessibility mask borrow Chapter~\ref{ch:03} vocabulary without satisfying Chapter~\ref{ch:04} realizability ownership
\cite{allen1992,harrington2001,belinfante1940}. Table~\ref{tab:ch3.deferred-ledger} assigns those proofs forward so Volume~I readers know the proof is intentionally absent here
\cite{stratton1941,devaney2004}.

\begin{table}[t]
\centering
\caption{Subsection~\ref{sec:ch3112} deferral acceptance: what Chapter~\ref{ch:03} may and may not claim.}
\label{tab:ch3.3112-deferral-acceptance}
\small
\begin{tabular}{@{}p{0.38\linewidth}p{0.52\linewidth}@{}}
\toprule
Claim type & Status in Ch.~3 \\
\midrule
VSH expansion of $\mathbf{v}_\omega$ & established (cite §3.4) \\
Source realizability of $\widehat{c}$ & deferred (Ch.~4) \\
Love equivalent currents & deferred (Ch.~4) \\
Fredholm inversion uniqueness & deferred (Vol.~II) \\
Truncation remainder norm & established (cite §3.6.2) \\
Truncation scaling exponent & deferred (Ch.~4) \\
\bottomrule
\end{tabular}
\end{table}

\paragraph{Love equivalence and aperture currents (expanded deferral).}
Love/Huygens constructions replace impressed sources by equivalent surface currents on apertures \(\mathcal{A}\) while preserving exterior fields on \(\mathcal{R}_\omega\)
\cite{stratton1941,harrington2001}. Chapter~\ref{ch:03} expands those exterior fields in Eq.~\eqref{eq:ch3.normalized-coefficient-resolution} but does not prove that a given \(\widehat{c}\) vector arises from a realizable \(\mathbf{J}_{\mathrm{eq}}\) on \(\mathcal{A}\)
\cite{devaney2004,harrington2001}. That inverse map is Chapter~\ref{ch:04} §4.2 ownership; Volume~I readers should cite boundary taxonomy Eq.~\eqref{eq:ch3.boundary-scattering-mixing-matrix} when symmetry breaking is present, not when realizability is claimed
\cite{jackson1999,stratton1941,hansen1988}.

\paragraph{Operator rank and measurement null spaces (expanded deferral).}
Singular-value chains on \(\mathcal{R}_\omega\) restricted to finite supports determine which angular harmonics are stably measurable
\cite{devaney2004,harrington2001}. Chapter~\ref{ch:03} supplies forward harmonic coordinates and gauge screens; Chapter~\ref{ch:04} §4.3 computes rank and nullity as functions of source support and measurement operator \(\mathbf{G}\) from Eq.~\eqref{eq:ch1.modal-projection-matrix}
\cite{harrington2001,stratton1941}. Subsection~\ref{sec:ch394} identifiability bound Eq.~\eqref{eq:ch3.modal-identifiability-bound} is the interpretation-layer preview of that rank discussion, not its proof
\cite{devaney2004,hansen1988,belinfante1940}.

\paragraph{Fredholm inversion and Volume~II boundary.}
Full Fredholm inversion, uniqueness, and superresolution scaling exponents exceed Volume~I scope \cite{devaney2004,harrington2001}. Chapter~\ref{ch:03} ends with normalized charts so Volume~II can assume Eq.~\eqref{eq:ch3.continuous-parseval-flux-fixed}, Eq.~\eqref{eq:ch3.gauge-qualified-eigenfield-report}, and Eq.~\eqref{eq:ch3.reportable-eigenfield-observable-set} as fixed infrastructure
\cite{hansen1988,stratton1941,belinfante1940}. Table~\ref{tab:ch3.3112-deferral-acceptance} is the quick reference for reviewers who encounter inversion claims in Chapter~\ref{ch:03} citations
\cite{harrington2001,devaney2004}.

\paragraph{Forward pointer without proof (Volume~II).}
Fredholm stability and superresolution uniqueness are named in Table~\ref{tab:ch3.deferred-ledger} so readers know the proof is intentionally absent in Volume~I
\cite{devaney2004,harrington2001}. Chapter~\ref{ch:03} nevertheless fixes the forward harmonic infrastructure those theorems will assume: Parseval meters, gauge screens, and reportable observable vocabulary
\cite{hansen1988,belinfante1940,stratton1941}. Citing Chapter~\ref{ch:03} for inversion theorems without citing Volume~II ownership is a category error Subsection~\ref{sec:ch3112} prevents by explicit deferral
\cite{harrington2001,devaney2004}.

\paragraph{Anti-confusion table (established vs deferred).}
Readers entering Chapter~\ref{ch:04} should memorize one distinction: \emph{established} objects have equation anchors in Table~\ref{tab:ch3.results-ledger}; \emph{deferred} objects have owners in Table~\ref{tab:ch3.deferred-ledger}
\cite{stratton1941,harrington2001,devaney2004}. Subsection~\ref{sec:ch3112} exists to prevent category errors---especially treating truncation remainders as realizability theorems or treating LG labels as gauge-invariant eigenfield observables
\cite{allen1992,belinfante1940,hansen1988}. Volume~II Fredholm items are named only as pointers, not as proved results
\cite{devaney2004,harrington2001}.

Subsection~\ref{sec:ch3113} states the explicit bridge map to Chapter~\ref{ch:04}
\cite{stratton1941,harrington2001}.


\subsection{Bridge to Chapter 4}
\label{sec:ch3113}
% TOC tag: [DEFER]

Chapter~\ref{ch:04} begins from the harmonic and spectral coordinates assembled in Sections~\ref{sec:ch30}--\ref{sec:ch39}:
normalized VSH expansions, continuous angular spectra, basis-change matrices, and gauge/inference screens. These exports are
imported by reference; Chapter~\ref{ch:04} may develop realizability, equivalence apertures, and singular-value source chains
without reopening Chapter~\ref{ch:03} proofs. Compatibility is encoded by
\begin{equation}
\label{eq:ch3.bridge-inheritance-map}
\mathcal{D}_{\mathrm{bridge}}^{(3\to 4)}(\mathcal{Y})=1
\iff
\left(
\mathcal{C}_{\mathrm{ch3}}(\mathcal{Y})=1\
\wedge\
\mathcal{R}_{\mathrm{rec}}=0
\right),
\end{equation}
where \(\mathcal{C}_{\mathrm{ch3}}(\mathcal{Y})=1\) means \(\mathcal{Y}\) is compatible with \(\mathcal{R}_{\mathrm{ch3}}\) and \(\mathcal{R}_{\mathrm{rec}}=0\) forbids recursive reproof of Chapter~\ref{ch:03} results inside Chapter~\ref{ch:04},
screened by Eq.~\eqref{eq:ch3.inheritance-admissibility} and Eq.~\eqref{eq:ch2.bridge-inheritance-map}
\cite{stratton1941,harrington2001}.

\paragraph{VSH basis and eigenfield export.}
Chapter~\ref{ch:04} inherits the harmonic construction of Section~\ref{sec:ch34} by reference: eigenfield definition
Eq.~\eqref{eq:ch3.eigenfield-VSH-def}, \(\mathbf{M}\)/\(\mathbf{N}\) templates Eqs.~\eqref{eq:ch3.Mlm-def}--\eqref{eq:ch3.Nlm-def},
simultaneous angular-momentum system Eq.~\eqref{eq:ch3.J-simultaneous-eigenfield-system}, and normalized resolution
Eq.~\eqref{eq:ch3.normalized-coefficient-resolution} supply the coordinate chart on \(\mathcal{F}_{\mathrm{rad}}\)
\cite{hansen1988,stratton1941,jackson1999}. Completeness Theorem~\ref{thm:ch3.vsh-completeness} is cited, not reproved; realizability
arguments in Chapter~\ref{ch:04} therefore treat \(\widehat{c}_{\ell m}^{(\sigma)}\) as given projections once outgoing selection and
flux meters are declared \cite{kato1995,harrington2001}. Mathematically, the export is a Riesz basis on the radiative Hilbert space;
physically, it is the watt-diagonal language in which finite sources will be compared to measured far-zone patterns
\cite{poynting1884,stratton1941}.

\paragraph{Normalization and Parseval meters.}
Watt closure enters Chapter~\ref{ch:04} through Section~\ref{sec:ch35}: unit-flux normalization
Eq.~\eqref{eq:ch3.VSH-unit-flux-normalization}, energy norm Eq.~\eqref{eq:ch3.VSH-unit-energy-norm}, and mode power law
Eq.~\eqref{eq:ch3.mode-power-from-coefficient} fix how impressed-current models map to reportable \(\mathcal{P}_{\mathrm{rad}}\)
\cite{poynting1884,harrington2001,stratton1941}. Continuous Parseval Eq.~\eqref{eq:ch3.continuous-parseval-flux-fixed} and discrete DOS
Eq.~\eqref{eq:ch3.cumulative-disc-DOS} remain the dual ledgers for plane-wave and VSH charts; Chapter~\ref{ch:04} may compare source
spectra to those meters but does not redefine them \cite{hansen1988,jackson1999}. Any singular-value chain on finite support presumes
the same normalization tuple archived with Subsection~\ref{sec:ch3103} acceptance thresholds \cite{harrington2001,devaney2004}.

\paragraph{Truncation, boundaries, and accessibility.}
Truncation remainder Eq.~\eqref{eq:ch3.truncation-remainder-energy}, boundary mixing matrix
Eq.~\eqref{eq:ch3.boundary-scattering-mixing-matrix}, and accessible watt budget Eq.~\eqref{eq:ch3.accessible-watt-budget} form the
symmetry-breaking taxonomy imported from Section~\ref{sec:ch36} \cite{harrington2001,stratton1941}. Chapter~\ref{ch:04} uses these objects to
separate mathematical channel count from probe-visible count without reopening image-theory derivations \cite{jackson1999,hansen1988}.
Subsection~\ref{sec:ch3104} is the arithmetic template: \(\mathbf{S}_{b}\), Gram matrix \(\mathbf{G}\), and coupled powers
Eq.~\eqref{eq:ch3.wp3104-coupled-power} must be cited when realizability claims involve ground planes or mutual coupling
\cite{harrington2001,stratton1941}.

\paragraph{Angular-spectrum and basis-change bridges.}
Discrete--continuous bridge Eq.~\eqref{eq:ch3.discrete-continuous-spectrum-bridge}, plane-wave resolution
Eq.~\eqref{eq:ch3.plane-wave-basis-resolution}, and basis-change matrix Eq.~\eqref{eq:ch3.basis-change-matrix} connect lateral FFT,
\(S^{2}\), and VSH charts \cite{jackson1999,hansen1988,stratton1941}. Chapter~\ref{ch:04} Fourier--angular bandlimit arguments inherit
Subsection~\ref{sec:ch381} Jacobian Eq.~\eqref{eq:ch3.plane-wave-direction-jacobian} and power invariance
Eq.~\eqref{eq:ch3.power-invariance-basis-change} as fixed transformation rules \cite{harrington2001,devaney2004}. Source geometry
determines which chart is primary, not which measure is correct \cite{stratton1941,jackson1999}.

\paragraph{Gauge and inference screens.}
Gauge-qualified eigenfield report Eq.~\eqref{eq:ch3.gauge-qualified-eigenfield-report}, paraxial validity
Eq.~\eqref{eq:ch3.paraxial-reduction-condition}, and reportable observable set Eq.~\eqref{eq:ch3.reportable-eigenfield-observable-set}
close the interpretation layer of Section~\ref{sec:ch39} \cite{belinfante1940,harrington2001,allen1992}. Chapter~\ref{ch:04} must not promote
Debye or LG labels to realizability claims without those screens; Subsection~\ref{sec:ch391} Debye-to-VSH promotion
Eq.~\eqref{eq:ch3.debye-vsh-promotion-screen} is the operational gate \cite{tai1994,hansen1988,stratton1941}.

\paragraph{Near-to-far composition without altering exports.}
Chapter~\ref{ch:04} couples impressed currents to \(\mathcal{R}_\omega\) through Eq.~\eqref{eq:ch2.near-to-far-operator-composition} \cite{stratton1941,devaney2004}.
That composition presumes the Chapter~\ref{ch:03} harmonic coordinates are fixed: changing normalization conventions mid-derivation would break
Eq.~\eqref{eq:ch3.mode-power-from-coefficient} and invalidate Subsection~\ref{sec:ch3103} Parseval audits \cite{poynting1884,harrington2001,hansen1988}.
Realizability proofs therefore begin by citing \(\mathcal{N}_{\mathrm{mtr}}=1\) and the VSH resolution Eq.~\eqref{eq:ch3.normalized-coefficient-resolution},
then ask which source supports generate a target coefficient vector \cite{harrington2001,stratton1941,devaney2004}.

\paragraph{Regime gates carried forward.}
Chapter~\ref{ch:02} regime gates and Chapter~\ref{ch:01} gauge screens remain mandatory at the Chapter~\ref{ch:04} entry point \cite{belinfante1940,harrington2001,stratton1941}.
A realizability claim that skips Eq.~\eqref{eq:ch3.inheritance-admissibility} fails Eq.~\eqref{eq:ch3.bridge-inheritance-map} even if the source construction is
locally correct \cite{harrington2001,devaney2004}. Figure~\ref{fig:ch3.3113-bridge} summarizes the one-way import: coordinates from Chapter~\ref{ch:03},
operators from Chapter~\ref{ch:02}, conservation certificates from Chapter~\ref{ch:01} \cite{stratton1941,hansen1988,belinfante1940}.

\paragraph{Anti-recursion examples (what Chapter~\ref{ch:04} must not do).}
Chapter~\ref{ch:04} must not re-prove Theorem~\ref{thm:ch3.vsh-completeness}, re-derive Eq.~\eqref{eq:ch3.plane-wave-basis-resolution}, or re-open
Belinfante symmetrization \cite{kato1995,belinfante1940,hansen1988}. It may \emph{use} those results to bound realizability error, construct equivalence
apertures, or truncate source expansions \cite{devaney2004,harrington2001,stratton1941}. Subsection~\ref{sec:ch3104} is the template for boundary
problems: cite \(\mathbf{S}_{b}\) and \(\mathbf{G}\), do not re-derive image theory \cite{jackson1999,harrington2001}.

\paragraph{Chapter~\ref{ch:04} entry checklist (executable).}
Subsection~\ref{sec:ch3113} closes Volume~I representation with an entry checklist encoded by
\begin{equation}
\label{eq:ch3.ch4-entry-checklist}
\mathcal{E}_{4}=1
\iff
\left(
\mathcal{D}_{\mathrm{bridge}}^{(3\to4)}=1\
\wedge\
\mathcal{N}_{\mathrm{mtr}}=1\
\wedge\
\mathcal{L}_{\mathrm{decl}}=1\
\wedge\
\mathcal{R}_{\mathrm{rec}}=0
\right),
\end{equation}
with \(\mathcal{L}_{\mathrm{decl}}=1\) denoting a declared accessible label set and \(\mathcal{R}_{\mathrm{rec}}=0\) forbidding reproof of \(\mathcal{R}_{\mathrm{ch3}}\), screened by Eq.~\eqref{eq:ch3.bridge-inheritance-map} and Eq.~\eqref{eq:ch3.inheritance-admissibility}
\cite{stratton1941,harrington2001,devaney2004}. Equation~\eqref{eq:ch3.ch4-entry-checklist} is the operational distillation of Figure~\ref{fig:ch3.3113-bridge}: Chapter~\ref{ch:04} begins only when harmonic coordinates, meters, and accessibility are frozen imports
\cite{hansen1988,poynting1884,belinfante1940}.

\paragraph{Export bundle index for quick citation.}
The five export bundles named in Subsection~\ref{sec:ch3113} map to Section~\ref{sec:ch34} eigenfields, Section~\ref{sec:ch35} meters, Section~\ref{sec:ch36} truncation/boundary taxonomy, Sections~\ref{sec:ch37}--\ref{sec:ch38} spectrum bridges, and Section~\ref{sec:ch39} interpretation screens
\cite{hansen1988,jackson1999,belinfante1940,harrington2001}. Chapter~\ref{ch:04} authors should cite the bundle anchor first, then state the realizability construction, never the reverse
\cite{stratton1941,devaney2004}. Subsection~\ref{sec:ch310} FIG+PROB equations remain the numeric templates attached to bundles three and four
\cite{harrington2001,hansen1988}.

Chapter~\ref{ch:04} couples these coordinates to
\(\mathcal{R}_\omega\) on finite supports without altering export content
Eq.~\eqref{eq:ch2.near-to-far-operator-composition} \cite{stratton1941,devaney2004}. Regime gates from
Chapter~\ref{ch:02} and gauge screens from Chapter~\ref{ch:01} remain mandatory before any modal label is promoted to
realizability or transport claim \cite{belinfante1940,harrington2001,stratton1941}.

\begin{figure}[t]
\centering
\ChOneFigBlock{%
\begin{tikzpicture}[ch1fig, node distance=7mm and 9mm]
  \node[ch1box] (ch2) {Ch.~\ref{ch:02} $\mathcal{R}_\omega$};
  \node[ch1box, right=10mm of ch2] (ch3) {Ch.~\ref{ch:03} coords};
  \node[ch1box, right=10mm of ch3] (ch4) {Ch.~\ref{ch:04} realizability};
  \node[ch1box, below=10mm of ch3] (eq) {Eq.~\eqref{eq:ch3.bridge-inheritance-map}};
  \draw[ch1flow] (ch2) -- (ch3) -- (ch4);
  \draw[ch1flow] (eq) -- (ch3);
\end{tikzpicture}%
}
\caption{Subsection~\ref{sec:ch3113} bridge: Chapter~\ref{ch:03} supplies harmonic coordinates; Chapter~\ref{ch:04} supplies source realizability without recursive proofs.}
\label{fig:ch3.3113-bridge}
\end{figure}

\paragraph{First-use citation map for Chapter~\ref{ch:04} authors.}
On first appearance of \(\widehat{c}_{\ell m}^{(\sigma)}\), cite Eq.~\eqref{eq:ch3.normalized-coefficient-resolution}; on first watt sum, cite Eq.~\eqref{eq:ch3.mode-power-from-coefficient}; on first truncation claim, cite Eq.~\eqref{eq:ch3.truncation-remainder-energy}; on first boundary geometry, cite Eq.~\eqref{eq:ch3.boundary-scattering-mixing-matrix}; on first OAM headline, cite Eq.~\eqref{eq:ch3.reportable-eigenfield-observable-set}
\cite{harrington2001,stratton1941,hansen1988,belinfante1940,devaney2004}. Subsection~\ref{sec:ch3113} encodes that map so Chapter~\ref{ch:04} prose cannot drift into recursive re-derivation
\cite{harrington2001,jackson1999,kato1995}.

\paragraph{Compatibility tests before Chapter~\ref{ch:04} entry.}
Any Chapter~\ref{ch:04} development must satisfy Eq.~\eqref{eq:ch3.bridge-inheritance-map} on its first use of Chapter~\ref{ch:03} symbols:
coefficient expansions cite Eq.~\eqref{eq:ch3.normalized-coefficient-resolution}; watt sums cite Eq.~\eqref{eq:ch3.mode-power-from-coefficient};
truncation claims cite Eq.~\eqref{eq:ch3.truncation-remainder-energy}; boundary problems cite Eq.~\eqref{eq:ch3.boundary-scattering-mixing-matrix};
interpretation claims cite Eq.~\eqref{eq:ch3.reportable-eigenfield-observable-set} \cite{harrington2001,stratton1941,belinfante1940,hansen1988}.
Failure at any row is a recursion violation, not a new theorem \cite{stratton1941,harrington2001}. Figure~\ref{fig:ch3.3113-bridge} is the schematic
of that one-way import: Chapter~\ref{ch:02} supplies \(\mathcal{R}_\omega\); Chapter~\ref{ch:03} supplies harmonic coordinates; Chapter~\ref{ch:04}
supplies realizability on finite supports \cite{devaney2004,harrington2001,jackson1999}.

\paragraph{Volume~I representation closure before realizability.}
Subsection~\ref{sec:ch3113} closes Chapter~\ref{ch:03} with a single scientific claim: outgoing radiative structure on \(\mathcal{F}_{\mathrm{rad}}\) is \emph{labeled, normalized, transformed, screened, and arithmetic-checked} in harmonic coordinates ready for source-side analysis
\cite{hansen1988,stratton1941,harrington2001,poynting1884,belinfante1940}. Chapter~\ref{ch:04} inherits that claim through Eqs.~\eqref{eq:ch3.bridge-inheritance-map}, \eqref{eq:ch3.ch4-entry-checklist}, and \eqref{eq:ch3.inheritance-admissibility} without altering export content
\cite{devaney2004,harrington2001,jackson1999}. Readers entering realizability theory should carry Table~\ref{tab:ch3.results-ledger} and Table~\ref{tab:ch3.deferred-ledger} as complementary maps: what is established here versus what is intentionally owned forward
\cite{stratton1941,harrington2001}.

\paragraph{Section terminal summary.}
Chapter~\ref{ch:03} now stands as the representation and harmonic-coordinate basis for Volume~I: outgoing fields are expanded in
VSH and allied charts, normalized for watts, screened for gauge and inference, and verified in Section~\ref{sec:ch310}.
Chapter~\ref{ch:04} inherits this base by reference and develops source realizability and spectral truncation without
re-deriving Chapter~\ref{ch:03}, Chapter~\ref{ch:02}, or Chapter~\ref{ch:01} conservation proofs
\cite{stratton1941,harrington2001,hansen1988,devaney2004,belinfante1940}.

\paragraph{Handoff sentence to Chapter~\ref{ch:04}.}
Volume~I representation theory ends here; source realizability begins in Chapter~\ref{ch:04} under Eq.~\eqref{eq:ch3.ch4-entry-checklist} with frozen imports from Eqs.~\eqref{eq:ch3.bridge-inheritance-map} and~\eqref{eq:ch3.inheritance-admissibility}
\cite{harrington2001,devaney2004,stratton1941}. Every Chapter~\ref{ch:04} section header should be traceable to a row of Table~\ref{tab:ch3.results-ledger} or Table~\ref{tab:ch3.deferred-ledger}; headers that are traceable to neither require revision before drafting continues
\cite{hansen1988,belinfante1940,jackson1999,kato1995}.

\bibliographystyle{IEEEtran}
\bibliography{references}

